% !TEX program = xelatex
\documentclass[UTF8,a4paper,zihao=5,openany,twocolumn,fontset=windows]{ctexbook}

% ==================== 页面边距与排版设置 ====================
% 用户要求：双栏布局、9pt字号、7.5mm超窄边距、1.1行距、主要用宋体
\usepackage{geometry}
\geometry{
    left=7.5mm,
    right=7.5mm,
    top=1.8cm,
    bottom=1.8cm,
    headheight=14pt,
    headsep=10pt,
    footskip=0.8cm,
    columnsep=16pt
}
\let\cleardoublepage\clearpage

% ==================== 字体与行距 ====================
\usepackage{fontspec}
\setmainfont{Times New Roman}
\setsansfont{Arial}
\setmonofont{Consolas}
\linespread{1.1}

% ==================== 常用数学与表格宏包 ====================
\usepackage{amsmath,amssymb,amsfonts}
\usepackage{booktabs,longtable,tabularx,multirow,array}
\usepackage{etoolbox}
\usepackage{dblfloatfix}
\usepackage{float}
\usepackage{makecell}
\usepackage{cuted}
\newcolumntype{Y}{>{\raggedright\arraybackslash}X}

% 保证 longtable 在单栏全宽显示，结束后恢复双栏
\makeatletter
\BeforeBeginEnvironment{longtable}{%
    \if@twocolumn
        \clearpage
        \onecolumn
    \fi
}
\makeatother
\AfterEndEnvironment{longtable}{\clearpage\twocolumn}

% 让目录保持双栏排版
\makeatletter
\patchcmd{\tableofcontents}{\@restonecoltrue\onecolumn}{\@restonecolfalse}{}{}
\patchcmd{\tableofcontents}{\clearpage\if@restonecol\twocolumn\fi}{}{}{}
\makeatother

% ==================== 颜色与超链接 ====================
\usepackage[table]{xcolor}
\definecolor{primaryblue}{RGB}{20, 50, 95}
\definecolor{deepcyan}{RGB}{0, 102, 128}
\definecolor{darkgray}{RGB}{60, 60, 60}
\definecolor{boxbg}{RGB}{245, 248, 253}
\definecolor{boxframe}{RGB}{180, 205, 235}
\definecolor{taskbg}{RGB}{250, 250, 252}
\definecolor{taskframe}{RGB}{100, 140, 190}

\usepackage{hyperref}
\hypersetup{
    colorlinks=true,
    linkcolor=primaryblue,
    citecolor=deepcyan,
    urlcolor=primaryblue,
    pdfauthor={ACM / IEEE-CS / AAAI Joint Task Force},
    pdftitle={计算机科学课程体系规范 CS2023 (中文全译本)},
    pdfsubject={Computer Science Curricula 2023 (Version Beta)},
    pdfkeywords={CS2023, 计算机教育, 课程规范, ACM, IEEE-CS, AAAI}
}

% ==================== 页眉页脚 ====================
\usepackage{fancyhdr}
\pagestyle{fancy}
\fancyhf{}
\fancyhead[LE,RO]{\songti\small\textcolor{darkgray}{\thepage}}
\fancyhead[LO]{\songti\small\textcolor{darkgray}{计算机科学课程大纲 CS2023（中文全译本）}}
\fancyhead[RE]{\songti\small\textcolor{darkgray}{\leftmark}}
\renewcommand{\headrulewidth}{0.4pt}
\renewcommand{\footrulewidth}{0pt}

% ==================== 列表紧凑化 ====================
\usepackage{enumitem}
\setlist{nosep, leftmargin=1.5em, labelsep=0.5em, itemsep=2pt, topsep=2pt}

% ==================== 彩色卡片环境 ====================
\usepackage{tcolorbox}
\tcbuselibrary{skins,breakable}

% 序言与变革卡片
\newtcolorbox{preamblebox}[1][]{
    enhanced,
    breakable,
    colback=boxbg,
    colframe=primaryblue,
    arc=3pt,
    boxrule=1pt,
    left=8pt,
    right=8pt,
    top=6pt,
    bottom=6pt,
    fonttitle=\heiti\bfseries\fontsize{8pt}{9.6pt}\selectfont,
    coltitle=white,
    title=#1
}

% 真实能力任务卡片 (Competency Specification Task)
\newtcolorbox{taskbox}[1][]{
    enhanced,
    colback=taskbg,
    colframe=taskframe,
    arc=2pt,
    boxrule=0.8pt,
    left=7pt,
    right=7pt,
    top=5pt,
    bottom=5pt,
    fonttitle=\heiti\bfseries\fontsize{8pt}{9.6pt}\selectfont,
    coltitle=white,
    title=#1
}

% 核心重点与技能提示盒
\newtcolorbox{infobox}[1][]{
    enhanced,
    breakable,
    colback=white,
    colframe=deepcyan,
    arc=2pt,
    boxrule=0.8pt,
    left=7pt,
    right=7pt,
    top=5pt,
    bottom=5pt,
    fonttitle=\heiti\bfseries\fontsize{8pt}{9.6pt}\selectfont,
    coltitle=white,
    title=#1
}

% ==================== 标题格式定制 ====================
\ctexset{
    chapter = {
        format = \heiti\zihao{3}\bfseries\raggedright,
        name = {第,章},
        number = \arabic{chapter},
        beforeskip = 12pt,
        afterskip = 10pt
    },
    section = {
        format = \heiti\zihao{4}\bfseries\raggedright,
        beforeskip = 10pt,
        afterskip = 5pt
    },
    subsection = {
        format = \heiti\zihao{-4}\bfseries\raggedright,
        beforeskip = 6pt,
        afterskip = 3pt
    },
    subsubsection = {
        format = \heiti\zihao{5}\bfseries\raggedright,
        beforeskip = 5pt,
        afterskip = 2pt
    }
}

% ==================== 文档正文 ====================
\begin{document}

% 封面前置
\begin{titlepage}
    \centering
    \vspace*{3cm}
    {\zihao{-0}\bfseries\songti 计算机科学课程体系 2023\par}
    \vspace{0.8cm}
    {\zihao{2}\textbf{Computer Science Curricula 2023 (CS2023)}\par}
    \vspace{2.5cm}
    
    {\zihao{4}\songti
    \textbf{计算机课程联合工作组 编著}\par
    \vspace{0.4cm}
    美国计算机学会（ACM）\par
    电气电子工程师学会计算机学会（IEEE-CS）\par
    国际人工智能促进协会（AAAI）\par
    }
    
    \vfill
    
    {\zihao{3}\songti\textbf{Minsecrus 译}\par}
    \vspace*{3cm}
\end{titlepage}

\frontmatter
\fontsize{8pt}{9.6pt}\selectfont
\tableofcontents

\mainmatter
\fontsize{8pt}{9.6pt}\selectfont

% 引入各章节
\chapter{导言与总体框架}

\section{CS2023 的研制背景与工作历程}

\begin{preamblebox}[历史沿革与演进]
半个多世纪以来，随着计算机科学学科的飞速发展与广泛渗透，国际权威计算学术组织不断更新计算机科学本科课程指导纲要：
\begin{itemize}
    \item \textbf{Curriculum 68}：50 多年前由美国计算机学会（ACM）发布的首部计算机科学本科课程规范指导。
    \item \textbf{Curriculum 78}：对课程体系进行了系统性重构，首次确立了“核心课程”与“选修课程（选修）”的分类组织模式。
    \item \textbf{Computing Curricula 1991 (CC1991)}：ACM 首次携手电气电子工程师学会计算机学会（IEEE-CS），联合研制了全新指导方针。
    \item \textbf{Computing Curricula 2001 (CC2001)}：首次明确将关注点纯化为计算机科学本体，将计算机工程、软件工程等分支独立成各自专有的课程规范。
    \item \textbf{Computer Science Curriculum 2008 (CS2008)}：作为对 CC2001 的阶段性重大修订版本。
    \item \textbf{Computer Science Curricula 2013 (CS2013)}：ACM 与 IEEE-CS 联合发布的里程碑式规范，确立了三级认知层次与 Tier~1/Tier~2 核心学时划分。
    \item \textbf{CS2023}：作为计算机科学课程体系的最新一代纲领性规范，由 ACM、IEEE-CS 联合发起，并首次正式吸纳国际人工智能促进协会（AAAI）作为共同主办单位。
\end{itemize}
\end{preamblebox}

\subsection{联合工作组的建立与运作架构}
ACM、IEEE-CS 与 AAAI 联合工作组组织架构如下：
\begin{enumerate}
    \item \textbf{指导委员会}：共 17 位领域权威专家，由 ACM 共同主席 Amruth N. Kumar 与 IEEE-CS 共同主席 Rajendra K. Raj 领衔，成员涵盖来自高校、工业界（如 Google 等）与专业教育委员会的代表。
    \item \textbf{17 个知识领域专家委员会}：每个知识领域均由指导委员会成员带头，吸纳全球 5 至 10 位本领域的顶尖学者与行业专家组成。
\end{enumerate}

\subsection{研制原则}
CS2023 工作组在研制过程中遵循以下五大核心原则：
\begin{itemize}
    \item \textbf{全球深度协作}：每个知识领域均由国际产学研联合委员会研制，汇聚多元学术视角与行业最新实践。
    \item \textbf{持续同行评审}：每一版草案均由至少三位外部独立资深专家进行匿名评审。
    \item \textbf{广泛社群参与}：通过十余个 ACM SIG 专委会、SIGCSE 年会等渠道面向全球计算机教育界征集反馈。
    \item \textbf{实据驱动革新}：充分调研了全球各高校（包括美国本土与国际高校）以及 865 家工业界企业的反馈数据。
    \item \textbf{全程透明公开}：通过官方平台面向公众开放全部过程文档与草案，吸纳全行业建设性意见。
\end{itemize}

\section{CS2023 的核心变革与重大创新}

相较于 CS2013 及以往单纯以“知识模型”为主的框架，CS2023 引入了范式级的升级与扩充。

\subsection{知识模型与能力模型的深度双轨融合}
自 Curriculum 68 到 CS2013，计算机课程规范均属于经典的“知识模型”——自顶向下拆解为知识领域 (KA)、知识单元 (KU)与具体主题，并针对主题指定学习成果。

随着国际工程教育对毕业生真实职业能力要求的提升，CC2020、IT2017 等前沿规范引入了“能力模型”。能力被严格定义为：\textbf{在具体任务情境中，融合应用知识、技能与专业素养特质的综合能力}。

CS2023 并没有盲目废弃知识模型，而是创新性地实现了两者的双轨互补：
\begin{itemize}
    \item \textbf{知识模型面向教学组织}：为教师搭建系统、连贯、层次清晰的知识传授阶梯。
    \item \textbf{能力模型面向产业实战}：以行业真实工程任务为牵引，反向推导培养学生所需的知识、技能与职业思维。
\end{itemize}

\subsection{知识领域的优化与重组}
根据近十年来计算技术的巨变，CS2023 对各大知识领域进行了重大重命名与战略调整：
\begin{table*}[htbp]
\centering\small
\begin{tabularx}{\textwidth}{lXl}
\toprule
\textbf{CS2013 领域名称} & \textbf{CS2023 领域名称} & \textbf{变革主旨说明} \\
\midrule
离散结构 & 数学与统计学基础 (\textbf{MSF}) & 强化概率、现代统计学及数据分析数学基础 \\
图形学与可视化 & 图形学与交互技术 (\textbf{GIT}) & 涵盖现代实时渲染、XR 交互与沉浸式体验 \\
信息保障与安全 & 计算机安全 (\textbf{SEC}) & 聚焦全栈安全防护、威胁建模与系统韧性 \\
信息管理 & 数据管理与数据库 (\textbf{DM}) & 全面涵盖非关系型数据库、大数据处理管道与治理 \\
智能系统 & 人工智能 (\textbf{AI}) & 全面拥抱深度学习、强化学习、大语言模型与伦理 \\
基于平台的开发 & 专用平台开发 (\textbf{SPD}) & 涵盖移动端、云原生、物联网与嵌入式平台生态 \\
编程语言 & 编程语言理论基础 (\textbf{FPL}) & 突出类型理论、形式化语义、并发语义基础 \\
社会问题与专业实践 & 社会、伦理与专业素养 (\textbf{SEP}) & 提升至核心贯穿层面，各知识领域均增设 SEP \\
计算科学 & \emph{移除独立 KA} & 其核心方法论下沉融入数学与算法交叉模块 \\
\bottomrule
\end{tabularx}
\end{table*}

\subsection{社会与职业伦理全面渗透机制}
在现代计算无处不在的今天，社会伦理不能再作为孤立的一门课程“轻描淡写”。CS2023 设立了革命性的要求：\textbf{全书所有知识领域均设立专有的 SEP 单元}，例如在算法中探讨“算法歧视与差分隐私”，在操作系统中探讨“隐私监控与资源隔离伦理”，在人工智能中深入“可解释性、对齐与公平性”。

\subsection{对齐布鲁姆教育分类学的四级技能标准}
CS2023 摒弃了过往模糊的 Familiarity/Usage/Assessment 分级，全面采用修订版布鲁姆认知目标分类学（Bloom's Taxonomy）：
\begin{enumerate}
    \item \textbf{解释}：识记与理解。行为动词包括：定义、描述、列举、辨析、阐述、总结。
    \item \textbf{应用}：实操与执行。行为动词包括：计算、配置、调试、部署、实验、测量、测试、追踪、虚拟化。
    \item \textbf{评估}：分析与评判。行为动词包括：比较、鉴别、归类、权衡、评审、防御、推荐、排序。
    \item \textbf{开发}：综合与创造。行为动词包括：设计、构建、合成、组合、重构、开发、重构、生产。
\end{enumerate}

\subsection{核心学时与“向日葵模型”}
CS2013 采用 Tier~1 (165学时) 与 Tier~2 (143学时) 的刚性要求，导致总核心学时逐年膨胀，给地方院校和新兴交叉学科带来巨大排课负担。

CS2023 提出了兼具通用性与灵活度的\textbf{“向日葵”模型}：
\begin{itemize}
    \item \textbf{花盘中心：计算机核心}：所有受规范认可的计算机专业毕业生\textbf{必须百分之百掌握}的核心基础底座。
    \item \textbf{花瓣外延：领域核心}：若高校根据自身办学特色开设某特定专业方向（如系统方向、应用方向、软件工程方向），则在涉猎该知识领域时需达成的深度核心标准。
\end{itemize}
这种设计允许高校在保障核心扎实的前提下，根据区域产业需求或办学特长，在某些“花瓣”上做深做强，而在其余花瓣保持基准，极大增强了教学弹性。

\section{计算机专业毕业生核心素养画像}

CS2023 指出，面对日新月异的技术变革，当代计算机科学本科毕业生应具备以下核心特质：
\begin{itemize}
    \item \textbf{算法问题求解能力}（Algorithmic Problem-solver）：能够在适宜的抽象层次上，针对常见问题设计良好的算法解决方案。
    \item \textbf{合格的程序设计与实现能力}（Competent Programmer）：具备规范的代码编写、调试与重构能力，掌握基本软件工程原则。
    \item \textbf{具备计算的心智模型}（Mental Model of Computation）：对计算机底层硬件、系统抽象及程序运行时机制具备清晰的认知与心智模型。
    \item \textbf{终身学习与适应力}（Life-long Learner \& Adaptable）：面对快速涌现的新技术与不确定性，具备自主学习与适应新技术栈的能力。
    \item \textbf{多层次抽象思维}（Think at Multiple Levels of Abstraction）：理解从微架构、指令集、操作系统到应用层及分布式系统的分层体系，具备跨层次分析与故障定位意识。
    \item \textbf{对抗性思维与安全防范直觉}（Adversarial Thinker）：具备基本的系统安全性视角、防御意识与严密的逻辑分析能力。
    \item \textbf{社会责任与伦理意识}（Socially Responsible \& Ethical Behavior）：理解代码与计算技术对社会、隐私、安全及伦理规范的现实影响，恪守职业准则。
    \item \textbf{跨学科协作与有效沟通}（Collaborative \& Effective Communication）：能够理解非计算机学科的基本背景，与跨领域团队成员协作并清晰表达技术方案。
\end{itemize}

\section{高校计算机教育面临的时代挑战}

报告深刻剖析了全球高校在推进计算机课程建设时普遍面临的严峻挑战：
\begin{enumerate}
    \item \textbf{生成式人工智能（AIGC）对编程教育的重构}：自动化代码生成工具对传统作业与考试评估机制带来根本性冲击，教育者亟需将 AI 深度融入核心教学并重构评价指标。
    \item \textbf{生源暴增与师资供需失衡}：计算机专业学生人数持续爆发，师资招聘与留存成本高企。
    \item \textbf{学科边界细分与碎片化挑战}：数据科学、人工智能、网络空间安全纷纷独立建院或设系，如何保持计算机大类核心基础的凝聚力成为课题。
    \item \textbf{计算机通识教育需求膨胀}：如何为全校非计算机专业学生提供高质量、易吸收的现代计算思维教育。
\end{enumerate}

\chapter{算法基础与复杂性}

\section{领域导言与演进背景}

\begin{preamblebox}[知识领域概述与变革]
算法与数据结构是计算机科学与软件工程不可动摇的基石。任何理论计算模型与现实工程软件，其本质均是在具备特定物理与逻辑结构的数据元素上执行算法操作。面对复杂的现实工程挑战，选择并设计出兼具计算效率、资源节约（包括能耗与碳足迹）与社会公平性的算法，要求从业者不仅理解算法的经验表现，更要洞察其内在的理论极限与计算复杂性。这种理论素养能使计算机专业人员超越具体的编程语言、编程范式或特定底层硬件架构，直击计算问题的本质规律。

\textbf{自 CS2013 以来的核心变革：}
\begin{enumerate}
    \item \textbf{重命名与回归理论根基}：从 CS2013 的“算法与复杂性”拓展为“算法基础”，重新强调所有计算机专业毕业生必须牢固掌握计算理论基础（自动机、形式语言、可计算性与不可判定性）。
    \item \textbf{增设“算法与社会”核心单元}：面对算法在现代社会决策中的无所不在，专门增设该单元，系统涵盖算法偏见、公平性度量、差分隐私与绿色低能耗计算等前沿社会议题。
    \item \textbf{与软件开发基础（SDF）的协同映射}：明确划分了初级数据结构应用与高级算法分析的边界，形成循序渐进的进阶路径。
\end{enumerate}
\end{preamblebox}

\section{核心学时分配建议}

本知识领域在“向日葵模型”下的学时推荐如下表所示：

\begin{table}[H]
\centering
\setlength{\tabcolsep}{4pt}\fontsize{8pt}{9.6pt}\selectfont\begin{tabularx}{\columnwidth}{Y c c}
\toprule
\textbf{知识单元} & \textbf{计算机核心} & \textbf{领域核心} \\
\midrule
基础数据结构与算法 & 36 & 7 \\
算法策略与范式 & 4 & -- \\
计算复杂性分析 & 4 & 2 \\
计算模型与形式理论 & 9 & 23 \\
算法与社会伦理 & 4 & -- \\
\midrule
\textbf{合计} & \textbf{57 学时} & \textbf{32 学时} \\
\bottomrule
\end{tabularx}
\end{table}

\section{知识单元与详细大纲}

\subsection{基础数据结构与算法}
\textbf{【学时建议】}：36 计算机核心学时，7 领域核心学时。

本单元构建起所有计算机科学专业学生必须掌握的数据结构与核心算法底座，与 SDF（软件开发基础）和 MSF（离散数学）紧密交织。

\subsubsection{计算机核心主题}
\begin{itemize}
    \item \textbf{抽象数据类型 (ADT)}：包（Bag）、容器、字典（Dictionary）、列表（List）与集合（Set）的代数规范与行为不变量。
    \item \textbf{数组与线性表}：一维与多维数组存储映射；线性查找与二分查找算法及其边界条件。
    \item \textbf{链表结构}：单向链表、双向链表、循环链表及其指针操作安全性与内存局部性特征。
    \item \textbf{栈与队列}：后进先出（LIFO）与先进先出（FIFO）机制、双端队列、基于堆的优先队列。
    \item \textbf{哈希表与散列技术}：哈希函数设计原理、冲突消解策略（线性探测、二次探测、链地址法、一致性哈希初步与动态重哈希）。
    \item \textbf{字符串与文本处理}：高效模式匹配算法（KMP 算法、Boyer-Moore 算法）。
    \item \textbf{树形结构与遍历}：
    \begin{itemize}
        \item 二叉树、多叉树、二叉搜索树（BST）及其前序、中序、后序与层序遍历。
        \item 平衡二叉树与多路平衡树原理：AVL 树、红黑树、2-3 树、B 树/B+ 树的核心旋转与重平衡机制。
        \item 堆结构（最大堆/最小堆）与其在优先队列中的工程实现。
        \item 编码与贪心应用：霍夫曼树与霍夫曼编码（Huffman Coding）。
    \end{itemize}
    \item \textbf{经典排序算法族}：
    \begin{itemize}
        \item $O(n^2)$ 经典排序：选择排序、插入排序、冒泡排序及其稳定性和原地性分析。
        \item $O(n \log n)$ 高效排序：快速排序（基准元划分策略与最坏情况退化规避）、归并排序、堆排序。
        \item 线性时间排序：计数排序、桶排序、基数排序及其应用边界。
    \end{itemize}
    \item \textbf{图结构与图遍历算法}：
    \begin{itemize}
        \item 图的拓扑属性：有向/无向、有环/无环（DAG）、连通性、带权图。
        \item 图的机器表示：邻接矩阵与邻接表的时空权衡。
        \item 基础遍历：广度优先搜索（BFS）及其在无权最短路中的应用；深度优先搜索（DFS）及其在环检测、连通分量与拓扑排序中的应用。
        \item 最小生成树（MST）：Prim 算法与 Kruskal 算法（并查集 Union-Find 的应用）。
        \item 最短路径算法族：单源最短路径 Dijkstra 算法（贪心策略）、Bellman-Ford 算法（负权边检测）；全源最短路径 Floyd-Warshall 算法。
        \item 传递闭包与 Warshall 算法。
    \end{itemize}
    \item \textbf{现代核心交叉技术}：
    \begin{itemize}
        \item 密码学基础散列与非对称基础（SHA-256、RSA 的数论支撑与实现逻辑）。
        \item 差分隐私（Differential Privacy）在聚合统计算法中的注入机制。
        \item 算法不变量的形式化证明与前置/后置条件。
    \end{itemize}
\end{itemize}

\subsubsection{领域核心与选修主题}
\begin{itemize}
    \item \textbf{领域核心}：
    \begin{itemize}
        \item 分布式共识算法（Consensus Algorithms，如 Paxos、Raft、区块链工作量证明 PoW 与权益证明 PoS 机制）。
        \item 计算几何基础（凸包 Graham 扫描法、最近点对分治求解）。
        \item 运筹与线性规划（Simplex 单纯形法与对偶原理）。
        \item 近似算法（针对背包问题、旅行商 TSP 的常数比例近似）。
        \item 随机化算法（MaxCut 局部搜索、球与箱模型 Balls and Bins、跳跃表 SkipList）。
    \end{itemize}
    \item \textbf{拓展选修}：
    \begin{itemize}
        \item 量子计算算法（Deutsch-Jozsa、Shor 大数质因数分解、Grover 数据库量子搜索）。
        \item 信号处理算法（快速傅里叶变换 FFT、离散卷积计算）。
        \item 遗传算法与演化计算（Genetic Algorithms、群体智能优化）。
    \end{itemize}
\end{itemize}

\subsubsection{示范性学习成果与能力目标}
\begin{enumerate}
    \item \textbf{[解释]} 针对本单元列出的所有基础 ADT 与数据结构，精准阐述其数学定义、不变量、多维底层存储表征及支持的增删查改原子操作。
    \item \textbf{[应用]} 给定具体的工业级应用场景（如高频交易订单撮合、海量日志检索），选择并实现最优的数据结构，证明其时空渐进复杂度的优越性。
    \item \textbf{[评估]} 全面对比各种排序、图搜索算法在最坏、最好、平均情况下的时间开销与内存局部性，评估真实硬件缓存环境对理论性能的影响。
    \item \textbf{[开发]} 针对复杂现实需求（例如基于多因子约束的路径规划系统），组合设计出混合数据结构，并撰写工程可行性与健壮性论证报告。
\end{enumerate}

\subsection{算法策略与范式}
\textbf{【学时建议】}：4 计算机核心学时。

本单元系统提炼通用的算法设计模式与高级思维方法：
\begin{itemize}
    \item \textbf{穷举与蛮力法}：组合搜索空间生成与状态遍历。
    \item \textbf{分治法}：大问题独立子问题解构、合并技术（典型范例：二分查找、快速排序、归并排序、Strassen 矩阵乘法）。
    \item \textbf{减治法}：每次将问题规模递减常数或常数因子（典型范例：插入排序、二叉查找、拓扑排序）。
    \item \textbf{动态规划}：最优子结构与重叠子问题、备忘录自顶向下法与制表自底向上法（典型范例：最长公共子序列 LCS、背包问题、Floyd 算法）。
    \item \textbf{贪心策略}：贪心选择性质与无后效性证明（典型范例：哈夫曼编码、Dijkstra 最短路、Kruskal 最小生成树）。
    \item \textbf{回溯法与分支限界法}：状态剪枝、Alpha-Beta 剪枝与组合优化边界估算。
    \item \textbf{时空权衡与化归策略 (Space-Time Tradeoff \& Reduction)}：以空间换时间（散列表、预处理排序）与问题规约转化（2-3树转换、图问题向网络流转化）。
\end{itemize}

\subsection{计算复杂性分析}
\textbf{【学时建议】}：4 计算机核心学时，2 领域核心学时。

本单元建立算法效能的严谨数理评估框架：
\begin{itemize}
    \item \textbf{渐进记号体系}：$O$（渐进上界）、$\Omega$（渐进下界）、$\Theta$（紧确界）、$o$ 与 $\omega$ 的严谨数学定义及极限证明。
    \item \textbf{基础时间复杂度阶}：常数阶 $O(1)$、对数阶 $O(\log n)$、线性阶 $O(n)$、线性对数阶 $O(n \log n)$、多项式阶 $O(n^k)$、指数阶 $O(2^n)$、阶乘阶 $O(n!)$。
    \item \textbf{递推关系式求解}：主定理三大分支条件推导与应用、代入归纳法、递归树分析法。
    \item \textbf{均摊分析}：聚合分析、记账法、势能法（Potential Method，如动态扩容数组分析）。
    \item \textbf{难解性与计算复杂性理论}：
    \begin{itemize}
        \item 判定问题与搜索问题的多项式时间可解性（类 P）。
        \item 非确定性多项式时间（类 NP）与多项式时间验证机概念。
        \item NP-完全性与 Cook-Levin 定理、经典 21 个 NPC 问题规约链条（SAT、3-SAT、团问题、顶点覆盖、哈密顿回路、TSP、子集和）。
        \item P 对 NP 问题的哲学意义与计算局限。
    \end{itemize}
\end{itemize}

\subsection{计算模型与形式理论}
\textbf{【学时建议】}：9 计算机核心学时，23 领域核心学时。

回归计算理论核心，使学生掌握计算能力的本质边界：
\begin{itemize}
    \item \textbf{有穷自动机与正则语言}：确定性有穷自动机 (DFA)、非确定性有穷自动机 (NFA)、子集构造法等价转化、正则表达式与正则文法、泵引理 及其不可正则性证明。
    \item \textbf{上下文无关文法与下推自动机}：上下文无关文法产生式、语法分析树与歧义性、下推自动机 (PDA) 的栈式存储模型、非上下文无关语言的泵引理。
    \item \textbf{图灵机与丘奇-图灵论题 (Turing Machines \& Church-Turing Thesis)}：标准单带图灵机形式化定义、多带/非确定性图灵机等价性、通用图灵机、丘奇-图灵论题的深刻内涵。
    \item \textbf{可计算性与不可判定性 (Computability \& Undecidability)}：图灵可识别与图灵可判定语言、对角线法证明不可数性、停机问题 的不可判定性严密反证法、波斯特对应问题 (PCP) 与莱斯定理 (Rice's Theorem)。
\end{itemize}

\subsection{算法与社会伦理}
\textbf{【学时建议】}：4 计算机核心学时。

\begin{itemize}
    \item \textbf{算法偏见与歧视}：训练数据历史偏见与算法反馈循环；招聘筛选、司法量刑、信贷审批算法中的不公诱因与防范机制。
    \item \textbf{算法公平性度量标准}：人口统计学一致性、均等化赔率、预测均等 及其内在的数学互斥性定理。
    \item \textbf{透明度、可解释性与知情权 (Explainability \& Right to Explanation)}：黑箱算法决策的审计追踪、模型决策证据生成机制。
    \item \textbf{绿色计算与算法能耗}：大型模型训练与高并发图计算的碳排放评估、算力节约型与能耗感知型算法设计。
\end{itemize}

\section{专业素养与思维倾向}

在算法领域的学习与实践中，重点塑造学生的以下专业特质：
\begin{itemize}
    \item \textbf{严密分析性}：面对模糊混沌的现实需求，能剥离表象，建立精确的形式化数学模型。
    \item \textbf{精益求精与严谨性}：深入推敲极端边界情况、死循环风险与内存开销溢出。
    \item \textbf{坚韧不拔的探索精神}：在攻坚复杂 NP-难问题或多约束系统时，能持续尝试近似、启发式或松弛技术。
    \item \textbf{负责任的工程伦理}：时刻警惕算法可能引发的社会排斥、隐私泄露与能耗过载风险。
\end{itemize}

\section{数学基础要求}

算法领域对离散数学与应用数学有着高度依赖：
\begin{itemize}
    \item \textbf{必须具备的离散数学核心}：集合论、谓词逻辑、二元关系（等价关系、偏序关系）、图论（度数、回路、树的同构）、初等组合计数、形式化数学证明（直接反证法、良序原理、超限数学归纳法）。
    \item \textbf{概率论与数理统计}：期望与方差的线性性质、条件概率、大数定律、贝叶斯定理、马尔可夫不等式与切比雪夫不等式（用于随机算法与哈希性能分析）。
    \item \textbf{线性代数与微积分}：矩阵运算、特征值、函数的渐近极限比较（L'H\^opital 法则在复杂度比较中的应用）。
\end{itemize}

\section{课程设计与开课规划建议}

联合工作组推荐两套典型开课方案：

\subsection{三门课贯穿方案 (面向通用型计算机专业)}
\begin{enumerate}
    \item \textbf{CS1：计算机编程基础 I}（涵盖基础数组、字符串查找、基础对象与迭代遍历，渗透初阶时空直觉）。
    \item \textbf{CS2：面向对象与数据结构}（15 学时覆盖单/双链表、栈、队列、哈希表、二叉平衡树、快速排序、归并排序、SHA-256 算法及 2 学时算法社会学）。
    \item \textbf{CS3：算法分析与计算理论综合课 (Algorithms \& Computation)}（以 16 学时高阶图论、动态规划、NPC 理论结合 9 学时计算模型与自动机基础）。
\end{enumerate}

\subsection{四门课深化方案 (面向研究型与深层次理论方向)}
将理论计算模型单独开设为一门 32 学时完整的《计算理论形式化方法》（涵盖自动机、下推系统、图灵机、不可判定性与复杂性类深入），而在《高级算法设计与分析》中腾出学时专攻均摊分析、网络流、近似算法与几何算法。

\section{真实能力任务规范}

\begin{taskbox}[任务 1：急救车辆智能紧急调度路由系统]
\textbf{【任务情境】}：设计一个在复杂城市路网拥堵与突发施工事件下的救护车智能调度路由模块，以毫秒级响应将危重病人送往最优急救医院。\\
\textbf{【能力陈述】}：检索、评估并融合现存高效图算法（如动态加权 Dijkstra 算法、A* 启发式导航算法结合收缩层次结构 Contraction Hierarchies），解决带实时权值变动的计算问题。\\
\textbf{【能力领域】}：软件工程 \quad \textbf{【能力单元】}：系统设计与评估\\
\textbf{【核心要求】}：计算机核心 \quad \textbf{【达标等级】}：\textbf{评估}
\end{taskbox}

\begin{taskbox}[任务 2：高精度慈善募捐个性化推荐与资产配置]
\textbf{【任务情境】}：根据用户输入的多元化价值偏好与税务属性，从数千家合规慈善机构中自动权衡并推荐最佳组合。\\
\textbf{【能力陈述】}：综合运用多目标贪心策略、带权二分图最大权匹配与多维背包动态规划模型，设计并实现一套高吞吐、抗并发的算法求解器。\\
\textbf{【能力领域】}：软件工程 \quad \textbf{【能力单元】}：系统开发\\
\textbf{【核心要求】}：计算机核心 \quad \textbf{【达标等级】}：\textbf{开发}
\end{taskbox}

\begin{taskbox}[任务 3：遗留客户服务系统的算法性能重构]
\textbf{【任务情境】}：某企业老旧工单处理系统在千万级数据量下出现严重查询阻塞，需彻底消除 $O(n^2)$ 级隐藏瓶颈。\\
\textbf{【能力陈述】}：剖析系统瓶颈来源，使用红黑树平衡索引、布隆过滤器与缓存哈希表替换原生遍历，将核心响应时间降低数个数量级。\\
\textbf{【能力领域】}：软件工程 \quad \textbf{【能力单元】}：系统重构与性能改进\\
\textbf{【核心要求】}：计算机核心 \quad \textbf{【达标等级】}：\textbf{开发}
\end{taskbox}

\begin{taskbox}[任务 4：跨机构医疗患者全流程状态跟踪自动化]
\textbf{【任务情境】}：构建一套覆盖门诊预检、手术排期、病房流转直至出院随访的全生命周期严密跟踪算法。\\
\textbf{【能力陈述】}：基于确定性有穷状态机 (DFA) 与层次状态图建立无死锁、无异常旁路的业务状态转移模型，并做全分支可达性形式化验证。\\
\textbf{【能力领域】}：软件工程 \quad \textbf{【能力单元】}：需求分析与架构设计\\
\textbf{【核心要求】}：计算机核心 \quad \textbf{【达标等级】}：\textbf{开发}
\end{taskbox}

\begin{taskbox}[任务 5：多维度复杂排班优化问题的难解性论证]
\textbf{【任务情境】}：客户提出一个兼顾多班次工时、请假优先级、技能准入与奖惩成本的万能实时排班需求。\\
\textbf{【能力陈述】}：撰写专业技术白皮书，通过多项式时间规约证明该问题为强 NP-难（规约到 3-SAT 或精确覆盖问题），指出其精确最优求解在多项式时间内不可行，并推荐高效模拟退火或遗传近似求解方案。\\
\textbf{【能力领域】}：理论计算 \quad \textbf{【能力单元】}：可行性评估与技术论证\\
\textbf{【核心要求】}：计算机核心 \quad \textbf{【达标等级】}：\textbf{解释与应用}
\end{taskbox}

\begin{taskbox}[任务 6：数据最优无损压缩极限的理论不可判定性论证]
\textbf{【任务情境】}：某企业试图研发一种“对任意输入数据均能产生理论极限最优压缩率且保证无损还原”的单体通用算法。\\
\textbf{【能力陈述】}：撰写严谨的理论学术报告，通过反证法证明该最优性判定问题可归约为图灵机停机问题（柯尔莫哥洛夫复杂性不可计算定理），以严密逻辑澄清科学边界，避免企业投入无谓研发资本。\\
\textbf{【能力领域】}：理论计算 \quad \textbf{【能力单元】}：技术论证\\
\textbf{【核心要求】}：计算机核心 \quad \textbf{【达标等级】}：\textbf{评估}
\end{taskbox}

\begin{taskbox}[任务 7：全民点赞与建议评级推荐算法的社会与伦理审计]
\textbf{【任务情境】}：社区自治平台拟推行基于热度点赞的市民建议自动排序与采纳机制。\\
\textbf{【能力陈述】}：撰写算法社会影响评估报告，揭示马太效应与沉默螺旋效应，设计差分隐私噪声注入与长尾代表性补偿模型，防止少数群体利益被算法遮蔽。\\
\textbf{【能力领域】}：软件与社会 \quad \textbf{【能力单元】}：社会议题调适\\
\textbf{【核心要求】}：计算机核心 \quad \textbf{【达标等级】}：\textbf{解释与评估}
\end{taskbox}

\begin{strip}
\section{附录：核心主题与技能要求等级对照表}
\noindent\begin{tabularx}{\textwidth}{p{0.18\textwidth} Y c c c}
\toprule
\textbf{知识单元} & \textbf{核心知识点与具体主题} & \textbf{认知等级} & \textbf{分类} & \textbf{学时} \\
\midrule
\multirow{4}{*}{\makecell[l]{算法/基础数据结构\\与算法}} 
& 常用 ADT（字典、映射、集合）、基础线性表与树 & 解释 & 计算机核心 & 12 \\
& 经典查找、快速排序、归并排序、堆排序实现与优化 & 应用 & 计算机核心 & 12 \\
& 图的连通性、Dijkstra 最短路、Prim/Kruskal 最小生成树 & 应用 & 计算机核心 & 12 \\
& 线性规划单纯形法、凸包几何算法、近似与随机算法 & 评估 & 领域核心 & 7 \\
\midrule
\multirow{2}{*}{\makecell[l]{算法/算法策略与范式}} 
& 分治法、减治法、动态规划、贪心法、回溯与分支限界原理 & 解释 & 计算机核心 & 2 \\
& 针对工程问题选择或重构适宜的算法范式进行求解 & 应用 & 计算机核心 & 2 \\
\midrule
\multirow{3}{*}{\makecell[l]{算法/复杂性分析}} 
& 渐进符号体系（大 O、大 $\Omega$、大 $\Theta$）与主定理推导 & 解释 & 计算机核心 & 2 \\
& 递归与非递归算法的时空开销度量、经验评测 & 应用 & 计算机核心 & 2 \\
& P 与 NP 问题概念、NPC 经典问题规约链条、均摊分析 & 评估 & 领域核心 & 2 \\
\midrule
\multirow{3}{*}{\makecell[l]{算法/计算模型与\\形式理论}} 
& 正则语言、DFA 与 NFA 等价性、正则表达式转化 & 应用 & 计算机核心 & 5 \\
& 上下文无关文法、下推自动机、图灵机定义与停机问题 & 解释 & 计算机核心 & 4 \\
& 泵引理证明、非确定图灵机、莱斯定理、不可判定性层级 & 开发 & 领域核心 & 23 \\
\midrule
\multirow{2}{*}{\makecell[l]{算法/算法与社会}} 
& 算法偏见来源、公平性数学度量（统计均等 vs 均等化赔率） & 解释 & 计算机核心 & 2 \\
& 差分隐私扰动注入、绿色计算与算法能耗感知优化 & 评估 & 计算机核心 & 2 \\
\bottomrule
\end{tabularx}
\end{strip}

\chapter{计算机体系结构与组织}

\section{领域导言与演进背景}

\begin{preamblebox}[知识领域概述与变革]
在计算机学科的工程实践中，专业人员通常花费大量时间编写高能效代码，以解决特定应用领域的复杂计算问题。随着摩尔定律（Moore's Law）与登纳德缩放定律面临物理极限与终结，几乎所有计算系统（包括最普及的商用消费级硬件）均日益依赖硬件系统层面的并行性来满足持续增长的计算性能需求。这种从传统的纯顺序处理向高度并行计算的根本范式转变，要求软件从业者对底层计算机体系结构建立更为深刻透彻的理解。体系结构已不再能被视作一个可以任意套用抽象原则的“黑盒”；相反，现代程序员必须能够“透视黑盒内部”，主动利用特定的底层硬件架构组件与微架构特性，以系统性提升计算性能并大幅优化能源效率。

“计算机体系结构与组织”（Architecture and Organization, AR）知识领域旨在培养学生对现代一切计算基石——底层硬件运行环境，以及硬件向高层软件栈所提供各类接口的深层次理解。本领域涵盖的目标硬件跨度广泛，从低功耗物联网嵌入式处理器，直至面向超算与云端的大规模企业级多处理器系统。

深入学习本领域的知识主题，将使学生全面领会现代计算机系统的核心体系结构设计原则，特别是如何攻克驾驭硬件并行性以维持未来持续性能与能效提升的重大挑战。此外，深刻理解计算机体系结构是开发兼具高性能与超低能耗程序的必由之路。本知识领域将协助计算机科学专业学生彻底告别“黑盒化”思维，建立对底层计算机系统的全局敏锐洞察，并充分发掘与释放特定硬件架构所能实现的极致效能。

\textbf{自 CS2013 以来的核心变革：}
\begin{enumerate}
    \item \textbf{核心知识单元全面重构与更新}：重点重塑了“体系结构/存储体系”（体系结构/层次存储系统）与“体系结构/性能与能效”（体系结构/性能与能效）两个知识单元，深度融合了片上多级缓存优化、高带宽存储以及能耗感知计算的最新前沿进展。
    \item \textbf{新增“异构架构”核心单元（体系结构/异构架构）}：直面计算机体系结构领域的新兴突破，系统覆盖存算一体 (PIM)、领域专用架构（Domain-Specific Architectures, DSA，如神经网络处理器 NPU/TPU）、片上互联与先进封装技术（3DIC / Chiplets 芯粒）等前沿架构。
    \item \textbf{新增“量子体系结构”先锋单元（体系结构/量子体系结构）}：鉴于量子计算与量子处理器架构研究的飞速推进，专门增设该单元作为前沿“工具箱”，系统涵盖量子比特物理实现、量子线路、量子门、基础量子算法及量子处理单元（QPU）的微架构组织。
    \item \textbf{整合交叉重叠单元}：将原 CS2013 中的“多处理与替代架构”（体系结构/Multiprocessing and Alternative Architectures）整合并入全新的“异构架构”单元中，使知识逻辑更加紧凑清晰。
\end{enumerate}
\end{preamblebox}

\section{核心学时分配建议}

本知识领域在“向日葵模型”下的学时推荐如下表所示：

\begin{table}[H]
\centering
\setlength{\tabcolsep}{3.5pt}\fontsize{8pt}{9.6pt}\selectfont\begin{tabularx}{\columnwidth}{Y c c}
\toprule
\textbf{知识单元} & \textbf{计算机核心} & \textbf{领域核心} \\
\midrule
数字逻辑与数字系统 & -- & 3 \\
机器级数据表示 & 1 & -- \\
汇编级机器组织 & 1 & 2 \\
层次存储系统 & 6 & -- \\
接口与通信 & 1 & -- \\
功能部件组织 & -- & 2 \\
性能与能效 & -- & 3 \\
异构架构 & -- & 3 \\
量子体系结构 & -- & 3 \\
\midrule
\textbf{合计} & \textbf{9 学时} & \textbf{16 学时} \\
\bottomrule
\end{tabularx}
\end{table}

\section{知识单元与详细大纲}

\subsection{数字逻辑与数字系统}
\textbf{【学时建议】}：3 领域核心学时。

本单元建立现代计算系统的底层数字电路与硬件抽象基础：

\subsubsection{领域核心主题}
\begin{itemize}
    \item \textbf{体系结构历史与演进概述}：计算技术的发展历程（电子管、晶体管、集成电路至超大规模集成电路 VLSI），冯·诺依曼结构的发展演变。
    \item \textbf{组合逻辑与时序逻辑 / 现场可编程门阵列 (FPGA)}：
    \begin{itemize}
        \item 基础组合逻辑构建块：逻辑门、多路选择器、译码器、优先编码器与算术逻辑单元（ALU）。
        \item 时序逻辑基本构建块：双稳态锁存器、触发器（Flip-Flop，D 触发器与 JK 触发器）、时序寄存器、计数器与有限状态机（FSM）。
    \end{itemize}
    \item \textbf{多层硬件与软件抽象体系结构}：晶体管层、门级电路、寄存器传输级（RTL）、指令集体系结构（ISA）、操作系统内核及应用程序层。
    \item \textbf{计算机辅助设计 (CAD) 工具链}：处理硬件表示与体系结构模拟的综合、布局布线与仿真工具。
    \item \textbf{高层次综合}：
    \begin{itemize}
        \item 寄存器传输级表示法。
        \item 硬件描述语言（HDL，如 Verilog、VHDL 与现代化 Chisel 语言）。
    \end{itemize}
    \item \textbf{片上系统 设计流程}：软硬件协同设计、IP 核复用、时钟树综合与物理验证流程。
    \item \textbf{物理与电气工程约束}：
    \begin{itemize}
        \item 门延迟与建立/保持时间。
        \item 扇入与扇出能力。
        \item 能量与功耗约束（动态开关功耗与静态漏电功耗）。
        \item 信号传播物理极限（光速在硅介质中的传播延迟）。
    \end{itemize}
\end{itemize}

\subsubsection{示范性学习成果与能力目标}
\begin{enumerate}
    \item \textbf{[解释]} 评述计算机底层组件技术的历史演进——从真空电子管到现代超大规模集成电路（VLSI），从传统单体大型机架构演化至现代仓储级计算机 (WSC)的分布式数据中心集群架构。
    \item \textbf{[评估]} 深入分析现代异构计算机体系结构中组件之间及组件内部的并行性形态与数据依赖约束。
    \item \textbf{[解释]} 详尽阐释“功耗墙”（Power Wall）产生的物理根源，剖析其为何终结了单核频率的提升，并成为全面发掘与驾驭硬件并行性的核心瓶颈。
    \item \textbf{[开发]} 提出计算机基础构建模块的严密设计方案：算术逻辑单元 ALU（门级实现）、寄存器与寄存器堆（门级实现）、中央处理器 CPU（寄存器传输级 RTL 实现）以及各级存储器（寄存器传输级 RTL 实现）。
    \item \textbf{[评估]} 评估并验证简易计算机硬件设计的基础构建模块（如算术逻辑单元、寄存器组及内部总线数据流动）的时序与逻辑正确性。
    \item \textbf{[评估]} 验证简易处理器的时序图行为特征，精准识别并定位微架构中的数据依赖与时序冒险风险。
\end{enumerate}

\subsection{机器级数据表示}
\textbf{【学时建议】}：1 计算机核心学时。

本单元建立机器底层二进制数据存储、编码与精度约束的数理理解：

\subsubsection{计算机核心主题}
\begin{itemize}
    \item \textbf{位、字节与字}：二进制位、字节与机器字长的体系结构定义，自然对齐与非对齐访存。
    \item \textbf{数值数据表示与进位制系统}：
    \begin{itemize}
        \item 定点数表示法。
        \item 浮点数表示法（Floating-point representation，IEEE 754 标准：阶码、尾数、规格化数、非规格化数、NaN 与无穷大值）。
    \end{itemize}
    \item \textbf{有符号数与补码表示}：原码、反码与补码（Two's-complement）表示法，算术溢出判定与符号位扩展机制。
    \item \textbf{非数值数据的位级表示}：ASCII 与 Unicode（UTF-8、UTF-16）字符集编码规范，图像、音频等多媒体数据的位级量化表征。
    \item \textbf{复合数据结构的内存布局}：结构体/记录的字节对齐与填充规则，多维数组在内存中的行优先/列优先连续排布。
\end{itemize}

\subsubsection{示范性学习成果与能力目标}
\begin{enumerate}
    \item \textbf{[解释]} 深刻阐释现代计算机体系中“一切皆为数据”（包括程序指令自身亦以二进制数据形态存储执行）的冯·诺依曼同构存储哲学及其技术影响。
    \item \textbf{[应用]} 结合数值表示位分布图，标定与剖析由于定长二进制表示局限而引发精度损失、舍入误差、数值截断及上下溢出的敏感区间。
    \item \textbf{[解释]} 评述原码与补码在存储负整数时的底层机制差异，阐明补码在统一加减法硬件运算与消除双零歧义上的计算优势。
    \item \textbf{[解释]} 结合实际工程场景，合理论证并解释不同数值精度格式（如 FP64、FP32、FP16、BF16 及 INT8）在科学计算与深度学习模型推理中的选型依据。
    \item \textbf{[解释]} 阐述非数值数据（如字符、字符串、结构体记录与多维数组）在底层内存中的位级排布、对齐填充与指针寻址机制。
    \item \textbf{[应用]} 熟练完成不同数制（二进制、八进制、十进制、十六进制）与数值表示格式之间的精确手工与程序化相互转换。
\end{enumerate}

\subsection{汇编级机器组织}
\textbf{【学时建议】}：1 计算机核心学时，2 领域核心学时。

本单元建立从高级语言到机器底层指令执行与微架构映射的核心桥梁：

\subsubsection{计算机核心主题}
\begin{itemize}
    \item \textbf{冯·诺依曼计算机体系结构}：存储程序概念、程序计数器（PC）、总线瓶颈与核心功能子系统（计算、控制、存储、输入、输出）划分。
    \item \textbf{控制单元与指令周期}：取指 (IF)、译码 (ID)与执行 (EX)的微周期节拍机制。
    \item \textbf{费林分类法 (Flynn's Taxonomy) 导论}：单指令流单数据流（SISD）、单指令流多数据流（SIMD）与多指令流多数据流（MIMD）的体系结构特征与演进。
    \item \textbf{共享内存多处理器 / 多核组织架构}：多核对称多处理器（SMP）、片上互联与共享总线微架构（与操作系统 OS 交叉）。
\end{itemize}

\subsubsection{领域核心主题}
\begin{itemize}
    \item \textbf{指令集体系结构 (ISA)}（涵盖典型 CISC 与 RISC 架构代表：x86、ARM 与开源开放的 RISC-V）：
    \begin{itemize}
        \item 指令编码格式（操作码 Opcode、寄存器操作数、立即数）。
        \item 数据传送、算术与逻辑运算、程序控制流指令及基本 I/O 指令。
        \item 寻址模式（直接寻址、立即寻址、间接寻址、基址变址寻址与 PC 相对寻址）。
        \item 机器语言编程与机器码格式。
        \item 汇编语言程序设计技术。
    \end{itemize}
    \item \textbf{子程序调用与返回机制}：调用约定、参数传递、返回值保存、栈帧平衡与寄存器保护机制（调用者保存 vs 被调用者保存，交叉引用 PL/编程语言翻译与执行，与 OS 交叉）。
    \item \textbf{I/O 与中断处理机制}：中断请求（IRQ）、中断向量表（IVT）与保护现场（与 OS 交叉）。
    \item \textbf{程序运行时内存段划分}：堆区、静态数据区、栈区与代码段（与 OS 交叉）。
\end{itemize}

\subsubsection{示范性学习成果与能力目标}
\begin{enumerate}
    \item \textbf{[解释]} 在现代嵌入式系统语境下（尤其是紧耦合片上内存与高延迟片外存储系统）分析并定位经典冯·诺依曼功能单元的物理映射与作用。
    \item \textbf{[解释]} 详细阐述经典冯·诺依曼架构下指令执行的微周期阶段，以及针对多线程并发、多处理器硬件同步机制与 SIMD 向量化执行的体系结构拓展支持。
    \item \textbf{[评估]} 在包含指令级并行（ILP）与流水线冒险（结构冒险、数据冒险与控制冒险）的示例微架构图上进行专业标注，评述现代处理器流水线化解冒险的关键技术手段。
    \item \textbf{[解释]} 阐明机器指令在二进制机器码层级以及符号汇编器语境下的表示与转换映射规则。
    \item \textbf{[应用]} 将常见的高级程序语言控制结构（如 if-else 条件分支、for/while 循环、函数调用与对象成员访问）映射并翻译为等价的汇编/机器指令序列。
    \item \textbf{[评估]} 评述与对比不同指令格式的设计权衡，例如单指令操作数地址数（零地址、一地址、二地址、三地址架构）以及定长指令集（RISC）与变长指令集（CISC）的优劣。
    \item \textbf{[解释]} 结合栈帧剖析图，详细阐述汇编语言层面子程序调用的压栈、参数传递、局部变量分配、返回地址维护与现场恢复的全过程。
    \item \textbf{[解释]} 深刻阐释硬件中断机制、中断向量表（IVT）与 I/O 端口操作的基本工作原理。
    \item \textbf{[开发]} 编写一段高效的汇编语言程序，实现严密底层的字符串遍历、模式查找与文本内存操作。
\end{enumerate}

\subsection{层次存储系统}
\textbf{【学时建议】}：6 计算机核心学时。

本单元是计算机体系结构的核心支柱，深入解析现代多级存储系统如何化解“访存墙”瓶颈：

\subsubsection{计算机核心主题}
\begin{itemize}
    \item \textbf{存储体系原理}：层次存储金字塔，时间局部性与空间局部性的决定性意义（与 OS 交叉）。
    \item \textbf{主存储器组织与操作}：动态随机存取存储器（DRAM）微架构、电容电荷刷新机制、行缓冲、双倍数据速率（DDR）技术与多通道总线配置。
    \item \textbf{持久化与非易失性存储}：固态硬盘（SSD）、NAND Flash 闪存存储原理、磨损均衡与传统机械磁盘（HDD）组织（与 OS 交叉）。
    \item \textbf{存储性能度量}：访存延迟、周期时间、峰值与持续带宽以及多体交叉存取。
    \item \textbf{高速缓存 架构原理}（与 OS 交叉）：
    \begin{itemize}
        \item 地址映射策略：直接映射、全相联与组相联。
        \item 缓存块大小的性能权衡与末级缓存（LLC）设计。
        \item 替换策略：最近最少使用（LRU）、先进先出（FIFO）、伪 LRU（PLRU）与随机替换。
        \item 写策略：写直达与写回，写分配与非写分配。
    \end{itemize}
    \item \textbf{多处理器缓存一致性}：假共享（False Sharing）、总线监听协议、基于目录的协议与 MESI / MOESI 状态机模型。
    \item \textbf{虚拟内存 体系结构支持}：页表结构、多级页表、转译后备缓冲器（TLB 快表）、缺页异常处理流程与内存保护机制（与 OS 交叉）。
    \item \textbf{硬件故障处理与可靠性保证}（与 OS 交叉）。
    \item \textbf{硬件级可靠性技术}（交叉引用 系统基础/基于冗余的可靠性）：
    \begin{itemize}
        \item 差错控制编码：奇偶校验位、ECC 内存校验、海明码与单错校正双错检测（SEC-DED）。
        \item 硬件辅助数据压缩。
        \item 端到端数据完整性验证。
    \end{itemize}
    \item \textbf{非冯·诺依曼架构探索}：
    \begin{itemize}
        \item 存算一体 / 内存计算与近存计算。
    \end{itemize}
\end{itemize}

\subsubsection{示范性学习成果与能力目标}
\begin{enumerate}
    \item \textbf{[解释]} 借助计算机存储层次金字塔图，辨析主流半导体与非易失存储技术（如片上寄存器、SRAM、DRAM、NAND Flash、持久内存）的物理特性、相对每比特成本、容量密度与访存延迟差异。
    \item \textbf{[评估]} 精确测量与量化不同层级存储访问延迟（从 L1 Cache 到主存及持久磁盘）对真实应用软件端到端执行时间的影响。
    \item \textbf{[解释]} 系统阐明具备硬件虚拟内存管理单元（MMU）的现代系统所提供的关键功能（地址空间隔离、动态内存保护、虚拟空间超额分配及按需分页）。
    \item \textbf{[应用]} 熟练推导并计算在多样化 Cache 与内存微架构配置、不同指令与数据引用混合比场景下的平均访存时间。
\end{enumerate}

\subsection{接口与通信}
\textbf{【学时建议】}：1 计算机核心学时。

本单元探讨处理器与外设之间的高速互联、数据交互与通信控制：

\subsubsection{计算机核心主题}
\begin{itemize}
    \item \textbf{输入/输出 (I/O) 基础原理}（与 OS 及 SPD 交叉）：
    \begin{itemize}
        \item 异步握手协议与多级 FIFO 缓冲机制。
        \item 程序控制 I/O（Programmed I/O / 轮询 Polling）。
        \item 中断驱动 I/O（Interrupt-driven I/O）。
    \end{itemize}
    \item \textbf{中断系统微架构}：向量中断与优先级中断仲裁、中断确认握手与多重中断嵌套（与 OS 交叉）。
    \item \textbf{外存储物理组织与通信驱动}：磁盘驱动器物理结构、固态存储控制器通信协议（如 NVMe、SATA）（与 OS 交叉）。
    \item \textbf{系统总线基础技术}（与 OS 及 SPD 交叉）：
    \begin{itemize}
        \item 总线时序协议：同步总线与异步总线。
        \item 总线仲裁机制：集中式仲裁（菊花链、独立请求）与分布式仲裁。
        \item 直接内存访问（DMA: Direct-Memory Access）：DMA 控制器微架构、突发传输模式与周期窃取。
    \end{itemize}
    \item \textbf{片上网络}：多核片上互联拓扑结构（2D Mesh、Torus、Crossbar）、虚通道流控机制与微路由器仲裁。
\end{itemize}

\subsubsection{示范性学习成果与能力目标}
\begin{enumerate}
    \item \textbf{[解释]} 参照中断控制微架构时序图，系统剖析硬件中断如何协调 I/O 外设控制器与 CPU 之间的数据传送与突发事件响应。
    \item \textbf{[解释]} 详细列举现代计算机系统中的各类互联总线与接口标准（如 PCIe、USB、AXI 片上总线），对比其时钟频率、通道宽度与有效带宽特性。
    \item \textbf{[评估]} 系统梳理传统机械硬盘（HDD）的优劣特征，并从寻道延迟、随机读写 IOPS、抗震性、耐用磨损及单位存储成本等维度与固态硬盘（SSD）进行深度对比评估。
\end{enumerate}

\subsection{功能部件组织}
\textbf{【学时建议】}：2 领域核心学时。

本单元深入剖析处理器内部数据通路构建、流水线调度与控制单元设计：

\subsubsection{领域核心主题}
\begin{itemize}
    \item \textbf{基础数据通路 的设计与实现}：ALU 算术逻辑单元、寄存器堆内部总线互联、时钟周期控制、单周期与多周期数据通路划分。
    \item \textbf{控制单元 的设计与实现}：
    \begin{itemize}
        \item 硬布线逻辑实现：组合逻辑门阵列、有限状态机时序生成器。
        \item 微程序控制器实现：控制存储器、微指令格式（水平型微指令与垂直型微指令）、微地址转移逻辑与微解释器。
    \end{itemize}
    \item \textbf{指令流水线技术}：经典 RISC 五级流水线（取指 IF、译码 ID、执行 EX、访存 MEM、写回 WB）及其时空图。
    \item \textbf{流水线冒险及其硬件消解}：
    \begin{itemize}
        \item 结构冒险：硬件资源冲突与多端口存储器解决方案。
        \item 数据冒险：写后读（RAW）、读后写（WAR）、写后写（WAW），旁路转发机制与硬件流水线暂停（Stalling / 流水线气泡 Bubble）。
        \item 控制冒险：分支跳转延迟、分支延迟槽与静态分支预测。
    \end{itemize}
    \item \textbf{指令级并行 导论}：多发射流水线（Superscalar 超标量架构基础）。
\end{itemize}

\subsubsection{示范性学习成果与能力目标}
\begin{enumerate}
    \item \textbf{[评估]} 验证并评述现代计算机微架构中单周期、多周期与经典多级流水线数据通路实现的优缺点与关键时钟频率约束。
    \item \textbf{[开发]} 分别采用硬布线控制器组合逻辑与微程序微指令控制字方案，针对两整型操作数相加并写回寄存器的微操作序列，设计出一组精确完整的硬件控制信号集合。
    \item \textbf{[解释]} 深刻阐释利用流水线发掘指令级并行的基本机理，系统分析运行过程中可能产生的各类严重流水线冒险及其硬件规避对策。
    \item \textbf{[开发]} 独立设计一个包含完整数据通路与控制单元、支持基础 RISC 指令子集的简易可综合处理器微架构。
    \item \textbf{[评估]} 针对给定的处理器流水线架构与层次内存延迟参数，精确推导并计算程序的平均每条指令周期数（CPI）及处理器的实际指令执行吞吐量。
\end{enumerate}

\subsection{性能与能效}
\textbf{【学时建议】}：3 领域核心学时。

本单元深入研究微架构前沿性能挖掘技术、能效优化策略与芯片物理约束：

\subsubsection{领域核心主题}
\begin{itemize}
    \item \textbf{性能与能耗综合评估导论}：计算性能三要素公式（$T = \text{IC} \times \text{CPI} \times \tau$）、动态功耗（$P_{\text{dynamic}} \propto C V^2 f$）与静态漏电功耗、内存访问能耗与片间通信开销。
    \item \textbf{高级微架构推测与乱序执行}：
    \begin{itemize}
        \item 分支预测技术：静态预测、动态 2 位饱和计数器、局部/全局自适应分支预测器与锦标赛预测器。
        \item 推测执行与精确异常处理。
        \item 乱序执行（Out-of-Order Execution）：Tomasulo 算法、保留站、重排序缓冲区 (ROB)与寄存器重命名。
    \end{itemize}
    \item \textbf{硬件预取技术}：指令预取、基于步长的数据预取器。
    \item \textbf{向量处理器与 GPU 架构增强}：单指令多数据（SIMD）向量扩展指令集、单指令多线程（SIMT）模型、张量加速核心。
    \item \textbf{多线程的硬件支持}：
    \begin{itemize}
        \item 细粒度、粗粒度多线程与同步多线程（SMT / 超线程技术）。
        \item 硬件级竞态条件防范。
        \item 硬件锁机制的实现：原子比较并交换（CAS）、加载链接/条件存储（LL/SC）与硬件自旋锁。
        \item 点对点硬件同步与硬件屏障同步机制。
    \end{itemize}
    \item \textbf{体系结构可扩展性定律}：阿姆达尔定律（Amdahl's Law）与古斯塔夫森定律（Gustafson's Law）。
    \item \textbf{非传统替代架构}：超长指令字（VLIW / EPIC 显式并行指令计算）、专用硬件加速器与协处理器。
    \item \textbf{动态电压与频率调节 (DVFS: Dynamic Voltage and Frequency Scaling)}。
    \item \textbf{“暗硅”挑战}：登纳德缩放定律失效后芯片热耗散极限（TDP）与不能全速同时点亮所有晶体管的物理困境。
\end{itemize}

\subsubsection{示范性学习成果与能力目标}
\begin{enumerate}
    \item \textbf{[解释]} 深刻评述衡量现代计算机系统能效与计算性能的核心基准测试与度量指标（如 SPEC 基准套件、MIPS/MFLOPS、每瓦性能比 Performance-per-Watt、能耗延迟积 EDP）。
    \item \textbf{[评估]} 跟踪推测执行与乱序微架构状态转移图，系统撰写并剖析在分支预测正确提交与预测失配重滚刷新的关键判定决策。
    \item \textbf{[开发]} 搭建并绘制针对特定 GPU 硬件架构的性能-功耗基准测试测评体系图。
    \item \textbf{[开发]} 编写一个底层多线程并发程序（如基于 C++ / Python），实现两高维整型向量的并行高效相加，并分析 CPU 核心绑定与访存开销。
    \item \textbf{[开发]} 针对特定应用计算场景（如边缘端低功耗推理或高并发网络包解析），提出一组关于替代性硬件架构的权衡选型方案。
    \item \textbf{[解释]} 深入阐释动态电压与频率调节（DVFS）的核心物理原理、功耗立方/平方降比效应及在现代操作系统调度器中的应用机制。
    \item \textbf{[评估]} 精确评测并量化神经网络模型在采用 8 位整型（INT8）量化相较于传统 32 位浮点（FP32）格式时所获得的内存带宽压缩率与硬件能耗节约幅度。
\end{enumerate}

\subsection{异构架构}
\textbf{【学时建议】}：3 领域核心学时。

本单元面向后摩尔时代前沿，系统解构多类异构加速器与先进微系统集成技术：

\subsubsection{领域核心主题}
\begin{itemize}
    \item \textbf{异构并行体系结构形态}：SIMD 向量处理器、通用图形处理器（GPGPU）、张量处理器（TPU）与神经网络处理器（NPU）微架构解构。
    \item \textbf{异构存储体系结构}：
    \begin{itemize}
        \item 共享内存与分布式内存混合拓扑。
        \item 易失性与非易失性混合持久存储。
        \item 异构缓存一致性协议（CXL, Compute Express Link; CCIX 互联规范与一致性总线）。
    \end{itemize}
    \item \textbf{领域专用架构}：
    \begin{itemize}
        \item 机器学习专用硬件加速器（脉动阵列 Systolic Array、数据流数据通路 Dataflow Datapath）。
        \item 网络内计算：可编程交换芯片（P4 架构）与智能网卡。
        \item 面向新兴前沿应用的低功耗专用嵌入式系统。
        \item 类脑/神经拟态计算架构（Neuromorphic Computing，脉冲神经网络 SNN、存算一体交叉阵列与忆阻器芯片）。
    \end{itemize}
    \item \textbf{先进半导体封装与异构集成方案}：三维集成电路（3DIC）、芯粒技术、硅中介层与通用芯粒互连标准。
\end{itemize}

\subsubsection{示范性学习成果与能力目标}
\begin{enumerate}
    \item \textbf{[评估]} 跟踪并研读包含典型异构并行微架构（如 CPU-SIMD、GPU 流多处理器 SM、NPU 脉动阵列）的现代异构系统集成框图，精准标注其数据流动与控制模型的核心差异。
    \item \textbf{[解释]} 深刻阐述在多处理器与多加速器异构系统中独有而在单处理器单核系统中不存在的存储管理难题（如非一致内存访问 NUMA 效应、跨设备虚拟地址翻译、异构缓存一致性风暴），并说明软硬件协同应对策略。
    \item \textbf{[评估]} 深入论证系统背板总线、处理器-内存直连通道以及基于高速网络的远程直接内存访问（RDMA）在访存延迟与吞吐量上的阶跃差异及其对分布式并行程序性能的直接制约。
    \item \textbf{[评估]} 针对给定的实际算法负载特征（如计算密度比 Arithmetic Intensity、访存模式规则性与控制流分支发散度），科学论证在何种场景下应当引入领域专用加速器（DSA）而非依赖通用多核 CPU。
    \item \textbf{[解释]} 系统剖析并阐述向量处理单元与标量处理单元在寄存器组织、指令流水与内存 Gather/Scatter 机制上的体系结构设计哲学差异。
    \item \textbf{[评估]} 全面评估存算一体（PIM）体系结构的工程优缺点：系统权衡其在粉碎“冯·诺依曼访存墙”、降低数据搬运能耗等方面的巨大收益，以及在模拟计算噪声、热耗散集中、制造良率与现有编译软件栈兼容性方面面临的技术挑战。
\end{enumerate}

\subsection{量子体系结构}
\textbf{【学时建议】}：3 领域核心学时。

本单元面向未来颠覆性计算前沿，构建量子物理基础与量子微架构体系的认知桥梁：

\subsubsection{领域核心主题}
\begin{itemize}
    \item \textbf{量子计算物理基础原理}：
    \begin{itemize}
        \item 微观粒子的波粒二象性原理。
        \item 双缝干涉实验中的海森堡测不准/不确定性原理。
        \item 什么是量子比特？量子态叠加与测量坍缩；以单光子作为物理量子比特的实现。
        \item 双量子比特系统：量子纠缠、贝尔基态与无超光速通信定理。
        \item 量子力学公理系统：态矢量叠加公理、投影测量公理、闭系统酉演化变换。
        \item 量子线路模型的单量子比特逻辑门：Pauli-X 门（非门）、Pauli-Z 门（相位翻转门）、Hadamard 门（H 门，创建叠加态）。
        \item 双量子比特受控门（如 CNOT 门）与张量积运算；量子门矩阵代数表示。
        \item 量子不可克隆定理；量子隐形传态协议。
    \end{itemize}
    \item \textbf{经典与基础量子算法}：
    \begin{itemize}
        \item 基础预言机算法：Bernstein-Vazirani 算法、Simon 周期查找算法。
        \item 利用马赫-曾德尔干涉仪光学实现 Deutsch-Jozsa 算法。
        \item 量子质因数分解算法（Shor's Algorithm）核心流程。
        \item 量子无序数据库加速搜索算法（Grover's Algorithm）。
    \end{itemize}
    \item \textbf{量子硬件实现与微架构工程}：
    \begin{itemize}
        \item 量子比特的主流物理实现路线（当前主流九大物理实现载体：超导量子电路、捕获离子阱、半导体量子点硅自旋、光量子芯片、中性冷原子阵列、拓扑准粒子、金刚石氮-空位色心 NV、液氦表面电子等）。
        \item 量子处理单元（QPU）的经典微电子控制架构（低温同轴线缆、微波脉冲生成、高精度 DAC/ADC 控制逻辑与实时反馈回路）。
        \item 量子误差缓解与量子纠错：表面码、含噪声中等规模量子（NISQ）时代的现实局限与容错通用量子计算远景。
    \end{itemize}
    \item \textbf{新兴前沿应用方向}：
    \begin{itemize}
        \item 抗量子后量子密码学（Post-quantum cryptography / PQC，如基于格密码的算法过渡）。
        \item 量子互联网（The Quantum Internet：量子中继器、长距离纠缠分发与量子密钥分发 QKD）。
        \item 绝热量子计算 (AQC)与量子退火。
    \end{itemize}
\end{itemize}

\subsubsection{示范性学习成果与能力目标}
\begin{enumerate}
    \item \textbf{[解释]} 深入阐释微观量子物体的波粒二象性运动图景：以粒子形态产生，以波函数形式在空间中传播演化，并在测量探测时以遵从波函数模平方概率分布的粒子形态呈现。
    \item \textbf{[解释]} 借助实验时序图，深入评述微观量子系统的内禀概率本质与观测介入导致的系统态坍缩特性。
    \item \textbf{[评估]} 严谨表述对量子纠缠的科学认知：纠缠可用于建立非定域的强非经典关联，但依据“无信号定理”，物理上绝对无法利用量子纠缠进行超光速的有效信息通信。
    \item \textbf{[解释]} 深刻阐明量子并行性的运行机理，并针对量子测量的概率特性，分析相长干涉与相消干涉在筛选与放大正确答案概率中的核心设计机制。
    \item \textbf{[应用]} 在给定的 Shor 算法微代码或线路图上，准确标注量子傅里叶采样（QFT）在提取量子叠加态相位与模周期中的关键运算角色。
    \item \textbf{[开发]} 在开源量子计算仿真平台（如 Qiskit / Cirq）中编写 Shor 算法仿真程序，并撰写工程报告，深入辨析经典计算模块（数论模幂规约、连分数展开）与量子加速核心子模块在算法整体中的协同边界。
    \item \textbf{[解释]} 系统梳理并详细阐明当前各主流量子比特实现路线（如超导、离子阱、硅自旋、光量子、中性原子、拓扑等）的物理工作机制、相干保真度指标与可扩展微加工壁垒。
    \item \textbf{[评估]} 深入辨析绝热量子计算（AQC）/量子退火与基于逻辑门的通用线路模型量子计算在硬件体系结构、控制机制与适用计算问题类别（如组合优化与退火基态搜索 vs 通用代数运算）上的本质差异。
    \item \textbf{[评估]} 深刻论证与评述以下关键体系结构命题：“在现代异构计算系统的宏观视角下，量子处理单元（QPU）本质上是一种由经典宿主计算机通过特定总线进行控制与卸载计算的异构专用协处理器微架构（类似于 FPGA 或 GPU）”。
\end{enumerate}

\section{专业素养与思维倾向}

在计算机体系结构与组织的研习与工程实践中，重点塑造学生以下维度的专业心智素养与职业特质：
\begin{itemize}
    \item \textbf{自主求知与深度钻研}：具备日益增强的自主求知欲与主动性，能自觉主动研读繁复的芯片体系结构白皮书、底层汇编规范、微架构参考手册以及工业级硬件设计文档，持续拓宽软硬件交叉知识边界。
    \item \textbf{前瞻洞察与跨界融通}：展现积极主动且独立探索的工程素养，能跨越传统学科界限，无缝整合体系结构、编译优化、操作系统调度与算法分析等多领域知识，精准洞察与驾驭计算机系统的底层运行逻辑。
    \item \textbf{突破创新的架构思维}：勇于超越平凡直觉与常规解决方案，面对苛刻的硬件算力、带宽瓶颈与功耗墙约束，能因地制宜、创造性地发掘并利用底层硬件专有微架构特性（如 SIMD 向量指令、片上共享存储、异步 DMA 流水等）以达成极致能效。
    \item \textbf{严谨敬业与工程伦理操守}：秉持高度的职业审慎态度与道德责任感。现代计算机系统（特别是直接贴附或植入人体的生物传感器芯片，如实时血糖监测与心脏起搏设备）与人身安全直接相连，工程师必须在硬件芯片与微架构底层将用户隐私、物理安全性与系统可靠性置于至高无上的优先级。
\end{itemize}

\section{数学基础要求}

计算机体系结构领域对离散数学、线性代数及统计物理有着多层次的数学依赖：

\subsection{必须掌握的核心数学}
\begin{itemize}
    \item \textbf{离散数学}：集合论、二元关系、布尔逻辑与谓词运算（用于数字逻辑综合与状态机设计）、初等数论（同余定理、原根与模算术，用于计算机算术、差错控制编码与硬件密码加速）。
    \item \textbf{线性代数}：算术基础运算、矩阵乘法与多维张量操作（用于 GPU/TPU 脉动阵列加速器微架构设计）、矩阵特征值与酉算符分解（用于量子逻辑门矩阵代数）。
\end{itemize}

\subsection{进阶推荐数学背景}
\begin{itemize}
    \item \textbf{面向量子计算的数理基础}：基础概率论与随机变量、三角函数、复数向量空间与复数内积运算、欧拉公式（$e^{i\theta} = \cos\theta + i\sin\theta$）。
    \item \textbf{系统性能评估与试验设计方法}：数理统计基础、正态分布与置信区间、析因实验设计与方差分析（用于多微架构参数调优、性能瓶颈敏感度消融实验）。
\end{itemize}

\section{交叉主题与跨学科共享概念}

\subsection{跨学科共享概念}
\begin{itemize}
    \item \textbf{虚拟内存}：体系结构/层次存储系统 与 操作系统/内存管理（操作系统/内存管理）深度共享。
    \item \textbf{差错控制编码、数据压缩与数据完整性}：体系结构/层次存储系统 与 系统基础/Software Reliability through Redundancy（系统基础/基于冗余的可靠性）深度共享。
    \item \textbf{通信互联与片上网络}：体系结构/接口与通信 与 网络通信/网络导论与体系结构（网络与通信导论）深度共享。
    \item \textbf{处理器微架构设计}：体系结构/异构架构 与 专用平台/移动平台（移动平台）及 专用平台/Industrial（工业平台）交叉。
    \item \textbf{领域专用加速架构}：体系结构/异构架构 与 人工智能/Advanced Machine Learning（高级机器学习）交叉。
    \item \textbf{新兴硬件用户接口}：体系结构/异构架构 与 人机交互/Human Factors \& Security（人机交互/人类因素与安全）交叉。
    \item \textbf{底层并行计算系统}：体系结构/汇编级机器组织 与 并行分布式/并行体系结构（并行计算/并行体系结构）交叉。
\end{itemize}

\subsection{跨学科横切主题}
\begin{itemize}
    \item \textbf{多层级并行性}：指令级并行（ILP）、数据级并行（SIMD/向量化）、线程级并行（TLP）与请求级并行（RLP）。
    \item \textbf{系统资源分配、硬件调度与空间隔离}。
    \item \textbf{性能与能源效率的永恒权衡}。
    \item \textbf{底层计算系统的社会与生态环境影响}：数据中心碳排放、高能耗芯片制冷挑战、暗硅现象与电子硬件生命周期管理。
\end{itemize}

\section{课程设计与开课规划建议}

联合工作组针对不同层次的培养目标推荐两套典型的开课规划方案：

\subsection{方案一：计算机体系结构入门课程 (建议约 42 学时)}
本方案面向本科通用计算机专业，构建完整的软硬件交界面认知：
\begin{itemize}
    \item 社会伦理/计算机发展简史 [2 学时]
    \item 体系结构/机器级数据表示 [2 学时]
    \item 体系结构/汇编级机器组织 [2 学时]
    \item 体系结构/层次存储系统 [10 学时]
    \item 操作系统/内存管理 [10 学时]
    \item 体系结构/接口与通信 [4 学时]
    \item 体系结构/异构架构 [5 学时]
    \item 并行分布式/程序与执行 [4 学时]
    \item 社会伦理/伦理分析方法 [3 学时]
\end{itemize}
\textbf{【先修要求】}：离散数学（集合、二元关系、逻辑运算、初等数论）。\\
\textbf{【能力目标】}：完成本课程的学生应能全面理解现代计算机系统的底层基础架构，深刻领会 Cache 高速缓存、流水线冲突以及硬件内存管理机制所面临的系统级挑战。

\subsection{方案二：高级计算机系统综合课程 (建议约 38 学时)}
本方案面向计算机系统方向的高年级拔尖创新人才与研究型培养：
\begin{itemize}
    \item 社会伦理/计算机发展简史 [2 学时]
    \item 社会伦理/隐私保护 [3 学时]
    \item 软件工程/软件架构与设计 [4 学时]
    \item 并行分布式/并行算法与应用 [4 学时]
    \item 体系结构/异构架构 [4 学时]
    \item 操作系统/角色与定位 [3 学时]
    \item 体系结构/层次存储系统 [7 学时]
    \item 体系结构/性能与能效 [5 学时]
    \item 网络通信/网络化应用与套接字编程 [5 学时]
\end{itemize}
\textbf{【先修要求】}：离散数学（集合、二元关系、逻辑运算、初等数论）。\\
\textbf{【能力目标】}：完成本课程的学生应能系统洞察现代计算硬件体系结构的核心设计哲学，具备在后摩尔时代驾驭硬件并行性与异构算力、持续提升系统性能与能效的高维系统工程素养。

\section{真实能力任务规范}

\begin{taskbox}[任务 1：应用程序中的高速缓存性能瓶颈评估与优化]
\textbf{【任务情境】}：某企业核心高频业务系统的延迟指标出现严重长尾波动，开发团队怀疑是由底层的 Cache 缺失与颠簸导致。\\
\textbf{【能力陈述】}：批判性地剖析应用程序在高速缓存层面的访存性能瓶颈（命中率、行冲突、假共享），运用微架构性能计数器（PMU）进行量化剖析，并撰写总结核心瓶颈成因与重构建议的权威评估报告。\\
\textbf{【能力领域】}：系统 / 应用 \quad \textbf{【能力单元】}：系统开发 / 部署 / 评估 / 改进\\
\textbf{【关联知识单元】}：体系结构/层次存储系统、体系结构/性能与能效、操作系统/内存管理\\
\textbf{【达标等级】}：\textbf{评估}
\end{taskbox}

\begin{taskbox}[任务 2：边缘嵌入式设备上机器学习模型部署的“每瓦性能比”评估]
\textbf{【任务情境】}：在电池供电的低功耗移动机器人上部署云端训练完成的目标检测神经网络模型，需在帧率与续航之间达成最优平衡。\\
\textbf{【能力陈述】}：在目标嵌入式硬件设备上实测并评估运行轻量化机器学习模型的端到端推断吞吐量、芯片动态功耗与热耗散表现，计算其“每瓦性能比”（Performance-per-Watt），指导模型量化与微架构剪枝策略。\\
\textbf{【能力领域】}：系统 / 应用 \quad \textbf{【能力单元】}：部署 / 评估 / 改进\\
\textbf{【关联知识单元】}：体系结构/性能与能效、专用平台/移动平台\\
\textbf{【达标等级】}：\textbf{评估与开发}
\end{taskbox}

\begin{taskbox}[任务 3：将 CPU 应用程序重构移植至专用硬件加速器]
\textbf{【任务情境】}：传统在 x86 CPU 上运行的大规模科学模拟或图像处理程序算力不足，需将其高性能关键计算循环迁移至现代硬件加速器上。\\
\textbf{【能力陈述】}：综合运用系统架构设计知识，将通用 CPU 应用程序的核心密集运算代码重构并移植至专用硬件加速器（GPU、TPU 或 NPU），合理规划设备端内存分配与主机-设备数据传输流水线，并精确评估计算加速比。\\
\textbf{【能力领域】}：软件 / 系统 \quad \textbf{【能力单元】}：需求 / 设计 / 开发 / 测试 / 技术文档\\
\textbf{【关联知识单元】}：体系结构/异构架构、并行分布式/并行体系结构、系统基础/系统设计\\
\textbf{【达标等级】}：\textbf{评估与开发}
\end{taskbox}

\begin{taskbox}[任务 4：构建性能基准测试工具以诊断并消除程序 I/O 吞吐瓶颈]
\textbf{【任务情境】}：大规模数据处理程序在吞吐量扩展时遭遇未知物理阻滞，亟需确定系统瓶颈究其在于 CPU 计算能力还是外部 I/O 通信通路。\\
\textbf{【能力陈述】}：开发或熟练配置一套系统级性能基准测试与性能剖析工具，实时监测磁盘总线、网络套接字与内存 DMA 数据流，量化识别 I/O 阻塞节点，并评估通过消除 I/O 瓶颈所获得的系统级性能加速收益。\\
\textbf{【能力领域】}：软件 / 系统 \quad \textbf{【能力单元】}：开发 / 评估 / 系统维护\\
\textbf{【关联知识单元】}：体系结构/接口与通信、体系结构/性能与能效、系统基础/资源分配与调度、系统基础/系统性能\\
\textbf{【达标等级】}：\textbf{评估与解释}
\end{taskbox}

\begin{taskbox}[任务 5：评估 CPU-GPU 异构统一内存管理中的缺页异常及其性能损耗]
\textbf{【任务情境】}：在现代异构统一虚拟内存架构中，过多的主机与加速器间隐式页面迁移导致应用程序性能剧烈退化。\\
\textbf{【能力陈述】}：运用操作系统内存管理与计算机体系结构原理，深入剖析 CPU 与 GPU 之间统一虚拟内存的缺页异常触发机制、页面迁移开销与 PCIe/NVLink 总线拥堵效应，并提出显式预取与提示优化方案。\\
\textbf{【能力领域】}：系统 \quad \textbf{【能力单元】}：系统设计 / 开发 / 评估 / 改进\\
\textbf{【关联知识单元】}：体系结构/层次存储系统、体系结构/性能与能效、体系结构/异构架构、操作系统/内存管理\\
\textbf{【达标等级】}：\textbf{应用与评估}
\end{taskbox}

\begin{taskbox}[任务 6：连续体征监测医疗设备在硬件与系统级的数据隐私与安全影响评估]
\textbf{【任务情境】}：研发一款贴附于患者身体、具备高频连续生理体征数据采集能力的穿戴式医疗设备，需规避硬件侧信道泄露与隐私窃取。\\
\textbf{【能力陈述】}：应用底层计算系统的专业技术知识，权衡嵌入式传感芯片的数据采集精度、片上数据最小化处理策略与硬件安全信任根机制，撰写关于医疗设备硬件与固件级数据隐私合规的技术白皮书。\\
\textbf{【能力领域】}：系统 / 应用 \quad \textbf{【能力单元】}：需求分析 / 系统设计 / 技术文档 / 社会议题调适\\
\textbf{【关联知识单元】}：体系结构/异构架构、社会伦理/隐私保护、编程语言/嵌入式计算与硬件接口\\
\textbf{【达标等级】}：\textbf{评估}
\end{taskbox}

\begin{taskbox}[任务 7：传感器硬件集成的软硬件协同驱动接口与模块设计]
\textbf{【任务情境】}：为新型高精度工业物联网传感硬件开发低延迟、高可靠的板级支持包（BSP）与操作系统驱动框架。\\
\textbf{【能力陈述】}：设计并实现面向传感器芯片的高层软件抽象 API 以及底层基于中断与总线寄存器的硬件交互驱动模块，达成无缝的软硬件系统集成与容错处理。\\
\textbf{【能力领域】}：系统 \quad \textbf{【能力单元】}：需求工程 / 架构设计 / 模块开发\\
\textbf{【关联知识单元】}：体系结构/接口与通信、体系结构/异构架构、操作系统/文件系统 API 与底层实现、软件开发基础/基本编程概念、软件工程/API 设计原则、专用平台/嵌入式系统\\
\textbf{【达标等级】}：\textbf{开发与评估}
\end{taskbox}

\begin{taskbox}[任务 8：系统软硬件架构选型论证与技术方案设计规范撰写]
\textbf{【任务情境】}：面向高可靠严苛工控领域，企业拟研发下一代综合计算平台，需全面论证处理器微架构选型与系统拓扑方案。\\
\textbf{【能力陈述】}：全面收集梳理复杂应用场景的计算需求、功耗预算与散热条件，权衡通用处理器与专用加速芯片配置，提出系统软硬件体系结构综合设计方案，并撰写高质量的技术决策规范文档。\\
\textbf{【能力领域】}：系统 / 应用 \quad \textbf{【能力单元】}：需求工程 / 架构设计 / 技术文档\\
\textbf{【关联知识单元】}：体系结构/机器级数据表示、体系结构/性能与能效、体系结构/异构架构、社会伦理/专业沟通与表达、软件工程/需求工程与产品规约\\
\textbf{【达标等级】}：\textbf{评估与解释}
\end{taskbox}

\begin{taskbox}[任务 9：自动化制造单元中高精度感知-执行机械臂系统的跨学科协同研发]
\textbf{【任务情境】}：在现代智能制造柔性装配线中，跨学科研发团队正在攻坚具备毫秒级视觉反馈与高精度伺服控制的智能协作机械臂系统。\\
\textbf{【能力陈述】}：在多学科复合型研发团队中高效沟通协作，负责机械臂底层嵌入式微控制器、总线通信协议与实时加速芯片的软硬件联合设计、集成调试与系统性能验证。\\
\textbf{【能力领域】}：软件 / 系统 / 应用 \quad \textbf{【能力单元】}：需求 / 设计 / 开发 / 测试 / 集成 / 评估\\
\textbf{【关联知识单元】}：体系结构/异构架构、专用平台/机器人平台、软件工程/团队协作\\
\textbf{【达标等级】}：\textbf{应用与开发}
\end{taskbox}

\begin{taskbox}[任务 10：深度利用底层多核 CPU 与层次存储系统的高性能并发程序开发]
\textbf{【任务情境】}：针对金融量化高频交易或超大规模网格计算，需彻底压榨现代多核多级缓存服务器硬件的每一周期算力。\\
\textbf{【能力陈述】}：透彻理解现代 CPU 核心流水线、乱序执行、SIMD 向量扩展与层次内存架构，设计并编写底层高性能并发程序，消除伪共享、优化缓存局部性、精细化调度多核线程以达成极致吞吐。\\
\textbf{【能力领域】}：系统 \quad \textbf{【能力单元】}：需求 / 设计 / 开发 / 性能调优\\
\textbf{【关联知识单元】}：体系结构/汇编级机器组织、体系结构/层次存储系统\\
\textbf{【达标等级】}：\textbf{评估与开发}
\end{taskbox}

\section{指导委员会成员名单}

\begin{itemize}
    \item \textbf{主席}：Marcelo Pias，巴西里奥格兰德联邦大学, Brazil
    \item \textbf{工作组成员}：
    \begin{itemize}
        \item Brett A. Becker，爱尔兰都柏林大学学院, Ireland
        \item Mohamed Zahran，美国纽约大学, USA
        \item Monica D. Anderson，美国阿拉巴马大学, USA
        \item 向乔，厦门大学
        \item Adrian German，美国印第安纳大学布卢明顿分校, USA
    \end{itemize}
\end{itemize}

\section{附录：核心主题与技能要求等级对照表}

\begin{longtable}{p{0.18\textwidth} p{0.48\textwidth} c c c}
\toprule
\textbf{知识单元} & \textbf{核心知识点与具体主题} & \textbf{认知等级} & \textbf{分类} & \textbf{学时} \\
\midrule
\endfirsthead
\toprule
\textbf{知识单元} & \textbf{核心知识点与具体主题} & \textbf{认知等级} & \textbf{分类} & \textbf{学时} \\
\midrule
\endhead
\bottomrule
\endfoot
\multirow{2}{*}{\makecell[l]{体系结构/数字逻辑与\\数字系统}} 
& 体系结构演进历史；组合逻辑与时序逻辑基本构建块；硬件与软件多层抽象体系结构 & 解释 & 领域核心 & \multirow{2}{*}{3} \\
\cmidrule{2-4}
& 硬件 CAD 工具链；高层次综合（RTL、Verilog/VHDL/Chisel）；SoC 设计流程；物理与电气约束（延迟、扇入扇出、功耗、光速） & 评估 & 领域核心 & \\
\midrule
\makecell[l]{体系结构/机器级\\数据表示} 
& 位、字节与字；定点数与浮点数（IEEE 754）；原码与补码运算及溢出；非数值字符编码；记录/结构体对齐与多维数组内存排布 & 应用 & 计算机核心 & 1 \\
\midrule
\multirow{2}{*}{\makecell[l]{体系结构/汇编级\\机器组织}} 
& 经典冯·诺依曼架构；控制单元与指令周期；SIMD 与 MIMD 导论及 Flynn 分类；共享内存多核多处理器组织 & 解释 & 计算机核心 & 1 \\
\cmidrule{2-5}
& 指令集体系结构（ISA：x86/ARM/RISC-V）；寻址模式；机器码与汇编程序设计；子程序调用机制与栈帧；I/O 与中断；内存段划分 & 开发 & 领域核心 & 2 \\
\midrule
\multirow{2}{*}{\makecell[l]{体系结构/层次存储系统}} 
& 时间与空间局部性；主存组织（DRAM/DDR）；持久化存储（SSD/HDD）；访存延迟、周期、带宽与交叉存取；虚拟内存支持；硬件可靠性与纠错码；存算一体 & 解释 & 计算机核心 & \multirow{2}{*}{6} \\
\cmidrule{2-4}
& 高速缓存 Cache（地址映射、块大小、替换与写策略）；多处理器缓存一致性（监听/目录协议、MESI/MOESI 状态机） & 评估 & 计算机核心 & \\
\midrule
\makecell[l]{体系结构/接口与通信} 
& I/O 基础原理（握手、缓冲、轮询、中断驱动）；向量与优先级中断结构；外存储物理组织；系统总线协议、仲裁与 DMA；片上网络 & 解释 & 计算机核心 & 1 \\
\midrule
\multirow{2}{*}{\makecell[l]{体系结构/功能部件组织}} 
& 数据通路设计；控制单元实现（硬布线逻辑与微程序控制器）；五级指令流水线实现；流水线冒险（结构、数据、控制冒险）检测与消解 & 开发 & 领域核心 & \multirow{2}{*}{2} \\
\cmidrule{2-4}
& 深入指令流水线（时空图）；指令级并行（ILP）核心概念与超标量架构基础 & 解释 & 领域核心 & \\
\midrule
\multirow{2}{*}{\makecell[l]{体系结构/性能与能效}} 
& 性能与能耗综合评估导论；分支预测技术；推测执行；乱序执行与 Tomasulo 算法；硬件预取技术 & 评估 & 领域核心 & 2 \\
\cmidrule{2-5}
& 向量处理器与 GPU 增强；多线程硬件支持（硬件锁、同步与屏障）；可扩展性定律；VLIW 替代架构；DVFS 动态能效调节；暗硅挑战 & 解释 & 领域核心 & 1 \\
\midrule
\makecell[l]{体系结构/异构架构} 
& SIMD/MIMD 异构架构；异构存储（共享 vs 分布式、CXL 缓存一致性）；领域专用架构 DSA（机器学习加速器、网络内计算、类脑计算）；先进封装（3DIC 与 Chiplets 芯粒） & 解释 & 领域核心 & 3 \\
\midrule
\makecell[l]{体系结构/量子体系结构} 
& 量子物理原理（波粒二象性、测不准原理、Qubit、叠加态、纠缠、公理体系、单/双量子比特门、不可克隆定理）；量子算法（Bernstein-Vazirani、Simon、Shor、Grover）；硬件实现（九大物理路线、经典控制、NISQ）；新兴应用（抗量子密码 PQC、量子互联网、量子退火） & 解释 & 领域核心 & 3 \\
\bottomrule
\end{longtable}

\chapter{人工智能}

\section{领域导言与演进背景}

\begin{preamblebox}[知识领域概述与变革]
人工智能 (AI)致力于研究与解决那些难以或无法通过传统算法有效求解的计算问题。这类问题往往需要模拟人类智能，其解决方案通常能够跨类别泛化至更广泛的同构问题。如今，人工智能技术在计算领域无处不在，深度赋能电子邮件、社交媒体、智能摄影、金融量化交易以及智能虚拟语音助理（如 Siri、Alexa 等）等日常生活应用。此外，这些技术也被广泛应用于自主智能体的设计与分析中，使其能够自主感知外部环境并进行理性交互，例如自动驾驶车辆与各类自主移动机器人。

从传统上看，人工智能融合了符号学派与亚符号/联结学派的方法。其求解机制依赖于一套兼具通用性与领域专用性的知识表示体系、问题求解机制与数值优化技术。这些技术涵盖了感知（如语音识别、自然语言理解、计算机视觉）、问题求解（如启发式搜索、自动规划、约束优化）、动作执行（如机器人学、任务自动化、智能控制），以及支撑上述能力的系统架构（如单智能体架构与多智能体系统）。学习人工智能旨在培养学生具备以下专业能力：针对特定现实问题准确研判是否适用 AI 方法；科学选择与设计适宜的知识表示与推理机制；高质量实现工程原型；并在算法计算性能与更广泛的社会影响两个维度上进行严谨的系统评估。

在过去十年中，在一系列高影响力的机器学习重大突破推动下，“人工智能”一词已在企业界、新闻媒体与日常话语中广泛普及。这些跨越式进展得益于海量大数据的可用性、计算算力（如 GPU/TPU）的迅猛飞跃以及算法理论的持续演进。特别显著的是，人工智能的研究范式发生了重大转向——从以往高度依赖人工工程化设计的特征与知识表示，转向通过大规模数据集优化自适应学习得到的特征表示。由此催生的技术变革使“神经网络”和“深度学习”成为大众熟知的流行语。企业纷纷将基于 AI 的技术方案作为核心服务增值亮点进行推广，使得“人工智能”既是一个严谨的技术学术术语，也是一个备受瞩目的商业营销关键词。与此同时，生物学、艺术创作、建筑设计和金融工程等跨学科领域也日益深度引入 AI 技术以求解其学科前沿问题。人类历史上首次迎来广泛大众能够直接接触并使用复杂精尖的 AI 驱动工具的新时代，包括根据提示词自动生成长篇论文或优美诗歌的生成式工具、依据文本描述生成高品质摄影或艺术图像的扩散工具，以及能够极其逼真地合成特定人物图像与音视频的深度伪造技术。AI 技术现已深度渗透至金融股票交易、新闻与社交网络信息流推荐算法、招聘求职者自动化多维筛选、临床医疗病症辅助诊断，乃至通过再犯预测算法深度介入司法量刑判决。因此，人工智能技术具备极为深远的社会影响力，在研发与落地应用的全生命周期中必须对其潜在风险进行审慎考量与科学治理。

\textbf{自 CS2013 以来的核心变革：}
\begin{enumerate}
    \item \textbf{领域命名规范化更替}：将 CS2013 中的“智能系统”正式更名为“人工智能”，以精准呼应学科内部最通用的学术规范及全社会最广泛的行业共识。
    \item \textbf{强化神经网络与表示学习，平衡符号与联结范式}：顺应深度学习与表示学习的技术浪潮，大幅提升了神经网络与特征学习的教学权重。图搜索作为 AI 的经典基石依然得到坚守，但适度精简了部分传统符号演算方法，以腾出学时深入剖析亚符号方法与学习型表征。与此同时，课程体系坚决保留基于知识的系统与符号推理，因为符号方法具有独到的可解释性与严密逻辑推理能力，在工程实践中依然发挥关键作用，且最新前沿“神经符号人工智能”正在将学习表征与符号知识深度融合。
    \item \textbf{突出跨领域工程落地与深度生成式模型}：全面深化 AI 在医疗健康、环境可持续发展、社交网络治理等多元垂直领域的实战应用。特别将以 ChatGPT、Bard/Gemini 为代表的大语言模型（LLM）以及以 DALL-E、Midjourney、Stable Diffusion 为代表的深度图像扩散生成模型纳入核心教学视野，系统讲授其宏观技术脉络、落地场景、幻觉缺陷与潜在风险。
    \item \textbf{全面融入社会影响、伦理道德与可解释性}：深刻体现理解并客观评估 AI 技术与应用对人类社会深远影响的重要性，将人工智能伦理、算法公平性、公众信任度与可解释性作为核心必修议题深度织入各知识模块。
    \item \textbf{构筑与数据科学及数据管理的紧密协同链接}：强化人工智能知识领域与数据管理（DM）及数据科学基础设施的跨学科交叉映射。
    \item \textbf{确立面向所有计算专业学生的人工智能素养与批判性思维目标}：鉴于 AI 技术已广泛渗透至计算机科学的各行各业，明确将培养每位计算机科学专业学生的基础 AI 素养与批判性思辨能力作为核心培养目标。
\end{enumerate}
\end{preamblebox}

\section{核心学时分配建议}

本知识领域在“向日葵模型”架构下的学时推荐如下表所示：

\begin{table}[H]
\centering
\setlength{\tabcolsep}{3.5pt}\fontsize{8pt}{9.6pt}\selectfont\begin{tabularx}{\columnwidth}{Y c c}
\toprule
\textbf{知识单元} & \textbf{计算机核心} & \textbf{领域核心} \\
\midrule
基础议题 & 2 & 1 \\
搜索算法与问题求解 & 2 (+3 AL)$^*$ & 4 \\
基础知识表示与推理 & 1 (+1 MSF)$^\dagger$ & 2 \\
机器学习 & 4 & 4 \\
应用与社会影响 & 2 & 2 \\
逻辑表示与推理 & -- & 选修 \\
概率表示与推理 & -- & 选修 \\
规划系统 & -- & 选修 \\
智能体系统 & -- & 选修 \\
自然语言处理 & -- & 选修 \\
机器人学 & -- & 选修 \\
感知与计算机视觉 & -- & 选修 \\
\midrule
\textbf{合计} & \textbf{11 (+4 跨领域共享)$^\ddagger$ 学时} & \textbf{13 学时} \\
\bottomrule
\end{tabularx}
\begin{flushleft}\footnotesize
$^*$注：搜索算法单元包含 5 个计算机核心学时，其中 3 个学时（无信息图搜索部分）与算法基础（算法/基础数据结构与算法）交叉共享。\\
$^\dagger$注：基础知识表示与推理单元包含 2 个计算机核心学时，其中 1 个学时（贝叶斯概率推理基础）与数学与统计学基础交叉共享。\\
$^\ddagger$注：计算机核心总学时为 11 学时（由 AI 独立承担），加上由 AL 和 MSF 共享承载的 4 学时，该领域覆盖的计算机核心内容总量共计 15 学时。
\end{flushleft}
\end{table}

\section{知识单元与详细大纲}

\subsection{人工智能基础议题}
\textbf{【学时建议】}：2 计算机核心学时，1 领域核心学时，包含选修主题。

本单元系统阐述人工智能的核心概念、发展全景与基本哲学命题，界定智能行为与智能体的基本内涵，为学生构建起涵盖问题建模、环境属性分类与社会影响洞察的全局认知视野。

\subsubsection{计算机核心主题}
\begin{itemize}
    \item \textbf{人工智能全貌与前沿进展}：人工智能求解问题的本质范式；近期标志性成功工程应用案例（如大语言模型、视觉生成模型、AlphaFold 蛋白质结构预测、高阶自动驾驶等）。
    \item \textbf{智能体的概念与典型范式}：智能体的形式化定义；反应式智能体与深思式智能体的核心特征与对比。
    \item \textbf{智能行为的本质思考}：什么是智能行为？智能行为的计算模拟途径。
    \item \textbf{图灵测试与心智思辨}：图灵测试的提出背景、测试协议及其理论缺陷；经典哲学反思与认知争议。
    \item \textbf{多模态交互}：多模态输入与输出（文本、视觉、语音、传感器时序等维度的融合感知与多模态生成）。
    \item \textbf{推理范式辨析}：理性推理与非理性/类人推理的本质边界与目标差异。
    \item \textbf{问题与环境特征分类体系}：
    \begin{itemize}
        \item 完全可观察 vs. 部分可观察；
        \item 单智能体 vs. 多智能体；
        \item 确定性 vs. 随机性；
        \item 静态 vs. 动态；
        \item 离散 vs. 连续。
    \end{itemize}
    \item \textbf{智能体本质与自主性特征}：
    \begin{itemize}
        \item 完全自主、半自主与混合主动自主；
        \item 条件反射式、基于目标与基于效用的智能体架构；
        \item 不确定性环境与不完全信息条件下的决策制定机制；
        \item 环境感知与动态物理交互的核心重要性；
        \item 基于学习的智能体演进机制；
        \item 具身智能体核心三要素：传感器、系统动力学与执行器。
    \end{itemize}
    \item \textbf{人工智能的应用扩张与社会冲击}：AI 技术的爆发式增长及其在经济结构、社会就业与伦理治理层面的广泛影响。
\end{itemize}

\subsubsection{领域核心主题}
\begin{itemize}
    \item \textbf{环境与问题特征深度剖析}：结合复杂现实工程实例（如高频量化交易、开放道路自动驾驶、智能电网调度），深入解析各类环境维度的复合交叉特征。
    \item \textbf{智能体系统架构进阶}：结合典型工业级系统，深入分析分层控制、意图驱动与效用函数在复杂智能体中的架构实现。
    \item \textbf{人工智能深远影响专题研讨}：全面评估 AI 产业应用带来的宏观经济效益、劳动力结构重塑、数据隐私边界与伦理合规挑战。
\end{itemize}

\subsubsection{选修主题}
\begin{itemize}
    \item \textbf{人工智能哲学与认知科学前沿}：强人工智能与弱人工智能的思辨；功能主义、机器意识与意向性理论；机器伦理学与机器权利哲学探讨。
\end{itemize}

\subsubsection{示范性学习成果与能力目标}
\begin{enumerate}
    \item \textbf{阐释}图灵测试的基本协议机制，并结合约翰·塞尔的“中文屋”思想实验，深入剖析形式符号操作与真正语义理解之间的本质鸿沟。
    \item \textbf{辨析}最优理性推理/行为与拟人化推理/行为的根本差异，能够依据特定业务需求确定系统设计应追求的目标范式。
    \item \textbf{界定}任意现实工程问题与应用场景在可观察性、随机性、动态性、离散/连续性以及智能体协作模式等维度上的核心环境特征。
\end{enumerate}

\subsection{搜索算法与问题求解}
\textbf{【学时建议】}：5 计算机核心学时（其中 3 学时与 AL 算法的无信息搜索交叉），4 领域核心学时，包含丰富选修主题。

本单元系统讲授人工智能中基于状态空间搜索的问题求解方法，涵盖盲目图搜索、启发式图搜索、博弈对抗搜索及现代局部/随机优化搜索技术。

\subsubsection{计算机核心主题}
\begin{itemize}
    \item \textbf{问题的状态空间表示}：
    \begin{itemize}
        \item 状态、初始状态、目标测试与动作算子的规范化形式定义；
        \item 状态的特征因子化表示与假设空间；
        \item 基于图搜索的问题求解本质：将图视作潜在的状态空间，将树遍历视作对该状态空间的系统性探索；
        \item 隐式状态图的动态展开机制（状态图并非预先显式完整存储，而是通过后继生成函数动态按需构建）。
    \end{itemize}
    \item \textbf{无信息图搜索技术}：
    \begin{itemize}
        \item 广度优先搜索 (BFS)及其在无权最优路径求解中的应用；
        \item 深度优先搜索 (DFS)及其内存优势；
        \item 迭代加深深度优先搜索 (IDS)的完备性与时空折中；
        \item 一致代价搜索（Uniform Cost Search, UCS / 广义 Dijkstra 算法）。
    \end{itemize}
    \item \textbf{启发式/有信息图搜索技术}：
    \begin{itemize}
        \item 启发函数的数学构建及其可采纳性与一致性；
        \item 局部搜索与爬山法；
        \item 局部极小值与搜索地形特征（搜索高地、山脊与全局最优解的辨析）；
        \item 贪婪最佳优先搜索；
        \item A* 搜索算法的核心原理、评价函数 $f(n) = g(n) + h(n)$ 及其最优性定理。
    \end{itemize}
    \item \textbf{搜索算法时空复杂度分析}：各经典搜索算法的时间复杂度与空间复杂度理论推导、完备性与最优性证明。
\end{itemize}

\subsubsection{领域核心主题}
\begin{itemize}
    \item \textbf{双向搜索}：正向搜索与逆向搜索的交汇原理与边界条件。
    \item \textbf{束搜索}：受限束宽贪婪剪枝搜索机制。
    \item \textbf{双人对抗博弈搜索}：
    \begin{itemize}
        \item 极小化极大搜索算法与递归博弈树展开；
        \item Alpha-Beta 剪枝技术（$\alpha$-$\beta$ Pruning）及其最高剪枝效率（最优走法排序下的渐近加速）；
        \item 搜索深度截断与启发式静态估值函数设计。
    \end{itemize}
    \item \textbf{A* 搜索算法的工程优化实现}：利用优先队列、开放列表与闭合列表优化大规模状态搜索性能。
\end{itemize}

\subsubsection{选修主题}
\begin{itemize}
    \item \textbf{深入洞察复杂搜索空间}：
    \begin{itemize}
        \item 搜索树构建原理与分支因子导致的组合爆炸；
        \item 动态不确定搜索空间建模；
        \item 复杂搜索地形拓扑（鞍点 Saddle Points、高维平坦高原 Plateau、狭窄山脊 Ridges 等）。
    \end{itemize}
    \item \textbf{约束满足问题 (CSP)}：约束图表示、回溯搜索、前向检验与弧相容算法（AC-3）。
    \item \textbf{禁忌搜索}：短时禁忌表与特赦准则。
    \item \textbf{内存受限与高阶 A* 变体}：IDA*（迭代加深 A*）、SMA*（简化内存受限 A*）、RBFS（递归最佳优先搜索）。
    \item \textbf{双人博弈高级专题}：地平线效应与静态搜索；开局库与残局完备解库；博弈弱解与强解的数学定义（如西洋跳棋 Checkers 已被完全弱解）。
    \item \textbf{算法实战优化}：极小化极大算法与束搜索的高性能工程实现。
    \item \textbf{期望极大搜索}：含随机概率机会节点的博弈树展开与马尔可夫决策过程求解初阶。
    \item \textbf{随机化与演化搜索算法}：
    \begin{itemize}
        \item 模拟退火算法与 Metropolis 接受准则、降温调度表设计；
        \item 遗传算法：染色体编码、轮盘赌选择、交叉、变异与适应度函数；
        \item 蒙特卡洛树搜索 (MCTS)：选择、扩展、模拟与反向传播（Backpropagation / UCB1 策略）。
    \end{itemize}
\end{itemize}

\subsubsection{示范性学习成果与能力目标}
\begin{enumerate}
    \item \textbf{设计}经典组合谜题（如 N 皇后问题、滑块拼图或水桶分水问题）的形式化状态空间表示方案。
    \item \textbf{选择并实现}适宜的无信息搜索算法（如 BFS、DFS、IDS）求解具体问题（如迷宫遍历或井字棋），并严格推导其时间与空间复杂度。
    \item \textbf{设计}具备启发性的估值函数，并在特定问题场景下（如移动机器人在二维网格地图中的避障路径规划）\textbf{编码实现}有信息启发式搜索。
    \item \textbf{论证并评估}给定启发函数是否满足可采纳性与一致性，判断其能否严格保证搜索得到最优解。
    \item \textbf{ 应用}极小化极大搜索算法求解双人零和对抗博弈（如四子棋 Connect Four），在达到指定搜索深度截断时调用静态启发估值函数进行分数回溯传递。
    \item \textbf{设计并实现}遗传算法，求解复杂的组合优化或排班调度问题。
    \item \textbf{设计并实现}模拟退火冷却计划表，帮助局部搜索过程有效跳出局部极小值陷阱。
    \item \textbf{设计并编码实现} A* 搜索与束搜索算法求解现实工程问题，并在解的最优性、扩展节点数及运行开销等维度上与其他搜索算法进行定量基准评测。
    \item \textbf{应用} Alpha-Beta 剪枝技术大幅削减双人博弈树的搜索节点规模，验证其与朴素 Minimax 搜索结果的一致性。
    \item \textbf{对比分析}遗传算法与经典基于树/图的搜索策略及数值梯度优化算法（如梯度下降法）的异同，阐明在何种特定问题场景下应优先采用演化算法。
    \item \textbf{综合权衡}各类启发式搜索策略在面对不同规模、复杂地形与算力约束问题时的有效性与工程适用度。
\end{enumerate}

\subsection{基础知识表示与推理}
\textbf{【学时建议】}：2 计算机核心学时（其中 1 学时与 MSF 离散概率交叉），2 领域核心学时。

本单元系统讲授知识表示的核心范式，从传统形式逻辑符号表示延伸至现代连续亚符号表征，并以贝叶斯理论为基础构建不确定性知识推理与理性决策的坚实理论底座。

\subsubsection{计算机核心主题}
\begin{itemize}
    \item \textbf{知识表示的核心类型}：
    \begin{itemize}
        \item 符号化与逻辑表示：将自然语言业务问题准确形式化为符号逻辑陈述；
        \item 学习驱动的亚符号表示：低维稠密向量与连续嵌入表示；
        \item 概率图模型初阶：朴素贝叶斯模型与贝叶斯信念网络直观拓扑。
    \end{itemize}
    \item \textbf{概率推理回顾与贝叶斯定理}：先验概率、后验概率、似然度；贝叶斯定理（Bayes' Theorem）的形式化陈述（与 MSF 离散概率交叉）。
    \item \textbf{贝叶斯推理基础}：基于观测证据进行假设后验概率计算与置信度更新。
\end{itemize}

\subsubsection{领域核心主题}
\begin{itemize}
    \item \textbf{随机变量与概率分布}：离散/连续随机变量、联合概率分布、边缘概率分布与条件概率分布。
    \item \textbf{概率公理体系}：柯尔莫哥洛夫公理及推论。
    \item \textbf{概率推断机制}：基于联合概率分布全表展开推断及其指数级复杂度瓶颈。
    \item \textbf{贝叶斯法则的数学推导}：从条件概率定义严谨推导贝叶斯法则。
    \item \textbf{复杂情境下的贝叶斯推理}：多证据源、序贯观测下的后验概率链式更新计算。
    \item \textbf{独立性与条件独立性}：边际独立性定义；条件独立性的数学定义及其对因果解耦的核心价值。
    \item \textbf{效用理论与理性决策}：偏好关系、效用函数与最大期望效用 (MEU)决策准则。
\end{itemize}

\subsubsection{示范性学习成果与能力目标}
\begin{enumerate}
    \item 给定自然语言描述的实际工程业务需求，将其准确\textbf{形式化编码}为符号逻辑或结构化知识表示。
    \item 结合贝叶斯定理与效用理论，深入\textbf{阐明}智能体如何在不确定性与信息不完全环境下做出最优理性决策。
    \item 运用贝叶斯定理在现实工程场景（如垃圾邮件识别、疾病筛查检测）中\textbf{执行概率推断}，准确推算在观测到特定证据条件下的假设后验概率。
    \item 熟练\textbf{应用}贝叶斯法则推导并求解多条件下的逆向概率推断问题。
    \item 准确\textbf{计算}各类随机事件的概率分布，并定量严格\textbf{检验}变量或事件之间是否满足统计独立性或条件独立性。
\end{enumerate}

\subsection{机器学习}
\textbf{【学时建议】}：4 计算机核心学时，4 领域核心学时，包含丰富选修主题。

本单元系统讲授机器学习的核心范式、基本理论与工程全流程，涵盖经典统计学习、深度神经网络、模型泛化评估以及数据偏见与算法伦理，培养学生构建工业级端到端预测模型与严谨评测系统的专业能力。

\subsubsection{计算机核心主题}
\begin{itemize}
    \item \textbf{机器学习任务的核心范式与定义}：
    \begin{itemize}
        \item 监督学习：分类（Classification，离散标签预测）与回归（Regression，连续数值预测）；
        \item 强化学习：智能体与环境通过“状态-动作-奖励”闭环试错交互学习最优策略；
        \item 无监督学习：聚类分析、数据降维与密度估计。
    \end{itemize}
    \item \textbf{机器学习基本哲学与理论基石}：
    \begin{itemize}
        \item 没有免费的午餐定理及其对算法普适性的启示；
        \item 不可判定性与学习理论的固有边界；
        \item 机器学习系统预测误差的本质来源（固有噪声、模型偏差与数据方差）。
    \end{itemize}
    \item \textbf{经典统计监督学习模型初阶}：
    \begin{itemize}
        \item 朴素贝叶斯分类器与决策树；
        \item 重点聚焦于直观机制与工程调用，不陷入底层繁复的数值优化推导，确保学生能准确理解并正确使用工业级开源实现。
    \end{itemize}
    \item \textbf{过拟合问题与模型复杂度调控}：
    \begin{itemize}
        \item 过拟合与欠拟合的现象及成因；
        \item 正则化与剪枝的直观调控逻辑；
        \item 偏差-方差权衡的核心内涵。
    \end{itemize}
    \item \textbf{真实数据工程与数据处理}：
    \begin{itemize}
        \item 数据清洗与预处理的工程重要性及常见踩坑陷阱；
        \item 缺失值处理策略：统计插补补全与显式标记缺失的技术差异及其对下游模型行为的深远影响；
        \item 类别变量编码（One-Hot 独热编码、标签序数编码）与连续实值特征变换；
        \item 数据归一化与标准化（Min-Max 缩放与 Z-score 标准化）；
        \item 坚决强调面对真实世界脏数据的处理实战，而非理想化的人工教材玩具数据。
    \end{itemize}
    \item \textbf{特征与假设表示}：
    \begin{itemize}
        \item 假设空间及其容量复杂度；
        \item 简单的基特征扩展（Basis Feature Expansion，如单变量二次方项特征构造）；
        \item 自动学习得到的特征表示（Learned Feature Representations / 特征表示学习）。
    \end{itemize}
    \item \textbf{严谨的机器学习评估方法论}：
    \begin{itemize}
        \item 训练集、验证集与独立测试集的严格隔离规范与防数据泄漏机制；
        \item 分类器常用评估指标体系（准确率 Accuracy、精确率 Precision、召回率 Recall、F1 分数等）；
        \item 基于未参与训练的保留测试集进行泛化性能无偏评估；
        \item 基于验证集进行模型超参数调优；
        \item 深入理解模型真实决策机制、明晰模型适用边界/盲区漏洞及其现实决策风险的关键意义。
    \end{itemize}
    \item \textbf{基础神经网络导论}：
    \begin{itemize}
        \item 人工神经元、多层感知机（MLP）拓扑结构、激活函数的非线性映射；
        \item 前向传播与反向传播模型训练机制的直观理解，不展开繁复矩阵微积分推导。
    \end{itemize}
    \item \textbf{机器学习工程伦理与偏见防范}：
    \begin{itemize}
        \item 基于真实商业场景与典型社会案例研究；
        \item 数据集偏见、算法设计偏见与评估基准偏见的识别。
    \end{itemize}
\end{itemize}

\subsubsection{领域核心主题}
\begin{itemize}
    \item \textbf{机器学习形式化优化建模}：
    \begin{itemize}
        \item 以最小二乘线性回归与对数几率回归为例建立优化数学模型；
        \item 目标损失函数的构建；
        \item 梯度下降法与随机梯度下降（SGD）的迭代更新机理；
        \item 抑制过拟合的正则化数学公式表达（$L_2$ 岭回归与 $L_1$ Lasso 正则项）。
    \end{itemize}
    \item \textbf{模型集成机制}：简单的加权多数表决法。
    \item \textbf{深度学习进阶}：
    \begin{itemize}
        \item 深度前馈神经网络的深层表达优势（直观机理分析）；
        \item 卷积神经网络 (CNN)的局部感受野、权重共享与池化下采样机理；
        \item 深度网络各层自动学习学到的特征表征的可视化分析。
    \end{itemize}
    \item \textbf{高阶性能评估与验证}：
    \begin{itemize}
        \item 分类高级指标（精确率、召回率、PR 曲线、Macro/Micro 平均）；
        \item 回归模型评估指标（均方误差 MSE、均方根误差 RMSE、平均绝对误差 MAE、$R^2$ 判定系数）；
        \item 混淆矩阵的构建与诊断；
        \item 交叉验证（Cross-Validation，如 $k$ 折交叉验证）；
        \item 结合交叉验证的网格搜索与随机搜索超参数调优。
    \end{itemize}
    \item \textbf{强化学习全景概览}：状态空间、动作空间、即时奖励、折扣因子、价值函数与 Q 学习直观机理。
    \item \textbf{机器学习重大跨行业工程落地剖析}：深入解析至少两个真实行业落地案例（如智慧医疗疾病诊断、量化金融风控、机器视觉质检、自然语言文本摘要、机器人运动控制、棋盘游戏博弈等）。
    \item \textbf{机器学习伦理深化}：真实场景案例研究；用户隐私保护（差分隐私初步）；算法公平性量化度量（统计均等、机会均等等算法防歧视准则）。
\end{itemize}

\subsubsection{选修主题}
\begin{itemize}
    \item \textbf{统计学习理论与参数估计}：极大似然估计（MLE）、最大后验估计（MAP）、期望最大化（EM）算法。
    \item \textbf{经典监督学习算法族深入探究}：
    \begin{itemize}
        \item 决策树划分准则（信息增益、信息增益率、基尼系数）与后剪枝算法；
        \item $k$ 近邻算法（$k$-NN）及其高维距离度量；
        \item 多层感知机反向传播链式求导法则的严密数学推导；
        \item 线性回归正规方程解析解与数值优化解对比；
        \item 逻辑回归似然函数求解与几率比分析；
        \item 支持向量机（SVM）理论：最大间隔超平面、对偶问题转换、支持向量判定与核技巧（高斯 RBF 核、多项式核）；
        \item 高斯过程非参数贝叶斯建模与核协方差函数。
    \end{itemize}
    \item \textbf{过拟合与高维学习理论}：维数灾难；$L_1$ 产生稀疏解与 $L_2$ 权重平滑收缩的严密几何与优化解析。
    \item \textbf{科学实验设计与高阶评测}：
    \begin{itemize}
        \item 分层抽样、时序分割与特征穿越防御；
        \item 受试者工作特征曲线（ROC 曲线）与曲线下面积（AUC）数学内涵。
    \end{itemize}
    \item \textbf{贝叶斯学习高阶}：朴素贝叶斯与广义线性模型的代数等价性；贝叶斯信念网络结构学习与参数估计；共轭先验分布与后验更新。
    \item \textbf{生成式模型理论}。
    \item \textbf{深度学习前沿架构与理论}：
    \begin{itemize}
        \item 深度前馈网络、神经正切核 (NTK)理论与过参数化网络训练动力学；
        \item 现代 CNN 经典拓扑（残差连接 ResNet、特征金字塔 FPN）；
        \item 自编码器与变分自编码器（VAE）；
        \item 循环神经网络（RNN、LSTM、GRU）与时间反向传播（BPTT）；
        \item 特征表示迁移、预训练微调与领域自适应；
        \item 对抗训练与生成对抗网络（GAN）。
    \end{itemize}
    \item \textbf{高阶特征表示工程}：手工特征提取算子；高维核基函数展开；深度自监督表征学习。
    \item \textbf{无监督聚类前沿}：$K$-均值聚类（$K$-Means）及其 $K$-Means++ 初始化；高斯混合模型（GMM）；自组织特征映射网络（SOM）。
    \item \textbf{半监督学习与弱监督学习}：伪标签技术与一致性正则化。
    \item \textbf{概率图模型进阶}：马尔可夫随机场（MRF）与因果推断图。
    \item \textbf{集成学习高级范式}：自举汇聚法、提升法（Boosting / AdaBoost、GBDT、XGBoost、LightGBM）、随机森林与门控专家混合模型。
    \item \textbf{统计学习理论极限}：为什么机器能够实现泛化；VC 维（Vapnik-Chervonenkis Dimension）；Rademacher 复杂度；概率近似正确（PAC）理论框架与泛化误差界。
    \item \textbf{强化学习前沿理论}：
    \begin{itemize}
        \item 探索与利用权衡（Exploration vs. Exploitation Trade-off，多臂老虎机模型）；
        \item 马尔可夫决策过程（MDP）严格数学规范与贝尔曼最优方程；
        \item 动态规划求解：价值迭代与策略迭代；
        \item 无模型强化学习：时间差分学习 TD($\lambda$)、Q-Learning、SARSA；
        \item 策略梯度方法（Policy Gradient Methods / REINFORCE 算法）；
        \item 深度强化学习（Deep Q-Networks / DQN、Actor-Critic 架构、PPO 近端策略优化算法）。
    \end{itemize}
    \item \textbf{推荐系统架构}：基于用户的协同过滤、矩阵分解（SVD）、图神经推荐算法。
    \item \textbf{机器学习专用计算硬件}：GPU 体系结构与 CUDA 并行加速、TPU 张量处理器、混合精度训练与分布式数据并行（DDP）。
    \item \textbf{多行业实战专题}：医学影像诊断、量化经济分析、智能辅导教育、计算机视觉目标检测、机器翻译 NLP、具身机器人、复杂博弈对抗。
    \item \textbf{数据挖掘深度交叉}：关联规则挖掘（Apriori 算法）、高维时序序列模式挖掘。
    \item \textbf{可信与可解释性人工智能 (Trustworthy \& Explainable AI, XAI)}：
    \begin{itemize}
        \item 数据集、算法与评估偏见的深度数学审计；
        \item 差分隐私联邦学习安全保护框架；
        \item 模型可解释性分析方法（SHAP 归因值、LIME 局部可解释替代模型、注意力权重热力图）。
    \end{itemize}
\end{itemize}

\subsubsection{示范性学习成果与能力目标}
\begin{enumerate}
    \item 深入\textbf{阐释}监督学习、强化学习与无监督学习三大核心学习范式在数据形式、优化目标与交互反馈上的根本区别。
    \item 准确\textbf{辨析}人工智能（AI）、机器学习（ML）与深度学习（DL）的层级从属关系与概念内涵。
    \item 将一个真实业务应用（如从手写像素网格图像中识别字母字符）\textbf{形式化建模}为规范的分类问题，明确界定可用输入特征维度与目标输出标签。
    \item \textbf{应用}两种或以上的经典统计学习算法（如 $k$ 近邻与逻辑回归）求解特定分类任务，并定量评估分类器的准确率与误差。
    \item 结合训练与验证学习曲线准确\textbf{识别}模型是否存在过拟合现象，并提出针对性的正则化、早停或剪枝等解决策略。
    \item \textbf{解释}机器学习算法如何在高维参数假设空间中转化为数学损失函数的搜索与连续优化迭代过程。
    \item 独立\textbf{编码实现}一个统计学习算法及其配套的梯度优化求解流程，完成模型参数训练，并在全新不可见测试样本上获得预测结果。
    \item \textbf{阐明}深度神经网络的前向计算与反向传播训练全过程，解释层级化特征表征的自动抽象机制。
    \item 深入\textbf{解释}规范严密的机器学习实验评测流程，辨析训练误差与泛化测试误差的根本差异，指出导致评测性能失真虚高的典型实验错误（如数据泄漏）。
    \item 针对基准数据集，从零\textbf{编码实现}包含数据清洗、特征预处理与无偏评估的方法管线，客观\textbf{对比}两种机器学习算法的泛化效能。
    \item 运用主流开源可视化套件（如 TensorBoard、WandB）\textbf{实时可视化}神经网络的训练损失曲线，并对网络隐藏层学到的特征表征进行降维可视化展示。
    \item 独立\textbf{编码实现}涵盖监督学习、强化学习与无监督学习的基础代表性算法原型。
    \item 针对特定的现实垂直领域问题，科学研判并\textbf{选择}最适宜的机器学习学习范式与技术路线。
    \item 深入\textbf{对比剖析}决策树、逻辑回归、朴素贝叶斯、多层神经网络与贝叶斯信念网络，结合典型应用案例说明各类算法的相对优势与最佳应用场景。
    \item 在包含缺失与噪声的复杂真实世界数据集上，独立部署并客观\textbf{评估}一个端到端预测系统的综合性能指标。
    \item 准确\textbf{概括}现代统计学习理论的学术前沿现状，客观评价其理论突破（如泛化误差界）及其在过参数化现代深度网络面前的理论局限。
    \item 深入\textbf{阐明}过拟合的发生机理，熟练掌握并\textbf{应用}交叉验证、早停、权重衰减（$L_2$ 正则化）与 Dropout 等综合防治策略。
    \item 深入\textbf{解析}机器学习中“假设空间容量”、“训练集样本规模”与“模型泛化精度”之间的三维制衡关系。
\end{enumerate}

\subsection{人工智能应用与社会影响}
\textbf{【学时建议】}：2 计算机核心学时，2 领域核心学时（与 SEP 交叉）。

\begin{infobox}[重要教学整合说明]
将人工智能的实际应用与伦理道德、算法公平性、公众信任度及可解释性议题，与具体的算法模型和数学理论有机交织、协同讲授，其教学成效远优于将伦理剥离为一门孤立专设的独立课程。在课程实施中，应尽可能将这些议题无缝嵌入到各技术单元中。例如，在讲授决策树算法时，可引导学生探讨如何将其应用于国土空间规划与土地资源可持续分配这一现实挑战，进而深入研讨自动化决策背后的环境生态公平性与伦理道德影响。

在核心学时分配上：\textbf{计算机核心} 必须至少深度覆盖一个具体工业级现实应用，并对 人工智能/ML 带来的广泛社会伦理议题建立系统概览；\textbf{领域核心} 则需进一步纵深拓展至一个或多个附加重大应用领域，深入剖析前沿深度生成式模型的工作机理，并对复杂的社会治理难题开展批判性剖析与研讨。
\end{infobox}

\subsubsection{核心主题大纲}
\begin{itemize}
    \item \textbf{人工智能在多元重大跨学科领域的实战赋能}：
    \begin{itemize}
        \item 深度覆盖智慧医疗与生命健康、环境与可持续计算、社交网络内容治理、数字金融与量化经济、智慧教育、特种机器人等广泛场景（CS 核心要求精讲至少 1 个场景，KA 核心要求拓展至少 1 个额外深度场景）；
        \item 将特定现实业务需求规范抽象并定义为 AI 形式化求解问题；
        \item 现实业务数据的可获取性、稀疏性与纯净度评估；
        \item 基础数据清洗与特征工程预处理准则；
        \item 数据集固有偏见：历史标注偏差、采样代表性失真；
        \item 算法模型偏见：目标函数损失惩罚倾斜导致的结构性失衡；
        \item 评测基准偏见：评测数据集与实际分布脱节导致的性能假象。
    \end{itemize}
    \item \textbf{前沿深度生成式模型}：
    \begin{itemize}
        \item 深度图像生成扩散模型（如截至 2023 年以来的 DALL-E、Midjourney、Stable Diffusion 等）：高层生成扩散原理、工业级商业赋能场景、版权争议、伪造滥用风险及其生成失效短板；
        \item 大语言模型（Large Language Models, LLMs，如截至 2023 年以来的 ChatGPT、Bard/Gemini 等）：自回归预训练与指令微调原理、工业级应用落地、事实性幻觉、安全对齐缺陷与逻辑推理脆弱性。
    \end{itemize}
    \item \textbf{人工智能深远社会影响与多维治理}：
    \begin{itemize}
        \item 人工智能伦理：算法决策责任归属、自主武器系统的伦理禁区、人类自主性与尊严捍卫；
        \item 算法公平性：不同群体间统计均等、均等化赔率等算法公平性数学指标，招聘与司法再犯预测中的偏见规避；
        \item 信任度与可解释性：深度黑盒模型的信任危机、决策可解释性（XAI）技术、全流程透明度与可审计性。
    \end{itemize}
\end{itemize}

\subsubsection{示范性学习成果与能力目标}
\begin{enumerate}
    \item 给定具体垂直行业的真实业务需求与复杂痛点，设计出规范严密的人工智能端到端解决方案，完整\textbf{明确}所需输入数据、清洗预处理管线、特征表示机制、核心 AI 算法选型以及全流程评估指标体系。
    \item 针对一个或多个在现实社会中已广泛部署落地的 AI 系统（如自动招聘初筛系统、智能信贷审批、人脸安防或司法量刑辅助系统），深入\textbf{分析}其广泛社会影响，精准识别并撰写关于伦理底线、算法公平性、数据偏见、用户信任与可解释性维度的专业风险评估报告。
    \item 深入\textbf{描述}当前主流文本大语言模型或图像扩散生成模型的核心失效模式（如事实幻觉、常识推理断裂、版权侵权与刻板偏见放大），客观\textbf{论述}这些局限性如何制约并影响其在高安全、高可靠关键任务场景中的工程应用与边界把控。
\end{enumerate}

\subsection{逻辑表示与推理 [选修]}
\textbf{【学时建议】}：拓展选修主题。

本单元系统涵盖基于经典逻辑与非经典逻辑的确定性知识表示与形式化推理体系，讲授本体工程、语义网以及常识推理和动作规划的形式化逻辑基础。

\subsubsection{选修主题详细大纲}
\begin{itemize}
    \item \textbf{命题逻辑与谓词逻辑回顾}：语法、语义模型、可满足性 SAT 与有效性（与 MSF 离散数学逻辑部分交叉）。
    \item \textbf{归结原理与定理自动证明}：
    \begin{itemize}
        \item 命题逻辑子句化与归结反演；
        \item 前向链接与后向链接推理算法。
    \end{itemize}
    \item \textbf{知识表示的核心工程难题}：表示充分性、推理效率与表达能力的平衡。
    \item \textbf{描述逻辑}：概念、角色、TBox（术语公理）与 ABox（断言公理）。
    \item \textbf{本体工程}：领域本体的设计原则、概念层级构建方法学。
    \item \textbf{语义网技术栈}：资源描述框架（RDF）、Web 本体语言（OWL）与大规模知识图谱。
    \item \textbf{非单调推理}：非经典逻辑、缺省逻辑与封闭世界假说 (CWA)。
    \item \textbf{形式化论辩理论}：抽象论辩框架与攻击关系。
    \item \textbf{关于动作与变化的形式化推理}：情景演算、事件演算与经典框架问题。
    \item \textbf{时间与空间推理}：阿兰时间区间代数（Allen's Interval Algebra）、定性空间拓扑关系。
    \item \textbf{逻辑程序设计}：Prolog 语言推理机制、答案集编程 (ASP)。
    \item \textbf{基于规则的专家系统}：产生式系统与 Rete 匹配算法。
    \item \textbf{语义网络} 与概念图。
    \item \textbf{基于模型的推理与基于案例的推理 (Model-based \& Case-based Reasoning, CBR)}。
\end{itemize}

\subsubsection{示范性学习成果与能力目标}
\begin{enumerate}
    \item 将自然语言表述的复杂业务陈述准确\textbf{翻译转化}为严格的一阶谓词逻辑公式。
    \item 将任意给定的命题逻辑或一阶逻辑公式正确\textbf{转换}为规范的合取范式子句集。
    \item 熟练\textbf{应用}归结反演原理对给定的一组逻辑子句进行机器推导，以形式化回答目标查询。
    \item 深入\textbf{对比剖析}结构化知识表示的主流模型（如框架系统、语义网络、本体论与描述逻辑），全面评述各自在知识表达能力与计算开销上的优势与局限。
    \item 深刻\textbf{阐明}非单调推理的核心要素，解释其作为人类常识信念系统与动态知识库维护机制的必要性。
    \item \textbf{对比分析}处理现实不确定性知识的各类形式化逻辑与定性技术手段。
    \item \textbf{对比分析}空间拓扑与定性物理关系的经典知识表示方法。
    \item \textbf{应用}情景演算或事件演算，对包含动态物理动作执行与环境状态变迁的复杂系统进行严密的形式化建模。
    \item 深入\textbf{阐释}时间推理与空间推理的本质特征、代数结构以及两者的时空协同交织关系。
    \item 深入\textbf{阐明}基于规则的推理、基于案例的推理（CBR）以及基于系统深层因果模型的推理机制的差异与应用场景。
    \item 给出形式化规划系统的严谨\textbf{定义}，阐明规划问题求解与经典状态图启发式搜索的本质区别。
    \item 深入\textbf{描述}基于状态搜索的规划、基于算子的规划与命题规划的技术路线，指明各自适用的任务环境。
    \item 深入\textbf{解释}单调推理与非单调推理的形式化边界，阐明新证据的引入在何种条件下会导致既往推论失效。
\end{enumerate}

\subsection{概率表示与推理 [选修]}
\textbf{【学时建议】}：拓展选修主题。

本单元系统讲授在不确定性、含噪声与不完全信息环境下构建概率知识表示与执行复杂概率推断的数学理论与工程实现。

\subsubsection{选修主题详细大纲}
\begin{itemize}
    \item \textbf{条件独立性回顾}：条件独立性代数公理化表述与图拓扑含义。
    \item \textbf{不确定性知识表示形式}：
    \begin{itemize}
        \item 贝叶斯网络：有向无环图结构与条件概率表（CPT）紧凑表示；
        \item 精确概率推断及其计算复杂度：变量消去法、团树算法以及精确推断的 NP-难性证明；
        \item 马尔可夫毯与有向分离定理（$d$-Separation）；
        \item 随机采样与近似推断：直接采样、拒绝采样、似然加权采样、马尔可夫链蒙特卡洛算法（MCMC，如 Gibbs 采样）；
        \item 马尔可夫网络（Markov Networks / 马尔可夫随机场 MRF）：无向图模型、势函数与吉布斯分布；
        \item 关系概率模型（Relational Probability Models / 统计关系学习）；
        \item 隐马尔可夫模型 (HMM)：前向-后向算法、Viterbi 最优路径解码算法与 Baum-Welch 参数学习算法。
    \end{itemize}
    \item \textbf{决策论高级专题}：
    \begin{itemize}
        \item 风险偏好公理与效用函数构造；
        \item 最大期望效用原则 (MEU)；
        \item 影响图（Influence Diagrams / 决策网络）。
    \end{itemize}
\end{itemize}

\subsubsection{示范性学习成果与能力目标}
\begin{enumerate}
    \item 给定具体贝叶斯网络的拓扑结构与条件概率分布参数，在观测到特定证据变量的条件下，准确\textbf{计算}待测假设变量的后验边缘概率分布。
    \item 深入\textbf{阐明}条件独立性断言如何切断图节点的稠密全连接，从而将全联合概率分布的指数级存储与计算开销急剧降低。
    \item 针对具体的复杂不确定性推理场景，举例并\textbf{构建}相应的概率图知识表示模型。
    \item 严格陈述贝叶斯网络精确概率推断的最坏情况计算复杂度，并系统\textbf{列举并对比}克服计算瓶颈的主流蒙特卡洛近似推断方法。
    \item 针对具体的现实不确定性决策场景，\textbf{设计并编程实现}至少一种概率图知识表示与推理求解引擎。
    \item 深入\textbf{剖析}时序概率推理系统（如动态贝叶斯网络与滤波模型）随时间推移面临的状态爆炸与推断复杂度挑战。
    \item \textbf{设计并编程实现}一个完整的隐马尔可夫模型（HMM），实现时序滤波估计或运用 Viterbi 算法解码最可能的隐藏时序状态序列。
    \item 深入\textbf{阐述}决策主体的主观偏好排序与实值效用函数之间的严密映射公理与数理逻辑关系。
    \item 深入\textbf{解释}如何将效用理论与概率图模型推理有机结合，以驱动自主智能体在不确定性环境下做出最大化期望效用的理性决策。
\end{enumerate}

\subsection{规划系统 [选修]}
\textbf{【学时建议】}：拓展选修主题。

本单元系统讲授自动化规划的核心理论，涵盖经典状态规划、偏序规划、分层任务网络（HTN）以及动态不确定环境下的闭环重规划机制。

\subsubsection{选修主题详细大纲}
\begin{itemize}
    \item \textbf{命题与一阶逻辑基础回顾}：规划状态的形式化表示语言。
    \item \textbf{规划算子与世界状态规范}：STRIPS 算子体系（前置先决条件 Preconditions、添加列表 Add List、删除列表 Delete List）。
    \item \textbf{全序规划}：状态空间前向搜索与目标后向搜索。
    \item \textbf{偏序规划}：因果链概念、动作定序约束与因果威胁消解。
    \item \textbf{规划图与 GraphPlan 算法}：命题层、动作层、互斥约束传播与逆向无冲突解抽取。
    \item \textbf{分层任务网络规划}：复合任务分解、非基本动作演化。
    \item \textbf{现代标准化规划描述语言}：规划领域定义语言 (PDDL)规范。
    \item \textbf{多智能体协同规划}：分布式动作协同与资源互斥。
    \item \textbf{基于马尔可夫决策过程的规划}。
    \item \textbf{闭环规划与动态重规划体系}：
    \begin{itemize}
        \item 条件规划（Conditional Planning / 感知-动作协同交织）；
        \item 连续规划；
        \item 概率规划。
    \end{itemize}
\end{itemize}

\subsubsection{示范性学习成果与能力目标}
\begin{enumerate}
    \item 针对给定的实际工程规划问题（如物流仓储包裹搬运或无人机多航点巡航），准确\textbf{构建}其状态表示、目标规范以及包含先决条件与动作效果的规划算子集。
    \item 熟练运用 PDDL 语言\textbf{形式化编写}规划领域的领域文件与问题实例文件，并调用现代化规划求解器（如 Fast Downward）求解出可行或最优动作计划。
    \item 给定初始状态、目标状态与算子集，规范绘制偏序规划因果图，准确识别因果威胁冲突，并通过添加定序因果约束\textbf{成功消解所有冲突}。
    \item 针对具体的规划问题实例，严密\textbf{构建完整的 GraphPlan 规划图}，系统推演命题层与动作层的互斥约束关系，并演示如何逆向抽取有效执行计划。
\end{enumerate}

\subsection{智能体系统 [选修]}
\textbf{【学时建议】}：拓展选修主题（与 人机交互/Collaboration and Communication 交叉）。

本单元系统讲授自主智能体与多智能体系统 (MAS)的核心架构、博弈协作机制与群体智能算法。

\subsubsection{选修主题详细大纲}
\begin{itemize}
    \item \textbf{智能体经典架构体系}：简单反应式架构、分层包容架构、认知架构（Cognitive Architectures，如 BDI、ACT-R、SOAR）。
    \item \textbf{智能体形式化理论}：信念-愿望-意图（BDI）多模态逻辑系统、智能体状态的形式化定义。
    \item \textbf{理性决策与博弈论}：纳什均衡、占优策略、零和博弈与非零和博弈。
    \item \textbf{决策论智能体与马尔可夫决策过程 (MDP)}。
    \item \textbf{软件智能体与智能助手}：个人虚拟助手、信息检索与过滤智能体。
    \item \textbf{协作智能体与信息搜集智能体}。
    \item \textbf{拟真与情感智能体}：虚拟合成角色、智能体情绪建模与拟人化交互。
    \item \textbf{学习型智能体体系}：学习元件、执行元件、评判元件与问题生成器。
    \item \textbf{多智能体系统}：
    \begin{itemize}
        \item 协作多智能体系统与智能体团队编队协作；
        \item 竞争性智能体交互机制：拍卖机制（Auctions，如英式拍卖、荷式拍卖、维克里第二价格拍卖）、投票协议与社会选择理论；
        \item 群体智能系统与生物启发式模型（Swarm Systems，如蚁群算法、粒子群优化算法 PSO）；
        \item 多智能体强化学习 (MARL)。
    \end{itemize}
    \item \textbf{人-智能体交互与协作}：混合主动性协同决策。
    \item \textbf{智能体通信方法学}：知识查询与操作语言（KQML）、FIPA 智能体通信语言（ACL）、合同网协议。
    \item \textbf{工程实践挑战与应用落地}：自动化金融高频交易智能体、智能分布式供应链协调调度。
\end{itemize}

\subsubsection{示范性学习成果与能力目标}
\begin{enumerate}
    \item 深入\textbf{阐释并对比}主流智能体架构（简单反射式、基于模型、基于目标、分层包容及认知架构）的控制机理与环境适用性。
    \item 深入\textbf{描述}智能体理论在纯软件自动化脚本、智能个人数字助手以及游戏可信拟真虚拟化身中的典型工程落地路径。
    \item 详尽\textbf{阐明}自主学习型智能体的四大核心构件及其相互之间的演进迭代机制。
    \item 结合生动具体的现实工程案例（如分布式传感器网络协作或分布式电力市场撮合交易），详尽\textbf{论证}多智能体系统如何通过交互协议与黑板机制实现高效协同。
    \item 熟练运用成熟的认知计算架构（如 ACT-R 或 SOAR），针对特定的复杂推理或人机交互任务，独立\textbf{设计并构建}一个具备长期记忆与技能学习能力的智能体原型系统。
\end{enumerate}

\subsection{自然语言处理 [选修]}
\textbf{【学时建议】}：拓展选修主题（与 人机交互/New Interactive Technologies 交叉）。

本单元系统讲授计算机处理、理解与生成人类自然语言的理论体系，涵盖经典文法解析、统计 NLP、信息检索以及基于 Transformer 与大语言模型（LLM）的现代深度学习前沿。

\subsubsection{选修主题详细大纲}
\begin{itemize}
    \item \textbf{文法体系}：确定性文法与随机/统计文法。
    \item \textbf{句法分析算法}：
    \begin{itemize}
        \item 上下文无关文法与图表句法分析器（Chart Parsers，如 CYK 动态规划算法）；
        \item 概率上下文无关文法与带权 CYK 语法分析树推导。
    \end{itemize}
    \item \textbf{语义表示与语义理解}：
    \begin{itemize}
        \item 基于形式逻辑的知识表示与语义演算；
        \item 语义角色标注 (SRL)与深层语义依存分析；
        \item 时序逻辑语义表示；
        \item 对话理解中的信念、愿望与意图（BDI）语义映射。
    \end{itemize}
    \item \textbf{语料库驱动的统计 NLP 方法}：
    \begin{itemize}
        \item $N$-元文法语言模型（$N$-grams）与隐马尔可夫模型（HMM）；
        \item 概率平滑与回退策略（Backoff，如 Good-Turing 平滑、Kneser-Ney 平滑）；
        \item 核心应用实战：词性标注与形态学词法分析。
    \end{itemize}
    \item \textbf{信息检索与抽取}：
    \begin{itemize}
        \item 向量空间模型；
        \item 词频-逆文档频率（TF-IDF）权重度量与余弦相似度计算；
        \item 检索系统评估指标：精确率、召回率、平均准确率均值（MAP）；
        \item 信息抽取技术（命名实体识别 NER、实体关系抽取）；
        \item 机器翻译（统计机器翻译 SMT 与神经机器翻译 NMT）。
    \end{itemize}
    \item \textbf{文本分类与主题建模}：词袋模型与高维稀疏特征表示。
    \item \textbf{自然语言深度学习前沿}：
    \begin{itemize}
        \item 词向量表示（Word2Vec 连续词袋/Skip-gram 模型、GloVe 词嵌入）；
        \item 循环神经网络（RNN、LSTM、GRU）的时序语义建模；
        \item 变换器自注意力架构（Transformers / 多头自注意力 Multi-head Attention、位置编码 Positional Encoding）；
        \item 预训练语言模型演进（BERT 双向自编码器、GPT 自回归解码器、T5 统一序列到序列架构）；
        \item 多模态图文联合嵌入（Multi-modal Embeddings，如 CLIP 跨模态表征对齐）；
        \item 生成式大语言模型与提示词工程。
    \end{itemize}
\end{itemize}

\subsubsection{示范性学习成果与能力目标}
\begin{enumerate}
    \item 准确界定并\textbf{对比剖析}确定性文法与统计随机文法，结合自然语言中广泛存在的句法结构歧义，举例说明引入概率文法模型的重要性与必要性。
    \item 模拟、推导并\textbf{编程实现}经典确定性句法分析算法或概率 CYK 动态规划算法，实现输入自然语言句子的层次化语法树解析。
    \item 深入\textbf{识别并阐述}人类自然语言的多重歧义性、语境高度依赖性及隐含常识对计算机实现深层语义结构化表示所构成的核心理论与工程挑战。
    \item 详述在自然语言处理中采用标准化大型基准语料库的核心价值，并系统\textbf{列举}针对语法分析、命名实体识别、阅读理解等任务的主流标准语料集。
    \item 熟练\textbf{梳理并对比}现代信息检索、神经机器翻译与大规模文本自动分类的核心算法管线与技术范式。
    \item 从零编码\textbf{实现 TF-IDF 特征提取模块}，对大规模非结构化文本语料进行特征工程降维与数值向量化，并训练调用现成的机器学习分类器构建一个端到端的高性能文本分类系统。
\end{enumerate}

\subsection{机器人学 [选修]}
\textbf{【学时建议】}：拓展选修主题（与 HCI 及 专用平台/Industrial Platforms 交叉）。

本单元系统讲授自主机器人系统的核心软硬件原理，涵盖多源传感感知、空间运动学变换、定位建图（SLAM）、运动路径规划以及人机协作安全控制。

\subsubsection{选修主题详细大纲}
\begin{itemize}
    \item \textbf{机器人学全貌与演进}：核心科学问题、技术演进脉络与尖端前沿进展。
    \item \textbf{尖端机器人传感与执行系统硬件体系}：
    \begin{itemize}
        \item 多源传感器配置：激光雷达（LIDAR）、超声波声呐、单目/双目立体视觉相机、RGB-D 深度相机、仿生事件相机、麦克风阵列、六维力敏触觉传感器等；
        \item 机械执行机构：轮式差速/全向驱动底盘、多自由度连续机械臂、仿生多指灵巧手抓具；
        \item 传感原始信号预处理、标定与数字滤波。
    \end{itemize}
    \item \textbf{机器人控制架构演进}：慎思式控制 vs. 反应式控制、混合分层控制、布雷滕伯格车辆模型。
    \item \textbf{世界模型与物理环境表征}：占用网格地图、拓扑图、八叉树地图与三维语义场景图。
    \item \textbf{物理世界固有不确定性}：传感器测量噪声、机械传动间隙与打滑、非结构化环境干扰。
    \item \textbf{空间几何描述与坐标变换}：
    \begin{itemize}
        \item 二维与三维空间位姿表示；
        \item 齐次变换矩阵、旋转矩阵与四元数运算。
    \end{itemize}
    \item \textbf{构型空间与环境地图}：构型空间维度、障碍物在 C-空间的膨胀映射、工作空间向构型空间的转换。
    \item \textbf{不确定性传感数据融合与状态估计}：卡尔曼滤波、扩展卡尔曼滤波（EKF）、粒子滤波（Particle Filter / 蒙特卡洛定位 MCL）。
    \item \textbf{自主移动机器人关键导航技术}：
    \begin{itemize}
        \item 同时定位与建图 (SLAM)；
        \item 正向运动学 (FK)与逆向运动学 (IK)；
        \item 运动路径规划与轨迹平滑优化：基于采样的规划算法（概率路图法 PRM、快速探索随机树 RRT 及渐近最优算法 RRT*）、人工势场法；
        \item 关节伺服动力学建模与顺应性阻抗控制；
        \item 基于机器视觉的伺服闭环控制。
    \end{itemize}
    \item \textbf{多机器人协作与群体机器人}：多机通信拓扑、编队控制、分布式任务分配与拍卖协商。
    \item \textbf{人-机协作与共融交互}：
    \begin{itemize}
        \item 物理人机共享工作空间安全监控与防碰撞机制；
        \item 人类作业意图预测与行为轨迹推断；
        \item 显式与隐式多模态人机交互（语音指令、肢体手势、视线跟踪）；
        \item 人机交互信任度建模与评估。
    \end{itemize}
\end{itemize}

\subsubsection{示范性学习成果与能力目标}
\begin{infobox}[教学实现模式说明]
鉴于实体工业级机器人硬件较为昂贵，本单元的所有实践与考核成果均可依托高保真仿真环境（如 Gazebo、Isaac Sim、Webots）或低成本开源教育机器人平台完成。
\end{infobox}
\begin{enumerate}
    \item 系统梳理并\textbf{归纳列举}当今顶尖机器人系统的核心能力与物理局限性，详述各类核心车载传感器（LIDAR、深度视觉、IMU）的工作原理及其不可或缺的数据处理流水线。
    \item 针对特定的自动化作业任务，综合集成外设传感器、驱动执行器与底层控制软件，独立\textbf{搭建一套完整的软硬件闭环机器人系统}。
    \item 熟练运用慎思式、反应式或混合分层控制架构，为移动机器人\textbf{编写控制程序}以自主实现简单作业任务。
    \item 深入理解构型空间几何映射机制，并在高维构型空间中\textbf{实现基础的运动路径规划算法}（如 A*、RRT 或 Bug 导航算法）。
    \item 定量刻画常见工业级传感器与执行机构的物理不确定性与噪声漂移，并熟练\textbf{制定并实现}针对这些不确定性的算法滤波与状态估计补偿方案。
    \item 深入\textbf{辨析}机器人在其内部表示外部物理世界时的各类拓扑与几何表示方案（栅格地图、点云地图、拓扑图），客观评述各自的优势与不足。
    \item 详尽\textbf{对比剖析}在已知地图与完全未知探索环境下机器人自主导航的主流技术路线（全局拓扑导航、局部实时避障 DWA 等），剖析各自的适用环境与鲁棒性边界。
    \item 深入\textbf{阐释}至少一种协调多台自主机器人协同感知与同步动作以合力完成复杂搬运或搜救任务的分布式协同算法。
    \item 深入\textbf{对比}“多机协同”与“人机共融协作”在动力学安全规范、意图辨识与突发交互容错上的核心差异，并将差异根源追溯至人机物理与认知环境的不确定性本质。
\end{enumerate}

\subsection{感知与计算机视觉 [选修]}
\textbf{【学时建议】}：拓展选修主题。

本单元系统讲授计算机从图像、视频、音频及多模态物理传感器信号中抽取结构化知识的高级理论，涵盖传统图像几何处理、特征工程、音频信号解析与现代深度感知学习网络。

\subsubsection{选修主题详细大纲}
\begin{itemize}
    \item \textbf{计算机视觉成像与基础表示}：数字图像采集几何光学原理、空间采样与灰度量化、多通道矩阵表示、色彩空间（RGB、HSV、Lab）变换与图像底层属性。
    \item \textbf{图像底层特征工程与几何分析}：形状表示、边缘检测算子（Sobel、Canny）、角点检测（Harris、FAST）、局部尺度不变特征描述符（SIFT、SURF、ORB）。
    \item \textbf{图像分割技术}：自适应阈值分割（Otsu 算法）、基于边缘跟踪的分割、区域生长算法与图割。
    \item \textbf{运动分析与光流估计}：光流约束方程、Lucas-Kanade 稀疏光流与 Horn-Schunck 稠密光流算法。
    \item \textbf{生成式视觉前沿}：扩散概率模型、神经辐射场、三维高斯泼溅（3D Gaussian Splatting）。
    \item \textbf{音频与语音识别基础}：连续声学信号数字化采样、短时傅里叶变换（STFT）、梅尔频率倒谱系数（MFCC）、端到端语音识别声学模型。
    \item \textbf{触觉与本体感受}：触觉阵列传感器信号解析、机器人自身运动力矩感知。
    \item \textbf{其他多感官跨模态感知}（如电子鼻气味传感等）。
    \item \textbf{感知模式识别范式与模块化架构}（与 人工智能/机器学习 深度交叉）。
    \item \textbf{分类与目标检测评估指标体系}：Top-1/Top-5 准确率、交并比、平均精度均值。
    \item \textbf{统计模式识别技术}。
    \item \textbf{深度感知学习前沿架构}：现代卷积神经网络（ResNet、ConvNeXt）、视觉变换器、端到端目标检测网络（YOLO 系列、DETR）、实例分割与大模型分割范式（Mask R-CNN、Segment Anything Model, SAM）。
\end{itemize}

\subsubsection{示范性学习成果与能力目标}
\begin{enumerate}
    \item 深入\textbf{阐述}图像与目标识别在现代人工智能技术生态中的基石地位，并系统列举其在工业自动化质检、自动驾驶、智慧医疗影像筛查等关键领域的重大产业应用。
    \item 详尽\textbf{列举}至少三种图像分割的核心算法范式（阈值分割、边缘算子分割、区域生长/图割法），剖析各自的算法数学定义、计算瓶颈与适用场景。
    \item \textbf{编码实现}基于几何轮廓边缘描述或二值区域特征的二维工业零件几何匹配与目标识别算法。
    \item 严格\textbf{辨析}声音事件检测、连续语音转写与特定说话人声纹确认在感知目标上的本质差异，指明原始音频时序信号在上述三类任务中预处理与特征抽取的根本不同。
    \item 举出至少两个在跨感官模态之间实现感知数据域转换的工程范例（例如将指尖高密度微触觉力敏阵列信号映射转化为单通道二维灰度压力分布图像进行视觉化卷积解析）。
    \item 在真实物理场景数据上独立\textbf{编程实现特征提取算子}（如实现高精度 Canny 边缘算子、Harris 亚像素角点检测器，或提取音频短时切片的快速傅里叶变换 FFT 频谱特征向量）。
    \item 编写算法将底层原始基元特征有效组合凝聚为高阶感知结构（例如将离散像素边缘连接拟合为凸多边形或平滑样条曲线，或将音频声学特征向量聚类映射为候选音素假设）。
    \item \textbf{编程实现}一个完备的模式分类器，将输入的连续感知特征向量正确划分至目标类别空间，并运用混淆矩阵对其泛化分类性能进行客观定量评估。
    \item 运用严格的控制变量对照实验，定量\textbf{评估}底层特征提取算法相较于其他备选特征提取机制对最终端到端模式识别准确率的独立贡献度与鲁棒性表现。
    \item 深入\textbf{阐释}至少三种经典感知分类算法的数学原理、先验假设前提、技术优势以及在复杂高维非凸地形下的性能短板。
    \item 针对计算机视觉中的复杂感知任务（如大规模通用物体识别、开集检测或复杂街景多目标语义分割），设计、微调并部署一个基于现代深度神经网络（如先进 CNN 或 ViT 架构）的工业级感知解决方案，并完成高标准的定量评估与算力推理时延调优。
\end{enumerate}

\section{专业素养与思维倾向}

在人工智能领域的理论研究与工程实践中，学生应在专业思维与职业素养方面重点培育以下三种核心倾向：

\begin{itemize}
    \item \textbf{严谨细致}：在实现人工智能和机器学习算法时，必须对细微之处保持高度敏锐与专注。数据分布的细微偏移、微小的超参数变动、张量维度的不对齐或概率平滑因子的微调，均可能引发灾难性的模型震荡或隐蔽的系统漏洞。学生必须具备对数学推导、数据管道及底层代码实现精益求精、一丝不苟的科学态度。
    \item \textbf{坚韧执着}：人工智能技术往往运行在部分可观察、高噪声与非平稳的复杂开放物理环境中，优化算法在多轮梯度迭代中极易遭遇误差级联放大、梯度消失或陷入鞍点。要使一个 AI 系统具备确定性、高鲁棒的工程落地表现，通常需要经历漫长的参数调优、假设验证、失败归因与反复试错。这种极具挑战性的技术特性要求学生在面对复杂阻碍与模型失效时，具备强大的心理韧性与持之以恒的攻关毅力。
    \item \textbf{责任与伦理担当}：人工智能的应用已深远渗透至社会生产与日常生活的各个角落，其输出结果直接影响微观个体的合法权益、人身财产安全以及宏观社会的公平与稳定。这要求学生深刻洞悉 AI 技术带来的跨学科社会影响，始终秉持对人类福祉高度负责的伦理底线，审慎决策何时何地适宜应用 AI，并自觉落实算法透明度、数据隐私保护与无歧视的伦理规范。
\end{itemize}

\section{数学基础要求}

人工智能作为一门理论深厚、高度跨学科的硬核计算机科学分支，要求学生具备扎实、严密的数理科学底座：

\subsection{必修数学基础}
\begin{itemize}
    \item \textbf{离散数学}：
    \begin{itemize}
        \item 集合论（集合代数、笛卡尔积、幂集）、二元关系（等价关系、偏序关系）、映射与函数理论；
        \item 命题逻辑与一阶谓词逻辑语法与语义、逻辑推理规则、形式化数学归纳法与反证法。
    \end{itemize}
    \item \textbf{线性代数}：
    \begin{itemize}
        \item 矩阵代数系统：矩阵乘法、转置、行列式、伴随矩阵、矩阵求逆、秩与线性方程组求解；
        \item 向量空间理论：线性无关性、正交基底、子空间投影、特征值与特征向量分解、奇异值分解（SVD）。
    \end{itemize}
    \item \textbf{概率论与数理统计}：
    \begin{itemize}
        \item 基础概率公理系统、条件概率计算、全概率公式、事件独立性；
        \item 贝叶斯定理及其在逆向推断中的严谨应用；
        \item 随机变量数学期望（均值）、方差、协方差、高阶矩；
        \item 经典概率分布族（高斯正态分布、伯努利分布、二项分布、泊松分布、多项分布）；
        \item 描述性统计量与推断统计学基础（假设检验、$p$-值检验、置信区间估计）；
        \item \textbf{【教学特别强调】}：上述所有概率统计理论必须紧密围绕真实工程决策场景、运用真实世界“脏数据”开展案例教学，坚决摒弃脱离现实工程本质的纯人工玩具例题。
    \end{itemize}
\end{itemize}

\subsection{进阶与推荐数学基础}
\begin{itemize}
    \item \textbf{基于微积分的高等概率与连续统计}：多维连续型随机变量联合概率密度函数、多重微积分变换、概率生成函数与特征函数。
    \item \textbf{高级推断与统计实验设计}：科学实验设计方法学、数据重采样与非参数方法（Bootstrap 自助抽样法、Jackknife 刀切法、排列检验）。
    \item \textbf{数值优化理论}：凸集与凸函数理论、多变量连续函数无约束梯度优化、拉格朗日乘子法与 KKT 最优性条件、对偶理论、牛顿法与拟牛顿法（L-BFGS）。
\end{itemize}

\section{交叉主题与跨学科共享概念}

\subsection{跨领域共享概念}
\begin{itemize}
    \item \textbf{搜索算法}：人工智能/搜索算法与问题求解 单元与 算法/Algorithms and Complexity（算法基础与复杂性）中的高级图算法和搜索策略深度融合。
    \item \textbf{数据管理与数据工程}：机器学习全流程中的高维数据治理、特征仓库、分布式数据读取与 数据管理/Data Management 及数据科学基础设施紧密贯通。
    \item \textbf{伦理与社会规范}：人工智能/人工智能应用与社会影响 单元与 社会伦理/Society, Ethics and Professionalism（社会、伦理与专业素养）构建了统一的科技治理视角。
    \item \textbf{机器人学与嵌入式硬件}：人工智能/机器人学 单元与 专用平台/Specialized Platform Development（专用平台开发中的机器人工业控制平台）构成软硬件一体化的协同映射。
\end{itemize}

\subsection{贯穿式核心维度}
\begin{itemize}
    \item \textbf{计算效率}：贯穿于高维状态搜索的剪枝加速、图模型推断的因果解耦，以及大规模深度神经网络的分布式算力并行优化中。
    \item \textbf{伦理道德}：贯穿于从数据采集标注、特征工程设计，到模型训练偏见控制、可解释性审查及司法/医疗落地全生命周期的伦理约束。
    \item \textbf{系统建模}：贯穿于将开放非结构化现实业务难题抽象转化为状态图搜索、确定性逻辑命题、概率图模型以及损失函数优化数学模型的建模全过程。
    \item \textbf{工程编程}：强调将复杂的数学优化方程转化为高内聚、低耦合、高吞吐的工业级高质量生产代码实现。
\end{itemize}

\section{课程设计与开课规划建议}

联合工作组针对不同培养侧重，提出三套典型的人工智能课程体系设计方案：

\subsection{方案一：《人工智能基础导论》（38 学时综合基石课程）}
面向全校或计算机大类专业通识进阶设置，旨在让学生系统建立对人工智能各核心领域的全景认知与扎实求解能力：
\begin{itemize}
    \item \textbf{学时构成}：
    \begin{itemize}
        \item 人工智能/人工智能基础议题（基础议题）：4 学时
        \item 人工智能/搜索算法与问题求解（搜索算法与问题求解）：9 学时
        \item 人工智能/知识表示与推理基础（基础知识表示与推理）：4 学时
        \item 人工智能/Probabilistic Knowledge Representation and Reasoning（概率表示与推理）：5 学时
        \item 人工智能/机器学习（机器学习）：12 学时
        \item 人工智能/人工智能应用与社会影响（应用与社会影响）：4 学时（强调整合贯穿全课程各讲）
    \end{itemize}
    \item \textbf{先修课程要求}：CS2（数据结构与面向对象编程）、离散数学、概率论基础。
    \item \textbf{能力目标陈述}：完成本课程的学生能够系统掌握人工智能各大核心子领域的基本范式与方法论，具备在各方向独立理解、开发并应用代表性算法的核心能力；能够运用状态搜索技术、基础贝叶斯推理与经典机器学习方法求解复杂的现实工程问题；能够全面理解 AI 技术的广泛应用潜力，并能对其引发的法律合规、算法偏见与社会伦理挑战进行理性的审视与权衡。
\end{itemize}

\subsection{方案二：《机器学习原理与实战》（44 学时深度算法课程）}
面向数据科学、智能科学与技术或研究型计算机专业方向专设，聚焦于数据驱动的学习系统理论与工业级实战：
\begin{itemize}
    \item \textbf{学时构成}：
    \begin{itemize}
        \item 人工智能/机器学习（机器学习与深度网络前沿）：32 学时
        \item 人工智能/知识表示与推理基础（知识表示与贝叶斯推断）：4 学时
        \item 人工智能/自然语言处理（自然语言处理精选专题：TF-IDF、词袋模型、文本分类与 Transformer 初步）：4 学时
        \item 人工智能/人工智能应用与社会影响（应用与伦理公平审计）：4 学时（贯穿全生命周期实验教学）
    \end{itemize}
    \item \textbf{先修课程要求}：CS2、离散数学、概率论、数理统计、线性代数（强烈推荐）。
    \item \textbf{能力目标陈述}：完成本课程的学生能够透彻理解、自主设计、独立开发并部署先进的监督学习与强化学习模型；能够针对特定工程问题准确选择最适宜的算法范式，严谨执行无泄漏的数据预处理与特征工程，设计规范科学的评估基准，并深入解释产出模型的决策逻辑与其内在缺陷；能够敏锐识别并消除有偏数据集及各类系统误差来源，在将机器学习应用于实际工程问题时严格践行符合伦理道德与社会责任的标准。
\end{itemize}

\subsection{方案三：《自主机器人系统与智能控制》（42 学时机电一体化软硬件课程）}
面向智能装备、物联网与具身智能领域专设，强调软硬件深度交织的物理交互与环境闭环导航：
\begin{itemize}
    \item \textbf{学时构成}：
    \begin{itemize}
        \item 人工智能/机器人学（机器人学核心理论与实践）：25 学时
        \item 专用平台/机器人平台（专用机器人硬件平台、物理约束与嵌入式软件架构）：4 学时
        \item 人工智能/搜索算法与问题求解（机器人路径搜索与避障规划，如 A*、RRT 算法）：4 学时
        \item 人工智能/机器学习（机器学习在机器人感知与视觉目标识别中的应用）：6 学时
        \item 人工智能/人工智能应用与社会影响（自主机器人系统社会伦理与安全影响）：3 学时（融入全流程安全论证）
    \end{itemize}
    \item \textbf{先修课程要求}：CS2、线性代数与空间解析几何。
    \item \textbf{能力目标陈述}：完成本课程的学生能够系统掌握先进机器人学技术体系；熟练运用多源传感器感知物理外部世界，基于提取的环境特征与高精度地图实现鲁棒的自主实时定位（SLAM），并在简单及中等复杂度的现实场景中完成运动路径规划与自主移动闭环导航；能够熟练运用基础计算机视觉处理算法、机械臂正逆运动学解析求解技术及空间轨迹优化方法。
\end{itemize}

\section{真实能力任务规范}

真实能力规范是 CS2023 引入的重大教育创新，通过高还原度的真实工程挑战，综合考量学生的知识融通、高阶认知与专业素养。

\begin{taskbox}[任务 1：销售库存智能优化决策工具]
\textbf{【任务情境】}：某全国性连锁商超或跨国零售企业的区域销售团队面临数万种 SKU 商品在多门店与中心仓之间的动态调配难题，需要在满足复杂货架库容、冷链保鲜时效及运输履约成本多重严苛物理制约的前提下，实现库存调配方案的全局动态优化。\\
\textbf{【能力陈述】}：将该业务库存调度问题规范抽象为状态空间图搜索模型，准确界定显式与隐式业务边界约束，设计具有严密理论可采纳性的高品质启发函数，编码实现智能化图搜索算法（如 A* 搜索结合约束满足 CSP 求解器），高效求解出满足各项严苛约束的最优或近优库存调配执行方案。\\
\textbf{【能力领域】}：应用 \quad \textbf{【能力单元】}：系统设计 / 软件开发 / 性能评估\\
\textbf{【所需知识单元】}：人工智能/搜索算法与问题求解\\
\textbf{【认知与技能要求】}：\textbf{应用 / 开发} \quad \textbf{【核心要求】}：计算机核心
\end{taskbox}

\begin{taskbox}[任务 2：不确定环境下双人博弈策略智能体]
\textbf{【任务情境】}：设计一个在复杂双人零和对抗博弈（如包含随机抽卡或掷骰子机会事件的棋盘游戏，或带有突发价格波动与信息不完全的战略采购博弈）中能够自主做出全局最优决策的博弈智能体。\\
\textbf{【能力陈述】}：熟练构建带有 Alpha-Beta 剪枝的极小化极大算法（Minimax with $\alpha$-$\beta$ Pruning），引入处理随机不确定性概率动作的期望机会节点，并在指定搜索深度处设计合理的启发式局面估值函数，开发并部署一个能在不确定对抗环境下具备高胜率对抗水平的智能博弈决策引擎。\\
\textbf{【能力领域】}：应用 \quad \textbf{【能力单元】}：需求分析 / 架构设计 / 软件开发 / 仿真测试 / 生产部署\\
\textbf{【所需知识单元】}：人工智能/搜索算法与问题求解、人工智能/知识表示与推理基础\\
\textbf{【认知与技能要求】}：\textbf{应用 / 开发} \quad \textbf{【核心要求】}：领域核心
\end{taskbox}

\begin{taskbox}[任务 3：基于多源体检数据的疾病患病概率贝叶斯预测系统]
\textbf{【任务情境】}：智慧医疗系统需要基于病人体检报告中呈现的非完美化验指标、不完全既往病史及流行病学先验统计，对特定高危慢性病或急性并发症的发病风险给出具备统计置信度的概率推断。\\
\textbf{【能力陈述】}：将现实临床诊断中的多维不确定性因果关联建模为概率贝叶斯推理问题，熟练应用贝叶斯法则计算后验概率，编码实现精确或采样概率推断算法，根据输入的临床观测证据实时推导出疾病发生率并给出具备统计说服力的合理解释。\\
\textbf{【能力领域】}：应用 \quad \textbf{【能力单元】}：系统开发 / 实验测试\\
\textbf{【所需知识单元】}：人工智能/搜索算法与问题求解、人工智能/逻辑表示与推理、数学统计/概率与统计\\
\textbf{【认知与技能要求】}：\textbf{应用 / 开发} \quad \textbf{【核心要求】}：计算机核心
\end{taskbox}

\begin{taskbox}[任务 4：现代复杂物流调度与分拣规划求解器]
\textbf{【任务情境】}：大型现代化智能物流枢纽面临海量多目的地包裹的自动化交叉分拣、机械臂连续抓取位姿衔接与无人引导搬运车（AGV）无碰撞路径协同的时空连续规划挑战。\\
\textbf{【能力陈述】}：将该复杂物流调度任务形式化规范为规划初始状态、目标状态以及包含先决条件与动作效果的规划算子集，运用 PDDL 语言建模并采用偏序规划或分层任务网络规划算法，求解出无资源死锁且全局搬运效率最优的动作序列。\\
\textbf{【能力领域】}：应用 \quad \textbf{【能力单元】}：系统设计 / 理论阐释\\
\textbf{【所需知识单元】}：人工智能/逻辑表示与推理、人工智能/规划系统\\
\textbf{【认知与技能要求】}：\textbf{解释 / 应用} \quad \textbf{【核心要求】}：拓展选修
\end{taskbox}

\begin{taskbox}[任务 5：表格化用户消费行为监督学习预测流水线]
\textbf{【任务情境】}：某大型零售或金融科技企业积累了海量高维表格化交易日志数据（如用户历史消费明细、APP 浏览频次、人口统计学标签等），需挖掘关键消费趋势并对高价值客户的未来流失概率进行精准预测。\\
\textbf{【能力陈述】}：熟练运用工业级机器学习库（如 Scikit-Learn、XGBoost），自主构建端到端的数据预处理流水线（缺失值工程补全、高维类别特征 One-Hot 与 Target 编码、特征无偏缩放），严格隔离数据集以防御数据泄漏，实施分层交叉验证与超参数网格搜索调优，构建出具备高泛化能力的监督学习分类模型并进行专业质量评估。\\
\textbf{【能力领域】}：应用 \quad \textbf{【能力单元】}：需求分析 / 架构设计 / 算法研发 / 效能评估 / 生产部署\\
\textbf{【所需知识单元】}：人工智能/机器学习、人工智能/人工智能应用与社会影响\\
\textbf{【认知与技能要求】}：\textbf{应用 / 开发} \quad \textbf{【核心要求】}：计算机核心
\end{taskbox}

\begin{taskbox}[任务 6：基于深度迁移微调的商品货架图像识别系统]
\textbf{【任务情境】}：无人零售便利店与智能结算货架需要通过高分辨率工业相机拍摄的杂乱货架照片，实时高精度识别上百种摆放各异、存在光照变化与局部遮挡的零售商品品类。\\
\textbf{【能力陈述】}：掌握现代深度学习与迁移学习技术体系，基于先进的预训练视觉大模型（如 ResNet 或 Vision Transformer Backbone），设计任务专用的特征分类头，利用标注业务样本完成模型微调，使用严谨的混淆矩阵、Top-1 准确率与 mAP 指标全面评估模型质量，并将模型工程化导出部署至边缘计算设备进行实时推理。\\
\textbf{【能力领域】}：应用 \quad \textbf{【能力单元】}：需求分析 / 架构设计 / 算法研发 / 系统测试 / 效能评估 / 生产部署\\
\textbf{【所需知识单元】}：人工智能/机器学习、人工智能/人工智能应用与社会影响\\
\textbf{【认知与技能要求】}：\textbf{应用 / 开发} \quad \textbf{【核心要求】}：领域核心
\end{taskbox}

\begin{taskbox}[任务 7：视频游戏中的深度强化学习智能对抗体]
\textbf{【任务情境】}：在复杂即时战略或动作类视频游戏中，开发团队需要构建一个具备高水平博弈意识、能够自主适应玩家战术流派演进的拟人化智能非玩家角色（AI 对手）。\\
\textbf{【能力陈述】}：深刻理解并应用现代深度强化学习 (DRL)算法框架（如 PPO、DQN 或 Actor-Critic 架构），将原始游戏画面像素与物理状态向量构建为环境状态，设计多目标奖励函数，完成策略网络的端到端训练与超参数收敛性调优，成功部署并在复杂视觉输入下展现出高胜率竞技水平。\\
\textbf{【能力领域】}：应用 \quad \textbf{【能力单元】}：需求定义 / 算法设计 / 策略开发 / 仿真评估 / 游戏部署\\
\textbf{【所需知识单元】}：人工智能/机器学习\\
\textbf{【认知与技能要求】}：\textbf{应用 / 开发} \quad \textbf{【核心要求】}：领域核心
\end{taskbox}

\begin{taskbox}[任务 8：社交媒体多模态/文本情绪与情感分析系统]
\textbf{【任务情境】}：数字舆情监测平台需要对主流社交媒体上每日产生的数百万条用户短文本发帖与图文评论进行实时多维情感极性（正面、负面、中性及具体细粒度情绪）精准识别与趋势预警。\\
\textbf{【能力陈述】}：构建非结构化文本处理流水线，执行文本清洗、分词、停用词过滤与现代特征提取（如 TF-IDF 向量化结合密集词嵌入，或微调预训练 Transformer），训练并调优监督分类器完成文本情感分类，输出包含置信度与情感极性权重的可视化分析报告。\\
\textbf{【能力领域】}：应用 \quad \textbf{【能力单元】}：需求分析 / 系统设计 / 算法研发 / 效能评估 / 生产部署\\
\textbf{【所需知识单元】}：人工智能/机器学习、人工智能/自然语言处理、人工智能/人工智能应用与社会影响\\
\textbf{【认知与技能要求】}：\textbf{应用 / 开发} \quad \textbf{【核心要求】}：领域核心
\end{taskbox}

\begin{taskbox}[任务 9：已部署机器学习系统的数据集偏差与算法偏见审计校正]
\textbf{【任务情境】}：某大型金融机构或公共服务部门在生产环境中运行的自动化信贷审批与人才简历筛选模型，被外部公众或监管机构质疑存在针对特定性别、年龄层或少数族裔群体的系统性偏见。\\
\textbf{【能力陈述】}：深刻理解、准确定位并量化审计生产机器学习模型中的各类偏见根源（历史数据代表性缺失偏差、特征标注偏差、算法损失惩罚倾斜及评测基准失真），撰写符合合规标准的专业审计报告，设计算法层面的缓解修正策略（如反事实再加权、对抗去偏训练、决策阈值多重校准），显著降低不公概率并向治理委员会做技术汇报。\\
\textbf{【能力领域】}：应用 \quad \textbf{【能力单元】}：沟通协作与专业审计\\
\textbf{【所需知识单元】}：人工智能/机器学习、人工智能/人工智能应用与社会影响\\
\textbf{【认知与技能要求】}：\textbf{理解与阐述} \quad \textbf{【核心要求】}：计算机核心
\end{taskbox}

\begin{taskbox}[任务 10：基于大语言模型（LLM）的一线智能客服系统设计与风控对齐]
\textbf{【任务情境】}：跨国电商或电信运营商拟引入基于商用大语言模型（LLM）的一线智能客服系统，实现 7$\times$24 小时自然交互，但系统必须严防事实幻觉误导用户、敏感数据泄露或被提示注入攻击。\\
\textbf{【能力陈述】}：深刻理解现代生成式深度大模型的内在工作机理与潜在失效模式（如事实性幻觉 Hallucination、逻辑脆弱性、越狱诱导脆弱性等），合理设计检索增强生成（RAG）管道与精细化提示工程，搭建安全围栏与人工兜底回退机制，在充分释放大模型高情商对话能力的同时严控业务安全与法律合规风险。\\
\textbf{【能力领域】}：应用 \quad \textbf{【能力单元】}：沟通论证与技术方案设计\\
\textbf{【所需知识单元】}：人工智能/人工智能应用与社会影响\\
\textbf{【认知与技能要求】}：\textbf{理解与阐述} \quad \textbf{【核心要求】}：计算机核心
\end{taskbox}

\section{工作组委员会名单}

\textbf{工作组主席}：
\begin{itemize}
    \item Eric Eaton，宾夕法尼亚大学 (USA)
\end{itemize}

\textbf{工作组委员}：
\begin{itemize}
    \item Zachary Dodds，哈维穆德学院 (USA)
    \item Susan L. Epstein，纽约市立大学亨特学院与研究生中心（Hunter College \& The Graduate Center of CUNY, New York, NY, USA）
    \item Laura Hiatt，美国海军研究实验室 (USA)
    \item Amruth N. Kumar，新泽西拉马波学院 (USA)
    \item Peter Norvig，谷歌公司 / 斯坦福大学
    \item Meinolf Sellmann，通用电气全球研发中心 (USA)
    \item Reid Simmons，卡耐基梅隆大学 (USA)
\end{itemize}

\textbf{贡献专家}：
\begin{itemize}
    \item Claudia Schulz，汤森路透
\end{itemize}

\clearpage
\onecolumn
\section{附录：核心主题与技能要求等级对照表}

本附录完整汇总了 CS2023 人工智能知识领域中所有计算机核心与领域核心知识点、具体主题内容、布鲁姆认知技能达标要求等级及推荐核心学时分配：

\begingroup
\zihao{-5}
\renewcommand{\arraystretch}{0.92}
\setlength{\tabcolsep}{3pt}
\noindent\begin{tabularx}{\textwidth}{>{\raggedright\arraybackslash}p{0.14\textwidth} Y >{\centering\arraybackslash}p{0.15\textwidth} >{\centering\arraybackslash}p{0.09\textwidth} >{\centering\arraybackslash}p{0.06\textwidth}}
\toprule
\textbf{知识单元} & \textbf{核心知识点与具体主题} & \textbf{认知等级} & \textbf{分类} & \textbf{学时} \\
\midrule
\multirow{4}{*}{\makecell[l]{人工智能/基础议题}} 
& AI 问题概述、近期重大成功应用案例；智能体定义与典型范式（反应式、深思式）；什么是智能行为；图灵测试及其理论缺陷；多模态输入与输出；智能行为的计算模拟；理性推理 vs. 非理性/类人推理 & 解释 & 计算机核心 & \multirow{3}{*}{2} \\
\cline{2-4}
& 问题特征分类（完全/部分可观察、单/多智能体、确定/随机、静态/动态、离散/连续）；智能体本质与自主性特征（完全/半/混合自主，反射/目标/效用，不确定性与不完全信息决策，感知与环境交互核心地位，基于学习的智能体，具身智能体三要素：传感器、动力学、执行器） & 评估 & 计算机核心 & \\
\cline{2-4}
& 人工智能产业应用、爆发式增长与多维深远影响（经济、社会、法律与伦理治理） & 解释 & 计算机核心 & \\
\cline{2-5}
& 结合复杂现实工程实例对环境与问题特征的深度剖析；结合工业级案例对各类智能体架构与行为特性的深度分析；AI 发展趋势与宏观经济、社会、伦理冲击的深入研讨 & 评估 & 领域核心 & 1 \\
\midrule
\multirow{5}{*}{\makecell[l]{人工智能/搜索算法\\与问题求解}} 
& 状态空间表示（状态、目标与动作算子规范，假设空间状态因子化，基于图搜索的问题求解，图的隐式动态按需展开） & 解释 & 计算机核心 & \multirow{4}{*}{\makecell[c]{5\\{\footnotesize (含 3h}\\{\footnotesize 共享 AL)}}} \\
\cline{2-4}
& 无信息图搜索求解（广度优先搜索 BFS、深度优先搜索 DFS 及其迭代加深 IDS、一致代价搜索 UCS） & 开发/应用 & 计算机核心 & \\
\cline{2-4}
& 启发式图搜索求解（启发函数构建与可采纳性/一致性、爬山法、局部极小值与搜索地形、贪婪最佳优先搜索、A* 搜索） & 开发/应用 & 计算机核心 & \\
\cline{2-4}
& 图搜索算法的时间复杂度与空间复杂度理论推导、完备性与最优性严密评估 & 评估 & 计算机核心 & \\
\cline{2-5}
& 双向搜索、束搜索；双人对抗博弈（极小化极大搜索 Minimax、Alpha-Beta 剪枝技术、深度截断与启发式静态估值）；A* 算法的高效工程优化实现 & 开发/应用 & 领域核心 & 4 \\
\midrule
\multirow{3}{*}{\makecell[l]{人工智能/基础知识\\表示与推理}} 
& 知识表示类型（符号与逻辑表示、从自然语言抽象形式逻辑、学习驱动的亚符号连续向量嵌入、概率图模型初阶如朴素贝叶斯与贝叶斯网） & 解释 & 计算机核心 & \multirow{2}{*}{\makecell[c]{2\\{\footnotesize (含 1h}\\{\footnotesize 共享 MSF)}}} \\
\cline{2-4}
& 概率推理回顾、贝叶斯定理（与 MSF 离散概率交叉）；贝叶斯推断与证据置信度更新计算 & 应用 & 计算机核心 & \\
\cline{2-5}
& 随机变量与概率分布、概率公理体系、概率推断机制、贝叶斯法则数学推导、复杂场景下的贝叶斯推理与证据更新、边际独立性与条件独立性、效用理论与最大期望效用决策 & 应用 & 领域核心 & 2 \\
\midrule
\multirow{5}{*}{\makecell[l]{人工智能/机器学习}} 
& 机器学习主流任务范式（监督学习：分类与回归；强化学习；无监督学习：聚类）；理论哲学（没有免费的午餐定理、不可判定性、系统误差来源） & 应用/开发/评估 & 计算机核心 & \multirow{3}{*}{4} \\
\cline{2-4}
& 简明统计监督学习（朴素贝叶斯、决策树原理与调用）；过拟合与复杂度控制（正则化、剪枝直观理解，偏差-方差权衡）；真实数据工程（清洗、缺失值填补/标记、特征编码、归一化/标准化、真实脏数据处理）；特征表示（假设空间、基特征扩展、表示学习）；模型严谨评测（训练/验证/测试集隔离、分类评估指标、保留测试集性能估计、验证集超参数调优、模型局限洞察）；基础神经网络（前向与反向传播直观理解） & 应用/开发/评估 & 计算机核心 & \\
\cline{2-4}
& 机器学习工程伦理（真实场景案例、数据集偏见、算法偏见、评估偏见识别） & 解释/评估 & 计算机核心 & \\
\cline{2-5}
& 形式化优化建模（最小二乘线性回归、逻辑回归、损失函数、梯度下降法、过拟合正则化数学推导）；模型集成（加权多数表决）；深度学习（深度前馈网络、CNN 局部感受野/权重共享直观机制、深度网络学到特征表征的可视化）；高级性能评估（分类高级指标、回归指标、混淆矩阵、交叉验证、超参数网格/随机搜索）；强化学习全景概览；重大跨行业工程落地剖析 & 应用/开发/评估 & 领域核心 & \multirow{2}{*}{4} \\
\cline{2-4}
& 机器学习伦理深化（真实案例研究、差分隐私与数据隐私保护、算法公平性度量） & 解释/评估 & 领域核心 & \\
\midrule
\multirow{4}{*}{\makecell[l]{人工智能/应用与\\社会影响}} 
& 跨学科重大领域实战应用（智慧医疗、可持续发展、社交网络、金融经济、教育、机器人等，抽象为 AI 求解问题、数据可用性与清洗、数据/算法/评估偏见防范） & 解释/评估 & 计算机核心 & \multirow{2}{*}{2} \\
\cline{2-4}
& 已部署前沿深度生成式模型（深度图像扩散生成模型、大语言模型 LLM 的机制、落地赋能与事实幻觉/逻辑脆弱性失效短板）；AI 深刻社会影响（伦理、算法公平性、信任度与决策可解释性） & 解释/评估 & 计算机核心 & \\
\cline{2-5}
& 深入剖析一个或多个额外重大应用垂直领域；大语言模型与深度生成模型的工业实战前沿与高安全场景落地边界；社会治理议题（伦理责任、防歧视公平治理、可信度、因果可解释性）的纵深批判性研讨 & 解释/评估 & 领域核心 & 2 \\
\bottomrule
\end{tabularx}
\endgroup
\clearpage
\twocolumn

\chapter{数据管理与数据库}

\section{领域导言与演进背景}

\begin{preamblebox}[知识领域概述与时代演进]
计算机科学的各个细分领域都可以精炼地概括为“研究面向特定目标的算法与数据结构”。对于数据管理 (DM)领域而言，这一目标正是：“\textbf{高效处理与组织持久化数据集——且这些数据集的规模往往过于庞大，无法完全容纳于计算机主存储器（内存）之中}”。

自 20 世纪 70 年代中期以来，本领域的专业教学与学术研究几乎排他性地聚焦于关系数据库系统。根据各高校不同的办学定位与培养侧重，以往学生在不同比例上系统修读了以下四个经典支柱：
\begin{enumerate}
    \item \textbf{数据建模与数据库设计}：例如实体-联系数据模型（E-R 模型）、关系模型、函数依赖与模式规范化范式理论。
    \item \textbf{声明式查询构建}：例如形式化关系代数、SQL 结构化查询语言。
    \item \textbf{查询处理与执行优化}：例如底层物理索引结构（B+ 树、静态与动态哈希）、核心查询物理执行算法（外部归并排序、选择、投影、多种表连接算法）、基于代价的查询优化器（等价代数变换、物理访问路径与索引选型）。
    \item \textbf{数据库管理系统（DBMS）内核实现}：例如并发控制与两阶段封锁协议、事务管理与 ACID 特性、内存缓冲池置换与脏页调度管理。
\end{enumerate}

\textbf{现代计算机专业人才培养定位的深刻变革：}\\
今天的行业与学术界对绝大多数计算机专业毕业生的核心期待，已明显转向要求其具备扎实的“\textbf{DBMS 高级用户与数据应用架构师}”素养，而非人人去从零编写单机数据库内核。这些核心素养首要包括数据建模与高效查询构建能力：即能够面对现实业务中纷繁杂乱、非结构化的原始数据集合，运用现代 DBMS 建立严谨规范的数据模式，并通过编程式接口与高级查询语言实现亚秒级的低延迟检索与一致性更新。

此外，现代计算机专业毕业生必须深入研读以下两大新兴维度：
\begin{itemize}
    \item \textbf{数据在现代组织与社会治理中的全生命周期角色}：
    \begin{itemize}
        \item \textbf{数据生命周期}：数据生成与采集$\to$ 处理与清洗$\to$ 审查与报表呈现$\to$ 留存、归档与检索$\to$ 销毁与销户清理。
        \item \textbf{数据资产的社会、伦理与法律维度}：数据爆炸带来的规模效应、数据隐私权益、内生合规与隐私设计、去标识化与匿名化技术、数据所有权与数据资产边界、数据可靠性与防篡改审计、全栈数据库安全，以及算法应用中预期之外的次生歧视或伦理偏见。
    \end{itemize}
    \item \textbf{超越传统单机关系系统的现代存储与大数据分析前沿技术}：
    \begin{itemize}
        \item 正在显著增强、补充甚至在特定高并发场景下替代传统关系系统的 NoSQL 数据库（如 JSON/XML 文档存储、高性能键值存储、图数据库等）；
        \item 云原生弹性分布式数据库；
        \item MapReduce 分布式海量批处理范式；
        \item 面向现代数据科学与分析工程的数据帧与自动化数据治理管道。
    \end{itemize}
\end{itemize}

\textbf{数据管理生态中的多元专业工程角色：}\\
大纲充分认识到参与现代数据管理链条的多元化分工角色：
\begin{itemize}
    \item \textbf{产品功能工程师}：在业务系统一线同时驾驭 SQL 关系数据库与 NoSQL 操作型数据库，支撑高并发核心业务功能。
    \item \textbf{数据分析工程师 / 数据工程师}：精通分析型 SQL、Python 和 Scala 编程，构建高质量数据仓库、数据湖资产与端到端 ETL/ELT 自动化数据管道。
    \item \textbf{业务分析师}：熟练使用电子表格、敏捷商业智能（BI）工具与低代码报表进行业务数据多维洞察。
    \item \textbf{数据基础架构工程师}：专注研发与维护高可靠、高吞吐的联机事务处理（OLTP）引擎、分布式存储基座与云原生查询加速器。
    \item \textbf{“全员共识” (Everyone)}：每一位生产或消费数据的计算专业人员，都必须深刻理解数据背后的法律、伦理、安全与社会责任。
\end{itemize}

\textbf{超越一切分工的核心担当——“数据托管人”：}\\
过去，数据仅仅被视为企业内部可供调配与增值的“普通资产”（信息系统管理视角）。而在今天，数据被赋予了不可剥夺的人权与法律属性：关于客户的数据本质上属于客户个体自身（或由客户保留所有权）。现代计算专业人员必须树立清醒共识：合乎法律伦理地存储、防护和使用机构数据，是从业者作为一名负责任的“数据托管人”的基本底线。

\textbf{科学理论探索与工程实践导向之间的张力调和：}\\
大纲深刻指出了计算机课程体系中“职业工程技能培养”与“作为科学理论严密探索”之间的内在张力。这一矛盾在数据管理领域尤为尖锐。例如，形式化证明函数依赖理论中阿姆斯特朗公理系统（Armstrong's Axioms）的完备性与可靠性，是经典数据库理论的学术瑰宝；然而，绝大多数毕业生在漫长的工程研发职业生涯中可能永远不会直接书写此类数学证明。类似的情况也存在于关系规范化高阶理论之中。反之，如果毕业生只会机械地运行自动化建模工具将模式切分至 BCNF 范式，只会编写低效嵌套的初级 SQL，却对底层物理索引在加速查询执行中的决定性作用一无所知，这无疑是计算机高等教育的重大失职。

为此，CS2023 在学时规划上推陈出新：将全专业通用的计算机核心学时精简为 9 学时，而将大量底层进阶机制划入领域核心（20.5 学时）。这种结构赋予了各高校灵活定制的自适应空间：对于偏向系统底层的培养方案，可深入钻研 OLTP 存储引擎、事务可串行化与故障恢复机制；而对于偏向应用开发与数据科学的方案，则可侧重既有大型分布式数据库的编程式访问、现代数据建模与高级查询优化。

\textbf{计算机高等教育的历史反讽：}\\
长期以来，计算机专业学生在本科期间的大部分时间里，都在学习通过平面文本文件读写或控制台交互式 I/O 来处理数据；然而在象牙塔外的真实产业世界中，绝大多数真实数据（以绝对压倒性的体量）都是通过应用程序编程接口（API、数据库驱动与 ORM）从数据库中访问与持久化的。CS2023 强烈呼吁高校应在程序设计早期尽早引入编程式数据库访问，并在整个大学课程体系中持续强化这一现代工程习惯。

\textbf{与网络安全 (SEC) 及社会伦理 (SEP) 的深度交织与原点地位：}\\
尽管数据管理在学科分类上与网络安全（SEC）及社会、伦理与专业素养（SEP）正交，但现实中它恰恰是安全与伦理的“风暴原点”。在设计持久化存储系统时，“什么数据应当被采集和存储”必须首先经受严格的法律合规与伦理道德审查：是否存在隐私侵蚀？数据安全防护体系是否固若金汤？
\end{preamblebox}

\section{核心学时分配建议}

本知识领域在“向日葵模型”下的核心学时推荐如下表所示：

\begin{table}[H]
\centering
\setlength{\tabcolsep}{3.5pt}\fontsize{8pt}{9.6pt}\selectfont\begin{tabularx}{\columnwidth}{Y c c}
\toprule
\textbf{知识单元} & \textbf{计算机核心} & \textbf{领域核心} \\
\midrule
数据的作用与治理 & 2 & -- \\
数据库系统核心概念 & 2 & 0.5 \\
数据建模 & 2 & 3 \\
关系数据库与规范化理论 & 1 & 3 \\
查询构建与 SQL 语言 & 2 & 4 \\
查询处理与执行优化 & -- & 4 \\
数据库管理系统内核实现 & -- & 4 \\
NoSQL 系统与非关系型存储 & -- & 2 \\
数据安全与隐私保护 & 待定 & 待定 \\
数据分析与大数据处理 & 待定 & 待定 \\
分布式数据库与云计算 & \emph{拓展选修} & \emph{拓展选修} \\
半结构化与非结构化数据库 & \emph{拓展选修} & \emph{拓展选修} \\
\midrule
\textbf{合计} & \textbf{9 学时} & \textbf{20.5 学时$^*$} \\
\bottomrule
\end{tabularx}
\begin{flushleft}\footnotesize
\vspace{-2pt}
$^*$注：大纲文本总计标记为 25 学时，其中包含为制定中的“数据安全与隐私”及“数据分析”模块预留的 4.5 个核心学时。
\end{flushleft}
\end{table}

\section{知识单元与详细大纲}

\subsection{数据的作用与治理}
\textbf{【学时建议】}：2 计算机核心学时。

本单元系统阐述数据在现代社会与企业组织中的战略资产地位，全面梳理数据的完整生命周期，并深入剖析数据采集、留存与使用过程中的法律、伦理与合规边界。

\subsubsection{计算机核心主题}
\begin{itemize}
    \item \textbf{数据生命周期}：
    \begin{itemize}
        \item 数据生成与采集：业务产生、传感器采集、日志埋点与跨系统汇聚。
        \item 处理与清洗：数据验证、异常值清洗、格式转换与丰富增强。
        \item 审查与报表呈现：商业智能展现、审计追踪与数据质量校验。
        \item 留存、归档与检索：冷热数据分级存储、合规归档与高效索引检索。
        \item 销毁与销户清理：安全删除机制、加密擦除、被遗忘权与不可逆物理擦除。
    \end{itemize}
    \item \textbf{数据集合的社会与法律伦理维度}：
    \begin{itemize}
        \item 数据规模化风险：海量数据集中带来的攻击面扩大与系统性风险。
        \item 数据隐私权益：个人身份信息（PII）保护、隐私泄漏防范。
        \item 内生合规与隐私设计：在数据库模式设计之初内嵌合规约束（如 GDPR、HIPAA 等）。
        \item 匿名化与去标识化：泛化、置换、差分扰动及抗再标识攻击能力。
        \item 数据所有权与资产边界：用户数据所有权、平台托管权与数据资产变现合法性。
        \item 数据可靠性与可信度：数据真实性、防篡改机制、数据溯源。
        \item 预期用途与非预期次生风险：数据二次利用、算法画像歧视与监控风险防范。
    \end{itemize}
\end{itemize}

\subsubsection{示范性学习成果}
\begin{enumerate}
    \item \textbf{[解释]} 阐明数据生命周期各核心阶段的划分逻辑，并列举确保各阶段数据完整性与合规性的工程措施。
    \item \textbf{[评估]} 针对给定的实际业务应用，全面评估其在海量数据采集与长期持久化留存中的隐私、法律与伦理风险，撰写合规防范白皮书。
\end{enumerate}

\subsection{数据库系统核心概念}
\textbf{【学时建议】}：2 计算机核心学时，0.5 领域核心学时。

本单元奠定数据库系统最本质的理论概念与体系结构认知，剖析数据库相对于传统操作系统的文件系统在并发、持久性与逻辑抽象层级上的革命性跨越。

\subsubsection{计算机核心主题}
\begin{itemize}
    \item \textbf{数据库系统的用途与核心优势}：消除数据冗余、保障数据一致性、提供原子并发控制、强故障恢复能力与声明式数据独立性。
    \item \textbf{数据库系统总体架构与组件划分}：存储管理器（文件组织、缓冲池管理）、查询处理器（解析器、优化器、执行引擎）、事务管理器（并发控制、恢复管理）。
    \item \textbf{DBMS 核心功能模块设计原理}：查询执行机制、事务一致性维护、缓冲池置换策略、底层物理访问路径。
    \item \textbf{三级模式体系结构与数据独立性}：外模式（视图层）、模式（逻辑层）与内模式（物理层）；物理数据独立性（物理存储变化不影响逻辑模式）与逻辑数据独立性（逻辑模式改变不破坏外部视图）。
    \item \textbf{声明式查询语言的核心思想}：关注“要获取什么数据”（What）而非“如何计算与获取数据”（How）。
    \item \textbf{事务管理基本概念}：ACID 特性（原子性 Atomicity、一致性 Consistency、隔离性 Isolation、持久性 Durability）的工程意义。
    \item \textbf{规范化理论初步}：消除重复与冗余数据的动因，更新异常、插入异常与删除异常的本质。
    \item \textbf{大规模海量数据管理策略}：水平扩展理念、NoSQL 系统破局思路、MapReduce 分布式并行处理范式。
    \item \textbf{针对纯增删改查业务应用的支持模式}：轻量级接口封装与高效对象映射。
    \item \textbf{分布式与云原生数据库系统概览}：弹性计算与存储分离架构、跨机房复制基础。
    \item \textbf{结构化、半结构化与非结构化数据库的分类与特征对比}。
\end{itemize}

\subsubsection{领域核心主题}
\begin{itemize}
    \item \textbf{流数据与现代流式数据库系统}：支持持续结构化流式内容摄入、滑动窗口聚合与亚秒级响应的流数据库架构。
\end{itemize}

\subsubsection{示范性学习成果}
\begin{enumerate}
    \item \textbf{[解释]} 辨析并准确陈述使用专业数据库系统相较于操作系统平面文件系统所提供的至少四项决定性优势。
    \item \textbf{[解释]} 列举关系数据库系统的核心架构组件，并清晰描述各组件的职责边界。
    \item \textbf{[应用]} 完整追踪一条典型 SQL 查询在数据库管理系统各组件之间解析、校验、优化、生成物理计划与检索磁盘的执行全过程。
    \item \textbf{[评估]} 论证数据独立性（尤其是物理数据独立性）对现代大型软件系统长期解耦演进的核心价值。
    \item \textbf{[开发]} 在 SQL 中熟练编写基础的多表“选择-投影-连接 (SPJ)”查询。
    \item \textbf{[解释]} 详细阐明确保事务正确执行的核心 ACID 四大特性及其违背时可能产生的破坏性后果。
    \item \textbf{[解释]} 阐述消除数据库中重复冗余数据的深层优势（存储优化、更新开销降低与数据不一致规避）。
    \item \textbf{[解释]} 概述 MapReduce 如何利用“数据不动计算动”与数据局部性原理实现海量数据的高效并行处理。
    \item \textbf{[评估]} 权衡对比结构化、半结构化与非结构化数据在数据模式灵活性、查询表达能力与存储效率上的本质差异。
\end{enumerate}

\subsection{数据建模}
\textbf{【学时建议】}：2 计算机核心学时，3 领域核心学时。

数据建模是将现实世界复杂的业务实体及其相互联系抽象转化为计算机可精确理解与操作的逻辑结构的关键桥梁。

\subsubsection{计算机核心主题}
\begin{itemize}
    \item \textbf{数据建模基础理论}：概念模型、逻辑模型与物理模型的划分与映射转换过程。
    \item \textbf{关系数据模型}：关系（表）、元组（行/记录）、属性（列/字段）、域（取值范围）、关系模式与关系实例的形式化定义。
\end{itemize}

\subsubsection{领域核心主题}
\begin{itemize}
    \item \textbf{概念数据模型构建}：
    \begin{itemize}
        \item 实体-联系模型（E-R 模型）：实体集、弱实体集、属性（单值/多值/复合/派生）、联系集及其度数（一元/二元/多元）。
        \item 映射基数与参与约束：$1:1$、$1:n$、$n:m$ 关联关系以及完全参与/部分参与约束。
        \item 统一建模语言（UML）类图在概念数据建模中的表达与转换。
    \end{itemize}
    \item \textbf{半结构化数据模型}：利用模式定义语言（如 DTD、XML Schema 或 JSON Schema）对半结构化文档进行严格的类型约束与层级结构规范化定义。
\end{itemize}

\subsubsection{拓展选修主题}
\begin{itemize}
    \item \textbf{电子表格数据模型}：二维网格表格计算模型与公式依赖拓扑。
    \item \textbf{面向对象数据模型}：对象标识符（OID）、封装性、继承体系与对象关系映射（ORM，关联 PL 领域）；GraphQL 统一图模式定义与图查询。
    \item \textbf{国际 SQL 标准演进前沿}：跟踪 SQL:2023 及后续标准最新引入的多维数组、属性图查询（SQL/PGQ）等新特性。
    \item \textbf{专用领域数据建模专题}：
    \begin{itemize}
        \item 时序数据建模：时间戳对齐、时序聚合、窗口多路连接与时间衰减归档。
        \item 图数据建模：节点属性、边属性、异构图拓扑建模与路径遍历。
        \item 物化视图与高维近似数据结构：物化视图维护策略；专用概率统计数据结构（如用于超大海量基数估计的 HyperLogLog、位图 Bitmap 等）。
        \item 性能架构警示：深入剖析为何直接对海量“原始高频时序数据”执行日均价、日去重 UV 统计、日活跃度等查询极易导致严重性能雪崩，掌握预聚合与分层汇总模式。
        \item 地理空间数据建模：点、线、多边形拓扑表示，空间参照系统与地理信息空间索引建模。
    \end{itemize}
\end{itemize}

\subsubsection{示范性学习成果}
\begin{enumerate}
    \item \textbf{[解释]} 准确阐述关系数据模型的核心构成要素（关系、元组、属性、域）及其代数本质。
    \item \textbf{[开发]} 熟练运用关系模型，通过主外键设计与中间关联表，准确建模现实业务中的 $1:1$、$1:n$ 和 $n:m$ 联系。
    \item \textbf{[解释]} 阐明实体-联系（E-R）模型及其他非关系概念模型的核心语义构件与基数约束表达机制。
    \item \textbf{[开发]} 面对给定的真实复杂业务环境，抽象出精确的概念数据模型（绘制标准 E-R 图或 UML 类图）。
    \item \textbf{[开发]} 针对高灵活性或松散结构的业务场景，构建规范的基于文档存储或键值存储的数据模式（如编写 JSON Schema）。
\end{enumerate}

\subsection{关系数据库与规范化理论}
\textbf{【学时建议】}：1 计算机核心学时，3 领域核心学时。

关系数据库以严密的数学集合论和一阶谓词逻辑为支撑。本单元聚焦关系完整性约束、模式映射以及运用函数依赖进行规范化设计的核心理论。

\subsubsection{计算机核心主题}
\begin{itemize}
    \item \textbf{关系完整性约束规则}：
    \begin{itemize}
        \item 实体完整性：主键非空且唯一。
        \item 参照完整性：外键引用有效性、级联操作。
        \item 键的形式化层次：超键、候选键与主键的定义与判定。
    \end{itemize}
    \item \textbf{关系数据库逻辑设计基础}：关系模式初步建立与约束声明。
\end{itemize}

\subsubsection{领域核心主题}
\begin{itemize}
    \item \textbf{概念模式向关系模式的映射转换}：强实体集、弱实体集、二元关联、多元联系、ISA 继承层次的标准化表转换算法。
    \item \textbf{物理数据库设计}：底层文件结构（堆文件、顺序文件、定长与变长记录）、块存储组织与记录指针映射。
    \item \textbf{函数依赖理论导引}：函数依赖形式化定义（$X \to Y$）、逻辑蕴涵与平凡/非平凡依赖。
    \item \textbf{关系规范化理论}：
    \begin{itemize}
        \item 关系模式分解准则：无损连接分解的判定定理；保持函数依赖的形式化检验。
        \item 核心范式体系：第一范式（1NF，属性不可分）、第二范式（2NF，消除非主属性对码的部分依赖）、第三范式（3NF，消除非主属性对码的传递依赖）与 Boyce-Codd 范式（BCNF，消除任何属性对码的部分与传递依赖）。
        \item 工业反规范化设计：针对高频读性能瓶颈，审慎引入受控冗余表、宽表与冗余字段的工程权衡与同步一致性开销。
    \end{itemize}
\end{itemize}

\subsubsection{拓展选修主题}
\begin{itemize}
    \item \textbf{高级函数依赖理论}：
    \begin{itemize}
        \item 阿姆斯特朗公理系统（自反律、增广律、传递律及其派生规则）的严密数学证明。
        \item 属性集闭包（$X^+$）的高效线性时间计算算法及其在超码判定中的应用。
        \item 最小函数依赖集 / 正规覆盖的极小化构造算法。
    \end{itemize}
    \item \textbf{高阶规范化理论与多值依赖}：
    \begin{itemize}
        \item 多值依赖 (MVD)的形式化定义与第四范式（4NF）。
        \item 连接依赖 (JD)与投射-连接范式（PJNF / 第五范式 5NF）。
        \item 关系数据库模式表示理论。
    \end{itemize}
\end{itemize}

\subsubsection{示范性学习成果}
\begin{enumerate}
    \item \textbf{[解释]} 深入阐明关系数据模型背后的本质数学特征（关系即元组集合的笛卡尔积子集）。
    \item \textbf{[解释]} 深入剖析外键、候选键与超键在唯一性保障和多表关联上的本质差异。
    \item \textbf{[解释]} 全面列举关系数据库所支持的完整性约束类型（域约束、主键约束、参照完整性、用户自定义 CHECK 约束）。
    \item \textbf{[开发]} 将包含 $1:1$、$1:n$ 和 $n:m$ 联系的概念数据模式严谨转换为无歧义的关系模式集合。
    \item \textbf{[应用]} 针对给定的关系表与物理索引，配置匹配的底层文件存储组织形式（如堆文件或聚簇文件）。
    \item \textbf{[解释]} 阐明函数依赖理论如何将“超键与候选键”的概念进行严谨的形式化推广。
    \item \textbf{[评估]} 运用代数算法（如表格判定法）严格证明或否定一个模式分解是否具备无损连接性与保持依赖性。
    \item \textbf{[评估]} 检测给定关系模式所属的最高范式等级（1NF、2NF、3NF 还是 BCNF），并指出潜在的数据更新异常。
    \item \textbf{[评估]} 深入论述在特定高吞吐工程场景下对关系表实施反规范化宽表设计的合理动机、收益及其数据同步风险。
\end{enumerate}

\subsection{查询构建与 SQL 语言}
\textbf{【学时建议】}：2 计算机核心学时，4 领域核心学时。

SQL 是与关系数据库进行人机和程序化交互的国际通用语言。本单元涵盖关系代数形式化基础、复杂查询表达、数据操纵与数据库编程式访问。

\subsubsection{计算机核心主题}
\begin{itemize}
    \item \textbf{SQL 声明式查询语句规范与构建}：
    \begin{itemize}
        \item 单表与多表查询：SELECT 子句、FROM 子句、WHERE 谓词过滤。
        \item 分组与聚合分析：GROUP BY、HAVING 过滤谓词、聚集函数（COUNT、SUM、AVG、MIN、MAX）。
        \item 多表连接技术：内连接、左/右/全外连接、自连接。
        \item 集合运算：UNION、INTERSECT、EXCEPT/MINUS。
        \item 嵌套子查询：标量子查询、集合成员测试、量化比较（ALL/ANY）、存在性测试。
    \end{itemize}
\end{itemize}

\subsubsection{领域核心主题}
\begin{itemize}
    \item \textbf{关系代数理论基础}：
    \begin{itemize}
        \item 经典集合运算：并（$\cup$）、交（$\cap$）、差（$-$）、广义笛卡尔积（$\times$）。
        \item 专门关系运算：选择（$\sigma$）、投影（$\pi$）、更名（$\rho$）、自然连接（$\bowtie$）、$\theta$-连接、条件等值连接、除法运算（$\div$）。
        \item 关系代数表达式的形式化推导与等价变换。
    \end{itemize}
    \item \textbf{SQL 全体系语言规范}：
    \begin{itemize}
        \item 数据定义语言（DDL）：CREATE/ALTER/DROP TABLE，主外键约束、唯一约束（UNIQUE）、CHECK 检查条件、非空约束与默认值。
        \item 数据更新语言（DML）：INSERT 插入、UPDATE 批量更新、DELETE 条件删除与 MERGE 操作；处理事务控制与隐式提交。
        \item 视图管理：视图定义、更新条件判定与可更新视图约束。
    \end{itemize}
\end{itemize}

\subsubsection{拓展选修主题}
\begin{itemize}
    \item \textbf{关系演算}：元组关系演算（TRC）与域关系演算（DRC）；安全关系演算表达式；Codd 等价定理（关系代数与安全关系演算的表达能力等价性）。
    \item \textbf{图形化与第四代查询环境}：示例查询 (QBE)与敏捷数据构建环境。
    \item \textbf{宿主语言调用非过程化查询的多种途径}：嵌入式 SQL（Embedded SQL）、动态 SQL、CLI/ODBC/JDBC 驱动标准、现代对象关系映射（ORM）框架的利弊权衡（关联 PL 领域）。
    \item \textbf{其他主流专用查询语言导引}：半结构化 XML 查询语言、语义网与资源描述框架图查询语言（SPARQL）、图数据库属性查询语言。
    \item \textbf{数据库端过程化扩展编程}：存储过程、用户自定义函数（UDF）、行级/语句级触发器及其级联调试。
\end{itemize}

\subsubsection{示范性学习成果}
\begin{enumerate}
    \item \textbf{[开发]} 熟练编写涵盖多表连接、分组聚合、集合并/交/差以及除法语义的高级 SQL 查询语句。
    \item \textbf{[评估]} 准确辨别一个嵌套子查询是否属于“相关子查询”，并评估其潜在的重复计算性能开销。
    \item \textbf{[应用]} 在通用宿主编程语言（如 Python、Java 等）中，通过标准数据库驱动连接池编程式执行参数化 SQL 查询，并高效迭代处理返回的元组结果集。
    \item \textbf{[开发]} 纯熟运用 SQL DDL 语法完整定义关系模式，涵盖所有主外键依赖、CHECK 业务约束以及级联更新/删除触发器。
    \item \textbf{[开发]} 编写精准、具备事务原子性保障的目标元组更新与安全批量变更 SQL 脚本。
\end{enumerate}

\subsection{查询处理与执行优化}
\textbf{【学时建议】}：4 领域核心学时。

数据库的高性能源自其底层复杂精密的物理存储索引、查询操作符算法以及基于代价的查询优化器（CBO）。本单元揭示查询引擎“黑盒”之下的物理运转机理。

\subsubsection{领域核心主题}
\begin{itemize}
    \item \textbf{核心物理索引结构}：
    \begin{itemize}
        \item B+ 树索引结构：节点结构（分支节点与叶子节点）、高扇出、自平衡分裂与合并算法、范围扫描的高性能保障。
        \item 哈希索引：静态哈希、动态可扩展哈希与线性哈希；点查优势与无法支持范围查询的物理局限。
        \item SQL 索引生命周期管理：普通索引、唯一索引、复合联合索引的最左前缀原则。
    \end{itemize}
    \item \textbf{查询操作符的物理实现算法}：
    \begin{itemize}
        \item 外部归并排序：两路/多路归并、内存缓冲区页帧划分、运行段生成与合并 I/O 代价。
        \item 选择操作的物理实现：全表顺序扫描vs 基于聚簇/非聚簇索引的高效点查与范围检索。
        \item 投影操作的物理实现：基于排序与基于哈希的重复记录消除机制。
        \item 自然连接与等值连接算法族：简单嵌套循环连接、块嵌套循环连接、排序归并连接、Grace 哈希连接。
        \item 算法复杂度与效率量化分析：以磁盘 I/O 块传输次数为度量标准的精确代价估算模型。
    \end{itemize}
    \item \textbf{查询代数等价变换规则}：选择操作尽早下推、投影操作提前消除冗余属性、连接顺序重排等启发式优化规则。
    \item \textbf{基于代价的查询优化器}：
    \begin{itemize}
        \item 物理访问路径选择。
        \item 物理执行计划生成：左深连接树与动态规划搜索空间剪枝（System R 算法）。
        \item 统计信息与选择率估算：直方图、数据基数与均匀性独立假设。
        \item 仅索引扫描计划（Index-only Plans / 覆盖索引 Covering Index）：免除回表物理 I/O 开销的终极优化方案。
    \end{itemize}
    \item \textbf{数据库系统性能调优与索引选型}：
    \begin{itemize}
        \item 索引对系统性能的双刃剑效应：加速只读查询吞吐 vs 显著拖慢高频插入/更新/删除操作并占用磁盘空间的权衡。
    \end{itemize}
\end{itemize}

\subsubsection{示范性学习成果}
\begin{enumerate}
    \item \textbf{[解释]} 阐明 B+ 树与动态哈希索引的内部拓扑组织结构，解释为什么关系数据库普遍采用 B+ 树作为默认主索引。
    \item \textbf{[开发]} 熟练编写 SQL DDL 语句创建高选择性的聚簇索引与复合索引，优化特定业务查询负载。
    \item \textbf{[解释]} 详细步进式剖析外部归并排序、去重投影、排序归并连接、散列哈希连接以及块嵌套循环连接的完整运行流程。
    \item \textbf{[评估]} 针对上述查询物理算法，根据分配的可用内存页数与关系表数据块规模，精确推导并计算其总磁盘 I/O 请求代价。
    \item \textbf{[应用]} 运用关系代数等价变换公理重写复杂 SQL 逻辑代数树，使其转化为适合流水线执行的物理执行计划。
    \item \textbf{[评估]} 根据字段统计直方图与谓词条件，精确计算单表过滤及两表等值连接的选择率估计值。
    \item \textbf{[开发]} 阐述如何通过在复合索引中追加包含字段，使高频查询能够直接在索引树上完成所有字段投影（覆盖索引扫描）。
    \item \textbf{[评估]} 针对复杂现实业务场景下的读写负载特征，权衡并决策出支持全套高频查询所需保留的最精炼索引组合。
\end{enumerate}

\subsection{数据库管理系统内核实现}
\textbf{【学时建议】}：4 领域核心学时。

DBMS 内核是系统软件设计的集大成者。本单元剖析内存缓冲池置换、并发事务的隔离性与调度可串行化理论，以及保障持久性与原子性的故障恢复子系统。

\subsubsection{领域核心主题}
\begin{itemize}
    \item \textbf{数据库缓冲池管理}：
    \begin{itemize}
        \item 内存缓冲区页帧组织、脏页标记与固定计数器。
        \item 经典置换策略：最近最少使用（LRU）、时钟算法、LRU-K 以及扫描抗击防护机制。
        \item 脏页异步后台刷盘调度与预读优化。
    \end{itemize}
    \item \textbf{事务处理与并发语义}：
    \begin{itemize}
        \item 事务隔离级别全解析：读未提交、读已提交、可重复读与可串行化。
        \item 并发读写异常本质：脏读、不可重复读与幻读。
        \item 事务 ACID 四大特性的底层技术支撑映射。
        \item 可串行化理论：调度、冲突等价、冲突可串行化判定准则与拓扑排序优先图。
    \end{itemize}
    \item \textbf{并发控制机制}：
    \begin{itemize}
        \item 封锁技术：排他锁（X 锁）与共享锁（S 锁）、意向锁（意向共享 IS、意向排他 IX）。
        \item 两阶段封锁协议（2-Phase Locking, 2PL）：生长阶段与收缩阶段，严格两阶段锁（Strict 2PL）防范级联回滚，强两阶段锁（Rigorous 2PL）。
        \item 死锁应对策略：死锁预防（Wait-Die 与 Wound-Wait 时间戳算法）、死锁检测（等待图 Wait-For Graph 周期环检测）与超时回滚。
    \end{itemize}
    \item \textbf{恢复管理器与缓冲池交互}：
    \begin{itemize}
        \item 缓冲池刷盘管理策略：STEAL/NO-STEAL（是否允许未提交事务将脏页刷盘）与 FORCE/NO-FORCE（是否要求提交事务在返回前将所有脏页完全落盘）。
    \end{itemize}
\end{itemize}

\subsubsection{拓展选修主题}
\begin{itemize}
    \item \textbf{高级并发控制协议}：
    \begin{itemize}
        \item 乐观并发控制 (OCC)：读取阶段、有效性确认验证阶段与写入阶段。
        \item 时间戳排序协议 (T/O)与多版本并发控制（MVCC，快照隔离 Snapshot Isolation 的底层版本链与 GC）。
    \end{itemize}
    \item \textbf{故障恢复子系统深度实现}：
    \begin{itemize}
        \item 预写日志技术 (WAL)：日志先行核心原则，日志序列号（LSN，PageLSN 与 FlushedLSN 判定）。
        \item 工业级 ARIES 故障恢复算法体系：
        \begin{itemize}
            \item 分析阶段：从检查点重构活跃事务表（ATT）与脏页表（DPT）。
            \item 重做阶段：从最早未落盘脏页开始重放历史，恢复系统至崩溃瞬间状态。
            \item 撤销阶段：反向回滚未提交事务，生成补偿日志记录（CLR）防止恢复期间再次崩溃。
        \end{itemize}
    \end{itemize}
\end{itemize}

\subsubsection{示范性学习成果}
\begin{enumerate}
    \item \textbf{[解释]} 阐明 DBMS 缓冲池管理器的内存分页管理机制、缓冲置换算法与页固定机制。
    \item \textbf{[解释]} 深入阐释一个正确的事务管理器所必须提供的 ACID 核心特性，以及数据库内核如何通过并发控制与 WAL 日志保障这四项特性。
    \item \textbf{[解释]} 概述可串行化的数学理论基石，利用事务读写冲突规则构建优先图，准确判定一个调度是否具备冲突可串行化。
\end{enumerate}

\subsection{系统与非关系型存储}
\textbf{【学时建议】}：2 领域核心学时。

互联网高并发、海量弹性扩展及半结构化数据的爆发催生了 NoSQL 运动。本单元探讨关系模型以外的替代性数据组织形态及其在分布式环境下的设计权衡。

\subsubsection{领域核心主题}
\begin{itemize}
    \item \textbf{为何需要 NoSQL 系统？}：
    \begin{itemize}
        \item 传统关系模型与现代应用程序对象模型之间的“阻抗失配”。
        \item 传统单机 RDBMS 在面对超大规模海量读写时的横向水平扩展物理瓶颈。
        \item 动态无模式或松散模式对敏捷业务迭代的强力支撑。
    \end{itemize}
    \item \textbf{两大主流非关系型数据模型}：
    \begin{itemize}
        \item 键值数据模型：极简的 PUT/GET/DELETE 原语、极高读写吞吐与内存化加速（如 Redis）。
        \item 文档数据模型：以 JSON/BSON 为核心组织单元的层次嵌套半结构化表达（如 MongoDB）。
    \end{itemize}
\end{itemize}

\subsubsection{拓展选修主题}
\begin{itemize}
    \item \textbf{非关系存储底层引擎实现}：基于 LSM 树（Log-Structured Merge-tree，MemTable、SSTable、Compaction）的高吞吐写优化存储引擎。
    \item \textbf{分布式分布模型}：
    \begin{itemize}
        \item 数据分片：范围分片、哈希分片与一致性哈希路由。
        \item 数据副本复制：单主多从同步/异步复制、无主对等复制。
    \end{itemize}
    \item \textbf{分布式一致性模型与理论极限}：
    \begin{itemize}
        \item CAP 定理：一致性、可用性、分区容错性的三者权衡边界。
        \item BASE 最终一致性模型：基本可用、软状态、最终一致。
        \item Quorum 一致性读写仲裁协议（$N, W, R$ 参数配置与强/弱一致性切换）。
    \end{itemize}
    \item \textbf{分布式并行处理范式}：MapReduce 分布式计算范式、多阶段与增量 MapReduce。
    \item \textbf{工业案例研究与工程反模式}：
    \begin{itemize}
        \item 云原生对象存储系统（如 AWS S3 等海量非结构化资产底座）。
        \item 图数据库系统（Graph Databases，复杂社交拓扑与风控网络）。
        \item 深度工程反思：“\textbf{何时绝不应使用 NoSQL}”——识别缺乏 ACID 强一致性、涉及多表强关联金融级结算场景中的 NoSQL 滥用反模式。
    \end{itemize}
\end{itemize}

\subsubsection{示范性学习成果}
\begin{enumerate}
    \item \textbf{[解释]} 结合真实的工程业务场景，清晰论证在特定需求下选择使用 NoSQL 替代传统 RDBMS 的决定性用例与权衡动因。
    \item \textbf{[解释]} 深入阐释键值存储与文档数据库背后的本质特征、数据组织拓扑结构及其在读写吞吐上的核心优势。
\end{enumerate}

\subsection{数据安全与隐私保护}
\textbf{【学时建议】}：学时待定（CS2023 框架正与 SEC 知识领域深度统筹规划）。

数据管理是现代网络空间安全与个人隐私保护的最核心阵地。本单元涵盖从物理静态存储到逻辑访问层的全链路安全加固。

\subsubsection{大纲规划重点议题}
\begin{itemize}
    \item \textbf{数据加密与完整性审计}：静态透明数据加密 (TDE)、传输层 TLS/SSL 加密、密钥生命周期管理与列级敏感字段端到端加密。
    \item \textbf{细粒度访问控制体系}：基于角色的访问控制（RBAC）、基于属性的动态访问控制（ABAC）、数据库行级安全 (RLS)与动态数据脱敏。
    \item \textbf{数据库隐私保护计算前沿}：差分隐私（Differential Privacy）在统计聚合查询中的拉普拉斯/高斯噪声注入机制；多方安全计算（SMPC）与同态加密在分布式联合查询中的应用。
    \item \textbf{数据匿名化工程标准}：$k$-匿名（$k$-Anonymity）、$l$-多样性（$l$-Diversity）与 $t$-贴近度（$t$-Closeness）在防止准标识符重识别攻击中的实现。
    \item \textbf{数据库抗攻击与审计防线}：SQL 注入漏洞的编程式根除（强参数化查询与预编译语句）、数据库操作审计日志与防篡改区块链存证。
\end{itemize}

\subsection{数据分析与大数据处理}
\textbf{【学时建议】}：学时待定（CS2023 框架正与 ACM 数据科学课程体系全面对齐）。

从联机事务处理（OLTP）走向联机分析处理（OLAP）与智能决策支持，是现代数据技术演进的另一条主线。

\subsubsection{大纲规划重点议题}
\begin{itemize}
    \item \textbf{数据仓库与联机分析处理 (OLAP)}：星型模型、雪花模型与事实表/维度表设计；多维数据集与上卷、下钻、切片和切块。
    \item \textbf{面向海量分析的物理存储革命}：列式存储格式（Columnar Storage，如 Parquet、ORC）、列式压缩算法（字典编码、Run-Length 游程编码、位图编码）与向量化查询执行引擎。
    \item \textbf{现代数据工程治理与管道}：自动化 ETL/ELT 管道架构、数据湖与湖仓一体架构；数据血缘追踪与数据质量度量。
    \item \textbf{现代计算抽象与分布式框架}：基于 DataFrame 的声明式数据转换、分布式内存并行分析引擎（如 Apache Spark）。
\end{itemize}

\subsection{分布式数据库与云计算}
\textbf{【学时建议】}：拓展选修。

面对全球跨地域容灾与海量伸缩需求，现代数据库已经全面拥抱分布式与云原生架构。

\subsubsection{拓展选修主题}
\begin{itemize}
    \item \textbf{分布式数据库系统}：
    \begin{itemize}
        \item 分布式数据存储：水平分段、垂直分段、混合分段与全局目录管理。
        \item 分布式查询处理与优化：分布式查询重写、半连接算法与最小化网络传输开销策略。
        \item 分布式事务原子性协议：两阶段提交协议 (2PC)的协调者/参与者状态转移、单点阻塞故障与三阶段提交（3PC）演进。
        \item 同构与异构分布式数据库集成：联邦数据库系统与跨异构源中间件。
        \item 客户端-服务器分布式计算拓扑（跨领域关联 系统基础/计算范式）。
    \end{itemize}
    \item \textbf{并行数据库系统架构}：
    \begin{itemize}
        \item 经典物理架构分类：共享内存、共享磁盘与无共享架构的伸缩性与容错性对比。
        \item 性能扩展度量：加速比与扩展比的理论上限。
        \item 数据多副本高可用同步与弱一致性容灾模型。
    \end{itemize}
\end{itemize}

\subsection{半结构化与非结构化数据库}
\textbf{【学时建议】}：拓展选修。

现代人工智能与多媒体应用的繁荣带来了对高维向量、复杂嵌套文档与空间数据的爆发式存储检索需求。

\subsubsection{拓展选修主题}
\begin{itemize}
    \item \textbf{高维向量化非结构化数据与向量存储}：
    \begin{itemize}
        \item 非结构化文本与多媒体特征提取：基于 N-gram 的 TF-IDF 文本特征向量化、Word2Vec 词向量嵌入。
        \item 专用向量数据库与多维数组数据库：支持高维向量近似最近邻搜索（ANN，如 HNSW、IVF-PQ 算法）。
    \end{itemize}
    \item \textbf{半结构化文档数据深度管理（以 JSON 为核心）}：
    \begin{itemize}
        \item 物理存储引擎：嵌套半结构化数据类型的高效二进制编码与列存压缩（如 JSONB、BSON）。
        \item 专用索引加速技术：
        \begin{itemize}
            \item 跳跃索引与布隆过滤器加速块级过滤。
            \item 倒排索引与位图压缩索引。
            \item 面向半结构化地理空间数据的空间填充曲线索引（Space Filling Curve，如 Z-order 曲线、Hilbert 空间填充曲线）。
        \end{itemize}
        \item 混合负载（HTAP）下的查询处理权衡：插入吞吐、范围点查、就地更新/删除之间的工程权衡。
        \item 经典前沿系统工程剖析：PostgreSQL 的 JSONB 扩展支持、MongoDB 纯文档架构与 Snowflake 云数仓半结构化 JSON 混合查询。
    \end{itemize}
\end{itemize}

\section{专业素养与思维倾向}

数据管理不仅是一门关于算法与存储工程的技术学科，更直接承载着现代数字社会的信任基石。本领域着重塑造学生以下四项关键专业思维特质：
\begin{itemize}
    \item \textbf{专业严谨}：这是数据管理中最为核心的职业素养。尽管全球不同国家与法域对“数据所有权”的法律界定存在差异，但任何现代企业与计算专业人员对于严密、合规、专业化管理数据的敬畏之心与职业操守绝不可动摇。
    \item \textbf{责任担当}：在专业处理海量（个人敏感）数据时，承担高度社会责任至关重要。从业者必须具备两重底线：既要构筑坚如磐石的技术防线抵御未经授权的非法访问与网络攻击，更要从伦理高度坚决杜绝“看似合法合规但本质上不负责任、侵蚀公众利益的数据滥用行为”。此外，数据托管人必须确保数据在面对意外崩溃、物理灾难与不可抗力风险时能够完好无损地复原。
    \item \textbf{协作精神}：数据库管理员、架构师、安全审计员与业务数据消费者之间必须建立高度透明、无缝对接的协作机制，确保在合规授权的前提下让正确的数据以合乎规范的方式流通并赋能业务。
    \item \textbf{需求导向与敏捷响应}：数据的存储模式、保留周期与索引结构选型绝非闭门造车，而必须紧密敏捷地响应组织真实的业务需求与合规监管变化。
\end{itemize}

\section{数学基础要求}

为深入理解与推导关系代数、函数依赖与范式分解理论，学生必须牢固掌握以下核心数学前置知识：
\begin{itemize}
    \item \textbf{离散朴素集合论}：
    \begin{itemize}
        \item 集合的基本概念、元素属于与子集包含关系；
        \item 集合代数运算：并集、交集、差集；
        \item 广义笛卡尔积与元组映射；
        \item 幂集概念及其在属性集闭包计算中的应用。
    \end{itemize}
\end{itemize}

\section{交叉主题与跨学科共享概念}

数据管理领域与计算机科学各大知识领域存在着深厚而频繁的交叉共振：
\begin{itemize}
    \item \textbf{社会、伦理与专业素养 (SEP)}：数据隐私立法、GDPR/CCPA 时代合规审计、被遗忘权与算法公平性。
    \item \textbf{程序设计语言 (PL/FPL)}：声明式语法解析、类型系统、对象关系阻抗失配、ORM 映射与内存对象持久化。
    \item \textbf{系统基础 (SF)}：现代计算机系统的存储层次结构（Cache、RAM、NVMe SSD、HDD）、页式内存管理与 I/O 调度。
    \item \textbf{算法与复杂性 (AL)}：B+ 树平衡检索、多路归并外部排序、哈希冲突消解与动态规划查询优化。
    \item \textbf{并行与分布式计算 (PDC)}：数据分片路由、多副本复制协议、分布式事务提交（2PC）、分布式锁与 MapReduce 并行批处理。
    \item \textbf{专用平台开发 (SPD)}：移动端嵌入式数据库、物联网边缘时序存储。
    \item \textbf{计算机安全 (SEC)}：静态存储加密（TDE）、细粒度行级访问控制（RLS）、SQL 注入攻击防范与审计溯源。
    \item \textbf{数据挖掘与数据科学}：海量特征工程与湖仓一体分析（系统参见 ACM 数据科学核心规范）。
\end{itemize}

\section{课程体系规划与开课建议}

针对我国高校不同的办学定位与专业特色，建议采取以下三种典型的开课与课程规划路径：
\begin{enumerate}
    \item \textbf{标准计算机专业必修单门模式（“数据库系统原理与设计”，4 学分 / 64 学时）}：
    全面覆盖全部 9 个计算机核心学时以及 20.5 个领域核心内容。理论部分涵盖概念数据建模、关系代数与 SQL、函数依赖与 BCNF/3NF 规范化理论、B+ 树与哈希索引、基于代价的查询优化、事务并发控制（2PL）与 WAL 恢复机制；实验部分必须包含真实数据库应用系统的端到端设计开发，并包含针对复杂查询的索引性能调优分析。
    \item \textbf{高阶双轨递进模式（“基础数据库系统” $+$ “海量数据管理与分布式存储”）}：
    适用于高水平研究型大学与卓越工程师培养计划。第一门课专注经典关系理论与 SQL 开发应用；第二门进阶选修课系统探讨 NoSQL 存储引擎、分布式一致性协议、分布式事务（2PC）、列式存储与现代云原生数据仓库。
    \item \textbf{先行贯通融合模式}：
    破除传统“前两年完全用纯文本读写文件，直到大三才接触数据库”的脱节弊端。强烈建议在大学第一年《程序设计基础》（SDF）课程中尽早引入轻量级嵌入式数据库（如 SQLite）的编程式访问，培养学生使用标准 SQL 进行数据持久化存储的现代工程素养。
\end{enumerate}

\section{真实工程能力任务规范}

本领域制定了 8 个涵盖系统开发、风险审计、架构设计、性能调优与故障恢复的典型真实能力任务：

\begin{taskbox}[任务 1：为业务目标实现数据库编程式访问与数据处理]
\textbf{【任务情境】}：现代软件系统需要通过应用程序代码连接数据库，执行自动化数据检索与业务计算，而非依赖人工交互式控制台操作。\\
\textbf{【能力陈述】}：能够通过标准编程式接口（如 JDBC/ODBC、ORM 或原生数据库驱动）安全连接数据库，构造参数化查询，并在宿主编程语言中遍历返回的关系元组流，执行相应的数据处理与业务计算逻辑。\\
\textbf{【能力领域】}：软件工程、系统基础 \quad \textbf{【能力单元】}：系统开发、测试验证\\
\textbf{【关联合约/前置知识】}：数据管理/查询构建、社会伦理/社会语境、社会伦理/公平多元与包容、社会伦理/伦理分析方法、社会伦理/专业素养与工程伦理\\
\textbf{【核心要求】}：计算机核心 \quad \textbf{【达标等级】}：\textbf{开发}
\end{taskbox}

\begin{taskbox}[任务 2：评估新老应用数据收集与留存的社会及伦理风险（白皮书撰写）]
\textbf{【任务情境】}：组织拟上线一款涉及大规模用户敏感数据的新系统（或升级既有业务系统），需在立项之初评估其法律合规与伦理影响。\\
\textbf{【能力陈述】}：撰写专业决策白皮书，全面识别多方利益相关者，深入评估数据采集与长期持久化留存的长期潜在后果；审慎分析非预期数据滥用及潜在数据泄露事件所带来的危害与法律风险，并提出合规缓解方案。\\
\textbf{【能力领域】}：系统管理、理论基础 \quad \textbf{【能力单元】}：系统评估、数据治理、社会议题调适\\
\textbf{【关联合约/前置知识】}：社会伦理/社会语境、社会伦理/伦理分析方法、社会伦理/隐私与公民自由、社会伦理/专业伦理、社会伦理/安全法规与计算机犯罪、社会伦理/公平多元与包容、数据管理/数据的作用、安全隐私/安全基础\\
\textbf{【核心要求】}：计算机核心 \quad \textbf{【达标等级】}：\textbf{评估与解释}
\end{taskbox}

\begin{taskbox}[任务 3：面向新应用场景评估并决策底层数据存储架构模型]
\textbf{【任务情境】}：面对新型业务形态（如高并发社交互动、物联网大规模时序流水或金融级强事务支付），需在关系模型与非关系模型间做出科学架构选型。\\
\textbf{【能力陈述】}：综合剖析应用的数据读写特征（读写比、ACID 强一致性需求 vs BASE 最终一致性、水平扩展性诉求与模式可变性），在传统关系数据库、键值存储、文档数据库、图数据库之间进行定量权衡，选定最适合的数据模型并提供详尽的技术论证。\\
\textbf{【能力领域】}：应用架构、理论基础 \quad \textbf{【能力单元】}：系统评估、架构设计\\
\textbf{【关联合约/前置知识】}：数据管理/数据的作用、数据管理/数据库系统核心概念、数据管理/数据建模、数据管理/关系数据库、数据管理/NoSQL 系统、社会伦理/安全法规与计算机犯罪\\
\textbf{【核心要求】}：计算机核心与领域核心 (计算机核心 / 领域核心) \quad \textbf{【达标等级】}：\textbf{评估与解释}
\end{taskbox}

\begin{taskbox}[任务 4：基于索引调优与查询计划优化数据库应用执行性能]
\textbf{【任务情境】}：企业核心业务数据库在流量高峰期响应缓慢，关键业务报表与核心高频交易查询发生严重 I/O 阻塞。\\
\textbf{【能力陈述】}：分析数据库执行计划与查询物理访问路径，评估既有索引的利用率与更新维护代价；精准创建适宜的高效索引结构（如复合 B+ 树索引、动态哈希索引、覆盖索引/仅索引扫描）或移除低效冗余索引，显著降低高频重要查询的 I/O 次数与响应延迟。\\
\textbf{【能力领域】}：软件工程、应用开发 \quad \textbf{【能力单元】}：系统设计、系统重构与性能改进\\
\textbf{【关联合约/前置知识】}：数据管理/查询处理、数据管理/DBMS 内核实现、并行分布式/并行算法分析与编程\\
\textbf{【核心要求】}：领域核心 \quad \textbf{【达标等级】}：\textbf{评估与开发}
\end{taskbox}

\begin{taskbox}[任务 5：基于视图与访问控制机制实现数据安全防护与权限隔离]
\textbf{【任务情境】}：跨部门业务系统需共享底层数据库表，但必须严格防范越权访问与未经授权的敏感隐私字段泄露。\\
\textbf{【能力陈述】}：根据最小特权原则与数据治理合规要求，设计并创建精细化数据库视图，结合行级/列级访问控制策略和用户权限角色划分，从数据访问层彻底杜绝非授权访问。\\
\textbf{【能力领域】}：系统管理 \quad \textbf{【能力单元】}：运维管理、系统管理\\
\textbf{【关联合约/前置知识】}：数据管理/数据的作用、数据管理/关系数据库、数据管理/查询处理、社会伦理/安全法规与计算机犯罪、社会伦理/专业伦理、社会伦理/隐私与公民自由、安全隐私/安全基础\\
\textbf{【核心要求】}：领域核心 \quad \textbf{【达标等级】}：\textbf{开发}
\end{taskbox}

\begin{taskbox}[任务 6：在系统灾难或突发中断（如断电/崩溃）后快速恢复数据库至安全状态]
\textbf{【任务情境】}：生产数据库服务器突发断电、操作系统宕机或硬件故障，导致内存缓冲池未刷盘事务中断，需保障数据一致性并迅速重新上线。\\
\textbf{【能力陈述】}：掌握数据库恢复管理器工作机制，基于预写日志（WAL）与检查点技术，熟练运行恢复与还原工具，执行事务的 REDO（重做）与 UNDO（撤销）流程，将数据库迅速恢复到故障前最新的持久化一致安全状态。\\
\textbf{【能力领域】}：系统基础 \quad \textbf{【能力单元】}：运维管理、系统管理\\
\textbf{【关联合约/前置知识】}：数据管理/DBMS 内核实现\\
\textbf{【核心要求】}：领域核心 \quad \textbf{【达标等级】}：\textbf{应用与实践}
\end{taskbox}

\begin{taskbox}[任务 7：面向新业务应用完成从概念建模到关系规范化存储设计]
\textbf{【任务情境】}：组织启动一项新业务应用系统的研发，需要从零构建结构清晰、无数据异常的高质量底层数据存储方案。\\
\textbf{【能力陈述】}：深入理解业务领域需求，运用实体-联系模型（E-R 图）或 UML 类图完成概念数据建模；将概念模式规范映射为关系模式（准确定义主键、候选键、外键与完整性约束），并应用函数依赖与规范化理论将其分解至 BCNF/3NF 范式，消除插入、删除与更新异常。\\
\textbf{【能力领域】}：应用架构 \quad \textbf{【能力单元】}：系统设计、模式改进\\
\textbf{【关联合约/前置知识】}：数据管理/数据的作用、数据管理/数据建模、社会伦理/社会语境、社会伦理/伦理分析方法、社会伦理/隐私与公民自由、社会伦理/专业伦理、社会伦理/安全法规与计算机犯罪、社会伦理/公平多元与包容、安全隐私/安全基础\\
\textbf{【核心要求】}：计算机核心与领域核心 (计算机核心 / 领域核心) \quad \textbf{【达标等级】}：\textbf{应用与开发}
\end{taskbox}

\begin{taskbox}[任务 8：端到端完整实现新应用数据库系统的工程构建与部署落地]
\textbf{【任务情境】}：新系统从立项到落地需要一体化贯通数据存储需求分析、合规治理、模型设计、系统部署与索引优化的全生命周期。\\
\textbf{【能力陈述】}：综合执行端到端数据库工程构建全流程：完成概念与逻辑数据建模；系统审计数据采集与存储的伦理法律风险；科学决断关系型与 NoSQL 存储选型；编写全套 DDL 脚本定义数据表、完整性约束与触发器；针对核心查询负载合理配置 B+ 树与哈希索引，完成整套生产级数据库实例的构建与验证。\\
\textbf{【能力领域】}：架构设计、软件工程、系统基础 \quad \textbf{【能力单元】}：系统开发、测试验证\\
\textbf{【关联合约/前置知识】}：数据管理/数据的作用、数据管理/数据库系统核心概念、数据管理/数据建模、数据管理/关系数据库、数据管理/NoSQL 系统、数据管理/DBMS 内核实现、SEP 各单元、安全隐私/安全基础\\
\textbf{【核心要求】}：计算机核心与领域核心 (计算机核心 / 领域核心) \quad \textbf{【达标等级】}：\textbf{开发}
\end{taskbox}

\section{指导委员会名单}

\textbf{委员会主席}：
\begin{itemize}
    \item Mikey Goldweber, 泽维尔大学
\end{itemize}

\textbf{委员会委员}：
\begin{itemize}
    \item Sherif Aly, 开罗美国大学
    \item Sara More, 约翰霍普金斯大学
    \item Mohamed Mokbel, 明尼苏达大学
    \item Rajendra Raj, 罗切斯特理工学院
    \item Avi Silberschatz, 耶鲁大学 —— \emph{注：国际知名计算机科学家，经典著作《操作系统概念》与《数据库系统概念》资深作者}
    \item Min Wei, 微软公司
    \item 向乔，厦门大学
\end{itemize}

\begin{strip}
\section{附录：核心主题与技能要求等级对照表}

\noindent\begin{tabularx}{\textwidth}{p{0.18\textwidth} Y c c c}
\toprule
\textbf{知识单元} & \textbf{核心知识点与具体主题} & \textbf{认知等级} & \textbf{分类} & \textbf{学时} \\
\midrule
\multirow{2}{*}{\makecell[l]{数据管理/数据的作用\\与治理}} 
& 数据生命周期（创建、处理、审查、归档、销毁）与全链路合规 & 评估 & 计算机核心 & \multirow{2}{*}{2} \\
& 数据集合的社会与法律伦理维度（规模效应、隐私设计、所有权） & 评估 & 计算机核心 & \\
\midrule
\multirow{3}{*}{\makecell[l]{数据管理/数据库系统\\核心概念}} 
& 数据库系统核心价值与架构组件；三级模式结构与数据独立性 & 解释 & 计算机核心 & \multirow{2}{*}{2} \\
& 事务管理、规范化消除冗余、海量数据 MapReduce 批处理机制 & 解释 & 计算机核心 & \\
& 声明式查询语言使用范式与基础 SPJ 复杂查询构建 & 开发 & 计算机核心 & \\
\cline{2-5}
& 现代结构化流数据与流式计算数据库系统 & 解释 & 领域核心 & 0.5 \\
\midrule
\multirow{2}{*}{\makecell[l]{数据管理/数据建模}} 
& 现实世界数据建模流程；关系数据模型（关系、元组、属性） & 开发 & 计算机核心 & 2 \\
\cline{2-5}
& 概念数据建模（E-R 图、UML 类图）；半结构化模型 & 解释 & 领域核心 & 3 \\
\midrule
\multirow{2}{*}{\makecell[l]{数据管理/关系数据库与\\规范化理论}} 
& 实体完整性、参照完整性（外键约束）；超键与候选键定义 & 解释 & 计算机核心 & 1 \\
\cline{2-5}
& 概念模式向关系模式转换；底层物理文件与存储结构设计；函数依赖理论；规范化范式（BCNF/3NF 无损连接与保持依赖分解）；反规范化设计 & 开发 & 领域核心 & 3 \\
\midrule
\multirow{2}{*}{\makecell[l]{数据管理/查询构建与\\SQL 语言}} 
& SQL 声明式查询高级构建（单表/多表过滤、分组聚合、集合操作） & 开发 & 计算机核心 & 2 \\
\cline{2-5}
& 关系代数（并、交、差、笛卡尔积、选择、投影、连接、除法）；SQL DDL 模式完整性定义；SQL DML 更新子语言 & 开发 & 领域核心 & 4 \\
\midrule
\multirow{2}{*}{\makecell[l]{数据管理/查询处理与\\执行优化}} 
& 物理索引结构（B+ 树、哈希索引）；查询操作符算法（外部排序、选择、去重投影、连接算法族）；等价代数变换；基于代价的优化器（访问路径、选择率估算、仅索引扫描） & 解释 & 领域核心 & \multirow{2}{*}{4} \\
& 数据库物理性能调优：工作负载分析与索引选型决策（权衡读写开销） & 开发 & 领域核心 & \\
\midrule
\multirow{2}{*}{\makecell[l]{数据管理/DBMS\\内核实现}} 
& 缓冲池管理（页置换、固定机制）；事务处理与 ACID 特性保障 & 解释 & 领域核心 & \multirow{2}{*}{4} \\
& 事务隔离级别与可串行化理论；两阶段封锁协议（2PL）与死锁处理；恢复管理器与缓冲池交互机制 & 解释 & 领域核心 & \\
\midrule
\multirow{2}{*}{\makecell[l]{数据管理/NoSQL 系统与\\非关系型存储}} 
& 为何需要 NoSQL？CRUD 应用与关系数据库的阻抗失配；水平扩展动因 & 解释 & 领域核心 & \multirow{2}{*}{2} \\
& 键值数据模型与基于 JSON 的文档数据模型核心特征与适用边界 & 解释 & 领域核心 & \\
\midrule
数据管理/数据安全与隐私 & 规划中：静态加密、细粒度行级访问控制、差分隐私、脱敏防泄露 & 制定中 & 规划中 & 待定 \\
\midrule
数据管理/数据分析 & 规划中：OLAP 架构、列式存储、数据湖仓一体与现代数据管道 & 制定中 & 规划中 & 待定 \\
\bottomrule
\end{tabularx}
\end{strip}

\chapter{编程语言理论基础}

\section{领域导言与演进背景}

\begin{preamblebox}[知识领域概述与变革]
编程语言理论基础 (FPL)知识领域为理解复杂多变的现代编程语言提供了深植于离散数学与数理逻辑的坚实基石，涵盖编程语言的核心基础、工程实现机制以及严格的形式化描述方法。尽管不同编程语言因其设计范式与所面向的问题域而形态各异，且伴随社会需求演进与计算技术突破而持续迭代，但它们在深层次上共享着高度统一的计算抽象模型与程序演化逻辑。即便在处理器底层硬件与高层编程工具链日益紧密交织、系统复杂度呈指数级跃升的今天，这一底层规律依然坚如磐石。深刻领悟这些共享的核心抽象与编程范式，能够赋予计算机专业人员在面对全新编程语言与技术生态时快速洞察本质、触类旁通的卓越学习能力。

本知识领域聚焦于清晰阐明编程语言的核心概念与设计法则、语言语法规范与程序动态运行行为的形式化界定、现代语言编译系统与运行时环境的具体实现原理、系统性对比各主流编程范式的优势与局限，以及精细刻画编程语言与底层操作系统、处理器硬件之间的接口交互机理。本领域涵盖的理论法则具有普适性，超越任何具体语言的语法表象；掌握这些核心原理，不仅能帮助工程师在异构多语言技术栈之间无缝切换，更能指导开发者立足具体现实工程问题，科学权衡并选用最优的编程范式与语言工具。

在本章末尾，联合工作组给出了两套推荐的示范课程方案以展示本领域内容的教学组织：第一门是兼顾计算机核心与领域核心的基础导论课程，重点聚焦于主流编程范式的系统学习与实践对比，使学生具备足够的辨析洞察力，能够在面对不同业务场景时准确判定最适用的范式；第二门则是聚焦于编译系统工程实现、编程语言形式语义描述以及程序动态行为形式化证明的高级进阶课程。

过去十年间，上述两门课构成了该领域的经典教学模式，但这绝非唯一的实施路径。各高校完全可以根据自身培养定位进行灵活重构：例如，可在学时较短的单门课程中仅涵盖 28 学时的计算机核心内容；亦可将本领域的计算机核心知识与软件工程（SE）等相近领域的关键核心深度融合打包开课。

编程语言是程序员精确描述思想概念、构建复杂算法并对系统解决方案进行数理推演的核心媒介。在漫长的职业生涯中，计算机科学家必然会面对并驾驭大量异构语言。软件架构师与研发人员必须深入领会底层编程模型、前沿语法构造，并在支持多重互补范式的现代混合语言中做出具备前瞻性的设计决策。面对层出不穷的新型编程语言与语言特性，从业者必须掌握语言特性的形式化定义、正交组合与底层工程实现原理，以保障软件系统的执行效能与长期可维护性。要真正高效利用编程语言并洞悉其能力边界，还必须建立关于语言编译翻译、程序静态分析、运行时系统行为、内存精细管理，以及并发进程间通过消息传递、共享内存与同步原语进行协同通信的完整认知体系。此外，对于致力于设计全新领域特定语言（DSL）的系统研发者而言，熟练掌握这些语言理论基础更是不可或缺的必备素养。

\textbf{自 CS2013 以来的核心变革：}
\begin{enumerate}
    \item \textbf{学科重命名与聚焦基础理论}：将原“编程语言”正式更名为“编程语言理论基础”，旨在深刻彰显本知识领域聚焦于支撑语言演进的本质规律与底层基础理论，而非局限于特定工程语言的语法细节。
    \item \textbf{学时层级体系的重构与精细化调整}：将原 CS2013 的一级核心（Core Tier-1）与二级核心（Core Tier-2）全面重构升级为计算机核心与领域核心。所有计算机科学专业毕业生均必须掌握 计算机核心 规定的核心学时；而致力于在系统软件或理论计算机方向深化的学生，需进一步修满 领域核心 学时；其余主题则归为拓展选修。根据过去十年软件产业界与学术界的技术重心迁移，核心学时进行了针对性调整：
    \begin{itemize}
        \item \textbf{面向对象编程 (OOP)}：+1 计算机核心学时，-2 领域核心学时
        \item \textbf{函数式编程}：+1 计算机核心学时
        \item \textbf{逻辑式编程}：+3 核心学时（强化面向数据逻辑与 AI 的谓词基础）
        \item \textbf{事件驱动与响应式编程}：+2 计算机核心学时
        \item \textbf{并行与分布式计算}：+1 计算机核心学时，+2 领域核心学时
        \item \textbf{类型系统}：+1 计算机核心学时，-1 领域核心学时
        \item \textbf{语言翻译与执行}：+4 计算机核心学时，-3 领域核心学时
        \item \textbf{程序表示}：+2 领域核心学时
    \end{itemize}
    \item \textbf{知识单元更名与整合}：
    \begin{itemize}
        \item “静态分析”更名为“程序分析与分析器”，突出工具链与分析框架的工程落地；
        \item “并发与并行性”更名为“并行与分布式计算”；
        \item “运行时系统”更名为“运行时行为与系统”；
        \item 将原“基础类型系统”与“高级类型系统”合并为高度凝聚的单一知识单元“类型系统”。
    \end{itemize}
    \item \textbf{面向 2030 年代新增四大前沿知识单元}：
    \begin{itemize}
        \item \textbf{脚本编程}：新增 2 计算机核心学时，强化现代自动化流水线与数据处理脚本开发能力；
        \item \textbf{形式化开发方法学}：涵盖交互式定理证明器、模型检验与程序验证；
        \item \textbf{嵌入式计算与硬件接口}：面向边缘计算与异构加速硬件的语言接口；
        \item \textbf{编程语言基础与社会、伦理及专业素养}：探讨英语中心主义、代码可访问性、神经多样性与技术术语包容性演进。
    \end{itemize}
    \item \textbf{多范式深度融合与跨领域协同趋势}：
    \begin{itemize}
        \item 现代语言主流趋势呈现出命令式、面向对象、函数式、逻辑式以及响应式事件流的深度多范式杂交（如现代 C++、Rust、Kotlin、Scala）；
        \item 云计算与多核体系使“分布式、并发、并行与代码动态迁移”成为通用语言必备的内建范式；
        \item 面对现代海量数据管理与人工智能浪潮，基于一阶逻辑与谓词约束的逻辑式编程成为数据查询优化与智能推理的重要基础；
        \item 针对移动互联网、人机交互与普适物联网，事件驱动与嵌入式硬件接口在系统软件开发中发挥着不可替代的支撑作用。
    \end{itemize}
\end{enumerate}
\end{preamblebox}

\section{核心学时分配建议}

本知识领域在“向日葵模型”下的学时推荐如下表所示（计算机核心共计 24 学时，领域核心共计 21 学时，核心总学时 45 学时）：

\begin{table*}[htbp]
\centering\small
\setlength{\tabcolsep}{2.5pt}\footnotesize\begin{tabularx}{0.88\textwidth}{Y c c}
\toprule
\textbf{知识单元} & \textbf{计算机核心} & \textbf{领域核心} \\
\midrule
面向对象编程 & 6 & 3 \\
函数式编程 & 4 & 4 \\
逻辑式编程 & -- & 3 \\
脚本编程 & 2 & -- \\
事件驱动与响应式编程 & 2 & 2 \\
并行与分布式计算 & 3 & 2 \\
类型系统 & 3 & 4 \\
语言翻译与执行 & 4 & -- \\
程序表示 & -- & 3 \\
语法分析 & \emph{拓展选修} & \emph{拓展选修} \\
编译器语义分析 & \emph{拓展选修} & \emph{拓展选修} \\
程序分析与分析器 & \emph{拓展选修} & \emph{拓展选修} \\
代码生成 & \emph{拓展选修} & \emph{拓展选修} \\
运行时行为与系统 & \emph{拓展选修} & \emph{拓展选修} \\
高级编程构造 & \emph{拓展选修} & \emph{拓展选修} \\
语言实用学 & \emph{拓展选修} & \emph{拓展选修} \\
形式语义学 & \emph{拓展选修} & \emph{拓展选修} \\
形式化开发方法学 & \emph{拓展选修} & \emph{拓展选修} \\
嵌入式计算与硬件接口 & \emph{拓展选修} & \emph{拓展选修} \\
编程语言与社会、伦理及专业素养 & \emph{拓展选修} & \emph{拓展选修} \\
\midrule
\textbf{合计} & \textbf{24 学时} & \textbf{21 学时} \\
\bottomrule
\end{tabularx}
\end{table*}

\section{知识单元与详细大纲}

\subsection{面向对象编程}
\textbf{【学时建议】}：6 计算机核心学时，3 领域核心学时。

面向对象编程范式是现代工业级软件工程最广泛采纳的体系之一，核心依托于将系统状态与计算行为打包为高内聚的对象抽象。

\subsubsection{计算机核心主题}
\begin{itemize}
    \item \textbf{面向对象设计}：
    \begin{itemize}
        \item 职责分离与对象解耦：将系统分解为承载内部私有状态并对外暴露业务行为的离散对象。
        \item 面向领域建模的类继承层次结构顶层设计。
    \end{itemize}
    \item \textbf{类的定义机制}：内部私有状态字段、方法与构造函数的定义与生命周期管理。
    \item \textbf{继承与重写}：子类化、单继承与多重继承模型、方法重写与覆盖。
    \item \textbf{动态分派}：方法调用的晚期绑定语义规范与动态执行机理。
    \item \textbf{异常处理机制}：面向对象的结构化异常抛出、捕获、堆栈展开与清理契约。
    \item \textbf{封装与信息隐藏惯用法}：
    \begin{itemize}
        \item 类成员可见性控制（Private、Protected、Public）、数据隐藏与私有不变量维护。
        \item 仅暴露方法调用签名的纯抽象接口。
        \item 抽象基类、特质与混入组合。
    \end{itemize}
    \item \textbf{动态属性 vs 静态属性}：静态成员/方法与动态实例属性的时空开销权衡。
    \item \textbf{对象复用原则}：组合优于继承（Composition vs. Inheritance）的深度设计准则。
\end{itemize}

\subsubsection{领域核心主题}
\begin{itemize}
    \item \textbf{子类型化理论}：
    \begin{itemize}
        \item 子类型多态及强类型系统中的隐式安全向上转型。
        \item 行为替换原则 (LSP)：子类型在任何语义上下文中均必须能够无缝充当超类型的行为替代。
        \item 子类型多态（类型系统分类学）与类继承（代码实现复用机制）的深层本质区别与内在联系。
    \end{itemize}
    \item \textbf{通用库构件}：标准泛型集合类库、迭代器模式及常用基础设施构件。
\end{itemize}

\subsubsection{示范性学习成果与能力目标}
\begin{enumerate}
    \item \textbf{[创建/开发]} 通过需求分析、面向对象设计、代码实现与单元测试，构建高质量的类以严格满足预设的功能与行为规约。
    \item \textbf{[应用]} 构建清晰简洁的类继承层次结构，在保障代码复用性的同时确保各类子类行为语义的独立性与正确性。
    \item \textbf{[分析/评估]} 准确预测、追踪并测试验证包含多态与动态分派的程序控制流执行轨迹。
    \item \textbf{[分析]} 深入对比与权衡：a) 过程式/函数式范式（为每种业务操作定义独立函数，函数体内部以模式匹配处理每种数据变体）；b) 面向对象范式（为每种数据变体定义独立类，类定义内部为每种操作实现具体方法）。深刻理解两者本质上是在对“操作--数据变体”二维矩阵进行正交分解（表达式问题）。
    \item \textbf{[分析]} 深入辨析并评估类继承（子类化）与面向对象组合（尤其是结合高阶函数的动态委托组合）在代码复用、系统耦合度与长期维护成本上的得失。
    \item \textbf{[解释]} 深刻阐释面向对象继承机制（代码复用与方法重写）与类型系统子类型化（子类型对象可在任何期望超类型的环境中安全替代）之间的数学与工程差异。
    \item \textbf{[应用]} 熟练运用现代面向对象的严格封装机制（如接口隔离、包内可见性与私有数据隐藏）防御性编写高健壮代码。
    \item \textbf{[枚举/应用]} 在多种主流面向对象编程语言中，熟练定义并使用聚合容器对象上的迭代器及高阶变换操作（包括以函数为实参的方法），并严格遵循各语言最自然的编码范式与规范惯例。
\end{enumerate}

\subsection{函数式编程}
\textbf{【学时建议】}：4 计算机核心学时，4 领域核心学时。

函数式编程范式建立在数学函数映射与 Lambda 演算基础之上，以不可变性、表达式求值与高阶抽象为核心基石。

\subsubsection{计算机核心主题}
\begin{itemize}
    \item \textbf{Lambda 表达式与求值机制}：
    \begin{itemize}
        \item 变量绑定与静态词法作用域规则。
        \item 形参实参求值传递机制。
        \item 嵌套 Lambda 表达式抽象与归约次序（Reduction Order，如正规序 vs 应用序）。
    \end{itemize}
    \item \textbf{无副作用编程}：
    \begin{itemize}
        \item 纯函数调用无隐藏副作用，具备极佳的引用透明性与逻辑组合式推理优势。
        \item 不可变数据与值复制语义 vs 表达式按步归约。
        \item 递归计算 vs 状态循环 vs 管道流式处理（Map / Filter / Reduce 变换流）。
    \end{itemize}
    \item \textbf{结构化代数数据处理}：
    \begin{itemize}
        \item 针对结构化复合数据（如代数数据类型、树形拓扑），基于模式匹配为每种数据变体提供独立处理分支。
        \item 基于子构件上的映射函数递归组合定义高阶复合数据操作。
        \item 持久化数据结构与结构共享原理。
    \end{itemize}
    \item \textbf{高阶函数}：将函数作为头等公民进行传递、组合、存储与动态返回。
\end{itemize}

\subsubsection{领域核心主题}
\begin{itemize}
    \item \textbf{函数闭包}：
    \begin{itemize}
        \item 闭包的数学定义与语义：函数对象在运行时捕获并绑定其外层封闭词法环境。
        \item 工业级经典惯用法：事件回调函数、集合迭代器实参、基于函数参数的泛化逻辑解耦。
        \item 运用闭包在内存环境中实现高内聚私有数据的严密封装。
        \item 惰性求值与及早求值的语义与开销对比。
    \end{itemize}
    \item \textbf{高阶函数复合与生成}：定义并实现能够接受函数作为实参、进行函数管道组合或动态构造返回新函数的复杂高阶系统。
\end{itemize}

\subsubsection{拓展选修主题}
\begin{itemize}
    \item \textbf{图归约机与按需调用}：基于图归约与按需调用的无副作用求值引擎。
    \item \textbf{惰性求值的底层工程实现}：重写 Thunk 悬挂计算节点与记忆化展开。
    \item \textbf{与逻辑式编程的融合}：等式逻辑、窄化推导、留数机制与语义合一。
    \item \textbf{跨范式整合}：函数式范式与命令式、面向对象范式的多范式协同设计。
\end{itemize}

\subsubsection{示范性学习成果与能力目标}
\begin{enumerate}
    \item \textbf{[创建/开发]} 开发避免向可变状态赋值且不依赖底层引用指针相等性判定的高可靠纯函数式算法。
    \item \textbf{[创建/开发]} 熟练开发接受其他函数作为参数并生成返回全新函数的高阶函数构件。
    \item \textbf{[分析]} 深入辨析过程式/函数式数据变体分解方案与面向对象类分解方案在系统演进时的扩展性得失。
    \item \textbf{[解释]} 结合 SECD 抽象机（Stack, Environment, Code, Dump 抽象模型）演示并详述 Lambda 表达式求值过程中内存环境的分配、检索与清理回收。
    \item \textbf{[分析]} 在包含深度嵌套闭包的代码中正确解析自由变量绑定关系与词法作用域链的动态生命周期。
    \item \textbf{[应用]} 熟练运用函数式封装技术（如环境闭包与高阶模块接口）实现模块解耦。
    \item \textbf{[分析]} 深入对比有状态执行与无状态执行在并发安全性与内存吞吐上的性能取舍。
    \item \textbf{[枚举/应用]} 在多种编程语言中定义并运用函数式高阶数据处理管道，以最高表达力的地道惯用法解决复杂数据流转换。
    \item \textbf{[解释]} 针对包含共享子表达式的复杂 $\lambda$-表达式，手绘展示图归约机的归约图变换过程。
    \item \textbf{[解释]} 借助 ABC 抽象图归约机演示嵌套 $\lambda$-表达式的编译展开与图遍历求值步骤。
    \item \textbf{[解释]} 借助经典示例演示窄化、留数与语义合一算法在混合求解中的实际推演。
    \item \textbf{[应用]} 编写函数式并发示例（如管道并发、无锁共享队列与顺序流混合控制），展示函数式无副作用在并行调度中的天然优势。
\end{enumerate}

\subsection{逻辑式编程}
\textbf{【学时建议】}：3 领域核心学时，其余为拓展选修。

逻辑式编程将计算视为在形式逻辑系统中的受控推理，程序由事实与规则构成，依托合一算法与自动回溯完成目标证明。

\subsubsection{领域核心主题}
\begin{itemize}
    \item \textbf{形式量词体系}：全称量词（$\forall$）与存在量词（$\exists$）的语义界定。
    \item \textbf{逻辑理论基础}：一阶谓词逻辑（FOPL）与高阶逻辑的表达力对比。
    \item \textbf{关系复合表达}：运用逻辑联结词（合取、析取、蕴涵、否定）与原子谓词构建高维复杂实体关系。
    \item \textbf{霍恩子句体系}：霍恩子句形式定义、事实声明、推导规则、顶层目标与求解子目标。
    \item \textbf{合一机制}：一阶合一原理与最一般合一子（MGU）求解算法；深刻辨析“合一”与“断言声明”及“表达式代数求值”的本质差异。
    \item \textbf{关系与函数的协同}：谓词关系逻辑与函数表达式的相互映射与混合编程。
    \item \textbf{非确定性控制}：非确定性推导执行模型、深搜回溯机制与剪枝运算符（Cut `\texttt{!}'）的控制语义。
    \item \textbf{世界假设理论}：封闭世界假设（Closed-World Assumption, 失败即否定）与开放世界假设的哲学与逻辑差异。
\end{itemize}

\subsubsection{拓展选修主题}
\begin{itemize}
    \item \textbf{内存开销优化}：处理深度递归与迭代逻辑程序时逻辑变量绑定环境复制的内存开销与尾递归优化。
    \item \textbf{中间状态记忆化}：制表记忆化与搜索树剪枝启发式控制。
    \item \textbf{逻辑与函数式杂交}：等式逻辑约束、窄化求解与语义合一框架。
    \item \textbf{前沿逻辑范式}：约束逻辑编程（CLP，如实数约束 CLP(R)、有限域约束 CLP(FD)）、归纳逻辑编程（ILP）与高阶逻辑式编程。
    \item \textbf{高级逻辑结构}：差分列表、自定义项结构与集合推导谓词（\texttt{setof/3}, \texttt{findall/3}）。
\end{itemize}

\subsubsection{示范性学习成果与能力目标}
\begin{enumerate}
    \item \textbf{[应用]} 运用纯逻辑式编程语言（如 Prolog）实现经典的图搜索或排序算法。
    \item \textbf{[应用]} 运用子句事实、递归关系规则与剪枝操作符，在逻辑语言中构建包含隐式搜索与约束剪枝的状态空间求解器。
    \item \textbf{[应用]} 借助清晰示例，形式化阐明一阶谓词逻辑公式与基于霍恩子句的逻辑程序之间的精确映射关系。
    \item \textbf{[应用]} 演示最一般合一算法在查询目标归约推导过程中的参数双向绑定流向与执行步骤。
    \item \textbf{[应用]} 编写交织关系逻辑、递归调用与数值运算的逻辑程序（如斐波那契序列、图拓扑通路求解），准确追踪其回溯树展开过程。
    \item \textbf{[分析/改进]} 分析简单朴素逻辑程序（如指数级递归求解）中的大量重复回溯开销，引入制表记忆化或累加器消除冗余，大幅提升推演性能。
\end{enumerate}

\subsection{脚本编程}
\textbf{【学时建议】}：2 计算机核心学时。

脚本编程关注快速胶水粘合、任务自动化批处理、操作系统交互与灵活的文本模式解析，具有极高的工程生产力。

\subsubsection{计算机核心主题}
\begin{itemize}
    \item \textbf{任务解构范式}：分而治之、系统构件粘合与规约处理。
    \item \textbf{多任务并发机制}：进程分叉、管道协同与多线程任务派发。
    \item \textbf{异常与防御性处理}：脚本执行环境中的错误捕获、返回码检查与健壮异常防护。
    \item \textbf{标准 I/O 与管道重定向}：输入输出流重定向（\texttt{<}, \texttt{>}, \texttt{2>\&1}）与跨进程数据流管道（\texttt{|}）。
    \item \textbf{系统级调度与环境变量}：调用宿主操作系统底层系统命令、子进程拉起与系统环境变量的读取和动态注入。
    \item \textbf{文件系统探针}：文件类型检查运算符、权限判定与路径元数据检索。
    \item \textbf{内置高阶动态容器}：动态数组、切片操作、列表推导式与关联字典。
    \item \textbf{文本模式工程}：基于正则表达式的高级文本过滤、提取与批量模式替换。
    \item \textbf{动态类型系统与运行时解析}：动态弱/强类型绑定规则、函数对象头等化与动态代码对象求值。
\end{itemize}

\subsubsection{示范性学习成果与能力目标}
\begin{enumerate}
    \item \textbf{[应用]} 编写并执行自动化运维与系统管理脚本，高效实现批处理任务调度、日志归档与开发流程流水线集成。
    \item \textbf{[应用]} 熟练运用脚本语言内置的高级数据结构与正则表达式，快速解决海量非结构化文本清洗与数据流转换问题。
\end{enumerate}

\subsection{事件驱动与响应式编程}
\textbf{【学时建议】}：2 计算机核心学时，2 领域核心学时。

本单元既可独立讲授，亦可与并发、异步及多线程知识单元紧密融合，用于深刻对比“事件循环驱动”与“抢占式多线程”两种执行模型。

\subsubsection{计算机核心主题}
\begin{itemize}
    \item \textbf{执行范式演进}：经典过程式顺序控制 vs 响应式事件驱动；响应式模型在捕获异步外界刺激与高并发 I/O 时的显著架构优势。
    \item \textbf{事件架构核心基石}：事件源、事件流信号、事件监听器与分发器、事件载荷对象、协议适配器与事件回调处理器。
    \item \textbf{行为建模}：事件驱动架构下基于单线程事件循环与非阻塞队列的行为演进模型。
    \item \textbf{经典落地场景}：现代图形用户界面（GUI）、移动终端交互、机器人实时传感交互、高吞吐网络代理服务。
    \item \textbf{形式化抽象}：将响应式系统严格形式化建模为状态转移系统。
\end{itemize}

\subsubsection{领域核心主题}
\begin{itemize}
    \item \textbf{响应式框架工程实践}：
    \begin{itemize}
        \item 编写工业级事件处理器与异步数据流管道监听器。
        \item 事件发布者与事件参数的泛型参数化解耦。
        \item 外部硬件中断触发事件与软件运行时合成派发事件的协同处理。
    \end{itemize}
    \item \textbf{架构解耦分层}：模型--视图--控制器（MVC）及其现代演进架构 (MVI)在事件驱动框架下的职责隔离。
\end{itemize}

\subsubsection{示范性学习成果与能力目标}
\begin{enumerate}
    \item \textbf{[应用]} 在真实响应式工程（如 GUI 交互或实时数据流监控系统）中设计并实现低延迟、无阻塞的健壮事件处理器。
    \item \textbf{[分析]} 深入论证为何在面向不可预测外部交互的软件领域中，事件驱动响应式编程是最自然、扩展性最高的设计范式。
    \item \textbf{[应用]} 熟练运用工业级响应式流编程框架（如 ReactiveX、Flow）构建复杂异步事件处理管道。
    \item \textbf{[解释]} 运用 MVC 架构模式精确绘制交互式系统的状态演化流向与各组件协作职责划分。
\end{enumerate}

\subsection{并行与分布式计算}
\textbf{【学时建议】}：3 计算机核心学时，2 领域核心学时。

编程语言对并发与分布式的原生支持是现代软件系统性能与弹性的核心瓶颈，本单元与体系结构及操作系统深度交织。

\subsubsection{计算机核心主题}
\begin{itemize}
    \item \textbf{并发正确性基础}：安全性（Safety / 坏事绝不发生）与活性（Liveness / 好事终将发生）形式化属性。
    \item \textbf{并发隐患}：数据竞争、死锁、活锁与饥饿。
    \item \textbf{数据与控制依赖}：因果前置条件依赖与终止性证明。
    \item \textbf{主流并发编程模型}：
    \begin{itemize}
        \item Actor 消息驱动无共享模型（如 Erlang/Akka）。
        \item 过程式与响应式异步编程模型。
        \item 现代语言中的同步/异步抽象（\texttt{async/await}、Promise/Deferred）。
        \item 数据并行性抽象。
    \end{itemize}
    \item \textbf{并发语义学准则}：计算交换律、事件因果序、执行独立性、一致性等级、原子性事务与分布式共识。
    \item \textbf{底层通信与协调原语}：消息传递通道、共享内存锁机制、管程、\texttt{cobegin-coend} 协同原语、守卫语句。
\end{itemize}

\subsubsection{领域核心主题}
\begin{itemize}
    \item \textbf{高级异步凭据}：Future（延时承诺对象）的代数结构与链式复合。
    \item \textbf{数据并行语言级支持}：\texttt{forall} 并发循环、自动循环展开、Map-Reduce 原语的硬件向量化展开。
    \item \textbf{硬件内存一致性模型}：顺序一致性、释放一致性、弱内存模型对编程语言语义规范与编译器优化重排的深刻影响。
    \item \textbf{分布式服务集成}：RESTful 架构契约与微服务远程方法调用。
    \item \textbf{并发基础设施开销}：消息序列化与深拷贝开销、上下文切换代价、多核缓存一致性伪共享（False Sharing）。
    \item \textbf{并发粒度权衡}：有效挖掘计算性能的最佳任务粒度，避免通信开销压倒并发增益。
    \item \textbf{范式交叉融合}：纯函数式不可变性与无锁并发编程的天然协同。
\end{itemize}

\subsubsection{示范性学习成果与能力目标}
\begin{enumerate}
    \item \textbf{[解释]} 深入阐明在存在数据竞争时现代语言运行时为何无法保证顺序一致性，以及开发者必须如何遵循语言内存模型契约以消除未定义行为。
    \item \textbf{[应用]} 熟练运用多种并发编程模型（共享内存互斥、Actor 消息队列、Future 异步链、数据并行原语）编写无死锁的并发程序。
    \item \textbf{[应用]} 运用无共享消息传递模型设计、建模并形式化分析一套分布式系统网络交互协议。
    \item \textbf{[应用]} 在多线程高并发系统中正确运用管程或同步锁机制实现临界资源的安全互斥调度。
    \item \textbf{[应用]} 运用变量读写依赖图对复杂计算过程的数据依赖关系进行形式化分析。
    \item \textbf{[应用]} 运用控制流分支与循环控制对复杂系统的前置条件与控制依赖进行结构化建模。
    \item \textbf{[解释]} 阐明 RESTful API 与无状态设计原则如何支撑现代互联网系统的高弹性横向扩展与业务自动化集成。
    \item \textbf{[解释/评估]} 系统评估分布式与多核并行计算相较于传统单线程顺序计算的理论算力收益、物理扩展瓶颈与系统工程挑战。
\end{enumerate}

\subsection{类型系统}
\textbf{【学时建议】}：3 计算机核心学时，4 领域核心学时，其余为拓展选修。

类型系统是现代计算机科学中最严谨、可扩展的轻量级程序形式化规约与静态验证工具链，直接捍卫软件的内存安全与业务逻辑正确性。

\subsubsection{计算机核心主题}
\begin{itemize}
    \item \textbf{类型的数学本质}：类型作为合法值的集合以及施加于该集合上的一组良定义操作算子集合。
    \begin{itemize}
        \item 原生基元类型（Primitive Types，如定宽整数、浮点数、布尔值）。
        \item 复合代数类型（Compound Types，如积类型记录 Records、和类型联合 Unions、数组、列表、函数空间、可变引用）。
    \end{itemize}
    \item \textbf{类型绑定体系}：静态类型推导/显式标注与变量、函数形参、返回值及对象字段的关联。
    \item \textbf{类型安全性}：类型安全作为程序逻辑正确性不可动摇的核心防线；未定义类型错误及其引发的内存崩溃灾难。
    \item \textbf{类型检查时序}：静态类型（编译期）与动态类型（运行期）的底层机制与权衡。
    \item \textbf{类型等价性判准}：名义等价vs 结构等价的理论与工程差异。
    \item \textbf{泛型与参数多态}：
    \begin{itemize}
        \item 多态体系四大分支：参数多态（泛型）、子类型多态（继承包含）、特设多态（Ad-hoc / 函数重载）与强制类型转换。
        \item 单态类型与多态类型的表达力对比。
        \item 泛型容器类库与标准模板架构。
        \item 描述性多态 vs 规定性多态。
        \item 泛型多态的底层编译器实现模型：类型擦除（Type Erasure，如 Java）vs 模板单态化膨胀（Monomorphization，如 C++/Rust）。
    \end{itemize}
\end{itemize}

\subsubsection{领域核心主题}
\begin{itemize}
    \item \textbf{静态类型与动态类型的深度工程互补}：
    \begin{itemize}
        \item 早期编译拦截 vs 延迟运行时抛出；
        \item 长期演进中利用类型不变量进行大规模重构保证 vs 敏捷原型期推迟定型并容纳高度灵活的异构集合；
        \item 严防代码误用与错误参数调用 vs 追求高自由度代码多态复用；
        \item 拦截未完全实现的代码 vs 允许不完备软件敏捷增量试错。
    \end{itemize}
    \item \textbf{形式化定型规则}：
    \begin{itemize}
        \item 函数类型（$T_1 \to T_2$）、积类型（$T_1 \times T_2$）与和类型（$T_1 + T_2$）的公理化定型推导规则。
        \item 基于一阶合一的最一般类型推导算法。
        \item 包含 Let-多态的 Lambda 演算系统的数学推导。
    \end{itemize}
    \item \textbf{与静态分析及可计算性的理论联结}：静态类型检查作为保守近似的数据流分析；莱斯定理与类型系统可判定性边界。
    \item \textbf{前沿高级工业类型系统}：以 Rust 为代表的仿射类型、所有权转移、借用检查器与静态生命周期系统。
\end{itemize}

\subsubsection{拓展选修主题}
\begin{itemize}
    \item \textbf{高阶组合类型构造器}：递归类型、全称量化类型（System F / $\forall X. T$）、存在量化类型（抽象数据类型理论 / $\exists X. T$）。
    \item \textbf{子类型化理论深入}：逆变与协变规则、行为替换模型、子类型与代码复用解耦。
    \item \textbf{类型健全性数学证明}：赖特-费莱森范式——类型保持性定理与进展性定理。
    \item \textbf{柯里-霍华德同构}：命题即类型，证明即程序；直觉主义逻辑合取对应积类型、析取对应和类型、蕴涵对应函数类型。
    \item \textbf{依赖类型理论}：值决定类型；全称量化对应依赖函数类型（$\Pi$ 类型），存在量化对应依赖对偶类型（$\Sigma$ 类型）。
\end{itemize}

\subsubsection{示范性学习成果与能力目标}
\begin{enumerate}
    \item \textbf{[解释]} 针对基元类型与各类代数复合类型，精确定义其数学值域空间及其合法操作约束。
    \item \textbf{[解释]} 阐明静态类型系统在编译期所能静态拦截并严厉禁止的非法操作集合。
    \item \textbf{[解释]} 举例阐述各类被现代高级类型系统在开发阶段直接阻断的经典运行时系统缺陷。
    \item \textbf{[枚举]} 在多种主流语言中，清晰甄别出在编译期静态检查的属性与在运行时动态校验的系统属性。
    \item \textbf{[解释]} 编写一个在特定静态语言中无法通过类型检查、但在实际运行时却绝对安全且无任何错误的合法程序范例，以此证明类型系统的保守近似性。
    \item \textbf{[应用]} 熟练运用精准的类型标注及编译器给出的丰富类型错误诊断报告进行代码编写与高精度除错。
    \item \textbf{[解释]} 阐明形式化定型规则如何通过前置判断严谨界定一种类型所支持的全部合法操作。
    \item \textbf{[枚举]} 形式化列举支配特定复合类型（如函数类型、积类型、和类型）构造与解构的核心定型演绎规则。
    \item \textbf{[解释]} 基于停机问题的不可判定性，严密论证为何任何健全的静态类型系统在判定程序语义时都必须采取保守近似策略。
    \item \textbf{[枚举/应用]} 熟练定义并运用泛型（参数多态）构件（如泛型算法、泛型集合类库）构建高度通用的可复用模块。
    \item \textbf{[分析]} 深入研讨并辨析参数多态（泛型）、子类型多态与特设多态（重载）在语言语义与底层编译器实现上的本质区别。
    \item \textbf{[解释]} 系统论述静态类型系统在工业级大型软件重构、接口协议维护与团队协同调试中所赋予的核心价值。
    \item \textbf{[枚举]} 形式化定义一套具备严密演绎规则的抽象类型系统。
    \item \textbf{[枚举]} 针对各类高级基础类型构造器，精确指出其所描述的数学值域及其强制维护的系统不变量。
    \item \textbf{[解释]} 精确陈述一套健全类型系统在动态求值演化中所永恒保持的语义不变量。
    \item \textbf{[评估]} 运用结构归纳法为一门核心微型语言严谨证明其类型安全性（证明类型保持性与进展性双定理）。
    \item \textbf{[应用]} 独立编写并实现基于 Robinson 合一算法的 Hindley-Milner 泛型类型推断求解器。
    \item \textbf{[解释]} 深入阐释静态重载解析算法与多态分派机制对程序动态执行性能与控制流的深远影响。
\end{enumerate}

\subsection{语言翻译与执行}
\textbf{【学时建议】}：4 计算机核心学时，其余为拓展选修。

本单元系统剖析从高级语言文本向机器指令翻译的完整编译流水线与运行时内存物理布局。

\subsubsection{计算机核心主题}
\begin{itemize}
    \item \textbf{执行范式三元谱系}：源码解释执行 vs 编译为宿主本地机器码 vs 编译为跨平台中间表示（如 JVM Bytecode、LLVM IR）。
    \item \textbf{现代语言端到端翻译流水线}：词法分析、上下文无关语法分析、静态语义检查与类型标注、中间代码生成与机器无关优化、目标机器代码生成、符号解析与链接、加载与物理执行。
    \item \textbf{形式文法规约}：上下文无关文法（CFG）的巴科斯-诺尔范式（BNF）及扩展 BNF（EBNF）精确规范。
    \item \textbf{语法分析树}：针对算术表达式与分支选择结构的推导展开与树形表征。
    \item \textbf{文法歧义性消解}：因缺乏明确优先级与结合性规范引发的语法二义性及其文法改写消歧技术。
    \item \textbf{底层执行环境}：本地硬件裸机执行 vs 托管虚拟机环境执行。
    \item \textbf{高级运行时编译策略}：模块动态链接加载与即时编译（JIT）代码生成。
    \item \textbf{控制流建模}：基于条件分支与循环迭代的程序控制流图（CFG）构造。
    \item \textbf{翻译与执行的核心数据结构}：物理调用栈、动态堆内存、别名指针、变长索引数组与不可变字符串。
    \item \textbf{内存寻址物理模型}：直接寻址、间接寻址与变址寻址机制。
    \item \textbf{数据抽象的运行时表征}：标量变量、多维矩阵、动态切片向量、结构体记录、基于指针的链表/树形构件与面向对象实例。
    \item \textbf{抽象低级执行模型}：用于推演语言编译与执行的精简抽象指令集、执行栈与堆模型。
    \item \textbf{进程运行时物理内存布局}：
    \begin{itemize}
        \item 活动记录/栈帧（Activation Record / Stack Frame，包含返回地址、动态调用链、静态词法链指针、局部变量分配）；
        \item 全局静态数据段、代码只读段、动态调用栈与动态增长堆区。
    \end{itemize}
    \item \textbf{控制与函数翻译机制}：
    \begin{itemize}
        \item 将选择与循环控制结构编译为控制流跳转图与底层指令；
        \item 循环展开、通用递归函数执行与尾调用消除优化；
        \item 函数/过程调用的底层汇编级协议：调用栈分配、现场保存、传参约定与调用返回。
    \end{itemize}
    \item \textbf{现代内存精细管理}：
    \begin{itemize}
        \item 高级复杂复合数据在底层的物理对齐与内存分配；
        \item 过程调用返回时栈指针平移对局部变量的零成本自动回收；
        \item 手工堆内存管理：分配、释放与内存碎片重用；
        \item 自动内存管理：基于对象图可达性判定的现代垃圾回收（GC）机制。
    \end{itemize}
\end{itemize}

\subsubsection{拓展选修主题}
\begin{itemize}
    \item \textbf{高级核心特性的底层硬件表示}：面向对象多态虚方法表的指针寻址；头等函数闭包对象在内存中的物理对齐与环境指派。
    \item \textbf{安全编译器工程}：防御栈溢出破坏（Canary 保护）、地址空间布局随机化（ASLR）与控制流完整性（CFI）加固机制。
\end{itemize}

\subsubsection{示范性学习成果与能力目标}
\begin{enumerate}
    \item \textbf{[解释]} 深刻辨析编程语言的抽象规范定义（语义内涵）与具体工程实现（编译器、解释器、JIT、内存对象排布）之间的本质区别。
    \item \textbf{[解释]} 严格厘清词法/语法解析与深层语义解释/指令求值之间的理论与技术分界。
    \item \textbf{[应用]} 熟练运用 BNF 与 EBNF 为具体语言特性（如控制流分支、模式匹配、类型声明）编写严格的语法产生式规约。
    \item \textbf{[应用]} 针对给定的算术表达式或嵌套控制语句，准确手绘并推导其语法解析树。
    \item \textbf{[应用]} 熟练将直观可视化的语法图谱精确转换为等价的 BNF/EBNF 产生式集合。
    \item \textbf{[应用]} 针对经典的悬空 Else（Dangling-else）等语法二义性难题，通过文法改写与引入显式结合律消解歧义。
    \item \textbf{[应用]} 精确手绘现代高级语言核心数据与控制抽象在底层抽象机中的物理运行时内存布局图。
    \item \textbf{[解释]} 详述语言运行时如何将物理虚拟内存划分为代码段、静态全局区、堆区与栈区，并解释递归函数与动态对象如何映射到该内存架构中。
    \item \textbf{[应用]} 借助调试工具与内存检测工具（如 Valgrind、AddressSanitizer），精准排查、定位并修复由野指针、悬空指针与内存泄露引发的系统崩溃漏洞。
    \item \textbf{[解释/评估]} 深入论述自动垃圾回收相较于手动内存管理的技术利弊，并基于根对象集阐明基于图可达性遍历的核心回收机制。
\end{enumerate}

\subsection{程序表示}
\textbf{【学时建议】}：3 领域核心学时。

程序不仅是供机器执行的指令序列，更是元程序（如编译器、分析器、优化器）作为输入进行分析与变换的复杂数据结构。

\subsubsection{领域核心主题}
\begin{itemize}
    \item \textbf{形式语言表征工具}：BNF、EBNF 与正则表达式的表达力边界。
    \item \textbf{以程序为输入的元程序系统}：解释器、编译器、静态类型检查器、代码重构引擎、自动化文档生成器。
    \item \textbf{语言构成要素}：字母表、词法分隔符、词法记号、良构语句、具体语法与深层语义。
    \item \textbf{形式语义学流派概览}：操作语义、公理语义、指称语义与行为语义；抽象语法树（AST）的定义与遍历；抽象语法 vs 具体语法的解耦。
    \item \textbf{程序的形式化本质}：程序作为一组无歧义且蕴含明确计算语义的良构语句集合。
    \item \textbf{核心通用编程抽象}：常量、变量、嵌套声明、命令语句、表达式、赋值、选择分支、有界/无界循环、迭代器、函数、过程、模块、异常处理。
    \item \textbf{可变性与状态管理}：可变变量 vs 不可变变量；原地覆盖复用内存的性能收益与调试复杂性 vs 复制创建新值保留历史快照的推理确定性与内存开销；中间状态记忆化 vs 重复计算的工程权衡。
    \item \textbf{左值与右值}：可变变量名到内存左值地址（L-value）的映射，不可变绑定到右值具体内容（R-value）的映射。
    \item \textbf{变量的作用域与生存期}：静态全局变量、局部栈变量、非局部自由变量（闭包捕获变量）、全局共享变量；非局部与全局变量引发的耦合隐患与管控策略。
    \item \textbf{作用域消解法则}：静态词法作用域vs 动态作用域；变量可见性遮蔽；隐藏副作用。
    \item \textbf{环境与存储}：运行时环境映射（Environment: Identifier $\to$ Location）与内存存储映射（Store: Location $\to$ Value）的二元数学解耦模型。
    \item \textbf{模块间信息交互机制}：非局部变量访问、全局共享变量、形参实参传递、显式模块导出/导入接口。
    \item \textbf{参数传递机制全景对比}：传值调用、传引用调用、传值-结果调用、传名调用（Call by name / 替换求值）、按需调用（Call by need / 惰性记忆化求值）。
    \item \textbf{代码的多级中间表示}：用于代码分析、优化与迁移的内部数据结构（抽象语法树、三地址码、控制流图、静态单赋值 SSA 形式）。
    \item \textbf{底层指令表征}：虚拟机操作码、汇编指令与二进制 ELF/PE 机器码布局。
    \item \textbf{Lambda 演算理论基础}：变量绑定、自由变量与约束变量、$\alpha$-变换（重命名）、$\beta$-归约与 $\eta$-变换。
    \item \textbf{移动计算代码迁移}：分布式系统基于字符串与序列化字节码的代码动态迁移与远程执行。
\end{itemize}

\subsubsection{示范性学习成果与能力目标}
\begin{enumerate}
    \item \textbf{[应用]} 编写示例程序，精确演示变量的静态词法作用域边界、嵌套遮蔽效应与生命周期可见性。
    \item \textbf{[应用]} 借助清晰程序实例，深入对比传值、传引用、传值-结果与传名调用在参数求值时序与副作用表现上的本质差异。
    \item \textbf{[解释]} 详述滥用全局变量与外层非局部变量如何引发难以排查的竞态隐患，并提出工程重构与参数化解耦方案。
    \item \textbf{[解释]} 深入阐释元编程系统（如编译器、代码审计器）如何将其他程序源代码作为结构化输入数据流进行解析与变换。
    \item \textbf{[解释/设计]} 为一门自定义微型领域语言规范其上下文无关文法产生式，并设计对应的抽象语法树（AST）数据结构。
    \item \textbf{[解释]} 阐述使用结构化树形抽象语法表示程序相较于使用原始源代码纯文本字符串所带来的不可替代的分析与变换优势。
    \item \textbf{[应用]} 编写一个能够遍历并处理某种代码结构化中间表示的实用工具，如表达式解释器、AST 常量折叠优化器或 API 接口文档生成器。
\end{enumerate}

\subsection{语法分析}
\textbf{【学时建议】}：拓展选修。

语法分析是编译前端的核心组件，涵盖从字符流到词法记号、进而构建语法分析树的形式化推导过程。

\subsubsection{拓展选修主题}
\begin{itemize}
    \item \textbf{文法理论分类}：正则文法与上下文无关文法（CFG）的理论表达力极限对比。
    \item \textbf{词法扫描与语法分析协同架构}：基于语言规格说明书的词法扫描与语法解析流水线。
    \item \textbf{词法分析理论}：基于正则表达式的词法规范定义、DFA/NFA 状态机转换与最小化。
    \item \textbf{词法记号}：Token 的属性值载荷及其向解析器的传递协议。
    \item \textbf{语法分析策略全景}：
    \begin{itemize}
        \item 自顶向下解析：递归下降预测分析、LL(1) 与 LL(k) 分析表构建；
        \item 自底向上解析：移入-规约冲突消解、LR(0)、SLR(1)、LR(1)、LALR(1) 与广义上下文无关解析（GLR）。
        \item 前瞻查找表在自底向上状态转移中的数学构建原理。
    \end{itemize}
    \item \textbf{形式语言层次理论}：乔姆斯基文法体系；最左推导、最右推导与文法二义性判定；左递归消除与提取左公因子变换。
    \item \textbf{工业级错误恢复机制}：恐慌模式、短语级恢复与错误产生式配置。
    \item \textbf{解析器自动化生成工具}：运用声明式工具（如 Lex/Flex, Yacc/Bison, ANTLR）自动生成工业级词法解析器与语法分析器。
\end{itemize}

\subsubsection{示范性学习成果与能力目标}
\begin{enumerate}
    \item \textbf{[应用]} 熟练运用形式文法为一门全新领域特定语言（DSL）精确界定其完备语法产生式。
    \item \textbf{[解释]} 结合具体解析场景，详述前瞻符号表在驱动自底向上 LR 解析器进行移入/规约决策时的关键作用。
    \item \textbf{[应用]} 熟练运用现代声明式编译器前端工具（如 ANTLR）自动构建具有良好错误提示能力的健壮解析器。
    \item \textbf{[枚举]} 全面归纳语法规范定义中必须妥善解决的三大核心语法问题：二义性消解、算符结合性规则以及运算优先级阶梯。
\end{enumerate}

\subsection{编译器语义分析}
\textbf{【学时建议】}：拓展选修。

语义分析确保源代码符合语言的上下文敏感约束，是连接语法解析树与目标代码生成的桥梁。

\subsubsection{拓展选修主题}
\begin{itemize}
    \item \textbf{高级程序树形表示}：抽象语法树（AST）与具体解析树（CST）的解耦；AST 的定义、遍历（访问者模式 Visitor Pattern）与结构重写。
    \item \textbf{作用域与符号表管理}：多层嵌套作用域符号表体系构建与标识符绑定消解。
    \item \textbf{静态语义检查规约}：类型系统相容性检查；先声明后使用规则；静态断言与扩展静态检查框架。
    \item \textbf{值范畴语义}：左值（L-value）与右值（R-value）语义判定。
    \item \textbf{调用与分派语义}：函数调用语义、重载函数解析与多态动态绑定签名匹配。
    \item \textbf{形参实参绑定语义}：形参与实参传递协议的静态约束校验。
    \item \textbf{声明式语义分析体系}：属性文法、综合属性与继承属性在处理上下文敏感特性中的工程实现。
\end{itemize}

\subsubsection{示范性学习成果与能力目标}
\begin{enumerate}
    \item \textbf{[解释]} 为一门具体的小型语言精确定义并绘制其抽象语法树（AST）数据结构节点层级规范。
    \item \textbf{[应用]} 独立实现上下文敏感的源码级静态语义分析器，完成符号表管理、作用域查找与类型匹配检查。
    \item \textbf{[解释]} 运用属性文法为语言语法产生式标注语义规则，并演示综合属性自底向上与继承属性自顶向下的求值传播过程。
\end{enumerate}

\subsection{程序分析与分析器}
\textbf{【学时建议】}：拓展选修。

程序分析通过形式化方法推演程序运行时行为性质，是现代编译器高级优化、漏洞挖掘与软件验证的核心底座。

\subsubsection{拓展选修主题}
\begin{itemize}
    \item \textbf{核心中间图表示}：基本块、控制流图（CFG）、定义-使用链、静态单赋值形式（SSA）。
    \item \textbf{可计算性根本边界}：莱斯定理（Rice's Theorem）与停机问题不可判定性对程序静态分析的根本理论制约。
    \item \textbf{流无关分析}：类型推断检查、高可扩展性指针分析与别名分析（如 Andersen、Steensgaard 算法）。
    \item \textbf{流敏感分析}：基于半格理论与最大固定点算法（MFP）的前向/后向单调数据流分析框架（活跃变量分析、可用表达式、到达定值）。
    \item \textbf{路径敏感分析}：符号执行（Symbolic Execution）、软件模型检验与抽象解释。
    \item \textbf{工业级分析平台框架}：现代静态分析基础设施（如 LLVM, Soot, WALA, Infer）。
    \item \textbf{静态分析指导代码优化}：常量传播、死代码消除、循环不变式外提与并发数据依赖分析。
    \item \textbf{安全性漏洞挖掘}：空指针解引用、缓冲区溢出、资源泄露的静态模式检测与形式化验证。
    \item \textbf{并行化专门分析}：面向多核的自动并行化分析、数据竞争与潜在死锁检测。
\end{itemize}

\subsubsection{示范性学习成果与能力目标}
\begin{enumerate}
    \item \textbf{[枚举]} 在格论与单调函数的数据流分析框架下，形式化定义一套具体的数据流分析算法。
    \item \textbf{[解释]} 深入阐明数据流图与控制流图在信息流动表征上的根本区别。
    \item \textbf{[解释]} 从图灵停机问题不可判定性出发，严格论证为何任何有价值的健全静态分析算法都必须采取保守近似策略。
    \item \textbf{[分析]} 形式化推导并论证一个具体的程序静态分析算法在数学上的完备健全性与算法必终止性。
    \item \textbf{[解释]} 阐明潜在指针别名如何严重破坏静态分析的精度，并分析别名分析技术如何缓解该瓶颈。
    \item \textbf{[应用]} 运用程序分析算法导出的数据流信息指导编译器后端实施高效代码优化，或在 CI/CD 流水线中进行源码级缺陷自动化挖掘。
\end{enumerate}

\subsection{代码生成}
\textbf{【学时建议】}：拓展选修。

代码生成是将高级中间表示高效映射为底层目标体系结构指令的关键编译后端阶段。

\subsubsection{拓展选修主题}
\begin{itemize}
    \item \textbf{目标体系结构模型}：目标指令集架构（ISA，如 x86-64, ARM64, RISC-V）的寄存器组织与指令约束。
    \item \textbf{控制流映射}：将控制流图分支映射为条件跳转与无条件分支指令序列。
    \item \textbf{运行时内存与布局约束}：目标平台内存对齐限制与系统调用规约。
    \item \textbf{过程调用约定 (ABI)}：应用二进制接口（ABI）、调用者/被调用者保存寄存器划分、栈帧分配与虚方法分派开销。
    \item \textbf{分离编译与链接}：多模块独立编译、符号重定位、静态链接与动态共享库链接。
    \item \textbf{指令选择技术}：树重写模式匹配与动态规划最优指令覆盖（Maximal Munch 算法）。
    \item \textbf{指令调度技术}：利用硬件流水线、规避乱序执行冒险与减少等待气泡。
    \item \textbf{寄存器分配理论}：干涉图构建、肯普图着色启发式算法（Chaitin-Briggs 算法）与线性扫描分配。
    \item \textbf{机器级后端优化}：窥孔优化、公共子表达式消除与死代码清除。
\end{itemize}

\subsubsection{示范性学习成果与能力目标}
\begin{enumerate}
    \item \textbf{[枚举]} 全面归纳将高级语言中间表示自动化转换为底层汇编或目标二进制指令的关键流水线阶段。
    \item \textbf{[创建/开发]} 为现代高级语言中的函数调用生成严格符合目标平台 ABI 规范的底层机器指令序列（包含调用栈维护与寄存器保护）。
    \item \textbf{[解释]} 阐释为何跨模块的分离编译体系强制要求整个软件生态严格遵守统一的二进制调用约定（ABI）。
    \item \textbf{[解释]} 论述分离编译在带来独立构建优势的同时，因跨编译单元调用的黑盒未知效应给全局跨过程优化带来的制约。
    \item \textbf{[解释/评估]} 深入研讨朴素代码翻译所引入的大量机器指令冗余，并系统分析现代编译器利用指令选择、指令调度与图着色寄存器分配进行深度优化的核心原理。
\end{enumerate}

\subsection{运行时行为与系统}
\textbf{【学时建议】}：拓展选修。

运行时系统在幕后支撑高级语言特性的运转，涵盖栈帧恢复、动态内存池、动态加载以及沙箱隔离等机制。

\subsubsection{拓展选修主题}
\begin{itemize}
    \item \textbf{运行时进程模型}：基于调用栈与堆的活动记录（Activation Record / 栈帧）生命周期；借助帧指针与返回地址实现函数调用环境恢复。
    \item \textbf{面向对象运行时方法查找}：基于哈希表的动态方法检索、虚方法表与接口方法表的索引分派机制。
    \item \textbf{物理对象布局}：对象头、对齐填充与栈帧物理内存布局。
    \item \textbf{堆对象分配策略}：空闲链表、伙伴系统、线程本地分配缓冲区（TLAB）。
    \item \textbf{面向对象高级动态行为}：虚基类指针偏移量计算与内联缓存加速。
    \item \textbf{作用域链与非局部访问}：嵌套子函数通过静态作用域链跨级访问外层词法作用域变量的栈内存分配权衡。
    \item \textbf{现代动态内存管理技术}：
    \begin{itemize}
        \item 手工管理：\texttt{malloc/free} 内存池；
        \item 自动垃圾回收（GC）：标记-清除、复制算法、引用计数与分代回收；
        \item 基于生命周期区域的内存分配。
    \end{itemize}
    \item \textbf{即时编译与动态优化}：JIT 编译执行、动态重编译与基于运行时 Profiling 数据推测去优化。
    \item \textbf{操作系统底层交互}：程序冷启动初始化、信号处理与系统调用封装。
    \item \textbf{跨语言互操作性 (FFI)}：跨语言二进制参数传递协议、C-ABI 兼容性、大小端对齐（Big-Endian vs. Little-Endian）与多维数组存储对齐。
    \item \textbf{虚拟机核心服务}：类动态加载器、轻量级绿色线程/协程调度器与安全边界检查。
    \item \textbf{安全沙箱隔离机制}：不可信移动代码与插件容器的运行时内存保护与系统调用拦截。
\end{itemize}

\subsubsection{示范性学习成果与能力目标}
\begin{enumerate}
    \item \textbf{[分析]} 结合内存碎片率、空间局部性、时间局部性以及元数据空间开销等核心指标，全面对比不同动态内存管理策略的性能得失。
    \item \textbf{[解释/研讨]} 系统论述自动化内存回收相较于手动管理的软件工程收益，以及垃圾回收导致的“全局停顿”（Stop-The-World）对实时系统的挑战。
    \item \textbf{[解释]} 深入阐明在对象与活动记录的运行时表征中元数据（如类元指针、数组边界长度、帧指针、返回地址）的关键作用。
    \item \textbf{[分析]} 深刻对比与辨析程序数据元素在静态数据区、栈式自动分配与堆式动态分配三种模式下的生命周期与访问开销。
    \item \textbf{[解释]} 从数据生命周期超越调用上下文的角度，解释为何某些数据必须从栈提升至堆分配，从而揭示垃圾回收的必然需求。
    \item \textbf{[解释]} 全面论述即时编译（JIT）与动态重编译技术在执行吞吐、热点代码优化与预热时延上的综合取舍。
    \item \textbf{[解释]} 阐述在不可信移动代码与插件架构中运用沙箱技术实现系统隔离与权限防御的核心机制。
    \item \textbf{[枚举]} 系统归纳现代高级编程语言运行时系统为上层应用所提供的核心基础设施服务列表。
\end{enumerate}

\subsection{高级编程构造}
\textbf{【学时建议】}：拓展选修。

本单元涵盖现代主流语言为突破传统抽象瓶颈而引入的前沿控制结构、高级元编程与领域特定语言机制。

\subsubsection{拓展选修主题}
\begin{itemize}
    \item \textbf{高级封装机制}：模块系统、不透明类型与包私有隔离。
    \item \textbf{惰性求值与数据流}：惰性计算流与无限序列（Infinite Streams / 生成器模式）。
    \item \textbf{求值策略对比}：惰性求值 vs 及早求值的内存语义与执行开销深入权衡。
    \item \textbf{形式计算机制对比}：逻辑合一机制 vs 断言机制 vs 表达式求值的本质对照。
    \item \textbf{高阶控制抽象}：结构化异常处理机制、头等延续（First-Class Continuations / \texttt{call/cc}）、单子（Monads，如 State, IO, Maybe 单子）。
    \item \textbf{面向对象高阶结构}：多重继承菱形问题及其变体、混入组合、特质、多重分派/多分派方法。
    \item \textbf{现代元编程体系}：编译期宏系统（卫生的语法宏 Hygienic Macros）、生成式编程、基于模型的全自动代码生成。
    \item \textbf{高级模式匹配}：基于模式匹配与解构赋值的高效结构化数据处理。
    \item \textbf{动态代码求值}：运行时动态字符串代码求值（\texttt{eval}）与反射机制的威力与安全漏洞隐患。
    \item \textbf{契约式设计支持}：语言级运行时断言、类不变量、前置条件与后置条件的静态/动态检验。
    \item \textbf{领域特定语言 (DSL)}：数据库查询语言、数据科学分析语言、同步响应式系统语言、硬件描述语言。
    \item \textbf{大规模高性能计算语言}：面向超大规模集群与多核异构的高性能并行语言模型。
\end{itemize}

\subsubsection{示范性学习成果与能力目标}
\begin{enumerate}
    \item \textbf{[应用]} 熟练且正确地在工业级复杂软件工程中运用各类高级编程构造与设计模式。
    \item \textbf{[认知]} 深入研讨各类高级编程构造如何有效改善大型软件架构的解耦度、提升软件质量并成倍释放开发者生产力。
    \item \textbf{[认知]} 深入研讨引入高级语言构造（如延续、单子、多继承、动态反射）如何与其他已有语言特性的形式化定义及运行时编译器实现产生复杂的交互与干扰。
\end{enumerate}

\subsection{语言实用学}
\textbf{【学时建议】}：拓展选修。

语言实用学探究技术演进、应用场景需求以及人因工程对编程语言设计哲学与工程实现产生的实际约束。

\subsubsection{拓展选修主题}
\begin{itemize}
    \item \textbf{需求驱动的语言演进}：计算平台变迁与软件工程规模爆炸对编程语言诞生与演进的深层驱动机制。
    \item \textbf{问题领域与范式匹配}：特定行业应用特征与过程式、面向对象、函数式、逻辑式主流编程范式的匹配原则。
    \item \textbf{优秀语言设计准则}：正交性（Orthogonality / 独立概念互不影响）、合理的默认配置、高可读性与表达力、简洁性。
    \item \textbf{语法构造规范}：现代控制语句与迭代语法糖的精巧设计。
    \item \textbf{超大型系统工程化}：命名空间隔离、依赖包管理器与大型软件工程模块化体系。
    \item \textbf{求值语义确定性}：表达式求值顺序、运算符优先级与结合性规则对语义确定性的影响。
    \item \textbf{及早求值 vs 延迟求值}：语言层面对求值时机的折中设计与语法体现。
    \item \textbf{外部函数与系统库集成}：外部函数接口（FFI）调用与底层系统库集成对语言实现与编译器的深远影响。
\end{itemize}

\subsubsection{示范性学习成果与能力目标}
\begin{enumerate}
    \item \textbf{[解释]} 深入研讨在语言设计中“正交性原则”与“合理的默认配置”如何显著降低开发者心智负担并提升语言规范的优雅度。
    \item \textbf{[应用]} 运用客观、多维度的技术评估准则，对现实中的语言设计决策进行深入评测与架构选型。
    \item \textbf{[应用]} 编写具体的示例程序，直观展示该程序在不同求值顺序、不同优先级或不同结合性规则下所产生的迥异运行结果。
    \item \textbf{[应用]} 演示延迟求值技术的高级工程应用场景（如自定义短路控制结构）。
    \item \textbf{[解释]} 论述现代语言允许无缝调用外部原生函数与操作系统级库的现实必要性，并剖析其对跨语言垃圾回收及内存安全带来的系统实现挑战。
\end{enumerate}

\subsection{形式语义学}
\textbf{【学时建议】}：拓展选修[交叉引用: 软件工程/Formal Methods]。

形式语义学为程序语言提供严格的数学定义，是进行程序正确性形式化证明、编译器验证与类型安全论证的数理基石。

\subsubsection{拓展选修主题}
\begin{itemize}
    \item \textbf{语法与语义的严格解耦}：具体语法产生式 vs 抽象数学语义解释。
    \item \textbf{形式语义学三大主流流派体系}：
    \begin{itemize}
        \item \textbf{公理语义学}：基于霍尔逻辑的霍尔三元组 $\{P\} C \{Q\}$，最弱自由前置条件推导，循环不变量的形式化构造与程序部分/完全正确性证明。
        \item \textbf{操作语义学}：结构化操作语义（SOS / 小步操作语义）与自然语义（大步操作语义），基于环境与存储抽象转移关系的机器状态转换推演；符号执行技术与约束求解器协同。
        \item \textbf{指称语义学}：语法实体向数学语义域的函数映射，偏序集、完全偏序（CPO）与斯科特拓扑，基于塔斯基/克尼斯特不动点理论解析递归与循环的数学内涵。
        \item \textbf{基于类型的语义学}。
    \end{itemize}
    \item \textbf{Lambda 演算数学推导}：纯无类型与类型化 Lambda 演算的形式化归约规则。
    \item \textbf{结构归纳法证明}：基于语言归约规则与语法结构的超限数学归纳法。
    \item \textbf{类型系统形式化证明}：类型安全作为“类型保持性”与“进展性”双定理的严格结晶 [交叉引用: 编程语言/类型系统]。
    \item \textbf{柯里-霍华德同构}：直觉主义命题逻辑与类型化演算的数学等价映射（蕴涵映射为函数，逻辑合取映射为积类型，逻辑析取映射为和类型）。
    \item \textbf{高级类型理论}：依赖类型系统（$\Pi$ 类型与 $\Sigma$ 类型）的语义模型。
    \item \textbf{参数多态性理论}：雷诺兹参数多态定理与免费定理（Theorems for free!）。
\end{itemize}

\subsubsection{示范性学习成果与能力目标}
\begin{enumerate}
    \item \textbf{[创建/开发]} 为一门具体的小型命令式或函数式语言构建严谨的小步操作语义或公理语义系统。
    \item \textbf{[应用]} 编写纯 Lambda 演算表达式，并通过严格的 $\beta$-归约与 $\alpha$-变换演算推导其至标准正规形式。
    \item \textbf{[解释]} 系统研讨并深刻辨析操作语义、指称语义与公理语义在设计初衷、数学表达与工程适用性上的本质区别。
    \item \textbf{[应用]} 运用数学结构归纳法严格证明目标语言中所有良构程序所共同遵循的确定性或终止性定理。
    \item \textbf{[应用]} 运用结构归纳法证明：凡在形式化类型系统中通过类型检查的良构程序，在操作语义演进中均满足类型安全性（保持性与进展性）。
    \item \textbf{[应用]} 运用参数多态性的数学定理，仅凭借一个泛型多态函数的类型签名推导并锁定该函数所必须遵守的唯一外在行为。
\end{enumerate}

\subsection{形式化开发方法学}
\textbf{【学时建议】}：拓展选修[交叉引用: 软件工程/Formal Methods]。

形式化开发方法学涵盖利用形式逻辑、规格说明语言与机器辅助证明工具进行高安全、完全可信软件开发的前沿工程路径。

\subsubsection{拓展选修主题}
\begin{itemize}
    \item \textbf{形式化规格说明语言}：工业级形式化建模规范（如 Z 语言、B 方法、TLA+、Alloy、VDM）。
    \item \textbf{定理证明与辅助工具}：交互式定理证明助手（Interactive Proof Assistants，如 Coq, Isabelle/HOL, Lean, Agda）及底层高阶逻辑体系。
    \item \textbf{自动约束求解器}：SAT/SMT 自动化约束求解器（如 Z3, CVC5）。
    \item \textbf{依赖类型形式化系统}：基于类型理论的程序即证明验证范式。
    \item \textbf{完全验证软件构建}：基于前置/后置契约、精化类型的完全形式化规约与机器证明释放。
    \item \textbf{形式化建模与逐步精化}：软件系统的高层形式化建模与自顶向下阶梯式手工/辅助程序精化。
    \item \textbf{符号化测试与模糊测试}：符号执行与导向型模糊测试（Fuzzing）在高可靠验证流水线中的集成。
    \item \textbf{模型检验}：基于状态机遍历、Kripke 结构与线性时序逻辑（LTL）/ 计算树逻辑（CTL）的并发协议自动化穷尽验证。
    \item \textbf{形式化方法工程落地策略}：在航空航天飞控、高可靠微内核、区块链智能合约及底层密码库研发中评估并规划形式化方法的最佳应用边界。
\end{itemize}

\subsubsection{示范性学习成果与能力目标}
\begin{enumerate}
    \item \textbf{[应用]} 运用形式化建模技术（如 TLA+ 或 Alloy）对分布式共识协议或复杂系统架构进行形式化建模与全状态空间验证。
    \item \textbf{[应用]} 运用基于证明助手的编程语言（如 Coq 或 Lean）开发携带机器可检验数学正确性证明的高可靠软件构件。
    \item \textbf{[应用]} 充分利用主流编程语言中内置的契约规格支持与形式化验证器，严格验证核心系统属性与内存安全。
    \item \textbf{[分析]} 将符号执行与形式化属性验证工具链无缝嵌入持续集成与日常软件开发流水线中。
    \item \textbf{[解释]} 深入研讨在现代大型工业软件工程全生命周期中，究竟在何时、以何种粒度引入形式化方法能最有效地降低系统性灾难风险。
\end{enumerate}

\subsection{嵌入式计算与硬件接口}
\textbf{【学时建议】}：拓展选修。

本单元聚焦于编程语言在面对极端物理资源受限、直接控制硬件以及保障确定性时间约束时的语言机制与接口交互。

\subsubsection{拓展选修主题}
\begin{itemize}
    \item \textbf{面向定制硬件架构的专用语言}：GPU 大规模并行计算语言、FPGA 高层次综合硬件描述、异构多核加速语言。
    \item \textbf{嵌入式系统软件基石}：微控制器（MCU）裸机与 RTOS 编程、硬件中断机制与闭环反馈、高级语言中的中断服务程序（ISR）设计与现场上下文保护。
    \item \textbf{中断分类与异常陷阱}：硬中断、软中断、异常捕获与陷阱退出。
    \item \textbf{底层物理交互}：与外部传感器、执行器及外设的直接内存映射 I/O（MMIO）交互。
    \item \textbf{高能效计算}：面向超低功耗电池驱动环境的能效感知型编程。
    \item \textbf{松散时序与同步}：松散时序编码与硬件事件同步。
    \item \textbf{硬件抽象层设计}：硬件抽象层（HAL）与板级支持包（BSP）适配器设计。
    \item \textbf{实时系统核心理论}：
    \begin{itemize}
        \item 硬实时与软实时系统的严谨界定；
        \item 确定性时效性；
        \item 实时调度理论（单调速率调度 RMS、最早截止时间优先 EDF）与任务集可调度性数学分析；
        \item 优先级反转隐患与优先级继承/置顶协议；
        \item 确定性执行时延与计算抖动的严格抑制。
    \end{itemize}
    \item \textbf{底层物理内存控制}：高级语言变量与物理固定地址/MMU 虚拟页面的映射；共享物理内存段；裸机环境下的无堆手工静态分配；嵌入式环境下的实时确定性垃圾回收。
\end{itemize}

\subsubsection{示范性学习成果与能力目标}
\begin{enumerate}
    \item \textbf{[创建/开发]} 设计并实现能够直接与底层传感器、执行器等硬件设备进行高可靠交互与实时控制的系统级软件。
    \item \textbf{[创建/开发]} 为具有严苛时间约束的硬实时嵌入式系统设计高确定性的任务调度与控制逻辑。
    \item \textbf{[分析]} 评估实时任务调度策略并进行严谨的数学可调度性分析，论证任务集在最坏情况下的时限保证。
    \item \textbf{[分析]} 运用形式化方法对嵌入式控制系统的精确时间约束与物理响应属性进行形式化规约与安全性验证。
\end{enumerate}

\subsection{编程语言基础与社会、伦理及专业素养}
\textbf{【学时建议】}：拓展选修。

本单元审视编程语言设计对多元包容性、国际化文化差异、残障人士可访问性以及行业技术伦理所产生的深刻社会影响。

\subsubsection{拓展选修主题}
\begin{itemize}
    \item \textbf{以英语为中心的编程语言反思}：现代主流编程语言以英语关键词为核心对全球非英语母语开发者造成的认知摩擦、数字鸿沟与文化冲击。
    \item \textbf{提升残障人士开发者的可访问性与包容性}：语言语法对屏幕阅读器等辅助技术的原生友好性支持。
    \item \textbf{语言可用性与人因工程学}：
    \begin{itemize}
        \item 语法复杂度与标点符号设计对无障碍开发工具与语音输入识别的影响；
        \item 支持全球多元文化差异（本地化货币单位、十进制小数分隔符、公历/农历日期时间格式）；
        \item 对神经多样性群体（如诵读障碍、自闭症谱系开发者）友好的语法设计与 IDE 高对比度呈现。
    \end{itemize}
    \item \textbf{编程语言历史术语的认识论反思}：反思历史遗留技术术语（如 Master/\allowbreak Slave 主从模型、Blacklist/\allowbreak Whitelist 黑白名单）背后的压迫隐喻与历史创伤，积极倡导中立、包容且精确的技术术语演进（如 Primary/\allowbreak Replica、Allowlist/\allowbreak Blocklist 等替代术语）。
    \item \textbf{跨语言与国际化编码支持}：语言核心全面拥抱 Unicode（UTF-8）标准，原生支持非拉丁字符进行变量命名与字符串多语言文本处理。
\end{itemize}

\subsubsection{示范性学习成果与能力目标}
\begin{enumerate}
    \item \textbf{[应用]} 在设计全新编程语言、领域专用语言或制定团队编码规约时，具备高度的伦理自觉与包容意识，主动选用包容、无冒犯性且符合全球无障碍标准的技术术语体系与语言特性。
\end{enumerate}

\section{专业素养与思维倾向}

在编程语言理论基础领域的系统学习与工程实践中，重点培育并塑造学生的以下专业特质与思维倾向：
\begin{itemize}
    \item \textbf{专业水准}：在使用编程语言特性与形式化方法构建关键系统时，严格恪守最高的工业专业标准与工程规范，确保系统完全满足安全、合规与预期目的。
    \item \textbf{创新探索}：深刻领会程序设计与形式化数学证明本质上是高度富于创造性的探索过程，善于打破常规，针对复杂工程瓶颈提出具有革新性的算法与范式解决方案。
    \item \textbf{审慎严谨}：在运用语言高级特性、类型系统标注以及形式化证明方法时展现出对细微边界、死角条件与类型推导的一丝不苟。
    \item \textbf{担当负责}：作为软件开发者，对程序在所有极端边界输入下的运行行为负有全权责任，必须具有前瞻性地预见各种潜在故障并设计完备的防御性解决方案。
    \item \textbf{坚韧执着}：在面对复杂语言实现难题、深层编译器缺陷或繁琐的形式化定理证明时，展现坚韧不拔的攻坚品质，通过反复推导、测试与迭代精化直至逼近最优解。
    \item \textbf{问责自律}：对自己在架构选型、语言特性取舍与系统实现中所做出的每一项工程决策承担完全责任。
\end{itemize}

\section{数学基础要求}

编程语言理论高度依赖离散数学与数理逻辑的深厚积淀，其核心数学前置需求如下：
\begin{itemize}
    \item \textbf{离散数学}：布尔代数、严格数学证明技术（反证法、数学归纳法、结构归纳法）、数理命题、集合论及其代数运算、映射理论、函数与二元关系（等价关系、偏序关系与良序集）、状态机与系统不变量、图论与树拓扑分析、初等组合计数、递推关系式求解、有穷状态自动机（FSM）与正则文法。
    \item \textbf{数理逻辑}：命题逻辑（否定、合取、析取、条件与双条件逻辑演算）、一阶谓词逻辑（量词、项、公式与自由变量）、逻辑推理范式（归纳推理 Induction、演绎推理 Deduction、溯因推理 Abduction）。
\end{itemize}

\section{交叉主题与跨学科共享概念}

本知识领域与计算机科学其他核心领域存在广泛而深刻的理论交织与概念共享：
\begin{itemize}
    \item \textbf{编程语言/事件驱动与响应式编程} 与 \textbf{算法/自动机与计算复杂性} 在状态机演化模型与状态转移图谱上高度重合；
    \item \textbf{编程语言/并行与分布式计算} 与 \textbf{并行分布式/并行与分布式计算} 及 \textbf{系统基础/并行性基础} 在并发安全性、活性与通信原语上高度互补；
    \item \textbf{编程语言/程序表示} 与 \textbf{体系结构/数据的机器级表示} 在指令编码与二进制表示上紧密交织；
    \item \textbf{编程语言/语言翻译与执行} 与 \textbf{体系结构/存储系统架构}、\textbf{操作系统/内存管理} 及 \textbf{安全隐私/基础安全} 在活动记录栈帧、堆分配以及防御缓冲区溢出上深度协同；
    \item \textbf{编程语言/编译器语义分析} 与 \textbf{软件开发基础/基础编程构造} 在标识符作用域消解与左值/右值语义上自然衔接；
    \item \textbf{编程语言/代码生成} 与 \textbf{体系结构/汇编级机器组织}、\textbf{体系结构/存储系统组织} 及 \textbf{操作系统/内存管理} 在目标指令发射与 ABI 调用约定上严密对接；
    \item \textbf{编程语言/形式语义学} 与 \textbf{软件工程/形式化方法} 在公理语义推导与程序正确性证明上完全贯通；
    \item \textbf{编程语言/形式化开发方法学} 与 \textbf{软件工程/形式化方法} 在证明助手使用与高安全系统完全验证上共享理论基础；
    \item \textbf{编程语言/程序分析与分析器} 与 \textbf{并行分布式/并行性基础} 在静态数据依赖分析与循环自动并行化上深度交叉。
\end{itemize}

\section{课程设计与开课规划建议}

联合工作组提供了两套典型开课打包方案，以指导各高校计算机专业合理规划教学内容：

\subsection{基础导论课程规划建议 (面向全体计算机专业，涵盖 计算机核心 与 领域核心)}
本门课程旨在系统讲授主流编程范式，确保学生深刻理解各范式核心差异，能够在现实工程中针对具体业务场景选用最适配的语言与范式。课程总学时约 50 学时，内容包含：
\begin{itemize}
    \item 编程语言/面向对象编程：5 计算机核心学时 + 3 领域核心学时
    \item 编程语言/函数式编程：4 计算机核心学时 + 4 领域核心学时
    \item 编程语言/逻辑式编程：3 核心学时
    \item 编程语言/事件驱动与响应式编程：2 计算机核心学时 + 2 领域核心学时
    \item 编程语言/脚本编程：2 计算机核心学时
    \item 编程语言/并行与分布式计算：3 计算机核心学时 + 2 领域核心学时
    \item 编程语言/类型系统：3 计算机核心学时 + 4 领域核心学时
    \item 编程语言/语言翻译与执行：4 计算机核心学时
    \item 编程语言/程序表示：3 领域核心学时
    \item 编程语言/高级编程构造：4 拓展选修学时
    \item 编程语言/硬件接口：2 拓展选修学时
    \item 编程语言/编程语言与社会伦理：1 拓展选修学时
\end{itemize}
\textbf{先修课程要求}：离散数学（布尔代数、证明技术、集合与映射、函数与关系、状态与不变量、图与树、自动机基础）。

\subsection{高级进阶课程规划建议 (面向系统软件与理论方向，专攻编译与形式语义)}
本门课程专为深化编程语言底层工程实现、高级编译器架构、运行时系统与形式语义描述而设计，内容包含：
\begin{itemize}
    \item 编程语言/语言翻译与执行：3 选修学时
    \item 编程语言/语法分析：3 选修学时
    \item 编程语言/类型系统：3 选修学时
    \item 编程语言/编译器语义分析：5 选修学时
    \item 编程语言/程序分析器：5 选修学时
    \item 编程语言/代码生成：5 选修学时
    \item 编程语言/运行时系统：4 选修学时
    \item 编程语言/语言实用学：3 选修学时
    \item 编程语言/硬件接口：3 选修学时
    \item 编程语言/形式语义学：5 选修学时
    \item 编程语言/形式化开发方法学：5 选修学时
\end{itemize}
\textbf{先修课程要求}：离散数学、数理逻辑（命题与一阶谓词逻辑）、编程语言导论基础课、熟练的高级语言编程技能（熟练掌握类型声明、数组/结构体、类/子类继承、作用域规则、分支/循环与函数调用）、数据结构（抽象数据类型、链表、树、哈希表、栈与队列）、计算机系统基础与体系结构（数字逻辑、总线、寄存器、缓存、内存层次、CPU/GPU）、操作系统基础知识（中断、线程与并发、调度与优先级、内存碎片与延迟）。

\section{真实能力任务规范}

\begin{taskbox}[任务 1：为特定工程应用科学选型适宜的编程语言与编程范式]
\textbf{【任务情境】}：面对一个涵盖高并发高吞吐网络 I/O、复杂业务状态机以及高度可维护性要求的大型商业软件工程，为研发团队进行权威的架构选型论证。\\
\textbf{【能力陈述】}：综合运用对多种编程范式（过程式、面向对象、函数式、逻辑式、事件响应式）的深厚理论认知，剖析各范式及具体编程语言在目标应用场景下的核心优势与潜在短板，科学选型并确定最契合工程需求的技术路线与范式组合。\\
\textbf{【能力领域】}：软件/应用 \quad \textbf{【能力单元】}：设计/开发/评估\\
\textbf{【所需知识单元】}：编程语言/面向对象编程、编程语言/函数式编程、编程语言/逻辑式编程、编程语言/事件驱动与响应式编程、编程语言/类型系统、编程语言/语言翻译与执行、编程语言/语言语用与运行时架构、编程语言/嵌入式计算与硬件接口、编程语言/高级编程结构\\
\textbf{【技能达标等级】}：\textbf{解释/评估}
\end{taskbox}

\begin{taskbox}[任务 2：运用脚本语言高效实现研发运维任务的自动化编排]
\textbf{【任务情境】}：针对持续集成/持续交付（CI/CD）流水线中的重复性跨平台构建、自动化测试、日志聚合清洗与微服务状态巡检任务进行自动化重构。\\
\textbf{【能力陈述】}：编写并运行高效、高鲁棒性的自动化控制脚本，灵活调用系统命令、管道重定向与正则表达式模式匹配，全自动调度和纳管各类复杂系统任务与开发运维流程。\\
\textbf{【能力领域】}：软件/系统/应用 \quad \textbf{【能力单元】}：工程开发\\
\textbf{【所需知识单元】}：编程语言/脚本编程、操作系统/文件系统 API 与底层实现\\
\textbf{【技能达标等级】}：\textbf{开发}
\end{taskbox}

\begin{taskbox}[任务 3：向研发团队深度阐释事件驱动与响应式系统的工作机理]
\textbf{【任务情境】}：团队正在重构一套基于响应式事件流的高并发物联网数据采集或复杂交互界面系统，部分初级工程师对非阻塞回调与事件循环机制存在混淆。\\
\textbf{【能力陈述】}：运用事件驱动与响应式编程的核心理论，向团队成员深入阐述事件源、事件队列、事件循环、异步监听器与背压控制的工作机制，清晰指导如何编写健壮的无死锁响应式代码。\\
\textbf{【能力领域】}：软件/应用 \quad \textbf{【能力单元】}：设计/开发\\
\textbf{【所需知识单元】}：编程语言/事件驱动与响应式编程、编程语言/嵌入式计算与硬件接口、专用平台/嵌入式平台\\
\textbf{【技能达标等级】}：\textbf{解释/评估}
\end{taskbox}

\begin{taskbox}[任务 4：撰写专业技术白皮书论证并行与分布式计算的技术收益与工程挑战]
\textbf{【任务情境】}：某企业核心业务系统面临从单机顺序批处理向多核集群并发与分布式微服务架构跃迁的关键决策节点。\\
\textbf{【能力陈述】}：撰写详实严谨的技术论证白皮书，从理论极限（阿姆达尔定律）、内存一致性代价、数据竞争防御、分布式网络延迟与节点故障容错等多维度，系统论证并行化改造带来的吞吐增益与工程实施挑战。\\
\textbf{【能力领域】}：软件/系统/应用/理论 \quad \textbf{【能力单元】}：设计/评估\\
\textbf{【所需知识单元】}：编程语言/并发与分布式计算、并行分布式/并发与分布式计算、系统基础/系统级并行\\
\textbf{【技能达标等级】}：\textbf{解释/评估}
\end{taskbox}

\begin{taskbox}[任务 5：撰写技术白皮书深度剖析源码编译翻译与运行时内存执行的全流程机制]
\textbf{【任务情境】}：团队负责构建高可靠系统软件，需要全体成员深刻理解高级代码在底层硬件上的真实足迹，以根治隐蔽的内存踩踏与悬空指针故障。\\
\textbf{【能力陈述】}：精确绘制语言核心数据与控制抽象在底层的运行时内存布局图；系统解析内存如何划分为全局数据段、只读代码段、动态堆区与栈帧活动记录；阐明递归调用与内存分配的物理映射；熟练运用调试工具定位并彻底根治内存泄露与悬空指针解引用。\\
\textbf{【能力领域】}：软件/系统 \quad \textbf{【能力单元】}：系统评估与分析\\
\textbf{【所需知识单元】}：编程语言/语言翻译与执行\\
\textbf{【技能达标等级】}：\textbf{解释/评估/应用}
\end{taskbox}

\begin{taskbox}[任务 6：深度应用现代编程语言类型系统构建高安全健壮的软件工程]
\textbf{【任务情境】}：在金融支付核心系统或高可靠航空航天控制软件研发中，杜绝任何因空指针异常、越界类型转换或并发数据竞争引发的未定义崩溃。\\
\textbf{【能力陈述】}：深度利用目标语言的静态与动态定型规则（包含代数数据类型、参数多态、所有权与生命周期借用检查），在编译阶段构建严密的不变量护栏，从数学和架构上证明系统核心模块的类型健全性与运行正确性。\\
\textbf{【能力领域】}：软件/应用 \quad \textbf{【能力单元】}：系统开发\\
\textbf{【所需知识单元】}：编程语言/类型系统\\
\textbf{【技能达标等级】}：\textbf{开发}
\end{taskbox}

\begin{taskbox}[任务 7：设计并实现将高级语言转换为底层可执行形式的编译器/解释器子系统]
\textbf{【任务情境】}：为企业内部自研领域特定语言（DSL）或游戏脚本引擎构建一套完整的端到端代码编译与执行解析工具链。\\
\textbf{【能力陈述】}：分析输入源代码文本，构建词法与上下文无关语法解析器，验证静态语义与类型正确性，生成结构化抽象语法树（AST）及三地址码中间表示，进而实施代码优化并生成高效底层机器/字节码指令。\\
\textbf{【能力领域】}：软件/理论 \quad \textbf{【能力单元】}：系统设计与开发\\
\textbf{【所需知识单元】}：编程语言/语言翻译与执行、编程语言/程序表示、编程语言/静态分析、编程语言/语法与静态分析、编程语言/程序分析与分析器、编程语言/代码生成、编程语言/基于规则的系统\\
\textbf{【技能达标等级】}：\textbf{解释/应用/评估/开发}
\end{taskbox}

\begin{taskbox}[任务 8：运用数理逻辑与形式化方法论证程序针对功能规约的严格正确性]
\textbf{【任务情境】}：为关键安全系统（如轨道交通互锁控制核心或智能合约金库）的关键算法模块提供毋庸置疑的数学逻辑完备性背书。\\
\textbf{【能力陈述】}：运用一阶谓词逻辑、公理语义学（霍尔逻辑三元组与循环不变量）以及形式化开发规约，对程序算法代码针对预设业务规格说明书进行严格的形式化数学演绎证明。\\
\textbf{【能力领域】}：软件/系统/应用/理论 \quad \textbf{【能力单元】}：开发/文档化/工程验收\\
\textbf{【所需知识单元】}：编程语言/语言语用学、编程语言/形式语义、编程语言/形式化开发方法\\
\textbf{【技能达标等级】}：\textbf{解释}
\end{taskbox}

\begin{taskbox}[任务 9：构建多语言混合系统并实现异构组件间的高效互操作与协同]
\textbf{【任务情境】}：构建一套现代化大型企业应用，底层核心采用高性能系统级语言（如 C++/Rust），业务层采用面向对象语言（如 Java），胶水调度采用脚本语言（如 Python），且融合了函数式流处理组件。\\
\textbf{【能力陈述】}：综合运用对多编程范式及底层数据物理内存排布的深刻理解，通过统一的 C-ABI、外部函数接口（FFI）以及跨语言内存对齐映射，开发出跨语言无缝协同、内存零拷贝且执行极速的高效混合多范式软件系统。\\
\textbf{【能力领域】}：软件/应用 \quad \textbf{【能力单元】}：系统设计与开发\\
\textbf{【所需知识单元】}：编程语言/面向对象编程、编程语言/函数式编程、编程语言/逻辑式编程、编程语言/事件驱动与响应式编程、编程语言/高级编程结构、编程语言/类型系统、编程语言/语言翻译与执行、编程语言/语言语用与运行时架构、编程语言/嵌入式计算与硬件接口、编程语言/并发与分布式计算、编程语言/脚本编程、编程语言/程序表示\\
\textbf{【技能达标等级】}：\textbf{解释/应用/评估/开发}
\end{taskbox}

\begin{taskbox}[任务 10：撰写高安全软件技术白皮书详述现代语言特性对安全性与健壮性的保障机制]
\textbf{【任务情境】}：安全关键行业审计部门要求企业提交全面技术论证，详细阐述自研软件系统为何能够完全杜绝高危内存破坏与并发竞态漏洞。\\
\textbf{【能力陈述】}：撰写高安全软件架构白皮书，从强静态类型系统、静态语义分析、内存自动安全回收机制、类型保持性与进展性证明等多个理论与实现维度，严密论证该软件系统如何全方位借助语言特质抵御安全风险。\\
\textbf{【能力领域】}：软件/系统/应用 \quad \textbf{【能力单元】}：设计/开发/重构改进\\
\textbf{【所需知识单元】}：编程语言/面向对象编程、编程语言/函数式编程、编程语言/逻辑式编程、编程语言/事件驱动与响应式编程、编程语言/高级编程结构、编程语言/类型系统、编程语言/语言翻译与执行、编程语言/语言语用与运行时架构、编程语言/嵌入式计算与硬件接口、编程语言/并发与分布式计算、编程语言/脚本编程、编程语言/程序表示\\
\textbf{【技能达标等级】}：\textbf{应用/评估}
\end{taskbox}

\begin{taskbox}[任务 11：撰写系统级资源优化白皮书解析程序在内存与算力占用上的高效执行机制]
\textbf{【任务情境】}：面对海量高并发请求引发的云端服务器成本高企，需对底层系统的 CPU 占用与内存驻留集进行系统性画像与极限能效重构。\\
\textbf{【能力陈述】}：撰写系统性能剖析与能效优化技术报告，从数据结构的内存紧凑对齐、避免闭包捕获内存膨胀、尾调用递归消除、数据局部性缓存命中，以及避免跨语言上下文切换开销等角度，详尽阐释并实施系统能效优化方案。\\
\textbf{【能力领域】}：软件/系统/应用 \quad \textbf{【能力单元】}：设计/开发/重构改进\\
\textbf{【所需知识单元】}：编程语言/面向对象编程、编程语言/函数式编程、编程语言/逻辑式编程、编程语言/事件驱动与响应式编程、编程语言/高级编程结构、编程语言/语言翻译与执行、编程语言/语言语用与运行时架构、编程语言/嵌入式计算与硬件接口、编程语言/并发与分布式计算、编程语言/脚本编程、编程语言/程序表示\\
\textbf{【技能达标等级】}：\textbf{应用/评估}
\end{taskbox}

\section{指导委员会与主要贡献者}

\begin{itemize}
    \item \textbf{工作组主席}：
    \begin{itemize}
        \item Michael Oudshoorn（海波因特大学，美国北卡罗来纳州）
    \end{itemize}
    \item \textbf{委员会核心成员}：
    \begin{itemize}
        \item Annette Bieniusa（凯泽斯劳滕工业大学，德国）
        \item Brijesh Dongol（萨里大学，英国）
        \item Michelle Kuttel（开普敦大学，南非）
        \item Doug Lea（纽约州立大学奥斯威戈分校，美国纽约州）
        \item James Noble（惠灵顿维多利亚大学，新西兰）
        \item Mark Marron（微软研究院，美国华盛顿州）
        \item Peter-Michael Osera（格林内尔学院，美国艾奥瓦州）
        \item Michelle Mills Strout（亚利桑那大学，美国亚利桑那州）
    \end{itemize}
    \item \textbf{特邀贡献者}：
    \begin{itemize}
        \item Alan Dearle（圣安德鲁斯大学，苏格兰）
    \end{itemize}
\end{itemize}

\section{附录：核心主题与技能要求等级对照表}

下表详细汇总了编程语言理论基础（FPL）知识领域的核心知识单元、对应主题、布鲁姆认知技能达标等级、核心学时分类以及推荐学时分配：

\begin{longtable}{p{0.18\textwidth} p{0.48\textwidth} c c c}
\toprule
\textbf{知识单元} & \textbf{核心知识点与具体主题} & \textbf{认知等级} & \textbf{核心分类} & \textbf{学时} \\
\midrule
\endfirsthead
\toprule
\textbf{知识单元} & \textbf{核心知识点与具体主题} & \textbf{认知等级} & \textbf{核心分类} & \textbf{学时} \\
\midrule
\endhead
\bottomrule
\endfoot
\multirow{4}{*}{\makecell[l]{编程语言/面向对象编程}} 
& 面向对象设计：状态与行为解耦、面向领域建模的类层次结构设计；类定义（字段、方法、构造函数）；子类化、单/多继承与方法重写；动态分派（方法调用晚期绑定）；异常处理机制 & 开发 & 计算机核心 & 3 \\
\cline{2-5}
& 面向对象封装惯用法：成员可见性与数据隐藏、仅暴露方法签名的纯接口抽象、抽象基类、特质与混入；动态属性 vs 静态属性；组合优于继承的设计原则 & 开发 & 计算机核心 & 3 \\
\cline{2-5}
& 子类型化理论：子类型多态、隐式向上类型转换、里氏行为替换原则（LSP）、子类型多态与继承实现的本质区别；通用泛型集合容器库、迭代器模式及标准构件 & 开发 & 领域核心 & 3 \\
\midrule
\multirow{3}{*}{\makecell[l]{编程语言/函数式编程}} 
& Lambda 表达式与求值机制：变量绑定、词法作用域、参数传递、嵌套 $\lambda$ 表达式与归约序；无副作用纯函数编程、引用透明性、不可变数据、递归 vs 循环 vs 管道流；结构化代数数据处理与持久化数据结构；头等与高阶函数 & 开发 & 计算机核心 & 4 \\
\cline{2-5}
& 函数闭包机制：闭包数学定义、捕获并绑定外层词法环境；闭包经典用法（回调函数、迭代器实参、通用逻辑封装）；运用闭包实现私有环境数据封装；惰性求值 vs 及早求值；高阶函数复合与动态函数生成 & 解释 & 领域核心 & 4 \\
\midrule
\makecell[l]{编程语言/逻辑式编程} 
& 逻辑式编程基础：全称与存在量词、一阶谓词逻辑（FOPL）与高阶逻辑；运用逻辑联结词表达复合关系；霍恩子句、事实、规则、目标与子目标定义；合一机制与合一算法（MGU）；合一 vs 断言 vs 算术求值；关系与函数混合；非确定性、回溯深搜与剪枝操作符；封闭世界假设 vs 开放世界假设 & 解释 & 领域核心 & 3 \\
\midrule
\makecell[l]{编程语言/脚本编程} 
& 脚本编程与自动化：分治与构件组合、进程分叉与并发、健壮错误与异常处理、I/O 流与管道重定向、系统底层命令调用与环境变量配置、文件系统测试运算符；动态容器（切片、列表推导式、关联字典）、正则表达式文本清洗、动态类型系统与动态代码求值 & 开发 & 计算机核心 & 2 \\
\midrule
\multirow{2}{*}{\makecell[l]{编程语言/事件驱动与\\响应式编程 (EDR)}} 
& 响应式编程演进：过程式顺序控制 vs 响应式事件驱动模型；响应式架构核心构件（事件源、事件流信号、监听器与分发器、事件载荷对象、适配器、回调处理器）；基于事件循环与非阻塞队列的行为演化；经典场景（GUI、移动端、机器人、高吞吐服务器）；响应式系统状态转移建模 & 开发 & 计算机核心 & 2 \\
\cline{2-5}
& 响应式框架工程落地：事件处理器与异步数据流监听器编写；事件发布者与参数泛型解耦；外部硬件事件与程序内部合成事件协同；模型--视图--控制器（MVC）及变体分层架构解耦 & 开发 & 领域核心 & 2 \\
\midrule
\multirow{2}{*}{\makecell[l]{编程语言/并行与\\分布式计算 (PDC)}} 
& 并发与分布式基础：安全性与活性属性；数据竞争、前置条件与终止性判定；并发模型（Actor 模型、过程式/响应式并发、异步编程抽象 async/await、Promise、数据并行）；并发语义（交换律、因果序、独立性、一致性、原子性、分布式共识）；通信与协调（消息传递通道、共享内存互斥、管程 Monitors、守卫语句） & 开发 & 计算机核心 & 3 \\
\cline{2-5}
& 高级并发支持与系统开销：Future 异步凭据抽象与链式复合；数据并行语言级支持（forall 循环、自动展开、Map-Reduce 向量化）；硬件内存一致性模型（顺序一致性、弱内存模型）对语义与编译优化的制约；RESTful API 微服务架构；消息传递序列化开销、通信与计算粒度平衡；函数式无锁并发 & 解释 & 领域核心 & 2 \\
\midrule
\multirow{2}{*}{\makecell[l]{编程语言/类型系统}} 
& 类型作为值的集合及操作集合；基元类型与复合代数类型（积类型、和类型、数组、函数空间、引用）；类型与变量绑定；类型安全作为程序正确性防线；静态类型 vs 动态类型；结构等价 vs 名义等价；静态类型的目标与局限（早期拦截缺陷）；泛型与参数多态（多态四分类、单态 vs 多态、模板单态化 vs 类型擦除） & 开发 & 计算机核心 & 3 \\
\cline{2-5}
& 静态与动态类型的互补工程价值；形式化定型规则（函数、积类型、和类型的定型演绎推导）；基于合一的最一般类型推导算法（Hindley-Milner 系统）；包含 Let-多态的 Lambda 演算；可判定性边界与静态分析；工业级前沿类型系统（Rust 所有权、借用检查与生命周期系统） & 开发 & 领域核心 & 4 \\
\midrule
\makecell[l]{编程语言/语言翻译与\\执行 (LTE)} 
& 执行三元谱系（解释执行 vs 本地机器码 vs IR 字节码）；端到端翻译流水线（词法、语法、语义检查、优化、生成、链接、加载、执行）；上下文无关文法 BNF/EBNF 规范；语法分析树推导；文法二义性消歧；本地机器执行 vs 虚拟机环境；动态加载与 JIT 即时编译；控制流图（CFG）构造；运行时内存物理布局（全局静态区、代码段、调用栈、堆区）；活动记录栈帧（返回地址、帧指针、动态链、静态链）；函数调用汇编协议；循环展开与尾调用消除；物理内存管理（栈自动回收、手工堆分配、基于对象图可达性的自动垃圾回收 GC） & 解释 & 计算机核心 & 4 \\
\midrule
\makecell[l]{编程语言/程序表示} 
& 形式语言表征工具（BNF、EBNF、正则表达式）；以程序为输入的元程序体系（编译器、解释器、分析器、文档生成器）；语言构成要素（字母表、记号、语句、语法与语义）；抽象语法树（AST）与具体语法解耦；核心通用编程抽象（变量、控制语句、模块）；可变性权衡（原地覆盖 vs 纯值复制、记忆化）；左值（L-value）与右值（R-value）；变量作用域与生存期（静态词法作用域 vs 动态作用域、可见性遮蔽、副作用）；环境与存储二元数学模型；参数传递全景对比（传值、传引用、传值-结果、传名调用、按需调用）；多级中间表示（AST、SSA 形式）；汇编与机器指令表示；Lambda 演算基础（绑定、$\alpha$-变换、$\beta$-归约） & 解释 & 领域核心 & 3 \\
\end{longtable}
\chapter{图形学与交互技术}

\section{领域导言与演进背景}

\begin{preamblebox}[知识领域概述与变革]
计算机图形学是利用计算机生成与处理数字图像的核心学科，其本质可以被视为“以计算为媒介实现视觉感知与人机信息沟通的科学”。从最初的线框显示到当今照片级拟真画质，计算机图形学的应用已渗透至现代计算的方方面面：机器学习、医学成像、工程辅助设计、科学与信息数据可视化、三维动画电影、电影视觉特效（VFX）、沉浸式训练仿真器以及高交互数字游戏等。

在传统的计算机科学本科教育中，计算机图形学主要聚焦于三维渲染管线、线性代数与空间变换、几何光学物理现象学模拟等经典主题。然而在现代高级培养阶段，全球顶尖高校越来越倾向于开设多门高度分化且面向前沿的进阶专向课程，例如游戏工程引擎架构、高级角色动画、科学可视化、物理/实体计算以及涵盖虚拟现实、增强现实与混合现实的沉浸式计算。该领域与其他计算机专业核心领域存在极其深厚的交叉与重叠，尤其是：人工智能（AI）、人机交互（HCI）、并行与分布式计算（PDC）、专用平台开发（SPD）以及社会、伦理与专业素养（SEP）。

为了使计算机科学专业学生能够娴熟地运用并开发计算机图形系统，教学必须深入剖析一系列具体的工程与底层实现机制，包括人类视觉感知与认知心理学、数字图像与几何文件编码规范、底层硬件抽象接口以及现代图形渲染应用程序接口。通过理论推导与工程实践的紧密结合，学生将具备作为高阶用户和底层开发者来理解、评估、重构并实现前沿图形与交互技术的核心能力。鉴于该领域技术演进一日千里，CS2023 工作组在制定本大纲时特意避免过度限定具体商业 API 或特定编程语言，文中所提及的开发接口、图形系统与开发工具代表了截至 2022 年前后的典型工业界与学术界主流范式。

\textbf{自 CS2013 以来的核心演进与重大变革：}
\begin{enumerate}
    \item \textbf{领域命名与 ACM SIGGRAPH 国际标准对齐}：为更契合现代前沿计算形态，CS2023 将 CC2013 中的“图形学与可视化 (GV)”正式更名为“图形学与交互技术 (GIT)”，实现与计算机图形学全球顶级学术组织 ACM SIGGRAPH 的命名统一。
    \item \textbf{移动端普及与算力下沉驱动的全民视觉化}：自 2013 年以来，机器学习、计算机视觉、数据科学与 AI 的爆发，以及从智能汽车到家用物联网设备广泛内嵌的传感器，深刻改变了图形交互的生态。如今移动智能手机的算力飞跃，使全球数十亿普通大众既是高画质图形与交互的消费者，更是创作者。2013 年仅能在高性能桌面图形工作站上运行的复杂动画骨骼蒙皮、流体物理仿真、实时全局光照与沉浸式应用，如今已能在手持设备上流畅运行。
    \item \textbf{泛在图像与文化产业数字化重塑}：十年来全球数字数据爆发式增长，新闻媒体与科研机构每天都在利用高维信息可视化揭示自然与社会规律；移动数字游戏的全球产值已经超越电影与音乐两大传统娱乐工业的总和；而在现代电影工业中，计算机生成图像（CGI）已成为几乎所有商业影视的不可或缺之基石。
    \item \textbf{全面扩充五大前沿核心知识单元}：为捕捉学科外延的剧烈拓展，在原有知识单元（基本概念、基础渲染、几何建模、高级着色、计算机动画、数据可视化）之上，全新增设了以下 5 个知识单元：
    \begin{itemize}
        \item \textbf{沉浸式技术 (Immersion: MR, AR, VR)}：系统建立空间计算、动光延迟抑制与虚实融合机制。
        \item \textbf{交互技术}：涵盖事件驱动范式、多模态融合、力反馈外骨骼与脑机接口前沿。
        \item \textbf{数字图像处理}：融合形态学、空频域滤波、图像修复以及深度卷积与 Transformer 视觉前沿。
        \item \textbf{实体与物理计算}：贯通传感器采集、执行器控制、数控数字制造（CAD/CAM/3D打印）与物联网系统。
        \item \textbf{计算机物理仿真}：确立面向图形视觉可信度的质点弹簧、SPH 粒子流体、欧拉网格烟火与碰撞动力学。
    \end{itemize}
    \item \textbf{深度审视图形与视觉技术的伦理困境}：人工智能生成内容（AIGC）与高保真图形合成技术的结合，带来了前所未有的伦理考验。师生必须直面在自动驾驶机器视觉感知失效、深空遥感伪影、计算摄影学“脑补”失真、基于人脸与身形的深度伪造（Deepfakes）以及人脸识别在不同种族发型上的算法偏见与不公等现实伦理痛点。
\end{enumerate}
\end{preamblebox}

\section{核心学时分配建议}

在 CS2023“向日葵模型”架构下，图形学与交互技术领域各单元的核心推荐学时分配如下表所示：

\begin{table}[H]
\centering
\setlength{\tabcolsep}{3.5pt}\fontsize{8pt}{9.6pt}\selectfont\begin{tabularx}{\columnwidth}{Y c c}
\toprule
\textbf{知识单元} & \textbf{计算机核心} & \textbf{领域核心} \\
\midrule
基本概念 & 4 & -- \\
基础渲染技术 & -- & 18 \\
几何建模 & -- & 6 \\
高级着色与进阶渲染 & -- & 6 \\
计算机动画 & -- & 6 \\
数据与信息可视化 & -- & 6 \\
沉浸式技术 & -- & 6 \\
交互技术 & -- & 6 \\
数字图像处理 & -- & 6 \\
实体与物理计算 & -- & 6 \\
计算机物理仿真 & -- & 6 \\
\midrule
\textbf{合计} & \textbf{4 学时} & \textbf{72 学时} \\
\bottomrule
\end{tabularx}
\end{table}

\section{知识单元与详细大纲}

\subsection{计算机图形学基本概念}
\textbf{【学时建议】}：4 计算机核心学时。

对于几乎所有计算机专业毕业生与软件开发者而言，理解人机视觉交互的底层规律都是必不可少的素养。本单元的核心内容既可作为独立图形学课程的先导篇章，亦可融入大学计算机导论或软件开发基础（SDF）课程中讲授。

\subsubsection{计算机核心主题}
\begin{itemize}
    \item \textbf{娱乐、工商业与科学应用图景}：视觉特效（VFX）、机器学习与计算机视觉交叉、现代图形用户界面（GUI）、视频编辑管线、游戏引擎架构、计算机辅助设计与制造（CAD/CAM）、空间与时变数据可视化、虚拟/增强/混合现实（VR/AR/MR）。
    \item \textbf{人类视觉生理系统与感知模型}：
    \begin{itemize}
        \item 视网膜感光锥状细胞与三刺激响应（RGB 三原色理论）。
        \item 经典“眼球作为摄像机”的透视投影成像几何。
        \item 视觉暂留特性、临界闪烁融合频率、帧率与动态模糊。
        \item 对比度感知机制：边缘增强特性、马赫带效应、色彩抖动与空间走样。
        \item 人眼对光强的非线性对数响应：高动态范围（HDR）与色调映射必要性。
        \item 双目立体视觉（Binocular Vision 与双目视差 Stereo）。
        \item 视觉无障碍包容性设计：红绿色盲等色觉辨认障碍、频闪光敏性癫痫规避、单目弱视交互补偿（与 SEP、HCI 深度交叉）。
    \end{itemize}
    \item \textbf{模拟物理信号的数字化机制}：
    \begin{itemize}
        \item 光栅图元表征与矢量数学表征的本质对立与时空权衡：多边形面片、体积标量场与连续实体对象。
        \item 分辨率度量：显示屏幕的像素密度（PPI）与激光打印设备的点密度（DPI）。
        \item 连续空间的离散空间采样与色彩强度量化。
    \end{itemize}
    \item \textbf{标准媒体编码与格式}：
    \begin{itemize}
        \item 光栅图像格式：无损压缩编码（如 PNG、无损 TIFF）与有损压缩编码（如 JPEG 离散余弦变换）。
        \item 矢量数据格式规范（如 SVG、PDF 路径）。
    \end{itemize}
    \item \textbf{经典色彩模型与空间}：加色模型（RGB，用于自发光显示器）、减色模型（CMYK，用于印刷油墨染料）、面向人类感知直觉的色彩模型（HSV/HSL，色相、饱和度、明度）。
    \item \textbf{存储空间占用与动态重算时延的权衡}：光栅高保真缓存与矢量数学解析重绘的计算复杂度对比。
    \item \textbf{动画的离散本质}：连续时空运动在离散时序静止画面帧序列中的表征机理。
    \item \textbf{伦理、版权与社会议题 (SEP)}：换脸与深度伪造（Deepfakes）技术、人脸识别中的偏见、个人生物特征隐私合规与视觉作品知识产权保护。
\end{itemize}

\subsubsection{示范性学习成果与能力目标}
\begin{enumerate}
    \item \textbf{[识别]} 列举并系统剖析向人类用户呈现数字信息的通用技术方案（包括计算机图形学与声学技术）。
    \item \textbf{[解释]} 阐述连续模拟物理信号如何通过离散采样合理表征为数字信号（例如连续光学图像如何量化表征为二维像素阵列）。
    \item \textbf{[计算]} 给定目标分辨率、颜色深度与通道数，精确计算未压缩全彩光栅图像所占用的物理内存容量。
    \item \textbf{[开发]} 绘制或编程构建图表，演示人类视觉感知系统的生理极限如何直接制约模拟信号数字化采样参数的工程决策。
    \item \textbf{[设计]} 针对标准人机交互界面，设计一套能够适应红绿色盲或色弱用户的替代色彩与形状双编码无障碍界面方案。
    \item \textbf{[开发]} 调用行业标准图形 API（如 HTML5 Canvas、Qt 或基础 OpenGL），构建一个包含交互图元的简易图形用户界面。
    \item \textbf{[权衡]} 论证在何种工程情境下应当选用 JPG、PNG、MP3、MP4 或 GIF 文件格式，并从压缩机理与感知失真角度给出深度技术论据。
    \item \textbf{[评估]} 指出通用媒体格式中采用的典型有损与无损压缩算法实例，并分析其香农信息论与心理物理学依据。
    \item \textbf{[解释]} 系统阐释常见色彩空间（RGB、CMYK、HSV）的坐标基底，并说明其在物理显示输出设备中的硬件映射关系。
    \item \textbf{[分析]} 深入辨析在图像存储中“完整持久化光栅数据”与“仅存储关键几何信息并在渲染时动态重建”之间的时空权衡边界。
    \item \textbf{[计算]} 给定时长 $n$ 秒、帧率 $f$ 帧/秒及色彩分辨率 $r$，计算未压缩视频流的实时内存带宽与物理存储开销。
    \item \textbf{[解释]} 阐释人类视觉生理机制中从离散静态定格画面序列生成连续平滑运动感知的基本过程（临界闪烁融合机理）。
    \item \textbf{[解释]} 举例说明对现实模拟物理世界进行数字离散采样时可能出现的视觉虚假走样效应（如莫尔条纹、车轮倒转效应）。
\end{enumerate}

\subsection{数据与信息可视化}
\textbf{【学时建议】}：6 领域核心学时。

可视化是为抽象数据建立视觉图形表征的核心过程。本单元与人机交互（HCI）及计算科学（CN）存在极强的内在关联，涉及目标受众的认知特征评估与交互界面的效能测度。

\subsubsection{领域核心主题}
\begin{itemize}
    \item \textbf{数据可视化与信息可视化辨析}：科学计算数据场（强空间绑定）与抽象结构化信息（非物理空间实体）的可视化差异。
    \item \textbf{多模态数据类型的可视化表征}：
    \begin{itemize}
        \item 二维/三维标量场数据：等值线、等值面、伪彩色映射。
        \item 矢量场与流场数据：方向线、流线、脉线、迹线与粒子轨迹跟踪。
        \item 时变时序多维动态数据场。
        \item 超高维多元特征数据：散点图矩阵（SPLOM）、平行坐标轴。
        \item 非空间关系型数据：层次树状图、力导向网络拓扑图、词云与文本语义演化流。
    \end{itemize}
    \item \textbf{高阶可视化变换技术}：高维降维投影（PCA 主成分分析、t-SNE、UMAP）、多元数据空间剖分与交互联动刷选（Brushing \& Linking）。
    \item \textbf{直接体绘制算法}：光线投射算法、传输函数的色彩与不透明度分配、体素分割与边缘增强。
    \item \textbf{通用科学与空间数据存储规范}：HDF5、netCDF、GeoTIFF、原始二进制连续流、结构化 CSV、规范化文本序列。
    \item \textbf{工业界与学术界主流可视化工具链}：R 语言（ggplot2）、Processing、D3.js、地理信息系统（GIS）、MATLAB、IDL、Python 可视化生态（Matplotlib、Seaborn、VTK）。
    \item \textbf{驱动视觉抽象的认知心理学底座}：视觉感知信道理论（位置、长度、角度、面积、颜色饱和度）、贝特定律、格式塔心理学原则、色彩感知调色板（感知均匀色阶如 Viridis、区分性色阶）。
    \item \textbf{端到端可视化工程设计原则}：
    \begin{itemize}
        \item 目标导向性：科研探索与公众科学传播的区别。
        \item 受众分层定位：专业领域科学家与大众公众。
        \item 负责任与道德规范的可视化：防范虚假视觉修辞（坐标轴刻度截断欺骗、过度平滑消除关键极值、数据清洗阶段的恶意填补插值）；识别真数据引起的视觉误导（如高频运动引起的频闪混叠导致风扇叶片视觉停转）。
    \end{itemize}
    \item \textbf{可视化有效性评估方法}：任务完成时间、准确率眼动追踪分析与主观认知负荷评估量表。
    \item \textbf{可视化认知偏差与偏见}：知觉锚定效应、框架效应与视觉突显偏差。
\end{itemize}

\subsubsection{示范性学习成果与能力目标}
\begin{enumerate}
    \item \textbf{[开发]} 编程实现经典基础可视化算法（如基于移动立方体 Marching Cubes 的等值面提取、离散标量场传递函数着色映射）。
    \item \textbf{[评估]} 在计算时间开销、显存带宽占用与视觉呈现保真度之间，量化评估不同可视化算法的工程权衡。
    \item \textbf{[设计]} 针对特定领域的数据拓扑特征、具体分析任务与目标受众类型，制定兼具美学价值与认知严谨性的可视化设计系统方案。
    \item \textbf{[分析]} 剖析给定工业级或科学可视化图表在执行特定决策任务时的有效性与认知效率，指出潜在设计缺陷。
    \item \textbf{[评估]} 设计一套严谨的用户体验实验流程，以定量与定性手段科学评测某可视化算法或软件系统的实际效用。
    \item \textbf{[识别]} 调研并归纳可视化技术在现代医学断层成像诊断、气象流场预警、天体物理模拟及金融风险空间地理分析中的应用范式。
\end{enumerate}

\subsection{基础渲染技术}
\textbf{【学时建议】}：18 领域核心学时。

本单元涵盖计算机图形学的核心物理与数学基础，是所有高阶计算机图形学、游戏开发与三维工程软件的共同理论基石。

\subsubsection{领域核心主题}
\begin{itemize}
    \item \textbf{现代计算机图形渲染管线}：输入装配、顶点着色、几何与曲面细分处理、图元剪裁与屏幕映射、光栅化、片元/像素着色、深度测试与模板测试、Alpha 混合及帧缓冲输出。
    \item \textbf{自然界光传输物理学与数值积分}：辐射通量、辐照度与辐亮度；光源辐射、介质散射、能量守恒定律；渲染方程的半球立体角积分形式及其向离散数值积分的转换。
    \item \textbf{正向渲染与反向渲染架构}：基于物体空间投影的光栅化与基于图像空间像素追踪的光线投射之范式对比与计算复杂度对偶性。
    \item \textbf{三维多边形网格表征}：三角网格拓扑、顶点缓冲区（VBO）、顶点索引缓冲区（EBO）、顶点法向量与切线空间。
    \item \textbf{几何光学透视与针孔投影模型}：相似三角形原理、透视除法、视场角（FOV）与平截头体构建。
    \item \textbf{仿射变换与三维齐次坐标系系统}：
    \begin{itemize}
        \item 局部坐标系、世界坐标系、观察/相机坐标系、裁剪坐标系与规范化设备坐标系（NDC）。
        \item 平移、旋转（Rodrigues 公式）、缩放矩阵以及 $4\times 4$ 齐次变换矩阵的非交换性链式复合。
    \end{itemize}
    \item \textbf{光线追踪基础算法}：主光线生成、光线与几何图元（球体、平面、三角形）求交、阴影感知射线检测、理想镜面反射与折射（Snell 定律与 Fresnel 方程）。
    \item \textbf{可见性计算与遮挡消隐机制}：深度缓冲算法（$Z$-buffer / Depth Buffering）、画家算法（Painter's Algorithm 深度排序局限性）、基于光线穿透的消隐判断。
    \item \textbf{三角形基元光栅化算法}：基于平面半空间测试与重心坐标的图元覆盖率判定与属性线性透视校正插值。
    \item \textbf{基于现代可编程着色器 API 的工程渲染}：现代可编程管线架构（如 OpenGL 3.3+ Core Profile、Vulkan 或 DirectX），着色语言编程与Uniform缓冲对象。
    \item \textbf{纹理映射与采样滤波理论}：UV 纹理坐标映射、纹理包裹模式、双线性插值与双三次插值、纹理放大与缩小走样、三线性多级渐远纹理映射与各向异性过滤。
    \item \textbf{空间加速数据结构在渲染中的应用}：轴对齐包围盒（AABB）、包围盒层次结构（BVH）、八叉树与 KD-Tree 在光线求交与视锥剔除中的应用。
    \item \textbf{采样理论与反走样技术}：香农采样定理、奈奎斯特采样频率极限、超采样反走样（SSAA）、多重采样反走样（MSAA）与时间序列抗锯齿（TAA）初步。
    \item \textbf{层次化场景图模型}：树状拓扑层次关系、局部变换向全局变换的递归遍历累积。
\end{itemize}

\subsubsection{示范性学习成果与能力目标}
\begin{enumerate}
    \item \textbf{[解释]} 绘制光线能量在物理场景中从光源发射、表面复杂多重散射并最终被虚拟镜头接收的传输图解，阐明其与半球数值积分的数理对应关系。
    \item \textbf{[描述]} 详细叙述现代硬件图形渲染管线的各个核心阶段及其在 GPU 硬件流水线中的协作逻辑。
    \item \textbf{[对比]} 深入对比正向渲染（光栅化管线）与反向渲染（光线追踪）在处理海量动态光源与复杂几何遮挡时的时空算力开销差异。
    \item \textbf{[开发]} 编写独立程序，在图形窗口中高效绘制二维动态矢量几何形状。
    \item \textbf{[开发]} 基于可编程渲染管线，编写并运行能够加载并显示交互式三维网格模型的渲染系统。
    \item \textbf{[推导]} 利用相似三角形几何原理，推导出将三维坐标点 $(x, y, z)$ 投射映射到投影平面点 $(x/z, y/z, 1)$ 的线性透视除法数学公式。
    \item \textbf{[计算]} 给定由平移、旋转与缩放组成的复合仿射矩阵，精确计算二维与三维几何顶点经矩阵变换后的新坐标。
    \item \textbf{[应用]} 掌握三维笛卡尔坐标系统，推演将二维变换算法扩展至三维齐次变换矩阵所需的维数扩展与基底变换法则。
    \item \textbf{[解释]} 深刻阐释纹理映射的基本原理，以及欠采样与过采样诱发走样的物理机理和 MIP-mapping 的消除策略。
    \item \textbf{[评估]} 在处理复杂三维几何遮挡与可见性判断问题时，系统横向评测光线追踪与 $Z$-buffer 算法的优劣。
    \item \textbf{[开发]} 编程实现对二维几何图元的平移、旋转、缩放复合变换以及齐次视锥剪裁（如 Cohen-Sutherland 或 Liang-Barsky 算法）。
    \item \textbf{[开发]} 基于现代着色器接口（如现代 OpenGL / WebGL），利用顶点缓冲（VBO）、索引缓冲（EBO）与编写 GLSL 顶点/片元着色器，实现一个完整的实时着色器渲染器。
    \item \textbf{[对比]} 深入对比 Phong 经验光照模型、Gouraud 顶点光照插值与基于物理的 Microfacet 材质渲染在真实感上的阶跃演进。
    \item \textbf{[计算]} 根据屏幕显示分辨率、每像素颜色位深与通道格式，精确计算双缓冲图形输出所需的物理显存带宽与容量。
    \item \textbf{[计算]} 根据显示器垂直刷新率（如 60Hz、144Hz）与几何三角形吞吐量，量化推导光栅化渲染管线在单帧时间预算（如 16.6ms、6.9ms）内的算力瓶颈。
\end{enumerate}

\subsection{几何建模}
\textbf{【学时建议】}：6 领域核心学时。

几何建模涵盖二维形状与三维立体几何实体的数学表征、程序化创建、拓扑交互与编辑操作。

\subsubsection{领域核心主题}
\begin{itemize}
    \item \textbf{空间几何原语计算与邻近性测试}：点-线-面三维求交、射线求交、距离场计算、包围盒快速相交与分离轴判定定理（SAT）。
    \item \textbf{三维表面表征与网格几何模型}：
    \begin{itemize}
        \item 几何细分与镶嵌（Tessellation，曲面多边形化）。
        \item 网格拓扑数据结构（半边结构 Half-Edge、翼边结构 Winged-Edge）、网格光顺与平滑算法。
        \item 经典网格生成技术：平面点集 Delaunay 三角剖分、三维隐式等值面提取移动立方体算法。
        \item 参数化多项式曲线与曲面：Hermite 曲线、Bézier 曲线与曲面特征多边形网格、B-样条与非均匀有理 B-样条（NURBS）。
        \item 隐式几何曲面（Implicit Surfaces，代数曲面方程、融合电位球 Metaballs）。
        \item 空间二叉分割树（BSP 树）与场景几何划分。
    \end{itemize}
    \item \textbf{体素与体积几何模型}：
    \begin{itemize}
        \item 三维规则体素网格与三维点云数据模型。
        \item 有符号距离场 (SDF)解析与网格交互融合。
        \item 层次化稀疏体素网格架构（如 OpenVDB 拓扑与存储优化）。
        \item 构造实体几何 (CSG)：正方体、球体、圆柱体等基本几何实体的布尔交、并、差集合运算。
    \end{itemize}
    \item \textbf{过程式与生成式几何建模}：
    \begin{itemize}
        \item 分形几何（Fractals，自相似性 Mandelbrot 集合、分形地形生成）。
        \item 文法驱动的生成式系统：L-系统（Lindenmayer Systems，植物分支形态生成，与编程语言形式文法交叉）。
        \item 基于深度学习的生成式几何模型（如 NeRF 神经辐射场与 3D 高斯泼溅 Gaussian Splatting 初步）。
    \end{itemize}
    \item \textbf{弹性变形与自由形态变形}：
    \begin{itemize}
        \item 物理准静态几何变形方法。
        \item 双调和曲面形变场捕获。
    \end{itemize}
    \item \textbf{多分辨率网格模型与几何简化}：多细节层次（LOD）、基于二次误差度量（QEM）的边折叠网格简化算法。
    \item \textbf{三维反求几何重建技术 (3D Reconstruction)}：基于激光雷达点云与多视角立体视觉（MVS）的表面拓扑重构算法。
\end{itemize}

\subsubsection{示范性学习成果与能力目标}
\begin{enumerate}
    \item \textbf{[对比]} 深入对比参数化曲面（如 Bézier / NURBS）与隐式代数曲面（如 SDF / CSG）在几何求交、渲染投射与拓扑编辑上的优劣与计算特性。
    \item \textbf{[开发]} 通过多边形表面曲面细分与半边拓扑缝合，独立构建基础多面体几何模型。
    \item \textbf{[开发]} 编程实现移动立方体算法，从三维隐式等值面场或体素场中提取闭合的连续三角多边形网格。
    \item \textbf{[开发]} 运用过程式规则或分形算法（如分形布朗运动 fBm 或细分菱形-方形算法），程序化生成高起伏逼真的分形山脉地形。
    \item \textbf{[开发]} 从工业级激光三维扫描仪采集的散乱无序空间点云中，编程实现泊松表面重构算法生成光滑闭合网格。
    \item \textbf{[应用]} 基于基础几何图元（如单位立方体、二次曲面球体与圆柱），通过 CSG 树状布尔交并差运算构建复杂机械零件模型。
    \item \textbf{[评估]} 全面对比多边形网格、隐式距离场与体素体数据在内存消耗、几何拓扑修改延迟与最终渲染光照质感上的技术权衡。
\end{enumerate}

\subsection{高级着色与进阶渲染}
\textbf{【学时建议】}：6 领域核心学时。

本单元深入探究真实物理光线传输机制的精确求解算法、微表面微观物理材质建模以及现代图形处理器（GPU）硬件架构的性能极值优化。

\subsubsection{领域核心主题}
\begin{itemize}
    \item \textbf{渲染方程的高阶数值求解与近似算法}：
    \begin{itemize}
        \item 分布式光线追踪与纯蒙特卡洛积分收敛原理。
        \item 经典单向蒙特卡洛路径追踪算法与重要性采样。
        \item 光子映射技术（Photon Mapping，处理焦散效应 Caustics 与漫反射间接光照）。
        \item 双向路径追踪 (BDPT)。
        \item 都会光线传输算法（Metropolis Light Transport, MLT，基于马尔可夫链蒙特卡洛 MCMC 状态空间探索）。
    \end{itemize}
    \item \textbf{真实物理成像维度与时间/镜头集成}：
    \begin{itemize}
        \item 时间域采样：相机快门开启过程中的时间积分与真实动态模糊。
        \item 空间光学透镜系统：有限薄透镜孔径成像、景深与对焦控制。
        \item 连续光学光谱频率：全光谱光传输与色散棱镜彩虹物理效果。
    \end{itemize}
    \item \textbf{高级阴影计算体系}：阴影贴图技术、采样自走样消除、透视阴影修正、级联阴影贴图 (CSM)与方差阴影贴图 (VSM)。
    \item \textbf{几何遮挡剔除技术}：硬件遮挡查询、潜在可见集（PVS）与软件光栅化剔除。
    \item \textbf{基于物理的双向散射分布函数 (BSDF) 与微表面理论}：
    \begin{itemize}
        \item 双向反射分布函数（BRDF）与双向透射分布函数（BTDF）。
        \item 经典微表面理论：法线分布函数（NDF，如 GGX / Trowbridge-Reitz）、几何遮蔽函数（Smith 几何项）与菲涅尔项（Schlick 近似公式）。
    \end{itemize}
    \item \textbf{次表面散射物理模拟}：偶极子扩散近似模型及其在皮肤、蜡质、玉石等半透明多孔介质渲染中的应用。
    \item \textbf{真实多面体与多边形区域面光源采样}：多边形光源的解析解（如基于 LTC 线性变换余弦）与软阴影半影区物理过渡。
    \item \textbf{层次化深度缓冲加速}：GPU 硬件 Early-Z 早期深度剔除流水线。
    \item \textbf{光场理论与基于图像的渲染技术}：全光函数、视差贴图、环境光照映射（IBL）与神经辐射场呈现机理。
    \item \textbf{非真实感艺术风格渲染}：卡通着色、轮廓边缘检测（Sobel 与法线/深度不连续性提取）、手绘风格与素描笔触合成。
    \item \textbf{现代 GPU 高性能计算与微架构解剖}：流式多处理器（SM）组织、Warp 线程束调度、显存访存合并、硬件光线追踪加速内核（如 RT Core BVH 遍历与求交引擎）。
    \item \textbf{人类视觉感知在高阶着色中的映射}：明暗适应性动态调节、人眼对光强高频噪声的掩蔽效应、频闪融合与高帧率刷新对真实感的塑造。
\end{itemize}

\subsubsection{示范性学习成果与能力目标}
\begin{enumerate}
    \item \textbf{[证明]} 通过数学推导，证明蒙特卡洛路径追踪算法如何通过对全场景表面立体角的随机光线采样，逼近渲染方程的真实解析解。
    \item \textbf{[证明]} 分析并严格证明某种渲染数值积分算法的理论属性（例如：完全性、统计一致性 Consistent 与无偏性 Unbiased）。
    \item \textbf{[分析]} 定量分析一个典型片元着色算法（如带软阴影过滤的 PBR 着色）在特定分辨率下的显存带宽消耗与浮点计算吞吐量需求。
    \item \textbf{[开发]} 基于现代实时渲染 API，独立编程实现一种非平凡的高阶着色算法（如卡通着色结合轮廓提取、四阶级联阴影贴图 CSM 或实时环境光遮蔽 SSAO）。
    \item \textbf{[应用]} 说明特定传统绘画艺术技法（如明暗对照法 Chiaroscuro、水彩晕染效果）如何在计算机着色器算法中得到精准的物理与数学映射。
    \item \textbf{[分析]} 给定一幅高拟真计算机合成图像，准确反向分析识别其在生成时所运用的特定核心图形学着色技术（如次表面散射、菲涅尔反射、全局间接焦散）。
    \item \textbf{[开发]} 脱离任何现代上层图形 API，在最基础的逐像素缓冲区纯 CPU 内存层面，手工编写代码实现光栅化与光照着色。
    \item \textbf{[开发]} 编写一个支持任意三角形三维场景的递归光线追踪渲染器，完整支持经典 Phong/Blinn-Phong BRDF 模型、理想镜面高阶反射与折射光线跟踪。
\end{enumerate}

\subsection{计算机动画}
\textbf{【学时建议】}：6 领域核心学时。

本单元涵盖二维与三维虚拟对象随时间演化发生几何形变、骨骼驱动与空间运动的核心数学规律、算法实现及艺术表现原则。

\subsubsection{领域核心主题}
\begin{itemize}
    \item \textbf{经典动画十二基本法则}：
    \begin{itemize}
        \item 挤压与拉伸（Squash and Stretch，维持体积恒定下的弹性形变）。
        \item 时间节奏与运动节拍。
        \item 预备动作（Anticipation，唤起受众对后续主要动作的心理预期）。
        \item 演出布局与视线引导。
        \item 跟随动作与重叠动作（Follow Through and Overlapping Action，动量滞后传递）。
        \item 逐帧绘制与原画关键帧推进。
        \item 慢入与慢出（Slow In and Slow Out，加减速物理缓冲惯性）。
        \item 弧线运动轨迹（Arcs，自然生物关节转动的几何规律）。
        \item 夸张手法与视觉感染力。
    \end{itemize}
    \item \textbf{关键帧动画技术与插值数学}：
    \begin{itemize}
        \item 关键帧标定与时序曲线编辑。
        \item 经典插值数学方法：线性插值、四元数球面线性插值、三次样条插值与贝塞尔加速度平滑曲线。
    \end{itemize}
    \item \textbf{正向运动学与逆向运动学}：
    \begin{itemize}
        \item 正向运动学 (FK)：骨骼层级树状结构、局部变换矩阵沿父子节点的向下递归累乘。
        \item 逆向运动学 (IK)：末端执行器约束求解、循环坐标下降法（CCD）、FABRIK 算法以及基于雅可比矩阵伪逆的数值优化求解。
    \end{itemize}
    \item \textbf{骨骼蒙皮算法族}：
    \begin{itemize}
        \item 蒙皮权重分配与顶点绑定。
        \item 线性混合蒙皮 (LBS)及其固有的关节点“糖果纸包装扭曲”与体积丢失缺陷。
        \item 双四元数蒙皮 (DQS)在刚性体积保持上的数学优势。
    \end{itemize}
    \item \textbf{骨骼角色装配绑定体系}：骨骼层级约束、反向控制器、极向量约束。
    \item \textbf{混合变形与表情系统}：Blendshapes（形态键/融合成形）、姿态空间形变（Pose Space Deformation, PSD，用于矫正肌肉挤压拉伸隆起）。
    \item \textbf{空间几何变换系统与旋转表征体系}：
    \begin{itemize}
        \item 三维平移与各向异性缩放/剪切。
        \item 三维空间旋转数学表征之比较：欧拉角及其固有的万向节死锁现象；四元数（Quaternions，单位四元数乘法与旋转向量旋转）；轴角表示法与李代数指数映射。
        \item 变换矩阵的链式乘法顺序对最终姿态的决定性影响（SRT 顺序 vs XYZ 轴顺序）。
    \end{itemize}
    \item \textbf{虚拟摄影机轨迹动画}：观察注视点（Look-at 矩阵构建与时序插值跟踪）、变焦与景深对焦动画。
    \item \textbf{动作捕捉技术}：
    \begin{itemize}
        \item 光学反光标记点追踪、惯性传感器（IMU）追踪系统设置与数据标定管线。
        \item 动作捕捉伦理与社会规范：捕捉数据的隐私知情权；坚决摒弃“默认通用行走步态”的狭隘工程设定——现实世界中不存在所谓“标准人类步态”，系统设计必须兼顾残障、老龄与多元生理群体的非标运动形态。
    \end{itemize}
\end{itemize}

\subsubsection{示范性学习成果与能力目标}
\begin{enumerate}
    \item \textbf{[计算]} 给定骨骼层级树与各关节点的局部旋转矩阵，利用正向运动学（FK）自顶向下递归计算指定肢体关节点的世界绝对空间坐标与朝向。
    \item \textbf{[计算]} 针对给定的肢体末端执行器目标空间位置与朝向，采用逆向运动学（IK）算法求解出各关节所需配置的角度姿态。
    \item \textbf{[对比]} 深入对比欧拉角、旋转矩阵、四元数与轴角在存储紧凑度、插值连续平滑性以及避免奇异性（万向节死锁）上的核心权衡。
    \item \textbf{[开发]} 编程实现三次样条插值或四元数球面线性插值算法，平滑生成连续的关键帧中间姿态与运动过渡。
    \item \textbf{[应用]} 熟练运用工业级动画建模软件（如 Blender 或 Maya），结合骨骼装配体系与隐式电位球或多边形网格，完成一个有机生命角色的基础动作装配与行走动画序列制作。
\end{enumerate}

\subsection{计算机物理仿真}
\textbf{【学时建议】}：6 领域核心学时。

计算机物理仿真在数理根基上与计算科学紧密相连。但在计算机图形学与交互技术语境下，仿真技术被赋予了截然不同的学科使命：它不再片面追求严苛的高精度物理预测模型，而是追求“物理真实感与艺术家创作意图之间的精妙平衡”。例如，科学计算关注“精确建模流体的微观表面张力偏微分方程”，而图形学则要求“受控地生成极具视觉美感与冲击力的皇冠状跌落水花”。

\subsubsection{领域核心主题}
\begin{itemize}
    \item \textbf{碰撞检测与动态碰撞响应 (Collision Detection \& Response)}：
    \begin{itemize}
        \item 基于有符号距离场（SDF）的快速内外穿透判断与梯度法向量反弹计算。
        \item 球-球相交、点-三角形空间包含判定、边-边最小距离分离轴测试。
        \item 连续碰撞检测（CCD）规避高速物体运动“穿模隧道效应”。
    \end{itemize}
    \item \textbf{基于随机噪声的过程式运动动画}：Perlin 噪声、Simplex 噪声在风场扰动与复杂海浪表面起伏中的程序化叠加。
    \item \textbf{粒子动力学系统}：
    \begin{itemize}
        \item 时间连续常微分方程的离散数值积分器：前向欧拉法（Forward Euler，条件不稳定）、中点法、蛙跳积分法、辛欧拉法与 Verlet / Velocity-Verlet 积分。
        \item 质点-弹簧阻尼网络（Mass-Spring Networks，用于布料织物与柔软弹性形变模拟）。
        \item 基于位置的动力学（Position-Based Dynamics, PBD，无条件数值稳定、快速约束投影求解）。
        \item 规则驱动的群体仿生学系统：Boids 经典鸟群三法则（分离、对齐、凝聚）与高密度人群疏散仿真。
        \item 刚体运动学与动力学（冲量积分法与角动量牛顿-欧拉方程）。
    \end{itemize}
    \item \textbf{基于欧拉网格的流体动力学}：
    \begin{itemize}
        \item 不可压缩流体纳维-斯托克斯方程。
        \item 半拉格朗日平流算法（Semi-Lagrangian Advection，无条件稳定的一阶回溯插值）。
        \item 亥姆霍兹-霍奇分解与泊松压力投影方程求解器（Pressure Projection，保持速度场无散度 $\nabla \cdot \mathbf{u} = 0$）。
    \end{itemize}
    \item \textbf{高度场与地形演化系统}：
    \begin{itemize}
        \item 地表沉积物输运、水力侵蚀与热风化侵蚀模型。
        \item 水波仿真：波动方程与浅水方程 (SWE)。
    \end{itemize}
    \item \textbf{基于离散规则的复杂系统}：L-系统（L-Systems，植物形态生长）、空间占优殖民系统（Space-colonizing Systems，叶脉与树枝向光性扩展）、康威生命游戏（Game of Life，元胞自动机）。
    \item \textbf{仿真必备高级几何数据结构}：
    \begin{itemize}
        \item 密集三维体素标量场/矢量场。
        \item 层次化稀疏体素场（Sparse Volumes，如 OpenVDB，有效节约烟雾火焰的巨大空白内存开销）。
        \item 自适应网格细分场（Adaptive Volumes，在激波与飞溅界面动态加密）。
        \item 携带多物理属性的粒子点集（Points with Attributes，质量、速度、密度、相态）。
        \item 非结构化三角形汤。
        \item 二维高程图。
    \end{itemize}
    \item \textbf{目标导向的仿真技术选型映射体系}：
    \begin{itemize}
        \item 粒子系统应用：粒子数值积分（前向欧拉、中点法、蛙跳法）。
        \item 刚体动力学：质点系统扩展、三角形-顶点/边-边相交碰撞检测。
        \item 织物布料模拟：基于质点-弹簧模型或 PBD 位置约束、自碰撞检测机制。
        \item 基于粒子的水体流体：光滑粒子流体动力学 (SPH)核函数插值、基于 SDF 的刚度障碍物碰撞反弹。
        \item 基于网格的烟雾与火焰：半拉格朗日平流、交错网格 MAC 离散化、泊松压力投影。
        \item 混合粒子-网格水体模拟：FLIP（Fluid-Implicit-Particle）与 PIC（Particle-In-Cell）混合算法。
    \end{itemize}
\end{itemize}

\subsubsection{示范性学习成果与能力目标}
\begin{enumerate}
    \item \textbf{[开发]} 基于经典牛顿力学定律，编程实现粒子动力学物理仿真引擎，支持辛欧拉法、中点法或 Verlet 积分器，完成诸如弹性蠕虫链、弹性摆球等物理现象的交互仿真。
    \item \textbf{[对比]} 深入对比不同流体仿真技术（如纯拉格朗日 SPH 粒子法、欧拉网格半拉格朗日法与混合 FLIP 算法）在表现弹道式飞溅水花、爆炸烟尘与高扩散火光时的物理直观性、数值耗散与算力开销。
    \item \textbf{[开发]} 独立设计并编程实现一个支持实时交互的二维/三维烟雾求解器，支持用户实时添加障碍物、扰动速度场并注入热源与烟尘浓度。
\end{enumerate}

\subsection{沉浸式技术}
\textbf{【学时建议】}：6 领域核心学时。

沉浸式技术重新定义了人与虚拟数字空间的交互边界。本单元涵盖虚拟现实（VR）、增强现实（AR）与混合现实（MR）所依托的显示光学、空间立体追踪、超低时延渲染管线、人因工程及空间隐私伦理。

\subsubsection{领域核心主题}
\begin{itemize}
    \item \textbf{扩展现实概念谱系辨析}：虚拟现实（VR）、增强现实（AR）、混合现实（MR）与扩展现实（XR）的技术边界；Milgram 提出的“现实-虚拟连续统”理论。
    \item \textbf{双目立体视觉显示技术}：左右眼视差生成、微显示器光机、菲涅尔透镜与薄煎饼透镜光学畸变预校正、色差补偿、眼动辐辏调节冲突 (VAC)。
    \item \textbf{观察者空间位姿追踪技术}：
    \begin{itemize}
        \item 追踪架构：由外向内追踪（Outside-In，利用外部红外基站或光捕相机）与由内向外追踪（Inside-Out，利用头显机载摄像机与 VIO 视觉惯导 SLAM 算法）。
        \item 空间 6 自由度（6-DoF）追踪维度：头部位置与朝向、双手控制器与裸手手势追踪、全身骨骼姿态捕捉与视线眼动追踪。
    \end{itemize}
    \item \textbf{沉浸式超低延迟渲染与可见性优化}：
    \begin{itemize}
        \item 动光延迟（Motion-to-Photon Latency，从用户头部物理转动到显示屏光子完全刷新出对应画面的物理延迟，须严格压制在 20ms 以内以防止生理眩晕）。
        \item 异步时间扭曲（ATW）、异步空间扭曲（ASW）与重投影技术。
        \item 注视点渲染（Foveated Rendering，利用眼动追踪在人眼中心凹视野区域提供全高清画质，周边外围视野大幅降采样）。
        \item 基于图像的全景 VR 系统与多细节层次（LOD）动态调度。
    \end{itemize}
    \item \textbf{网络化分布式虚拟现实}：多用户跨局域网/广域网的虚拟协同空间、空间状态同步、网络抖动补偿与空间音频定位渲染。
    \item \textbf{虚拟现实中的交互式三维空间建模}。
    \item \textbf{沉浸式技术行业深度赋能}：精准数字化外科手术导航、复杂战术对抗仿真推演、高端装备制造维修全息装配指引、沉浸式科学数据多维漫游。
    \item \textbf{沉浸式系统的人因安全与物理风险}：
    \begin{itemize}
        \item 晕动症/眩晕病的神经病理机制（前庭感官感知运动与视网膜视觉静止之间的感觉冲突）。
        \item 物理环境遮蔽安全隐患：沉浸式 VR 完全阻断现实视觉，导致室内跌倒摔伤、撞击实体家具与墙壁风险。
    \end{itemize}
    \item \textbf{沉浸式技术的无障碍普惠设计}：针对严重肢体运动障碍用户的注视交互适配、针对卧床不起病患的空间视锥倾斜补偿。
    \item \textbf{空间计算隐私与伦理风险 (交叉参考 SEP)}：
    \begin{itemize}
        \item 个人物理私密空间数字重构的泄露风险（头显机载相机对私人卧室、办公室空间的连续深度点云扫描与环境建模数据上传云端）。
        \item 生物特征数据滥用（步态、瞳孔缩放、身体骨骼尺寸比例）。
        \item 视觉强沉浸感导致的心理成瘾、现实脱节以及沉浸过程中用户对周边物理攻击的防御丧失。
    \end{itemize}
\end{itemize}

\subsubsection{示范性学习成果与能力目标}
\begin{enumerate}
    \item \textbf{[开发]} 编程渲染生成一对符合人类平均瞳距（IPD）几何视差的左/右眼立体匹配图像对。
    \item \textbf{[对比]} 详细阐明并对比由外向内与由内向外空间定位追踪技术在部署成本、追踪范围、遮挡敏感性与计算负荷上的核心优缺点。
    \item \textbf{[对比]} 深入对比基于纯几何建模的虚拟现实系统与基于全景立体图像重构的虚拟现实系统在自由漫游度与计算资源需求上的本质差异。
    \item \textbf{[评估]} 在高并发网络化多人协作沉浸式系统中，评估并设计保证动作同步实时性与多端数据状态强一致性的网络传输方案。
    \item \textbf{[设计]} 面向医院洁净手术室的心脏介入手术场景，撰写一套详尽的增强现实（AR）辅助导航应用程序功能需求与软硬件技术规范。
    \item \textbf{[评估]} 依据人机交互（HCI）与无障碍通用设计准则，对一款现存的沉浸式 XR 应用进行全维度的无障碍包容性评估并出具改进报告。
    \item \textbf{[识别]} 明确指出引发虚拟现实晕动病的核心关键软硬件技术指标（如动光延迟 MTP、双目视差校准度、画面帧率抖动），并论述抑制眩晕的系统工程优化手段。
\end{enumerate}

\subsection{交互技术}
\textbf{【学时建议】}：6 领域核心学时。

现代交互式计算机图形学构成了从桌面排版工具到大型多人在线游戏、元宇宙沉浸空间等所有高实时性交互软件的运转灵魂。本单元重点培养学生跨越底层硬件中断、事件驱动消息机制直至多模态感官交互全链路的系统构建能力。

\subsubsection{领域核心主题}
\begin{itemize}
    \item \textbf{事件驱动编程范式}：
    \begin{itemize}
        \item 鼠标指针移动、滚轮与物理触控点击事件。
        \item 键盘矩阵按键按下与释放状态机。
        \item 语音识别输入流解析与自然语言交互触发。
        \item 物理环境传感器数据触发（光感、陀螺仪、接近传感器）。
        \item 跨组件与跨线程的消息传递通信。
        \item 异步网络套接字事件轮询（如 Select/Epoll 事件模型）。
        \item 底层硬件中断事件处理与驱动级回调机制。
    \end{itemize}
    \item \textbf{单通道经典图形用户界面 (WIMP 范式)}：视窗、图标、菜单、定点设备的状态变迁与事件分发路由机制。
    \item \textbf{手势交互接口}：多点触控手势识别、空中自由手势几何匹配（如捏合、滑动）、无障碍替代方案设计（如为手部缺失或运动神经损伤者提供软件逻辑层面的“虚拟点赞”事件注入）。
    \item \textbf{触觉与力反馈接口}：
    \begin{itemize}
        \item 外部线性振动执行器（LRA 与 ERM 电机）。
        \item 具备压力与纹理触觉感知的力反馈数据手套。
        \item 限制关节运动的大负载力反馈机械外骨骼。
    \end{itemize}
    \item \textbf{多模态融合交互系统}：视线追踪、语音指令、手势捕捉与微表情识别的时空融合与多义性消歧。
    \item \textbf{基于人工智能的前沿沉浸式交互接口}：
    \begin{itemize}
        \item 脑机接口（BCI）：基于非侵入式脑电图（EEG）电极阵列的稳态视觉诱发电位（SSVEP）与运动想象意图解码。
        \item 集成眼动追踪的混合现实智能头显：凝视选择与眼部眨眼触发。
        \item 智能全息光波导 AR 眼镜的极简微交互。
    \end{itemize}
    \item \textbf{无障碍设计准则与通用包容性设计 (交叉参考 SEP)}：为肢体、听觉、认知等各类障碍群体提供无障碍等价交互通道。
\end{itemize}

\subsubsection{示范性学习成果与能力目标}
\begin{enumerate}
    \item \textbf{[开发]} 编写一个完整的二维图形小游戏，能够对单通道鼠标移动、点击以及键盘状态事件做出即时、平滑的逻辑响应。
    \item \textbf{[开发]} 编写嵌入式控制程序，使单片机电路能够根据外接可变电位器（旋转电阻）的模拟阻值变动，平滑调节虚拟图形对象的多维属性。
    \item \textbf{[开发]} 面向移动智能终端，开发一个能精准响应单指滑动、双指缩放及多点复杂手势的交互式移动应用程序。
    \item \textbf{[应用]} 利用计算机视觉库或手势追踪外设，编写程序实现通过手部空中手势对三维场景中虚拟物体的空间抓取与平移旋转。
    \item \textbf{[设计]} 设计并编程实现一个集成立体触觉反馈（如变频振动或力阻尼）的交互应用程序，向用户传递虚拟物体的表面粗糙度与刚度质感。
    \item \textbf{[开发]} 架构并实现一个复杂的响应式图形交互系统，使其能够稳定并发处理来自网络套接字、定时器硬件中断以及异步用户外设的多源事件触发器。
\end{enumerate}

\subsection{数字图像处理}
\textbf{【学时建议】}：6 领域核心学时。

数字图像处理涵盖面向多种工程目标的图像分析、滤波、增强、编码与几何变换，其核心主旨在于提升视觉感知品质与实现图像内容的智能操纵。该单元是计算机视觉（CV）不可动摇的技术基石。

\subsubsection{领域核心主题}
\begin{itemize}
    \item \textbf{二值数学形态学处理}：
    \begin{itemize}
        \item 像素连通分量标记（4-连通与 8-连通域分析）。
        \item 经典形态学算子：膨胀、腐蚀、开运算与闭运算。
        \item 连通区域几何属性统计提取：面积、周长、质心坐标、方向角、主轴长度与圆度。
    \end{itemize}
    \item \textbf{色彩直方图分析与对比度操纵}：直方图统计表征模型、直方图均衡化、局部自适应直方图均衡化（CLAHE）与动态范围归一化拉伸。
    \item \textbf{图像空域滤波与图像增强}：
    \begin{itemize}
        \item 二维离散空间卷积运算原理。
        \item 图像平滑与模糊算子：均值滤波、高斯低通滤波、双边滤波（Bilateral Filter，保边去噪）。
        \item 图像锐化算子：拉普拉斯算子、Sobel 梯度算子、反锐化掩模技术。
        \item 频域滤波技术：二维快速傅里叶变换（FFT）、理想低通/高通滤波器、巴特沃斯滤波器与高斯频域滤波。
    \end{itemize}
    \item \textbf{图像退化物理建模与图像复原}：加性高斯白噪声、椒盐噪声、运动模糊退化函数模型；逆滤波、维纳滤波；现代图像修复技术（Inpainting，基于偏微分方程与泊松图像编辑 Poisson Image Editing）。
    \item \textbf{数字图像编码与信息压缩理论}：
    \begin{itemize}
        \item 空间冗余、时间冗余与人眼视觉心理冗余。
        \item 无损熵编码：香农-哈夫曼编码（Huffman Coding）、算术编码与行程编码（RLE）。
        \item 正交变换编码：离散余弦变换（DCT，JPEG 核心）、离散小波变换（DWT，JPEG 2000 核心）。
        \item 奈奎斯特-香农采样定理与频谱混叠规避。
        \item 数字盲水印嵌入与数字版权管理（DRM）、抗几何攻击鲁棒性与隐蔽性伦理考量。
    \end{itemize}
    \item \textbf{前沿深度学习视觉技术}：卷积神经网络（CNN，特征提取层、池化层与残差连接）、视觉 Transformer 架构（ViT，自注意力机制在图像分块表征中的应用）、生成对抗网络（GAN）与前沿扩散去噪概率模型。
    \item \textbf{伦理、司法防伪与社会公平挑战 (SEP)}：换脸与深度伪造（Deepfakes）的自动取证与鉴伪技术；计算机视觉算法因训练样本失衡而在特定肤色、发型或少数族裔群体上出现的严重误识别、司法误判与社会歧视隐患。
\end{itemize}

\subsubsection{示范性学习成果与能力目标}
\begin{enumerate}
    \item \textbf{[应用]} 运用形态学膨胀与腐蚀算子，对带有毛刺噪声的二值图像进行平滑与孔洞填补，并准确分离粘连的微小图元。
    \item \textbf{[开发]} 编程提取图像的 HSV 色彩空间矩阵，在独立保持亮度与饱和度不变的前提下，操纵调整图像的色相通道。
    \item \textbf{[开发]} 编写程序实现高通滤波算子（在空间域利用卷积核，或在频域利用傅里叶高通掩模），实现对图像微弱边缘轮廓的高对比度增强。
    \item \textbf{[开发]} 采用高级图像修复补全算法（如泊松图像无缝融合 Poisson Editing 或基于样本块的修补技术），对被严重损毁或人为涂抹的图片进行无痕修复还原。
    \item \textbf{[开发]} 编程对二维图像进行傅里叶正变换，通过在频域平面内有选择性地滤除特定周期性高频条纹噪声，实现图像质量逆向增强。
    \item \textbf{[评估]} 深入阐明人脸识别算法在特定应用场景下的技术缺陷与伦理陷阱；针对“人脸识别能否以及如何做到合乎工程伦理的应用”撰写深度学术评述报告。
\end{enumerate}

\subsection{实体与物理计算}
\textbf{【学时建议】}：6 领域核心学时。

本单元与专用平台开发（SPD）以及人机交互（HCI）高度交织，旨在打破传统计算机“屏幕-键鼠”的物理禁锢，将计算能力拓展至真实物理世界的实体对象、智能空间、传感器网络与数字制造链条中。

\subsubsection{领域核心主题}
\begin{itemize}
    \item \textbf{与真实物理世界的通信与交互感知}：
    \begin{itemize}
        \item 物理环境数据感知采集：模拟信号与模拟-数字转换（ADC）、温度/光敏电阻/超声波测距/红外接近/加速度与陀螺仪传感器数据流处理。
        \item 物理世界驱动控制：脉冲宽度调制（PWM）技术驱动舵机、直流电机驱动器、高精度双极性步进电机、电磁阀与继电器。
    \end{itemize}
    \item \textbf{微控制器的事件驱动固件架构}：定时器硬件中断、外部引脚电平跳变中断服务函数（ISR）、环形串口缓冲队列与非阻塞式主循环调度。
    \item \textbf{与物理实体的制造连接}：
    \begin{itemize}
        \item 计算机辅助设计（CAD）：实体造型、参数化约束、装配体干涉检查。
        \item 计算机辅助制造（CAM）：刀具路径规划、G 代码与 M 代码生成。
        \item 现代数字制造工艺路线：
        \begin{itemize}
            \item 人机交互快速原型设计（HCI Prototyping，低成本纸板与激光亚克力快速成型）。
            \item 增材制造技术（3D Printing）：熔融沉积成型（FDM）、立体光固化成型（SLA）、选择性激光烧结（SLS）工艺特性对比。
            \item 减材制造技术：数控机床铣削、高精度激光切割与水刀切割。
            \item 成型制造工艺：真空热吸塑成型、金属板料冲压折弯。
        \end{itemize}
    \end{itemize}
    \item \textbf{物联网感知与互联体系}：
    \begin{itemize}
        \item 网络互联拓扑结构与轻量级协议栈（MQTT、CoAP）。
        \item 短程与低功耗无线通信标准：Wi-Fi (802.11)、低功耗蓝牙 (BLE)、Zigbee、LoRa 远程扩频通信与近场通信 (RFID/NFC)。
    \end{itemize}
    \item \textbf{实体计算安全与隐私伦理 (SEP)}：物联网终端物理逆向工程漏洞、家用常开麦克风与环境传感器的隐私嗅探防范、固件空中升级（OTA）防篡改机制。
\end{itemize}

\subsubsection{示范性学习成果与能力目标}
\begin{enumerate}
    \item \textbf{[开发]} 构建一个基础硬件控制电路（包含微控制器、限流电阻与物理按键开关），编写固件实现通过按键消抖触发点亮或关闭高亮发光二极管（LED）。
    \item \textbf{[开发]} 搭建并联调一个闭环实体机电控制系统，使舵机能够根据超声波或红外距离传感器的实时测距数值，做出精准角度旋转跟踪。
    \item \textbf{[开发]} 利用光敏传感器采集环境照度，通过编程动态驱动外部物理执行器（如平滑调节 LED 阵列的发光亮度或驱动显示屏幕中图形的虚拟色彩）。
    \item \textbf{[开发]} 在专业计算机辅助设计（CAD）工业软件中独立构建一个复杂的三维实体零件，达成以下指标：
    \begin{itemize}
        \item 在 CAD 系统中演示并论证空间仿射变换（平移、旋转、镜像）的底层数学实现。
        \item 演示对象的实例化复用机制。
        \item 制定完整的零件数字化制造工序计划，详细评估材料损耗与加工工时成本。
        \item 实际操控 3D 打印机或 CNC 机床完成实体加工，对比实测物理零件与理论设计图纸的公差差异，剖析加工误差来源。
    \end{itemize}
    \item \textbf{[开发]} 手工编写标准数控 G 代码与 M 代码，用于控制数控铣床自动雕刻一个具有特定拓扑连通性的三维迷宫零件，并利用 CAM 仿真软件检验刀路碰撞。
    \item \textbf{[评估]} 面向适老化智能医疗药盒的研发任务，系统论证应选用以太网、Wi-Fi、BLE、RFID/NFC 或蜂窝网络中的哪种方案实现云端互联，并独立制作可运行的硬件原型系统。
    \item \textbf{[对比]} 深入对比增材制造、减材制造与吸塑成型工艺在材料强度、几何自由度、单件打样成本与规模化量产时的经济技术适用边界。
\end{enumerate}

\section{专业素养与思维倾向}

在图形学与交互技术的系统研习中，必须高度注重塑造学生以下几维关键的工程师素养与思维倾向：
\begin{itemize}
    \item \textbf{积极主动与自发探索}：面对未知的底层图形 API、晦涩的物理着色方程与复杂的物理外设驱动，具有自主查阅前沿文献与工业规范的主动探索精神。
    \item \textbf{自我驱动与终身学习}：图形硬件与实时渲染技术日新月异，从业者需保持高度的自驱力，能够紧跟全球学术界（如 SIGGRAPH）与工业界引擎前沿演进。
    \item \textbf{目标导向与结果思维}：在算力、显存带宽与物理帧率强约束下，始终以“满足真实感与帧率要求”为最终导向开展架构设计。
    \item \textbf{职业操守与伦理自律}：在处理涉及人脸生物特征、空间点云扫描、数字版权与生成式深度伪造内容时，恪守工程道德，尊重知识产权，严格保护用户个人隐私。
    \item \textbf{敏捷适应与灵活变通}：面对跨平台图形后端（DirectX、Vulkan、Metal、WebGPU）以及异构计算环境，具备极强的代码架构迁移与兼容变通能力。
    \item \textbf{团队协作与跨界融合}：图形工程通常处于软件工程师、数学算法专家、三维数字艺术家、工业设计师以及产品经理的交叉枢纽，必须具备强大的跨学科团队协作素养。
    \item \textbf{敏锐响应用户反馈}：高度关注真实用户的感官体验、交互延迟以及生理不适，能敏锐依据用户体验数据与可用性测试迭代优化软件。
    \item \textbf{精益求精与严谨细致}：深入推敲每一个像素着色计算精度、浮点数舍入误差、网格法线奇异点以及内存缓冲区边界溢出，防范图形破面与渲染穿模。
    \item \textbf{突破定势与开拓创新}：面对极其耗时的物理级光线模拟计算，敢于跳出传统蛮力计算框架，创造性地运用微表面理论、降噪神经网络或启发式空间剖分技术进行性能突围。
    \item \textbf{勇于担当与后果意识}：深刻意识到图形交互算法对物理现实安全的直接影响（如自动驾驶机器视觉盲区、医疗 AR 手术导航定位偏差），对软件系统的一切潜在后果负有敬畏与担当之心。
    \item \textbf{规范的代码审查能力}：善于开展深入的同行代码审查，既能发现他人代码中的显存泄露与性能瓶颈，亦能以开放的心态吸纳同行改进建议。
    \item \textbf{多维立体的高效沟通素养}：
    \begin{itemize}
        \item \textbf{口头表达}：能面向非技术背景的艺术设计师清晰阐释复杂着色器技术参数的物理内涵。
        \item \textbf{书面表达}：能撰写结构严密、数学定义准确的技术白皮书、系统架构设计说明书与学术研究报告。
        \item \textbf{代码表达}：能编写结构优雅、自解释性强、注释规范且符合工业级编码风格规范的高质量代码。
    \end{itemize}
\end{itemize}

\section{数学与基础数据结构要求}

图形学与交互技术建立在扎实完备的现代应用数学基石之上：

\subsection{必须具备的数学核心}
\begin{itemize}
    \item \textbf{线性代数}：
    \begin{itemize}
        \item 空间点、笛卡尔坐标系、仿射空间与四维齐次坐标系统。
        \item 向量代数：向量加减、标量乘法、点积（Dot Product，用于夹角与漫反射余弦计算）、叉积（Cross Product，用于表面法向量求解与三角形朝向判定）。
        \item 矩阵代数：矩阵乘法、逆矩阵、转置矩阵以及伴随矩阵（用于法向量变换逆转置矩阵推导）。
        \item 仿射变换矩阵构造：平移、旋转、非均匀缩放、剪切与透视投影矩阵全流程推导。
    \end{itemize}
    \item \textbf{微积分}：标量函数与多元空间矢量函数的连续性概念、偏导数、空间梯度、定向导数与多重微积分基本概念。
\end{itemize}

\subsection{建议掌握的进阶数学}
\begin{itemize}
    \item \textbf{高阶线性代数}：
    \begin{itemize}
        \item 矩阵特征值与特征向量分解（Eigen Decomposition，用于主成分分析 PCA 与网格主曲率计算）。
        \item 高斯消元法与矩阵三角分解（LU 分解，用于刚体与流体线性泊松方程组快速求解）。
        \item 奇异值分解（SVD，用于三维刚性点云精确配准与最佳位姿拟合）。
    \end{itemize}
    \item \textbf{进阶代数与数值微积分}：
    \begin{itemize}
        \item 四元数代数体系（Quaternions，哈密顿算子、四元数插值 Slerp 与万向节死锁彻底消除）。
        \item 常微分方程数值解法与蒙特卡洛积分方法。
    \end{itemize}
\end{itemize}

\subsection{核心数据结构要求}
\begin{itemize}
    \item \textbf{必须熟练掌握的底层数据结构}：
    \begin{itemize}
        \item 有向无环图（DAG）：用于层次化场景图、着色器材质节点图。
        \item 定长元组：空间几何点、三维物理向量、固定阶数齐次变换矩阵。
        \item 密集多维张量数组（Dense 1D, 2D, 3D Arrays）：帧缓冲区、深度缓冲区、二维纹理采样图元、三维体素标量场。
    \end{itemize}
    \item \textbf{极具工程价值的进阶数据结构}：
    \begin{itemize}
        \item 结构体数组与数组结构体的底层访存布局对偶性：理解数据导向型设计（DOD）在 CPU 缓存局部性与 GPU SIMD 向量化并行中的极致性能影响。
        \item 高阶空间检索树结构：$K$-d 树、四叉树、八叉树、包围盒层次结构（BVH）以及用于图像压缩的哈夫曼树。
    \end{itemize}
\end{itemize}

\section{交叉主题与跨学科共享概念}

计算机图形学与交互技术并非孤立存在，它与计算机科学各分支学科紧密交融：
\begin{itemize}
    \item \textbf{沉浸式技术 (GIT) 与人机交互 (HCI)}：空间三维用户体验、凝视交互、人机认知负荷评估。
    \item \textbf{交互技术 (GIT) 与 人机交互/GUI 编程以及计算科学 (CN) 交互可视化}：复杂事件驱动模型、状态图迁移管理。
    \item \textbf{过程式生长建模 (GIT) 与编程语言理论 (FPL)}：L-系统文法与巴科斯-诺尔范式（BNF）形式文法的对偶性——利用文法规则派生三维几何分形图像。
    \item \textbf{物理仿真 (GIT) 与计算科学 (CN)}：连续动力学方程离散化与数值阻尼控制。
    \item \textbf{数据可视化 (GIT) 与人工智能 (AI) 及专用平台开发 (SPD)}：高维机器学习特征流形嵌入可视化分析。
    \item \textbf{数字图像处理 (GIT) 与 人工智能/机器视觉}：卷积特征提取、图像分割、双目立体匹配。
    \item \textbf{实体计算与计算机动画 (GIT) 与机器人操作系统及嵌入式平台 (专用平台/AI)}：正逆运动学控制、机械臂末端轨迹规划。
    \item \textbf{图形核心渲染管线 (GIT) 与并行分布式计算 (PDC)}：超大规模 GPU 细粒度数据级并行流式计算。
\end{itemize}

贯穿全书的核心横切主题包括：\textbf{效能}、\textbf{伦理}、\textbf{建模}、\textbf{编程}、\textbf{原型设计}、\textbf{可用性} 以及 \textbf{综合评估}。

\section{真实能力任务规范}

\begin{taskbox}[任务 1：多场景数字图像格式甄别与工程选型]
\textbf{【任务情境】}：某工业智能装备企业正在开发一套集产品外观全彩高精摄影、二维矢量卡通吉祥物绘制以及数控超高压高精度水刀金属板料切割轨迹规划于一体的多模态系统。\\
\textbf{【能力陈述】}：针对 1) 真实场景彩色摄影照片、2) 扁平化矢量卡通图形、3) 工业数控超高压水刀切割运动路径这三类性质截然不同的媒体数据，准确选定最优的数字图像存储文件格式，并从压缩机制、几何可缩放性与时空开销角度给出严密论证。\\
\textbf{【能力领域】}：应用与理论 \quad \textbf{【能力单元】}：系统评估\\
\textbf{【核心要求】}：图形交互/基础核心 \quad \textbf{【达标等级】}：\textbf{理解与解释}
\end{taskbox}

\begin{taskbox}[任务 2：现代图形渲染管线技术白皮书撰写]
\textbf{【任务情境】}：某新兴游戏引擎研发团队需要为新入职工程师编写一份关于现代 GPU 渲染底层架构的官方技术白皮书。\\
\textbf{【能力陈述】}：以权威、严密的工程语言，系统详尽地阐述从应用程序调用图形 API（如 OpenGL / Vulkan / DirectX）发出绘制调用开始，直至图元最终在物理帧缓冲区完成光栅化与色彩混合输出的完整渲染管线各物理与逻辑阶段。\\
\textbf{【能力领域】}：应用实践 \quad \textbf{【能力单元】}：技术沟通与文档规范\\
\textbf{【核心要求】}：图形交互/基本概念、基础渲染技术 \quad \textbf{【达标等级】}：\textbf{理解与解释}
\end{taskbox}

\begin{taskbox}[任务 3：遵循物理与艺术法则的真实弹跳球动画制作]
\textbf{【任务情境】}：动画工作室正在制作一部高品质三维动画短片，需要构建一个从空中自由跌落并与刚性地面多次反弹碰撞的充气橡胶皮球动态资产。\\
\textbf{【能力陈述】}：综合运用重力加速度、弹性恢复系数、以及经典动画十二原则中的“挤压与拉伸”、“预备动作”、“慢入与慢出（加减速缓冲）”与“时间节拍”，设计并编程或装配渲染出一个物理可信且极具艺术张力的逼真弹跳球动画序列。\\
\textbf{【能力领域】}：工程应用 \quad \textbf{【能力单元】}：需求分析/设计/开发/测试全流程\\
\textbf{【核心要求】}：图形交互/基本概念、基础渲染、计算机动画、高级着色 \quad \textbf{【达标等级】}：\textbf{应用、评估与开发}
\end{taskbox}

\begin{taskbox}[任务 4：基于反应-扩散模型的有机生物条纹生成系统]
\textbf{【任务情境】}：某数字艺术生成实验室需要程序化模拟热带斑马鱼或猎豹皮肤表面的自然有机条纹图案，并要求支持艺术家现场实时操纵化学扩散参数。\\
\textbf{【能力陈述】}：使用自主选定的现代编程语言，编程实现经典的图灵反应-扩散方程（Turing Reaction-Diffusion Model，如 Gray-Scott 偏微分系统），并在屏幕上实时渲染生成的有机条纹；设计并实现一个响应式图形用户界面，使系统的化学物质扩散率、反应速率等参数能够在毫秒级实时联动调节。\\
\textbf{【能力领域】}：应用与开发 \quad \textbf{【能力单元】}：需求分析/架构设计/开发/测试/部署\\
\textbf{【核心要求】}：图形交互/基础渲染、高级着色、物理仿真、交互技术；人机交互/系统设计、用户理解；软件开发基础/开发方法 \quad \textbf{【达标等级】}：\textbf{应用、评估与开发}
\end{taskbox}

\begin{taskbox}[任务 5：高交互可定制的粒子系统数字烟花仿真引擎]
\textbf{【任务情境】}：某虚拟城市文旅庆典平台需要上线一个供全网游客在线自主定制并燃放的 Web 虚拟烟花庆典互动模块。\\
\textbf{【能力陈述】}：编程实现一个包含重力加速度、空气阻力衰减、爆炸冲量分形扩散以及色彩渐变消隐的物理粒子系统烟花发生器；构建直观的人机交互界面，使用户能够实时定制粒子发射速率、色彩调色板、爆炸碎片半径与重力加速度等多维参数。\\
\textbf{【能力领域】}：工程应用 \quad \textbf{【能力单元】}：全周期软件开发与交付\\
\textbf{【核心要求】}：图形交互/基础渲染、物理仿真、交互技术；人机交互/系统设计；软件开发基础/开发方法 \quad \textbf{【达标等级】}：\textbf{应用、评估与开发}
\end{taskbox}

\begin{taskbox}[任务 6：基于高斯卷积核的人类屈光不正视力障碍仿真滤镜]
\textbf{【任务情境】}：眼科医学研究中心与无障碍设计团队需要一套能够向视力健全者直观仿真不同程度近视患者真实视错觉感知的视觉渲染系统。\\
\textbf{【能力陈述】}：设计并实现基于不同核尺寸与标准差 $\sigma$ 的二维离散高斯模糊卷积滤波算法，建立屈光度（近视度数）与卷积核弥散斑之间的精确数学对应关系，对输入的高清物理场景图像进行实时退化模糊仿真。\\
\textbf{【能力领域】}：工程应用 \quad \textbf{【能力单元】}：需求分析/算法设计/工程测试/用户验收\\
\textbf{【核心要求】}：图形交互/基本概念、基础渲染、数字图像处理；软件开发基础/开发方法 \quad \textbf{【达标等级】}：\textbf{应用、评估与开发}
\end{taskbox}

\begin{taskbox}[任务 7：低对比度退化图像的自适应色彩直方图归一化增强]
\textbf{【任务情境】}：恶劣低照度雾霾天气下安防监控摄像头采集的数字图像对比度极度匮乏，关键目标特征被严重掩蔽，亟待自动化提升图像视觉动态范围。\\
\textbf{【能力陈述】}：针对三通道 RGB 彩色图像，编程实现多通道色彩直方图统计、动态范围自适应线性拉伸以及直方图均衡化处理，在消除曝光不足的同时防止过饱和与高频色彩失真。\\
\textbf{【能力领域】}：工程应用 \quad \textbf{【能力单元】}：算法开发与性能验证\\
\textbf{【核心要求】}：图形交互/基本概念、基础渲染、数字图像处理 \quad \textbf{【达标等级】}：\textbf{应用、评估与开发}
\end{taskbox}

\begin{taskbox}[任务 8：跨移动双平台的多人井字棋交互应用工程实现]
\textbf{【任务情境】}：针对 6 岁及以上低龄儿童群体，开发一款界面友好、反应敏捷、兼顾平板与手机触控交互的跨平台益智井字棋游戏。\\
\textbf{【能力陈述】}：面向 iOS 与 Android 移动操作系统平台，完整设计并实现一套具备状态判定、胜负仲裁、悔棋及平局处理逻辑的井字棋移动应用；界面需针对低龄儿童认知特征进行图标与动效优化，支持低时延触控手势响应。\\
\textbf{【能力领域】}：移动软件工程 \quad \textbf{【能力单元】}：需求分析/交互设计/跨平台开发/测试交付\\
\textbf{【核心要求】}：图形交互/基本概念、基础渲染、交互技术；人机交互/系统设计、理解用户；专用平台/移动开发基础、移动平台架构；软件开发基础/开发方法 \quad \textbf{【达标等级】}：\textbf{应用、评估与开发}
\end{taskbox}

\begin{taskbox}[任务 9：区域历史气象气温时空动态演化可视化系统]
\textbf{【任务情境】}：国家气象科学数据中心需要将近半个世纪以来某特大城市群的地表温度时序栅格数据转换为直观生动的动态演化可视化成果，用于联合国气候大会公众科普。\\
\textbf{【能力陈述】}：导入多年连续气温时空观测数据集，设计并实现一套高质量的三维/二维温度时序动态演化动画系统；采用感知均匀的伪彩色色阶编码，叠加等温线动态拓扑形变，生动展示区域气温长期冷暖波动特征。\\
\textbf{【能力领域】}：应用与科学计算 \quad \textbf{【能力单元】}：全流程系统实现与科普可用性评估\\
\textbf{【核心要求】}：图形交互/基本概念、基础渲染、数据可视化；人机交互/系统设计、用户理解 \quad \textbf{【达标等级】}：\textbf{应用、评估与开发}
\end{taskbox}

\begin{taskbox}[任务 10：双区域 COVID-19 传染病时序扩散对比可视化]
\textbf{【任务情境】}：全球流行病预防控制机构急需对两个地理特征相似但防疫政策迥异的代表性行政区域的 COVID-19 历史感染扩散时序数据进行严密可视化对比。\\
\textbf{【能力陈述】}：基于流行病学每日新增感染与累计确诊数据集，设计并构建交互式双联对比动画系统；合理运用双 Y 轴折线、动态热力密度图与时间轴无缝联动刷选，使公共卫生决策者能一目了然洞察两地病毒传播峰值与平台期的动态差异。\\
\textbf{【能力领域】}：应用实践 \quad \textbf{【能力单元】}：需求分析/可视化设计/测试/决策评估\\
\textbf{【核心要求】}：图形交互/基本概念、基础渲染、数据可视化；人机交互/交互设计、受众认知 \quad \textbf{【达标等级】}：\textbf{应用、评估与开发}
\end{taskbox}

\begin{taskbox}[任务 11：弱光自适应智能光感微控制器控制终端系统]
\textbf{【任务情境】}：某智慧农业与仓储物流企业需要部署一批无需人工干预、可在暗光或夜间自动触发照明补偿的超低功耗智能终端节点。\\
\textbf{【能力陈述】}：使用工业标准微控制器（如 Arduino、ESP32 或 STM32）与外围电子元器件，独立搭建包含光敏电阻分压传感回路的硬件电路；编写嵌入式固件程序，实现利用 ADC 定期采样环境照度，并在光强跌破设定阈值时通过驱动外部三极管/MOSFET 稳定点亮高功率发光二极管（LED）。\\
\textbf{【能力领域】}：嵌入式与实体计算 \quad \textbf{【能力单元】}：硬件设计/电路焊接/固件编写/联调部署\\
\textbf{【核心要求】}：图形交互/基本概念、实体与物理计算；专用平台/嵌入式平台；人机交互/硬件设计；社会伦理/物联网隐私 \quad \textbf{【达标等级】}：\textbf{解释、应用、评估与开发}
\end{taskbox}

\section{课程设计与开课规划建议}

联合工作组针对不同专业定位的高校，推荐五套典型的专业课程组织与开课规划方案：

\subsection{方案 1：交互式计算机图形学}
\textbf{【总学时】}：44 学时。
\begin{itemize}
    \item \textbf{涵盖模块}：图形交互/基础渲染技术（40 学时）+ 社会伦理/工程伦理（4 学时）。
    \item \textbf{先修课程要求}：计算机编程基础 II (CS2)、线性代数（熟练掌握仿射变换）、三角学。
    \item \textbf{课程能力目标}：修完本课程的学生应深刻理解并能熟练调用现代图形 API 独立构建计算机图形渲染系统；精准掌握几何三维模型坐标变换、虚拟观察摄影机平截头体定位以及点光源/平行光源的光照计算。
\end{itemize}

\subsection{方案 2：实体与物理计算}
\textbf{【总学时】}：44 学时。
\begin{itemize}
    \item \textbf{涵盖模块}：图形交互/实体与物理计算（30 学时）+ 专用平台/嵌入式平台开发（10 学时）+ 社会伦理/隐私保护与无障碍包容性（4 学时）。
    \item \textbf{先修课程要求}：计算机编程基础 I (CS1)。
    \item \textbf{课程能力目标}：修完本课程的学生应能独立设计并装配基础电子电路，熟练编写微控制器固件；透彻掌握电路极性、欧姆定律、信号模数转换及电子仪器安全操作规程。
\end{itemize}

\subsection{方案 3：数字图像处理}
\textbf{【总学时】}：44 学时。
\begin{itemize}
    \item \textbf{涵盖模块}：图形交互/数字图像处理（30 学时）+ 图形交互/基础渲染技术（10 学时）+ 社会伦理/多元包容、数据隐私与知识产权法（4 学时）。
    \item \textbf{先修课程要求}：计算机编程基础 II (CS2)、线性代数。
    \item \textbf{课程能力目标}：修完本课程的学生应熟练掌握数字图像的数字化采集、空域与频域滤波、色彩增强、数学形态学变换、特征提取、有损/无损压缩编码与磁盘持久化全流程。
\end{itemize}

\subsection{方案 4：数据科学与信息可视化}
\textbf{【总学时】}：47 学时。
\begin{itemize}
    \item \textbf{涵盖模块}：图形交互/数据可视化（30 学时）+ 图形交互/基础渲染技术（10 学时）+ 人机交互/理解用户认知（3 学时）+ 社会伦理/数据伦理、算法偏见与隐私规约（4 学时）。
    \item \textbf{先修课程要求}：计算机编程基础 II (CS2)、线性代数。
    \item \textbf{课程能力目标}：修完本课程的学生应能独立清洗并甄别真实科学或社会学数据集，设计并实现一套符合认知心理学规律、严谨防范误导性修辞、兼具高可用性的交互式可视化软件系统。
\end{itemize}

\subsection{方案 5：计算机物理仿真系统}
\textbf{【总学时】}：44 学时。
\begin{itemize}
    \item \textbf{涵盖模块}：图形交互/计算机物理仿真（30 学时）+ 图形交互/基础渲染技术（10 学时）+ 社会伦理/工程伦理（4 学时）。
    \item \textbf{先修课程要求}：计算机编程基础 II (CS2)、线性代数、基础微积分。
    \item \textbf{课程能力目标}：修完本课程的学生应掌握时间连续物理方程向离散数值积分转化的方法论，能够独立编程构建支持用户实时交互、兼顾物理真实感与艺术可控性的物理与非物理动力学仿真引擎。
\end{itemize}

\vspace{10pt}
\noindent\textbf{【课程规范编写委员会】}
\begin{itemize}
    \item \textbf{主席}：Susan Reiser（北卡罗来纳大学阿什维尔分校，UNC Asheville，美国）
    \item \textbf{核心委员}：
    Erik Brunvand（犹他大学，University of Utah，美国）、
    Kel Elkins（NASA 戈达德太空飞行中心科学可视化实验室，NASA/GSFC，美国）、
    Jeff Lait（SideFX 软件公司，Toronto，加拿大）、
    Amruth Kumar（拉马波学院，Ramapo College，美国）、
    Paul Mihail（瓦尔多斯塔州立大学，Valdosta State University，美国）、
    Tabitha Peck（戴维森学院，Davidson College，美国）、
    Ken Schmidt（美国国家海洋和大气管理局 NCEI，NOAA NCEI，美国）、
    Dave Shreiner（Unity Technologies，旧金山，美国）。
    \item \textbf{主要贡献专家}：
    Greg Shirah（NASA/GSFC 科学可视化实验室）、
    AJ Christensen（NASA/GSFC 科学可视化实验室 / SSAI）、
    Barbara Mones（华盛顿大学，University of Washington）、
    Beatriz Sousa Santos（阿威罗大学，University of Aveiro，葡萄牙）、
    Ted Kim（耶鲁大学，Yale University）、
    Ginger Alford（南卫理公会大学，Southern Methodist University）。
\end{itemize}

\clearpage
\onecolumn
\section{附录：核心主题与技能要求等级对照表}

\noindent\begin{tabularx}{\textwidth}{p{0.15\textwidth} Y p{0.12\textwidth} p{0.10\textwidth} p{0.08\textwidth}}
\toprule
\textbf{知识单元} & \textbf{核心知识点与具体主题} & \textbf{认知等级} & \textbf{分类} & \textbf{学时} \\
\midrule
\multirow{4}{*}{\makecell[l]{图形交互/基本概念}} 
& 娱乐、工商业与科学应用领域的现代图形学落地生态 & 解释 & 计算机核心 & 1 \\
& 人类视觉生理系统、三刺激接收、视差与视觉无障碍 & 解释 & 计算机核心 & 1 \\
& 模拟数据数字化、采样量化、光栅图元 vs 矢量数学模型 & 解释 & 计算机核心 & 1 \\
& 标准媒体格式、加色/减色色彩空间、动画离散时序与伦理 & 解释 & 计算机核心 & 1 \\
\midrule
\multirow{3}{*}{\makecell[l]{图形交互/基础渲染}} 
& 现代渲染管线、仿射空间变换、物理光辐射与数值积分、纹理映射、MIP-mapping、走样与深度消隐 & 解释 & 领域核心 & 10 \\
& 正向与反向渲染范式、多边形网格拓扑、辐射度量、光线追踪、渲染方程、三角形光栅化、空间树结构 & 解释 & 领域核心 & 5 \\
& 基于现代可编程着色器标准 API（如 OpenGL）生成交互图像 & 实现 & 领域核心 & 3 \\
\midrule
\multirow{4}{*}{\makecell[l]{图形交互/数据可视化}} 
& 二维/三维标量场色彩映射、等值面与时变动态数据可视化 & 解释/实现 & 领域核心 & 3 \\
& 多元数据高维降维技术、平行坐标轴、流场与关系拓扑图 & 解释/实现 & 领域核心 & 2 \\
& 驱动视觉符号抽象的认知心理学底座与色彩学原理 & 解释 & 领域核心 & 1 \\
& 可视化设计认知偏见与伦理责任（防范视觉欺骗修辞） & 评估 & 领域核心 & 1 \\
\midrule
\multirow{3}{*}{\makecell[l]{图形交互/几何建模}} 
& 表面模型、细分多边形化、Delaunay 剖分、Marching Cubes 算法、参数多项式曲线曲面（Bézier/NURBS） & 解释/使用 & 领域核心 & 2 \\
& 体素与体积几何模型、SDF 距离场、稀疏体素网格与 CSG 构造实体几何集合布尔运算 & 解释/使用 & 领域核心 & 2 \\
& 过程式几何表征、分形拓扑、L-系统文法派生与生成式三维建模 & 解释/使用 & 领域核心 & 2 \\
\midrule
\makecell[l]{图形交互/高级着色} 
& 渲染方程蒙特卡洛高阶求解、时间动态模糊、镜头景深、级联阴影、遮挡剔除、面光源、Hi-Z、非真实感 NPR & 解释/使用 & 领域核心 & 6 \\
\midrule
\multirow{3}{*}{\makecell[l]{图形交互/计算机动画}} 
& 迪士尼经典动画十二原则（挤压拉伸、时间节奏等）、关键帧时间轴编辑与多维曲线插值（Lerp/Slerp/样条） & 解释/使用 & 领域核心 & 2 \\
& 骨骼装配体系、正向运动学（FK）与逆向运动学（IK）求解 & 解释/使用 & 领域核心 & 2 \\
& 空间几何变换复合链条、旋转多重表征、虚拟摄影机注视动画 & 解释/实现 & 领域核心 & 2 \\
\midrule
\multirow{3}{*}{\makecell[l]{图形交互/物理仿真}} 
& 粒子动力学系统、牛顿力学离散数值积分法（辛欧拉/Verlet） & 解释/实现 & 领域核心 & 2 \\
& 碰撞检测与动态物理响应算法（SDF 求交、分离轴定理） & 解释/实现 & 领域核心 & 2 \\
& 基于欧拉网格的流体动力学、半拉格朗日平流与泊松压力投影 & 解释/实现 & 领域核心 & 2 \\
\midrule
\multirow{3}{*}{\makecell[l]{图形交互/沉浸式技术}} 
& 扩展现实谱系定义辨析（VR/AR/MR）、医学与国防仿真应用 & 解释 & 领域核心 & 1 \\
& 双目立体显示光学、6-DoF 空间位姿追踪、动光延迟抑制 & 解释/实现 & 领域核心 & 3 \\
& 沉浸人因安全（晕动症规避）、无障碍通用设计与空间隐私伦理 & 解释/评估 & 领域核心 & 2 \\
\midrule
\multirow{2}{*}{\makecell[l]{图形交互/交互技术}} 
& 事件驱动编程范式、WIMP 单通道图形用户界面、力反馈外骨骼 & 解释/实现 & 领域核心 & 4 \\
& 图形用户界面无障碍等价通道设计与多模态脑机前沿评估 & 解释/评估 & 领域核心 & 2 \\
\midrule
\multirow{3}{*}{\makecell[l]{图形交互/数字图像处理}} 
& 空间离散卷积滤波核（高斯平滑、拉普拉斯边缘锐化） & 解释/实现 & 领域核心 & 2 \\
& 卷积神经网络与视觉 Transformer 架构（与 AI 领域交叉） & 解释/实现 & 领域核心 & 2 \\
& 色彩直方图均衡化、正交频域变换（傅里叶/离散余弦变换） & 解释/实现 & 领域核心 & 2 \\
\midrule
\multirow{4}{*}{\makecell[l]{图形交互/实体与物理计算}} 
& 与真实物理世界通信交互、传感器 ADC 采集与执行器 PWM 驱动 & 解释/实现 & 领域核心 & 1 \\
& 微控制器事件驱动固件架构、硬件定时器与外部中断响应 & 解释/实现 & 领域核心 & 1 \\
& 连接物理实体、CAD/CAM 数字制造（3D 打印/数控机床/吸塑成型） & 解释/评估/实现 & 领域核心 & 3 \\
& 物联网体系架构、轻量级网络通信协议与嵌入式物理安全 & 解释 & 领域核心 & 1 \\
\bottomrule
\end{tabularx}
\clearpage
\twocolumn
\chapter{人机交互}

\section{领域导言与演进背景}

\begin{preamblebox}[知识领域概述与变革]
计算系统不仅赋予用户解决现实复杂问题的能力，更促进了人与人之间的社会连接，全方位赋能人类多元化的实践与探索。因此，这些系统在与用户交互及解决问题时，必须恪守尊重个体尊严、维护社会正义、捍卫人类崇高价值与激发创造力的基本准则。人机交互 (HCI)正是一门从跨学科视角切入的核心领域，深度融汇了认知心理学、人因工程学、商业战略与交互设计原则。

每一个用户都是独一无二的。从人机交互的视角出发，任何与人类交互的系统设计，都必须前瞻性地预见并充分尊重这种多样性。这不仅涵盖无障碍辅助接入，更延伸至文化习俗与社会规范、神经多样性、多感官交互模态以及系统在用户心理与生理上激发的全方位反应。一个高效卓越的计算系统，应当在公平对待用户、严守隐私边界、构筑安全屏障并恪守伦理准则的同时，赢得用户的深层信任。

实现上述目标，必须确立以“设计为中心”的工程范式。该范式始于明确的设计意图，并深刻认识到设计是一个动态的迭代过程——需要针对系统的可用性及其对最终用户的深远影响进行持续且反复的评测与验证。此外，现代技术不仅通过其计算产出激发用户反馈，其感知环境与交互通信的多样化模态同样深刻塑造着人机体验。本知识领域致力于系统性提升计算机从业者对这些核心议题的敏感度与工程能力，理应深刻启迪与塑造每一位未来的计算机科学家与工程师。

\textbf{自 CS2013 以来的核心变革：}
\begin{enumerate}
    \item \textbf{知识单元重构与现代实践全面对接}：重新命名并系统重构了各知识单元，精准反映当前产业界面向敏捷交付、用户体验（UX）研究与多端融合的最新工程实践，并前瞻预见空间计算与生成式 AI 等前沿技术的演进趋势。
    \item \textbf{强化多样性本质与以用户为中心的设计}：更加突出用户群体在生理、认知与文化上的内在多样性本质，将“以用户为中心的设计 (UCD)”和“包容性设计”置于计算机专业教育的核心枢纽地位。
    \item \textbf{交互模态视角的现代化演进}：多模态交互（如文本、图形、语音、触觉与姿态）依然占据核心地位，但大纲降低了对孤立特定输入设备的机械强调，转而倡导更加与时俱进、富有同理心与情境感知的全通道交互工程方法。
    \item \textbf{深度关注系统社会影响与伦理可解释性}：课程体系高度强调系统评估计算软硬件对用户的长远影响，将可持续性、算法公平性、防范差别待遇、系统透明度、可信度与可解释性贯穿于交互全生命周期。
    \item \textbf{确立跨学科交叉领域的枢纽定位}：鉴于与计算机学科各分支的广泛交织，HCI 本身即是一个核心交叉领域，与人工智能（AI）、社会、伦理与专业素养（SEP）、软件开发基础（SDF）以及软件工程（SE）等知识领域形成了深度的协同互补。
\end{enumerate}
\end{preamblebox}

\section{核心学时分配建议}

本知识领域在“向日葵模型”下的学时推荐如下表所示：

\begin{table}[H]
\centering
\setlength{\tabcolsep}{3.5pt}\fontsize{8pt}{9.6pt}\selectfont\begin{tabularx}{\columnwidth}{Y c c}
\toprule
\textbf{知识单元} & \textbf{计算机核心} & \textbf{领域核心} \\
\midrule
理解用户 & 2 & 5 \\
设计中的问责与责任 & 2 & 2 \\
无障碍与包容性设计 & 2 & 2 \\
交互设计评估 & 1 & 2 \\
交互系统设计 & 1 & 5 \\
\midrule
\textbf{合计} & \textbf{8 学时} & \textbf{16 学时} \\
\bottomrule
\end{tabularx}
\end{table}

\noindent\textbf{说明}：知识单元 6（人机交互与社会伦理，HCI and SEP）属于贯穿整个计算机科学培养体系的跨领域交叉单元，其核心学时已与《社会、伦理与专业素养 (SEP)》深度有机融合，并在本领域各单元教学实践与期末综合项目中持续渗透。

\section{知识单元与详细大纲}

\subsection{理解用户：个体目标与人际交互}
\textbf{【学时建议】}：2 计算机核心学时，5 领域核心学时。

本单元奠定人本设计的核心认知基石。学生必须打破以自我为中心的工程师思维惯性，深刻理解真实用户的生理机能、认知模型与多元社会背景，掌握科学的用户需求挖掘与画像建模方法。

\subsubsection{计算机核心主题}
\begin{itemize}
    \item \textbf{以用户为中心的设计与评估方法}：
    \begin{itemize}
        \item “你不是用户”的核心哲学内涵及其对软件工程思维的重塑。
        \item 用户需求挖掘技术：探索性调研、定性观察与用户痛点定位。
        \item 形成性研究：在概念设计与早期原型阶段探寻需求规律。
        \item 核心调研手段：结构化/半结构化深度访谈、问卷调查法、基线可用性实测。
    \end{itemize}
\end{itemize}

\subsubsection{领域核心主题}
\begin{itemize}
    \item \textbf{用户中心设计方法论}：
    \begin{itemize}
        \item 用户画像构建与画像图谱：典型用户与极端边界用户的特征抽象与共情载体。
        \item 用户故事与叙事构建：敏捷用户故事模板、场景化故事板与叙事采集方法。
        \item 同理心地图：说（Says）、想（Thinks）、做（Does）、感受四象限建模。
        \item 进阶需求评估技术：民族志观察与情境调查。
        \item 用户体验旅程图：触点分析、情绪曲线映射与痛点识别（参见“交互设计评估”单元）。
    \end{itemize}
    \item \textbf{用户的生理机能与认知特征}：
    \begin{itemize}
        \item 指导交互设计的生理学基础：人机工程学、色彩感知生理机制（视网膜杆体/锥体细胞特性与色盲模拟）。
        \item 指导交互设计的认知心理学模型：注意力分配机制（选择性/持续性注意）、知觉感知与图形识别机制、运动控制模型（菲茨定律 Fitts' Law 初步）、人类记忆多存储模型（感觉记忆、工作记忆容量限制 $7\pm2$ 效应、长时记忆提取）。
        \item 社会与行为心理学课题：认知偏差（Confirmation Bias、锚定效应）、变化盲视与界面诱导心理。
    \end{itemize}
    \item \textbf{面向多元用户群体的设计}：
    \begin{itemize}
        \item 用户属性多元性：种族、生理能力、年龄阶层、性别、文化习俗、数字素养经验与受教育程度对交互体验与用户期望的决定性影响。
        \item 软件国际化 (Internationalization, i18n) 与跨文化设计：跨文化维度理论在界面布局、隐喻及信息架构中的映射。
        \item 有效设计评估面临的严峻挑战：被试抽样偏差、实验结论的普适泛化性局限；残障人士与功能受限者的真实交互体验感知。
        \item 通用设计哲学（参见“无障碍与包容性设计”单元）。
    \end{itemize}
    \item \textbf{协同交互与群体通信}：
    \begin{itemize}
        \item 多用户与组织情境下的用户认知特征：角色分工、权力结构与群组动力学。
        \item 同步群体通信系统：实时聊天、多方音视频会议协同、多人在线协作虚拟空间。
        \item 异步群体通信机制：电子邮件线程、主题论坛、协同文档版本流转、社交网络。
        \item 社交媒体、社交计算 与社交网络分析 (SNA) 初步。
        \item 在线协同、社会化协调机制与虚拟数字社区治理。
        \item 虚拟化身、数字化身交互与虚拟三维空间协同。
    \end{itemize}
\end{itemize}

\subsubsection{拓展选修与前沿主题}
\begin{itemize}
    \item 群体智能与协同众包计算系统架构；元宇宙沉浸式数字空间中具身虚拟化身的非语言交流与微表情映射。
\end{itemize}

\subsubsection{示范性学习成果与能力目标}
\begin{enumerate}
    \item \textbf{[应用]} 在一个真实工程软件开发项目中，规范组织并端到端实施“以用户为中心的设计（UCD）”全流程。
    \item \textbf{[解释]} 深入对比与辨析最终用户真实诉求与系统软件工程师/架构师先验技术偏见之间的本质分歧。
    \item \textbf{[解释]} 精准识别给定系统的代表性目标用户画像，并系统分析该系统可能对周边间接利益相关者产生的连锁影响。
    \item \textbf{[解释]} 阐明在现代交互设计全流程中，“同理心构建”与“科学评估”作为核心驱动要素的方法论意义。
    \item \textbf{[应用]} 独立组织针对真实目标用户的调研工作，并撰写符合工业规范的用户特征与多维度需求分析报告。
    \item \textbf{[开发]} 基于系统化的需求评估结果，归纳提炼出满足敏捷开发标准的高质量用户故事。
    \item \textbf{[开发]} 针对特定初始受众设计的产品方案进行可用性重构，使其精准适配文化背景或认知能力存在显著差异的新用户群体。
    \item \textbf{[评估]} 针对给定的实际设计难题，横向权衡并对比不同需求挖掘方法的适用场景、成本开销与信息捕获深度。
    \item \textbf{[评估]} 批判性反思并评估当前系统方案中引入低技术或无技术非数字组件所能带来的系统增益与体验改善。
    \item \textbf{[解释]} 深刻剖析多用户协作系统与社交计算情境对交互架构、并发响应与冲突消解带来的深层技术影响。
\end{enumerate}

\subsection{设计中的问责与责任：可持续性、安全、隐私、信任与伦理}
\textbf{【学时建议】}：2 计算机核心学时，2 领域核心学时。

软件交互并非价值中立。交互设计师与软件工程师做出的每一个界面与逻辑决策，均深刻影响着资源消耗、个人隐私、社会公平及人类身心健康。本单元引导学生建立负责任的工程伦理与问责意识。

\subsubsection{计算机核心主题}
\begin{itemize}
    \item \textbf{交互设计的系统性深远影响}：
    \begin{itemize}
        \item 绿色低碳与可持续性：算力与带宽感知的轻量化界面设计、减少数字废弃物与设备过度消耗。
        \item 社会包容性 与系统安全性 (Safety \& Security)：防范界面诱导导致的误操作与财产/人身安全风险。
        \item 数据隐私 与防范技术危害：知情同意权界定、敏感数据非必要不采集原则。
        \item 差别待遇与歧视性影响：算法推荐与界面交互对特定群体的隐性排斥。
    \end{itemize}
    \item \textbf{交互设计中的职业伦理}：
    \begin{itemize}
        \item 设计方法与工程解决方案中的伦理底线；人工智能技术介入交互系统（如生成式 AI、自主决策代理）时的边界规范。
        \item 统筹考量利益相关者多维影响与人因工程的责任担当。
        \item 坚持以人为本、以解决用户真实痛点为导向的设计使命，坚决抵制损害用户权益的病态 KPI 驱动设计。
    \end{itemize}
    \item \textbf{工程设计中的法规与合规约束}：
    \begin{itemize}
        \item 软件所有权归属与问责边界。
        \item 法律法规合规框架：通用数据保护条例（GDPR）、个人信息保护法、未成年人网络保护条例等合规基准。
        \item 超越用户即时表层欲望的长效考量：通过对工程问题本质的持续反思与迭代重构，防范成瘾机制与负面社会外部性。
    \end{itemize}
\end{itemize}

\subsubsection{领域核心主题}
\begin{itemize}
    \item \textbf{价值敏感设计}：
    \begin{itemize}
        \item 直接利益相关者 与间接利益相关者 的精准辨识与诉求映射。
        \item 系统梳理多元利益相关者的价值体系，建立兼顾尊严、自由、公正与信任的工程价值权衡模型。
    \end{itemize}
    \item \textbf{界面诱导与说服性设计}：
    \begin{itemize}
        \item 科学量化与评估交互界面中的说服性元素。
        \item 正向健康行为引导（如促进运动、助益专注）作为合法合规的设计目标。
        \item 批判性识别与抵制“暗黑模式”：诱导式收费、隐藏取消订阅通道、虚假紧迫感及情感绑架等恶意交互手段。
    \end{itemize}
\end{itemize}

\subsubsection{示范性学习成果与能力目标}
\begin{enumerate}
    \item \textbf{[评估]} 敏锐识别并批判性审视系统设计方案对社会及目标社群潜在的多维冲击，系统应对可持续性、包容性、安全性、隐私保护、潜在危害及差别对待等重大现实关切。
    \item \textbf{[解释]} 准确识别与剖析给定工程设计中潜在的人因工程学要素。
    \item \textbf{[解释]} 系统梳理并深刻洞察产品生态中的直接利益相关者与间接利益相关者群体。
    \item \textbf{[开发]} 拓展构建覆盖产品完整生命周期的情境剧本，前瞻预判超出初始设计意图之外的长周期社会与群体演变。
    \item \textbf{[评估]} 批判性审视系统设计中直接或间接影响利益相关者与更广泛社会的决定性因素（如算法透明度、系统可持续生存力、用户信任机制、AI 自主性边界等）。
    \item \textbf{[评估]} 严格审视产品界面与交互逻辑中的说服性元素及其设计初衷，客观评估其是否与用户的长远真实利益保持一致。
    \item \textbf{[评估]} 对照系统初始设计意图，客观评估系统在真实环境下的最终实施成效与衍生产出。
    \item \textbf{[解释]} 深刻领会关键架构与交互设计决策在社会、组织与个体层面引发的连锁反应与长远影响。
\end{enumerate}

\subsection{无障碍与包容性设计}
\textbf{【学时建议】}：2 计算机核心学时，2 领域核心学时。

信息无障碍是现代数字文明与人权保障的关键基石。计算机专业毕业生必须树立“无障碍设计使所有人受益”的崇高理念，精通国际主流无障碍规范标准，熟练掌握支持各类机能受限群体的辅助开发技术。

\subsubsection{计算机核心主题}
\begin{itemize}
    \item \textbf{社会责任与法律基准}：
    \begin{itemize}
        \item 保障残障群体数字平权的法律法规支撑与强制性义务（如中国《无障碍环境建设法》、美国 ADA 法案、欧盟欧洲无障碍法案 EAA）。
        \item 核心人本理念：“无障碍设计使所有人受益”——解决极端困难场景的设计往往能极大提升常规场景的用户体验。
    \end{itemize}
    \item \textbf{无障碍工程标准与技术规范}：
    \begin{itemize}
        \item 万维网联盟（W3C）Web 内容无障碍指南（WCAG 2.1/2.2 AA 级）核心技术原则：可感知性、可操作性、可理解性与健壮性。
    \end{itemize}
    \item \textbf{辅助技术与软硬件赋能产品}：
    \begin{itemize}
        \item 现代操作系统与浏览器内嵌的无障碍功能与辅助技术 (AT)：可访问性对象模型（AOM）、屏幕阅读器、高对比度模式、单手操作模式与物理开关控制。
    \end{itemize}
    \item \textbf{包容性设计框架}：
    \begin{itemize}
        \item 承认并正视个体能力维度的客观差异；通用设计七大核心原则。
    \end{itemize}
\end{itemize}

\subsubsection{领域核心主题}
\begin{itemize}
    \item \textbf{人口统计学与多元障碍形态谱系}：
    \begin{itemize}
        \item 三大典型机能障碍形态：永久性残障（Permanent，如全盲或截瘫）、暂时性机能受损（Temporary，如腕骨骨折或白内障术后）、情境性受限（Situational Disability，如强阳光下屏幕眩光或嘈杂环境中听觉受阻）。
        \item 残障认知的全球化视野与跨文化认知演进；消除技术研发中对特殊群体的隐性歧视与社会偏见。
    \end{itemize}
    \item \textbf{包容性用户体验研究与软件工程实践}：
    \begin{itemize}
        \item 邀请障碍人士参与的全流程包容性 UX 调研与测试方法。
        \item 前端与客户端工程层面的无障碍落地实践：WAI-ARIA 规范标签语义化注入、焦点顺序管理、键盘完全可达性导航、颜色对比度（$\ge 4.5:1$）自动化校验。
    \end{itemize}
    \item \textbf{先进辅助功能与适配技术}：
    \begin{itemize}
        \item 屏幕阅读器严格适配（如 NVDA、JAWS、TalkBack、VoiceOver）；图片多模态替代文本语义化生成；动态字体无级缩放与响应式重排防溢出处理。
    \end{itemize}
    \item \textbf{包容性协同流程与大社会效益}：
    \begin{itemize}
        \item 参与式设计与共同设计流程；以包容性驱动全社会数字鸿沟弥合。
    \end{itemize}
\end{itemize}

\subsubsection{拓展选修与前沿主题}
\begin{itemize}
    \item 认知除蔽与反思性质疑；跨学科残障学理论研究。
    \item 企业实施包容性设计的商业投资回报率量化评估模型。
    \item 用户敏感型包容性设计范式。
    \item 批判性人机交互研究范式：HCI 中的批判性种族理论、女权主义人机交互理论以及批判性残障理论。
\end{itemize}

\subsubsection{示范性学习成果与能力目标}
\begin{enumerate}
    \item \textbf{[解释]} 精准辨识视觉、听觉、肢体与认知等各类机能障碍人士面临的交互瓶颈，明确指出与之相匹配的无障碍与辅助技术解决方案。
    \item \textbf{[解释]} 阐明并应用适宜的包容性设计范式（如通用设计 Universal Design 与基于能力的设计 Ability-Based Design）。
    \item \textbf{[应用]} 熟练掌握并严格遵循主流软件与网络内容无障碍设计指南（如 WCAG 规范族）。
    \item \textbf{[解释]} 准确列举并理解适于无障碍交互软件开发的国内外法律法规与国家强制性标准。
    \item \textbf{[开发]} 灵活将包容性框架运用于视觉界面、语音交互界面（VUI）和触控交互界面的无障碍工程落地，确保全通道无障碍可达。
    \item \textbf{[应用]} 在与多元背景及障碍群体用户进行实地交流访谈与测试时，展现出高度专业、得体且尊重的跨群体沟通素养。
    \item \textbf{[评估]} 综合运用无障碍自动化检测工具与手工辅助技术走查，对现代 Web 页面及移动 App 进行严格的合规审计与无障碍漏洞定位。
    \item \textbf{[评估]} 探究在软硬件工程发展史上，针对残障、种族或性别的隐性偏见如何潜移默化地渗透到技术系统设计中，并从中提炼反思性指导法则。
    \item \textbf{[开发]} 独立构想并实施高阶用户体验研究方案，深度挖掘边缘弱势群体的潜在需求并转化为切实的设计洞察。
\end{enumerate}

\subsection{交互设计评估}
\textbf{【学时建议】}：1 计算机核心学时，2 领域核心学时。

未经科学评估的设计只是未经检验的技术假设。本单元培养学生掌握定性、定量与混合评估范式，能够严谨设计并实施可用性测试实验，借助数理统计工具与质性分析法对交互系统做出可辩护的科学推断。

\subsubsection{计算机核心主题}
\begin{itemize}
    \item \textbf{以用户为中心的评测核心范式}：
    \begin{itemize}
        \item 形成性评估：需求探索、早期纸质原型走查与概念纠偏。
        \item 总结性评估：核心功能验收、交互效率基准测试与可用性达标量化评测。
    \end{itemize}
    \item \textbf{交互系统核心评估维度}：
    \begin{itemize}
        \item 实用性：系统能否真正支撑目标任务的达成。
        \item 操作效率：熟练用户完成既定任务所需耗费的精力、时间与操作击键成本。
        \item 易学性：新手用户初次接触系统达到基线操作水平的学习曲线与门槛。
        \item 用户主观满意度：系统交互引发的用户心理认同与情感愉悦度。
    \end{itemize}
    \item \textbf{伦理审查要求}：
    \begin{itemize}
        \item 在开展涉及人类受试者的可用性测试与数据采集前，必须获取伦理委员会批准（IRB 审查），确保知情同意与隐私脱敏。
    \end{itemize}
\end{itemize}

\subsubsection{领域核心主题}
\begin{itemize}
    \item \textbf{定性、定量与混合研究方法矩阵}：
    \begin{itemize}
        \item 质性分析技术：测试会话录像转录、质性扎根编码与主题提炼分析。
        \item 定量评估方法：任务成功率、任务完成时长、错误率统计与标准心理量表（SUS 系统可用性量表、NASA-TLX 认知负荷量表）。
        \item 混合研究方法：结构化现场观察、出声思考法、深度半结构化访谈与受控对比实验。
        \item 专家走查法：启发式评估（Heuristic Evaluation，重点掌握 Nielsen 十大可用性启发式原则）、认知走查法。
        \item 专业呈现与交付：撰写高质量可用性工程评估白皮书、提炼严谨的用户画像。
        \item 评测方法论瓶颈与挑战：被试抽样偏差、实验生态效度与泛化性限制。
    \end{itemize}
    \item \textbf{可用性实验设计与严谨规划}：
    \begin{itemize}
        \item 明确研究目标，确立科学、可证伪的实验原假设（$H_0$）与备择假设（$H_1$）。
        \item 机构伦理审查委员会（IRB）与伦理合规备案审批全流程。
        \item 开放科学框架下的实验方案预注册机制。
        \item 经典实验设计范式比较：组内设计（Within-Subjects Design，平衡顺序效应与疲劳效应）与组间设计的优劣权衡与控制变量准则。
    \end{itemize}
    \item \textbf{设计产出的综合长效影响评测}：
    \begin{itemize}
        \item 系统对自然生态环境、物质资源消耗、社会结构伦理、网络空间安全与数据隐私构架的深层综合影响评估。
    \end{itemize}
\end{itemize}

\subsubsection{拓展选修与前沿主题}
\begin{itemize}
    \item 定量高级统计分析工具包：R 语言、Python Scipy/Statsmodels、JASP 等在交互数据中的实战；可视分析与探索性数据分析（EDA）。
    \item 统计假设检验进阶：单因素/多因素方差分析（ANOVA）、双样本独立/配对 $t$ 检验、事后多重检验校正（Bonferroni 校正）；参数检验与非参数检验（Wilcoxon 符号秩检验、Mann-Whitney $U$ 检验）的选型条件；效应量（Cohen's $d$、$\eta^2$）计算与统计功效分析。
    \item 评估研究数据合规管理：基于开放科学原则的数据长期安全归档与脱敏共享机制；受试者敏感个人可识别信息（PII）的去标识化与差分隐私注入。
\end{itemize}

\subsubsection{示范性学习成果与能力目标}
\begin{enumerate}
    \item \textbf{[解释]} 深入辨析形成性评估与总结性评估的本质特征、介入节点及其在驱动交互设计持续演进中的关键作用。
    \item \textbf{[应用]} 针对软件系统研发的不同生命周期阶段（从概念草图到发布后监控），精准选取并编排最适宜的形成性或总结性评测方法。
    \item \textbf{[解释]} 深入论证在交互评测中将定性洞察与定量数据相结合的互补优势。
    \item \textbf{[评估]} 针对给定的实际设计方案，全面评估其对社会伦理、环境消耗与用户社群带来的综合宏观影响。
    \item \textbf{[开发]} 独立策划针对特定交互界面的可用性评测实验方案，严谨论证其研究目标设定、实验假设推导及被试控制组设计的合理性。
    \item \textbf{[评估]} 规范组织并执行界面可用性评测，收集多模态实验数据，并在实验范式约束下得出可辩护、高置信度的评估结论。
    \item \textbf{[应用]} 运用专业统计软件对实验收集的行为日志与量表数据执行假设检验，科学解读显著性水平（$p$ 值）与置信区间。
    \item \textbf{[应用]} 完成一份严格符合学术与工业规范的可用性实验设计预注册报告。
\end{enumerate}

\subsection{交互系统设计}
\textbf{【学时建议】}：1 计算机核心学时，5 领域核心学时。

本单元系统阐述从概念草图到工业级软件界面的完整工程构建方法。重点涵盖低保真与高保真原型工具栈、现代前端与客户端软件架构模式、双钻设计流程、格式塔视觉心理学原则以及多通道先进交互技术。

\subsubsection{计算机核心主题}
\begin{itemize}
    \item \textbf{原型设计技术与工具栈}：
    \begin{itemize}
        \item 低保真原型设计：纸质线框图绘制、流程草图、低成本走查与概念试错。
        \item 快速原型制作 与探索废弃型原型。
        \item 原型的交互粒度控制：外观保真度、交互深度与数据动态性的合理平衡。
    \end{itemize}
    \item \textbf{交互系统设计模式}：
    \begin{itemize}
        \item 迭代设计闭环：需求分析 $\to$ 原型构思 $\to$ 实施开发 $\to$ 评估反馈 $\to$ 重构演进。
        \item 通用设计模式与现代交互模式：数据驱动设计与事件驱动交互体系。
    \end{itemize}
    \item \textbf{工程设计约束与权衡}：
    \begin{itemize}
        \item 跨平台异构环境约束（Web、移动端 iOS/Android、桌面端、车载系统、嵌入式终端）。
        \item 设备物理特性限制（屏幕视网膜像素密度、算力能耗配额、触控精度、网络带宽与离线状态）。
    \end{itemize}
\end{itemize}

\subsubsection{领域核心主题}
\begin{itemize}
    \item \textbf{软件架构模式与设计指南}：
    \begin{itemize}
        \item 客户端交互架构：MVC、MVVM、声明式 UI 响应式状态管理（如 React/Flutter 单向数据流）。
        \item 跨平台一致性设计系统：组件库、设计令牌与原子设计理论。
        \item 分布式与多设备协同交互环境下的状态一致性与实时同步挑战。
    \end{itemize}
    \item \textbf{先进设计思维与过程范式}：
    \begin{itemize}
        \item 参与式设计 与共同设计。
        \item 英国设计协会“双钻模型”：探索期发散思维 $\to$ 定义期收敛思维 $\to$ 研发期发散构思 $\to$ 交付期收敛验证。
    \end{itemize}
    \item \textbf{先进交互技术与通道向量}：
    \begin{itemize}
        \item 输入与输出多维通道向量：手势追踪、姿态检测、多点触控、语音指令、力敏反馈。
        \item 现代图形用户界面 (GUI) 精细化组件、物理硬件控制器、触觉反馈技术。
        \item 交互容错机制与优雅降级：防呆设计、异常状态回滚、高容错操作撤销。
    \end{itemize}
    \item \textbf{视觉交互界面设计心理学法则}：
    \begin{itemize}
        \item 色彩理论、对比度阶层与排版网格系统。
        \item 格式塔知觉心理学原理：邻近原则、相似原则、连续原则、闭合原则、主体与背景原则。
    \end{itemize}
\end{itemize}

\subsubsection{拓展选修与前沿主题}
\begin{itemize}
    \item 沉浸式交互空间：虚拟现实（VR）、增强现实（AR）、混合现实（MR）与扩展现实（XR）交互套件开发；空间音频声场定位。
    \item 3D 打印与智能快速成型制造（3D Printing and Digital Fabrication）。
    \item 异步交互协作模型。
    \item 创造力支持工具 (CST)设计原理。
    \item 会话式用户界面（CUI）与生成式智能体对话交互设计。
\end{itemize}

\subsubsection{示范性学习成果与能力目标}
\begin{enumerate}
    \item \textbf{[开发]} 针对特定的应用场景与任务特征，提出并构建量身定制的最优交互模式与系统架构设计方案。
    \item \textbf{[开发]} 严格践行“深度理解用户 $\to$ 增量式原型构建 $\to$ 严密评测反馈 $\to$ 指导下一轮迭代演进”的敏捷交互开发闭环。
    \item \textbf{[解释]} 深入阐述面对动态变化的工程约束（底层硬件变迁、算力限制、用户多元化需求、安全合规收紧）时，系统架构所必须做出的权衡取舍。
    \item \textbf{[评估]} 紧扣核心项目目标，系统评估并抉择主流前端与客户端软件架构模式的优劣。
    \item \textbf{[解释]} 准确定位分布式与多设备协同交互环境中面临的网络并发、状态同步不一致与体验延迟等核心工程挑战。
    \item \textbf{[评估]} 针对特定真实业务情景，深入评估并横向比对不同交互输入技术可能引发的个人数据与生物特征隐私风险。
    \item \textbf{[解释]} 深度运用格式塔等经典视觉认知心理学法则，对特定界面的视觉层级、元素布局与信息架构做出专业的技术原理解释。
    \item \textbf{[评估]} 深入评估空间计算与 VR/AR/MR 虚实融合情境下所暴露的用户空间注视点数据、生物识别及环境空间映射的隐私安全隐患。
\end{enumerate}

\subsection{人机交互与社会、伦理及专业素养的交叉融合}
\textbf{【说明】}：本单元为跨领域核心贯穿单元，与《社会、伦理与专业素养 (SEP)》深度锚定，致力于将人本设计思想与工程师社会责任深度融合，不单独列出静态学时，贯穿于整套课程教学。

\subsubsection{计算机核心主题}
\begin{itemize}
    \item \textbf{通用设计与以用户为中心的设计之伦理意涵}：将普遍可达性、用户尊严维护、反歧视伦理确立为软件工程需求分析与架构设计的第一性原则。
\end{itemize}

\subsubsection{领域核心主题}
\begin{itemize}
    \item \textbf{全流程参与式与包容性设计流程}：保障边缘群体与非技术利益相关者在软件全生命周期中的知情权与决策参与权。
    \item \textbf{交互设计的全方位外部性审计}：系统评测软硬件系统对自然生态环境、物质资源消耗、社会正义、数字安全、数据隐私、职业道德及长远文明构架的系统影响。
\end{itemize}

\subsubsection{示范性学习成果与能力目标}
\begin{enumerate}
    \item \textbf{[应用]} 将以人为本的通用设计规范内嵌于工业软件需求工程与全生命周期开发规范中。
    \item \textbf{[解释]} 深刻辨识计算机工程师在构建大众交互产品时所肩负的法律、伦理与社会正义责任。
    \item \textbf{[评估]} 批判性剖析业界近期发生的缺乏包容性设计或算法偏见的真实工程反面案例，深入论述其引发的恶劣社会后果，并提出针对性的架构整改方案。
    \item \textbf{[评估]} 针对给定的交互软件，系统识别并评估其对生态环境、社会弱势群体与长效文明构建带来的深层外部性影响。
    \item \textbf{[评估]} 全面评估沉浸式虚拟现实（VR/AR/MR）交互情境下的空间隐私与用户数字踪迹保护机制。
\end{enumerate}

\section{专业素养与思维倾向}

在人机交互领域的学习与工程实践中，重点塑造学生的以下专业素养与思维倾向：
\begin{itemize}
    \item \textbf{高度适应性}：在面对快速更迭的软硬件形态、不确定的用户群体以及动态变化的业务需求时，能够敏锐调整设计范式与技术选型。
    \item \textbf{严谨细致的一丝不苟}：在微交互动效、界面边界留白、字体排版层级、无障碍焦点顺序以及异常容错提示等每一个微小工程细节中保持极致追求。
    \item \textbf{深具同理心}：能够超越技术开发者固有的“专家视角”，真正换位思考，体察不同文化背景、不同生理认知条件及边缘群体的真实喜怒哀乐与交互痛点。
    \item \textbf{团队协作导向}：优秀的 HCI 实践者必须将团队的整体成功置于首位，善于在跨职能团队中架起工程师、产品经理、设计师、业务人员与终端用户之间的沟通桥梁。
    \item \textbf{富有创造力}：能够基于历史工程实践、目标受众诉求及 HCI 核心理论，通过非凡的创造力将技术可行性与人本体验完美融合，设计出既优雅又高效的破局方案。
\end{itemize}

\section{数学基础要求}

人机交互领域高度依赖应用数理统计学，以支撑对可用性评估实验数据、A/B 测试结果及海量用户行为日志的科学解析：
\begin{itemize}
    \item \textbf{必须具备的基础统计学核心}：
    \begin{itemize}
        \item 集中趋势度量：均值、中位数、众数在偏态分布与长尾数据下的选用边界。
        \item 离散程度与变异性度量：方差、标准差、四分位距（IQR）与极差。
        \item 频数分布与概率分布规律：正态分布、二项分布与泊松分布；理解长尾分布对用户体验基准测试的影响。
        \item 推断统计基础直觉：参数置信区间、显著性检验（$p$ 值）以及假设检验基本原理。
    \end{itemize}
\end{itemize}

\section{交叉主题与跨学科共享概念}

人机交互作为一门承上启下的枢纽学科，与计算机科学其他关键领域深度交织：
\begin{itemize}
    \item \textbf{与人工智能领域 (AI) 的共享概念}：
    \begin{itemize}
        \item 个人智能助理：多模态自然交互与意图理解。
        \item 机器人学：人机共融协作与人机信任机制。
        \item 增强智能：AI 作为人类智力放大器而非简单替代工具的人机协作范式。
        \item 自主系统中的人工干预与控制权覆盖机制：保障人类核心掌控权的“人在回路”设计。
    \end{itemize}
    \item \textbf{与社会、伦理与专业素养 (SEP) 的共享概念}：
    \begin{itemize}
        \item 数据隐私与行为追踪保护：用户界面中的透明授权与隐私合规。
        \item 数字包容性与社会平权：跨越数字鸿沟与保障弱势群体福祉。
    \end{itemize}
    \item \textbf{与软件工程 (SE) 的共享概念}：
    \begin{itemize}
        \item 产品需求工程与用户故事定义：以人为本的需求引出。
        \item 软件系统验证与可用性确认：形式化测试与真实用户评测的互补。
        \item 与多方利益相关者的协同对接：多角色诉求对齐。
        \item 现代化原型测试、界面协同设计与跨职能沟通工具链。
    \end{itemize}
\end{itemize}

\section{真实能力任务规范}

\begin{taskbox}[任务 1：现代用户界面开发框架横向技术选型与深度权衡对比]
\textbf{【任务情境】}：某企业拟启动大型跨平台数字化业务系统重构，技术委员会需要在主流客户端与前端开发框架（如声明式响应式框架、跨平台自绘引擎与原生平台套件）之间进行工程选型。\\
\textbf{【能力陈述】}：系统阐述特定 UI 框架在应对特定软件工程难题时的优势与局限，清晰证明“界面设计与开发是目标驱动且充满多维权衡”的深刻认知，综合平衡渲染性能、开发维护成本、多端一致性及原生无障碍 API 支持。\\
\textbf{【能力领域】}：软件工程 \quad \textbf{【能力单元】}：需求分析、架构设计、技术文档编制、系统评估与项目管理\\
\textbf{【核心要求】}：计算机核心 \quad \textbf{【达标等级】}：\textbf{评估}
\end{taskbox}

\begin{taskbox}[任务 2：企业核心业务系统可用性研究实验的全流程规划、文档编制与实施]
\textbf{【任务情境】}：某全国性在线医疗问诊挂号平台改版，需针对中老年及视障受试者群体的操作阻塞点，组织严密的可用性实测。\\
\textbf{【能力陈述】}：深刻理解反馈闭环在全生命周期人本设计中的关键驱动作用；根据系统所处的研发阶段与目标受众的认知特征，独立制定兼顾伦理审批、被试招募、任务设计的可用性测试方案，严密组织出声思考走查并产出高质量分析白皮书。\\
\textbf{【能力领域】}：应用工程 \quad \textbf{【能力单元】}：系统设计、技术文档、系统评估、项目管理、用户接受度验证、社会议题调适、系统持续改进\\
\textbf{【核心要求】}：计算机核心 \quad \textbf{【达标等级】}：\textbf{应用与评估}
\end{taskbox}

\begin{taskbox}[任务 3：公共政务系统的可用性、无障碍性与包容性综合审计与改进方案]
\textbf{【任务情境】}：审计某市政“一网通办”数字政务门户，系统排查各类残障人士、高龄老人及数字技能欠缺群体在表单填写与办事指南浏览中的交互障碍。\\
\textbf{【能力陈述】}：熟练运用多元评估透镜（可用性十大原则、W3C WCAG 2.1 AA 级标准与包容性指标），全面诊断面向公众的交互系统在易用性、屏幕阅读器兼容性及排版对比度上的缺陷，制定具有工程可行性的界面重构与代码整改技术方案。\\
\textbf{【能力领域】}：应用工程 \quad \textbf{【能力单元】}：系统评估、技术文档、社会议题调适、系统持续改进\\
\textbf{【核心要求】}：计算机核心 \quad \textbf{【达标等级】}：\textbf{评估}
\end{taskbox}

\begin{taskbox}[任务 4：跨职能敏捷团队协同推进智能交互系统的设计、评估与快速迭代]
\textbf{【任务情境】}：在一个由软件架构师、前端开发人员、产品经理、UI/UX 设计师及行业临床专家组成的敏捷跨职能团队中，协同打造面向复杂工作流的移动端协作系统。\\
\textbf{【能力陈述】}：系统整合应用人机交互各项知识单元，跨越专业术语壁垒进行高效技术沟通；将模糊的用户期望转化为明确的技术约束与接口规范，依托敏捷反馈闭环推动软件增量开发与持续重构，切实解决现实工业级难题。\\
\textbf{【能力领域】}：软件与应用工程 \quad \textbf{【能力单元】}：需求工程、系统设计、系统开发、技术文档、系统评估、项目管理、用户接受度、系统改进\\
\textbf{【核心要求】}：计算机核心 \quad \textbf{【达标等级】}：\textbf{开发}
\end{taskbox}

\begin{taskbox}[任务 5：基于高层设计原型编制高精度用户界面工程开发规格说明书]
\textbf{【任务情境】}：交互设计团队交付了业务系统高保真原型与设计规范系统，工程研发团队需要将其转化为严格、无歧义的前端与移动端软件工程实施规范。\\
\textbf{【能力陈述】}：紧扣可用性标准、用户真实需求与既定设计方案，编写详尽的软件用户界面需求规格说明书（SRS），精准界定页面状态机迁移、组件封装参数、网络弱网降级逻辑、响应式布局断点及前后端接口通信契约。\\
\textbf{【能力领域】}：软件与应用工程 \quad \textbf{【能力单元】}：需求工程、技术文档编制\\
\textbf{【核心要求】}：计算机核心 \quad \textbf{【达标等级】}：\textbf{开发}
\end{taskbox}

\begin{taskbox}[任务 6：支撑全球化出海业务系统的多语言与跨文化国际化 (i18n) 架构重构]
\textbf{【任务情境】}：某企业原本面向本土市场的企业管理与协同办公软件拟拓展至中东、欧洲等多国市场，系统出现布局溢出崩溃、文字截断、阅读流向错误与文化禁忌符号等严重问题。\\
\textbf{【能力陈述】}：系统化评估现有软件代码与界面实现，准确定位受文化、语言和地域习俗影响的硬编码文本、图标隐喻、日期/货币格式以及从右至左（RTL）排版依赖，提出全面的国际化与本地化架构重构规划。\\
\textbf{【能力领域】}：软件/系统/应用工程 \quad \textbf{【能力单元】}：系统设计、系统评估、社会议题调适、系统持续改进\\
\textbf{【核心要求】}：计算机核心 \quad \textbf{【达标等级】}：\textbf{评估}
\end{taskbox}

\begin{taskbox}[任务 7：实施形成性用户调研与实验，为全新交互系统提供战略级设计演进建议]
\textbf{【任务情境】}：某科技公司拟立项研发一款基于生成式大模型的专业程序员辅助编程助手，处于概念预研期，需摸清开发者的编码心智模型与核心交互痛点。\\
\textbf{【能力陈述】}：科学设计、组织并执行形成性用户体验研究（情境调查、出声思考走查、深度访谈），精准辨识交互原型的核心优势、认知阻碍与工程痛点，提炼出驱动系统架构长效演进的高价值设计策略建议。\\
\textbf{【能力领域】}：应用工程与理论计算 \quad \textbf{【能力单元】}：方案设计、测试验证、用户接受度评测、系统改进\\
\textbf{【核心要求】}：计算机核心 \quad \textbf{【达标等级】}：\textbf{评估}
\end{taskbox}

\begin{taskbox}[任务 8：跨部门解读、提炼并传达多源用户评估反馈，驱动研发版本敏捷迭代]
\textbf{【任务情境】}：某大型 SaaS 平台在灰度内测阶段，来自企业 IT 主管、一线业务操作人员、客服工单及系统埋点日志的多源反馈出现严重分歧与术语隔阂。\\
\textbf{【能力陈述】}：科学搭建反馈收集矩阵；结合不同角色评估者的立场视角、专业语境与优先级诉求，对海量质性与量化反馈进行系统归纳与根因溯源，撰写层次分明的技术摘要，制定指导软件敏捷迭代的整改路线图。\\
\textbf{【能力领域】}：软件工程 \quad \textbf{【能力单元】}：需求工程、测试验证、技术文档、用户接受度、社会议题调适、系统改进\\
\textbf{【核心要求】}：计算机核心 \quad \textbf{【达标等级】}：\textbf{评估}
\end{taskbox}

\begin{taskbox}[任务 9：面向特定敏感脆弱人群部署交互系统的社会伦理与技术风险综合审计]
\textbf{【任务情境】}：某智慧校园项目拟在中小学课堂全面部署基于人脸微表情识别与注视点追踪的注意力实时监测与情绪预警系统。\\
\textbf{【能力陈述】}：深入调研未成年人等特殊脆弱受众群体的身心发展特征；从数据隐私泄露、知情同意权缺位、心理胁迫焦虑、算法偏见与权力结构不对等维度，全面审计该系统潜在的多维负面社会冲击，出具严密的伦理风险审查与工程叫停/整改报告。\\
\textbf{【能力领域】}：应用工程 \quad \textbf{【能力单元】}：需求工程、系统评估、用户接受度审查、社会议题调适、系统改进\\
\textbf{【核心要求】}：计算机核心 \quad \textbf{【达标等级】}：\textbf{评估}
\end{taskbox}

\begin{taskbox}[任务 10：遵循工业级规范标准，端到端高保真构建可交互用户界面原型]
\textbf{【任务情境】}：根据产品规格说明书与主流工业级设计规范（如 Material Design / Apple HIG），为某物联网智能中控面板独立开发可交互、高保真的全功能软件界面原型。\\
\textbf{【能力陈述】}：严格依据软件设计规格书与专业工程师行为规范，熟练运用现代组件库、响应式布局引擎与状态管理模型，开发出视觉一致、交互灵敏、支持无障碍遍历且达到生产级工业标准的高保真交互原型系统。\\
\textbf{【能力领域】}：软件、应用与理论工程 \quad \textbf{【能力单元】}：需求转化、交互设计、系统开发\\
\textbf{【核心要求】}：计算机核心 \quad \textbf{【达标等级】}：\textbf{开发}
\end{taskbox}

\section{课程设计与开课规划建议}

联合工作组面向各高校计算机科学专业主修、辅修及跨学科方向，推荐三套典型示范课程方案：

\subsection{计算机专业主修与辅修核心课：《人机交互导论》（建议 30 学时）}
\begin{itemize}
    \item \textbf{学时结构分配}：
    \begin{itemize}
        \item 人机交互/理解用户 （单元 1）：7 学时
        \item 人机交互/设计中的问责与责任 （单元 2）：2 学时
        \item 人机交互/无障碍与包容性设计 （单元 3）：4 学时
        \item 人机交互/交互设计评估 （单元 4）：3 学时
        \item 人机交互/交互系统设计 （单元 5）：10 学时
        \item 人机交互/社会伦理交叉贯穿 （单元 6）：2 学时
    \end{itemize}
    \item \textbf{先修要求}：敏捷软件开发基础（或基础软件工程方法）。
    \item \textbf{能力目标陈述}：修完本课程的学生应能够透彻阐述以用户为中心的设计（UCD）原则，并将其系统应用于小型工程项目的全流程设计、实现与评估中。
    \item \textbf{课程特色与教学模式}：本课程采取“项目驱动、贯穿迭代”的教学范式。前 4 周奠定扎实的交互设计基础；在第 5--6 周，学生通过一个由教师给定的短期微型项目进行热身实操，激发探索深层技术概念的学习热情；在此基础上，学生团队自主对接真实社区或特定用户群体，运用快速原型迭代法完成完整的需求分析、交互设计、工程编码与可用性评估闭环。
    \item \textbf{15 周学期教学进度大纲}：
    \begin{enumerate}
        \item 交互设计导论与核心原则 （单元 1, 5）
        \item 认知心理学视角：思考、行动与评估三元模型 （单元 1, 4, 5）
        \item 人类记忆模型与操作失误防御设计 （单元 1, 3, 4）
        \item 交互设计原则与敏捷过程模型 （单元 5）
        \item[5--6.] 设计流程与现代软件工程开发的深度融合实战 （单元 5）
        \item[7.] 设计思维与启发式评估实战 （单元 1, 4）
        \item[8.] 无障碍设计规范（WCAG）与包容性技术 （单元 2, 3, 6）
        \item[9.] 视觉界面设计美学与用户画像进阶 （单元 1, 3, 4, 5）
        \item[10.] 期末大项目：用户同理心调研与需求精准界定 （单元 1～6）
        \item[11.] 期末大项目：头脑风暴构思与低保真原型制作 （单元 1, 4, 5）
        \item[12--13.] 期末大项目：交互系统工程实现与编码 （单元 5）
        \item[14.] 期末大项目：受试者可用性实测与量化评测 （单元 4）
        \item[15.] 期末大项目：成果路演汇报与评估文档编制 （单元 1～6）
    \end{enumerate}
\end{itemize}

\subsection{跨学科特色核心课：《数据可视化导论》（建议 47 学时）}
\begin{itemize}
    \item \textbf{学时结构分配}：
    \begin{itemize}
        \item 图形交互/数据可视化单元：30 学时
        \item 图形交互/基础渲染技术单元：10 学时
        \item 人机交互/理解用户单元：3 学时
        \item 社会伦理/多元性、隐私与数据伦理单元：4 学时
    \end{itemize}
    \item \textbf{先修要求}：面向对象程序设计基础 (CS2)、线性代数。
    \item \textbf{能力目标陈述}：修完本课程的学生应能够科学筛选与清理复杂数据集，验证数据的真实性、代表性与适用边界；设计、开发并测试一套直观揭示数据内在规律、高度易用且交互流畅的数据可视化分析系统。
\end{itemize}

\subsection{进阶实战专业课：《可用性工程与系统评测》}
\begin{itemize}
    \item \textbf{先修要求}：人机交互导论、社会科学研究方法、应用数理统计学。
    \item \textbf{课程特色与教学模式}：本课程为完全基于真实工程项目的进阶实战课，深度聚焦工业级交互产品的正式形式化评估。课程全面涵盖：评测目标科学确立、目标受试者分层抽样招募、测试任务与情境剧本构建、实验室环境搭建、测试方案编写与执行、实验日志与行为数据分析、评估白皮书撰写与整改建议答辩。使学生掌握针对复杂软硬件系统在多元真实场景下开展可用性工程评测的专业技术、规程与标准协议。
    \item \textbf{14 个核心专题模块}：
    \begin{enumerate}
        \item 可用性评测研究方案的系统规划
        \item 核心评测目标与科学假设确立
        \item 被试群体特征界定与代表性抽样策略
        \item 典型工程测试任务选取与情境剧本设计
        \item 可用性量化度量指标体系构建
        \item 测试指导手册、保密协议与知情同意书编制
        \item 测试软硬件设备配置与多机位实验室搭建
        \item 预实验与方案校准
        \item 正式可用性测试的主持与出声思考法引导
        \item 行为日志与量表数据制表统计与质性扎根编码
        \item 提出优先级分明的界面重构与代码整改建议
        \item 评估结论的专业化可视化沟通与呈现
        \item 制作典型操作阻塞点“高光剪辑录像”
        \item 推动产品架构持续演进与研发流程变革
    \end{enumerate}
    \item \textbf{核心能力产出}：
    \begin{itemize}
        \item 独立设计严密规范的工业级可用性测试方案；
        \item 依据严格抽样标准招募并管理多元代表性受试者；
        \item 熟练主持并执行无偏倚的高质量可用性测试会话；
        \item 科学统计实验数据，深度剖析根因并提出可辩护的重构方案；
        \item 面向工程团队与管理层清晰汇报评估发现与商业洞察；
        \item 撰写完全符合行业规范的软件可用性评估与改进工程白皮书。
    \end{itemize}
\end{itemize}

\section{委员会专家名单}

\begin{itemize}
    \item \textbf{主席}：
    \begin{itemize}
        \item Susan L. Epstein (纽约市立大学亨特学院及研究生中心，美国)
    \end{itemize}
    \item \textbf{委员}：
    \begin{itemize}
        \item Sherif Aly (开罗美国大学，埃及)
        \item Jeremiah Blanchard (佛罗里达大学，美国)
        \item Zoya Bylinskii (Adobe 研发实验室，美国)
        \item Paul Gestwicki (鲍尔州立大学，美国)
        \item Susan Reiser (北卡罗来纳大学阿什维尔分校，美国)
        \item Amanda M. Holland-Minkley (华盛顿与杰斐逊学院，美国)
        \item Ajit Narayanan (谷歌公司，美国)
        \item Nathalie Riche (微软全球研究院，美国)
        \item Kristen Shinohara (罗切斯特理工学院，美国)
        \item Olivier St-Cyr (多伦多大学，加拿大)
    \end{itemize}
\end{itemize}

\clearpage
\onecolumn
\section{附录：核心主题与技能要求等级对照表}

下表系统梳理人机交互知识领域各知识单元的核心知识点、认知等级（Bloom 分类）、知识范畴（计算机核心 / 领域核心）及推荐学时分配：

\begin{longtable}{p{0.18\textwidth} p{0.48\textwidth} c c c}
\toprule
\textbf{知识单元} & \textbf{核心知识点与具体主题} & \textbf{认知等级} & \textbf{分类} & \textbf{学时} \\
\midrule
\endfirsthead
\toprule
\textbf{知识单元} & \textbf{核心知识点与具体主题} & \textbf{认知等级} & \textbf{分类} & \textbf{学时} \\
\midrule
\endhead
\bottomrule
\endfoot
\multirow{4}{*}{\makecell[l]{人机交互/理解用户}} 
& 用户中心设计哲学（“你不是用户”）、需求挖掘、访谈与问卷调研 & 应用 & 计算机核心 & 2 \\
& 用户画像图谱、敏捷用户故事、同理心地图与用户体验旅程图构建 & 应用 & 领域核心 & 2 \\
& 用户的生理感知、认知注意力机制、记忆模型与操作认知偏差 & 解释 & 领域核心 & 2 \\
& 多元用户群体适配、国际化与跨文化设计、多用户群体协同通信计算 & 评估 & 领域核心 & 1 \\
\midrule
\multirow{2}{*}{\makecell[l]{人机交互/设计中的\\问责与责任}} 
& 系统设计影响（绿色可持续性、社会包容性、安全、隐私与歧视防范） & 评估 & 计算机核心 & 2 \\
& 价值敏感设计（VSD）、多方利益相关者体系、正向引导与防暗黑模式 & 评估 & 领域核心 & 2 \\
\midrule
\multirow{3}{*}{\makecell[l]{人机交互/无障碍与\\包容性设计}} 
& 无障碍法律法规、W3C WCAG 2.1/2.2 规范准则与系统级辅助技术支持 & 应用 & 计算机核心 & 2 \\
& 永久/暂时/情境障碍分类、屏幕阅读器适配与无障碍前端工程落地 & 应用 & 领域核心 & 1 \\
& 能力导向设计理念、包容性参与式研发流程与全人群数字鸿沟弥合 & 评估 & 领域核心 & 1 \\
\midrule
\multirow{2}{*}{\makecell[l]{人机交互/交互设计评估}} 
& 形成性评估与总结性评估辨析、可用性四大核心指标与伦理审查规范 & 解释 & 计算机核心 & 1 \\
& 质性编码主题分析、出声思考走查、启发式评估、对比实验与统计检验 & 评估 & 领域核心 & 2 \\
\midrule
\multirow{3}{*}{\makecell[l]{人机交互/交互系统设计}} 
& 低保真与快速原型制作技术、敏捷迭代设计闭环与多端硬件约束权衡 & 应用 & 计算机核心 & 1 \\
& 交互架构模式（MVC/MVVM/声明式 UI）、双钻设计流程、格式塔视觉法则 & 应用 & 领域核心 & 3 \\
& 多模态交互输入技术（手势/语音/触觉）、交互容错机制与优雅降级 & 开发 & 领域核心 & 2 \\
\midrule
\multirow{1}{*}{\makecell[l]{人机交互/社会伦理交叉}} 
& 通用设计人本伦理、全流程参与式包容研发与长效社会外部性审计 & 评估 & 领域核心 & 贯穿渗透 \\
\end{longtable}

\chapter{数学与统计学基础}

\section{领域导言与演进背景}

\begin{preamblebox}[知识领域概述与变革]
扎实的数学基础始终是计算机科学教育的基石。无论是理论还是应用领域——从算法设计、体系结构开发、物理世界建模到海量数据智能处理，数学形式化推理与数理分析均发挥着核心支撑作用。

作为 ACM CS2013 大纲中“离散结构”的承继与升级版本，本知识领域（MSF）旨在系统界定支撑现代计算机科学所必需的现代数学与统计学知识矩阵。

\textbf{自 CS2013 以来的核心变革：}
\begin{enumerate}
    \item \textbf{涵盖范围从离散数学走向现代连续数学}：CS2013 仅聚焦离散结构。随着人工智能、机器学习、机器人学、数据科学以及量子计算的爆发式增长，连续数学（微积分、多元优化、线性代数、连续概率与数理统计）已成为现代计算不可分割的基础。
    \item \textbf{全球师生与产业界双向调研驱动}：针对近 600 位高校计算机教师与 680 位工业界一线研发人员的深入调研显示，行业对微积分基础、线性代数、概率论与统计学的重视程度跃升至前五位。
    \item \textbf{正视学生数理准备度与课程衔接矛盾}：明确将初等代数与函数预备数学设为先修准入基准，倡导在计算机专业课中紧密结合具体编程案例来驱动数学概念传授，化解传统纯理论数学学习的脱节感。
\end{enumerate}
\end{preamblebox}

\section{核心学时分配建议}

\begin{table}[H]
\centering
\setlength{\tabcolsep}{3.5pt}\fontsize{8pt}{9.6pt}\selectfont\begin{tabularx}{\columnwidth}{Y c c}
\toprule
\textbf{知识单元} & \textbf{计算机核心} & \textbf{领域核心} \\
\midrule
集合、关系与函数 & 4 & -- \\
数理逻辑基础 & 9 & -- \\
形式化证明技术 & 10 & 1 \\
组合计数基础 & 5 & -- \\
图与树的数学理论 & -- & 1 \\
离散概率论 & 6 & 2 \\
微积分与连续数学基础 & \emph{贯穿渗透} & 领域核心支撑 \\
线性代数及其计算应用 & \emph{贯穿渗透} & 领域核心支撑 \\
应用统计学基础 & \emph{贯穿渗透} & 领域核心支撑 \\
\midrule
\textbf{合计} & \textbf{34 学时} & \textbf{4+ 学时} \\
\bottomrule
\end{tabularx}
\end{table}

\section{知识单元与详细大纲}

\subsection{微积分预备基础（先修准入要求）}
本模块作为计算机专业入学数学素养基线，不计入大学核心学时，但属于必修前置依赖：
\begin{itemize}
    \item \textbf{代数基础}：分式化简、幂与指数运算法则、单变量/双变量一元二次方程组求解。
    \item \textbf{函数分析}：函数形式化映射符号、单调性与函数图像性质。
    \item \textbf{指数与对数}：自然指数、常用对数、对数换底公式及其函数渐变图谱。
    \item \textbf{几何与三角学}：欧几里得空间距离度量、几何多边形特征、基本三角函数与单位圆。
\end{itemize}

\subsection{集合、关系与函数}
\textbf{【学时建议】}：4 计算机核心学时。

\begin{itemize}
    \item \textbf{集合论基础}：韦恩图、集合交/并/补/差运算、笛卡尔积、幂集、有限集合基数。
    \item \textbf{二元关系}：自反性、对称性、传递性、反对称性；等价关系与等价类划分；偏序关系与全序关系。
    \item \textbf{函数映射理论}：单射、满射、双射；反函数存在条件与函数复合。
\end{itemize}

\subsubsection{示范性学习成果}
\begin{enumerate}
    \item \textbf{[解释]} 结合具体计算问题（如关系数据库键定义、哈希表映射），阐明集合、偏序关系与双射函数的形式化概念。
    \item \textbf{[应用]} 熟练进行集合复合运算与关系性质验证，将真实系统的状态约束建模为等价类与偏序关系。
\end{enumerate}

\subsection{数理逻辑基础}
\textbf{【学时建议】}：9 计算机核心学时。

\begin{itemize}
    \item \textbf{命题逻辑}：逻辑联结词（$\neg, \land, \lor, \to, \leftrightarrow$）、真值表计算、析取范式 (DNF) 与合取范式 (CNF) 转换、公式有效性判定。
    \item \textbf{命题推理规则}：假言推理、拒取式、消解原理。
    \item \textbf{一阶谓词逻辑}：个体词、谓词、全称量词 ($\forall$) 与存在量词 ($\exists$)、约束变元与自由变元、前束范式。
    \item \textbf{形式逻辑的表达力极限}：命题逻辑与一阶逻辑的可判定性边界、表达力局限与不完备性概念。
\end{itemize}

\subsection{形式化证明技术}
\textbf{【学时建议】}：10 计算机核心学时，1 领域核心学时。

\begin{itemize}
    \item \textbf{逻辑命题形态}：充分条件、必要条件、逆命题、否命题、逆否命题及其等价证明。
    \item \textbf{直接证明与反例证伪 (Direct Proof \& Counterexamples)}：通用构造与奇偶性/整除性直接推演；通过最小反例推翻猜想。
    \item \textbf{反证法}：归谬推演与素数无穷性反证经典范例。
    \item \textbf{数学归纳法}：
    \begin{itemize}
        \item 自然数第一归纳法（弱归纳）与第二归纳法（强归纳）。
        \item 结构归纳法：在递归定义数据结构（二叉树、表达式文法）上的性质验证。
    \end{itemize}
    \item \textbf{良序原理}：良序公理及其与超限归纳法、循环终止性证明的对应关系。
\end{itemize}

\subsection{组合计数基础}
\textbf{【学时建议】}：5 计算机核心学时。

\begin{itemize}
    \item \textbf{基本计数原理}：加法原理、乘法原理、容斥原理、等差与等比级数求和。
    \item \textbf{鸽巢原理}：经典形态与广义鸽巢原理在哈希冲突与图存在性中的严格证明。
    \item \textbf{排列与组合}：二项式定理、帕斯卡恒等式（杨辉三角代数性质）、多重集合组合。
    \item \textbf{递推关系式求解}：斐波那契数列生成函数、线性齐次与非齐次递推方程特征根法求解。
    \item \textbf{初等同余与模算术}：同余类运算、欧拉定理、中国剩余定理基础。
\end{itemize}

\subsection{离散概率论}
\textbf{【学时建议】}：6 计算机核心学时，2 领域核心学时。

\begin{itemize}
    \item \textbf{概率空间与公理化度量}：样本空间、随机事件、柯尔莫哥洛夫概率公理。
    \item \textbf{条件概率与贝叶斯法则}：全概率公式、贝叶斯定理（Bayes' Theorem）及其在垃圾邮件过滤与分类决策中的应用。
    \item \textbf{独立性与随机变量}：事件独立性、离散型随机变量分布律（伯努利分布、二项分布、几何分布）。
    \item \textbf{期望与方差}：数学期望的线性性质（Linearity of Expectation，用于快速排序均摊复杂度推导）、方差与协方差、条件独立性。
\end{itemize}

\subsection{微积分、线性代数与统计学拓展}

\subsubsection{微积分}
极限与连续性；单变量微积分导数与链式法则；定积分与不定积分（在连续概率密度函数中的应用）；泰勒级数展开；多元微积分：偏导数、梯度向量、Hessian 矩阵、多元链式法则、凸优化理论与凸集判定、局部与全局极小值。

\subsubsection{线性代数}
矩阵与向量空间：矩阵线性方程组高斯消元法；向量组线性相关与线性无关、基与维数；正交投影、Gram-Schmidt 正交化、最小二乘法；特征值与特征向量、谱分解；计算机应用前沿：主成分分析 (PCA)、奇异值分解 (SVD)、PageRank 随机游走马尔可夫链稳态向量、三维图形学齐次坐标变换。

\subsubsection{应用数理统计学}
总体与样本、中心极限定理；点估计与区间估计（置信区间）；假设检验（$t$ 检验、卡方检验）；多元数据相关性分析与最小二乘回归；数据降维与平滑技术。

\section{专业素养与思维倾向}

\begin{itemize}
    \item \textbf{一丝不苟的严谨态度}：在符号演算、边界条件界定与逻辑闭环构建中保持高度审慎。
    \item \textbf{执着钻研的毅力}：面对抽象高深的数学定理与多步复杂推演，展现克服畏难情绪的坚韧品质。
\end{itemize}

\section{课程设计与教学实施建议}

联合工作组建议各高校结合办学定位采取弹性课程组织策略：
\begin{enumerate}
    \item \textbf{经典开课模式}：由数学系开设微积分、线性代数、概率统计与离散数学四门标准课程（共约 16 学分）。
    \item \textbf{现代“连续计算结构”整合课模式}：面向 AI 与数据时代，专门为计算机专业开设《连续数学基础》，融微积分梯度、多元凸优化、矩阵分解与统计学习于一体，通过 Python/NumPy 编程实战演示数学应用。
    \item \textbf{专业课内嵌式“即时教学”模式}：在机器学习、计算机视觉或密码学课程的前两周专门开设数学微模块，消除学生对纯理论数学脱离应用的认知障碍。
\end{enumerate}

\section{委员会专家名单}
\begin{itemize}
    \item \textbf{主席}：Rahul Simha (George Washington University)
    \item \textbf{委员}：Richard L. Blumenthal (Regis University)、Marc Deisenroth (University College London)、Michael Goldweber (Xavier University)、Cynthia Bailey Lee (Stanford University)、David Liben-Nowell (Carleton College)、Jodi Tims (Northeastern University)
\end{itemize}

\chapter{网络与通信}

\section{领域导言与演进背景}

\begin{preamblebox}[知识领域概述与变革]
计算机网络与现代通信技术是全球互联计算系统的核心中枢，深刻重构了数十亿人的生产与生活方式。从支撑数以亿计终端交互的公共互联网，到跨越全球版图的高性能专有云数据中心网络、内容分发网络（CDN）、近地低轨卫星宽带通信网，以及支撑物联网的近场终端自组织通信，现代网络架构正在朝着空天地海一体化、超低时延与智能化方向演进。

\textbf{自 CS2013 以来的核心变革：}
\begin{enumerate}
    \item \textbf{核心层（计算机核心）范围系统性强化}：突破以往仅聚焦网络导论与上层应用的局限，将可靠传输协议机制、路由寻址与单跳信道通信的核心原理纳为必修基础。
    \item \textbf{网络安全成为核心选修专章}：鉴于网络攻防与威胁态势的白热化，新增“网络安全”专有单元，覆盖零信任架构（ZTNA）、BGPsec 安全路由与加密信道。
    \item \textbf{以“前沿新兴议题”取代陈旧模块}：移除过时的社交网络浅层单元，增设“前沿新兴主题（Emerging Topics）”单元，全面涵盖软件定义网络（SDN）、数据中心网络虚拟化、网络中间件以及量子通信网络。
\end{enumerate}
\end{preamblebox}

\section{核心学时分配建议}

\begin{table}[H]
\centering
\setlength{\tabcolsep}{3.5pt}\fontsize{8pt}{9.6pt}\selectfont\begin{tabularx}{\columnwidth}{Y c c}
\toprule
\textbf{知识单元} & \textbf{计算机核心} & \textbf{领域核心} \\
\midrule
网络导论与基本原理 & 3 & -- \\
网络化应用与套接字编程 & 4 & -- \\
可靠传输保障机制 & -- & 6 \\
路由算法与分组转发机制 & -- & 4 \\
单跳信道与局域通信 & -- & 3 \\
移动通信与无线自组网 & -- & 4 \\
现代网络空间安全 & -- & 3 \\
网络新兴前沿与量子网络 & -- & 4 \\
\midrule
\textbf{合计} & \textbf{7 学时} & \textbf{24 学时} \\
\bottomrule
\end{tabularx}
\end{table}

\section{知识单元与详细大纲}

\subsection{网络导论与系统架构}
\textbf{【学时建议】}：3 计算机核心学时。

\begin{itemize}
    \item \textbf{现代计算中的网络地位与核心挑战}：全球异构互联、物理信道噪声、时延抖动与拥塞雪崩。
    \item \textbf{互联网组织分层架构}：最终用户、互联网服务提供商（ISP，Tier~1/2/3 架构）、对等互联、自治系统（AS）、内容提供商与 CDN 节点分布。
    \item \textbf{交换技术本质辨析}：电路交换与分组交换/包交换在统计复用、资源预留与传输延迟上的核心权衡。
    \item \textbf{五层与七层协议模型}：物理层、数据链路层、网络层、传输层与应用层的职责边界；协议数据单元（PDU）封装与解封装流程。
    \item \textbf{经典沙漏模型}：以 IP 协议为狭腰的核心设计哲学（End-to-End Principle 端到端原则）。
    \item \textbf{网络基础设施核心网元}：路由器、交换机、集线器、无线接入点（AP）与终端主机。
    \item \textbf{排队论与队列延迟直觉}：Little 定律直觉、排队延迟、丢包率与网络拥塞状态关联。
\end{itemize}

\subsection{网络化应用与套接字编程}
\textbf{【学时建议】}：4 计算机核心学时。

\begin{itemize}
    \item \textbf{网络命名与编址体系}：域名系统（DNS，层次递归与迭代解析根服务器体系）、IP 地址分配与统一资源标识符（URI/URL）。
    \item \textbf{分布式应用范式演进}：经典客户端/服务器（C/S）、对等计算（P2P）、云计算、边缘计算与雾计算。
    \item \textbf{典型应用层协议剖析}：HTTP/1.1、HTTP/2、HTTP/3（基于 QUIC）、SMTP/POP3/IMAP 邮件协议。
    \item \textbf{网络传输抽象与 Socket 编程接口}：流式套接字与数据报套接字、非阻塞 I/O 与并发网络服务器模型。
\end{itemize}

\subsection{可靠传输保障机制}
\textbf{【学时建议】}：6 领域核心学时。

\begin{itemize}
    \item \textbf{不可靠底座与尽力而为模型}：UDP 报文格式及其在实时音视频、DNS 中的适用性。
    \item \textbf{可靠数据传输协议原理}：差错检验（校验和 Checksum）、自动重传请求（ARQ）、停等协议、流水线协议（滑动窗口、回退 N 步 GBN、选择重传 SR）。
    \item \textbf{TCP 协议核心机制全解}：三次握手与四次挥手状态机演进、SYN Flood 防范、时间戳与 TIME\_WAIT 状态管理。
    \item \textbf{流量控制与拥塞控制 (Flow \& Congestion Control)}：
    \begin{itemize}
        \item 接收窗口 rwnd 与滑动窗口流量控制。
        \item 拥塞窗口 cwnd 与经典拥塞控制算法：慢启动、拥塞避免、快重传与快恢复（TCP Reno、Cubic）。
        \item 基于时延与带宽的拥塞控制（BBR 算法）。
        \item QUIC 协议在现代 Web 传输中的革新（解决队头阻塞 Head-of-Line Blocking）。
    \end{itemize}
\end{itemize}

\subsection{路由寻址与分组转发}
\textbf{【学时建议】}：4 领域核心学时。

\begin{itemize}
    \item \textbf{控制平面与数据平面解耦}：转发表生成与硬件流水线匹配。
    \item \textbf{分组转发机制}：最长前缀匹配（LPM）、三态内容寻址内存（TCAM）线速匹配。
    \item \textbf{域内路由协议}：距离向量算法（Bellman-Ford、RIP）、链路状态算法（Dijkstra、OSPF）、分层 OSPF 区域规划。
    \item \textbf{域间路由协议}：边界网关协议（BGP-4）、AS 路径属性、商业策略路由与路由振荡防范。
    \item \textbf{IP 地址扩展与过渡技术}：无类别域间路由（CIDR）、网络地址转换（NAT/NAPT）、IPv6 报头设计与双栈/隧道过渡机制。
\end{itemize}

\subsection{单跳信道与局域通信}
\textbf{【学时建议】}：3 领域核心学时。

\begin{itemize}
    \item \textbf{物理层基础}：信道带宽、奈奎斯特准则、香农定理、模拟调制与数字基带编码（曼彻斯特编码等）。
    \item \textbf{成帧与差错控制}：帧同步定界、位填充法、循环冗余校验（CRC-32）数学推导。
    \item \textbf{介质访问控制 (MAC)}：信道划分（FDMA、TDMA、CDMA）、随机访问（CSMA/CD、CSMA/CA 碰撞避免机制与指数退避算法）。
    \item \textbf{现代以太网与二层交换}：以太网帧结构、MAC 地址学习与转发表自学习、生成树协议（STP）消除二层环路、虚拟局域网（VLAN 802.1Q）。
    \item \textbf{现代数据中心网络拓扑}：Fat-Tree（胖树拓扑）、Spine-Leaf（脊叶架构）与无阻塞 ECMP 等价多路径路由。
\end{itemize}

\subsection{现代网络空间安全}
\textbf{【学时建议】}：3 领域核心学时。

\begin{itemize}
    \item \textbf{网络威胁全景}：DDoS 分布式拒绝服务攻击、ARP 欺骗、DNS 投毒与劫持、BGP 路由劫持、中间人攻击（MITM）、流量嗅探与重放攻击。
    \item \textbf{端到端与信道加密防御}：TLS 1.3 握手优化与证书信任链、IPsec 隧道/传输模式、虚拟专用网（VPN）。
    \item \textbf{网络纵深防御架构}：防火墙（状态检测防火墙、下一代防火墙 NGFW）、DMZ 隔离区配置、入侵检测与防御系统（IDS/IPS）、零信任网络访问（ZTNA）。
\end{itemize}

\subsection{移动计算与无线自组网}
\textbf{【学时建议】}：4 领域核心学时。

\begin{itemize}
    \item \textbf{蜂窝网络架构演进}：4G LTE 与 5G NR 核心网与无线接入网（RAN）解耦、网络切片。
    \item \textbf{无线局域网}：IEEE 802.11 协议族（Wi-Fi 6/6E/7）、信标帧、关联握手、MIMO 与波束成形技术。
    \item \textbf{多跳无线与自组网}：移动 Ad-hoc 网络（MANET）、车载自组网（VANET）、容迟网络（DTN）。
\end{itemize}

\subsection{前沿新兴议题}
\textbf{【学时建议】}：4 领域核心学时。

\begin{itemize}
    \item \textbf{网络中间件深入}：深度报文检测（DPI）、硬件负载均衡（L4/L7）、透明代理与内容分发智能缓存。
    \item \textbf{网络可编程与虚拟化}：软件定义网络（SDN，OpenFlow 与 P4 语言编程）、网络功能虚拟化（NFV）。
    \item \textbf{量子网络初探}：量子隐形传态、量子密钥分发（QKD，如 BB84 协议）、量子中继与量子互联网愿景。
\end{itemize}

\section{专业素养与思维倾向}

\begin{itemize}
    \item \textbf{全局工程严谨性}：在网络拓扑设计与协议参数调优中洞察端到端各层级性能关联。
    \item \textbf{敏锐的前瞻洞察力}：预判网络流量激增与突发拥塞，前瞻性部署弹性扩缩容与安全防御策略。
    \item \textbf{协同协作能力}：与分布式系统、应用开发及运维安全团队紧密跨界协作。
\end{itemize}

\section{课程设计与教学实施建议}

\begin{enumerate}
    \item \textbf{导论型计算机网络课 (Introductory Course, 45 学时)}：覆盖网络导论（9h）、网络应用与 Socket 编程（12h）、传输层与 TCP（6h）、IP 路由（4h）、以太网链路（3h）、无线与安全入门。侧重网络应用开发与基本抓包实战（Wireshark）。
    \item \textbf{高级进阶网络系统课 (Advanced Course, 45 学时)}：深入 BGP 与大规模路由（6h）、TCP 拥塞算法与 QUIC（8h）、数据中心 Spine-Leaf 与交换（5h）、移动蜂窝切片（5h）、零信任与网络攻防（5h）、SDN/P4 可编程数据平面（6h）。侧重协议设计与网络模拟实验（NS-3/Mininet）。
\end{enumerate}

\section{真实能力任务规范}

\begin{taskbox}[任务 1：企业混合云多分支网络架构设计白皮书]
\textbf{【任务情境】}：某跨国企业需在多地域分支机构、本地私有数据中心与公有云之间搭建安全、高可用且低时延的互联架构。\\
\textbf{【能力陈述】}：梳理业务与安全边界需求，设计融合 SD-WAN、BGP 多归属选路、IPsec VPN 隧道与零信任访问机制的企业级网络方案，撰写技术白皮书。\\
\textbf{【能力领域】}：系统架构 \quad \textbf{【能力单元】}：需求分析与方案设计\\
\textbf{【达标等级】}：\textbf{评估与开发}
\end{taskbox}

\begin{taskbox}[任务 2：高并发微服务场景下的自研轻量可靠传输协议]
\textbf{【任务情境】}：在跨机房弱网环境下开发高性能实时信令分发服务，原生 TCP 存在严重的队头阻塞，UDP 无法保障有序可靠性。\\
\textbf{【能力陈述】}：基于 UDP 设计并实现一套轻量级可靠传输协议，融合滑动窗口重传、前向纠错码（FEC）与拥塞反馈调优机制，并通过网络丢包模拟测试验证吞吐性能。\\
\textbf{【能力领域】}：系统与应用 \quad \textbf{【能力单元】}：协议研发与性能评测\\
\textbf{【达标等级】}：\textbf{开发}
\end{taskbox}

\begin{taskbox}[任务 3：数据中心大规模网络灰度故障与抖动诊断分析]
\textbf{【任务情境】}：生产环境数据中心遭遇隐蔽的“灰色故障”——交换机偶发微突发丢包未触发硬件告警，但导致分布式存储吞吐断崖式下跌。\\
\textbf{【能力陈述】}：部署端到端主动探测探针、收集 INT（带内网络遥测）报文流，多维度关联分析吞吐量、RTT 抖动与队列深度，准确定位故障物理交换机与受损链路。\\
\textbf{【能力领域】}：系统工程 \quad \textbf{【能力单元】}：系统监控与诊断评估\\
\textbf{【达标等级】}：\textbf{评估}
\end{taskbox}

\clearpage
\onecolumn
\section{附录：核心主题与技能要求等级对照表}

\noindent\begin{tabularx}{\textwidth}{p{0.18\textwidth} Y c c c}
\toprule
\textbf{知识单元} & \textbf{核心知识点与具体主题} & \textbf{认知等级} & \textbf{分类} & \textbf{学时} \\
\midrule
\multirow{2}{*}{\makecell[l]{网络通信/网络导论}} 
& 分组交换与电路交换本质辨析、五层模型与沙漏原则 & 解释 & 计算机核心 & 2 \\
& 自治系统 AS、ISP 架构与排队论核心概念 & 解释 & 计算机核心 & 1 \\
\midrule
\multirow{2}{*}{\makecell[l]{网络通信/网络化应用}} 
& DNS 体系结构与解析机制、HTTP/2/3 协议规范 & 解释 & 计算机核心 & 2 \\
& 套接字 Socket 编程接口、C/S 与 P2P 并发应用开发 & 应用 & 计算机核心 & 2 \\
\midrule
\multirow{3}{*}{\makecell[l]{网络通信/可靠传输}} 
& rdt 原理、滑动窗口协议、GBN 与选择重传算法 & 解释 & 领域核心 & 2 \\
& TCP 握手/挥手、状态机演变与重传定时器调优 & 应用 & 领域核心 & 2 \\
& 拥塞控制原理与算法（Reno、Cubic、BBR）分析 & 评估 & 领域核心 & 2 \\
\midrule
\multirow{2}{*}{\makecell[l]{网络通信/路由与转发}} 
& 最长前缀匹配算法、CIDR 与 NAT 转换机制 & 应用 & 领域核心 & 2 \\
& 域内（OSPF）与域间（BGP-4）路由协议与策略控制 & 评估 & 领域核心 & 2 \\
\midrule
\multirow{2}{*}{\makecell[l]{网络通信/单跳通信}} 
& 物理层香农定理、CSMA/CD 与 CSMA/CA 介质访问控制 & 解释 & 领域核心 & 2 \\
& 以太网交换机自学习算法、VLAN 划分与数据中心胖树结构 & 应用 & 领域核心 & 1 \\
\midrule
\makecell[l]{网络通信/网络安全} & DDoS 攻击防御、TLS 1.3 握手认证、零信任架构配置 & 评估 & 领域核心 & 3 \\
\midrule
\makecell[l]{网络通信/前沿新兴} & 软件定义网络 SDN/P4 编程、网络切片与量子通信网络 & 解释 & 领域核心 & 4 \\
\bottomrule
\end{tabularx}
\clearpage
\twocolumn

\chapter{操作系统}

\section{领域导言与演进背景}

\begin{preamblebox}[知识领域概述与变革]
操作系统 (OS)是向下安全接管底层物理硬件、向上为各类应用程序提供透明高效运行环境的核心系统软件集合。本知识领域聚焦于对计算处理、主存储器以及输入输出（I/O）外设进行抽象与虚拟化所需的软硬件协同机制与管理策略。在整个计算机系统各层级中被反复运用的顶层设计理念与方法论，在操作系统中得到了最为集中和深刻的体现（例如：轮询与中断驱动的抉择、多级分层缓存体系、通用灵活性与系统开销的内在权衡、贯穿于任务调度与页面置换算法中的对称设计思想等）。

一名合格的计算机科学专业毕业生，必须构建起自底向上的完整系统心智模型：从底层硬件流水线指令的微观执行，直至程序数据作用域在物理和虚拟内存空间中的精准布局。掌握操作系统核心原理，使学生能够将扎实的系统级直觉迁移应用至广阔的领域专用架构（DSA，如面向机器学习与大规模模型训练的 GPU 并行计算系统、高能效移动终端平台、安全攸关的高可靠嵌入式系统等）。由于现代所有上层软件都必须依托操作系统内核提供的服务运行，理解操作系统底层逻辑能够使学生在构建任何复杂应用或编写工程代码时，对系统吞吐效率、运行时额外开销以及各层级性能权衡进行严密推演。

对操作系统经典算法与系统机制的系统性研习，为学生评估和探索前沿软硬件技术构筑了坚实的认知基石。如果缺乏对进程沙箱隔离、可执行文件装载与动态链接、内核态特权执行边界的透彻理解，学生在面对底层安全漏洞挖掘、攻击向量防范以及高可靠防御架构设计时将陷入极为被动的境地。

\textbf{自 CS2013 以来的核心变革：}
\begin{enumerate}
    \item \textbf{文件系统与设备管理回归核心课程}：将原属于选修范畴的“文件系统 API 与底层实现”以及“设备管理”全面提升为计算机专业必修的计算机核心知识单元。这一调整深刻反映了现代数据密集型与并行计算应用中，持久化存储与外设 I/O 吞吐瓶颈对算法宏观性能的决定性制约。
    \item \textbf{核心学时平衡与高阶主题分流}：为将文件系统实现与设备管理平稳纳为 计算机核心 必修学时，将部分原有更为高阶的专业细分主题由 计算机核心 调整至领域核心，从而理顺由浅入深、梯次递进的知识架构。
    \item \textbf{“系统性能与评估”单元重组外迁}：将原操作系统知识体系中的“性能与评估”知识单元整合移交至《系统基础 (SF)》知识领域，以彰显性能建模与量化评测方法在硬件、操作系统、网络通信等多系统层级上的横向对称性与普适性。
    \item \textbf{现代虚拟化、容器与系统安全深度融合}：全面强化对现代系统虚拟化（Type-1/2 Hypervisor）、操作系统级轻量容器（Linux Namespaces、cgroups）、新型内核形态（微内核、单一内核/独核 Unikernel）、现代可执行文件布局（ELF、GOT/PLT）以及开源操作系统生态与生命周期维护的系统化讲授。
\end{enumerate}
\end{preamblebox}

\section{核心学时分配建议}

本知识领域在“向日葵模型”下的核心学时推荐如下表所示：

\begin{table}[H]
\centering
\setlength{\tabcolsep}{3.5pt}\fontsize{8pt}{9.6pt}\selectfont\begin{tabularx}{\columnwidth}{Y c c}
\toprule
\textbf{知识单元} & \textbf{计算机核心} & \textbf{领域核心} \\
\midrule
操作系统的角色与目标 & 2 & -- \\
操作系统原理与体系架构 & 2 & -- \\
并发与同步机制 & 2 & -- \\
处理器调度与资源分配 & 1 & -- \\
进程模型与执行环境 & 1 & -- \\
内存管理与虚拟内存 & 2 & -- \\
安全保护与故障隔离 & 2 & -- \\
设备管理与输入输出系统 & 1 & -- \\
文件系统 API 与底层实现 & 2 & -- \\
系统虚拟化与容器技术 & -- & 3 \\
实时与嵌入式系统 & -- & 2 \\
系统容错与高可用性 & -- & 3 \\
\midrule
\textbf{合计} & \textbf{15 学时} & \textbf{8 学时} \\
\bottomrule
\end{tabularx}
\end{table}

\noindent\textbf{说明}：上表列出了 CS2023 委员会界定的最低核心必修学时基准。在高校实际教学计划中，建议结合高级文件系统、开源社区生态与工程实验实操，开设约 37$\sim$48 学时的完整专业主干课程（详见第 8 节课程设计与教学实施建议）。

\section{知识单元与详细大纲}

\subsection{操作系统的角色与目标}
\textbf{【学时建议】}：2 计算机核心学时。

\subsubsection{计算机核心主题}
\begin{itemize}
    \item \textbf{操作系统作为软硬件中介}：操作系统作为通用底层物理硬件与特定领域应用软件之间的中介抽象层；通过标准应用程序接口（API）向上层软件提供统一的抽象虚拟机模型。
    \item \textbf{操作系统通用核心职能}：
    \begin{itemize}
        \item 交互与通信接口抽象（进程接口、用户命令行/图形界面、设备抽象文件接口等）。
        \item 面向特定应用软件的代码加载、地址空间创建与受控运行生命周期管理。
        \item 数据持久化存储管理与状态安全保障。
    \end{itemize}
    \item \textbf{扩展与专用操作系统职能}：
    \begin{itemize}
        \item 资源严苛环境下的嵌入式操作系统特殊机制。
        \item 面向高吞吐、高并发场景的专用服务器系统（文件存储服务器、Web 代理服务器、流媒体低时延服务器）。
        \item 系统引导加载程序与安全引导信任链机制。
    \end{itemize}
    \item \textbf{系统级设计权衡与冲突分析}：
    \begin{itemize}
        \item 效率、健壮性、灵活性、可移植性、安全性、向后兼容性、能耗预算以及人身与物理安全。
        \item 关键工程维度的两难权衡：深层边界错误检查与极致执行性能的冲突；高度模块化灵活性与执行性能开销的冲突；多层动态安全校验与系统访问时延的冲突。
    \end{itemize}
    \item \textbf{外部关键技术对操作系统的深远影响}：现代网络空间安全威胁、分布式与云计算架构、高通量多媒体流式计算、多核并行与异构计算协同。
    \item \textbf{安全与保护的全局贯穿性法则}：在操作系统软件栈的任何一个分层忽视安全，都将直接导致整机安全基准坍塌，从而为非法越权访问硬件资源洞开大门：
    \begin{itemize}
        \item 未加密硬盘介质在脱机移至其他物理主机后可被完全绕过权限直接读取。
        \item 在操作系统内核载入内存前，恶意固件代码渗透引导加载层可彻底颠覆内核安全体系。
        \item 系统 API 边界若缺乏严格细密的授权校验，将导致不同应用间的进程沙箱隔离失效。
    \end{itemize}
\end{itemize}

\subsubsection{示范性学习成果与能力指标}
\begin{itemize}
    \item \textbf{[解释] 理解目标与功能}：深刻阐明现代操作系统的根本设计目标及其支撑应用运转的三大核心虚拟化职能（计算、内存与外设）。
    \item \textbf{[评估] 多元场景架构权衡}：针对实时控制、智能手机移动端、云端高并发服务器等不同业务形态，严密权衡并评估操作系统在调度、内存管理与 I/O 上的设计分歧。
    \item \textbf{[解释] 便利性、效率与演进力}：全面分析当代操作系统在兼顾最终用户与开发者的交互便利性、底层硬件吞吐效率以及架构平滑演进能力上的实现路径。
    \item \textbf{[解释] 演进性与稳定性的对立统一}：深刻剖析操作系统在持续接纳新技术（演进性）与保持向后兼容及系统内核稳定（稳定性）之间既相互追求又内在对立的辩证矛盾。
\end{itemize}

\subsection{操作系统原理与体系架构}
\textbf{【学时建议】}：2 计算机核心学时。

\subsubsection{计算机核心主题}
\begin{itemize}
    \item \textbf{操作系统软件架构范式演进}：
    \begin{itemize}
        \item 宏内核：高性能、紧密集成与模块间缺乏故障隔离边界的缺陷。
        \item 分层架构：各层自底向上单向依赖的高内聚与调用性能折损。
        \item 模块化动态加载内核：在保持宏内核高性能的同时实现按需动态扩展。
        \item 微内核：将文件系统、设备驱动等非关键服务外迁至用户态，依托进程间通信（IPC）实现高安全与高健壮性。
        \item 单一内核 / 独核架构：将应用程序与精简内核库编译为单一地址空间镜像，专为云原生虚拟化容器提供极致轻量化与极小攻击面。
    \end{itemize}
    \item \textbf{核心系统级三大抽象}：进程（计算实体抽象）、地址空间与虚拟内存（存储资源抽象）、文件与 I/O 流（外设与通信管道抽象）。
    \item \textbf{系统调用与应用程序接口（API）}：
    \begin{itemize}
        \item 系统调用的本质：作为用户空间程序请求内核特权服务的受控入口。
        \item 高层任务通常需要内核底层多组细粒度系统调用的级联配合完成。
        \item 标准运行时 API（如 Win32、POSIX、Java Native Interface）的作用：屏蔽底层系统调用的繁琐细节与平台异构性，向上提供统一优雅的高级编程接口。
        \item 系统调用二进制接口设计哲学：Linux 的“永远不得破坏用户空间”永久稳定 ABI 原则与 OpenBSD/macOS 每次主要版本升级均可调整重构 ABI 的流派对比。
    \end{itemize}
    \item \textbf{硬件体系结构与操作系统协同演进的历史脉络}：双模式硬件支持、存储保护键、虚拟化扩展指令集对 OS 架构的深远驱动。
    \item \textbf{特权指令受控执行与资源保护}：
    \begin{itemize}
        \item 硬件强制保护某些特权机器指令（如修改页表寄存器、停机指令、屏蔽中断）。
        \item 严格禁止用户态应用任意访问非法物理内存地址或直接读写磁盘块物理扇区。
        \item 对协处理器、GPU、密码学加速芯片及高速网络接口卡进行受控保护与资源复用。
    \end{itemize}
    \item \textbf{硬件中断与异常处理机制}：
    \begin{itemize}
        \item 中断向量表（IVT）与中断服务例程（ISR）的硬件分发逻辑。
        \item 硬件可编程定时器中断：为抢占式调度器提供毫秒级时间片基准。
        \item 外设 I/O 中断：实现异步事件驱动，将等待数据的阻塞线程置入休眠等待队列，彻底消除 CPU 无谓轮询的能源与算力损耗。
    \end{itemize}
    \item \textbf{双模式运行与特权级跃迁}：用户态与内核态的硬件标志位界定；用户程序通过陷阱或系统调用指令安全陷入内核。
    \item \textbf{系统调用全周期动态机制}：系统调用号分发、用户态到内核态的特权模式切换、硬件上下文压栈保存、内核例程执行、结果寄存器返回及从中断安全返回用户态。
\end{itemize}

\subsubsection{示范性学习成果与能力指标}
\begin{itemize}
    \item \textbf{[解释] 软件架构对健壮性与可维护性的影响}：系统阐述宏内核、微内核、分层与独核等设计范式对操作系统抵御局部故障、热修补及长期工程维护的本质影响。
    \item \textbf{[解释] 系统调用分类与归纳}：熟练将系统调用按照功能维度（进程管理、文件操作、设备控制、状态查询、通信同步）进行精确分类。
    \item \textbf{[解释] 系统调用执行全流程追踪}：详述应用程序触发系统调用的完整软硬件协同过程（参数传递规范、陷阱门模式跃迁、内核栈切换、现场恢复）。
    \item \textbf{[评估] 应用层库与内核系统调用的职责边界}：针对给定的软件功能（如高精度数学计算 vs 磁盘写入缓冲），科学评估其应当直接在用户态库中实现，还是必须申请内核态系统调用完成。
    \item \textbf{[应用] 系统防护与吞吐量技术的横向迁移}：将操作系统中用于隔离、保护与防饥饿的通用机制（如进程防饥饿老化、磁盘调度电梯算法、信号量公平排队），熟练迁移应用于复杂分布式系统与网络调度设计中。
    \item \textbf{[解释] 双模式隔离的安全性与开销权衡}：深入阐明内核态与用户态硬件隔离屏障对防止整机崩溃的绝对必要性，以及频繁模式切换对高性能计算所带来的流水线清空开销。
    \item \textbf{[解释] 中断驱动并发的利弊分析}：全面论述中断机制在支撑多道程序并发运行中的关键价值，以及处理嵌套中断、中断延迟抖动带来的复杂性代价。
    \item \textbf{[评估] 底层安全威胁与防御特性审计}：系统分析操作系统内核可能遭遇的越权利用、内核栈溢出等攻击手法，并评估内核页表隔离（KPTI）、地址空间布局随机化（ASLR）等防御机制的有效性。
\end{itemize}

\subsection{并发与同步机制}
\textbf{【学时建议】}：2 计算机核心学时，3 领域核心学时（另见附录）。

\subsubsection{计算机核心主题}
\begin{itemize}
    \item \textbf{并发视角下的线程抽象}：线程作为 CPU 调度的最小执行单元；共享所属进程的虚拟地址空间与代码段，拥有独立的程序计数器（PC）、寄存器上下文与调用栈。
    \item \textbf{竞态条件与临界区}：
    \begin{itemize}
        \item 竞态条件的成因：多执行流对共享可变状态的非原子并发访问。
        \item 临界区定义及其访问原则（互斥、有空让进、有限等待）。
        \item 屏蔽硬件中断在保护临界区时的有效性局限（仅适用于单核系统内核关键路径，严禁暴露至用户态，无法解决多核并发冲突）。
    \end{itemize}
    \item \textbf{死锁与饥饿现象}：死锁（多个执行流因互相持有对方所需资源且不释放而陷入永久阻塞）与饥饿（低优先级线程长期无法获取所需调度资源）的直观物理表现。
    \item \textbf{多处理器架构下的并发新挑战}：
    \begin{itemize}
        \item 自旋锁的适用场景与多核总线竞争开销。
        \item 内核代码的可重入性与多核内存模型乱序可见性问题。
    \end{itemize}
\end{itemize}

\subsubsection{领域核心主题}
\begin{itemize}
    \item \textbf{线程管理全生命周期与数据结构}：线程的创建、就绪、运行、阻塞、终止状态机；线程控制块（TCB）的核心字段组成。
    \item \textbf{标准线程编程接口}：POSIX 线程标准的核心 API、系统级线程生命周期管控与线程池设计。
    \item \textbf{死锁理论形式化与综合治理}：
    \begin{itemize}
        \item 死锁产生的四大充分必要条件（Coffman 条件）：互斥条件、占有且等待、不可抢占、循环等待。
        \item 死锁预防（破坏四大必要条件之一，如按序申请资源消除环路）。
        \item 死锁避免（Dijkstra 银行家算法 Banker's Algorithm、系统安全状态判定）。
        \item 死锁检测与动态恢复（资源分配图 Resource Allocation Graph 的化简、死锁进程终止与资源回滚）。
    \end{itemize}
    \item \textbf{线程安全代码工程实现}：
    \begin{itemize}
        \item 互斥锁的加锁/解锁协议。
        \item 计数信号量与二进制信号量，Dijkstra P/V（wait/signal）操作规范。
        \item 条件变量机制（等待 wait、单播唤醒 signal、广播唤醒 broadcast）与规避伪唤醒的标准循环判断范式。
    \end{itemize}
    \item \textbf{共享内存环境下的竞态条件消除}：内存屏障、原子读-改-写原语在多核共享内存一致性中的关键角色。
\end{itemize}

\subsubsection{选修与前沿拓展主题}
\begin{itemize}
    \item \textbf{内核对象原子并发管理的高阶演进}：从粗粒度的大内核锁 (BKL)向细粒度锁的重构历史；免锁/无锁数据结构（Lock-free Queues、CAS 原子循环更新）；Linux 内核读-复制-更新机制在高并发只读场景下的无锁吞吐优势。
\end{itemize}

\subsubsection{示范性学习成果与能力指标}
\begin{itemize}
    \item \textbf{[解释] 并发机制的双刃剑效应}：辩证阐述并发机制作为操作系统核心支撑在显著提升算力利用率与响应性能的同时，所引入的调试困难与竞态不确定性风险。
    \item \textbf{[解释] 体系结构与并发竞态根源}：深入阐明微观体系结构特征（指令乱序发射、多级高速缓存不一致、编译器指令重排）如何直接诱发高级软件层面的竞态条件。
    \item \textbf{[解释] 多核多处理器并发挑战}：剖析在 SMP 对称多处理器环境中，多核心并发争抢共享数据结构所引发的总线风暴与自旋锁性能雪崩。
    \item \textbf{[解释] 操作系统级并发机制谱系}：系统比较自旋锁、互斥锁、读写锁、信号量、条件变量等机制在上下文切换开销与等待时延上的核心差异与选型场景。
    \item \textbf{[应用] 硬件原子指令与同步原语构建}：熟练运用底层硬件原子指令（如 CAS、Test-and-Set），在无锁环境下构建出正确的自旋锁与计数信号量原语。
    \item \textbf{[评估] 竞态条件与死锁诊断重构}：严密审查复杂的并发工程代码，精准识别隐藏的竞态条件、死锁隐患或优先级翻转缺陷，并设计出无死锁的高性能同步重构方案。
\end{itemize}

\subsection{处理器调度与资源分配}
\textbf{【学时建议】}：1 计算机核心学时，3 领域核心学时（另见附录）。

\subsubsection{领域核心主题}
\begin{itemize}
    \item \textbf{调度策略与抢占模式}：
    \begin{itemize}
        \item 非抢占式调度：进程持续运行直至主动让出 CPU 或进入阻塞，实现简单但可能导致系统无响应。
        \item 抢占式调度：基于时钟中断或高优先级就绪事件强制剥夺当前 CPU 执行权，保障系统交互式响应与公平性。
    \end{itemize}
    \item \textbf{经典调度策略与算法族}：
    \begin{itemize}
        \item 先来先服务（FCFS）：实现直观，但易受“护航效应”影响导致平均等待时间恶化。
        \item 短作业优先：理论最优的平均等待时间，但存在长作业饥饿风险且需要预测执行时间。
        \item 静态与动态优先级调度：结合任务关键等级进行调度，需引入“老化”机制规避低优先级饥饿。
        \item 时间片轮转 (RR)：时间片大小选取的工程权衡（过大退化为 FCFS，过小导致上下文切换频繁）。
        \item 多级反馈队列 (MLFQ)：动态调整任务优先级与时间片长度，兼顾短交互任务响应与长计算任务吞吐。
        \item 完全公平调度器 (CFS)：基于红黑树与虚拟运行时间度量的高性能公平调度。
    \end{itemize}
    \item \textbf{对称多处理（SMP）多核调度与缓存协同}：
    \begin{itemize}
        \item 处理器亲和性（Processor Affinity，软亲和性与硬亲和性）以最大限度复用 L1/L2 缓存。
        \item 多核负载均衡：任务推取与任务拉取对多核缓存一致性的扰动。
    \end{itemize}
    \item \textbf{高精度软件定时器虚拟化}：利用有限的物理硬件时钟中断发生器，基于最小堆或时间轮高效虚拟化出数以万计的应用层与内核层高精度定时器。
\end{itemize}

\subsubsection{选修与前沿拓展主题}
\begin{itemize}
    \item \textbf{能耗感知调度与绿色计算}：结合硬件动态电压频率调整（DVFS）与异构多核处理器拓扑（如 ARM big.LITTLE 大小核），按需迁移高算力与后台节能任务。
    \item \textbf{用户态协作调度与新型轻量化框架}：Linux futexes（快速用户空间互斥体）底层机制；Go 等现代语言运行时的 $M:N$ 协程调度与用户态调度技术。
\end{itemize}

\subsubsection{示范性学习成果与能力指标}
\begin{itemize}
    \item \textbf{[评估] 调度算法量化对比}：在吞吐量、平均周转时间、就绪等待时间以及公平份额维度上，严谨量化对比 FCFS、SJF、RR、MLFQ 等抢占与非抢占算法的性能表现。
    \item \textbf{[解释] 调度策略与业务应用匹配}：准确建立不同调度算法（批处理吞吐优先、桌面交互时延敏感、服务器公平共享）与实际应用业务特征的映射关联。
    \item \textbf{[解释] 多级调度器体系分工}：深入阐明系统中高级调度（作业准入调度）、中级调度（内存对换交换调度）、低级调度（CPU 进程分派）及 I/O 调度的协同工作机制。
    \item \textbf{[评估] 对称与非对称多处理架构适用性}：针对给定的多核工程计算场景，科学评估采用对称多处理（SMP）或非对称多处理（AMP）架构的性能边界与可行性。
    \item \textbf{[评估] 多进程与多线程架构选型}：在内存占用、故障隔离屏障、创建与销毁开销及上下文切换成本等维度上，客观评估选择多进程并发或多线程并发架构的工程得失。
    \item \textbf{[解释] 抢占与截止时间调度的必要性}：阐明在硬实时系统与防恶意拒绝服务攻击中，强行抢占机制与基于截止时间驱动调度的决定性作用。
    \item \textbf{[应用] 调度算法思想的跨界迁移}：将经典的操作系统 CPU 调度思维抽象提升，熟练迁移至网络数据包调度（如 WFQ）、磁盘块 I/O 请求调度以及企业工程项目进度规划等非计算领域中。
\end{itemize}

\subsection{进程模型与执行环境}
\textbf{【学时建议】}：1 计算机核心学时，3 领域核心学时（另见附录）。

\subsubsection{领域核心主题}
\begin{itemize}
    \item \textbf{虚拟化视角下的进程与线程}：进程作为受操作系统保护的独立虚拟地址空间与资源分配容器；进程状态（新建、就绪、运行、阻塞、终止）及其状态迁移图；基于进程控制块（PCB）维护的完整执行上下文。
    \item \textbf{进程内存镜像与段布局剖析}：
    \begin{itemize}
        \item 代码段：存放可执行机器指令，通常设置为只读与共享。
        \item 数据段：已初始化的全局变量与静态变量。
        \item BSS 段：未初始化或初始化为零的全局与静态变量，在磁盘文件中仅占位，加载时由内核清零分配物理页。
        \item 堆区：动态内存分配空间，自低地址向高地址动态生长。
        \item 栈区：局部变量、函数参数与返回地址，自高地址向低地址向下生长。
        \item 内存映射区：用于加载动态共享库与 \texttt{mmap} 匿名/文件映射。
    \end{itemize}
    \item \textbf{可执行程序的构建、装载与动态链接}：
    \begin{itemize}
        \item 源代码经过预处理、编译、汇编生成目标文件（\texttt{.o}），经静态链接器或运行时动态链接器生成进程镜像的完整流水线。
        \item 现代可执行与可链接格式（ELF）的剖析：ELF 头部、程序头表、节区头表。
        \item 动态链接机制、位置无关代码（PIC）、全局偏移表（GOT）与过程链接表（PLT）的延迟绑定实现细节。
    \end{itemize}
    \item \textbf{分派与上下文切换机制}：分派器的工作逻辑；上下文切换中寄存器现场保存到内核栈、更新 PCB 状态、刷新内存管理单元（MMU）页表基址寄存器（CR3）及其引起的 TLB 缓存失效开销。
    \item \textbf{进程间通信机制（IPC）全谱系}：
    \begin{itemize}
        \item 共享内存：最高吞吐的 IPC 形式，需配合信号量进行同步。
        \item 消息传递：系统内核提供消息队列，实现解耦异步通信。
        \item 管道技术：匿名管道（Pipes，仅适用于父子亲缘进程）与命名管道（FIFO，支持无亲缘跨进程）。
        \item 信号机制：内核向进程传递异步事件通知的标准软中断机制（如 \texttt{SIGINT}、\texttt{SIGSEGV}）。
        \item 环境变量与跨进程参数传递。
    \end{itemize}
\end{itemize}

\subsubsection{示范性学习成果与能力指标}
\begin{itemize}
    \item \textbf{[解释] 进程控制权虚拟化}：阐述操作系统如何借助进程与线程抽象以及底层时间片并发调度，实现物理硬件 CPU 控制权的透明虚拟化。
    \item \textbf{[解释] 中断与上下文切换的底层作用}：深入阐明时钟中断、陷入异常、分派器以及上下文切换在维系多任务并发多道程序环境中的核心软硬件机制。
    \item \textbf{[解释] 任务状态迁移与内核数据结构}：熟练绘制进程状态迁移图，并详述操作系统内核中 PCB/TCB 链表、运行就绪队列与阻塞等待队列的组织管理机制。
    \item \textbf{[应用] 编译、链接与程序生成实战}：熟练使用工具链（如 GCC、Clang、LD），将 C/C++ 源代码、静态归档库（\texttt{.a}）与动态共享目标（\texttt{.so}）编译链接为规范的可执行二进制文件。
    \item \textbf{[评估] 静态库与动态库的全面权衡}：在磁盘占用、主存复用率、启动加载时延、热升级维护便利性以及潜在依赖冲突等维度上，全面评估动静态链接方案的工程得失。
    \item \textbf{[解释] 任务内存分配机制对比}：阐述操作系统采用连续内存分配、分段、分页及段页式机制向进程分配物理内存的优缺点及内存碎片形成机制。
    \item \textbf{[应用] 场景驱动的高效 IPC 设计}：针对给定的软件系统协作需求（如高吞吐大文件处理 vs 低时延轻量控制指令），精准选择并编码实现最适宜的进程间通信（IPC）机制。
\end{itemize}

\subsection{内存管理与虚拟内存}
\textbf{【学时建议】}：2 计算机核心学时，4 领域核心学时（另见附录）。

\subsubsection{领域核心主题}
\begin{itemize}
    \item \textbf{物理内存与硬件地址转换机制}：物理寻址局限；内存管理单元（MMU）、基址/界限寄存器、重定位寄存器与页表基地址寄存器的硬件协同机制。
    \item \textbf{存储层次结构对操作系统策略的驱动}：
    \begin{itemize}
        \item 寄存器、L1/L2/L3 片上高速缓存、物理主存与外部持久化辅存的速度与容量金字塔。
        \item 针对现代多核系统的 CPU 亲和性调度与 Per-CPU 内存分配缓存，最大限度保证局部性并规避多核总线缓存一致性风暴。
    \end{itemize}
    \item \textbf{逻辑地址、物理地址与虚拟地址空间}：虚拟地址空间的本质；多进程各自拥有独立的平坦虚拟地址空间，通过页表实现非连续物理主存的透明映射与物理地址隔离。
    \item \textbf{分页机制与页面置换策略}：
    \begin{itemize}
        \item 虚拟页、物理页框、页表项（PTE，包含存在位、读写位、用户/内核位、脏位与访问位）。
        \item 多级页表与倒排页表；转换检测缓冲区（TLB）的高速地址转换缓存及其命中/缺页处理。
        \item 缺页中断处理全流程：从捕获硬件异常、陷入内核、自外存调入页面、更新 PTE 状态至重新执行发生缺页的指令。
        \item 典型页面置换算法：先进先出（FIFO，Belady 异常）、最优置换算法（OPT）、最近最少使用算法（LRU）、时钟置换算法（Clock / 二次机会算法）。
        \item 页面分配策略与系统抖动/颠簸：颠簸的成因（多道程序度过高导致频繁缺页与磁盘 I/O 堵塞）；基于工作集模型与缺页率控制（PFF）的颠簸根治机制。
    \end{itemize}
    \item \textbf{内核与用户态高效内存分配技术}：
    \begin{itemize}
        \item 伙伴系统：解决外部碎片的大块物理内存分配算法。
        \item Slab / Slub 分配器：针对内核频繁申请释放的小型数据对象缓存池。
        \item 空闲链表、大小分级分配器（Size Classes，如 TCMalloc、Jemalloc 思想）。
        \item 大页/巨页机制（Hugepages，2MB / 1GB 大页）：大幅降低大型内存应用的 TLB 未命中开销。
    \end{itemize}
    \item \textbf{主存缓存与多核缓存一致性}：高速缓存映射策略（直接映射、组相联）、缓存行填充与写回策略对系统调用性能的影响。
    \item \textbf{内存管理中的系统安全与隔离机制}：
    \begin{itemize}
        \item 进程沙箱与页表级严格读/写/执行（R/W/X）权限位控制。
        \item 硬件不可执行位：杜绝栈注入代码的非法执行。
        \item 地址空间布局随机化（ASLR）：随机化栈、堆和库的基地址，显著提高返回导向编程（ROP）攻击的试错门槛。
    \end{itemize}
\end{itemize}

\subsubsection{选修与前沿拓展主题}
\begin{itemize}
    \item \textbf{虚拟内存硬件的高阶服务赋能}：利用页表中断与保护位实现写时复制（Copy-on-Write, COW，如 \texttt{fork()} 的极速派生）、内存映射文件（\texttt{mmap} 零拷贝 I/O）、脏页追踪、用户态缺页异常处理接口（\texttt{userfaultfd}）在分布式虚拟机迁移与垃圾回收中的前沿应用。
\end{itemize}

\subsubsection{示范性学习成果与能力指标}
\begin{itemize}
    \item \textbf{[解释] 存储层次结构与成本-性能权衡}：详细阐明计算机存储体系各个层级的技术特征，并从系统工程角度论证容量、访问时延与制造成本之间的权衡规则。
    \item \textbf{[解释] 虚拟内存与分页机制原理}：高度总结虚拟内存结合底层分页机制实现物理主存弹性扩展、地址空间解耦与高效缓存的运行机理。
    \item \textbf{[评估] 存储容量与 CPU 性能协同权衡}：针对高负载计算环境，定量评估不同层级缓存容量与处理器主频对整体端到端计算吞吐的影响关系。
    \item \textbf{[解释] 硬件 Cache 机制与侧信道安全}：阐述多级高速缓存的设计初衷，并深入剖析微架构层缓存共享如何导致跨进程侧信道（Side-channel，如 Spectre、Meltdown 漏洞）攻击的安全边界隐患。
    \item \textbf{[开发] 局部性友好的高性能代码构建}：深入运用空间局部性与时间局部性原理，编写出兼顾页框分配与低缺页率的高效软件，彻底规避内存颠簸。
    \item \textbf{[解释] 硬件辅助内存虚拟化}：阐述底层硬件如何通过二级地址转换（SLAT，如 Intel EPT、AMD NPT）支持客户机物理地址到宿主机物理地址的高效并行翻译。
\end{itemize}

\subsection{安全保护与故障隔离}
\textbf{【学时建议】}：2 计算机核心学时，1 领域核心学时（另见附录）。

\subsubsection{计算机核心主题}
\begin{itemize}
    \item \textbf{操作系统安全机制全景概览}：操作系统作为安全信任根与参考监控器的核心地位；机密性、完整性与可用性（CIA 三元组）在系统资源调度中的具体体现。
    \item \textbf{典型对抗攻击与资源对抗}：
    \begin{itemize}
        \item 调度对抗与饥饿攻击（利用抢占与实时优先级耗尽系统算力）。
        \item 内存越权、特权提升与恶意拒绝服务。
    \end{itemize}
    \item \textbf{主流商业操作系统典型重大漏洞复盘}：栈缓冲区溢出、堆溢出、释放后使用（UAF）、整数溢出、文件系统检查-使用时间差（TOCTOU 条件竞争漏洞）。
    \item \textbf{操作系统级防御与缓解策略}：内核自防御机制、系统备份、增量镜像与崩溃转储离线取证审计。
\end{itemize}

\subsubsection{领域核心主题}
\begin{itemize}
    \item \textbf{策略与机制分离原则}：机制规定“如何做”（如内核提供支持访问控制的底层校验逻辑），策略规定“允许谁做什么”（如系统管理员配置的特定安全规则清单）。
    \item \textbf{体系结构级硬件安全设施}：
    \begin{itemize}
        \item 分级保护环：从 Multics 体系结构演化至现代 x86 的 Ring~0（内核态）、Ring~3（用户态）及支持硬件虚拟化的 Ring~-1（Hypervisor 根模式）。
    \end{itemize}
    \item \textbf{保护、访问控制与身份鉴别模型}：
    \begin{itemize}
        \item 访问控制矩阵模型。
        \item 访问控制列表（ACL）与基于权能的系统在微内核与现代操作系统中的工程落地。
        \item 多因子身份认证（MFA）在特权凭据授予中的防线作用。
    \end{itemize}
\end{itemize}

\subsubsection{示范性学习成果与能力指标}
\begin{itemize}
    \item \textbf{[解释] 安全保护机制的根本必要性}：阐明在现代多租户、多任务互联计算环境中，构建内生且不可旁路的操作系统安全保护机制的不可替代性。
    \item \textbf{[解释] 经典底层漏洞攻击向量剖析}：详细列举并深入剖析利用操作系统底层特权级漏洞进行渗透攻击的各类经典攻击向量。
    \item \textbf{[解释] 细粒度系统访问控制机制}：详述操作系统利用自主访问控制（DAC）、强制访问控制（MAC）与特权权能机制约束进程访问受限资源的具体途径。
    \item \textbf{[评估] 操作系统安全特性与防御局限审计}：客观总结操作系统的固有安全防护特性及其在防御供应链投毒、微架构侧信道等新兴高级威胁时的信任边界与局限性。
\end{itemize}

\subsection{设备管理与输入输出系统}
\textbf{【学时建议】}：1 计算机核心学时，2 领域核心学时（另见附录）。

\subsubsection{领域核心主题}
\begin{itemize}
    \item \textbf{I/O 数据缓冲技术与实现策略}：
    \begin{itemize}
        \item 缓解 CPU 极高计算速度与外设缓慢 I/O 数据传输之间的严重速度失配。
        \item 单缓冲、双缓冲（Double Buffering / Ping-Pong 缓冲）与内核循环环形缓冲区。
    \end{itemize}
    \item \textbf{设备控制机制与通信模式演进}：
    \begin{itemize}
        \item 程序直接控制（轮询 Polled I/O）：CPU 周期严重浪费。
        \item 中断驱动 I/O（Interrupt-Driven I/O）：异步事件驱动，但高频小包传输会诱发严重的中断风暴。
        \item 直接内存访问：DMA 控制器接管总线，直接在主存与设备缓冲区之间进行高速块数据传输，仅在整块传输结束时向 CPU 发出中断。
        \item 内存映射 I/O（MMIO）与端口独立编址 I/O（PMIO）的寻址差异；基于描述符环形队列的高性能 DMA 通信协议。
    \end{itemize}
    \item \textbf{持久化块存储设备介质演进与管理}：
    \begin{itemize}
        \item 机械硬盘（HDD）：磁头寻道时间、旋转延迟、扇区读写；传统磁盘调度算法（FCFS、SSTF 最短寻道优先、SCAN 电梯算法、C-SCAN 循环扫描算法）。
        \item 固态硬盘（SSD）：闪存物理特性（电荷捕获、只写一次、异地更新、整块擦除）；闪存转换层（FTL）、损耗均衡算法与操作系统 TRIM 块回收通知。
    \end{itemize}
\end{itemize}

\subsubsection{选修与前沿拓展主题}
\begin{itemize}
    \item \textbf{统一设备模型与硬件抽象层（HAL）}：设备树规范、字符设备、块设备与网络套接字设备的接口统一化抽象。
    \item \textbf{现代设备驱动程序工程挑战}：内核态驱动的设计原则、并发竞态、内核死锁隐患与驱动隔离沙箱技术（用户态驱动框架 UIO/VFIO）。
    \item \textbf{设备底层通信容错与高可用重试}：PCIe 总线链路层重传校验与错误报告机制。
\end{itemize}

\subsubsection{示范性学习成果与能力指标}
\begin{itemize}
    \item \textbf{[解释] 体系结构级设备控制与 OS 联动}：理解硬件设备控制器与接口协议（如 PCIe、NVMe）的微观运作，并将其与操作系统的 I/O 调度策略及中断例程紧密映射。
    \item \textbf{[解释] 操作系统分层 I/O 软件栈架构}：掌握操作系统设备管理分层模型（用户级 I/O 软件、设备独立性软件、设备驱动程序、中断处理程序、硬件控制器）。
    \item \textbf{[解释] 物理硬件与虚拟设备抽象映射}：深刻阐述操作系统如何将物理各异的异构硬件设备封装并抽象为标准的虚拟字符设备、块设备与套接字接口。
    \item \textbf{[解释] I/O 缓冲技术原理与选型}：详细阐述 I/O 数据缓冲的工作机理，并针对流式数据与块数据传输场景设计适宜的缓冲实现策略。
    \item \textbf{[解释] 直接内存访问（DMA）机制分析}：深入论述 DMA 传输的技术优势、总线仲裁机制，并明确界定其在高速大通量数据传输中的触发条件与适用边界。
    \item \textbf{[解释] 设备驱动程序构建规范与复杂度}：阐述编写安全可靠的设备驱动程序所面临的并发陷阱，并总结现代操作系统中设备驱动开发的最佳实践与静态验证方法。
\end{itemize}

\subsection{文件系统 API 与底层实现}
\textbf{【学时建议】}：2 计算机核心学时，2 领域核心学时（另见附录）。

\subsubsection{领域核心主题}
\begin{itemize}
    \item \textbf{文件的本质概念与核心抽象}：文件作为线性逻辑地址空间的字节序列集合；文件数据（Data 内容）与文件元数据（Metadata，如 inode 中存储的大小、类型、所有者、创建时间、访问权限及数据块指针）；文件基本操作接口（\texttt{open}、\texttt{read}、\texttt{write}、\texttt{seek}、\texttt{close}）与打开模式（只读、只写、追加、独占创建）。
    \item \textbf{文件系统的挂载机制}：将独立存储设备上的文件系统层次结构接入操作系统的统一全局目录树，形成无缝文件命名空间。
    \item \textbf{文件访问权限控制}：POSIX 标准权限位（所有者、所属组、其他人的读/写/执行权限）、访问控制列表（ACL）在文件层级的细粒度授权。
    \item \textbf{文件共享与多进程并发安全}：硬链接（Hard Links，共享相同 inode）与符号链接/软链接（Symbolic Links，独立 inode 包含路径字符串）；跨进程并发读写与文件锁（File Locking，劝告锁 vs 强制锁）。
    \item \textbf{文件块物理分配经典策略}：
    \begin{itemize}
        \item 连续分配：顺序读写极速，但易产生严重外部碎片且文件难以动态扩充。
        \item 链式分配：消除外部碎片，但随机寻道性能低下且指针断裂易损坏文件；显式链表文件分配表（FAT）思想。
        \item 索引分配：为每个文件配置索引节点；经典 UNIX 多级混合索引（直接指针块、一级间接块、二级间接块、三级间接块）的地址寻址极限。
    \end{itemize}
    \item \textbf{文件系统磁盘组织与元数据布局}：引导块、超级块（Superblock，存储文件系统全局配置与元数据总览）、inode 位图、数据块位图、inode 数组表及实际数据存储块；目录项结构（线性链表、散列表、B+ 树）及多级路径名递归解析机制。
    \item \textbf{磁盘块分配/释放技术对性能的影响}：内部碎片与外部碎片的成因；磁盘碎片整理与块分配局部性聚合算法。
    \item \textbf{外存空闲块管理技术}：位向量/位图法、空闲块链表法、成组链接法（UNIX 传统做法）。
    \item \textbf{目录分层实现与检索加速}：目录作为包含文件名到 inode 映射项的特殊文件；路径解析中的目录项缓存与极速检索优化。
\end{itemize}

\subsubsection{示范性学习成果与能力指标}
\begin{itemize}
    \item \textbf{[解释] 文件系统底层架构设计决策}：深入阐明在设计底层文件系统时，在寻道时延、磁盘利用率、元数据开销与故障恢复速度之间面临的根本设计决策。
    \item \textbf{[评估] 多元文件分配机制全面对比}：严格对比连续分配、链式分配与多级索引分配在顺序读写、随机点查、动态扩容及容错韧性上的优势与劣势。
    \item \textbf{[应用] 面向文件系统特性的高效 I/O 编程}：基于对块缓冲、预读机制及延迟写特性的深刻理解，编写出高吞吐且兼顾断电安全的应用层 I/O 软件。
\end{itemize}

\subsection{高级文件系统与现代存储}
\textbf{【学时建议】}：2 领域核心学时。

\subsubsection{领域核心主题}
\begin{itemize}
    \item \textbf{现代磁盘管理与虚拟文件系统（VFS）}：磁盘分区表（GPT/MBR）、逻辑卷管理（LVM）；虚拟文件系统（VFS）核心抽象架构：超级块对象、inode 对象、dentry 目录项对象与文件对象（\texttt{file\_operations} 虚函数指针表），向下统一纳管各类异构文件系统。
    \item \textbf{企业级文件系统前沿技术}：Ext4、XFS、ZFS 与 Btrfs 等现代文件系统的核心架构与高并发读写特性。
    \item \textbf{内存映射文件（mmap）机制}：将磁盘文件块直接映射至进程虚拟地址空间，读写文件转化为直接内存访问，消除系统调用与双重内核缓冲拷贝开销，支持多进程高效共享大文件。
    \item \textbf{特殊用途与虚拟伪文件系统}：\texttt{procfs}（暴露内核运行期数据结构与进程树）、\texttt{sysfs}（统一设备模型导出）、\texttt{tmpfs}/\texttt{ramfs}（高速纯内存文件系统）。
    \item \textbf{高级检索、加密与定期快照}：元数据快速索引、快照生成与异地灾备备份机制。
    \item \textbf{日志型与写时复制文件系统}：
    \begin{itemize}
        \item 日志型文件系统：预写式日志（WAL）机制；元数据日志与完整数据日志在崩溃一致性与吞吐量上的权衡。
        \item 日志结构文件系统与写时复制（COW）存储架构：所有更新以连续追加方式写入新位置，天然杜绝就地覆写引发的数据不一致，大幅提升随机写性能。
    \end{itemize}
\end{itemize}

\subsubsection{选修与前沿拓展主题}
\begin{itemize}
    \item \textbf{分布式文件系统与现代存储架构}：网络附加存储（NAS）、存储区域网络（SAN）；主流网络文件系统协议（NFS、SMB/CIFS）；超大规模分布式文件系统（Ceph、HDFS、Google GFS）；云原生对象存储体系（OSS / S3）。
    \item \textbf{全盘与文件级透明加解密技术}：静态数据安全防护。
    \item \textbf{崩溃一致性与应用层持久化保障}：\texttt{fsync()}、\texttt{fdatasync()} 屏障语义；现代 ACID 关系型数据库底层引擎对操作系统刷盘保障的深层依赖。
\end{itemize}

\subsubsection{示范性学习成果与能力指标}
\begin{itemize}
    \item \textbf{[解释] 硬件介质演进对文件系统的重构}：阐明高速非易失性介质（如 NVMe SSD、持久内存 PMEM）的发展如何颠覆性改变了文件系统的设计重心（从磁头寻道优化转向低锁争用与 CPU 周期节省）。
    \item \textbf{[应用] 文件系统多层映射与全景调用}：熟练将上层抽象的 POSIX 文件 API 逐步映射追踪至 VFS 抽象层、具体文件系统驱动、通用块层直至物理控制器设备驱动。
    \item \textbf{[解释] 挂载体系分类与应用实践}：清晰识别并归类本地分区挂载、网络远程挂载、回环设备挂载以及容器只读分层挂载的底层区别。
    \item \textbf{[解释] 专用场景文件系统特性选型}：针对海量微小图片存储或超大连续科学数据集等特定场景，选择匹配的专用文件系统并阐述其特性优势。
    \item \textbf{[解释] 日志技术与崩溃一致性恢复}：深入阐明文件系统日志机制如何通过事务提交与检查点，在系统断电或意外崩溃重启后，无需漫长全盘扫描（\texttt{fsck}）即可极速恢复数据一致性。
    \item \textbf{[解释] 分布式存储系统复杂性透视}：阐明分布式文件系统的体系设计初衷，分析网络时延、跨节点容错与一致性协议（如 Raft、Paxos）引入的复杂性瓶颈。
    \item \textbf{[解释] 分布式存储协议谱系与选型}：列举主流分布式存储系统与协议，分析其在云原生大数据处理中的典型应用场景。
    \item \textbf{[解释] 文件系统容错增强前沿机制}：系统掌握文件系统用于防止静默数据破坏的数据端到端校验和与自动静默自愈技术。
\end{itemize}

\subsection{系统虚拟化与容器技术}
\textbf{【学时建议】}：3 领域核心学时。

\subsubsection{领域核心主题}
\begin{itemize}
    \item \textbf{虚拟化与隔离的系统目标}：通过逻辑抽象将单台物理服务器划分为多个相互隔离的运行环境，实现强安全防护、动态资源分配与可预测的服务质量。
    \item \textbf{高级分页与二级内存虚拟化}：虚拟内存之上的再次虚拟化；传统纯软件影子页表的维护开销；现代硬件辅助二级地址转换（SLAT，如 Intel EPT、AMD NPT）的硬件级地址翻译流水线。
    \item \textbf{虚拟文件系统与虚拟设备抽象}：虚拟外设的仿真机制；标准化半虚拟化驱动框架在网络与块设备 I/O 中的低开销透传。
    \item \textbf{操作系统级轻量虚拟化——容器技术}：
    \begin{itemize}
        \item 容器相较于虚拟机的本质差异：共享宿主机操作系统内核，无需全量 Guest OS 仿真。
        \item Linux 命名空间：PID 进程隔离、Mount 挂载点隔离、IPC 进程间通信隔离、Net 网络栈隔离、User 用户特权隔离与 UTS 主机名隔离。
        \item Linux 控制组：实现容器对 CPU 核心份额、物理内存上限、块 I/O 读写带宽等资源的硬性限制与精细审计。
        \item 联合文件系统：支持只读镜像层叠加可写容器层的极速分层存储机制。
    \end{itemize}
    \item \textbf{虚拟化环境下的系统抖动}：内存过度承诺与各虚机/容器并发争抢物理内存引发的全系统性能塌陷机理。
\end{itemize}

\subsubsection{选修与前沿拓展主题}
\begin{itemize}
    \item \textbf{虚拟化技术分型全景}：硬件指令集纯仿真、硬件辅助全虚拟化、半虚拟化、操作系统级容器、服务与网络虚拟化（NFV）。
    \item \textbf{可移植虚拟化；指令级仿真与轻量隔离的对比}：指令仿真（如 QEMU 跨架构仿真）的高昂动态二进制翻译开销与内核隔离（如 Docker 容器）的原生性能对比。
    \item \textbf{虚拟化的系统级性能损耗剖析}：VM-Exit / VM-Entry 上下文切换开销、跨层 TLB 刷新与虚拟中断注入时延。
    \item \textbf{虚拟机逃逸与容器安全挑战}：Hypervisor 漏洞引发的虚拟机逃逸；容器共享内核导致的安全漏洞渗透宿主机威胁。
    \item \textbf{虚拟机监控器分类与架构}：
    \begin{itemize}
        \item Type-1（裸金属型）：直接掌控裸机硬件，无独立宿主操作系统（如 VMware ESXi、Xen）。
        \item Type-2（宿主型）：依托宿主内核调度提供虚拟化支持（如 QEMU/KVM、VirtualBox）。
    \end{itemize}
\end{itemize}

\subsubsection{示范性学习成果与能力指标}
\begin{itemize}
    \item \textbf{[解释] 硬件虚拟化扩展指令集支持}：深入阐明现代 CPU 硬件体系（如 Intel VT-x、AMD-V）如何通过根模式与非根模式为虚拟机提供硬件级提速与特权安全隔离。
    \item \textbf{[解释] 指令仿真与内核级隔离本质差异}：严格区分基于动态二进制翻译的指令级仿真与基于内核命名空间/控制组的轻量化隔离在运行机制与性能上的鸿沟。
    \item \textbf{[评估] 虚拟化架构全方位权衡选型}：针对云计算中心与边缘部署场景，客观评估裸金属虚机、宿主虚机与微服务轻量容器在资源利用率、部署弹性、隔离安全性与启动延迟上的全景得失。
    \item \textbf{[解释] Hypervisor 核心机理与分型演进}：系统掌握 Type-1 与 Type-2 虚拟机监控器的内部架构差异与主流工业界代表作的运行机理。
\end{itemize}

\subsection{实时与嵌入式系统}
\textbf{【学时建议】}：2 领域核心学时。

\subsubsection{领域核心主题}
\begin{itemize}
    \item \textbf{实时多任务模型与确定性调度}：实时系统核心目标是满足严格的时间约束；时钟驱动调度、优先级驱动调度；最坏情况执行时间（WCET）与确定性响应度量。
    \item \textbf{截止时间约束与实时性挑战}：任务释放时间、截止时间、时延抖动对工业控制回路稳定性的决定性作用。
    \item \textbf{硬实时与软实时系统的本质分水岭}：
    \begin{itemize}
        \item 硬实时：错过截止时间即判定为系统灾难性崩溃（如航空航天机载飞行控制、汽车防抱死制动系统 ABS）。
        \item 软实时：错过截止时间仅表现为服务质量退化或用户体验下降，系统仍可继续运行（如音视频流式多媒体通信、低延迟在线竞技游戏）。
    \end{itemize}
\end{itemize}

\subsubsection{选修与前沿拓展主题}
\begin{itemize}
    \item \textbf{实时环境下的内存与存储管理}：严禁不可预测的动态内存碎片化分配；页内存常驻物理内存锁定（\texttt{mlock}）以彻底消除缺页中断时延抖动。
    \item \textbf{系统瞬时故障、风险评估与确定性自愈恢复}：故障安全与高确定性降级运行模式。
    \item \textbf{安全攸关实时系统的核心陷阱与对策}：
    \begin{itemize}
        \item 优先级翻转现象剖析（火星探路者号经典案例）。
        \item 优先级继承协议 (PIP)与优先级天花板协议 (PCP)。
        \item 硬件看门狗定时器的自动超时重置与防死机机制。
    \end{itemize}
\end{itemize}

\subsubsection{示范性学习成果与能力指标}
\begin{itemize}
    \item \textbf{[解释] 实时操作系统的根本特征}：明确界定实时操作系统（RTOS）的根本特征：时间确定性与计算可预测性高于平均峰值吞吐量。
    \item \textbf{[解释] 软件时延来源与抖动分析}：深入剖析现代软件系统全链路中的延迟来源（中断屏蔽、关调度、锁竞争、内存缺页、Cache 未命中）及其长尾分布特征。
    \item \textbf{[解释] 实时系统特殊风险与对策}：系统阐述安全攸关实时系统面临的物理风险与优先级翻转隐患，并说明优先级继承协议的化解原理。
    \item \textbf{[解释] 专用 RTOS 架构与微内核特性}：详述典型专用实时操作系统（如 FreeRTOS、VxWorks、RT-Linux）的轻量内核架构与快速上下文切换机制。
\end{itemize}

\subsection{系统容错与高可用性}
\textbf{【学时建议】}：3 领域核心学时。

\subsubsection{领域核心主题}
\begin{itemize}
    \item \textbf{高可靠与高可用系统度量体系}：
    \begin{itemize}
        \item 可靠性：平均无故障时间（MTBF）。
        \item 可用性：系统正常运行时间占比（如“五个九”$99.999\%$ 可用度，年宕机时间小于 5.26 分钟）。
        \item 平均修复时间（MTTR）的系统级优化。
    \end{itemize}
    \item \textbf{软硬件协同容错技术——RAID 阵列详解}：
    \begin{itemize}
        \item 廉价磁盘冗余阵列（RAID）各级别机制：
        \item RAID 0：纯数据条带化，提升读写吞吐，无容错能力。
        \item RAID 1：磁盘镜像，完全空间冗余，高容错但磁盘利用率仅 $50\%$。
        \item RAID 5：分布式奇偶校验条带化，允许单盘故障损坏而不丢数据，利用率 $(N-1)/N$。
        \item RAID 6：双重分布式独立奇偶校验，允许两块磁盘同时损坏。
        \item RAID 10：条带化与镜像混合架构，兼备超高性能与高容错能力。
    \end{itemize}
\end{itemize}

\subsubsection{选修与前沿拓展主题}
\begin{itemize}
    \item \textbf{系统冗余设计三元方法论}：空间冗余（双机热备、三模冗余 TMR 多数表决）、时间冗余（异常指令/事务超时重试）与信息冗余（纠错码 ECC）。
    \item \textbf{硬件瞬态错误检测与在线纠错}：ECC 易失性内存单比特翻转自动纠错与双比特报错；针对磁盘静默数据破坏的端到端数据校验和机制。
    \item \textbf{文件系统一致性审计与修复}：\texttt{fsck} 工具的离线扫描与日志重放恢复原则。
    \item \textbf{容错技术工业级核心应用}：超大型金融交易数据库、安全关键航空航天机载软件。
    \item \textbf{操作系统核心服务自愈防御}：内核崩溃内核转储、核心服务隔离沙箱、独立微内核崩溃热重启与状态无损迁移。
\end{itemize}

\subsubsection{示范性学习成果与能力指标}
\begin{itemize}
    \item \textbf{[解释] 操作系统支撑容错与高可用的途径}：阐述操作系统如何通过内核基础设施、驱动隔离与状态持久化机制，为上层应用提供坚固的容错底座。
    \item \textbf{[解释] 容错实现技术方法谱系}：系统总结并分类在操作系统软硬件各层级中实现容错的经典技术手段。
    \item \textbf{[解释] 局部故障下的连续不间断服务}：深入分析在面临物理磁盘损坏、主存局部错误或外设断开等局部软硬件故障时，操作系统如何维持核心服务连续在线。
    \item \textbf{[评估] 容错机制的性能与开销权衡}：全面评估在操作系统中引入数据镜像、校验计算与状态快照对系统运行性能、存储开销与软件复杂度的折损与权衡。
    \item \textbf{[解释] 内核级容错深度议题}：详述操作系统内核崩溃保护、快速热重启与事务化状态机自愈的深度实现原理。
\end{itemize}

\subsection{社会、伦理与专业素养议题}
\textbf{【学时建议】}：4 领域核心学时。

\subsubsection{领域核心主题}
\begin{itemize}
    \item \textbf{操作系统领域的开源生态演进}：
    \begin{itemize}
        \item 开源操作系统内核（如 Linux Kernel）的代码审查、协作漏洞挖掘、协同负责任披露与全球补丁分发网络。
        \item 开源客户机操作系统与开源宿主机操作系统的协同繁荣。
        \item 商业化盈利模式的历史演化：从传统盒装软件授权许可，走向订阅技术支持、开源核心与云原生托管服务。
    \end{itemize}
    \item \textbf{操作系统生命周期终结 (EOL)的社会隐患}：
    \begin{itemize}
        \item 停服弃用老旧系统停发安全补丁所带来的系统性安全脆弱性。
        \item 交通基础设施、医疗设备与工业控制系统因历史技术债务继续依赖 EOL 操作系统所引发的公共安全与伦理风险。
    \end{itemize}
\end{itemize}

\subsubsection{示范性学习成果与能力指标}
\begin{itemize}
    \item \textbf{[解释] 开源内核协同安全审计的利弊}：辩证评估在开源操作系统内核中公开发现与修复漏洞的双重效应（代码透明带来的全球白帽安全强化 vs 零日攻击窗口期被恶意利用）。
    \item \textbf{[解释] Linux 开源生态的全球历史影响}：结合计算机产业演进脉络，客观评价 Linux 等开源操作系统对全球互联网基础设施与云计算现代化的革命性推动作用。
    \item \textbf{[解释] 依赖超期退役操作系统的社会风险}：详尽剖析银行金融、医疗健康等公共关键基础设施长期依赖超期操作系统（如 Windows XP/7、旧版 Linux）所引发的法律、经济与伦理危机。
    \item \textbf{[解释] 不同商业模式下的缺陷修复动力机制}：深入阐明开源社区驱动模式与闭源专有商业模式在漏洞响应速度、测试覆盖度以及利益驱动机制上的制度性差异。
\end{itemize}

\section{专业素养与思维倾向}

\begin{itemize}
    \item \textbf{责任担当}：对系统底层架构设计、特权模式配置、并发锁协议以及资源调度策略的决策后果保持高度负责的工程敬畏心，深切认识到内核级哪怕微小缺陷都将引发整机崩溃或系统级严重安全失陷。
    \item \textbf{极致严谨}：在涉及系统级或高性能软件项目研发中，极其审慎细密地考量操作系统微观机制（如虚拟内存缺页、硬件中断延迟、上下文切换开销、锁竞争与内存屏障）对宏观业务软件表现的深刻连锁影响。
\end{itemize}

\section{数学基础要求}

\begin{itemize}
    \item \textbf{必修数学素养}：离散数学
    \begin{itemize}
        \item 集合论与二元关系（用于形式化建模用户访问控制矩阵与进程资源所有权）。
        \item 图论基础（资源分配图 Resource Allocation Graph 构建、有向图环路检测算法在死锁判定中的核心应用）。
        \item 布尔代数与数理逻辑（硬件保护掩码、权限位运算、状态判断逻辑）。
    \end{itemize}
\end{itemize}

\section{交叉主题与跨学科共享概念}

操作系统作为底层系统软件的中枢，与计算机学科多个知识领域存在高度的横向交叉渗透：
\begin{itemize}
    \item \textbf{安全与保护} $\longleftrightarrow$ 计算机安全 (SEC)：内存越权防护、特权级最小化原则。
    \item \textbf{分布式文件系统} $\longleftrightarrow$ 网络与通信 (NC)：远程存储调用、网络块设备与网络延迟处理。
    \item \textbf{进程、线程与虚拟化映射} $\longleftrightarrow$ 计算机体系结构与组织 (AR)：硬件特权模式、寄存器现场与多线程硬件支持。
    \item \textbf{存储层次结构与高速缓存} $\longleftrightarrow$ 计算机体系结构与组织 (AR)：L1/L2/L3 硬件缓存与虚拟内存分页的局部性协同。
    \item \textbf{缓冲机制、DMA、轮询与内存映射 I/O} $\longleftrightarrow$ 计算机体系结构与组织 (AR)：硬件设备控制器交互标准。
    \item \textbf{设备驱动程序设计、抽象与实现} $\longleftrightarrow$ 计算机体系结构与组织 (AR)：底层总线协议驱动与寄存器读写。
    \item \textbf{系统状态模型与状态迁移图} $\longleftrightarrow$ 系统基础 (SF)：进程与线程生命周期状态机。
    \item \textbf{虚拟化与隔离机制} $\longleftrightarrow$ 系统基础 (SF)：硬件与软件隔离技术通则。
    \item \textbf{高可靠与高可用系统构建} $\longleftrightarrow$ 系统基础 (SF)：故障容错与 MTBF/MTTR 建模。
\end{itemize}

\section{真实能力任务规范}

\begin{taskbox}[任务 1：设计并开发多进程安全通信软件组件]
\textbf{【任务情境】}：在一个需要严格进程故障隔离的企业级核心数据处理平台中，多个独立执行的业务进程协同处理敏感高吞吐数据流，必须彻底杜绝死锁与数据竞争。\\
\textbf{【能力陈述】}：综合运用操作系统提供的管道、POSIX 共享内存、信号量或消息队列等底层 IPC 通信机制，设计并实现一套高吞吐、零死锁、抗竞态且具备完备权限隔离的进程间通信模块。\\
\textbf{【能力领域】}：软件工程 \quad \textbf{【能力单元】}：系统设计 / 软件开发 / 单元测试\\
\textbf{【关联知识单元】}：操作系统/操作系统的角色与定位、操作系统/操作系统基本原理、操作系统/并发机制、操作系统/处理器调度、操作系统/进程模型与调度、操作系统/内存管理、操作系统/安全防护与保护机制、并行分布式/通信与协同\\
\textbf{【核心要求】}：计算机核心 \quad \textbf{【达标等级】}：\textbf{开发与实现}
\end{taskbox}

\begin{taskbox}[任务 2：高可用应用在虚拟化与容器化运行时中的自动化部署与治理]
\textbf{【任务情境】}：企业核心业务服务需要从传统单机环境迁移部署至多租户云端虚拟化平台或轻量化容器集群中，需确保严格的跨租户隔离屏障、可预测的资源配额与故障自愈能力。\\
\textbf{【能力陈述】}：精准识别并消除目标运行环境中的宿主依赖与权限越权风险；编写标准化基础设施自动化配置脚本与容器清单文件；配置基于 Linux cgroups 的 CPU/内存资源硬限额，构建全生命周期的执行监控与异常告警机制。\\
\textbf{【能力领域】}：系统工程 \quad \textbf{【能力单元】}：系统评估 / 运维保障 / 持续改进\\
\textbf{【关联知识单元】}：操作系统/操作系统的角色与定位、操作系统/操作系统基本原理、操作系统/并发机制、操作系统/处理器调度、操作系统/进程模型与调度、操作系统/内存管理、操作系统/安全防护与保护机制、体系结构/汇编级机器组织\\
\textbf{【核心要求】}：计算机核心 \quad \textbf{【达标等级】}：\textbf{应用}
\end{taskbox}

\begin{taskbox}[任务 3：多租户高并发场景下操作系统资源竞争分析与性能调优]
\textbf{【任务情境】}：某多租户生产环境服务器在面对周期性高并发业务洪峰时，突发显著的系统吞吐量暴跌、长尾时延剧增与频繁缺页异常，亟需定位内核级系统性能瓶颈并完成重构优化。\\
\textbf{【能力陈述】}：熟练运用操作系统级性能分析观测工具（如 \texttt{perf}、eBPF 工具链、\texttt{vmstat}、\texttt{strace}），精准捕获系统瓶颈根因（内核锁严重争用、缓存伪共享、内存剧烈颠簸或调度延迟）；提出并实施针对性的 CPU 亲和性绑定、大页内存启用、I/O 调度器调整与内核参数调优方案。\\
\textbf{【能力领域】}：系统工程 \quad \textbf{【能力单元】}：系统评估分析 / 性能优化与改进\\
\textbf{【关联知识单元】}：操作系统/操作系统基本原理、操作系统/并发机制、操作系统/进程模型与调度、操作系统/内存管理、操作系统/处理器调度、操作系统/安全防护与保护机制\\
\textbf{【核心要求】}：领域核心 \quad \textbf{【达标等级】}：\textbf{评估与应用}
\end{taskbox}

\section{课程设计与教学实施建议}

联合工作组建议高校面向计算机专业本科生开设为期一学期的《操作系统导论》主干专业必修课（建议理论授课约 37$\sim$48 学时，并配套不少于 32 学时的上机系统级实验）：

\begin{itemize}
    \item \textbf{操作系统的角色与目标}：3 讲授学时。
    \item \textbf{操作系统原理与体系架构}：3 讲授学时。
    \item \textbf{并发与同步机制}：7 讲授学时（含多线程编程、死锁算法与无锁机制）。
    \item \textbf{处理器调度与资源分配}：3 讲授学时。
    \item \textbf{进程模型与执行环境}：3 讲授学时。
    \item \textbf{内存管理与虚拟内存}：4 讲授学时。
    \item \textbf{安全保护与故障隔离}：4 讲授学时。
    \item \textbf{设备管理与输入输出系统}：2 讲授学时。
    \item \textbf{文件系统 API 与底层实现}：2 讲授学时。
    \item \textbf{系统虚拟化与容器技术}：3 讲授学时。
    \item \textbf{高级文件系统与现代存储}：2 讲授学时。
    \item \textbf{实时与嵌入式系统}：1 讲授学时。
    \item \textbf{系统容错与高可用性}：1 讲授学时。
    \item \textbf{社会、伦理与专业素养议题}：4 讲授学时。
\end{itemize}

\noindent\textbf{先修课程准入依赖（来自计算机体系结构 AR 领域）：}
\begin{itemize}
    \item 汇编语言与机器组织
    \item 内存管理硬件基础
    \item 软件可靠性与硬件容错
    \item 接口技术与系统互联通信
    \item 计算机功能部件组织
\end{itemize}

\noindent\textbf{课程终结综合能力养成指标：}\\
修读完本课程的学生，应能深刻理解操作系统在处理器、内存、外设资源管理中的微观工作机制及其在性能吞吐与系统安全上的深刻权衡；能够熟练、安全、无死锁地设计与实现多进程/多线程并发控制与通信模块；能够客观评估并选型开源或专有操作系统方案；能够深刻掌握系统虚拟化、轻量容器与沙箱隔离技术作为现代安全操作系统基石的核心价值与实践部署方法。

\section{指导委员会工作组成员}

\begin{itemize}
    \item \textbf{主席}：Monica D. Anderson (阿拉巴马大学，University of Alabama, Tuscaloosa, AL, USA)
    \item \textbf{委员}：
    \begin{itemize}
        \item Renzo Davoli (意大利博洛尼亚大学，University of Bologna, Italy)
        \item Avi Silberschatz (美国耶鲁大学，Yale University, USA，《Operating System Concepts》作者)
        \item Marcelo Pias (巴西里奥格兰德联邦大学，Federal University of Rio Grande - FURG, Brazil)
        \item Mikey Goldweber (美国泽维尔大学，Xavier University, Cincinnati, USA)
        \item 向乔，厦门大学
    \end{itemize}
\end{itemize}

\section{附录：核心主题与技能要求等级对照表}

\begin{longtable}{p{0.20\textwidth} p{0.48\textwidth} c c c}
\toprule
\textbf{知识单元} & \textbf{核心知识点与具体主题} & \textbf{认知等级} & \textbf{分类} & \textbf{学时} \\
\midrule
\endfirsthead
\toprule
\textbf{知识单元} & \textbf{核心知识点与具体主题} & \textbf{认知等级} & \textbf{分类} & \textbf{学时} \\
\midrule
\endhead
\bottomrule
\endfoot
\multirow{4}{*}{\makecell[l]{操作系统/操作系统的\\角色与目标}} 
& 操作系统作为硬件与应用软件间的中介抽象层 [跨界 PL] & 解释 & 计算机核心 & \multirow{4}{*}{3} \\
& 通用操作系统功能（接口、应用执行生命周期、数据持久化） & 解释 & 计算机核心 & \\
& 专用系统功能（嵌入式、专用服务器、引导程序与安全启动） & 解释 & 计算机核心 & \\
& 系统设计权衡（效率、健壮性、可移植性、安全、功耗）；安全纵深防御关切 & 解释 & 计算机核心 & \\
\midrule
\multirow{4}{*}{\makecell[l]{操作系统/操作系统\\原理与体系架构}} 
& 软件设计范式（宏内核、分层、微内核、独核）；三大系统抽象 & 解释 & 计算机核心 & \multirow{4}{*}{3} \\
& 系统调用概念、API 桥接机制与 Syscall ABI 演进哲学 [共享 AR] & 解释 & 计算机核心 & \\
& 机器特权指令与系统资源的硬件级保护 [共享 AR] & 解释 & 计算机核心 & \\
& 硬件中断机制：ISR 与分派实现；双模式转换与系统调用全流程 [共享 AR] & 解释 & 计算机核心 & \\
\midrule
\multirow{4}{*}{\makecell[l]{操作系统/并发与\\同步机制}} 
& 线程抽象、竞态条件、临界区与关中断机制 & 解释 & 计算机核心 & \multirow{2}{*}{3} \\
& 死锁与饥饿现象、多处理器并发挑战（自旋锁、重入性） & 解释 & 计算机核心 & \\
\cline{2-5}
& 线程创建、生命周期状态机、TCB 数据结构与标准线程 API & 应用 & 领域核心 & \multirow{2}{*}{3} \\
& 死锁四大条件形式化与防范策略；同步原语（互斥锁、信号量、条件变量） [共享 PD] & 应用 & 领域核心 & \\
\midrule
\multirow{2}{*}{\makecell[l]{操作系统/处理器调度\\与资源分配}} 
& 抢占式与非抢占式调度；定时器虚拟化（硬件时钟到软件定时器树） & 解释 & 领域核心 & \multirow{2}{*}{3} \\
& 经典调度策略（FCFS、SJF、优先级、RR、MLFQ）；SMP 多核亲和性调度 [共享 AR] & 解释 & 领域核心 & \\
\midrule
\multirow{2}{*}{\makecell[l]{操作系统/进程模型与\\执行环境}} 
& 进程与线程的控制权虚拟化；进程内存布局段结构（Text、Data、Stack、Heap） & 解释 & 领域核心 & \multirow{2}{*}{3} \\
& ELF 可执行文件结构、动态链接与 GOT/PLT；上下文切换开销；IPC 通信机制 & 解释 & 领域核心 & \\
\midrule
\multirow{3}{*}{\makecell[l]{操作系统/内存管理与\\虚拟内存}} 
& 物理寻址与 MMU 硬件机制；存储层次结构对 OS 策略与 Per-CPU 缓存的影响 & 解释 & 领域核心 & \multirow{3}{*}{4} \\
& 分页机制、多级页表、TLB 缓存；页面置换算法（FIFO、LRU、Clock）与抖动防范 & 解释 & 领域核心 & \\
& 内存分配算法（伙伴系统、Slab、空闲链表、大页）；安全隔离（NX、ASLR） & 解释 & 领域核心 & \\
\midrule
\multirow{3}{*}{\makecell[l]{操作系统/安全保护与\\故障隔离}} 
& 操作系统安全保护机制全景；对抗攻击手法（调度饥饿、特权提升） & 应用 & 计算机核心 & \multirow{2}{*}{3} \\
& 真实商业 OS 经典安全漏洞复盘（栈溢出、UAF、TOCTOU）；防御与备份 & 应用 & 计算机核心 & \\
\cline{2-5}
& 机制与策略分离原则；分级保护环（Ring 0/3/-1）；访问控制列表（ACL）与多因子认证 & 应用 & 领域核心 & 1 \\
\midrule
\multirow{2}{*}{\makecell[l]{操作系统/设备管理与\\I/O 系统}} 
& I/O 缓冲技术策略（单缓冲、双缓冲、循环环形队列） & 解释 & 领域核心 & \multirow{2}{*}{2} \\
& 直接内存访问（DMA）、中断驱动 I/O、MMIO；块设备介质演进与调度（HDD/SSD） [共享 AR] & 解释 & 领域核心 & \\
\midrule
\multirow{2}{*}{\makecell[l]{操作系统/文件系统 API\\与底层实现}} 
& 文件本质抽象、挂载机制、权限访问控制、文件共享与文件锁 & 解释 & 领域核心 & \multirow{2}{*}{2} \\
& 基础分配方式（连续、FAT 链表、UNIX inode 索引）；目录项解析与空闲块位图 & 解释 & 领域核心 & \\
\midrule
\multirow{2}{*}{\makecell[l]{操作系统/高级文件系统\\与现代存储}} 
& 磁盘分区、LVM 卷管理、VFS 抽象框架与 \texttt{file\_operations} 虚函数表 & 解释 & 领域核心 & \multirow{2}{*}{2} \\
& 内存映射文件（mmap）；伪文件系统；日志型与写时复制系统 & 解释 & 领域核心 & \\
\midrule
\multirow{2}{*}{\makecell[l]{操作系统/系统虚拟化\\与容器技术}} 
& 虚拟化与隔离服务质量；二级内存虚拟化硬件加速（EPT/NPT）；Virtio & 解释 & 领域核心 & \multirow{2}{*}{2} \\
& 虚拟化过度承诺下的颠簸；轻量化容器底层技术（Linux Namespaces 与 cgroups） & 解释 & 领域核心 & \\
\midrule
\multirow{1}{*}{\makecell[l]{操作系统/实时与\\嵌入式系统}} 
& 实时多任务模型、截止时间约束、抖动；硬实时与软实时系统的本质分水岭 [跨界 PL] & 解释 & 领域核心 & 1 \\
\midrule
\multirow{1}{*}{\makecell[l]{操作系统/系统容错与\\高可用性}} 
& 高可靠（MTBF）与高可用（$99.999\%$）系统设计原理；RAID 0/1/5/6/10 机制详解 & 解释 & 领域核心 & 1 \\
\midrule
\multirow{2}{*}{\makecell[l]{操作系统/社会、伦理与\\专业素养议题}} 
& 开源操作系统生态演化、内核协同漏洞披露与打补丁机制；商业模式演进 & 解释 & 领域核心 & \multirow{2}{*}{4} \\
& 停服弃用与老旧系统过生命周期（EOL）引发的社会风险 [跨界 SE] & 解释 & 领域核心 & \\
\end{longtable}
\chapter{并行与分布式计算}

\section{领域导言与演进背景}

\begin{preamblebox}[知识领域概述与变革]
并行与分布式计算 (PDC)致力于组织、协调与控制跨越不同物理时空且并发发生的多种计算活动。伴随处理器内部晶体管规模的激增、单机中多核处理器的普及以及网络化计算节点的无处不在，并行性与分布式特性已成为现代计算系统的本质属性。相较于传统串行程序，并行与分布式架构不仅能够成倍释放硬件潜能、大幅提升计算性能，更是应对现实物理世界固有互联性、部件局部失效以及通信不可靠性的必然技术路径。通过打破传统串行编程中“所有计算步骤必须在单一位置按严格时序依次发生”的桎梏，并行与分布式计算揭示出贯穿现代计算系统各个层级的全套核心抽象、设计范式与分析方法。

在学术界与工业界的通用表述中：
\begin{itemize}
    \item \textbf{并行编程}：侧重于调度和组织多个计算活动在时间上的“协同并发发生”；
    \item \textbf{分布式编程}：侧重于协调发生在不同物理“空间位置”上的独立计算活动；
    \item \textbf{并发编程}：侧重于管理正在推进中的多项活动之间及其与外部环境的“交错交互”。
\end{itemize}
然而在现实工程中，这三者往往深度交织、密不可分。并行性天然蕴含某种程度的分布性，因为除依赖时间片上下文切换或量子效应外，无顺序依赖的多个并发操作必然发生于不同的物理硬件单元之上；反之，在缺乏明确通信约束时，发生于不同地理位置的操作彼此之间亦不存在绝对的先后次序。本知识领域既重点聚焦“主动显式引入并行与分布式”的技术语境，也涵盖“因数据或服务天然跨地域分布而不得不引入并发控制”的被动约束场景。

将并行与分布式计算讲深讲透极具教学挑战性。该领域已由早期的前沿高级选修课集群，演进为渗透至现代计算学科几乎每一个角落的通用核心知识体系。现实中，几乎所有具备串行解法的问题都存在并行或分布式解法，而更多前沿问题唯有在并发语境下才会真正涌现。历史上，该领域的各分支在不同计算子领域中呈现出非均衡的碎片化演进，有时甚至使用截然不同的术语、工程实践和课程载体，从而遮蔽了底层基本思想与核心技能的高度一致性。本知识领域的核心使命，正是提炼并讲授这一横跨软硬件体系的坚实知识基底。

\textbf{自 CS2013 以来的核心变革：}
\begin{enumerate}
    \item \textbf{知识体系抽象重构与原语凝练}：摒弃以往简单罗列特定多线程库（如 Pthreads、OpenMP）、消息传递库（如 MPI）或特定集群中间件（如 Hadoop）的做法，重塑为由“程序与执行”、“通信机制”、“协调与同步”、“软件工程”、“算法与应用领域”五大核心抽象原语构成的现代知识图谱。
    \item \textbf{精炼解耦的计算机核心}：将计算机专业通识必修的核心学时精炼为 10 学时（每个知识单元各 2 学时），聚焦于破除单机串行顺序执行假设的基本规律、Happens-before 偏序模型及非确定性并发直觉，不捆绑任何特定语言或特定硬件架构。
    \item \textbf{弹性模块化的领域核心}：领域核心（共 26 学时）采用“选修以下一项或多项”的拓展机制，允许各高校根据自身师资与专业定位，灵活构建高性能计算（HPC）、异构加速（GPU/SIMD）、多核共享内存系统、弹性云原生或去中心化容错存储等多元课程方案。
    \item \textbf{全面引入真实能力规范}：设置真实工程任务，重点考察学生在多核/异构/分布式环境下的并发重构、性能剖析（Work-Span 跨度分析与 Amdahl 定律应用）、死锁预防以及容错共识协议设计能力。
\end{enumerate}

\textbf{本领域课程起草委员会：}
\begin{itemize}
    \item \textbf{主席}：Doug Lea（纽约州立大学奥斯威戈分校）
    \item \textbf{委员}：Sherif Aly（开罗美国大学）、Michael Oudshoorn（高点大学）、项乔（Qiao Xiang，厦门大学）、Dan Grossman（华盛顿大学）、Sebastian Burckhardt（微软研究院）、Vivek Sarkar（佐治亚理工学院）、Maurice Herlihy（布朗大学）、Sheikh Ghafoor（田纳西理工大学）、Chip Weems（马萨诸塞大学阿默斯特分校）
    \item \textbf{主要贡献者}：Paul McKenney、Peter Buhr
\end{itemize}
\end{preamblebox}

\section{核心学时分配建议}

本知识领域在“向日葵模型”下的学时推荐如下表所示。计算机核心涵盖所有计算专业学生打破串行思维所必备的通识基础，领域核心则提供进阶深度选修指南：

\begin{table*}[htbp]
\centering\small
\setlength{\tabcolsep}{2.5pt}\footnotesize\begin{tabularx}{0.88\textwidth}{Y c c}
\toprule
\textbf{知识单元} & \textbf{计算机核心} & \textbf{领域核心} \\
\midrule
程序与执行 & 2 & 2 \\
通信机制 & 2 & 6 \\
协调与同步 & 2 & 6 \\
软件工程 & 2 & 3 \\
算法与应用领域 & 2 & 9 \\
\midrule
\textbf{合计} & \textbf{10 学时} & \textbf{26 学时} \\
\bottomrule
\end{tabularx}
\end{table*}

\textbf{学时说明与实施建议：}
\begin{itemize}
    \item \textbf{计算机核心}：共 10 学时，覆盖并行与分布式计算的核心理念。其主要教学目标是消除严格串行编程的思想禁锢，揭示并行与分布式系统及软件的本质结构与核心特性。学习成果侧重于开发小规模多活动并发程序，并分析其基础性质。该学时指标不包含特定编程语言、开发工具、计算框架与物理平台的具体细节教学（这些内容通常由后续具体专业课程结合实验项目承载）。
    \item \textbf{领域核心}：共 26 学时，各单元内容采用“选修以下一项或多项”的形式组织，为不同侧重点的专业选修课提供弹性空间。例如：
    \begin{itemize}
        \item 高性能计算（HPC）方向：聚焦异构数据并行与线性代数求解；
        \item 计算机系统方向：聚焦多核共享内存协调、中间件与系统资源管理；
        \item 分布式大数据方向：聚焦容错存储、分布式共识与大规模云服务。
    \end{itemize}
\end{itemize}

\section{知识单元与详细大纲}

\subsection{程序与执行}
\textbf{【学时建议】}：2 计算机核心学时，2 领域核心学时。

本单元奠定并行与分布式计算的核心抽象，重点讲授并行动作的定义、执行拓扑、时序关系以及启动控制机制。

\subsubsection{计算机核心主题}
\begin{itemize}
    \item \textbf{基本定义与核心属性}：
    \begin{itemize}
        \item \textbf{可并行化动作}：包括闭包、函数调用以及网络服务调用；复合动作、会话、并发任务与作用域；动作的实体命名与标识（如线程 ID、进程号、URI）；计算粒度与调度开销权衡。
        \item \textbf{动作先后定序与先行发生关系}：放宽严格串行时序约束；用串行/并行有向无环图建模程序组件的非确定性执行流动。
        \item \textbf{独立性判定}：基于操作的可交换性、数据流依赖性与前置条件，严格判定何时执行顺序互不相干。
        \item \textbf{一致性约定}：关于数值状态与谓词的协同协议；竞态条件、原子性与分布式共识直觉。
        \item \textbf{并行程序中的必要定序保障}：互斥锁、内存屏障、安全发布；实体间通信施加的时序制约（例如：消息的“发送”必然发生在“接收”之前）。
        \item \textbf{执行位置与执行实体}：承载并行动作的物理设备（硬件核、远程主机）；单一执行位置对多个并发实体的复用与调度；外部世界及人类用户作为并发实体的交互模式；基于时间片分时调度与虚拟化机制在少量处理器上仿真大量并行动作。
        \item \textbf{实体失效与通信故障}：不可信实体及不受本程序或管理域控制的网络协议所引发的系统失效；容错级别作为系统核心设计维度的折中考量。
    \end{itemize}
    \item \textbf{并行动作的启动机制}：
    \begin{itemize}
        \item \textbf{编排放置}：协调动作最终由指定物理实体或节点执行（从硬编码静态绑定、配置脚本到现代云平台的自动化弹性编排）；建立端到端通信与资源分配。
        \item \textbf{过程式启动}：在特定程序断点处派生并发动作（如创建新线程、Fork 子进程）；利用结构化作用域确保所有衍生动作全部结束时作用域方能退出。
        \item \textbf{反应式/事件驱动启动}：由外部事件触发并行动作（如注册事件回调处理器）；对动作具体启动和终止时间的非精确控制。
        \item \textbf{依赖驱动启动}：基于前序计算或前置任务的全部完结状态触发后续任务（如流式计算、延续触发 Continuations）。
    \end{itemize}
\end{itemize}

\subsubsection{领域核心与拓展主题}
\begin{itemize}
    \item \textbf{底层硬件映射与物理执行机制（选择以下一项或多项）：}
    \begin{itemize}
        \item CPU 数据级与指令级并行（ILP、流水线超标量、SIMD 向量扩展）。
        \item SIMD 向量化架构与异构数据并行（GPU、专用深度学习加速芯片）。
        \item 多核处理器中的抢占/协作调度并发、任务池与 Actor 消息传递模型。
        \item 计算集群、云计算基础设施与弹性算力资源编排。
        \item 参与者无边界、跨管理域的开放分布式系统。
        \item 新兴前沿计算模型：量子并行计算、DNA/分子计算。
    \end{itemize}
\end{itemize}

\subsubsection{示范性学习成果与能力目标}
\begin{enumerate}
    \item \textbf{[应用 (计算机核心)]} 针对复杂的复合数值表达式（例如 $a = (b + c) \times (d + e)$），绘制其并行计算依赖有向无环图（DAG），展示操作间的并发度与关键路径。
    \item \textbf{[评估 (计算机核心)]} 在给定的并发代码片段中准确排查并指出数据竞态隐患。
    \item \textbf{[解释 (计算机核心)]} 结合具体工程场景，深入阐述将既有串行程序并行化在多大程度上能够提升系统吞吐量或降低端到端时延，并客观评估由此对系统整体能效与热功耗产生的实际影响。
    \item \textbf{[开发 (领域核心)]} 编写一个高效统筹高并发事件（如传感器密集采样或网络数据包接收）的高性能计数统计函数。
    \item \textbf{[开发 (领域核心)]} 分别采用多种不同的编程风格（如基于线程池、函数式流式处理、Actor 消息传递等）实现经典 Filter-Map-Reduce（过滤-映射-规约）程序。
    \item \textbf{[开发 (领域核心)]} 编写一个并发网络服务程序，通过为每个新接入的客户端派生独立线程（或其他过程式激活机制）处理并返回所请求的 Web 页面。
\end{enumerate}

\subsection{通信机制}
\textbf{【学时建议】}：2 计算机核心学时，6 领域核心学时。

本单元系统剖析并发与分布式实体间交换信息的基础介质、信道抽象、内存一致性以及数据存储模型。

\subsubsection{计算机核心主题}
\begin{itemize}
    \item \textbf{通信基础要素}：
    \begin{itemize}
        \item \textbf{通信介质与形态}：管道/通道（Channels，消息传递或 I/O 流）、共享内存、异构显存、结构化数据存储；对底层物理硬件、网络拓扑、互联协议、语言级运行时以及仿真模拟机制的依赖。
        \item \textbf{通信通道}：显式端到端实体通信；寻址与命名机制；核心 API 体系（网络套接字 Sockets、基于体系结构或语言原生的通道构造如 Go Channel）。
        \item \textbf{内存通信模型}：实体仅通过指定物理/虚拟地址直接与内存进行信息交换的架构；内存一致性基础（按位原子性边界、缓存一致性 Cache Coherence、局部内存读写定序）；内存层次结构与空间/时间局部性（高速缓存 Cache、访问时延、伪共享 False Sharing 问题）；跨异构存储体的数据搬运（如 CPU 主存与 GPU 显存间的显式 DMA 传输）；多层级数据共享域、作用域与私有缓存。
        \item \textbf{协同数据存储}：由多个并发实体协同维护的结构化数据（映射表 Maps、集合 Sets 等并发抽象数据类型）；存储模式形态（独占式 Owned、共享式 Shared、分片式 Sharded、副本复制式 Replicated、不可变数据 Immutable、多版本并发控制 Versioned）。
    \end{itemize}
    \item \textbf{基于通信的并发编程范式}：
    \begin{itemize}
        \item 熟练运用网络通道、套接字（Socket）以及远程过程调用（RPC）API 构建分布式通信组件。
        \item 掌握特定编程语言中基于共享内存模型的安全数据交换构件与内存屏障用法。
    \end{itemize}
\end{itemize}

\subsubsection{领域核心与拓展主题}
\begin{itemize}
    \item \textbf{高级通信属性与扩展机制（选择以下一项或多项）：}
    \begin{itemize}
        \item \textbf{通信介质深度特性}：
        \begin{itemize}
            \item 网络拓扑结构：单播、组播、信箱、交叉开关矩阵与自适应路由算法。
            \item 并发交互属性：消息定序、传递一致性保证、操作幂等性、计算与通信重叠技术。
            \item 可靠性工程：物理信道噪声、传输校验错误、丢包重传与滑动窗口机制。
            \item 数据封送与序列化：跨平台数据表示格式（Protobuf、JSON）、数据封送。
            \item 协议设计理论：活性保障、通信死锁规避。
            \item 安全性保障：传输保密性、数据完整性、端到端身份认证与授权机制。
            \item 性能评估维度：端到端时延、有效带宽与吞吐量、信道资源争用、系统响应活性。
            \item 排队论在信道拥塞与性能容量预测中的建模应用。
        \end{itemize}
        \item \textbf{通道高级控制与策略}：
        \begin{itemize}
            \item 通道管理策略：通信端点生命周期、会话保持、队列缓冲容量、饱和应对策略（阻塞等待 vs 丢弃丢包）、流量与速率控制。
            \item 发送控制（通常为过程式）与接收控制（通常为响应式回调或 RPC 服务派发）的编程对称性分析。
            \item 数据打包、序列化反序列化、输入合法性校验、传输加密与数据压缩。
            \item 海量慢速 I/O 设备与高并发连接下的多路复用与解复用机制：基于完成态驱动与调度器驱动的异步技术；async/await 异步流、select 与 poll/epoll 事件轮询 API。
            \item 进程通信形式化理论与代数分析（例如通信顺序进程 CSP 模型）。
        \end{itemize}
        \item \textbf{高级共享内存体系}：
        \begin{itemize}
            \item 现代语言与硬件内存模型：顺序一致性 (SC)与获取-释放一致性。
            \item 高并发内存管理：跨核心共享数据的安全内存回收机制；引用计数与替代方案（如 Hazard Pointers、基于历元 Epoch 的回收）。
            \item 块级数据布局与批量迁移：减少跨节点通信流量与提升局部性；计算与数据搬运的异步流水化重叠；数据内存布局对缓存命中的颠覆性影响（结构体数组 AoS 与 数组结构体 SoA 的性能对比）。
            \item 共享内存仿真：分布式共享内存（DSM）、远程直接内存访问（RDMA）。
        \end{itemize}
        \item \textbf{分布式数据存储与共识}：
        \begin{itemize}
            \item 一致性模型谱系：原子性、线性一致性、事务可串行化、缓存一致性、因果一致性、并发冲突消解、最终一致性以及区块链账本技术。
            \item 节点故障与网络分区：多数派法定人数表决；经典容错分布式共识协议（Paxos、Raft）。
            \item 安全性与零信任：拜占庭容错机制（BFT）、工作量证明（PoW）及新型权益证明共识协议。
        \end{itemize}
    \end{itemize}
\end{itemize}

\subsubsection{示范性学习成果与能力目标}
\begin{enumerate}
    \item \textbf{[评估 (计算机核心)]} 针对给定的工业应用场景与系统架构约束，全面论证并判定采用共享内存模型还是消息传递模型更具性能与工程实现优势。
    \item \textbf{[应用 (计算机核心)]} 编写一个经典的生产者-消费者程序（其中一个组件生成数据，另一个组件计算其移动平均值）；分别在小规模标量数据和大尺寸高精度数值场景下测量系统的实际加速比。
    \item \textbf{[开发 (领域核心)]} 编写一个能够将大规模数据集的不同分片并发派发给多个 Worker 计算节点并汇聚各局部运算结果的分布式程序（例如分布式超大数组求和）。
    \item \textbf{[开发 (领域核心)]} 编写一个并发数据汇聚程序，从多个异地远程站点并发拉取异构数据并利用某种形式的分布式规约输出汇总统计。
    \item \textbf{[评估 (领域核心)]} 编写基准测试，精确对比有缓冲队列与无缓冲同步通道在不同并发压力下的吞吐量表现与延迟瓶颈。
    \item \textbf{[评估 (领域核心)]} 针对给定的业务通信方案，全面评估其在防窃听、抗重放攻击与身份校验等维度是否具备充足的安全保障。
    \item \textbf{[分析 (领域核心)]} 构造一个包含数据竞态的并发内存访问代码序列，并证明该执行历史不满足顺序一致性。
    \item \textbf{[分析 (领域核心)]} 针对网络消息通信场景，构造一个由同步阻塞式发送操作相互等待而导致系统死锁的经典对抗案例。
    \item \textbf{[应用 (领域核心)]} 阐述至少一种在涉及多个互斥锁获取时杜绝死锁与活性故障的通用设计模式（如固定全局加锁顺序、两阶段锁）。
    \item \textbf{[开发 (领域核心)]} 编写一段实验程序，利用多线程交错执行直观展示硬件层面的内存乱序重排或消息异步乱序到达现象。
    \item \textbf{[评估 (领域核心)]} 深入剖析在不同数据更新写争用频率下，乐观并发控制（OCC）与悲观并发控制（PCC）的相对优劣与吞吐量临界点。
    \item \textbf{[分析 (领域核心)]} 举例构造一个由于高并发写冲突严重导致乐观更新尝试无限重试、甚至永远无法提交成功（活锁或饥饿）的极端场景。
    \item \textbf{[开发 (领域核心)]} 重构既有的并发软件系统，使其平滑迁移至具备更高可扩展性、更高容错性与更高可用性的现代分布式数据存储引擎之上。
\end{enumerate}

\subsection{协调与同步}
\textbf{【学时建议】}：2 计算机核心学时，6 领域核心学时。

本单元聚焦多个并发实体在执行时序上的相互制约与协同推进，涵盖任务依赖、推进活性、原子性互斥以及高级同步机制。

\subsubsection{计算机核心主题}
\begin{itemize}
    \item \textbf{协调与同步基础}：
    \begin{itemize}
        \item \textbf{依赖性动作的执行控制}：当某个活动的启动或推进依赖于其他活动时的执行调度与控制；基于完成态的控制（屏障 Barriers、线程汇合 Joins）；基于数据就绪驱动的控制（生产者-消费者设计模式）；基于条件满足的控制（忙等待轮询 Polling、重试 Retrying、指数退避 Backoffs、协同协助 Helping、挂起等待 Suspension、等待队列排队、条件信号唤醒 Signaling、超时时限控制 Timeouts）；响应式控制（激活并触发 Continuation 延续动作）。
        \item \textbf{推进活性与无死锁保障}：循环等待依赖与死锁；系统推进条件的单调性。
        \item \textbf{原子性语义与互斥锁}：底层硬件原子指令（CAS、LL/SC、Fetch-and-Add）；强制局部内存访问次序；互斥锁与临界区互斥；死锁防范工程技术（锁层次全局定序、锁粒度粗化、随机退避重试、利用锁管理器 Lock Manager 进行安全封装）；并发常见严重缺陷（遗漏加锁或未配对解锁、持有锁时调用不受控的未知外部操作、嵌套加锁死锁）。
    \end{itemize}
    \item \textbf{协调与同步编程实践}：
    \begin{itemize}
        \item 掌握并发组件优雅退出与协同终止控制机制。
        \item 熟练运用特定语言中的互斥锁、屏障、信号量等经典同步原语，在彻底杜绝数据竞态的同时保障系统的活性。
        \item 掌握在特定软件框架中利用事务性 API（如软件事务内存 STM 或分布式事务接口）实现高可靠并发协作。
    \end{itemize}
\end{itemize}

\subsubsection{领域核心与拓展主题}
\begin{itemize}
    \item \textbf{高级协调机制与控制流交互（选择以下一项或多项）：}
    \begin{itemize}
        \item \textbf{推进活性进阶理论}：非阻塞算法活性等级（无锁 Lock-free、无等待 Wait-free、无障碍 Obstruction-free）；调度公平性与优先级调度反转规避；同步机制与数据一致性、系统可靠性之间的权衡；性能瓶颈剖析（锁争用开销、锁粒度、锁护航 Convoying 效应、可扩展性极限）；典型无锁并发数据结构与算法实现（如 Treiber 栈、Michael-Scott 队列）。
        \item \textbf{原子性与资源控制进阶}：数据所有权与无锁借用语义；高级锁变体（顺序锁 SeqLocks、读写锁 Read-Write Locks、可重入锁、票据锁 Ticket Locks）；基于事务的并发控制（乐观与悲观事务执行）；分布式锁机制与租约可靠性管理。
        \item \textbf{与其他程序控制流形态的交互}：
        \begin{itemize}
            \item 屏障替代方案：物理时钟与逻辑时钟（Lamport 逻辑时钟、向量时钟 Vector Clocks、虚拟时间）；数据流模型与延续传递风格（CPS）；期约机制与异步远程调用；基于共识驱动的状态机复制；规约器与收集器结果汇聚机制。
            \item 投机执行、分支裁决与协同取消：投机执行对微架构安全与系统可观测性的潜在风险（如侧信道漏洞）。
            \item 基于系统资源的同步：经典计数信号量与条件变量。
            \item 控制流拓扑：细粒度任务调度器、带动态选主的串行-并行混合循环、流水线与流式处理管道、嵌套分治并行。
            \item 异常与故障协同响应：分布式异常传播与捕获处理、心跳监测超时探测、自愈容错与多数派故障仲裁机制。
        \end{itemize}
    \end{itemize}
\end{itemize}

\subsubsection{示范性学习成果与能力目标}
\begin{enumerate}
    \item \textbf{[应用 (计算机核心)]} 针对一段包含典型同步缺陷或数据竞态的并发代码，展示如何引入合理的互斥机制将其修复为线程安全程序。
    \item \textbf{[开发 (计算机核心)]} 编写一个协调程序，确保当且仅当一组并发异步子任务全部安全完成计算后，主控程序方能正常关闭并优雅终止。
    \item \textbf{[开发 (领域核心)]} 实现一个能够高吞吐、低延迟安全聚合海量并发传感器输入或网络数据包计数的高性能非阻塞计数器。
    \item \textbf{[开发 (领域核心)]} 分别运用不同的并发协调范式（如屏障汇聚法、Actor 消息归约法）完整实现 Filter-Map-Reduce 处理引擎。
    \item \textbf{[开发 (领域核心)]} 编写一段多阶段流水化并行程序，严格保障前一组并行任务全部汇聚提交之后，后一组并行任务才能被触发分发。
    \item \textbf{[开发 (领域核心)]} 编写一个采用投机执行机制的并发解题搜索程序，一旦其中某个工作线程率先检索到有效答案，能够立即广播信号并安全取消终止其他所有线程的无谓搜索。
    \item \textbf{[开发 (领域核心)]} 编写一段健壮的并行协同处理程序，当某个并发执行流突发致命数值异常（如除以零或浮点溢出）时，能向所有协作实体广播故障信号并协同安全回滚退出。
    \item \textbf{[评估 (领域核心)]} 编写一个使多个分布式对等节点就当前统一时钟达成共识的简易协议，并深入剖析其与工业级标准网络时间协议（NTP）相比在网络抖动和时钟漂移下的本质局限。
    \item \textbf{[开发 (领域核心)]} 实现一个轻量级高并发网络服务器，为每个接入的新客户端创建独立并发上下文（如轻量级协程或虚拟线程）处理并响应 Web 请求。
\end{enumerate}

\subsection{并发与分布式软件工程}
\textbf{【学时建议】}：2 计算机核心学时，3 领域核心学时。

本单元系统讲授并发与分布式软件在规约制定、质量属性保障、缺陷排查以及性能评估方面的软件工程方法论。

\subsubsection{计算机核心主题}
\begin{itemize}
    \item \textbf{安全性、活性与性能指标需求}：
    \begin{itemize}
        \item 运用时态逻辑基础概念严格界定系统的核心不变量：区分“坏事绝不会发生”与“好事终将发生”。
        \item 关键工程度量指标体系：吞吐量、响应性、端到端延迟、服务可用性、系统能耗与碳足迹、水平/垂直可扩展性、硬件资源利用率、通信带宽开销、排队等待时间与速率限制、调度公平性指标、工业服务等级协议（SLA）指标规范。
    \end{itemize}
    \item \textbf{违规识别、测试验证与缺陷修复}：
    \begin{itemize}
        \item 并发软件常见故障分类识别：时序定序缺失（竞态错误 Race Conditions）、原子性破裂（如典型的检查-执行竞争 Check-Then-Act 漏洞）、非终止循环（活锁 Livelock、死锁 Deadlock）。
    \end{itemize}
\end{itemize}

\subsubsection{领域核心与拓展主题}
\begin{itemize}
    \item \textbf{并发与分布式软件的规范化与评测（选择以下一项或多项）：}
    \begin{itemize}
        \item \textbf{形式化规约与建模}：串行规约向并发语义的拓展（如线性一致性 Linearizability 规范）；协议交互契约、会话类型与跨节点事务规范；借助工业级形式化工具（如 UML 交互图、TLA+ 规范语言、程序逻辑推演）验证并发协议的不变量；恶意节点与拜占庭行为下的安全性与活性证明；通信通道隔离性要求（防跨层侧信道信息泄漏）、入口流量过滤与反 DoS 速率限制。
        \item \textbf{静态分析与算法建模}：面向并发正确性、吞吐量瓶颈、时延抖动、资源争用与能耗分析的静态代码审查技术；基于 DAG 图模型的并行算法效能分析（总工作量 Work $T_1$、关键路径跨度 Span $T_\infty$、理论并行度 $T_1/T_\infty$ 与贪心调度边界）。
        \item \textbf{经验性评估与动态测试}：并发测试与调试专用工具链（数据竞态动态检测器如 ThreadSanitizer、模糊测试工具 Fuzzers、锁依赖环检测工具、高负载压力测试与严苛环境折磨测试 Torture Tests、集成至 CI/CD 流水线的并发自动化测试生成器）；经验性多维度量化评测（实测并深度对比吞吐瓶颈、同步额外开销、线程阻塞等待、通信流量开销、数据局部性破坏以及突发流量/海量并发连接冲击下的系统退化表现）；针对特定应用领域的专有评测分析与调优技术。
    \end{itemize}
\end{itemize}

\subsubsection{示范性学习成果与能力目标}
\begin{enumerate}
    \item \textbf{[开发 (计算机核心)]} 在不破坏既有系统核心业务语义、不变量与核心功能的前提下，重构现有软件规约以引入并行处理与分布式解耦。
    \item \textbf{[解释 (计算机核心)]} 深入阐述为何在并发与分布式软件架构设计中，严格规避“非局部副作用”是一条不可动摇的黄金设计法则。
    \item \textbf{[应用 (计算机核心)]} 针对一个批量并行计算过程，精确指明并形式化定义在每个并发计算步骤中必须严格维持的一组状态不变量。
    \item \textbf{[开发 (计算机核心)]} 编写一个高灵敏度的自动化测试程序，能够高概率稳定暴露潜在的数据竞态错误（例如两个并发活动在未加同步的情况下同时对共享累加器执行递增操作导致写丢失）。
    \item \textbf{[评估 (领域核心)]} 针对某在线服务突发超出设计承载能力的超大规模请求洪峰，精确制定系统优雅降级规约并实测其过载防护表现。
    \item \textbf{[开发 (领域核心)]} 准确诊断并重构消除某分布式应用中由串行执行瓶颈（Amdahl 定律瓶颈）导致的性能瓶颈，使其重获线性加速能力。
    \item \textbf{[评估 (领域核心)]} 借助现有的工业级测试基准脚手架，实测并全面对比针对同一接口设计的两种不同并发底层实现版本在多核高并发下的吞吐效能与锁竞争损耗。
    \item \textbf{[开发 (领域核心)]} 诊断并有效修复因跨节点网络通信过于频繁或数据读取延迟过高而导致的系统整体吞吐暴跌故障（如引入数据预取、批量聚合与异步流水化）。
    \item \textbf{[开发 (领域核心)]} 准确定位并优化并发系统中因过度的线程频繁创建销毁、锁争用严重或小对象内存分配而导致的系统性能抖动与资源开销膨胀。
    \item \textbf{[开发 (领域核心)]} 针对分布式系统中单点宕机、网络偶发丢包或超时等典型故障，精准诊断并修复系统的可靠性或高可用性脆弱点。
\end{enumerate}

\subsection{并行算法与应用领域}
\textbf{【学时建议】}：2 计算机核心学时，9 领域核心学时。

本单元系统梳理主流并行与分布式计算范式，并深入探讨经典并行算法在各个核心工业与科研领域的落地实践。

\subsubsection{计算机核心主题}
\begin{itemize}
    \item \textbf{并行与分布式算法的表达与实现}：熟练掌握在现代主流编程语言与开发框架中落地并行与分布式算法的核心技能，包括派生并发执行活动（如线程、纤程、Actor）、安全访问共享内存数据结构，以及熟练调用通道、套接字（Socket）与远程过程调用（RPC）网络接口。
    \item \textbf{核心计算范式与系统架构全景速查}：
\end{itemize}

\begin{table*}[htbp]
\centering\small
\begin{tabularx}{\textwidth}{p{0.14\textwidth} p{0.16\textwidth} Y Y p{0.18\textwidth}}
\toprule
\textbf{系统范式分类} & \textbf{典型执行实体} & \textbf{典型通信机制} & \textbf{典型算法领域} & \textbf{核心工程目标} \\
\midrule
\textbf{多核系统} & 操作系统线程、轻量级任务 & 共享内存、硬件原子指令、互斥锁 & 系统资源调度、高并发实时数据处理 & 极致吞吐量、低时延、高能效比 \\
\midrule
\textbf{反应式系统} & 事件处理器、异步回调函数 & 非阻塞 I/O 通道、事件总线 & 在线微服务网关、实时交互控制系统 & 超低响应时延、高可用性 \\
\midrule
\textbf{数据并行系统} & GPU 核心、SIMD 向量单元、AI 加速器 & 异构内存空间、显存 DMA 批量搬运 & 稠密线性代数、计算机图形渲染、大数据分析 & 极致算力吞吐、能耗效率 \\
\midrule
\textbf{计算集群} & 托管物理计算节点、MPI 进程 & 高速以太网/InfiniBand 套接字、消息通道 & 复杂科学数值仿真、大规模分布式图挖掘 & 大规模批处理吞吐量 \\
\midrule
\textbf{云计算平台} & 弹性容器、虚拟机、Serverless 函数 & RESTful/gRPC 服务接口、分布式消息总线 & 大规模 Web 互联网应用、海量用户在线服务 & 弹性伸缩性、故障自动隔离 \\
\midrule
\textbf{开放分布式系统} & 跨地域跨管理域自治节点、P2P 节点 & 网络套接字、分布式去中心化数据存储 & 拜占庭容错存储、区块链共识系统 & 强容错性、去中心化信任 \\
\bottomrule
\end{tabularx}
\end{table*}

\subsubsection{领域核心与拓展主题}
\begin{itemize}
    \item \textbf{八大核心算法应用领域（选择以下一项或多项）：}
    \begin{itemize}
        \item \textbf{科学计算与线性代数}：向量与高维矩阵运算（BLAS 算子级并行）；数值计算精度与数值稳定性控制；线性代数在现代大数据分析、深度神经网络大模型训练与推理中的并行加速应用。
        \item \textbf{数据处理与检索}：并行排序（如双调排序 Bitonic Sort、并行样本排序）、分布式哈希查找与倒排索引检索、高吞吐并发数据结构设计。
        \item \textbf{图算法、搜索与组合优化}：并行图遍历与节点标记、大规模边并行化处理、搜索分支剪枝与限界分析、投机搜索并发化、基于海量真实网络拓扑的分布式图计算（如 PageRank、连通分量求解）。
        \item \textbf{建模与物理系统仿真}：常微分与偏微分方程数值并行解法（有限元法、有限差分法）；蒙特卡洛随机模拟；高并发 N 体问题天体引力仿真；遗传算法与群体智能演化计算。
        \item \textbf{形式逻辑与约束求解}：大规模并行命题逻辑可满足性求解器、并发逻辑程序设计范式。
        \item \textbf{计算机图形学与计算几何}：大规模坐标仿射几何变换、图形光栅化渲染管线、实时硬件光线追踪加速、高精度网格碰撞检测。
        \item \textbf{算力与系统资源管理}：CPU/GPU 处理单元、内存带宽、通信通道与主机节点的动态分配、空间放置、内存回收与负载均衡调度；独占资源与共享资源的并发调度冲突消解；静态排程、动态调度与弹性伸缩算法；任务批量合并、优先级反转规避、资源隔离分区；基于工作窃取去中心化任务队列的高性能负载均衡机制。
        \item \textbf{网络与金融服务系统}：高并发分布式 Web API 网关架构、电子加密货币协议、高频交易引擎与分布式 ACID 事务处理系统、大规模多人在线交互系统（如 MMO 游戏）的分布式状态同步引擎。
    \end{itemize}
\end{itemize}

\subsubsection{示范性学习成果与能力目标}
\begin{enumerate}
    \item \textbf{[开发 (计算机核心)]} 基于已知的经典并行算法（如并行归并、分治前缀和求值、并行 BFS），使用现代并发库开发出高可靠的并发/分布式组件。
    \item \textbf{[开发 (计算机核心)]} 编写一个数据并行程序（例如利用 SIMD 指令或 OpenMP 指令集并行计算百万级浮点数数组的均值与方差）。
    \item \textbf{[开发 (计算机核心)]} 对既有的单线程事件驱动串行程序进行多线程化改造，在事件回调函数（如 GUI 界面交互事件或网络请求处理器）中派生异步工作线程，彻底消除主事件循环的卡顿与阻塞。
    \item \textbf{[开发 (计算机核心)]} 通过识别计算瓶颈并引入并行化或分布式扩展技术，显著提升某个性能受限串行软件模块的实际执行效能。
    \item \textbf{[评估 (计算机核心)]} 针对给定的实际系统工程诉求，在多核共享内存、异构加速器、集群多机部署与无服务器云函数等多元并行/分布式方案中进行严谨的架构选型与权衡论证。
    \item \textbf{[开发 (领域核心)]} 在前述八大算法领域中挑选其一，独立设计、实现、分析并评估一套针对特定物理与计算约束的高性能并发系统组件（例如：开发一个用于自动化仓储调度排程的并行遗传算法引擎）。
    \item \textbf{[评估 (领域核心)]} 深入审查并客观评判一个现存复杂并行组件（或由同学开发的并发系统），精准指出其在架构设计、同步机制、数据局部性或异常容错方面的深层次缺陷。
    \item \textbf{[评估 (领域核心)]} 针对同一计算基准算法（如大矩阵乘法），分别开发多核 CPU、分布式集群（MPI）以及 GPU（CUDA/OpenCL）三个不同版本的工程实现，严谨实测并全面对比三者在纯算力吞吐量与能耗效率维度的量化表现。
\end{enumerate}

\section{专业素养与思维倾向}

在并行与分布式计算领域的学习与工程实践中，重点培育学生的以下专业心智模型与素养特质：
\begin{itemize}
    \item \textbf{严密分析性与并发思维 (Analytical \& Concurrency Mindset)}：能够彻底跳出初学者习惯的单线程、单步线性确定性执行思维定式，建立交错时空、部分有序与状态空间爆炸下的非确定性并发系统全局观。
    \item \textbf{细致防御性与对部分失效的敬畏 (Meticulous \& Fault-tolerant Awareness)}：深刻理解并发竞态隐蔽性强、偶发性高以及分布式系统网络延迟不可预测、局部节点故障不可避免的现实规律，在接口设计与代码实现中始终秉持防卫性设计哲学，不心存侥幸。
    \item \textbf{权衡与工程折中意识 (Pragmatic \& Trade-off Oriented)}：深刻洞察分布式世界中的核心理论边界（如布鲁尔 CAP 定理、时延-一致性 PACELC 定理、阿姆达尔定律），能够在强一致性与高可用性、硬件成本与能耗极限、开发复杂度与运行效率之间做出清醒、理性的工程折中。
    \item \textbf{社会责任感与服务韧性考量 (Responsible \& Resilient)}：高度关注支撑现代社会运转的金融、交通、医疗与公共基础设施分布式系统的鲁棒性，时刻警惕单点故障与雪崩效应可能对公众生活造成的灾难性影响，坚守容错冗余、优雅降级与数据容灾的职业道德底线。
\end{itemize}

\section{数学基础要求}

并行与分布式计算依赖扎实的形式化数学工具与应用数学素养：
\begin{itemize}
    \item \textbf{离散数学与图论 (Discrete Mathematics \& Graph Theory)}：偏序集、哈斯图、拓扑排序；有向无环图（DAG，用于描述计算任务因果依赖图与工作流）；状态转移系统与自动机可达性分析。
    \item \textbf{数理逻辑与时态逻辑 (Logic \& Temporal Logic)}：命题逻辑与谓词演算；线性时态逻辑（LTL）与计算树逻辑（CTL），用于形式化表述与证明并发系统的安全性不变量（“坏事绝不发生”）与活性属性（“好事终将发生”）。
    \item \textbf{概率论、随机过程与排队论 (Probability \& Queuing Theory)}：离散与连续概率分布、随机变量期望与方差；泊松过程、马尔可夫链状态转移；排队论模型（M/M/1、Little 定律），用于定量分析通信信道排队延迟、吞吐瓶颈与缓存命中率。
    \item \textbf{线性代数与微积分 (Linear Algebra \& Calculus)}：向量与矩阵乘法、张量运算、特征值分解；多元微积分与偏微分方程离散化求解，支撑高性能数值计算与深度学习算子并行加速。
\end{itemize}

\section{交叉主题与跨学科共享概念}

\subsection{先修概念与依赖关系}
学习并行与分布式计算领域计算机核心知识的前提是熟练掌握以下先修概念：
\begin{itemize}
    \item \textbf{SDF (软件开发基础)}：程序静态代码与动态执行实体的区别；规约与实现的解耦；变量生命周期与内存表示；顺序控制流（条件分支、迭代循环）；过程/函数/方法调用机制与调用栈；数组与基础连续存储。
    \item \textbf{MSF (数学与统计学基础)}：命题逻辑；离散数学结构（有向图、状态转移机）；初等线性代数与微积分基础直觉。
    \item \textbf{FPL (编程语言理论基础)}：事件驱动架构；面向对象封装与多态；函数式编程范式、高阶函数与不可变数据模型。
    \item \textbf{SF (系统基础)}：分层抽象系统架构；冯·诺依曼计算机体系结构及其存储层次。
\end{itemize}

\subsection{知识领域交叉与重叠}
并行与分布式计算与计算机科学其他知识领域存在深度交织，部分拓展内容可在相关专业课中协同讲授：
\begin{itemize}
    \item \textbf{SF (系统基础)}：远程过程调用（RPC）底层机制、系统基准性能测试与多级存储层次评估。
    \item \textbf{NC (网络与通信)}：网络传输协议栈、Socket 编程接口、流量拥塞控制与拓扑路由。
    \item \textbf{OS (操作系统)}：内核级进程与线程调度、虚拟内存映射、系统调用陷阱、互斥锁与信号量实现、内核级故障隔离。
    \item \textbf{FPL (编程语言理论基础)}：并发内存模型形式化语义（如 Java/C++ 内存模型）、语言级并发原语、类型系统对数据竞态的静态杜绝。
    \item \textbf{SE (软件工程)}：高并发系统非功能性需求分析、微服务设计模式、并发故障回归测试。
    \item \textbf{SEC (网络空间安全)}：跨租户容器侧信道隔离、拜占庭容错、安全多方计算与区块链加密共识。
\end{itemize}

\section{课程体系规划与开课建议}

由于并行与分布式计算的计算机核心仅规定了 10 学时的通识要求，高校通常\textbf{不需要也不建议}单独开设一门仅覆盖 计算机核心 的极短必修课。联合工作组推荐以下三种典型开课与融合方案：

\subsection{方案一：全课程矩阵分散渗透方案}
将 10 学时的 计算机核心 知识有机拆解并注入至各核心骨干课程之中：
\begin{itemize}
    \item 注入《程序设计语言》：讲解函数式不可变性、并发闭包与轻量级线程启动；
    \item 注入《操作系统》：深化讲解多核进程/线程调度、互斥同步原语与死锁防范；
    \item 注入《计算机体系结构》：讲授多核缓存一致性协议（MESI）、伪共享与 GPU SIMD 向量化；
    \item 注入《计算机网络》与《数据库系统》：讲授 Socket 通信、RPC、事务 ACID、多版本并发控制（MVCC）与 Raft 共识。
\end{itemize}

\subsection{方案二：独立专业核心课程方案}
开设一门 48--64 学时的独立专业课程《并行与分布式计算系统》，完整统揽 计算机核心（10 学时）与 领域核心（26+ 学时），将多核编程、GPU 异构加速与分布式系统三大支柱融会贯通，配合高强度的工程实验（如基于 OpenMP/MPI 的并行矩阵加速、基于 Raft 的分布式键值数据库实现）。

\subsection{方案三：多核与高性能计算聚焦型课程范例}
针对偏向系统结构与高性能计算方向的高校，工作组提供了一套经典的多核导向开课样板：
\begin{enumerate}
    \item \textbf{程序与执行}：计算机核心 + 多核环境下的线程池、细粒度任务、指令级并行（ILP）。
    \item \textbf{通信机制}：计算机核心 + 现代多核微架构、多级 Cache 层次、伪共享规避、并发无锁数据结构。
    \item \textbf{协调与同步}：计算机核心 + 阻塞与非阻塞同步、投机执行、分治数据并行。
    \item \textbf{软件工程}：计算机核心 + 基于 DAG 的 Work-Span 跨度分析、性能调优剖析器实战。
    \item \textbf{算法与应用}：计算机核心 + 基于工业项目的并发流式数据处理与工作窃取任务调度器研发。
\end{enumerate}

\section{真实能力任务规范}

\begin{taskbox}[任务 1：基于已知算法实现高性能并行/分布式软件组件]
\textbf{【任务情境】}：针对已有明确数学或串行逻辑定义的核心算法（如大矩阵相乘、分治快速排序、PageRank 图计算或分布式哈希聚合），将其重构为能够在现代多核多节点环境中高效并发运转的高性能工程组件。\\
\textbf{【能力陈述】}：深刻理解跨越不同编程语言与底层计算框架（如 C++ 标准并发库、OpenMP、MPI、Java JUC 或 CUDA）的并行与分布式核心概念，并能准确、高质量地将其映射为鲁棒的代码实现。\\
\textbf{【能力领域】}：软件工程 \quad \textbf{【能力单元】}：系统开发\\
\textbf{【支撑知识单元】}：并行分布式/并行算法与应用、编程语言/并发与分布式计算\\
\textbf{【核心要求】}：计算机核心 \quad \textbf{【达标等级】}：\textbf{应用与开发}
\end{taskbox}

\begin{taskbox}[任务 2：通过引入并行或分布式架构突破串行系统性能瓶颈]
\textbf{【任务情境】}：某企业级核心业务系统面临数据量激增与高突发访问，原有单机串行组件成为全系统的严重性能瓶颈，亟需通过并行化或分布式横向扩展进行架构升级。\\
\textbf{【能力陈述】}：准确评估在何种场景下、通过何种并行或分布式技术能够有效提升吞吐量与降低延迟（并清醒识别不可并行化的阿姆达尔瓶颈）；精准捕捉并发化突破口，完成系统重构，并通过量化基准测试实测验证加速比与扩展收益。\\
\textbf{【能力领域】}：软件 / 系统 / 应用\\
\textbf{【能力单元】}：系统设计 / 开发 / 评估 / 性能优化\\
\textbf{【支撑知识单元】}：并行分布式/并发软件工程、编程语言/并发与分布式计算\\
\textbf{【核心要求】}：计算机核心 \quad \textbf{【达标等级】}：\textbf{评估与开发}
\end{taskbox}

\begin{taskbox}[任务 3：在保障系统关键语义与正确性前提下修订规约以支持分布式扩展]
\textbf{【任务情境】}：将传统基于单机强一致性假设的旧版软件规范拓展为支持网络分布式部署与异步并发处理时，必须全面修订软件架构规约与交互协议。\\
\textbf{【能力陈述】}：在松弛串行顺序约束或引入跨网络异步通信机制时，能够运用时态逻辑与不变量推理，严密论证并确保系统不会引入数据竞态、死锁、活锁或一致性破裂等破坏性隐患，确保原有系统核心功能与业务安全语义始终成立。\\
\textbf{【能力领域】}：软件 / 系统 / 应用\\
\textbf{【能力单元】}：架构设计 / 系统维护\\
\textbf{【支撑知识单元】}：并行分布式/并发软件工程、编程语言/并发与分布式计算\\
\textbf{【核心要求】}：计算机核心 \quad \textbf{【达标等级】}：\textbf{解释与评估}
\end{taskbox}

\begin{taskbox}[任务 4：针对复杂工程需求在多种并行与分布式架构方案间进行评估与选型]
\textbf{【任务情境】}：面对一个融合了海量批处理、实时低延迟在线查询与海量并发写入的复杂工业系统建设需求，需要在异构计算与分布式架构之间做出总体技术选型。\\
\textbf{【能力陈述】}：全面权衡多核共享内存、异构数据加速（GPU/SIMD）、反应式事件驱动架构、集群分布式计算与弹性云服务等不同范式的相对优劣、硬件成本、研发复杂度与吞吐极限，为具体的工程问题制定最优的系统架构选型方案。\\
\textbf{【能力领域】}：软件 / 系统 / 应用\\
\textbf{【能力单元】}：架构设计与方案评估\\
\textbf{【支撑知识单元】}：并行分布式/并行算法与应用、编程语言/并发与分布式计算\\
\textbf{【核心要求】}：计算机核心 \quad \textbf{【达标等级】}：\textbf{评估}
\end{taskbox}

\section{附录：核心主题与技能要求等级对照表}

\begin{longtable}{p{0.18\textwidth} p{0.48\textwidth} c c c}
\toprule
\textbf{知识单元} & \textbf{核心知识点与具体主题} & \textbf{认知等级} & \textbf{分类} & \textbf{学时} \\
\midrule
\endfirsthead
\toprule
\textbf{知识单元} & \textbf{核心知识点与具体主题} & \textbf{认知等级} & \textbf{分类} & \textbf{学时} \\
\midrule
\endhead
\bottomrule
\endfoot
\multirow{3}{*}{\makecell[l]{并行分布式/程序与执行}} 
& 并行动作定义、Happens-before 先行发生关系、数据独立性、一致性、物理执行位置与通信故障 & 解释 & 计算机核心 & 1 \\
& 并行动作的启动机制：编排放置、过程式派生、事件驱动反应式、依赖驱动启动 & 解释 & 计算机核心 & 1 \\
& 底层映射与物理执行机制：CPU 数据/指令级并行、GPU/SIMD 异构加速、多核调度、集群/云弹性编排与新兴量子/分子计算 & 解释 & 领域核心 & 2 \\
\midrule
\multirow{3}{*}{\makecell[l]{并行分布式/通信机制}} 
& 通信基础要素：介质（通道、共享内存、异构显存、数据存储）、内存一致性与缓存局部性、协同数据存储抽象 & 解释 & 计算机核心 & 2 \\
& 基于通信的并发编程：网络通道、套接字（Socket）、远程过程调用（RPC）及特定语言共享内存原语运用 & 开发 & 计算机核心 & 1 \\
& 高级通信属性与扩展：单播/组播/路由拓扑、排队论建模、高级通道流控/复用、顺序与获取-释放内存模型、DSM/RDMA 仿真、Raft/Paxos 容错共识 & 解释 & 领域核心 & 6 \\
\midrule
\multirow{3}{*}{\makecell[l]{并行分布式/协调与同步}} 
& 协调与同步基础：依赖动作控制（屏障、汇合、生产者-消费者）、死锁条件与推进活性、硬件原子指令与互斥锁 & 解释 & 计算机核心 & 2 \\
& 协调与同步编程实践：并发任务生命周期终止控制、锁与同步器防竞态实战、框架事务性 API 编程 & 开发 & 计算机核心 & 1 \\
& 高级协调扩展：无锁与无等待算法、读写锁/顺序锁变体、逻辑/向量时钟、投机执行与协同取消、信号量与条件变量、工作窃取任务调度 & 解释 & 领域核心 & 6 \\
\midrule
\multirow{3}{*}{\makecell[l]{并行分布式/软件工程}} 
& 软件需求规范：基于时态逻辑表达安全性与活性、吞吐量/时延/可用性/能耗与 SLA 度量指标 & 解释 & 计算机核心 & 1 \\
& 违规识别与修复：数据竞态错误、检查-执行原子性破裂、活锁与死锁排查修复 & 评估 & 计算机核心 & 1 \\
& 规范化与经验评估：线性一致性形式化规约、TLA+ 建模、DAG 工作量-跨度效率分析、ThreadSanitizer 竞态检测与压力测试 & 评估 & 领域核心 & 3 \\
\midrule
\multirow{3}{*}{\makecell[l]{并行分布式/算法与应用}} 
& 并行与分布式算法表达与实现：在主流语言与并发框架中落地线程派生、共享内存同步、套接字与 RPC & 开发 & 计算机核心 & 1 \\
& 核心计算范式与系统架构总览：多核共享内存、反应式、数据并行、集群、云计算、开放分布式系统的全景权衡 & 解释 & 计算机核心 & 1 \\
& 八大核心算法应用领域：并行线性代数、大规模数据处理、图算法与组合搜索、偏微分仿真、并行 SAT 求解、图形学渲染、工作窃取资源调度、高并发金融微服务 & 开发 & 领域核心 & 9 \\
\end{longtable}

\chapter{计算机安全}

\section{领域导言与演进背景}

\begin{preamblebox}[知识领域概述与变革]
现代社会对计算基础设施的依赖已深入交通运输、金融通信、医疗卫生、电网能源等国家级关键基础设施。近年来，针对这些关键基础设施的勒索攻击与灾难性数据泄露频发，这清晰地表明：当代计算机科学专业的毕业生必须承担起设计、构建与运维高韧性、内在安全且保护用户隐私的软件系统的重大责任。

在 CS2023 体系中，“计算机安全 (SEC)”既是一个立足计算本质的独立知识领域，又代表着贯穿于全书所有知识领域（SDF、DM、OS、NC、PDC、SF、AI、SEP）的底层横切基石。安全思维必须深度内化为每一位毕业生的职业直觉。

\textbf{自 CS2013 以来的核心变革：}
\begin{enumerate}\setlength{\itemsep}{0pt}\setlength{\parskip}{0pt}
    \item \textbf{重命名与学科视角重塑}：将 CS2013 的“信息保障与安全 (IAS)”升级为“计算机安全”。面对网络空间安全已发展为独立学科（基于 CSEC 2017 规范）的现实，CS2023 聚焦于纯粹的\textbf{计算机科学视角}——侧重软件系统构建、底层架构防御、密码学原理与形式化安全验证。
    \item \textbf{确立网络空间安全的六大核心跨界主题}：机密性、完整性、可用性、风险评估、系统性思维，以及最为核心的\textbf{对抗性思维}。
    \item \textbf{安全内生与隐私设计原则}：系统化引入“设计即安全”、“设计即隐私”、“纵深防御”以及“安全保障”形式化评估体系。
\end{enumerate}
\end{preamblebox}

\section{核心学时与跨领域映射}

由于安全思维高度渗透在各个领域，SEC 领域的学时采用创新的“双轨制”统计：
\begin{itemize}\setlength{\itemsep}{0pt}\setlength{\parskip}{0pt}
    \item \textbf{SEC 独立核心学时}：由本领域专设的《安全基石》所占用的基础学时（3 计算机核心 + 1 密码算法 计算机核心 = 4 计算机核心）。
    \item \textbf{跨领域渗透学时}：已在其他 15 个知识领域中开设的安全专题（累计达 29 个跨界学时，下表中带括号展示）。
\end{itemize}

\begin{table}[H]
\centering
\setlength{\tabcolsep}{3pt}\fontsize{8pt}{9.6pt}\selectfont\begin{tabularx}{\columnwidth}{Y c c}
\toprule
\textbf{知识单元} & \textbf{计算机核心} & \textbf{领域核心} \\
\midrule
安全基石与对抗思维 & 3 & -- \\
算法与社会 / 基础密码学 & (2) + 1 & 5 \\
计算机体系结构硬件安全 & 1 & 3 \\
数据与数据库安全及审计 & (2) & 3 \\
沉浸式与交互图形安全 & -- & 2 \\
人本安全与可用性工程 & (1) & 2 \\
数论与密码学数学理论 & (1) & (4) \\
网络协议安全与零信任 & (1) & 2 \\
操作系统内存防护与访问控制 & (3) & 4 \\
并发竞态防御与分布式容错 & (3) & 2 \\
隐私法理、安全策略与网络犯罪 & (4) & 4 \\
软件防御性编程与安全重构 & (5) & 7 \\
专用平台生态与系统沙箱隔离 & (4) & 6 \\
\midrule
\textbf{SEC 领域专设净核心学时} & \textbf{4 学时} & \textbf{45 学时} \\
\textbf{全书贯穿安全总学时池} & \multicolumn{2}{c}{\textbf{33 CS 核心学时 + 45 领域核心学时}} \\
\bottomrule
\end{tabularx}
\end{table}

\section{核心知识单元大纲}

\subsection{现代计算机安全基石}
\textbf{【学时建议】}：3 计算机核心学时。

\begin{itemize}\setlength{\itemsep}{0pt}\setlength{\parskip}{0pt}
    \item \textbf{网络安全跨界元模型}：CIA 三元组（机密性、完整性、可用性）、真实性、不可抵赖性；系统性思维与对抗性思维。
    \item \textbf{威胁建模与攻击向量 (Threat Modeling \& Attack Vectors)}：STRIDE 威胁建模框架、攻击树、通用漏洞披露（CVE 与 CVSS 漏洞评分标准）。
    \item \textbf{身份鉴别与访问控制 (AAA)}：多因素身份认证（MFA）、自主访问控制（DAC）、强制访问控制（MAC）、基于角色的访问控制（RBAC）与基于属性的访问控制（ABAC）。
    \item \textbf{经典安全体系设计原则}：
    \begin{itemize}\setlength{\itemsep}{0pt}\setlength{\parskip}{0pt}
        \item 最小特权原则。
        \item 开放设计原则（Open Design / Kerckhoffs 原则：安全不应依赖隐匿）。
        \item 故障安全默认与全景检查。
        \item 纵深防御与隔离最小化机制。
    \end{itemize}
    \item \textbf{设计权衡冲突}：安全防护、用户隐私、执行性能与交互可用性之间的多维博弈与妥协。
    \item \textbf{可信计算与安全保障}：可信基（TCB）、硬件信任根、形式化安全属性验证。
\end{itemize}

\subsubsection{示范性学习成果}
\begin{enumerate}\setlength{\itemsep}{0pt}\setlength{\parskip}{0pt}
    \item \textbf{[解释]} 阐明最小特权原则、Kerckhoffs 原则以及对抗性思维在现代系统设计中的决定性价值。
    \item \textbf{[开发]} 针对包含外部公开接口的网络服务，绘制 STRIDE 威胁模型图，设计兼顾纵深防御与访问控制的多层隔离架构。
    \item \textbf{[评估]} 分析某历史重大安全事故（如 Log4j 漏洞或投毒供应链事件），评估其在认证、授权及输入验证阶段的失陷原因，提出系统化加固方案。
\end{enumerate}

\section{专业素养与思维倾向}

\begin{itemize}\setlength{\itemsep}{0pt}\setlength{\parskip}{0pt}
    \item \textbf{极致细密与审慎性}：安全无小事。在真实攻击中，哪怕一行未校验的输入边界或配置疏漏都将导致整套防御全盘瓦解。
    \item \textbf{自主追踪与终身学习}：黑客攻防对抗处于高度动态演进中，毕业生必须持续主动追踪最新零日漏洞与防御范式。
    \item \textbf{协同作战能力}：面对突发高烈度网络渗透与勒索攻击，善于跨部门紧密联动以实现秒级应急响应。
    \item \textbf{高度责任担当 (Responsible \& Accountable)}：对数据资产安全与公民个人隐私保持敬畏心，对编写的每一行代码承担伦理与工程责任。
\end{itemize}

\section{数学基础要求}

\begin{itemize}\setlength{\itemsep}{0pt}\setlength{\parskip}{0pt}
    \item \textbf{离散结构与近世代数}：群论、有限域（Galois Field $GF(2^n)$ 在 AES 对称加密中的代数基础）。
    \item \textbf{初等数论}：模幂运算、欧几里得扩展算法、费马小定理与欧拉函数（RSA 及椭圆曲线 ECC 基础）。
    \item \textbf{概率论与数理统计}：生日碰撞概率（在哈希抗碰撞中的应用）、伪随机数生成器检验。
\end{itemize}

\section{课程设计与教学实施建议}

联合工作组建议高校结合自身资源采取“双轨并行”方案：
\begin{enumerate}\setlength{\itemsep}{0pt}\setlength{\parskip}{0pt}
    \item \textbf{全谱系渗透模式}：在不增设计算机安全必修课的前提下，必须强制要求所有专业核心课严格落实安全教学：在《数据库》中讲注入与列加密，在《操作系统》中讲提权与栈溢出防御，在《网络》中讲 TLS 与 BGP 劫持。
    \item \textbf{独立综合课程模式}：开设一门 45 学时的《计算机系统安全与对抗工程》，集中讲授密码学工程应用、内存攻击防御（ASLR/DEP/ROP 链）、Web 全栈攻防、云原生与零信任网络。
\end{enumerate}

\section{真实能力任务规范}

\begin{taskbox}[任务 1：核心业务系统全面威胁建模与对抗性渗透防御分析]
\textbf{【任务情境】}：公司计划上线一套面向全球用户的开放金融开放平台 API 网关，面临复杂的外部黑客嗅探、重放与注入风险。\\
\textbf{【能力陈述】}：运用对抗性思维与系统化威胁分析方法（STRIDE），对网关暴露的所有网络端点与消息载荷进行穿透式威胁建模，设计防重放令牌、非对称验签、基于角色的最小特权 RBAC 及敏感字段数据脱敏架构，并输出专业级系统安全防护白皮书。\\
\textbf{【能力领域】}：软件工程与系统安全 \quad \textbf{【能力单元】}：系统设计与安全验证\\
\textbf{【达标等级】}：\textbf{评估与开发}
\end{taskbox}

\section{委员会专家名单}
\begin{itemize}\setlength{\itemsep}{0pt}\setlength{\parskip}{0pt}
    \item \textbf{主席}：Rajendra K. Raj (Rochester Institute of Technology)
    \item \textbf{委员}：Vijay Anand (Univ. of Missouri-St. Louis)、Diana Burley (American Univ.)、Sherif Hazem (Central Bank of Egypt)、Michele Maasberg (US Naval Academy)、Sumita Mishra (RIT)、Nicolas Sklavos (Univ. of Patras)、Blair Taylor (Towson Univ.)、Jim Whitmore (Distinguished IT Architect)
\end{itemize}

\chapter{社会、伦理与专业素养}

\section{领域导言与演进背景}

\begin{preamblebox}[知识领域概述与时代演进]
虽然技术议题始终是计算机课程体系的核心，但在更广阔的社会与人类文明语境中，仅有纯技术教学并不能构成一个完整的计算机高等教育体系。学生必须深入理解计算技术所植根的宏观社会背景，从而深刻把握相关的社会影响、伦理规范、法律法规及职业责任。将这些“非技术”议题系统性纳入计算机教学大纲的迫切需求，早在 1991 年 ACM/IEEE-CS 联合发布的《CS1991》课程指南中就已获得正式确认：
\begin{quote}
\small\songti
“本科生必须理解计算机学科所固有的基本文化、社会、法律与伦理问题。他们应当理解该学科的过去、现状与未来走向；理解个人在这一历史进程中的角色；并领会那些在学科发展中起到关键作用的哲学思辨、技术挑战与审美价值。

学生还需要培养针对计算技术的社会影响提出严肃质询、并批判性评估可能解答的能力。未来的从业者必须能够前瞻性地预见将某项产品引入特定环境所带来的冲击：该产品是会提升还是降低人类的生活质量？它对个人、群体乃至社会制度将产生何种深远影响？

最后，学生需要了解软硬件供应商和用户的基本法定权利，并理解构成这些权利基石的伦理价值观。未来的从业者必须清醒认识到自己肩负的社会责任，以及技术失败可能招致的严重后果。他们必须认识到自身的局限性以及所用工具的边界。所有从业者都必须做出长期承诺，在所选专业方向和整个计算机学科中保持终身学习与前沿跟进。”\footnote{引自 ACM/IEEE-CS Joint Curriculum Task Force, \textit{Computing Curricula 1991} (CS1991), ACM Press \& IEEE Computer Society Press.}
\end{quote}

随着现代信息技术的跨越式突破深刻重塑人类的生活、生产与社会治理模式，社会与伦理议题及专业职业实践的重要性和深远后果日益凸显；新型计算平台、大模型及自主决策系统每年都在提出更具挑战性的严峻难题。正是当下的学生，未来将带着审慎与担当步入工业界与学术界，承担起识别与化解这些复杂难题的崇高使命。

\textbf{课程实施路径与教学法建议：}\\
计算机教育界在讲授本领域内容时通常采取多元化模式：独立开课、在传统技术与理论课程中深度渗透、设立高阶专属研讨课、二者结合的混合模式，或作为毕业设计、工程实训和职业素养实践课程的专项模块。

本指导委员会强烈推荐\textbf{“一门必修独立核心课程 + 全程贯穿技术课程渗透”}的复合开课方案：
\begin{itemize}
    \item 一方面，本领域被列为“计算机核心”的诸多基础主题（如伦理分析哲学框架、知识产权法理、反垄断与科技经济学等）通常很难直接塞入传统算法、操作系统或编译原理等硬核技术课中。缺乏一门独立的专业核心课程，极难保证学生获得系统、完整和严谨的思辨训练。
    \item 另一方面，如果社会、伦理与职业准则仅局限于一门独立课程讲授，而在后续技术课程中绝口不提，将强化一种极度危险且错误的认知假象——即“纯粹的技术实现与伦理考量毫无瓜葛”。
\end{itemize}
因此，必须在软件工程、数据库、计算机图形学、计算机网络、系统安全、人工智能及计算机导论等多门传统技术骨干课中，以真实工程案例分析的形式，自然融入与该课程紧密相关的伦理、法律与社会思辨。例如：
\begin{itemize}
    \item 在数据库课程中探讨海量数据聚合、用户画像与数据挖掘带来的隐私侵蚀与知情同意边界；
    \item 在软件工程课程中剖析对商业客户的交付履约承诺与对终端用户、公共利益福祉之间可能产生的尖锐利益冲突；
    \item 在编程大作业中引入诸如“由非专业医师操作的高精度激光眼科手术机器人运动控制算法”等高可靠性案例，直观展示代码失误可能造成的生命伦理代价。
\end{itemize}
若专业教师对应用伦理学的教学内容与方法感到生疏，强烈建议充分利用 ACM、IEEE-CS、ACM 计算机与社会专业委员会（SIGCAS）等权威机构提供的丰富公开教学资源与案例库。

\textbf{伦理分析是贯穿所有分支的灵魂基准：}\\
伦理分析的实践构成了“社会、伦理与专业素养 (SEP)”知识领域每一个分支的底层骨架。《ACM 伦理与职业行为准则》（ACM Code of Ethics and Professional Conduct）\footnote{\url{https://www.acm.org/about/code-of-ethics}}、《IEEE 伦理准则》（IEEE Code of Ethics）\footnote{\url{https://www.ieee.org/about/corporate/governance/p7-8.html}}与《AAAI 伦理与职业行为准则》（AAAI Code of Ethics and Professional Conduct）\footnote{\url{https://aaai.org/Conferences/code-of-ethics-and-conduct.php}}共同提供了指导职业工程实践的行动规范与典型判例。其中的“一般道德律令”明确了计算机专业人员对个人责任、职业行为操守以及引领技术向善的责任基准。计算机专业教师应注重启发学生在计算机具体技术情境中的专业思辨，同时积极借助人文社科领域的优质通识课程形成学科交叉合力。

\textbf{自 CS2013 以来的核心重大变革：}
\begin{enumerate}
    \item \textbf{重磅增设“多元、公平与包容”独立核心知识单元}：自 2013 年以来，国际计算学界对正义、公平、多元和包容的认知与行动发生了深刻变革。计算机专业不仅必须为所有学生打造公平、多元与包容的受教育环境，更要求未来的计算机专业人才在投身工程实践时，确保其研发的软件、算法与系统以公平、公正与包容的方式造福全人类。为此，CS2023 正式将其升格为全领域的独立必修核心知识单元。
    \item \textbf{全面直面生成式 AI 与算法决策的社会伦理挑战}：系统涵盖大语言模型、生物特征识别、深度伪造、自动化行政裁量系统、预测性警务等新型技术在公民权利、偏见放大、劳动力替代和认知操纵方面的社会伦理冲击。
    \item \textbf{强调计算可持续性与绿色计算的紧迫性}：将硬件制造资源消耗、关键矿物提取、数据中心巨大能耗、区块链共识开销及全球电子废弃物危机纳入工程决策核心考量。
\end{enumerate}
\end{preamblebox}

\section{核心学时分配建议}

本知识领域建议的核心学时分配如下表所示。根据“向日葵模型”，全领域共涵盖 \textbf{17 个计算机核心学时}与 \textbf{14 个领域核心学时}，累计基准学时为 31 学时：

\begin{table}[H]
\centering
\setlength{\tabcolsep}{3.5pt}\fontsize{8pt}{9.6pt}\selectfont\begin{tabularx}{\columnwidth}{Y c c}
\toprule
\textbf{知识单元} & \textbf{计算机核心} & \textbf{领域核心} \\
\midrule
社会语境 & 3 & 2 \\
伦理分析方法 & 2 & 1 \\
职业伦理与专业准则 & 2 & 2 \\
知识产权与开源生态 & 1 & 1 \\
隐私与公民权利 & 2 & 1 \\
专业沟通与团队协作 & 2 & 1 \\
可持续性与绿色计算 & 1 & 1 \\
计算机发展简史 & 0 & 1 \\
计算经济学与产业生态 & 0 & 1 \\
安全策略、法律与网络犯罪 & 2 & 1 \\
多元、公平与包容 & 2 & 2 \\
\midrule
\textbf{合计} & \textbf{17 学时} & \textbf{14 学时} \\
\bottomrule
\end{tabularx}
\end{table}

\section{知识单元与详细大纲}

\subsection{社会语境}
\textbf{【学时建议】}：3 计算机核心学时，2 领域核心学时。

计算机与互联网技术在过去数十年间以超越以往任何工业发明的速度深刻重塑了整个人类社会：人类劳动生产率实现跃升，新闻资讯、大众娱乐与社交沟通渠道全面爆发，几乎在所有自然科学与工程技术领域均催生出基础性突破。社会语境单元构成了理解 SEP 领域其他所有知识单元（尤其是职业伦理）的认知根基。

\subsubsection{计算机核心主题}
\begin{itemize}
    \item \textbf{超网络化与人工智能高速演进世界的社会影响}：在人机高度互联、生成式 AI 与自主智能体加速演进的社会形态下，计算技术对劳动生产、组织协作与人类生存方式的根本性重塑。
    \item \textbf{社交媒体与人工智能对个体与意识形态的冲击}：算法推荐机制与生成内容对个体心理福祉、政治意识形态极化、信息茧房及文化认同的复杂塑造。
    \item \textbf{前沿计算技术融入公民生活与公共治理的影响}：人工智能、生物特征识别技术与算法决策系统在公民生活与行政监管中的应用剖析（如人脸识别系统、生物特征追踪标签、社会公共资源配给算法、预测性警务软件等）。
\end{itemize}

\subsubsection{领域核心主题}
\begin{itemize}
    \item \textbf{互联网、算力基础设施与人工智能的发展与控制权}：大型科技垄断巨头、跨国数据基础设施平台对公共话语场及技术治理规则的控制。
    \item \textbf{数字鸿沟及其社会外溢效应}：数字基础设施与计算资源获取能力在性别、阶层、种族、地理分布及发展中国家之间的显著失衡及其深远社会后果。
    \item \textbf{信息无障碍与暗黑设计模式}：无障碍技术的法律强制规范；以及利用心理认知陷阱诱导消费、窃取数据或阻碍注销的暗黑模式。
    \item \textbf{情境感知计算}：普适计算中设备基于环境与行为自动推断上下文所引发的社会介入与道德界限。
\end{itemize}

\subsubsection{示范性学习成果}
\begin{enumerate}
    \item \textbf{[解释]} 阐明计算机技术（计算机网络、移动计算、云计算平台）在个人交往与社会群体层面上调控和媒介人类互动的多重机制。
    \item \textbf{[解释]} 识别硬件与软件系统设计中潜藏的开发者价值预设与先验偏见，重点剖析其对边缘化群体和残障人士可用性造成的系统性障碍。
    \item \textbf{[评估]} 深度解读并阐释具体技术方案与其软硬件实现的社会语境关联。
    \item \textbf{[评估]} 运用实证经验数据，定量与定性评估特定技术设计与系统实现的社会有效性及副作用。
    \item \textbf{[解释]} 深刻阐明社交媒体算法对不同身份族群、文化群体与社会社群造成的结构性影响。
    \item \textbf{[解释]} 深入研讨互联网在促进公民、政府机构及社会各主体间沟通与互动中的赋能机制与异化风险。
    \item \textbf{[评估]} 深入剖析民主制度运行（如公共社会福利发放、电子投票系统）对计算机基础设施的依赖所带来的机遇与风险。
    \item \textbf{[解释]} 阐述历史上弱势少数族裔与女性在计算机专业领域代表性严重不足对行业文化和产品多元性造成的负面影响。
    \item \textbf{[解释]} 解释泛在普适计算系统中情境感知特性的社会伦理意涵。
    \item \textbf{[评估]} 评估不同社会形态下获取互联网和计算技术的权利鸿沟；论证为何以及如何在国际法理上将“互联网访问权”视作一项基本人权。
\end{enumerate}

\subsection{伦理分析方法}
\textbf{【学时建议】}：2 计算机核心学时，1 领域核心学时。

伦理学理论与道德哲学原则是进行理性伦理分析的基石，为复杂决策困境中的技术人员提供指引路径与思辨视角。不同的伦理学流派基于不同的哲学假设和论证方法来判定特定行为的合道德性。学生必须清醒认识到：在多元化的现实语境中，解决不同冲突可能需要借助不同的伦理学框架方能获得最佳方案；而何谓“最佳”，往往由特定语境与文化视角共同决定。应用伦理分析方法既要求深刻理解各分析工具背后的哲学原理与预设，也要求敏锐洞察决策所处的具体社会语境。源自西方哲学的传统伦理框架固然实用，但绝非放之四海而皆准；教育中必须积极引入去殖民化、本土原住民与历史上被边缘化群体的伦理视角。

\subsubsection{计算机核心主题}
\begin{itemize}
    \item \textbf{论证逻辑与逻辑谬误规避}：识别并规避技术与道德辩论中常见的非形式逻辑谬误（如偷换概念、诉诸人身、虚假对立、稻草人谬误）与误导性陈述。
    \item \textbf{伦理学理论与技术决策框架}：功利主义（后果论/效益主义）、义务论（康德道义论）、美德伦理学、关怀伦理学及社会契约论在计算技术决策中的推演与应用。
    \item \textbf{文化在技术认知中的决定性作用}：深刻理解文化背景如何塑造人们对计算技术的理解、接纳程度、产品设计与日常使用行为。
\end{itemize}

\subsubsection{领域核心主题}
\begin{itemize}
    \item \textbf{职业伦理自查清单}：工业界工程发布与算法审查中的伦理核查清单设计。
    \item \textbf{多维度评估量规}：技术方案社会效益与道德风险的量化与定性综合评估量规。
    \item \textbf{利益相关者分析}：全链路识别直接、间接利益相关者，剖析权责不对等与潜在受害者。
    \item \textbf{立场理论}：从被边缘化或受冲击最严重的弱势受众立场审视技术系统的内在缺陷。
\end{itemize}

\subsubsection{示范性学习成果}
\begin{enumerate}
    \item \textbf{[解释]} 识别并系统阐述特定文化语境与价值体系对技术工程决策产生的深刻制约与塑造。
    \item \textbf{[应用]} 熟练运用典型工程判例与哲学类比，开展严谨的技术伦理推导与论辩。
    \item \textbf{[应用]} 批判性剖析技术报告与公众论辩中的论述，准确识别并剔除常见的基础逻辑谬误。
    \item \textbf{[应用]} 严密解析复杂技术伦理争议，精准提取论辩的核心前提假设与推论逻辑链条。
    \item \textbf{[评估]} 全面评估在特定情境中所有利益相关者立场与其所处文化背景的复杂交互关系。
    \item \textbf{[开发]} 评估并量化技术决策中引入或积累“伦理债务”（Ethical Debt，即为了追求短期快速交付而推迟对伦理影响进行严谨评估所积攒的系统性潜在风险）的危害与应对策略。
\end{enumerate}

\subsection{职业伦理与专业准则}
\textbf{【学时建议】}：2 计算机核心学时，2 领域核心学时。

计算机职业伦理是应用哲学的一个重要分支，专门研究计算机从业者在专业技能与社会行为中应当遵循的决策规范。个人的职业决策主要受三重力量影响：1) 从业者个人的道德自律准则；2) 工作单位、行业协会或职业认证制度所要求的显性/隐性行为规范；3) 国际计算专业学会制定的正式职业道德与行为守则。

\subsubsection{计算机核心主题}
\begin{itemize}
    \item \textbf{社群公共价值与法律规制体系}：社会主流伦理准则、公共利益与现行法律法规的互动与统一。
    \item \textbf{专业职业素养的内涵实质}：专业能力、审慎尽责、受托人信托责任与学术/技术传承指导。
    \item \textbf{终身专业进阶与技能更迭}：保持对新兴开发工具、技术框架、法律规制与行业标准的动态跟进，具备批判性自我评估与持续职业学习能力。
    \item \textbf{专业协会职业道德准则与行为规范}：深入研习 ACM/IEEE-CS《软件工程职业道德准则和实践规范》（SE Code）、AITP、IFIP 及各国计算机学会的准则条款。
    \item \textbf{问责机制、责任归属与法律赔偿责任}：软件正确性、高可靠性与系统安全性的法律责任认定；网络安全工程师对数据隐私和防御失职的伦理保密边界。
    \item \textbf{技术哲学与社会学理论导论}：理解人类创造与使用技术的经典理论视角，包括工具主义、技术系统社会学、无障碍平权与障碍正义、技术中立论假说、实用主义、功利主义及去殖民化批判理论。
    \item \textbf{违规与道德风险的识别与举报机制}：制定策略与规范流程，敏锐识别并依法合规向上级或公众报告违反法律或职业道德准则的系统设计、软件漏洞或职业渎职行为。
\end{itemize}

\subsubsection{领域核心主题}
\begin{itemize}
    \item \textbf{计算从业人员在公共政策制定中的角色}：参与国家技术立法、标准制订、公众科普与反垄断咨询。
    \item \textbf{技术不良后果的动态前瞻与预警}：建立针对算法和系统次生灾害的持续监控与警觉机制。
    \item \textbf{伦理异议与“吹哨人”保护}：技术异议的合法提交通道、道德良知冲突下的合法检举制度及从业人员法律保护。
    \item \textbf{区域文化与伦理困境的冲突与调适}：跨国跨文化软件开发中遭遇的道德标准冲突（如言论边界、隐私感知差异）。
    \item \textbf{职场骚扰与歧视应对机制}：应对工作场所性别歧视、种族微歧视与职场霸凌的合规干预与维权途径。
    \item \textbf{职业认证与执业资质制度}：注册工程师与各专业认证的作用与局限。
    \item \textbf{工作场所计算资源的可接受使用策略 (AUP)}：雇主监控、个人隐私与企业资产安全防范边界。
    \item \textbf{人机工效学与健康计算环境}：职业性肌骨劳损、视力损害、心理重压与工效学健康标准。
    \item \textbf{上市时间、成本控制与工程质量标准的冲突}：面对商业“极速上线交付”与严格软件工程质量及安全标准冲突时的职业坚守。
\end{itemize}

\subsubsection{示范性学习成果}
\begin{enumerate}
    \item \textbf{[解释]} 敏锐识别软件需求分析、架构设计、编码实施与生产部署全流程中潜在的伦理陷阱。
    \item \textbf{[评估]} 综合权衡多方制约因素，制定妥善化解工程实践中复杂伦理困境的切实方案。
    \item \textbf{[解释]} 阐释确保软件正确性、高可靠性与人身系统安全背后的伦理与法律责任来源（如 ACM/IEEE 准则、行业强制标准、组织安全策略等）。
    \item \textbf{[解释]} 梳理专业从业者在职业生涯中持续跟踪道德准则、法律动态与最佳实践的机制与途径。
    \item \textbf{[评估]} 客观剖析相关职业道德准则的现实效力与内在局限性，评价其作为职业决策行动指南的优势与不足。
    \item \textbf{[评估]} 深入分析一个全球性计算伦理事件（如关键基础设施崩溃、重大数据泄露或波音 737 MAX 软件故障），系统剖析工程技术人员与监管官员在处置应对中的权责失衡。
    \item \textbf{[解释]} 阐明软件系统创造背后的哲学基础，及其如何决定我们选择以何种方式应用这些软件。
    \item \textbf{[解释]} 归纳计算专业从业者向政府公共政策建言献策的途径与有效形式。
    \item \textbf{[解释]} 阐明违反职业操守（如数据造假、蓄意留后门、知识产权侵权）对个人职业生涯、企业公信力及公共利益造成的灾难性后果。
    \item \textbf{[解释]} 熟知“吹哨人”道德检举机制，掌握在遭遇重大公共危害隐患时进行合规内部申诉与外部合法举报的行动指南。
    \item \textbf{[解释]} 列举并剖析区域文化差异与跨国软件工程伦理困境交互影响的真实工业案例。
    \item \textbf{[解释]} 探讨职场骚扰与歧视的各种隐蔽表现形式，并列举受害者获得组织援助与法律救济的合规通道。
    \item \textbf{[评估]} 考察国内外各类专业工程师执业资格认证制度在保障工程公众安全方面的实际效能。
    \item \textbf{[解释]} 阐释工作场所人机工效学设计与从业人员生理、心理健康之间的严密内在联系。
    \item \textbf{[解释]} 剖析工业界迫于“产品极速抢跑上市”商业压力，而妥协退让甚至践踏专业工程质量标准的严重弊端。
\end{enumerate}

\subsection{知识产权与开源生态}
\textbf{【学时建议】}：1 计算机核心学时，1 领域核心学时。

知识产权 (IP)是指人类在智力劳动成果（如计算机软件程序、算法创新、电路芯片设计等）上依法享有的无形专有权利体系。各国法律依据知识产权的具体客体属性提供不同模式的法定排他保护。在现代软件工程中，主要涉及四种截然不同的知识产权类型：专利权、版权/著作权、商业秘密以及商标权。每种权利所保障的法益对象、保护期限与排他效力各不相同。

\subsubsection{计算机核心主题}
\begin{itemize}
    \item \textbf{知识产权的哲学与法理基石}：洛克劳动财产权学说、人格权理论与功利主义激励理论对智力垄断与社会共享平衡的论证。
    \item \textbf{知识产权排他权利体系}：专有使用权、许可转让权、署名权、修改权及作品完整权保护。
    \item \textbf{无形数字知识产权 (IDIP)}：纯软件代码、数字化媒体资产、专有算法逻辑与虚拟资产的无形资产确权特征。
    \item \textbf{知识产权保护的法定框架与国际条约}：伯尔尼公约、TRIPS 协定、DMCA 法案以及各国软件著作权登记与专利审查标准。
\end{itemize}

\subsubsection{领域核心主题}
\begin{itemize}
    \item \textbf{数字版权管理 (DRM)}：技术保护措施（TPM）、加密水印与反规避条款的法理与技术博弈。
    \item \textbf{四项核心权利的边界与对比}：软件专利、软件著作权、商业秘密（如专有推荐算法源码）与商标的区别与联动防御策略。
    \item \textbf{学术与技术剽窃}：代码抄袭判定机制、开源合规审查与学术不端防范。
    \item \textbf{开源运动与开源许可协议族}：自由软件基金会（FSF）与开源倡议组织（OSI）理念；宽松型许可证（MIT、Apache 2.0、BSD）与传染型许可证（GPLv2/v3、AGPL）的法律约束与商业合规。
    \item \textbf{软件盗版与网络侵权}：全球软件盗版的经济成因、暗网盗版分发链条与反侵权救济。
\end{itemize}

\subsubsection{示范性学习成果}
\begin{enumerate}
    \item \textbf{[解释]} 结合具体国家或区域法域，深入辨析知识产权保护的哲学伦理基础（激励创新 vs 知识垄断弊端）。
    \item \textbf{[解释]} 系统阐明旨在遏制数字版权侵权行为的现行国际与国内法律法规。
    \item \textbf{[评估]} 批判性审视和评价数字版权相关法案在防范侵权与保障公众合理使用之间的失衡争议。
    \item \textbf{[应用]} 准确识别当今数字化前沿场景中涌现的各种无形数字知识产权形态（如 AI 生成物、数据集所有权、NFT）。
    \item \textbf{[评估]} 在具体工程研发情境下，合法合理论证对受版权保护素材进行“合理使用”的法律抗辩依据。
    \item \textbf{[评估]} 评估学术界与工业界各类自动化代码抄袭检测机制所引发的技术误判与潜在隐私伦理争议。
    \item \textbf{[应用]} 精准解读主流开源许可协议（如 GPL、Apache、MIT）的核心条款约束，评估其在商用产品集成中的合规风险。
    \item \textbf{[评估]} 全面权衡为软件算法申请专利保护与作为商业秘密保护之间的利弊得失与法律壁垒。
    \item \textbf{[解释]} 严谨区分并对比著作权、专利权、商业秘密和商标在保护计算机软件产品时的法律维度差异。
    \item \textbf{[解释]} 阐明开源运动的发展历史、核心价值观与推动现代软件基础设施繁荣的战略意义。
    \item \textbf{[解释]} 剖析全球范围内软件盗版与逆向工程的跨国治理困境与产业应对措施。
\end{enumerate}

\subsection{隐私与公民权利}
\textbf{【学时建议】}：2 计算机核心学时，1 领域核心学时。

电子信息的低成本存储、秒级跨网传播与海量聚合，使得在“保护个人隐私权益”与“促进信息自由流通和公共治理利用”之间寻求精巧平衡成为极其紧迫的时代课题。数字化数据采集的极度便捷，使个人隐私权与公民自由面临前所未有的制度性侵蚀，而不同文化传统与主权国家对隐私边界的定义与敏感度存在巨大差异。

\subsubsection{计算机核心主题}
\begin{itemize}
    \item \textbf{泛在数据采集的隐私侵蚀}：日常交易数据库、超大规模数据仓库、公共监控摄录系统、车联网传感器与云计算对个人全生命周期数字足迹的无死角记录。
    \item \textbf{差分隐私（Differential Privacy）数学原理与应用外溢}：隐私预算参数 $\varepsilon$ 的权衡、拉普拉斯噪声注入及在人口普查与用户画像聚合中的实践效应。
    \item \textbf{隐私保护技术增强方案}：同态加密、零知识证明、安全多方计算（SMPC）、端到端加密通信与联邦学习。
    \item \textbf{公民权利、言论空间与跨文化法理差异}：不同国家与文化背景下对信息知情权、被遗忘权、匿名权与公民言论自由的法理界定与实践冲突。
\end{itemize}

\subsubsection{领域核心主题}
\begin{itemize}
    \item \textbf{个人隐私权的哲学与宪法法理渊源}：隐私作为人格尊严核心要素的哲学论证，隐私权作为消极自由与自主控制权的法律演进。
    \item \textbf{重点法域的隐私保护立法框架}：欧盟通用数据保护条例（GDPR）、加州消费者隐私法案以及中国《个人信息保护法》（PIPL）的核心合规准则。
    \item \textbf{垂直专业领域的合规要求}：医疗健康数据（HIPAA）、金融交易信用数据、未成年人网络信息保护法等专项合规。
    \item \textbf{表达自由与有害言论规制的张力}：言论自由保护与遏制虚假信息、仇恨言论、网络诽谤之间的法理平衡界限。
\end{itemize}

\subsubsection{示范性学习成果}
\begin{enumerate}
    \item \textbf{[评估]} 针对高并发交易数据库与海量数据仓库系统中的数据泄露威胁与越权访问漏洞，提出完备的技术防御与审计方案。
    \item \textbf{[解释]} 系统阐明现代泛在监视系统（如 RFID 追踪、智慧城市人脸识别天网、ETC 轨迹追踪、智能手机地理围栏）背后的数据采集机制与社会控制隐患。
    \item \textbf{[解释]} 严密阐述差分隐私（Differential Privacy）的核心数学原理、隐私损失边界及在保护统计隐私中的实际成效与局限。
    \item \textbf{[评估]} 深入调研并评估新兴密码学方案与隐私增强技术（如联邦学习、同态加密）在工业级数据中台中的落地代价与防御成效。
    \item \textbf{[解释]} 探讨在不同司法管辖区及跨文化语境中，确立个人隐私法律保护的哲学法理渊源。
    \item \textbf{[评估]} 批判性审视各类现行隐私立法（如 GDPR、PIPL）在立法初衷、合规成本、遏制技术垄断与阻碍初创企业创新之间的多维博弈。
    \item \textbf{[应用]} 提出在设计社交内容与公共交流平台时，既保障用户正当言论自由又有效治理恶意违规内容的系统性工程策略。
\end{enumerate}

\subsection{专业沟通与团队协作}
\textbf{【学时建议】}：2 计算机核心学时，1 领域核心学时。

专业沟通的本质是向具有不同背景、专业知识储备和诉求目标的受众精准、清晰地传递复杂的技术与管理信息。高水平的专业技术沟通能力绝非与生俱来的天赋，而必须在整个大学本科课程体系中依托真实工程情境反复操练与持续雕琢。

\subsubsection{计算机核心主题}
\begin{itemize}
    \item \textbf{硬核技术资料的深度解读、提炼与综合}：深入研读复杂工业源码、架构规范文档、RFC 协议标准并精准撰写执行摘要与技术白皮书。
    \item \textbf{高质量工程技术文档编制}：编写规范严密的技术教程、开发者参考指南及生产级 OpenAPI/SDK 接口文档。
    \item \textbf{团队多模态高效沟通准则}：在口头汇报、书面备忘录、即时通讯工具与异步协作流中践行“清晰、礼貌、精炼”沟通规范。
    \item \textbf{受众感知与针对性沟通策略}：敏锐识别受众画像（商业客户、非技术管理层、资深架构师、法务合规人员），动态调整技术表达颗粒度与侧重点。
    \item \textbf{现代化协同开发工具链的规范运用}：Git/GitHub 协作流、Issue 跟踪系统、敏捷看板、知识库 Wiki 与远程协同评审工具的深度使用。
    \item \textbf{技术争议中的修辞谬误防范}：在代码评审与技术路线争辩中彻底摈弃人身攻击与情感绑架，依靠实测数据与逻辑推导化解分歧。
    \item \textbf{面向专业受众的无障碍与包容性表达}：图表可读性设计、色盲友好配色、屏幕阅读器可识别文本及无歧视中立技术词汇。
\end{itemize}

\subsubsection{领域核心主题}
\begin{itemize}
    \item \textbf{书面与口头交流中的跨文化能力}：在跨国、多语言、多元文化分布式团队中展现文化敏锐度，规避文化禁忌与语言歧义。
    \item \textbf{多维竞争价值观的权衡与精炼表达}：向高管与客户精准汇报技术指标、研发进度、系统质量、人力成本、市场窗口期与财务预算之间的妥协与权衡。
    \item \textbf{以专业写作推动职场问题解决与决策}：撰写正式建议书（RFC）、向管理层警示技术债务与伦理风险、提交针对系统易用性与无障碍改进的技术论证报告。
\end{itemize}

\subsubsection{示范性学习成果}
\begin{enumerate}
    \item \textbf{[解释]} 深刻理解遵循严格排版格式、准确绘制架构图表与规范引用权威文献对撰写专业技术文档的至高重要性。
    \item \textbf{[评估]} 从技术严谨性、表达精炼度、术语无歧义性及受众契合度等多重维度，对他人撰写的技术文档开展专业评审与同行反馈。
    \item \textbf{[开发]} 面向多元化专业受众，设计并完成一场结构严整、图文并茂、照顾无障碍需求且论点鲜明的正式技术汇报演说。
    \item \textbf{[应用]} 规划并主持团队协作活动（线上协作、线下碰头、跨时区异步文档共创），主动营造鼓励全员包容性参与、倾听并珍视不同创见的良性协作氛围。
    \item \textbf{[解释]} 归纳并践行高效团队沟通的黄金法则（明确行动项、降低沟通熵值、文档留痕与闭环反馈）。
        \item \textbf{[应用]} 探讨在多元化与跨文化工程团队中，促进团队高效协作并提升产出成果的有效方法。
    \item \textbf{[评估]} 客观评估自身在依托多元背景协同工作的远程分布式团队中的个人沟通优势与短板，提出针对性提升计划。
\end{enumerate}

\subsection{可持续性与绿色计算}
\textbf{【学时建议】}：1 计算机核心学时，1 领域核心学时。

联合国将“可持续发展”定义为：“既满足当代人的需求，又不对后代人满足其自身需求的能力构成危害的发展模式”；亦可概括为“生态环境、社会公平与经济效益之间的动态平衡”。随着计算技术全方位渗透人类社会，据权威测算，全球电力消耗中已有约 10\% 用于支撑各类计算与网络系统，且该比例仍在加速攀升。此外，消费级电子硬件的极速淘汰引发了对稀土元素与贵金属开采的畸形需求，并催生出触目惊心的全球电子废弃物生态危机。未来的计算机从业者必须牢固树立生态意识，在工程设计与架构决策中深刻预见计算基础设施给全球环境与地方社区带来的长期资源与碳足迹负担。

\subsubsection{计算机核心主题}
\begin{itemize}
    \item \textbf{践行可持续发展理念的工程实践}：在算法与系统实现中全方位权衡环境开销、社会影响、经济可行性与计算能耗。
    \item \textbf{计算设备全生命周期的生态足迹}：计算芯片制造过程中的高能耗与水资源消耗、稀土采掘污染、设备淘汰报废及跨国电子垃圾倾销对发展中国家脆弱生态的重创。
    \item \textbf{高能耗共识算法与能耗权衡}：以区块链比特币 PoW（工作量证明）算法与以太坊 PoS（权益证明）算法为典型范例，深入量化评估其计算能耗、算力军备竞赛与碳排放代价。
\end{itemize}

\subsubsection{领域核心主题}
\begin{itemize}
    \item \textbf{绿色 IT 与低碳软件设计标准规范}：绿色软件工程协会能效标准、能效评级指标（PUE）与软件碳强度（SCI）规范。
    \item \textbf{技术媒介复杂现象的系统性生态效应}：大规模居家远程办公、跨境算力外包、电商平台物流算法对整体社会能源消费结构的系统性级联影响。
    \item \textbf{普适计算与智慧环保赋能}：利用物联网智能电网、智慧交通路线规划、环境自动监测无线传感器网络、公民科学与环境激进主义系统推动社会实现低碳转型。
    \item \textbf{计算技术在解决前沿环境问题中的应用研究}：利用高性能计算与 AI 大模型攻关新能源开发、大气污染扩散模拟、碳循环捕集、循环经济资源调度与精准农业。
    \item \textbf{软件系统可持续性与社会制度的深度共生}：系统生命力与用户认知能力、企业组织流程及政府强制碳税政策之间的互赖机制。
\end{itemize}

\subsubsection{示范性学习成果}
\begin{enumerate}
    \item \textbf{[应用]} 掌握在日常软件工程与系统架构实践中践行“绿色与可持续工程师”的有效途径与衡量准则。
    \item \textbf{[评估]} 针对任意给定的计算工程项目（如大型数据中心上云、开发训练超大参数量模型、投产边缘物联网终端），系统列举并量化其全生命周期的环境与资源开销。
    \item \textbf{[解释]} 绘制并阐明全球电子废弃物跨国转移的生态危害链条，反思科技产品“计划性报废”商业模式的道德伦理失范。
    \item \textbf{[评估]} 客观列举并定量测算远程办公、在线电商或加密货币挖矿等现代数字化活动的综合能源增减效应。
    \item \textbf{[评估]} 深入分析算法时间复杂度、操作系统任务调度、分布式缓存、数据库查询优化等计算机内核设计决策对物理硬件功耗的直接影响。
    \item \textbf{[评估]} 调研全新大型软件系统架构投产对所在城市微电网、散热水资源等周边物理环境的长期负外部性。
    \item \textbf{[开发]} 依据国际绿色计算标准，编制一套适用于本企业软件研发与机房部署的绿色低碳开发操作指南。
    \item \textbf{[评估]} 调研智慧能源调度、智能物流供应链管理及环境遥感监测等计算系统在化解现实生态环境危机中的落地价值。
\end{enumerate}

\subsection{计算机发展简史}
\textbf{【学时建议】}：1 领域核心学时。

开设计算机科学史并非为了单纯记诵历史事实，而是为了让学生深刻领悟技术的高速演进如何在长周期内深刻塑造人类文明。科学史教学通常与“系统基础 (SF)”和“软件开发基础 (SDF)”紧密结合。研习历史使我们能够洞悉当今计算机系统架构为何如此设计、早期方案是在何种具体的社会需求与硬件技术瓶颈下诞生、以及历史的技术路径依赖如何持续在当今学术与产业界激荡回响。

\subsubsection{领域核心主题}
\begin{itemize}
    \item \textbf{第一纪元：史前文明与机械计算（1946 年 ENIAC 诞生前）}：古代模拟计算工具（巨石阵天文历法、安提基特拉机械仪、中世纪机械大钟）；欧几里得几何算法、阿达·洛夫莱斯与查尔斯·巴贝奇分析机、哥德尔不完备定理、阿隆佐·邱奇Lambda 演算、艾伦·图灵图灵机；机电式与继电器计算硬件。
    \item \textbf{第二纪元：现代早期数字电子计算}：ENIAC、UNIVAC、布莱切利庄园破译密码的密码机、晶体管替代电子管革命（TRADIC）、大型主机与最早的通用操作系统（IBM 操作系统/360）。
    \item \textbf{第三纪元：现代微型计算与个人电脑革命}：微处理器诞生、微型计算机（PC）普及、摩尔定律（Moore's Law）的半个世纪奇迹、施乐帕克研究中心对图形用户界面（GUI）与以太网的奠基、第一代 Apple 计算机与个人计算普及。
    \item \textbf{第四纪元：互联网时代与分布式互联}：ARPANET、TCP/IP 协议栈统一世界、万维网（WWW）诞生与浏览器大战、互联网技术标准委员会（IETF/W3C）、互联网商业泡沫与科技巨头垄断格局形成。
    \item \textbf{第五纪元：智能手机、移动互联与云计算时代}：智能手机生态普及（iOS、Android 移动革命）、集中式多租户云计算基础设施（AWS、Azure）、软件即服务、全网移动社交网络繁荣与隐私危机的全面交织。
    \item \textbf{第六纪元：AI 辅助与前沿智能时代}：深度学习复兴、基于算法的自动化决策支持系统、大规模推荐引擎、大语言模型（LLMs）与生成式人工智能对内容生产、工程研发与人机交互形态的全面重塑。
\end{itemize}

\subsubsection{示范性学习成果}
\begin{enumerate}
    \item \textbf{[解释]} 清晰辨析计算机学科发展史上几个关键历史时期的技术代际特征与核心演进驱动力。
    \item \textbf{[解释]} 准确阐述多位计算科学先驱（如图灵、冯·诺依曼、洛夫莱斯、香农、戴克斯特拉）对现代计算理论与工程架构的不可磨灭贡献。
    \item \textbf{[评估]} 深入探讨计算机发展史上重大里程碑事件的历史语境与深远意义（例如：电子管向晶体管与集成电路的过渡、操作系统/360 的软件工程危机、施乐 PARC 实验室的人机交互革命、第一款计算机病毒发端、WWW 协议开源、千年虫危机 Y2K、移动智能机浪潮）。
    \item \textbf{[评估]} 深刻对比个人电脑普及与互联网大爆炸前后，人类社会日常生活、工作协作与全球经济版图的天翻地覆变革。
\end{enumerate}

\subsection{计算经济学与产业生态}
\textbf{【学时建议】}：1 领域核心学时。

计算经济学不仅关乎软硬件服务提供商的生存壮大，也直接牵涉全社会各行业的经济运行，同时与数以亿计的专业及非专业终端用户的切实福祉紧密相连。计算机从业人员必须理解计算产业运行的经济学规律，方能在复杂的商业浪潮中做出清醒的技术选型与战略研判。

\subsubsection{领域核心主题}
\begin{itemize}
    \item \textbf{数字服务提供商的市场生态}：强监管与弱监管市场环境、自然垄断与开放竞争市场；大型科技跨国平台构建的“围墙花园”专有生态及其对跨平台创新的抑制。
    \item \textbf{知识经济与注意力经济}：数据要素资产化、算法诱导机制以争夺人类有限注意力作为核心变现模式的经济本质。
    \item \textbf{专业高技术人才供求平衡与技术产品质量的关系}：全球算力与软件人才缺口对系统架构质量、测试充分性与安全隐患的深层连锁反应。
    \item \textbf{计算领域的商业定价与商业模式演化}：买断制、软件订阅制、计划性报废（人为制造硬件换代需求）、软件使用许可授权、自由及开源软件商业化路径。
    \item \textbf{软件工程外包与离岸研发}：全球研发分工、离岸交付中心对发达国家本土就业与发展中国家技术生态的双向冲击。
    \item \textbf{全球化分工与科技自主的博弈}：跨国技术产业链分工与地缘政治科技壁垒对开发者生态和最终用户的交织影响。
    \item \textbf{计算资源获取鸿沟造成的经济不平等}：高端算力垄断与中小企业技术研发创新的两极分化。
    \item \textbf{前沿智能化对劳动力市场的颠覆性重构}：机器人自动化与生成式 AI 对白领、程序员及基础劳动岗位的结构性冲击与就业转型阵痛。
    \item \textbf{规模经济、科技创业孵化与科技慈善}：零边际复制成本、高固定研发投入特征下的风险投资（VC）模式与开源基金会慈善运作。
    \item \textbf{计算技术颠覆现代个人与社会金融体系}：分布式区块链与加密货币、移动支付与数字钱包、发展中国家基于短信（SMS）的普惠金融基础设施变革。
\end{itemize}

\subsubsection{示范性学习成果}
\begin{enumerate}
    \item \textbf{[解释]} 归纳并剖析科技巨头垄断、平台“围墙花园”封闭生态对数字自由竞争与开放 Web 生态的扼杀隐患。
    \item \textbf{[解释]} 探讨高水平专业技术人才严重短缺对企业级核心系统稳定性、安全性保障造成的直接拖累。
    \item \textbf{[解释]} 梳理计算软件与数字服务商业模式自早期纸带/软盘买断，走向订阅制、按用量计费与免费增值的历史演进脉络。
    \item \textbf{[解释]} 深入阐释“知识经济”与“注意力经济”的深层社会代价（如短视频沉迷、心理焦虑、深度虚假信息传播）。
    \item \textbf{[评估]} 总结科技全球化分工与近期部分国家盛行的逆全球化科技壁垒对全球开源协作与技术标准统一造成的双重影响。
    \item \textbf{[评估]} 深刻分析人工智能全自动化浪潮对人类劳动力市场结构的冲击，前瞻性预判传统技术开发岗位的演化方向。
    \item \textbf{[解释]} 详细阐明计算技术基础设施的升级如何彻底重塑现代企业的组织架构与管理范式。
    \item \textbf{[评估]} 系统梳理移动支付、区块链金融与数字法币对个人理财及国家金融安全带来的积极赋能与潜在洗钱、投机金融风险。
\end{enumerate}

\subsection{安全策略、法律与网络犯罪}
\textbf{【学时建议】}：2 计算机核心学时，1 领域核心学时。

虽然安全策略、法律规范与网络犯罪惩治是现代计算体系的核心支柱，但这些议题绝非单纯的技术对抗，必须植根于“知识产权”、“隐私与公民权利”、“社会语境”及“职业伦理”等先修知识的宏观视野中审视。计算机与互联网一方面为人类文明带来空前繁荣，另一方面也孕育出前所未有的隐私威胁、新型跨国有组织计算机犯罪、反社会网络暴力、企事业单位运营瘫痪以及关键基础设施的系统性高危集中风险。

\subsubsection{计算机核心主题}
\begin{itemize}
    \item \textbf{经典计算机犯罪典型案例与法律追责机制}：重大网络勒索攻击、分布式拒绝服务（DDoS）敲诈、特大跨国数据泄露及执法机构的跨国司法管辖追诉。
    \item \textbf{社会工程学、高技术欺诈与灾后响应恢复}：网络钓鱼、深度伪造欺诈、身份冒用以及被入侵后的应急灾备恢复机制。
    \item \textbf{计算机犯罪的法律界定范畴}：非法侵入计算机信息系统罪、越权获取数据、破坏计算机信息系统功能、非法利用信息网络协助犯罪。
    \item \textbf{网络恐怖主义与黑客地下黑产攻击动机及破坏后果}：国家级高级持续性威胁（APT）、关键市政能源基础设施破坏、地下黑客暗网数据买卖黑产链。
    \item \textbf{恶意软件谱系及其系统破坏机理}：计算机病毒、自主传播蠕虫、特洛伊木马、间谍软件、勒索加密软件与 Rootkit 的行为特征。
\end{itemize}

\subsubsection{领域核心主题}
\begin{itemize}
    \item \textbf{企事业单位安全防护与犯罪预防战略}：纵深防御架构、网络攻防演练（红蓝对抗）、零信任安全接入规范与安全意识全员培训。
    \item \textbf{企业级信息安全策略制定与合规审计}：最小权限准则、强密码生命周期管理、多因素认证策略、雇员合规网络行为审计准则。
\end{itemize}

\subsubsection{示范性学习成果}
\begin{enumerate}
    \item \textbf{[解释]} 梳理引发广泛社会震荡的经典计算机犯罪事件与社会工程学攻击案，总结其背后的技术与管理脆弱点。
    \item \textbf{[解释]} 准确列举并理解适用于惩治网络黑客入侵、数据窃取与网络诈骗的相关刑事法律与网络安全法条。
    \item \textbf{[解释]} 深入剖析网络恐怖主义与恶意破解者的心理动机、行为特征以及对国家安全与社会秩序的破坏性后果。
    \item \textbf{[评估]} 深刻审查因权限滥用、内部特权账号失控所导致的严重数据泄露事件中的法律违规与职业道德背叛。
    \item \textbf{[评估]} 深入探讨计算机专业人员在构筑信息系统安全防线中的关键角色，以及在“绝对安全”、“系统可用性”与“用户隐私”之间寻求权衡的职业取舍。
    \item \textbf{[应用]} 针对企事业单位及个人，系统设计一套抵御计算机勒索犯罪、防范身份冒用及防范跨国欺诈的综合防治方案。
    \item \textbf{[开发]} 为一家虚拟中型科技企业独立编制一份完备的企业级安全策略文档，涵盖口令密码策略、特权访问分级、员工行为合规监控及数据出境保护规程。
\end{enumerate}

\subsection{多元、公平与包容}
\textbf{【学时建议】}：2 计算机核心学时，2 领域核心学时。

自计算机学科建立伊始，多元性不足就是困扰该领域的长期顽疾。尽管计算机专业充满创造力、薪资待遇丰厚且就业空间广阔，但在种族、性别、身体机能及社会背景方面的代表性严重失衡现象却普遍存在。对于极其广泛的弱势群体而言，他们的第一门计算机先修课往往不幸成为他们的最后一门课。这种失衡背后交织着系统性的排斥壁垒、刻板印象、健全主义偏见、性别偏见与阶层固化等多重历史与制度成因，绝非依靠某种单一速效方案所能根治。

CS2023 的发起主办学会——ACM、IEEE-CS 与 AAAI，均将“多元、公平与包容（DEI）”确立为协会的核心准则。\footnote{参见 ACM DEI 准则：\url{https://www.acm.org/diversity-inclusion/about}；IEEE 多元化声明：\url{https://www.ieee.org/about/diversity-index.html}；AAAI 多元化政策：\url{https://aaai.org/Organization/diversity-statement.php}}计算机高等教育课堂不仅应当向学生传授这些原则，更应成为践行与倡导包容性文化的模范阵地：
\begin{itemize}
    \item \textbf{多元}：真诚拥抱并主动吸纳人类全方位的个体差异（包括性别认同、种族族群、身体功能状况、文化信仰、国籍背景等）；
    \item \textbf{包容}：致力打破物理与数字藩篱，消除无障碍壁垒，营造让每一位成员感受到被尊重、被倾听与拥有归属感的友好学术社群；
    \item \textbf{公平}：直面历史不公，精准倾斜教育资源，审查与校正无意识偏见，确保不同背景的人才均能获得平等的成就空间。
\end{itemize}
在计算机课程体系中全面贯穿 DEI 理念不仅具有坚实的道德与正义正当性，更是决定科技产业健康发展的关键能力。权威实证数据表明，背景多元化的研发团队在创造力与抗风险能力上显著超越同质化团队；许多震惊世界的重大技术溃败事故（如人脸识别在深肤色群体上的极高误判率、早期汽车安全气囊未适配女性与儿童体征致死悲剧），原本完全可以通过构建多元视角的研发与评测团队予以彻底规避。

\subsubsection{计算机核心主题}
\begin{itemize}
    \item \textbf{个体身份与计算技术环境的双向互塑}：个人身份特质如何深刻影响在学术与产业计算环境中的职业体验，以及技术系统反向对不同群体身份认同造成的赋能或标签化禁锢。
    \item \textbf{多元化研发团队的战略收益与同质化团队的重大盲区}：深入剖析团队成员背景多元化带来的算法盲区规避、产品健壮性提升以及同质化带来的技术偏见风险。
    \item \textbf{包容性语言与敏感偏见术语的净化规范}：掌握国际计算界推动的中立包容用语规范，规避带有历史压迫色彩的争议性词汇（如 Master/Slave 架构重构为 Primary/Replica、Blacklist/Whitelist 重构为 Denylist/Allowlist）。
    \item \textbf{包容性行为范式及其职场价值}：在项目评审、日常研发与决策讨论中主动践行倾听弱势声音、提供心理安全感的包容行为。
    \item \textbf{无障碍辅助技术与通用设计}：符合国际 W3C Web 内容无障碍指南（WCAG）规范的交互系统开发、辅助屏幕阅读器适配与无障碍软硬件基准。
    \item \textbf{计算机从业人员对社会正义的正负向杠杆作用}：从业者通过其编写的代码和算法，既可能加剧社会歧视与阶层鸿沟，也拥有强大的正向杠杆来促进社会公平与正义。
\end{itemize}

\subsubsection{领域核心主题}
\begin{itemize}
    \item \textbf{树立契合现实世界的多元化榜样}：在课程与教材中充分展示具有不同性别、种族和身体状况的杰出计算先驱、行业专家与高年级优秀学长。
    \item \textbf{多元化的系统性实证收益与同质化带来的高昂代价}：深入研讨多元化提升软件工程质量与规避重大事故的经典工业界实证调研与案例剖析。
    \item \textbf{技术霸权造成的历史边缘化与全球数字基建公平挑战}：西方技术霸权语境下边缘群体的技术失语症、全球算力基建分布不均对发展中国家公平获取技术红利的阻碍。
\end{itemize}

\subsubsection{示范性学习成果}
\begin{enumerate}
    \item \textbf{[解释]} 准确界定并辨析“公平”、“平等”、“多元”与“包容”在计算机工程语境中的深刻区别。
    \item \textbf{[解释]} 深刻阐明技术权力与特权结构在计算机行业文化、软硬件产品定义及社会阶层流动中扮演的角色。
    \item \textbf{[解释]} 总结能够营造包容性职场与团队氛围的日常表达用语、工程习惯与协作行为；主动践行 ACM “Words Matter”包容性用语指南。\footnote{参见 ACM Words Matter 倡议：\url{https://www.acm.org/diversity-inclusion/words-matter}}
    \item \textbf{[评估]} 依据 W3C WAI/WCAG 国际标准，独立开展本专业实验室、教学机房及学院官方主页的信息无障碍评测，精准定位可访问性障碍并提出修复方案。
    \item \textbf{[应用]} 具备与不同身份背景、文化信仰的团队伙伴精诚合作、相互尊重的团队协同能力。专业教师绝不能仅将学生随机分组了事，而应全程指导、监控并培养学生团队的包容性协作能力。
    \item \textbf{[评估]} 采集并比对所在高校计算机及 STEM 专业的人口统计学画像与全校整体学生画像的差异；深入剖析导致特定群体在计算机专业招生、留存与毕业深造中处于劣势的制度与心理成因。
    \item \textbf{[评估]} 比对所在高校或企业的人口结构与周边社会社区的人口多样性差异，评估科技高地是否存在与本地社区脱节的“精英技术孤岛”效应。
    \item \textbf{[解释]} 敏锐识别软硬件系统设计中暗含的开发者文化预设与潜意识偏见，评估其对不同文化背景和身体障碍用户的可用性杀伤力。
    \item \textbf{[解释]} 在课程技术案例中积极引用并致敬来自不同历史时期、跨文化背景的杰出计算科学家，打破行业单一刻板印象。
    \item \textbf{[解释]} 列举生动案例，证明研发团队多元化如何帮助企业成功开拓全球新兴市场，以及因缺乏多元视角而导致重大产品翻车的教训。
    \item \textbf{[开发]} 针对熟悉的日常场景（如《程序设计基础》初级核心课），提出一套切实可行的系统性改进方案，从教学案例选择、答疑辅导机制到助教培训全面促进多样、公平与包容。
\end{enumerate}

\section{专业素养与思维倾向}

在社会、伦理与专业素养领域，单纯掌握知识点不足以确保毕业生在职业生涯中做出合乎道德的技术行动。CS2023 强调必须内化以下六大核心素养特质：
\begin{itemize}
    \item \textbf{专业职业精神}：现代计算实践极少是孤立封闭的技术操练，高度依赖团队配合、客户协同、组织领导及同行交流。无论面对何种身份特质与文化背景的协作伙伴，均能展现出高度的同理心、职业严谨性与平等待人风范。
    \item \textbf{责任担当意识}：深怀对代码质量、系统可靠性与社会后果的敬畏之心，对微小架构漏洞与潜在恶性连锁反应保持极高的责任审慎感。
    \item \textbf{批判性自我反思}：具备深度审视自身行为、潜意识偏见、所享特权及深层动机的思辨反思能力。主动学习识别隐性偏见，并在算法建模和团队决策中持续纠偏。
    \item \textbf{环境敏捷响应性}：能够迅速、审慎地适应外部技术与社会环境的突变（如面对突发公共事件时从线下集中办公平稳切换至分布式远程研发模式），并主动重构对“专业性”内涵的认知。
    \item \textbf{前瞻主动性}：在职场中不仅被动执行规则，更能主动洞察新兴前沿趋势（如国际无障碍标准升级、绿色算法优化），并前瞻性地将最佳实践推行至团队现有研发流程中。
    \item \textbf{文化能力}：高度重视跨文化协同能力，在跨国分布式开发中规范使用包容性语言，警惕文化偏见，主动促进真诚、透明与平等的跨团队思想交锋。
\end{itemize}

\section{数学基础要求}

SEP 领域虽然重点关注社会学、法学与伦理学，但其定量评估部分与现代应用数学紧密相连：
\begin{itemize}
    \item \textbf{离散数学与形式逻辑}：论证有效性验证、逻辑谬误拆解、形式化安全策略验证与知识产权权利范围的严密界定。
    \item \textbf{概率论与数理统计}：算法公平性定量测算（如人口均等 Demographic Parity、均等化赔率 Equalized Odds）、差分隐私拉普拉斯扰动注入与隐私预算 $\varepsilon$ 评估、多元化实证调研抽样推断。
\end{itemize}

\section{交叉主题与跨学科共享概念}

\subsection{与全书其他知识领域的共享概念}
\begin{itemize}
    \item \textbf{社会正义、公平、多元与包容}：贯穿全书所有 17 个 KA，尤其是人工智能（AI）、人机交互（HCI）、软件工程（SE）和计算机系统（系统基础/OS）。
    \item \textbf{超网络化与 AI 高速演进世界的社会影响 (社会伦理/社会情境与计算影响)}：与 人工智能/人工智能基础议题、人机交互/交互基础、安全隐私/Fundamental Concepts 紧密相扣。
    \item \textbf{互联网的演进与控制权治理}：与 网络通信/网络导论与体系结构 互联网体系结构与组织管理共享。
    \item \textbf{情境感知计算}：与 人机交互/Design for non-mouse interfaces 普适与情境感知计算深度交叉。
    \item \textbf{职业认证与伦理行为规范}：与 安全隐私/Fundamental Concepts 伦理与合规要求互为支撑。
    \item \textbf{知识产权与数字专有权}：与 数据管理/Information Storage and Retrieval 知识产权与数字资产保护协同。
    \item \textbf{广泛数据采集带来的隐私侵蚀}：与 数据管理/Database Systems 数据独立性、数据管理/Data Mining 数据清洗与挖掘伦理直接关联。
    \item \textbf{基于技术的隐私保护解决方案}：与 安全隐私/Threats and Attacks 针对隐私与匿名性的对抗攻击互为攻防。
    \item \textbf{隐私权的哲学渊源}：与 人工智能/人工智能基础议题 哲学伦理思辨深度交织。
    \item \textbf{高效口头、书面与团队工程沟通 (社会伦理/专业沟通与表达)}：与 人机交互/Collaboration and Communication 群体沟通、软件工程/项目管理与度量 敏捷团队协作深度重叠。
    \item \textbf{现代计算发展史}：与 体系结构/Digital logic and digital systems 计算机体系结构演进史紧密映射。
    \item \textbf{计算机犯罪、社会工程学与安全策略}：与 安全隐私/Digital Forensics 证据采信规则、人机交互/Human Factors and Security 信任与欺骗建模、安全隐私/Security Policy and Governance 企业安全治理深度协同。
\end{itemize}

\subsection{贯穿全领域的横切主题}
\begin{itemize}
    \item \textbf{正义} \quad $\bullet$ \textbf{公平} \quad $\bullet$ \textbf{多元} \quad $\bullet$ \textbf{包容}
    \item \textbf{社会议题} \quad $\bullet$ \textbf{计算伦理} \quad $\bullet$ \textbf{专业职业素养}
\end{itemize}

\section{课程体系规划与开课建议}

联合工作组建议，各高校计算机系至少应在其他技术核心课程中全面涵盖 计算机核心 的学习目标；理想状态下，领域核心 的核心内容也应全方位渗透：

\subsection{中阶独立课程方案 (面向 17 计算机核心学时)}
在本科二年级或三年级设立一门 17--20 学时的必修独立课程《计算机职业素养与社会伦理导论》，紧凑覆盖所有 计算机核心 单元：
\begin{itemize}
    \item 社会语境：3 学时
    \item 伦理分析方法：2 学时
    \item 职业伦理与专业准则：2 学时
    \item 知识产权：1 学时
    \item 隐私与公民权利：2 学时
    \item 专业沟通：2 学时
    \item 可持续性与绿色计算：1 学时
    \item 安全策略、法律与网络犯罪：2 学时
    \item 正义、公平、多元与包容 (DEI)：2 学时
\end{itemize}

\subsection{中高阶独立完整深化方案 (面向 31 全核心学时)}
在本科高年级设立一门标准的 32 学时（2 学分）完整必修课《计算伦理、法律、社会与可持续性深度研讨》，完整统摄 计算机核心 与 领域核心：
\begin{itemize}
    \item 社会语境：5 学时
    \item 伦理分析方法：3 学时
    \item 职业伦理与专业准则：4 学时
    \item 知识产权与开源合规：2 学时
    \item 隐私与公民权利：3 学时
    \item 专业沟通与技术写作：3 学时
    \item 可持续性与绿色计算：2 学时
    \item 计算机发展简史：1 学时
    \item 计算经济学与产业生态：1 学时
    \item 安全策略、网络犯罪与法律法规：3 学时
    \item 多元、公平与包容 (DEI)：4 学时
\end{itemize}

\section{真实能力任务规范}

\begin{taskbox}[任务 1：评估部署特定 AI 驱动服务/软件/产品的伦理与社会影响]
\textbf{【任务情境】}：某企业拟将一套基于深度生成模型与算法决策的 AI 辅助招聘/信用审批服务投入商业化部署。\\
\textbf{【能力陈述】}：综合识别并评估该 AI 服务上线后将影响哪些社会利益相关群体，并精确判定其如何造成歧视性决策、隐私侵害或权益受损。\\
\textbf{【能力领域】}：理论与宏观分析 \quad \textbf{【能力单元】}：系统部署 / 方案评估 / 用户接纳 / 社会议题调适\\
\textbf{【支撑知识单元】}：人工智能/人工智能基础议题、社会伦理/社会情境与计算影响、社会伦理/伦理分析方法、社会伦理/隐私与公民自由、社会伦理/多样性、平等与包容 (DEI)\\
\textbf{【达标技能等级】}：\textbf{解释与评估} \quad \textbf{【核心要求】}：计算机核心
\end{taskbox}

\begin{taskbox}[任务 2：评估收集与使用客户/用户数据的法律与伦理意涵]
\textbf{【任务情境】}：研发团队拟在移动端 APP 和云端后端建立全埋点用户行为采集系统以支撑精准商业变现。\\
\textbf{【能力陈述】}：严密分析拟采集数据的敏感性属性，精准判定在该数据全生命周期采集、存储、聚合挖掘与流转过程中可能引发的所有法定合规风险与伦理侵权隐患。\\
\textbf{【能力领域】}：理论分析 \quad \textbf{【能力单元】}：方案评估 / 用户接纳 / 社会议题调适\\
\textbf{【支撑知识单元】}：数据管理/数据安全与隐私保护、社会伦理/伦理分析方法、社会伦理/隐私与公民自由、社会伦理/多样性、平等与包容 (DEI)\\
\textbf{【达标技能等级】}：\textbf{解释与评估} \quad \textbf{【核心要求】}：计算机核心
\end{taskbox}

\begin{taskbox}[任务 3：评估文化因素在技术产品采纳与社会公平中的决定性角色]
\textbf{【任务情境】}：某全球化跨国软件服务进军新兴发展中国家市场，遭遇用户使用习惯排斥与当地特定群体利益受损争议。\\
\textbf{【能力陈述】}：熟练运用严谨的伦理分析框架与立场理论，准确界定适用的社会文化语境，推演该技术引入对不同社会阶层、文化群体的深层伦理冲击与公平性影响。\\
\textbf{【能力领域】}：理论分析 \quad \textbf{【能力单元】}：系统部署 / 方案评估 / 社会议题调适\\
\textbf{【支撑知识单元】}：社会伦理/社会情境与计算影响、社会伦理/伦理分析方法、社会伦理/隐私与公民自由、社会伦理/多样性、平等与包容 (DEI)\\
\textbf{【达标技能等级】}：\textbf{解释与评估} \quad \textbf{【核心要求】}：计算机核心
\end{taskbox}

\begin{taskbox}[任务 4：编制企业发布技术服务所承担的问责、责任与法律责任规范文档]
\textbf{【任务情境】}：企业即将向公众正式发布一款高自主度自动驾驶系统或关键医疗辅助分析软件平台。\\
\textbf{【能力陈述】}：系统收集法务、工程质量、网络安全与行业规范信息，撰写一份条理严密、可实操的企业级责任声明与风险问责白皮书。\\
\textbf{【能力领域】}：应用与理论实践 \quad \textbf{【能力单元】}：需求工程 / 系统部署 / 文档编制\\
\textbf{【支撑知识单元】}：社会伦理/职业道德与规范、社会伦理/知识产权与开源许可、社会伦理/隐私与公民自由、社会伦理/专业沟通与表达、社会伦理/多样性、平等与包容 (DEI)\\
\textbf{【达标技能等级】}：\textbf{解释} \quad \textbf{【核心要求】}：计算机核心
\end{taskbox}

\begin{taskbox}[任务 5：制定团队跟进与企业战略相匹配的伦理、法律及职业议题策略]
\textbf{【任务情境】}：研发中心规模快速扩张，亟需建立规范机制确保工程团队在新兴技术研发中与企业伦理战略保持步调一致。\\
\textbf{【能力陈述】}：综合评估企业中长期战略目标，制定一套通过专业技术文档和日常沟通渠道向全体工程师持续宣贯、培训与追踪伦理、法律动态的行动方案。\\
\textbf{【能力领域】}：应用与工程管理 \quad \textbf{【能力单元】}：全周期管理 / 需求分析 / 方案评估 / 持续改进\\
\textbf{【支撑知识单元】}：社会伦理/专业沟通与表达、社会伦理/职业道德与规范\\
\textbf{【达标技能等级】}：\textbf{解释、评估与开发} \quad \textbf{【核心要求】}：计算机核心
\end{taskbox}

\begin{taskbox}[任务 6：将合法合规的隐私保护要求深度融入系统研发生命周期]
\textbf{【任务情境】}：架构团队正在设计全新企业级微服务分布式中台，需确保系统从代码编写、API 暴露到存储均满足 GDPR 与国家数据安全法。\\
\textbf{【能力陈述】}：识别并分解法定隐私与伦理合规指标，制定可执行的技术实施方案与团队沟通计划，确保隐私保护设计在系统开发全流程中落地。\\
\textbf{【能力领域】}：软件工程与理论实践 \quad \textbf{【能力单元】}：需求分析 / 架构设计 / 系统开发 / 流程管控 / 质量改进\\
\textbf{【支撑知识单元】}：社会伦理/职业道德与规范、社会伦理/隐私与公民自由、社会伦理/专业沟通与表达\\
\textbf{【达标技能等级】}：\textbf{解释与开发} \quad \textbf{【核心要求】}：计算机核心
\end{taskbox}

\begin{taskbox}[任务 7：在技术写作与口头汇报中展现敏锐的文化意识与文化能力]
\textbf{【任务情境】}：作为跨国开源基金会或大型跨国工程团队的负责人，需向来自不同大洲、多元文化背景的核心利益相关者发表技术演讲与撰写正式提案。\\
\textbf{【能力陈述】}：在书面技术文档与口头演说中高度展现文化敏锐度，自觉规避排他性用语与文化偏见，实现富有同理心且极具说服力的专业技术沟通。\\
\textbf{【能力领域】}：理论与人文素养 \quad \textbf{【能力单元】}：技术写作 / 社会议题调适\\
\textbf{【支撑知识单元】}：社会伦理/社会情境与计算影响、社会伦理/专业沟通与表达、社会伦理/多样性、平等与包容 (DEI)\\
\textbf{【达标技能等级】}：\textbf{解释} \quad \textbf{【核心要求】}：计算机核心
\end{taskbox}

\begin{taskbox}[任务 8：参与制定企业或研发团队的可持续发展与绿色低碳政策]
\textbf{【任务情境】}：科技企业响应国家双碳战略，计划在其全球数据中心运营与日常软件工程研发流程中推行全面碳减排。\\
\textbf{【能力陈述】}：有效将可持续计算与绿色计算知识与企业商业战略及组织使命深度融合，撰写具备行业领先水平的软件能耗优化与设备绿色循环政策。\\
\textbf{【能力领域】}：理论与工程战略 \quad \textbf{【能力单元】}：方案设计 / 系统构建 / 策略融合 / 文档编制 / 持续优化\\
\textbf{【支撑知识单元】}：社会伦理/专业沟通与表达、社会伦理/绿色计算与可持续发展\\
\textbf{【达标技能等级】}：\textbf{解释、评估与开发} \quad \textbf{【核心要求】}：计算机核心
\end{taskbox}

\begin{taskbox}[任务 9：阐释研发团队与用户多元化对企业创新与系统质量的深远收益及同质化危害]
\textbf{【任务情境】}：为企业高管层撰写决策咨询报告，论证在产研部门大力推进员工背景多元化、倾听弱势用户群体诉求对规避重大商业与公关危机的重要性。\\
\textbf{【能力陈述】}：运用严谨实证数据与典型工业惨痛事故教训，向决策层专业且令人信服地传递多元化对激发团队创新活力、消除技术偏见的决定性作用。\\
\textbf{【能力领域】}：理论与组织沟通 \quad \textbf{【能力单元】}：文档编制 / 团队管理 / 社会议题调适 / 组织改进\\
\textbf{【支撑知识单元】}：社会伦理/专业沟通与表达、社会伦理/多样性、平等与包容 (DEI)\\
\textbf{【达标技能等级】}：\textbf{解释} \quad \textbf{【核心要求】}：计算机核心
\end{taskbox}

\begin{taskbox}[任务 10：协助设计并推行旨在切实促进员工与用户群体 DEI 的系统性组织变革]
\textbf{【任务情境】}：高校计算机系或大型科技部门立项开展“多元、平等与包容（DEI）”系统性革新工程。\\
\textbf{【能力陈述】}：在具体的教学或研发组织情境中创造性应用 DEI 原理，针对人才选拔、代码评审、无障碍系统优化及职场关怀，设计并推行一套长效变革机制。\\
\textbf{【能力领域】}：应用与理论实践 \quad \textbf{【能力单元】}：需求分析 / 制度设计 / 用户接纳 / 社会议题调适 / 系统性改进\\
\textbf{【支撑知识单元】}：社会伦理/专业沟通与表达、社会伦理/多样性、平等与包容 (DEI)\\
\textbf{【达标技能等级】}：\textbf{解释、应用、评估与开发} \quad \textbf{【核心要求】}：计算机核心
\end{taskbox}

\section{工作组委员会与致谢名单}

\textbf{委员会主席：}\\
Brett A. Becker (爱尔兰都柏林大学 / University College Dublin, Ireland)

\textbf{指导委员会成员：}
\begin{itemize}
    \item Richard L. Blumenthal (美国瑞吉斯大学 / Regis University, USA)
    \item Michael Goldweber (美国泽维尔大学 / Xavier University, USA)
    \item James Prather (美国阿比林基督大学 / Abilene Christian University, USA)
    \item Susan Reiser (美国北卡罗来纳大学阿什维尔分校 / University of North Carolina Asheville, USA)
    \item Michelle Trim (美国马萨诸塞大学阿默斯特分校 / University of Massachusetts Amherst, USA)
    \item Titus Winters (美国纽约 Google 公司 / Google, Inc, USA)
\end{itemize}

\textbf{核心贡献者：}
\begin{itemize}
    \item Johanna Blumenthal (美国瑞吉斯大学 / Regis University, USA)
    \item MaryAnne Egan (美国锡耶纳学院 / Siena College, USA)
    \item Keith Quille (爱尔兰都柏林理工大学 / Technological University of Dublin, Ireland)
    \item Mehran Sahami (美国斯坦福大学 / Stanford University, USA)
    \item Mark Scanlon (爱尔兰都柏林大学 / University College Dublin, Ireland)
    \item Karren Shorofsky (美国旧金山大学法学院 / University of San Francisco School of Law, USA)
    \item Ellen Walker (美国海勒姆学院 / Hiram College, USA)
\end{itemize}

\section{附录：核心主题与技能要求等级对照表}

\begin{longtable}{p{0.20\textwidth} p{0.48\textwidth} c c c}
\toprule
\textbf{知识单元} & \textbf{核心知识点与具体主题} & \textbf{认知等级} & \textbf{分类} & \textbf{学时} \\
\midrule
\endfirsthead
\toprule
\textbf{知识单元} & \textbf{核心知识点与具体主题} & \textbf{认知等级} & \textbf{分类} & \textbf{学时} \\
\midrule
\endhead
\bottomrule
\endfoot
\multirow{2}{*}{\makecell[l]{社会伦理/社会语境}} 
& AI 演进世界的社会影响、社交媒体与算法决策融入公民生活 & 评估 & 计算机核心 & 3 \\
& 互联网演进与控制权、数字鸿沟、信息无障碍规范、情境感知计算 & 解释 & 领域核心 & 2 \\
\midrule
\multirow{2}{*}{\makecell[l]{社会伦理/伦理分析方法}} 
& 论证逻辑谬误规避、伦理学理论与技术决策框架、文化塑造作用 & 应用 & 计算机核心 & 2 \\
& 职业伦理自查清单、方案评估量规、利益相关者分析、立场理论 & 开发 & 领域核心 & 1 \\
\midrule
\multirow{2}{*}{\makecell[l]{社会伦理/职业伦理与\\专业准则}} 
& 公共价值与法律、专业素养与信托责任、学会守则、软件安全法律责任 & 评估 & 计算机核心 & 2 \\
& 公共政策参与、不良后果预警、“吹哨人”保护、跨文化道德困境、AUP 规范 & 解释 & 领域核心 & 2 \\
\midrule
\multirow{2}{*}{\makecell[l]{社会伦理/知识产权}} 
& 知识产权哲学法理基石、专有权利体系、无形数字知识产权 & 解释 & 计算机核心 & 1 \\
& DRM 技术法理、版权/专利/商密对比、代码抄袭剽窃判定、开源许可证 & 应用 & 领域核心 & 1 \\
\midrule
\multirow{2}{*}{\makecell[l]{社会伦理/隐私与公民\\权利}} 
& 事务数据库与普适监控隐私侵蚀、差分隐私数学原理、隐私增强技术 & 解释 & 计算机核心 & 2 \\
& 隐私权哲学法理渊源、重点法域隐私立法框架、表达自由 & 评估 & 领域核心 & 1 \\
\midrule
\multirow{2}{*}{\makecell[l]{社会伦理/专业沟通与\\团队协作}} 
& 源码文档提炼、API 手册编写、团队沟通规范、受众感知策略、无障碍表达 & 应用 & 计算机核心 & 2 \\
& 跨文化沟通能力、软件项目多维竞争价值观权衡、职场伦理倡导写作 & 应用 & 领域核心 & 1 \\
\midrule
\multirow{2}{*}{\makecell[l]{社会伦理/可持续性与\\绿色计算}} 
& 可持续工程实践、计算设备全生命周期生态影响、区块链共识能耗权衡 & 应用 & 计算机核心 & 1 \\
& 绿色 IT 与低碳软件设计标准、系统性生态效应、普适计算与智慧环保 & 评估 & 领域核心 & 1 \\
\midrule
\makecell[l]{社会伦理/计算机发展\\简史} 
& 机械计算前史、早期现代数字计算、微机革命、互联网与 Web、云计算与 AI & 解释 & 领域核心 & 1 \\
\midrule
\makecell[l]{社会伦理/计算经济学\\与产业生态} 
& 服务商生态与围墙花园、知识与注意力经济、人才供需、商业模式与定价、全球化分工与自动化冲击、数字金融变革 & 解释 & 领域核心 & 1 \\
\midrule
\multirow{2}{*}{\makecell[l]{社会伦理/安全策略、\\法律与网络犯罪}} 
& 计算机犯罪案例与法律追责、社会工程学欺诈、非法访问、黑客恶意软件 & 解释 & 计算机核心 & 2 \\
& 犯罪预防战略、红蓝对抗实操、编制企事业单位信息安全管理策略规范 & 应用 & 领域核心 & 1 \\
\midrule
\multirow{2}{*}{\makecell[l]{社会伦理/多元、\\公平与包容 (DEI)}} 
& 身份认同互塑、多元团队收益与同质化盲区、包容语言、无障碍开发 & 解释 & 计算机核心 & 2 \\
& 多元化行业榜样、多元化实证收益与缺乏多元化代价、技术霸权历史反思 & 评估 & 领域核心 & 2 \\
\bottomrule
\end{longtable}

\chapter{软件开发基础}

\section{领域导言与演进背景}

\begin{preamblebox}[知识领域概述与变革]
在整个计算机科学的知识版图中，熟练驾驭软件开发过程是每位学子立足的根本。为了最高效地利用计算机解决真实世界问题，学生必须能够流利地阅读、理解并运用多种主流编程语言编写稳健、规范的代码。然而，超越单纯的语法糖和基本编程技能，计算机专业人才更须具备敏锐的数据结构与算法选型能力，并熟练掌握现代软件工程开发与测试工具链。

SDF 知识领域汇聚了软件开发最为核心的基础概念与实操技能，主要面向计算机专业大一（第一学年）基础教学。它涵盖基本编程概念、内置数据结构的工程化应用、高质量编码规范与实践，以及算法对程序运行性能的直观影响。本领域包含的 43 学时核心内容构成了后续进阶专业课的基石。

\textbf{自 CS2013 以来的核心变革：}
\begin{enumerate}
    \item \textbf{强调现代语言原生数据结构赋能}：现代主流语言（如 Python、Java、C++、Rust 等）均内置了极高工业成熟度的关联容器（字典/映射、集合、动态切片）。CS2023 强烈倡导在早期即培养学生直接运用原生/库级高阶数据结构进行工程化解题的能力，避免过早陷入底层指针操作的泥潭。
    \item \textbf{全流程工程规范前置}：将测试驱动意识、单元测试框架、断点调试技术、代码静态审查与可读性重构要求，前置融入大一基础编程阶段。
    \item \textbf{明确划分 SDF 与 AL / SE 的协作边界}：SDF 聚焦于中小型程序的构建、常见数据结构使用与实证性能测试，高阶理论证明（复杂性下界、NPC、形式语义）留待 AL，大规模工业级生命周期管理留待 SE。
\end{enumerate}
\end{preamblebox}

\section{核心学时分配建议}

本领域内容全部为所有计算机专业学生所必需的计算机核心（计算机核心，共 43 学时）：

\begin{table*}[htbp]
\centering\small
\setlength{\tabcolsep}{2.5pt}\footnotesize\begin{tabularx}{0.88\textwidth}{Y c c}
\toprule
\textbf{知识单元} & \textbf{计算机核心} & \textbf{领域核心} \\
\midrule
基础编程概念与实践 & 20 & -- \\
基础数据结构应用 & 12 & -- \\
算法与实证效能 & 6 & -- \\
现代软件开发实践规范 & 5 & -- \\
\midrule
\textbf{合计} & \textbf{43 学时} & \textbf{0 学时} \\
\bottomrule
\end{tabularx}
\end{table*}

\section{知识单元与详细大纲}

\subsection{基础编程概念与实践}
\textbf{【学时建议】}：20 计算机核心学时。

\begin{itemize}
    \item \textbf{程序基础元素}：变量、作用域、生命周期、基本数据类型（整型、浮点型、布尔型、字符型）、表达式求值、赋值语义。
    \item \textbf{程序控制结构}：条件分支（if-else, switch/match 模式匹配）、循环控制、控制流跳转与循环不变量。
    \item \textbf{核心模块化抽象机制}：
    \begin{itemize}
        \item 函数/方法定义：参数传递机制（值传递、引用传递/指针传递、对象引用共享）、默认参数与可变参数。
        \item 作用域与命名空间：局部变量、全局变量、闭包上下文与变量遮蔽。
        \item 面向对象基础：类与对象的定义、数据封装、访问控制修饰符、简单继承与多态。
    \end{itemize}
    \item \textbf{现代复合数据类型}：序列结构（静态数组、动态数组/列表）、键值映射（字典/Map/哈希表）、集合（Set）、元组、不可变数据结构。
    \item \textbf{标准输入输出与外部交互}：控制台标准流 I/O、磁盘持久化文件读写、RESTful Web API 资源获取与 JSON 结构化解析。
    \item \textbf{递归思想与实现}：递归基准条件、递归工作栈、树状递归与递归向迭代转换。
    \item \textbf{运行时异常防御与处理}：异常抛出、捕获、异常层级与资源优雅释放。
    \item \textbf{代码心智模型与执行追踪}：堆栈内存布局、对象引用图谱、代码执行控制流跟踪。
\end{itemize}

\subsection{基础数据结构应用}
\textbf{【学时建议】}：12 计算机核心学时。

\begin{itemize}
    \item \textbf{标准抽象数据类型 (ADT) 运用}：列表（List）、栈（Stack，LIFO）、队列（Queue，FIFO）、双端队列、集合（Set）与字典（Dictionary）。
    \item \textbf{文本与字符串高级处理}：字符串不可变性、子串查找、正则表达式匹配与文本流切分。
    \item \textbf{针对业务场景的数据结构选型}：分析特定场景（高频写入 vs 频繁随机访问）下各种内置容器的渐进开销与内存布局。
    \item \textbf{实证性能度量}：编写基准评测代码，度量不同数据规模下容器操作的物理耗时。
\end{itemize}

\subsection{算法与实证效能}
\textbf{【学时建议】}：6 计算机核心学时。

\begin{itemize}
    \item \textbf{算法本质与效能度量}：算法定义、正确性验证直觉、时间与空间复杂性直观感知。
    \item \textbf{入门经典算法族}：
    \begin{itemize}
        \item 常用排序：选择排序、插入排序、快速排序思想与库级排序函数的使用。
        \item 经典查找：线性搜索、有序序列二分查找及其死循环规避。
        \item 树与图基础遍历：深度优先遍历（DFS）与广度优先遍历（BFS）的直观递归与队列实现。
    \end{itemize}
    \item \textbf{实测性能对比}：在真实硬件上对 $O(n)$ 与 $O(\log n)$、$O(n^2)$ 与 $O(n \log n)$ 算法进行规模压力评测。
\end{itemize}

\subsection{现代软件开发实践规范}
\textbf{【学时建议】}：5 计算机核心学时。

\begin{itemize}
    \item \textbf{工业级编码风格与可读性}：Clean Code 规范、自解释的变量与函数命名、模块职责单一原则、适度而精准的内联注释。
    \item \textbf{现代集成开发环境 (IDE) 与调试技术}：代码导航、符号查找、断点设置、条件断点、单步执行、变量监视与调用栈检查。
    \item \textbf{单元测试与测试用例设计}：基于现代测试框架（JUnit、pytest 等）编写断言；等价类划分、边界值分析与异常分支测试用例构造。
    \item \textbf{测试局限性反思}：深刻理解“测试可以证明程序存在缺陷，但永远无法证明程序没有缺陷”（Dijkstra 箴言）。
\end{itemize}

\section{专业素养与思维倾向}

\begin{itemize}
    \item \textbf{自主解决问题的韧性}：遇到语法或逻辑报错时，善于自主查阅官方文档、API 参考与技术社区排查方案。
    \item \textbf{实验探索精神}：勇于编写小 Demo 验证未知语言特性，熟练运用调试器单步推演探索程序行为。
    \item \textbf{对技术底层机理的好奇心}：不仅满足于代码能跑通，更渴望理解 IDE 内部运行机制、内存分配与编译流程。
    \item \textbf{严密的系统性与细密性 (Systematic \& Meticulous)}：在处理边界空指针、数组下标越界与输入非法校验时保持严谨。
\end{itemize}

\section{课程设计与教学实施建议}

联合工作组建议将 SDF 模块有机分布在高校第一学年的经典两门课中：
\begin{enumerate}
    \item \textbf{CS1：程序设计基础 I (Introductory Programming)}：
    覆盖 20 学时编程概念、12 学时数据结构应用、2 学时算法入门、2 学时测试调试实践及 2--4 学时工程伦理与规范。目标是使学生能独立开发包含函数抽象、对象交互与文件/网络 API 调用的百行级稳健程序。
    \item \textbf{CS2：数据结构与面向对象实战 (Data Structures \& OOP)}：
    深化承接 SDF 的后续 4 学时高级算法（与 AL 衔接）和 3 学时高阶软件工程规范（与 SE 衔接）。
\end{enumerate}

\section{真实能力任务规范}

\begin{taskbox}[任务 1：模块化功能组件的规范开发与单元测试]
\textbf{【任务情境】}：根据给定的业务功能规约（例如多格式日志解析器），设计并编写一个独立的函数或类。\\
\textbf{【能力陈述】}：编写高可读性、高聚合度的代码，覆盖所有入参边界，并编写单元测试验证所有分支覆盖。\\
\textbf{【能力领域】}：软件开发 \quad \textbf{【能力单元】}：系统实现与质量测试\\
\textbf{【达标等级】}：\textbf{开发}
\end{taskbox}

\begin{taskbox}[任务 2：复杂逻辑缺陷的断点调试与定位修复]
\textbf{【任务情境】}：某学生在线考试评分系统中存在偶发性并发计分错误和异常越界崩溃。\\
\textbf{【能力陈述】}：运用现代 IDE 的断点、调用栈溯源与变量跟踪功能，在没有源码注释的情况下快速定位 Bug 根源并重构修复。\\
\textbf{【能力领域】}：软件开发 \quad \textbf{【能力单元】}：系统调试与维护\\
\textbf{【达标等级】}：\textbf{评估与开发}
\end{taskbox}

\begin{taskbox}[任务 3：同侪代码评审 与代码异味排查]
\textbf{【任务情境】}：在敏捷团队中对同组开发者提交的代码拉取请求（PR）进行严格的同行评审。\\
\textbf{【能力陈述】}：阅读并理解第三方代码，准确指出其中的潜在逻辑漏洞、死循环边界隐患、不规范命名与可维护性缺陷，并给出具体的重构建议。\\
\textbf{【能力领域】}：软件工程与应用 \quad \textbf{【能力单元】}：代码评审与改进\\
\textbf{【达标等级】}：\textbf{评估}
\end{taskbox}

\begin{taskbox}[任务 4：基于第三方 RESTful API 的轻量级数据聚合程序]
\textbf{【任务情境】}：构建一个实时天气或交通状况聚合展示终端应用。\\
\textbf{【能力陈述】}：熟练调用网络请求库发起 HTTP 访问，安全处理网络超时、鉴权失败及 JSON 反序列化异常，将数据缓存至本地哈希表中并做格式化呈现。\\
\textbf{【能力领域】}：软件与应用 \quad \textbf{【能力单元】}：API 集成开发\\
\textbf{【达标等级】}：\textbf{开发}
\end{taskbox}

\section{附录：核心主题与技能要求等级对照表}

\begin{longtable}{p{0.20\textwidth} p{0.48\textwidth} c c c}
\toprule
\textbf{知识单元} & \textbf{核心知识点与具体主题} & \textbf{认知等级} & \textbf{分类} & \textbf{学时} \\
\midrule
\endfirsthead
\toprule
\textbf{知识单元} & \textbf{核心知识点与具体主题} & \textbf{认知等级} & \textbf{分类} & \textbf{学时} \\
\midrule
\endhead
\bottomrule
\endfoot
\multirow{2}{*}{\makecell[l]{软件开发基础/基础编程概念\\与实践}} 
& 变量、表达式、分支迭代控制、函数封装、异常与 I/O & 开发 & 计算机核心 & 18 \\
& 程序错误类型、断点调试技术、代码跟踪与阅读理解 & 应用 & 计算机核心 & 2 \\
\midrule
\multirow{2}{*}{\makecell[l]{软件开发基础/基础数据结构}} 
& 常用容器（列表、栈、队列、字典、集合）应用与字符串处理 & 开发 & 计算机核心 & 10 \\
& 容器选型权衡、时空性能实测与内存开销评估 & 评估 & 计算机核心 & 2 \\
\midrule
\multirow{2}{*}{\makecell[l]{软件开发基础/算法}} 
& 算法概念、基础查找与排序、树图简单遍历 & 解释 & 计算机核心 & 4 \\
& 算法时空效率的物理测量与规模敏感性实证评测 & 评估 & 计算机核心 & 2 \\
\midrule
\multirow{3}{*}{\makecell[l]{软件开发基础/软件开发实践}} 
& 工业级编码规范、代码风格与可读性重构 & 评估 & 计算机核心 & 1 \\
& 单元测试框架使用与测试用例设计（等价类与边界值） & 开发 & 计算机核心 & 2 \\
& 现代 IDE 工具链（断点调试、性能分析）熟练使用 & 应用 & 计算机核心 & 2 \\
\end{longtable}

\chapter{软件工程}

\section{领域导言与演进背景}

\begin{preamblebox}[知识领域概述与时代演进]
早在 20 世纪 70 年代初，软件工程先驱戴夫·帕纳斯（David L. Parnas）便提出了振聋发聩的经典定义：\textbf{“软件工程是多人协同构建多版本程序的过程”}。这一洞见揭示了软件工程的核心本质：如果说程序设计是决定我们编写单个程序能力的技能，那么软件工程则在\textbf{“时间”}与\textbf{“人员”}两个决定性维度上与单纯的编程产生了根本分野。

\textbf{维度一：人员协同}\\
软件工程项目本质上是一项团队协作事业。仅仅成为一名独来独往的编程高手是远远不够的。合格的软件工程师必须具备良好的沟通、倾听、同理心与协作能力。编程或许可以是一项高度个体化的智力活动，但软件工程从立项之初就是一项社会化协作活动，与专业操守、团队运作、跨职能交流和同侪评审紧密交织。

\textbf{维度二：跨越时间的生命周期（Multi-Version \& Lifespan）}\\
软件工程项目通常是“多版本演进”的。软件拥有其预期生命周期，往往需要在数月、数年甚至数十年中持续稳定运转。在此期间，为了满足不断变化的产品需求，功能特性会被频繁增删；底层的计算基础设施、操作系统、编译器与编程语言版本会不断升级演变，外部第三方依赖库亦会频繁迭代。与课堂上“写完即弃”的单人作业相比，这种直面时间流逝与环境变迁的现实是软件工程特有的挑战：\textbf{仅仅构建一个当前能跑通的程序是远远不够的，系统必须能够经受住时间的考验，在漫长的演化中始终保持正确运转}。当今技术界最棘手的一系列工程挑战，其共同根源皆在于“时间必然催生变化”：向后兼容性、版本偏差/漂移、依赖地狱与依赖管理、数据库模式演变以及通信协议演进。

\textbf{软件工程高等教育的特有挑战：}\\
在象牙塔的学术环境中开展软件工程教学面临着独特的结构性难题。鉴于软件工程与传统编程的核心差异在于“时间”与“团队”，在有限的学期周数内设计出既能充分展现团队协作摩擦、又能真实还原长周期演进风险的教学实验极其困难。此外，软件工程中的许多高阶主题，只有当学习者在真实工业场景——无论是前沿企业实习、深度参与大型开源社区、还是全职工程岗位——中亲身体验过大规模协作与长期代码维护后，才能产生最深刻的共鸣与领悟。哪怕仅一个月的全职实战经验，其有效工程工时往往就已超过一门普通大学课程的全学期学时。

因此，CS2023 软件工程课程规范明确聚焦于：\textbf{绝大多数计算机专业应届毕业生走向工业界首日即迫切需要的核心理念}。这些内容要么对于仅受过传统单人编程训练的学生而言是全新盲区，要么是工业界日常协作中习以为常、却极少被显式系统传授的抽象工程法则。涵盖的核心主题包括：
\begin{itemize}
    \item 工业级软件测试与质量验证
    \item 团队协作与冲突消解
    \item 专业化工程沟通与代码评审
    \item 软件架构、API 契约与详细设计
    \item 软件维护、重构与长周期演化
    \item 现代自动化工程工具链
\end{itemize}

\textbf{教学模式权衡：讲授与项目驱动的辩证融合}\\
部分理论内容非常适合采用标准课堂讲授或“讲授+实验”的形式进行剖析。深入探讨版本控制系统的底层理论模型（如 Git 的有向无环图 DAG）、工业界主流分支演进策略、静态代码分析与动态插桩分析的理论边界、以及依赖拓扑网络中的“菱形依赖问题”，能够帮助学生建立扎实的技术心智模型。

另一方面，软件工程的大量核心精髓必须在亲手实践中才能内化。长期以来，大作业项目驱动型课程是此类教学的通用载体。然而，CS2023 特别指出其中的潜在教学风险：\textbf{如果学生唯一的协作经验是与其他未经受过工程训练的新手组成团队去完成一个单一项目，他们极易对团队协作的真正复杂性与最佳实践产生认知偏差与错误归因}。因此，任课教师必须在理论沉淀与项目实践之间审慎权衡。规范强烈建议项目课程倡导“敏捷迭代与极速反馈”——重复演练若干小步快跑的规范任务（例如遵循敏捷模式构建），其教学效果远优于让新手在高度摩擦中狼狈应对一个庞大单体项目。若学校不具备开设长周期综合项目的条件，任何能够让学生体验协作与时间维度的轻量级设计都极具价值：例如向大型既有代码库贡献特性、不同小组分别开发系统的异构模块并通过严格契约对接、结对分别编写编解码器进行联调互测等。

\textbf{时代背景与团队工程的必然性：}\\
所有现实证据表明，软件在整个人类社会运行中的中枢地位仍将持续高速膨胀，而“车库里两个极客黑客敲出一套伟大软件”的英雄主义神话时代早已宣告终结。当今工业界所有重要的软件系统无一例外都是高度组织化团队协同的产物，绝大多数开发工作都是在既有庞大代码资产之上进行演进与复用。软件工程教学是计算机专业学生从象牙塔走向现实工程世界的决定性桥梁——我们必须让每位学子深刻认识到：在工业实践中，几乎没有任何现实软件是作为单打独斗的个人项目从零编写的。沟通、协作、规划、测试与工具链，才是支撑他们走出校园并在广阔现实世界中立足的真正核心基石。

\textbf{自 CS2013 以来的核心变革：}
\begin{enumerate}
    \item \textbf{高度聚焦现代工业软件工作流}：大幅强化版本控制、自动化测试套件、规范化同侪代码评审以及现代 CI/CD 自动化工具链的必修要求。
    \item \textbf{弱化传统宏观项目管理与团队领导力}：大幅削减以往偏向高级项目经理与团队主管的抽象管理说教（如大型组织架构规划、重型软件过程度量），使课程精准匹配应届毕业生的首要岗位诉求。
    \item \textbf{聚焦团队参与、专业沟通与协同编码}：将核心落脚点全面转向作为团队普通成员如何高效参与敏捷协作、开展无指责技术沟通、规范处理代码冲突与接口演进。
\end{enumerate}
\end{preamblebox}

\section{核心学时分配建议}

本知识领域在“向日葵模型”下的核心学时推荐如下表所示（计算机核心共 6 学时，领域核心共 21 学时，合计 27 核心学时）：

\begin{table}[H]
\centering
\setlength{\tabcolsep}{3.5pt}\fontsize{8pt}{9.6pt}\selectfont\begin{tabularx}{\columnwidth}{Y c c}
\toprule
\textbf{知识单元} & \textbf{计算机核心} & \textbf{领域核心} \\
\midrule
团队协作与沟通 & 2 & 2 \\
工具与开发环境 & 1 & 3 \\
产品需求工程 & -- & 2 \\
软件架构与系统设计 & 1 & 4 \\
软件构造与工程实现 & 1 & 3 \\
软件验证与确认 & 1 & 3 \\
代码重构与软件演进 & -- & 2 \\
软件可靠性工程 & -- & 2 \\
形式化方法 & -- & \emph{拓展选修} \\
\midrule
\textbf{合计} & \textbf{6 学时} & \textbf{21 学时} \\
\bottomrule
\end{tabularx}
\end{table}

\section{知识单元与详细大纲}

\subsection{团队协作与沟通}
\textbf{【学时建议】}：2 计算机核心学时，2 领域核心学时。

由于传统初级编程教学的个体化特征，绝大多数初修软件工程的学生极少或从未接触过软件工程的协作本质。单纯的实战练习（如课程大作业）虽有助益，但必须在课堂讲授与研讨中系统注入关于清晰、高效、专业沟通与团队协作核心价值的深度剖析。

\subsubsection{计算机核心主题}
\begin{itemize}
    \item \textbf{高效团队沟通机制}：同步沟通与异步沟通的权衡、技术备忘录撰写、沟通中的倾听与建设性反馈机制。
    \item \textbf{团队冲突根因与消解策略}：技术分歧与人际摩擦的本质辨析、对事不对人的专业态度（针对技术提案和代码实现提出批评，坚决避免人身攻击）、团队冲突仲裁与共识达成路径。
    \item \textbf{协同编程工程实践}：
    \begin{itemize}
        \item 结对编程：驾驶员与领航员角色轮换、心智协同与即时审查。
        \item 同侪代码评审：评审礼仪、建设性评注、评审清单、小粒度提交与及时响应。
    \end{itemize}
    \item \textbf{软件团队中的角色与职责划分}：
    \begin{itemize}
        \item 团队协同的规模放大优势与互补效应。
        \item 协作带来的沟通成本、同步阻碍与组织复杂性（人月神话与布鲁克斯定律直觉）。
    \end{itemize}
    \item \textbf{团队开发流程与敏捷过程规范}：
    \begin{itemize}
        \item 任务认领与责任追踪。
        \item 敏捷工作量估算与故事点。
        \item 高效站会与会议结构（Meeting Structure、议程管理、时间盒机制）。
        \item 工作排期与迭代节奏。
    \end{itemize}
    \item \textbf{团队多元化的战略价值}：认知背景多元、技能互补、视角包容度对软件设计鲁棒性与创新的巨大推动。
\end{itemize}

\subsubsection{领域核心主题}
\begin{itemize}
    \item \textbf{作为团队整体与外部利益相关者高效对接}：
    \begin{itemize}
        \item 管理层及其他非技术业务团队（产品、运营、法务、市场营销）。
        \item 商业客户（客户需求协商、合同边界澄清）。
        \item 最终用户（真实使用体验反馈、可用性痛点挖掘）。
    \end{itemize}
    \item \textbf{物理集中、全分布式、混合制与虚拟远程团队的系统性风险}：
    \begin{itemize}
        \item 跨时区与异步通信断层、认知脱节与孤岛效应。
        \item 组织架构脆弱性、单点人员故障与技术断层。
        \item 远程协作风险的缓解策略、应急恢复预案与团队心理健康维护。
    \end{itemize}
\end{itemize}

\subsubsection{示范性学习成果与能力目标}
\begin{itemize}
    \item 严格遵循高效团队专业沟通规范进行技术交流与工作协同。
    \item 阐明团队冲突的根源、潜在危害与建设性价值——特别聚焦于能够在不攻击他人人格的前提下，就技术思想或架构方案开展富有成效的良性争鸣。
    \item 在真实的团队工程环境中主导或推动行之有效的冲突化解策略。
    \item 在协同开发与结对编程过程中展现高水准的协作效能与默契度。
    \item 在软件开发团队中合理提议、协商并分配关键角色与工程职责。
    \item 针对团队技术评审或敏捷例会，撰写结构严谨的会议议程并在会议中严格执行。
    \item 深入参与团队综合项目，全面践行团队建设核心要素、建立健康包容的团队工程文化，并制定与推进团队工作执行计划。
    \item 深刻阐释并积极推动多元化与包容性为软件研发团队带来的深远工程益处与创新潜能。
    \item 阐明软件团队在技术与非技术双重层面与外部利益相关者建立定期沟通机制的重要性与核心策略。
    \item 系统梳理集中式、全分布式、混合制与虚拟远程团队所面临的特有工程风险与潜在故障点，并制定相应的风险缓释、容错与故障复盘机制。
\end{itemize}

\subsection{现代软件工程工具与开发环境}
\textbf{【学时建议】}：1 计算机核心学时，3 领域核心学时。

在过去的一代人时间里，整个软件工业界对软件工程工具链的依赖呈现爆发式增长：分布式版本控制系统已无处不在，自动化测试框架成为标配，静态代码分析与动态插桩检测融入日常，持续集成（CI）系统更是成为工程交付的生命线。与此相伴，日益强大的集成开发环境（IDE）提供了毫秒级代码语义检索、全局索引、自动化安全重构以及对各类外部工程工具的无缝集成。深刻理解这些工业级工具的运行机制与选型权衡——尤其是版本控制系统——是现代软件工程师不可或缺的底层素养。

\subsubsection{计算机核心主题}
\begin{itemize}
    \item \textbf{软件配置管理（SCM）与版本控制基础}：
    \begin{itemize}
        \item 代码与配置纳管：在版本控制中统一管理配置、可复现构建与环境即代码。
        \item 分支演进策略：功能开发分支vs 生产发布分支；工业级“基于主干的开发”范式。
        \item 代码分支集成策略：合并与变基的底层原理、适用场景与提交历史整洁性权衡。
    \end{itemize}
\end{itemize}

\subsubsection{领域核心主题}
\begin{itemize}
    \item \textbf{发布管理}：发布流水线、灰度发布、版本打标与发布分支维护。
    \item \textbf{专业测试工具链}：单元测试执行器、测试框架、静态代码分析器与动态内存/并发分析工具。
    \item \textbf{软件研发全流程自动化}：
    \begin{itemize}
        \item 构建系统：高速增量构建、密封隔离构建（Hermetic Builds，消除宿主机隐式污染）、构建缓存、现代主流构建工具体系对比（Make、Gradle、Bazel 等）。
        \item 持续集成 (CI)：利用自动化构建流水线与自动化测试套件作为工程“单一事实来源”，即时验证主干最新代码的构建可达性与基础测试通过率。
        \item 依赖管理：上游与第三方依赖库自动解析、依赖冲突消解、包管理器生态、语义化版本控制（SemVer，主版本号.次版本号.修订号）的严格契约。
    \end{itemize}
    \item \textbf{设计与协作沟通工具}：架构设计文档系统、图表绘制工具、通用建模语言规范（如 UML 类图、时序图、架构拓扑图）。
    \item \textbf{现代工具集成概念与机制}：Webhook、自动化钩子脚本、流水线编排协议与工具间上下文打通。
    \item \textbf{现代 IDE 高级工程套件}：全项目符号索引与毫秒级搜索、交互式断点调试、自动化代码重构、代码片段模板与辅助工具深度集成。
\end{itemize}

\subsubsection{示范性学习成果与能力目标}
\begin{itemize}
    \item 深入辨析集中式配置管理系统（如 SVN）与分布式版本控制系统（如 Git）在架构拓扑、协作效率与容灾能力上的核心差异。
    \item 阐述如何综合运用分支管理、标签与版本控制策略支持工业级软件发布管理。
    \item 在小型团队协作项目中准确识别关键配置项，并熟练使用分布式源代码版本控制工具完成协同开发。
    \item 阐明主流静态代码检查工具与动态测试工具如何无缝集成至开发环境与日常编码流水线中。
    \item 深刻理解持续集成（CI）系统作为团队共享主干代码健康状态唯一基准事实源的战略地位。
    \item 全面阐述在为特定软件系统进行工具链选型时的核心考量要素，涵盖需求追踪、架构建模、编码实现、构建自动化与自动化测试等各个环节。
    \item 展现熟练运用现代化软件工具链支持中等规模软件产品全流程开发的工程实践能力。
\end{itemize}

\subsection{产品需求工程}
\textbf{【学时建议】}：2 领域核心学时。

知道“如何构建系统”固然重要，但如果我们不知道“该构建什么”，一切技术投入都将付诸东流。产品需求（亦称需求工程、产品设计、需求诱导、产品需求规格说明 PRD 等）系统引导学生学习围绕定义新产品或新功能特性所必需的完整规范化流程。

\subsubsection{领域核心主题}
\begin{itemize}
    \item \textbf{功能需求的结构化与场景化表述}：
    \begin{itemize}
        \item 运用用例或用户故事（User Stories，如“作为……我希望……以便于……”格式）精确描述功能需求。
        \item 掌握至少一种规范化记录与组织功能需求的方法体系。
        \item 深刻理解需求如何作为后续架构设计、编码实现与验收测试的直接输入。
        \item 辩证评估特定需求描述方法的表达力优势与潜在局限。
    \end{itemize}
    \item \textbf{需求规约的核心质量属性}：一致性、合法有效性、完备性与工程可行性。
    \item \textbf{需求诱导与获取}：
    \begin{itemize}
        \item 需求来源辨析：最终用户、系统管理员、运维人员、合规审查方与业务支持人员。
        \item 需求采集方法：问卷调查、深度访谈、焦点小组、现场观察与埋点用户行为大数据分析。
    \end{itemize}
    \item \textbf{非功能需求与质量属性}：系统安全性、可用性/易用性、高并发性能、可靠性、可维护性与合规隐私要求（交叉 SEC 知识领域）。
    \item \textbf{需求阶段的风险识别与风险管理}：需求模糊度评估、范围蔓延防御与技术可行性预警。
    \item \textbf{需求规格说明书的专业沟通与形式化}：半形式化规范、验收准则与契约化签署。
\end{itemize}

\subsubsection{拓展选修主题}
\begin{itemize}
    \item \textbf{敏捷原型构建}：低保真线框图、高保真交互原型与最小可行产品（MVP），作为需求启发、推演与用户确认的核心工具。
    \item \textbf{产品需求演进与变更控制}：当业务需求发生不可逆变更时，系统化评估波及范围、依赖链路与需修改的系统模块。
    \item \textbf{敏捷工程估算}：
    \begin{itemize}
        \item 掌握科学评估开发任务工作量的经典技术（斐波那契故事点、计划扑克 Planning Poker、专家 Delphi 法）。
        \item 在实训中持续开展工作量预估，并与实际消耗工时进行复盘校准。
        \item 认识到软件估算的极高难度——即使学生初期存在显著偏差，亲历估算偏差并在复盘中修正认知模型亦极具教学价值。
    \end{itemize}
\end{itemize}

\subsubsection{示范性学习成果与能力目标}
\begin{itemize}
    \item 沿采集成本、覆盖广度、主观偏差等多个正交维度系统对比不同的需求获取方法。
    \item 深入辨析两种功能需求描述方法（如传统用例规范 vs 敏捷用户故事/客户访谈）的本质差异及其各自适配的工程场景。
    \item 给定一组明确的产品需求规范与一组候选系统行为，准确裁定哪些系统行为是被强制要求的、被允许的或被严格禁止的。
    \item 针对一个简易软件系统，独立完成全套业务需求的采集、分类与系统化整理。
    \item 给定现有系统的架构描述与一组新增业务需求，准确识别出系统中必须发生适应性调整的核心模块与代码范围。
    \item 从纷繁杂乱的需求集合中，精准辨识并分离功能需求与非功能需求（质量属性）。
    \item 针对待建软件系统快速构建原型（Mock-up 原型或 MVP），以实证方式验证关键业务需求集的可行性。
    \item 对一组研发任务预估工时，并在项目实施后将初始预估值与实际消耗工时进行量化复盘对比。
    \item 依据任务间的逻辑依赖关系制定合理的任务实现时序路线图，以“尽早消除核心技术风险”为第一准则组织开发。
    \item 为简易软件系统撰写专业、规范且具备可验证性的软件需求规格说明书（SRS/PRD）。
\end{itemize}

\subsection{软件架构与系统设计}
\textbf{【学时建议】}：1 计算机核心学时，4 领域核心学时。

产品需求聚焦于软件系统面向最终用户的“外部功能”，而软件设计则聚焦于面向软件工程师的“内部组件拓扑与结构逻辑”。本单元全面涵盖宏观的软件体系结构设计以及微观层面的 API 接口与数据契约设计。

\subsubsection{计算机核心主题}
\begin{itemize}
    \item \textbf{系统设计核心原则}：
    \begin{itemize}
        \item 抽象层次划分：宏观架构设计与微观详细设计。
        \item 关注点分离：业务逻辑、数据持久化与用户展现界面的解耦。
        \item 信息隐藏：隐藏内部数据表示与算法细节，仅暴露稳定抽象。
        \item 耦合度与内聚度：追求“高内聚、低耦合”的工程本质与量化直觉。
    \end{itemize}
    \item \textbf{经典软件体系结构与设计范式}：
    \begin{itemize}
        \item 核心设计范式：自顶向下功能分解/分层架构、面向数据架构、面向对象分析与设计（OOAD）、事件驱动架构。
        \item 经典标准体系结构：客户端/服务器（C/S）、多层企业架构、管道-过滤器、模型-视图-控制器（MVC）。
        \item 组件边界界定与依赖拓扑图谱识别。
    \end{itemize}
    \item \textbf{大规模编程与小规模编程的本质分野}：Programming in the Large（聚焦模块交互、契约制定、团队协作与边界治理）vs. Programming in the Small（聚焦单体算法逻辑、局部循环与数据结构操纵）。
    \item \textbf{代码异味与内在质量表征}：辨析代码“功能正确性”与“内在工程质量/可维护性”的本质区别（过长函数、上帝类、霰弹式修改、重复代码等）。
\end{itemize}

\subsubsection{领域核心主题}
\begin{itemize}
    \item \textbf{工业级 API 设计核心准则}：
    \begin{itemize}
        \item 一致性：接口命名、参数顺序与错误语义在系统内部严格统一，符合领域通用惯用法（降低学习成本，减少误用概率）。
        \item 可组合性：小而美的正交接口可通过流水线或组合轻松构建复杂能力。
        \item 契约式设计与文档规范：清晰界定接口对系统状态的变更影响而非泄露其内部实现细节；显式声明前置条件、后置条件与类不变量。
        \item 可扩展性：在不破坏现有调用的前提下为未来能力拓展预留平滑空间。
        \item 健壮的错误报告机制：错误语义明确、可预测且具备可操作性；非法入参必须显式报错拒绝；能够内部安全自愈的偶发故障无需过度上报。
    \end{itemize}
    \item \textbf{识别并形式化数据不变量与时间/状态不变量}：状态机转换边界、并发序列时序约束。
    \item \textbf{软件设计的结构模型与行为模型}：类拓扑图、调用关系网、状态迁移图。
    \item \textbf{数据架构设计}：核心数据结构选型、持久化存储引擎选型与交互协议（交叉 DM 知识领域）。
    \item \textbf{需求可追溯性}：建立业务需求项与架构设计组件之间的双向追溯矩阵。
\end{itemize}

\subsubsection{拓展选修主题}
\begin{itemize}
    \item \textbf{高级设计建模方法}：统一建模语言（UML）类图、实体-联系图（E-R 图）、时序图与组件部署图。
    \item \textbf{软件设计质量的量化度量与技术债务评估}：圈复杂度、模块不稳定性、抽象度指标。
    \item \textbf{安全设计与编码原则（Secure Design Principles，交叉 SEC 领域）}：
    \begin{itemize}
        \item 最小特权原则。
        \item 故障安全默认原则。
        \item 心理可接受度原则。
    \end{itemize}
    \item \textbf{系统设计多维权衡决策}：执行效率 vs. 可靠性、极致安全 vs. 操作便捷性等维度的科学决断。
\end{itemize}

\subsubsection{示范性学习成果与能力目标}
\begin{itemize}
    \item 给定一个高层架构方案，准确辨析其采用的标准软件体系结构模式（如 MVC、分层、微内核、事件驱动）。
    \item 针对简易软件系统选择并运用适宜的设计范式开展系统设计，并清晰阐明关注点分离、信息隐藏等设计原则在其中的贯彻落实。
    \item 识别并重构存在严重缺陷的系统设计方案，消除违反关注点分离与信息隐藏的设计弊端。
    \item 在给定的宏观架构设计中，精确识别所有软件组件之间的依赖拓扑与调用链条。
    \item 为大型软件系统中的独立组件设计生产级 API，包括形式化定义并归档其所有操作的不变量、调用契约与异常错误条件。
    \item 从一致性、可组合性与可扩展性等维度，严谨评估一份 API 设计规约的工程优劣。
    \item 对现有系统设计进行平滑扩展，在不破坏既有架构约束的前提下接纳全新业务功能模块。
    \item 设计一套兼顾访问效率与内存开销的高质量内部数据结构，以完美实现给定的公共 API 接口表面。
    \item 审核给定的软件架构设计，精确判定其完全满足了哪些原始业务需求项。
    \item 将自然语言描述的业务系统设计规约准确转化为规范的 UML 类图与交互时序图。
    \item 调整并加固存在安全缺陷的架构设计，使其严格遵循最小特权原则与故障安全默认原则。
    \item 从执行性能、资源消耗、可靠性与易用性等多元质量维度，系统对比与评估两种不同的软件设计方案。
\end{itemize}

\subsection{软件构造与工程实现}
\textbf{【学时建议】}：1 计算机核心学时，3 领域核心学时。

软件构造聚焦于直接指导高质量代码产出的工业实践：实用的测试驱动、规范的代码文档、严谨的编程规范、异常处理与防御性实现。进阶主题涵盖安全编码、依赖注入、任务时序与生产级既有代码库演进。

\subsubsection{计算机核心主题}
\begin{itemize}
    \item \textbf{实用小粒度单元测试实践}：
    \begin{itemize}
        \item 单元测试：针对独立函数或单一类型编写具备高分支覆盖率的自动化测试用例。
        \item 测试驱动开发 (TDD)：TDD 对学生具有不可替代的心理学教育价值——在编写实现之前，预先挑战自己去探寻目标 API 的刁钻边界、异常输入与极端测试用例，其心智阻力显著更小。如果先写完实现，开发者往往会本能地避开破坏自己心血作品的极端输入；而“测试先行”则让学生在发现潜在问题中获得智力满足，并在开发过程中亲眼目睹测试用例从红变绿，获得极大的工程成就感。
    \end{itemize}
    \item \textbf{工业级工程文档编写规范}：
    \begin{itemize}
        \item 接口契约文档：准确描述接口的功能约束、前后置条件与数据不变量。
        \item 实现级文档规范：重点聚焦于棘手、微妙且非直观的代码逻辑（无论是因为运用了高级语言特性还是业务逻辑本身高度复杂）。坚决杜绝在显而易见的简单语法或通用操作上添加无效注释；注释应致力于阐明数据流向与设计意图（代码在做什么、为什么这样做）；对于晦涩代码应优先通过重构使其自解释，必要时辅以精准注释点拨核心玄机。
    \end{itemize}
\end{itemize}

\subsubsection{领域核心主题}
\begin{itemize}
    \item \textbf{工业级编码风格与规范}：
    \begin{itemize}
        \item 团队代码规范（Style Guides，如 Google/PEP8 风格指南）。
        \item 规范化注释、代码格式化工具统一约束。
        \item 严谨自解释的标识符命名规范。
    \end{itemize}
    \item \textbf{高质量编码“最佳工程实践”}（交叉 安全隐私/防御性编程与 SDF 知识领域）：
    \begin{itemize}
        \item 防御性编程实践：输入合法性断言、边界严格截断、空值防御。
        \item 安全编码核心原则：抵御注入攻击、权限提升防御。
        \item 运用现代异常处理机制提升软件系统的健壮性与容错能力。
    \end{itemize}
    \item \textbf{代码调试技术}：符号调试器、条件断点、核心转储分析与调用栈溯源。
    \item \textbf{第三方成熟库与开源框架的工程化集成}：依赖评估、许可证合规、封装适配。
\end{itemize}

\subsubsection{拓展选修主题}
\begin{itemize}
    \item \textbf{更大规模的工程级测试技术}：
    \begin{itemize}
        \item 测试替身：桩代码、模拟对象、伪实现的本质差异与选型。
        \item 依赖注入 (DI)：通过控制反转解耦组件依赖，大幅提升单元测试的可测性。
    \end{itemize}
    \item \textbf{研发任务时序编排与里程碑规划}：
    \begin{itemize}
        \item 任务依赖识别：准确梳理复杂工程中不同任务间的上下游强依赖拓扑。
        \item 里程碑设定：设定具有明确功能闭环与交付价值的进度检查点。
        \item 风险前置消减：敏锐识别项目中风险最高的未知领域，优先安排排雷任务。
    \end{itemize}
    \item \textbf{典型底层安全漏洞防御}：缓冲区与整型溢出、数据竞争与并发竞态条件、特权滥用与未安全初始化、外部不可信输入校验。
    \item \textbf{自动化 API 文档生成技术}（如 Doxygen, Javadoc, Sphinx）。
    \item \textbf{开发工程上下文辨析}：全新“绿地开发”vs. 工业级“既有遗留代码库演进”——变更波及影响分析与渐进式变更落地。
    \item \textbf{软件发布管理工程流}。
\end{itemize}

\subsubsection{示范性学习成果与能力目标}
\begin{itemize}
    \item 针对小型软件组件（包含若干函数或单一实体类），编写完备、规范的自动化单元测试用例。
    \item 为小型软件组件撰写规范详尽的对外接口文档，并在代码内部关键复杂逻辑处编写精准的实现注释。
    \item 阐述为实现高可靠性、高执行效率与鲁棒容错设计所必需的技术、编码习惯用法与底层机制。
    \item 熟练运用语言级异常处理机制，编写具备弹性恢复能力的鲁棒代码。
    \item 详细阐述工业级安全编码规范与防御性编程的核心指导原则。
    \item 在小型软件开发项目中选择并严格贯彻执行统一的代码风格规范。
    \item 系统对比与辨析组件集成策略，包括自顶向下、自底向上以及三明治（混合）集成模式的优劣。
    \item 阐述在针对特定项目全新开发的定制代码库中，进行代码变更分析与规范化落地的全流程。
    \item 阐述针对工业界庞大复杂的既有遗留代码库进行代码变更分析与安全实施的工程规范。
    \item 重构一段存在隐患的简易程序，彻底根除其中的常见安全漏洞（如缓冲区溢出、整数溢出和竞态条件）。
    \item 编写一个能够执行非平凡复杂任务、且对外部非法输入和运行时偶发异常具备极高容忍弹性的软件组件。
\end{itemize}

\subsection{软件验证与确认 (软件测试)}
\textbf{【学时建议】}：1 计算机核心学时，3 领域核心学时。

软件验证与确认（V\&V）聚焦于如何最大化发挥测试的工程价值——深刻理解测试的战略定位、系统的失效模式、以及优秀测试与劣质测试之间的本质差异。

\subsubsection{计算机核心主题}
\begin{itemize}
    \item \textbf{验证与确认的哲学辨析}：
    \begin{itemize}
        \item 验证：“Are we building the thing right?”（我们是否以正确的方式构建了系统？系统是否符合预设规格说明？）。
        \item 确认：“Did we build the right thing?”（我们是否构建了正确的系统？系统是否真正解决了用户的实际需求？）。
    \end{itemize}
    \item \textbf{测试为何至关重要}：软件组件在漫长周期的持续代码演进中能否始终保持功能正常？测试是重构与演化的唯一安全气囊。
    \item \textbf{多维测试目标体系}：系统易用性、运行可靠性、规约一致性、系统吞吐与响应性能、全栈安全性（Security，交叉 SEC 领域）。
    \item \textbf{工业级测试层级划分}：单元测试、集成测试、确认/验收测试、端到端系统测试。
    \item \textbf{测试代码与生产代码的代码风格哲学差异}：
    \begin{itemize}
        \item DAMP（Descriptive And Meaningful Phrases，富有描述性与自解释性）vs. DRY（Don't Repeat Yourself，不要重复自己）。
        \item 在测试代码中，教条地消除一切重复反而会破坏用例的自包含性与可读性；为了使失败用例的上下文一目了然，适度的结构性代码重复在测试工程中是完全正当且值得提倡的。
    \end{itemize}
\end{itemize}

\subsubsection{领域核心主题}
\begin{itemize}
    \item \textbf{测试规划与用例设计生成}：
    \begin{itemize}
        \item 基于形式化模型与需求规约生成严密的测试用例。
        \item 多维测试覆盖度：测试矩阵；代码覆盖率（语句、分支、路径覆盖）；环境覆盖度（跨硬件 CPU 架构、异构 OS、主流浏览器与运行时版本）。
        \item 测试数据设计与边界输入生成（等价类划分、边界值分析）。
    \end{itemize}
    \item \textbf{现代测试开发技术体系}（交叉 软件开发基础/开发方法与 软件工程/软件构造）：
    \begin{itemize}
        \item 测试驱动开发（TDD）深度实践。
        \item 面向对象测试、测试打桩与依赖注入。
        \item 黑盒测试与白盒测试技术。
        \item 现代化测试工具集成：代码覆盖率探针、静态代码扫描、模糊测试（Fuzzing）。
    \end{itemize}
    \item \textbf{研发全生命周期中的 V\&V 闭环}：
    \begin{itemize}
        \item 同侪代码评审。
        \item 测试自动化与工具链全自动调度。
        \item 提交前门禁测试vs. 提交后持续集成测试。
        \item 测试深度覆盖率与测试执行吞吐量/排队延迟之间的工程权衡。
        \item 缺陷追踪、严重级别划分与缺陷重现机制。
    \end{itemize}
    \item \textbf{特定前沿领域的 V\&V 独特挑战}：高并发与性能基准测试、异步与分布式并行系统、高可靠安全攸关系统、高精度数值计算验证。
\end{itemize}

\subsubsection{拓展选修主题}
\begin{itemize}
    \item \textbf{高级 V\&V 工具与自动化进阶}：静态分析引擎、模糊测试（Fuzzing）、动态内存/竞态检测工具、缺陷日志自动化分析。
    \item \textbf{测试终止准则与缺陷密度估算}：基于缺陷播种的剩余故障估算、伪随机数在测试生成中的科学运用。
    \item \textbf{性能基准评测与剖析技术}：
    \begin{itemize}
        \item 吞吐量与延迟量化。
        \item 极限负载下的降级与压力测试（FIFO 与 LIFO 请求调度策略）。
        \item 加速比与强/弱可扩展性计算（阿姆达尔定律 Amdahl's Law、古斯塔夫森定律 Gustafson's Law）。
        \item 性能评估关键指标确立。
        \item 常见性能瓶颈诊断：计算受限、内存带宽受限、I/O 与网络延迟受限。
        \item 基准测试的统计学最佳实践：测量不确定度估算、置信区间、多轮采样消除噪声抖动、高精度硬件时钟计时挑战。
        \item 性能测试数据的专业可视化图表呈现。
    \end{itemize}
    \item \textbf{复杂异步、并行与分布式并发系统的测试方案}。
    \item \textbf{非代码工程工件（技术文档、培训手册、合规凭据）的验证与确认}。
\end{itemize}

\subsubsection{示范性学习成果与能力目标}
\begin{itemize}
    \item 深刻阐释测试在软件生命周期与长期演进中的核心战略价值。
    \item 严格辨析软件“验证”与“确认”的本质哲学与工程差异。
    \item 详细阐述软件测试涵盖的多维核心目标（可用性、可靠性、性能、安全等）。
    \item 系统对比不同类型与层级的测试策略（回归测试、单元测试、集成测试、系统测试与用户验收测试）。
    \item 描述制定完备测试计划与高效生成测试用例的核心工程方法。
    \item 针对一段中等规模代码，制定涵盖测试矩阵、边界测试数据与异常输入的完整测试计划。
    \item 落地执行该测试计划，编写并运行自动化测试套件。
    \item 阐明测试驱动开发（TDD）的核心原则，并解释自动化测试在 TDD 循环中扮演的核心支柱角色。
    \item 深入剖析面向对象软件测试所面临的特有挑战（多态分派、隐式状态、对象生命周期）。
    \item 阐述模拟测试与依赖注入（DI）的机制及其在隔离外部依赖中的具体应用。
    \item 作为团队协作成员，对同事提交的中等规模代码段执行专业严谨的同侪代码评审。
    \item 阐释现代化工具链在软件质量确认与合规性证明中所能发挥的关键作用。
    \item 为小型软件项目搭建全自动化测试执行与报表流水线。
    \item 深入分析提交前测试与提交后测试各自的定位、收益与算力开销权衡。
    \item 辩证研讨测试用例覆盖率深度与测试流水线执行吞吐量/反馈时延之间的工程权衡及其对敏捷开发的影响。
    \item 熟练运用工业级缺陷追踪系统在小型项目中闭环管理软件缺陷的生命周期。
    \item 探讨传统经验测试在并发、分布式与安全攸关等特定前沿领域中所暴露的理论与实践局限。
    \item 详细阐释并横向对比多种高级验证确认工具的技术原理与适用场景。
    \item 在小型软件项目中实现多种静态扫描与动态测试工具的自动化集成执行。
    \item 说明在测试用例生成中如何以及何时科学使用随机数，并规避测试不可复现风险。
    \item 阐明软件潜在残留故障估算的核心数学理论与工程方法。
    \item 基于缺陷播种法与缺陷密度历史数据，科学估算小型软件应用中的剩余缺陷数量。
    \item 清晰定义吞吐量与响应时延，并针对具体系统架构各列举典型实例。
    \item 解释计算加速比与不同形式的扩展性（强扩展与弱扩展）的计算模型与理论边界。
    \item 准确诊断并归纳常见的计算密集型、内存受限型与 I/O 延迟型性能瓶颈。
    \item 详述构建工业级性能基准测试时的统计学方法准则与实验设计规范。
    \item 阐明在高精度性能测量中高分辨率计时的技术方案与面临的系统抖动挑战。
    \item 为小型项目确立关键性能评价指标，构建并运行基准测试，完成统计学显著性分析并绘制专业可视化报表。
    \item 阐明测试异步、并发与并行软件系统的特有挑战与调试应对策略。
    \item 为包含异步或并发代码的中等规模模块编制包含测试矩阵与竞争诱发数据的测试计划。
    \item 阐述对非代码软件工程工件（如用户手册、架构白皮书）进行系统化验证与确认的评审流程。
\end{itemize}

\subsection{代码重构与软件演进}
\textbf{【学时建议】}：2 领域核心学时。

软件系统的长期维护与演化是软件工程区别于单纯编程的核心特征。本单元聚焦于系统演进中的核心规律、代码重构技艺、向后兼容性契约与工业级版本控制规范。

\subsubsection{领域核心主题}
\begin{itemize}
    \item \textbf{海勒姆定律（Hyrum's Law / 隐式接口定律）}：
    \begin{quote}
    \emph{“当一个 API 拥有足够多的使用者时，系统在契约中承诺的内容已不再重要：你所展现出的所有可被观察的行为（即使未在文档中声明、甚至是偶然出现的内部未定义行为），都会被某些下游用户所依赖。”}
    \end{quote}
    深刻理解这一定律对于大型系统接口设计、内部行为修改以及变更风险评估具有决定性指导意义。
    \item \textbf{向后兼容性的本质探析}：
    \begin{itemize}
        \item 兼容性并非某单一组件的内在单体属性，而是一种双方契约关系的相对属性。
        \item 向后兼容性必须在“提供者+ 消费者”的具体业务场景中，或者基于提供方明确承诺并达成的兼容性模型下进行科学评估。
    \end{itemize}
    \item \textbf{代码重构工程实践}：
    \begin{itemize}
        \item 经典重构模式：重命名、函数内联、函数提炼、抽取类/接口等。
        \item 现代 IDE 自动化重构工具的安全应用与原理。
        \item 静态分析工具在自动检测代码坏味道、提出重构建议以及批量生成安全补丁中的实战应用。
    \end{itemize}
    \item \textbf{版本控制与演进规范}：
    \begin{itemize}
        \item 语义化版本控制：Major.Minor.Patch 的严格契约与破坏性变更声明。
        \item 基于主干的持续演进与短生命周期特性分支。
    \end{itemize}
\end{itemize}

\subsubsection{拓展选修主题}
\begin{itemize}
    \item \textbf{大规模软件重构}：
    \begin{itemize}
        \item 当重构改动规模过于庞大无法原子提交（跨千万行代码库），或无法在提供者与所有消费者之间原子同步变更时（跨多代码仓库、下游存在私有闭源消费者）的工业级渐进式重构技术。
        \item 双轨并存机制：同时保留新旧两套 API 使之平滑共存。
        \item 微步演进：将庞大重构拆解为系列微小的行为保持变换。
        \item 区分“可直接感知的 API 消费者”与“不可见/隐匿的外部消费者（Consumers I can't see）”带来的不同发布安全策略。
    \end{itemize}
\end{itemize}

\subsubsection{示范性学习成果与能力目标}
\begin{itemize}
    \item 准确识别特定接口的显式契约行为与所有隐式可观察行为，并依据海勒姆定律深刻评估破坏隐式行为可能带来的严重工程崩溃风险。
    \item 综合静态代码分析工具的告警报告与软件设计准则，敏锐定位代码库中急需重构劣质代码。
    \item 严谨研判某项代码修改是否属于真正的“向后兼容”，并能给出关于哪些具体调用场景受支持、哪些不受支持的清晰声明。
    \item 在对既有下游调用方产生零影响或极低波及影响的前提下，重构接口的内部实现以显著改善代码清晰度、内聚度与可维护性。
    \item 结合当前项目所采用的版本规范策略（如 SemVer），严密评估某项提议的系统变更是否具备足够的发布安全性。
    \item 制定周密的多阶段平滑迁移方案，在超大规模代码库中安全地修改、替换并最终废弃某核心 API 的旧有默认行为。
\end{itemize}

\subsection{软件可靠性工程}
\textbf{【学时建议】}：2 领域核心学时。

软件系统的可靠性关乎现代社会基础设施的安全平稳运转。本单元系统剖析可靠性的数学本质、全生命周期缺陷防御、系统容错设计、以及现代站点可靠性工程（SRE）的先进理念。

\subsubsection{领域核心主题}
\begin{itemize}
    \item \textbf{可靠性核心概念体系}：失效率概率、平均无故障时间（MTBF）、平均修复时间（MTTR）；故障、错误与失效的演化因果链路。
    \item \textbf{面向不同领域软件的可靠性需求定级}（交叉 SEP 知识领域：消费级应用 vs. 医疗/航空/金融级关键系统）。
    \item \textbf{软件失效的缺陷根源与防御闭环}：软件失效根本上由代码缺陷诱发，因此高可靠性的目标是追求生产环境缺陷最少化——首要在源头减少缺陷注入（通过专业培训、系统化架构规划、高质量编码规范），其次在交付前最大化消除注入的缺陷（严密测试、静态扫描、同侪评审）。
    \item \textbf{软件可靠性、系统级可靠性与失效行为特征}（交叉 系统基础/基于冗余实现系统可靠性）。
    \item \textbf{缺陷注入与消除全周期}：各阶段缺陷检出率度量与消除手段选型。
    \item \textbf{“错误预算”与“零缺陷”理念的深度思辨}：SRE 现代错误预算理念（允许合理可控的服务中断以换取极速创新迭代）与传统绝对零缺陷理念的对比，准确划定二者各自的适用产业领域。
\end{itemize}

\subsubsection{拓展选修主题}
\begin{itemize}
    \item \textbf{软件可靠性增长模型}：基于失效历史数据的可靠性预测曲线。
    \item \textbf{软件容错架构与容错模型}：
    \begin{itemize}
        \item 容错上下文的深刻工程差异：在飞行控制等关键系统中，必须竭尽全力避免系统崩溃，通过冗余算路自愈；而在数据流批处理系统中，在损坏的脏数据落盘之前让系统主动崩溃、快速失败反而是最具价值的安全容错机制。
    \end{itemize}
    \item \textbf{软件可靠性工程（SRE）工业级实践}：混沌工程、故障注入演练、轻量级模型检验。
    \item \textbf{强相关与相互独立失效域识别}及其对分布式系统整体可用性的连锁影响。
    \item \textbf{基于度量与可观测性的可靠性分析}：全链路遥测、指标监控与智能告警、实时 SLO/SLA 仪表盘、版本发布准入量化指标。
\end{itemize}

\subsubsection{示范性学习成果与能力目标}
\begin{itemize}
    \item 阐述如何为特定业务场景下的软件系统科学测定与协商其所必需的可靠性级别（如服务等级目标 SLO）。
    \item 深刻阐明在真实工程中试图实现极高可靠性指标（如“五个九”99.999\% 可用性）所面临的指数级成本与物理阻碍。
    \item 掌握在软件全生命周期的每个关键阶段（需求、架构、编码、测试、交付、运维）最小化缺陷注入与最大化缺陷清除的系统化工程方法。
    \item 展现综合运用多种可靠性预测模型对特定软件系统的可靠性演进趋势进行量化评估的能力。
    \item 识别并设计能够达成特定可靠性级别指标的软件容错体系结构。
    \item 熟练运用多种维度的冗余机制（时间冗余、空间冗余、数据冗余、逻辑异构冗余）实现软件容错。
    \item 在系统架构设计中敏锐识别单点故障 (SPOF)依赖并提出根除方案。
\end{itemize}

\subsection{形式化方法 (拓展选修)}
\textbf{【学时建议】}：本单元内容为拓展选修，高度依赖离散结构领域的基础数学知识（特别是 DS/基础逻辑与 DS/证明技术）。

形式化方法运用数理逻辑与形式语言对软件规约与行为进行严格的数学证明与系统化模型分析，是保障关键系统万无一失的高级工程手段。

\subsubsection{核心主题}
\begin{itemize}
    \item \textbf{接口形式化规约}：
    \begin{itemize}
        \item 前置条件、后置条件与状态不变量的形式化定义。
        \item 编写与分析形式化前后置条件的专用规范语言与证明工具系统。
    \end{itemize}
    \item \textbf{形式化方法高度适用的核心工程问题域}：
    \begin{itemize}
        \item 无锁并发编程、微观数据竞争隐患。
        \item 复杂异步与分布式共识系统（死锁 Deadlock、活锁 Livelock、饥饿、状态空间爆炸）。
    \end{itemize}
    \item \textbf{形式化分析与经验缺陷检测手段的多维对比}：与传统单元测试、集成测试、模糊测试（Fuzzing）在覆盖完备性与工程投入产出比上的核心权衡。
    \item \textbf{现代软件形式化建模与自动化分析工具}：
    \begin{itemize}
        \item 模型检验器（Model Checkers，如 TLA+、SPIN）。
        \item 模型发现器（Model Finders，如 Alloy）。
    \end{itemize}
\end{itemize}

\subsubsection{示范性学习成果与能力目标}
\begin{itemize}
    \item 阐释形式化规约与自动化分析技术在大型复杂安全攸关软件系统研发中的关键角色，系统对比其与经验测试在验证和确认层面的本质优势与局限。
    \item 将形式化规约与形式化验证工具切实应用于低复杂度的软件核心设计、并发算法或关键业务协议之中。
    \item 深入剖析在工业工程实践中引入形式化规范语言所能带来的极高可靠性收益，及其伴随的高昂学习门槛与建模维护成本。
\end{itemize}

\section{专业素养与思维倾向}

软件工程的学科特质决定了从业者不仅需要掌握扎实的计算机科学技术，更必须具备特定的职业素养、思维方式与心理倾向：

\begin{itemize}
    \item \textbf{协作精神}：软件工程正日益被公认为一项高强度的“团队运动”——合格的软件工程师必须能够与他人高效协同。\textbf{谦逊、尊重与信任}是构筑健康团队协作关系的基石，也是在软件工程领域取得长远成功的根本前提。
    \item \textbf{专业操守}：软件工程师构建的数字化系统深刻重塑着全球数十亿人的日常生活与经济活动。清醒意识到软件系统在现代社会中的重大影响力、恪守崇高的工程伦理准则，并在团队内部与外部日常跨部门交互中始终展现专业、体面与受人尊重的行为风范，是从业者的立身之本。
    \item \textbf{敏锐沟通}：在庞大复杂的现代软件工程项目中，没有任何单一工程师能够通晓系统的所有技术细节。成功的软件项目无不依赖工程师之间高频、清晰、坦诚与严谨的沟通，以实现高保真度的工程协调与敏捷推进。
    \item \textbf{严谨细致}：软件工程对细节的高度敏感性要求团队中的每一位成员都展现出高度一致的严谨态度。在该领域的卓著成就无一例外得益于一丝不苟的工程习惯——对系统规约的全面理解、对规范化流程的严格遵循、以及对投机取巧、偷工减料的本能抵制。
    \item \textbf{责任担当}：软件工程的高度协作性进一步放大了个人担当的核心价值。推诿扯皮、不兑现交付承诺、对阻塞性问题隐瞒不报，是导致团队摩擦升级、技术债务激增乃至软件项目彻底流产的经典祸根。
\end{itemize}

\section{数学基础要求}

\textbf{建议先修内容：}
\begin{itemize}
    \item \textbf{应用概率统计基础}：用于系统多轮性能测试数据对比、A/B 测试实验评估、用户调研问卷结果统计推断、以及性能基准测试中的测量不确定度与置信区间分析。
\end{itemize}

\section{交叉主题与共享概念}

\subsection{共享概念}
\begin{itemize}
    \item \textbf{职业素养与团队协作}：与社会、伦理与专业素养（SEP）知识领域深度共享。
    \item \textbf{数据建模}：与数据管理（DM）知识领域深度共享。
    \item \textbf{安全编程与防御性开发}：与计算机安全（SEC）知识领域深度共享。
\end{itemize}

\subsection{跨学科贯穿主题}
\begin{itemize}
    \item \textbf{工程伦理与社会责任}：深刻贯穿于需求分析、合规审查、错误预算制定与系统可靠性权衡的全生命周期中。
    \item \textbf{工业级编程实现}：作为一切软件架构设计落地与单元测试演练的代码级载体。
\end{itemize}

\section{课程体系规划与开课建议}

针对高校计算机科学专业高年级开设的《高级软件工程》专业核心课程，规范建议至少包含以下 30 学时的进阶教学模块：
\begin{itemize}
    \item 团队协作与沟通 —— 4 学时
    \item 工具与开发环境 —— 4 学时
    \item 产品需求工程 —— 2 学时
    \item 软件架构与系统设计 —— 5 学时
    \item 软件构造与工程实现 —— 4 学时
    \item 代码重构与软件演进 —— 2 学时
    \item 软件可靠性工程 —— 2 学时
    \item 融合来自 SEP 领域的“工程伦理与专业沟通”模块 —— 7 学时
\end{itemize}

\textbf{先修课程要求}：
\begin{itemize}
    \item \textbf{软件开发基础 (SDF)}：学生必须具备足够扎实的个人代码编写流利度，唯有如此，他们才能将心智精力从繁琐的基础语法调试中解放出来，全力聚焦于团队协作、规范流程、架构权衡与工程演进。
\end{itemize}

\textbf{核心能力目标陈述}：\\
成功修完本门课程的合格毕业生，应当能够熟练为团队同事开展高质量的同侪代码评审（尤其注重专业友善的沟通礼仪与团队心理安全感建设），能够独立编写结构完备的自动化单元测试套件，娴熟驾驭现代软件工程工具链（工业级 IDE、分布式版本控制系统、静态分析与覆盖率检查工具），并全面胜任软件工程团队对应届初级工程师所期望的一切核心专业工程交付活动。

\section{真实能力任务规范}

\begin{taskbox}[任务 1：研发中新增业务代码的自动化测试套件构建]
\textbf{【任务情境】}：在敏捷迭代中为正在开发的新核心业务功能编写自动化测试，确保新代码在当前及未来长期维护中始终保持功能正确。\\
\textbf{【能力陈述】}：剖析目标 API 的预期行为与设计契约，将其系统分解为涵盖典型输入、极端边界输入与非法异常输入的测试用例集，构建完整的自动化单元测试套件，确保该功能在后续长期代码演进中持续保持正确性与稳健性。\\
\textbf{【能力领域】}：软件 / 应用 \quad \textbf{【能力单元】}：需求分析 / 系统设计 / 编码实现 / 质量测试 / 技术文档\\
\textbf{【知识单元关联】}：软件开发基础/开发方法、软件工程/软件构造 \quad \textbf{【达标等级】}：\textbf{应用与开发}
\end{taskbox}

\begin{taskbox}[任务 2：为团队同事执行规范化同侪代码评审]
\textbf{【任务情境】}：团队成员提交了一项核心功能实现的代码拉取请求（PR），需要你在合并入主干前开展严格的技术把关。\\
\textbf{【能力陈述】}：秉持专业素养与同理心，以清晰、建设性且富有协作精神的方式与提交代码的同事沟通，就代码的逻辑正确性、测试完备性、潜在安全隐患、代码异味与架构可读性提供详尽切实的评审反馈。\\
\textbf{【能力领域】}：软件 / 应用 \quad \textbf{【能力单元】}：系统设计 / 技术文档 / 代码改进\\
\textbf{【知识单元关联】}：软件开发基础/开发方法、软件工程/团队协作、软件工程/软件验证与确认 \quad \textbf{【达标等级】}：\textbf{应用}
\end{taskbox}

\begin{taskbox}[任务 3：工业级 API 接口文档编制与契约形式化]
\textbf{【任务情境】}：为系统关键业务模块提取通用服务接口，供团队内外其他研发团队进行跨模块调用集成。\\
\textbf{【能力陈述】}：准确梳理目标组件的 API 表面接口，编写结构严谨、表意精准的对外技术文档，明确界定参数类型、前置条件、后置条件、系统状态副作用与异常错误码规范，确保接口被正确理解与安全调用。\\
\textbf{【能力领域】}：软件工程 \quad \textbf{【能力单元】}：需求分析 / 系统设计 / 编码实现 / 技术文档\\
\textbf{【知识单元关联】}：软件开发基础/基础编程概念、软件工程/软件构造、软件工程/软件设计 \quad \textbf{【达标等级】}：\textbf{应用}
\end{taskbox}

\begin{taskbox}[任务 4：遗留与第三方代码库的缺陷排查与调试修复]
\textbf{【任务情境】}：线上运行的遗留业务系统出现复杂异常崩溃，该模块由前人编写且缺少完整注释，需紧急定位并修复。\\
\textbf{【能力陈述】}：认识到在工业界“阅读理解代码的频率远高于编写全新代码”的现实规律，能够熟练通读并透彻理解他人编写的陌生代码，运用现代调试工具链与调用栈追踪技术进行严密审视排查，准确定位隐蔽 Bug 根源并安全重构修复。\\
\textbf{【能力领域】}：软件 / 应用 \quad \textbf{【能力单元】}：系统实现 / 评估诊断 / 软件维护 / 代码改进\\
\textbf{【知识单元关联】}：软件开发基础/开发方法、软件工程/软件构造 \quad \textbf{【达标等级】}：\textbf{应用}
\end{taskbox}

\begin{taskbox}[任务 5：针对研发、设计与调试场景合理进行现代工具选型]
\textbf{【任务情境】}：团队启动一项全新复杂分布式服务的研发，需要搭建兼顾研发效能、代码质量与协作透明度的工具链平台。\\
\textbf{【能力陈述】}：全面洞察现代主流软件工程工具族（静态分析器、动态插桩检测器、分布式版本控制、代码覆盖率分析器、自动化重构工具等）的原理与特性，针对具体的工程痛点与开发瓶颈，科学评估并选用最适宜的工具组合解决问题。\\
\textbf{【能力领域】}：软件 / 应用 \quad \textbf{【能力单元】}：系统实现 / 工具集成\\
\textbf{【知识单元关联】}：软件开发基础/开发方法、软件工程/工具与环境、软件工程/软件构造、软件工程/软件验证与确认 \quad \textbf{【达标等级】}：\textbf{解释}
\end{taskbox}

\begin{taskbox}[任务 6：生产环境核心 API 的平滑演进与安全重构]
\textbf{【任务情境】}：生产环境中已被海量下游系统广泛依赖的基础核心 API 需要调整底层架构并变更部分默认行为。\\
\textbf{【能力陈述】}：在充分考虑海勒姆定律的前提下，精准评估语法变更与语义行为变动的潜在破坏风险，制定周密的多阶段迁移与双轨平滑演化方案，运用现代化重构工具与语义版本机制安全落地重构，向利益相关者清晰传达迁移指引。\\
\textbf{【能力领域】}：软件 / 应用 \quad \textbf{【能力单元】}：系统设计 / 软件维护\\
\textbf{【知识单元关联】}：软件工程/代码重构与软件演进 \quad \textbf{【达标等级】}：\textbf{应用}
\end{taskbox}

\begin{taskbox}[任务 7：自动化代码工程格式化与规范统一约束]
\textbf{【任务情境】}：在大型团队多人协同开发中，解决不同成员编码格式不一引发的代码冲突与 Review 噪音。\\
\textbf{【能力陈述】}：熟练掌握行业通用代码规范与自动化代码格式化工具，在团队流水线与本地 Git Hooks 中落地统一格式约束，深刻践行工程一致性与整洁代码的核心价值。\\
\textbf{【能力领域】}：软件 / 应用 \quad \textbf{【能力单元】}：系统实现 / 软件维护\\
\textbf{【知识单元关联】}：软件开发基础/基础编程概念、软件工程/软件构造 \quad \textbf{【达标等级】}：\textbf{应用}
\end{taskbox}

\begin{taskbox}[任务 8：复杂业务组件的高内聚低耦合 API 架构设计]
\textbf{【任务情境】}：根据给定的复杂业务需求说明，设计一套对外交互的高内聚、低耦合工业级 API。\\
\textbf{【能力陈述】}：综合权衡多维质量属性（执行性能、扩展性、安全性、易用性），设计一套逻辑清晰、命名一致、具备高度可组合性的公共 API 接口集，并给出形式化数据不变量与契约保障。\\
\textbf{【能力领域】}：软件工程 \quad \textbf{【能力单元】}：需求分析 / 架构设计\\
\textbf{【知识单元关联】}：软件工程/软件设计 \quad \textbf{【达标等级】}：\textbf{开发}
\end{taskbox}

\begin{taskbox}[任务 9：关键系统架构设计决策与权衡技术文档撰写]
\textbf{【任务情境】}：面对技术选型存在重大分歧的架构设计，撰写一份正式的架构决策记录供全员评审。\\
\textbf{【能力陈述】}：深刻理解软件设计本质上是针对多种质量属性（运行性能、可用性、维护成本、代码可读性等）的权衡艺术，撰写严谨规范的技术文档，系统阐释重大设计决策背后的技术考量、放弃的备选方案与根本决策依据。\\
\textbf{【能力领域】}：软件 / 应用 \quad \textbf{【能力单元】}：需求分析 / 架构设计 / 技术文档\\
\textbf{【知识单元关联】}：软件工程/软件设计 \quad \textbf{【达标等级】}：\textbf{开发}
\end{taskbox}

\begin{taskbox}[任务 10：作为高效成员参与敏捷团队协作与全周期协同]
\textbf{【任务情境】}：作为初级研发工程师加入敏捷开发团队，参与长周期的产品迭代开发。\\
\textbf{【能力陈述】}：成为团队中高产且可靠的团队成员，展现优秀的团队沟通技巧与长期心理建设意识，主动分担任务，积极参与敏捷站会与复盘，在团队面临工期压力与摩擦时维持良性人际互动与高质量工程交付。\\
\textbf{【能力领域】}：工程应用 \quad \textbf{【能力单元】}：技术文档 / 团队社会议题调适 / 过程改进\\
\textbf{【知识单元关联】}：软件工程/团队协作 \quad \textbf{【达标等级】}：\textbf{应用}
\end{taskbox}

\begin{taskbox}[任务 11：面向开放式综合工程项目的跨角色协同设计与敏捷交付]
\textbf{【任务情境】}：跨专业组队承接来自非技术外部客户提出的模糊开放式需求项目，需在限定预算与工期内完成交付。\\
\textbf{【能力陈述】}：与团队同侪及非技术客户紧密协同，梳理并澄清模糊需求，在严格的时间表和有限资源约束下，完成端到端的系统架构设计、编码实现、持续集成与部署验证，成功达成既定项目目标。\\
\textbf{【能力领域】}：软件 / 应用 \quad \textbf{【能力单元】}：需求工程 / 架构设计 / 技术文档\\
\textbf{【知识单元关联】}：软件工程/团队协作、软件工程/软件设计、软件工程/软件构造 \quad \textbf{【达标等级】}：\textbf{应用}
\end{taskbox}

\section{工作委员会成员名单}

\textbf{委员会主席}：
\begin{itemize}
    \item Titus Winters, 谷歌公司 —— \emph{注：《Google软件工程：源自世界顶级代码库的实践经验》资深领衔作者}
\end{itemize}

\textbf{委员会委员}：
\begin{itemize}
    \item Brett A. Becker, 都柏林大学
    \item Adam Vartanian, Cord 公司
    \item Bryce Adelstein Lelbach, 英伟达公司
    \item Patrick Servello, CIWRO 气象雷达研究所
    \item Pankaj Jalote, 德里因陀罗普拉斯塔信息技术学院
    \item Christian Servin, 埃尔帕索社区学院
\end{itemize}

\textbf{学术与工业界特别贡献者}：
\begin{itemize}
    \item Hyrum Wright, 谷歌公司 —— \emph{注：软件工程著名“海勒姆定律（Hyrum's Law）”命名提出者}
    \item Olivier Giroux, 苹果公司
    \item Gennadiy Civil, 谷歌公司
\end{itemize}

\clearpage
\onecolumn
\section{附录：核心主题与技能要求等级对照表}

\begin{longtable}{p{0.20\textwidth} p{0.46\textwidth} c c c}
\toprule
\textbf{知识单元} & \textbf{核心知识点与具体主题} & \textbf{认知等级} & \textbf{分类} & \textbf{学时} \\
\midrule
\endfirsthead
\toprule
\textbf{知识单元} & \textbf{核心知识点与具体主题} & \textbf{认知等级} & \textbf{分类} & \textbf{学时} \\
\midrule
\endhead
\bottomrule
\endfoot
\multirow{2}{*}{\makecell[l]{软件工程/团队协作\\与沟通}} 
& 有效沟通、团队冲突根因与消解策略、结对编程、同侪代码评审、角色分工、敏捷流程、团队多元化 & 评估 & 计算机核心 & 2 \\
\cmidrule{2-5}
& 与外部利益相关者（管理层/客户/用户）对接、物理与远程虚拟团队风险剖析 & 解释 & 领域核心 & 2 \\
\midrule
\multirow{2}{*}{\makecell[l]{软件工程/工具与\\开发环境}} 
& 软件配置管理与版本控制、可复现构建、基于主干开发的分支策略、合并与变基 & 评估 & 计算机核心 & 1 \\
\cmidrule{2-5}
& 发布管理、静态/动态测试工具、构建系统、CI 持续集成流水线、SemVer 依赖管理、现代 IDE 高级能力 & 解释 & 领域核心 & 3 \\
\midrule
\makecell[l]{软件工程/产品需求工程} 
& 用例与用户故事、需求一致性/完备性/可行性、需求诱导采集、非功能质量属性、需求风险管理与规约形式化 & 应用 & 领域核心 & 2 \\
\midrule
\multirow{2}{*}{\makecell[l]{软件工程/软件架构\\与系统设计}} 
& 系统设计核心原则（抽象、关注点分离、信息隐藏、高内聚低耦合）、经典架构模式（C/S、分层、MVC）、大/小规模编程辨析、代码异味 & 解释 & 计算机核心 & 1 \\
\cmidrule{2-5}
& API 设计原则（一致性、可组合性、契约设计、错误报告）、数据与状态不变量、结构与行为模型、数据架构设计、需求可追溯性 & 应用 & 领域核心 & 4 \\
\midrule
\multirow{2}{*}{\makecell[l]{软件工程/软件构造\\与工程实现}} 
& 实用小型测试（单元测试、测试驱动开发 TDD 心理学机制）、规范化接口文档与微妙实现注释 & 应用 & 计算机核心 & 1 \\
\cmidrule{2-5}
& 工业级编码风格与规范指南、高质量最佳实践（防御性编程、安全编码、异常容错）、断点调试、第三方成熟库集成 & 应用 & 领域核心 & 3 \\
\midrule
\multirow{2}{*}{\makecell[l]{软件工程/软件验证\\与确认 (V\&V)}} 
& 验证与确认概念辨析、测试战略意义、测试多维目标与测试层级、DAMP vs DRY 测试代码哲学 & 解释 & 计算机核心 & 1 \\
\cmidrule{2-5}
& 测试计划与用例生成、测试矩阵与覆盖率度量、TDD 实践、面向对象测试/Mocking/DI、代码评审、提交前后测试权衡、缺陷追踪、前沿领域验证挑战 & 解释 & 领域核心 & 3 \\
\midrule
\makecell[l]{软件工程/代码重构\\与软件演进} 
& 海勒姆定律（隐式接口定律）、向后兼容性契约本质、经典代码重构模式与 IDE 自动化重构、语义化版本控制与主干演进 & 解释 & 领域核心 & 2 \\
\midrule
\makecell[l]{软件工程/软件可靠性\\工程} 
& 可靠性数学概念、领域可靠性需求定级、缺陷注入与清除全周期、系统级冗余容错行为、现代 SRE“错误预算”与“零缺陷”对比思辨 & 解释 & 领域核心 & 2 \\
\end{longtable}

\chapter{专用平台开发}

\section{领域导言与演进背景}

\begin{preamblebox}[知识领域概述与变革]
专用平台开发 (SPD)知识领域旨在培养学生针对特定行业场景、特定硬件体系与严苛业务需求，构建专用软件平台与行业解决方案的核心能力。在现代数字化社会中，专用平台已全面渗透至医疗健康、金融科技、交通运输、智能制造、具身智能及元宇宙交互等诸多关键垂直领域。专用平台的工程构建涵盖了从系统需求分析、体系架构设计、异构软硬件协同开发，到全生命周期自动化部署与高可用维护的完整链路。

\textbf{自 CS2013 以来的核心变革：}
\begin{enumerate}
    \item \textbf{新兴计算前沿与垂直领域的爆发式增长}：数据科学与大数据分析依赖多端跨平台感知数据采集；网络空间安全要求在软硬件底层强化数据隔离与信道防御；人工智能与机器学习催生了具身机器人与无人机的自主协同控制；交互式计算平台迎来了大语言模型（LLMs，如 OpenAI ChatGPT、Codex 与 GitHub Copilot）的革命性融合，极大地拓展了代码生成、交互式探索与创意计算的技术边界。
    \item \textbf{产业界工程师核心能力诉求剧增}：行业对精通移动端、现代 Web、机器人、嵌入式硬件及游戏引擎的平台级全栈工程师呈现迫切需求。包括：在碎片化设备与操作系统间打造体验高度一致的移动应用、面向头戴式显示器（HMD）与混合现实（MR）评估无障碍交互体验、基于云原生与现代全栈语言构建高安全 Web 业务底盘，以及利用微控制器（Raspberry Pi、Arduino）和开源机器人操作系统（ROS）实现软硬件互联控制。
    \item \textbf{知识领域全面升级与重命名}：由 CS2013 中的“基于平台的开发 (PBD)”正式重命名为“专用平台开发 (SPD)”，精准反映了平台开发从通用中间件走向垂直专用软硬件生态的时代转型。
    \item \textbf{计算机科学核心学时（计算机核心）实现历史性突破}：在 CS2013 中本领域核心学时为 0；而在 CS2023 中，计算机科学核心学时大幅提升至 \textbf{8$\sim$9 学时}。大纲将 Web 与移动端划分为通识核心基础与专用平台进阶模块，为全球各高校依据自身办学定位灵活配置特色专业方向（如游戏开发、智能硬件、云计算全栈）提供了强大的弹性支撑。
    \item \textbf{吸纳前沿知识单元并强化社会伦理责任}：全面设立“机器人平台”、“嵌入式平台”、“游戏平台”、“交互式计算与大模型平台”，并将“社会、伦理与专业素养（SEP）”深度融入硬件适老化、无障碍设计及端侧数据隐私合规的全流程中。
\end{enumerate}
\end{preamblebox}

\section{核心学时分配建议}

本领域在课程体系中采用“核心基石 + 弹性专精（领域核心/选修）”的双轨学时机制。Web 与移动开发的核心基础构成了全体计算机专业毕业生的通用必修基底（计算机核心 累计 8$\sim$9 学时），而机器人、嵌入式、游戏引擎与交互计算平台则作为各高校计算机专业深化方向的核心选修（领域核心 累计 15 学时）：

\begin{table}[H]
\centering
\setlength{\tabcolsep}{3.5pt}\fontsize{8pt}{9.6pt}\selectfont\begin{tabularx}{\columnwidth}{Y c c}
\toprule
\textbf{知识单元} & \textbf{计算机核心} & \textbf{领域核心} \\
\midrule
跨平台共性技术与关注点 & 2 & -- \\
Web 平台全栈开发与云原生架构 & 3 & -- \\
移动计算平台与端侧智能应用 & 3 & -- \\
机器人系统与具身智能开发平台 & -- & 3 \\
嵌入式系统与低功耗物联网平台 & -- & 3 \\
游戏引擎架构与实时交互互动平台 & -- & 4 \\
交互式计算平台与大模型应用 & -- & 2 \\
平台社会责任、伦理合规与专业素养 & -- & 3 \\
\midrule
\textbf{合计} & \textbf{8 学时} & \textbf{15 学时} \\
\bottomrule
\end{tabularx}
\end{table}

\noindent\textbf{【学时弹性说明】}：部分办学方案中可将 Web 基础（2 学时）与移动端基础（1 学时）抽离为先导基础（计算机核心 共 3 学时），再结合高校特色在各专用平台（*）中选配 6 学时深度实践，总计覆盖 9 门类核心计算机学时。

\section{知识单元与详细大纲}

\subsection{跨平台共性技术与架构关注点}
\textbf{【学时建议】}：2 计算机核心学时（计算机核心）。

本单元旨在系统梳理面向多样化专用计算平台进行软件构建时的通用方法论与核心架构约束。引导学生深入理解不同平台（移动端、Web、游戏主机、工业自动化、嵌入式设备）的软硬件异构特性，掌握专用应用编程接口（API）与跨平台编程语言的高效应用方法。

\textbf{【核心知识点与专题大纲】}
\begin{itemize}
    \item \textbf{现代专用平台生态全景概览}：Web 浏览器平台、移动终端平台、主机与掌机游戏平台、工业自动化与机器人系统平台。
    \begin{itemize}
        \item \textbf{输入与感知控制机制}：多点触控、空间惯性传感器（IMU）、计算机视觉跟踪（Inside-out/Outside-in）、触觉力反馈。
        \item \textbf{严苛平台资源约束分析}：有限算力与热耗限制；受限板载内存与持久化存储容量；非对称与弱网无线通信带宽；系统高可用性、容错韧性与工业合规安全性。
        \item \textbf{输出与受控执行机构}：高刷新率与动态色域显示引擎、物理执行机构、高精度力矩伺服电机与震动反馈装置。
    \end{itemize}
    \item \textbf{专用平台 API 驱动式编程 vs. 传统通用软件构造}：硬件抽象层（HAL）与操作系统内核沙箱的交互契约；平台级生命周期事件驱动模型 vs. 传统主控制台执行流。
    \item \textbf{跨平台现代开发语言图谱与选型}：Kotlin、Swift、C\#、C++、Java、JavaScript、TypeScript 以及 HTML5 Web 标准。
    \item \textbf{平台约束环境下的工程构建}：受限交叉编译工具链、真机远程联合调试、内存基线分析与安装包体积削减策略。
    \item \textbf{专用领域语言与新兴平台的快速掌握范式}：语言运行时特性的迁移学习、平台官方架构指南研读法、SDK 核心范式的逆向抽象。
\end{itemize}

\textbf{【示范性学习成果与能力指标】}
\begin{itemize}
    \item 列举并深度剖析移动端开发在内存占用、并发能耗与交互时延方面的核心约束。
    \item 阐明 Web 开发经典三层架构（表示层、业务逻辑层、持久化数据层）在现代富客户端形态下的演化机制。
    \item 解释在 Web 无状态通信协议（HTTP）下利用 Cookie、Session、Token（JWT）及分布式缓存维护用户会话状态的底层原理。
    \item 辨析动态脚本语言与强类型编译型语言（C++/Rust）在运行效率、内存安全及跨平台部署维度的核心权衡。
\end{itemize}

\subsection{平台全栈开发与云原生架构}
\textbf{【学时建议】}：3 计算机核心学时（计算机核心）。

本单元聚焦于现代 Web 生态的体系结构、全栈编程技术、安全防御体系及云原生服务治理。涵盖富客户端现代框架、SaaS 架构、无障碍标准、端到端数据管理以及弹性伸缩机制。

\textbf{【核心知识点与专题大纲】}
\begin{itemize}
    \item \textbf{现代 Web 全栈开发语言体系}：HTML5 语义化标准、CSS3 弹性布局与网格排版系统、现代 ECMAScript（JavaScript）、TypeScript 静态类型系统及后端开发语言（Node.js、PHP、Go、Python）。
    \item \textbf{Web 应用平台与现代前端元框架}：组件化虚拟 DOM 体系、服务端渲染（SSR）、静态站点生成（SSG）、增量静态再生（ISR）以及全栈集成框架。
    \item \textbf{软件即服务设计范式}：多租户架构隔离（数据隔离与模式隔离）、订阅式计量与权限网关。
    \item \textbf{Web 国际标准与无障碍体验}：W3C 文档对象模型（DOM）规范、无障碍网页内容指南（WCAG 2.1）、语义化 ARIA 标签规范。
    \item \textbf{Web 平台全栈安全与隐私纵深防御}：跨站脚本攻击（XSS）、跨站请求伪造（CSRF）、SQL 注入防御；内容安全策略（CSP）、HTTP 严格传输安全（HSTS）、CORS 跨域资源共享策略。
    \item \textbf{Web 业务系统需求分析与架构权衡}：用户高并发峰值预测、首屏加载时间（FPT）、累计布局偏移（CLS）与全链路核心体验指标。
    \item \textbf{现代云端算力基础设施与服务治理}：主流公有云（AWS、Azure、阿里云等）计算生态。
    \begin{itemize}
        \item \textbf{云端托管与容器化}：虚拟机、Docker 容器化与 Kubernetes 编排集群。
        \item \textbf{系统弹性与伸缩}：自动伸缩组、多可用区负载均衡（ALB/SLB）、CDN 边缘智能缓存。
        \item \textbf{云服务容量规划与成本工程}：按需计费、预留实例分析及 FinOps 成本优化模型。
    \end{itemize}
    \item \textbf{全栈数据生命周期治理}：
    \begin{itemize}
        \item \textbf{数据驻留与属地合规}：跨国数据合规路径、跨境数据流动加密与主权审计。
        \item \textbf{数据完整性与合规销毁}：数据多副本强一致性容灾机制、用户注销后的彻底可信擦除。
    \end{itemize}
    \item \textbf{现代软件架构范式演进}：
    \begin{itemize}
        \item \textbf{单体架构 vs. 微服务架构}：服务解耦、领域驱动设计（DDD）、分布式事务（Saga 模式）。
        \item \textbf{微前端架构}：独立构建与部署、模块联邦、运行时沙箱隔离。
        \item \textbf{事件驱动架构（EDA） vs. RESTful API 架构}：异步消息队列高吞吐解耦 vs. 同步轻量契约通信的优劣权衡。
        \item \textbf{无服务器计算}：函数即服务、冷启动优化、事件源触发与按需弹性计费。
    \end{itemize}
    \item \textbf{持久化存储方案综合选型}：
    \begin{itemize}
        \item \textbf{关系型数据库（RDBMS）}：ACID 事务、索引 B+ 树优化、读写分离与主从复制。
        \item \textbf{NoSQL 与 NewSQL 体系}：文档型数据库、键值数据库（Redis 缓存防穿透/击穿/雪崩）、分布式列存数据库。
    \end{itemize}
\end{itemize}

\textbf{【示范性学习成果与能力指标】}
\begin{itemize}
    \item 遵循微服务或无服务器架构规范，独立设计并实现一套具备弹性伸缩能力的 Web 业务原型系统。
    \item 全面剖析 Web 运行平台带来的机遇与瓶颈，重点阐述容器托管、动态扩缩容与跨域跨端安全通信的实现路径。
    \item 对比分析 Web 敏捷开发范式与传统单机通用软件开发的异同点（在发布频率、状态持久化与客户端异构性层面）。
    \item 阐明 SaaS 架构在持续交付、多租户计费及弹性容量治理上相较于传统买断式桌面软件的核心优势。
    \item 针对现有线上大型 Web 系统执行 WCAG 2.1 无障碍合规性与 W3C 标准审计，提出优化整改技术方案。
\end{itemize}

\subsection{移动计算平台与端侧智能应用}
\textbf{【学时建议】}：3 计算机核心学时（计算机核心）。

本单元专注于移动计算终端（iOS/Android/鸿蒙）的开发技术、底层运行约束及新型交互能力。深入探讨移动端本地端侧计算、专用 NPU 神经网络加速、硬件沙箱安全隔离以及高并发异步流式响应。

\textbf{【核心知识点与专题大纲】}
\begin{itemize}
    \item \textbf{现代移动工程研发技术栈}：
    \begin{itemize}
        \item \textbf{主流移动开发语言}：Swift（现代函数式与强类型特性）、Kotlin（协程高并发与空安全）、Dart。
        \item \textbf{官方集成开发环境}：Xcode、Android Studio 及其真机性能调优工具链。
        \item \textbf{原生开发 vs. 跨平台技术框架}：Native 深度控制力 vs. 跨平台框架（Flutter 统一渲染管道、React Native 桥接架构）的开发效率与运行开销权衡。
        \item \textbf{移动端主流架构模式}：MVC、MVP、MVVM（数据绑定）与现代单向数据流架构（MVI）。
    \end{itemize}
    \item \textbf{移动平台硬性软硬件约束与优化}：
    \begin{itemize}
        \item \textbf{用户界面与触控交互设计}：响应式屏幕适配、沉浸式手势导航、深色模式与动态字体无障碍支持。
        \item \textbf{移动 App 与 Web 页面用户体验差异}：离线可用性、平滑转场动画（60/120 FPS 无卡顿刷新）与微交互设计。
        \item \textbf{端侧系统安全防护}：应用进程沙箱隔离、运行时动态权限最小化申请、安全飞地与 Keyring/Keystore 加密密钥存储。
        \item \textbf{功耗与计算性能权衡}：电量消耗模型、后台唤醒频率对齐、动态降频应对发热降频。
    \end{itemize}
    \item \textbf{移动数据互联与轻量接口设计}：
    \begin{itemize}
        \item \textbf{移动专属 API 端点优化}：针对弱网（高时延、高丢包）的批量请求合并、增量数据同步、BFF（Backend For Frontend）架构。
        \item \textbf{本地缓存与离线数据同步}：本地轻量 SQLite/Room/CoreData 存储与断网重连自动补偿机制。
    \end{itemize}
    \item \textbf{移动计算特有使能技术}：
    \begin{itemize}
        \item \textbf{空间位置感知计算}：高精度 GNSS、基站三角定位、Wi-Fi 室内指纹与地理围栏低功耗监测。
        \item \textbf{多模态传感器融合计算}：三轴陀螺仪、加速度计、气压计、健康体征传感器（光学心率/血氧）数据采集与异常过滤。
        \item \textbf{实时通讯与系统级通道}：移动电话通信、VoIP 低时延流媒体、系统级推送机制。
        \item \textbf{移动端增强现实（AR）}：基于 ARKit/ARCore 的平面检测、六自由度（6DoF）空间定位与虚拟物体虚实遮挡融合。
    \end{itemize}
    \item \textbf{工程规范、机型适配与自动化测试}：
    \begin{itemize}
        \item 碎片化屏幕与多版本 OS 兼容性云测、自动化 UI 测试、端侧灰度发布与崩溃埋点上报。
    \end{itemize}
    \item \textbf{移动端异步计算与响应式编程}：
    \begin{itemize}
        \item 异步非阻塞模型与同步阻塞模型的底层区别；主线程（UI 线程）避免掉帧（ANR 规避）。
        \item 现代异步控制范式：回调机制、Promise、async/await 以及结构化并发（Kotlin 协程 Coroutines / Swift Concurrency Tasks）。
        \item 异步错误捕获传播链与网络超时异常降级策略。
        \item 异步并发代码测试难题：竞态条件排查、时间依赖测试与内存悬空引用检测。
    \end{itemize}
\end{itemize}

\textbf{【示范性学习成果与能力指标】}
\begin{itemize}
    \item 实现一款基于位置感知（LBS）的移动应用，能够利用后台地理围栏与云端数据 API 实现动态信息实时推送。
    \item 构建一款传感器驱动的高性能移动应用，高效采集陀螺仪与体征数据并以低功耗模式同步至云端服务器。
    \item 结合原生通信能力与 WebRTC/WebSocket 协议，编写一款支持即时文本与低时延音视频通话的移动应用。
    \item 辨析原生开发与基于 Flutter/React Native 的跨平台技术方案，在首屏冷启动速度、动画渲染性能及团队维护成本上的核心利弊。
\end{itemize}

\subsection{机器人系统与具身智能开发平台}
\textbf{【学时建议】}：3 领域核心学时（领域核心）。

本单元探讨针对现有实体与仿真机器人平台进行系统级软件构建与具身智能部署的技术。内容涵盖机器人软硬件架构、专用开发环境与库（如机器人操作系统 ROS/ROS2）、空间运动学与动力学、自主导航与路径规划，以及物理世界人机协作安全规范。

\textbf{【核心知识点与专题大纲】}
\begin{itemize}
    \item \textbf{现代机器人平台与设备类型}：轮式移动机器人（AGV/AMR）、工业多轴机械臂、足式四足/双足人形机器人、无人飞行器（UAV）。
    \item \textbf{感知、嵌入式板载计算与受控执行器}：激光雷达、RGB-D 深度相机、超声波测距传感器；无刷伺服电机、高精度减速机（谐波/RV）与末端电磁抓夹。
    \item \textbf{机器人专用开发语言、框架与生态库}：C++ / Python 系统级接口、URDF/Xacro 机器人物理模型描述文件、MoveIt 运动规划框架、OpenCV 视觉与 PCL 点云处理库。
    \item \textbf{机器人平台物理约束与工程考量}：高精度实时时序约束、非完整约束、动力矩饱和限制与电池续航极限。
    \item \textbf{虚实融合交互与仿真引擎互联}：基于 Gazebo、Webots 或 NVIDIA Isaac Sim 建立高保真物理仿真模型并实现零成本迁移。
    \item \textbf{机器人学核心基础算法}：
    \begin{itemize}
        \item \textbf{软件体系架构}：机器人操作系统 ROS/ROS2 核心通信机制（节点 Node、话题 Topic 异步发布/订阅、服务 Service 同步请求、动作 Action 目标打断反馈机制）。
        \item \textbf{正向运动学}：空间位姿描述、齐次变换矩阵、DH 参数法。
        \item \textbf{逆运动学}：几何解析法与数值迭代法（雅可比矩阵求逆与阻尼最小二乘法）。
        \item \textbf{刚体动力学}：拉格朗日力学与牛顿-欧拉迭代动力学方程、转动惯量张量计算。
        \item \textbf{自主导航与路径规划}：激光/视觉 SLAM 同步定位与建图、全局规划（$A^*$ 算法、Dijkstra、RRT*）与局部避障控制（TEB/DWA 算法）。
        \item \textbf{末端操控与抓取规划}：空间轨迹插补（样条曲线）、末端六自由度抓取姿态估计与力控自适应抓取。
        \item \textbf{人机协作物理安全规范}：急停联锁保护机制、碰撞力矩检测顺应控制与 ISO 10218/TS 15066 安全标准。
    \end{itemize}
\end{itemize}

\textbf{【示范性学习成果与能力指标】}
\begin{itemize}
    \item 在指定机器人软硬件平台（如 ROS 连接的实体移动底盘或高保真物理仿真器）上独立设计并实现一套自主导航任务流。
    \item 基于微控制器（如 Arduino/树莓派）组装一套双轮差速寻迹机器人，编写底层嵌入式控制代码使其自主穿越未知迷宫。
    \item 结合计算图架构，深入阐释 ROS 发布/订阅机制与通用多线程操作系统进程间通信（IPC）在时间确定性与资源开销上的异同。
    \item 针对给定的工业物流搬运任务，使用 ROS2 架构设计一套完整的软硬件协同系统，明确定义各功能节点及其 Topic 与 Action 接口规范。
\end{itemize}

\subsection{嵌入式系统与低功耗物联网平台}
\textbf{【学时建议】}：3 领域核心学时（领域核心）。

本单元专注于资源严重受限的嵌入式微处理器、片上系统以及广泛泛在计算的物联网边缘节点系统开发。涵盖硬实时确定性、外设总线驱动、超低功耗休眠管理与软硬件协同测试验证。（并发、死锁与底层调度等通用议题可参照 PDC 与 OS 领域）。

\textbf{【核心知识点与专题大纲】}
\begin{itemize}
    \item \textbf{嵌入式系统的本质特征与分类}：
    \begin{itemize}
        \item 硬实时、软实时与非实时系统的时延确定性界定。
        \item 严苛资源受限环境：有限 RAM 堆栈空间、无虚拟内存 MMU 支持、严格硬性截止时间。
    \end{itemize}
    \item \textbf{系统功能安全分析与可靠性设计}：故障安全与故障容错架构；危险与可操作性分析（HAZOP）、故障树分析（FTA）及安全编码规范（如 MISRA C/C++）。
    \item \textbf{物理感知与执行外设硬件接口}：通用输入输出（GPIO）、模拟-数字转换（ADC/DAC）、脉宽调制（PWM）；常见工业总线协议（UART、SPI、I2C、CAN 总线及 Modbus）。
    \item \textbf{嵌入式核心开发范式}：裸机前后台轮询架构、实时操作系统内核（FreeRTOS、RT-Thread、Zephyr RTOS 任务调度）；交叉编译体系与链接器脚本内存分区映射。
    \item \textbf{实时资源调度与并发管理}：速率单调调度（RMS）、最早截止时间优先（EDF）；优先级反转风险与优先级天花板/优先级继承协议（PIP/PCP）；中断服务例程（ISR）短小快设计原则与延迟处理机制。
    \item \textbf{嵌入式系统验证与底层调试技术}：硬件在环仿真（HIL）、静态代码分析、JTAG/SWD 硬件在线断点调试、逻辑分析仪与示波器总线时序抓取。
    \item \textbf{低功耗物联网系统工程设计}：微安级深度休眠唤醒模式、看门狗定时器自动故障自愈、边缘轻量智能计算（TinyML）推理。
\end{itemize}

\textbf{【示范性学习成果与能力指标】}
\begin{itemize}
    \item 针对特定硬件微控制器平台（如 STM32、ESP32），独立设计并开发一套具备传感器采样与低功耗休眠唤醒功能的物联网节点应用。
    \item 深入对比嵌入式实时操作系统（RTOS）与通用分时操作系统（如 Linux）在任务上下文切换、中断响应延迟及内存管理策略上的根本差异。
    \item 编写底层外设驱动代码，通过 SPI/I2C 总线实现对惯性测量单元或温湿度传感器的数据采集，并驱动 OLED 显示。
    \item 熟练运用硬件调试器与逻辑分析仪排查嵌入式系统通信总线偶发数据丢失、时钟毛刺或栈溢出导致的 HardFault 死锁故障。
    \item 结合算力、能耗与物料成本（BOM），评估微控制器（MCU）、数字信号处理器（DSP）与现场可编程门阵列（FPGA）在特定嵌入式应用中的选型合理性。
\end{itemize}

\subsection{游戏引擎架构与实时交互互动平台}
\textbf{【学时建议】}：4 领域核心学时（领域核心）。

本单元专注于在严苛硬件算力约束下构建高性能实时交互图形渲染与互动游戏系统的核心软件工程原理。全面覆盖现代 GPU 硬件渲染管线、物理仿真、着色器编程、主流游戏引擎架构、跨平台移植与人机工效无障碍体验。

\textbf{【核心知识点与专题大纲】}
\begin{itemize}
    \item \textbf{游戏硬件与软件计算平台的演进谱系}：
    \begin{itemize}
        \item \textbf{平台演进历程}：从早期的 Brown Box、家用主机时代走向现代元宇宙、云游戏与数字孪生；CPU 与 GPU 协同异构计算的飞速跃迁；平台融合与移动化趋势。
        \item \textbf{典型游戏平台形态}：高性能个人电脑（PC）、家用主机（PlayStation、Xbox、Switch）、便携掌机、街机、智能电视大屏、移动智能手机、平板电脑、VR/AR 空间计算一体机、沉浸式互动环幕（CAVE 洞穴系统）、Web 浏览器、云端流式渲染平台。
        \item \textbf{不同游戏平台的底层硬件特性与严苛约束}：内存带宽与拓扑结构（DDR vs. GDDR 高带宽高延迟统一内存架构）；GPU/CPU 时钟频率与核心集群调度；专用协处理器（物理加速卡 PPU、向量协处理器 SIMD）；总线带宽（PCIe 通道数）；网络吞吐与往返时延（RTT）；最大函数调用栈深度限制；发热温控降频机制与整机功耗预算；字节序（大端/小端 Endianness）差异。
        \item \textbf{典型交互传感器、控制器与触觉执行机构}：键鼠、游戏手柄摇杆与线性扳机、空间体感可穿戴设备、触控交互表面；计算机视觉空间定位跟踪（Inside-out 与 Outside-in 跟踪）；震动与触觉力反馈系统。
    \end{itemize}
    \item \textbf{游戏平台的社会、法律、伦理与无障碍考量}：
    \begin{itemize}
        \item \textbf{可用性工程与人机工效}：玩家交互直觉、认知负荷控制、游戏人机交互启发式评估方法。
        \item \textbf{无障碍支持与包容性设计}：色盲色弱显示增强滤镜、字幕与屏幕朗读辅助、全键位无障碍重映射适配、相关法规遵从性评估（如 CVAA 法案）。
        \item \textbf{绿色计算与软硬件可持续性}：低碳渲染模式设计、平台功耗限制、设备生命周期环保回收与反计划报废设计。
    \end{itemize}
    \item \textbf{实时物理仿真与现代图形渲染系统}：
    \begin{itemize}
        \item \textbf{现代 CPU 与 GPU 体系架构}：费林分类法（Flynn's Taxonomy，SIMD/SIMT 并行架构）；GPU 图形计算阵列、显存控制器、PCIe 显卡总线、显示扫描输出引擎。
        \item \textbf{渲染流水线与硬件光栅化架构}：移动端分块延迟渲染 (TBDR)与桌面端立即渲染模式 (IMR)的内存带宽开销权衡。
        \item \textbf{核心算法、空间数据结构与数学模型}：确定性游戏主循环；空间加速层次树划分（八叉树 Octree、包围盒层次 BVH、KD-Tree）；视锥体裁剪与细节层次（LOD）动态调度；快速碰撞检测（分离轴定理 SAT、GJK 算法）与刚体碰撞响应；NPC 行为树与分层有限状态机（HFSM）；程序化内容生成（PCG）。
        \item \textbf{多媒体虚拟世界资产计算与表达}：空间多声道三维音频、精灵图集、3D 网格模型与物理基础材质（PBR 材质双向反射分布函数 BRDF）、骨骼动画与逆运动学插值。
    \end{itemize}
    \item \textbf{游戏工程研发技术栈与主流工具链}：
    \begin{itemize}
        \item \textbf{游戏编程语言体系}：C++（高性能引擎核心）、C\#（Gameplay 逻辑）、Lua/Python（轻量热更新脚本）。
        \item \textbf{现代着色器语言与可视化着色器}：HLSL、GLSL、MSL 及基于节点的 ShaderGraph 可视化渲染材质系统。
        \item \textbf{底层现代图形 API}：DirectX 12、Vulkan、Metal、WebGL / WebGPU 底层显存显卡控制。
        \item \textbf{主流商业与开源游戏引擎}：Unreal Engine、Unity、Godot、CryEngine 等现代引擎架构；版本控制工具（处理海量二进制资产的 Git LFS、Perforce）；GPU 性能剖析分析器（RenderDoc、PIX、NVIDIA Nsight）。
        \item \textbf{关键开发工程方法}：快速原型验证、可执行构建流水线编译、真机帧率与显存性能调优、跨平台架构抽象移植、网络联机架构（状态同步、帧同步与预测回滚 Rollback 算法）。
    \end{itemize}
    \item \textbf{现代游戏设计工程体系与泛在应用}：
    \begin{itemize}
        \item \textbf{设计语言与模型}：MDA 理论框架（机制 Mechanics、动态 Dynamics、美感 Aesthetics）、心流体验模型。
        \item \textbf{以玩家为中心的设计思维}：双钻设计模型、迭代原型开发、玩家社区参与。
        \item \textbf{核心流派架构特征}：角色扮演（RPG）、第一人称射击（FPS）、即时战略（RTS）、多人在线战术竞技（MOBA）。
        \item \textbf{受众细分与玩家行为学}：巴特尔玩家分类模型（成就者、探索者、社交者、杀手）、游戏动机量表。
        \item \textbf{游戏技术的跨行业溢出与泛在应用}：严肃游戏、工业数字孪生、虚拟制片、全沉浸式模拟器与医学仿真培训。
    \end{itemize}
\end{itemize}

\textbf{【示范性学习成果与能力指标】}
\begin{itemize}
    \item 阐述现代 GPGPU 架构的核心特征，说明 SIMT 线程分支发散对游戏渲染效率的负面影响。
    \item 梳理图形光栅化渲染流水线从顶点处理、图元装配、光栅化、片元着色到模板/深度测试与输出合并的完整阶段。
    \item 运用现代着色器语言，按照光照反射模型要求手写实现一个自定义材质片元着色器。
    \item 利用现代图形 API 或工业级引擎，独立实现一个具备物理碰撞检测、空间四叉树分区与相机漫游的 3D 虚拟场景。
    \item 运用 GPU Profiler 性能分析工具，对一款游戏原型进行每帧 CPU 耗时、Draw Call 数量、显存带宽及过度绘制的专项诊断与优化。
    \item 针对跨平台游戏发布场景，评估一套能兼顾桌面 PC 与便携移动终端的高可移植性软件架构设计方案。
\end{itemize}

\subsection{交互式计算平台与大模型应用}
\textbf{【学时建议】}：2 领域核心学时（领域核心）。

本单元探讨支持人类进行探索性数据分析、实时可视化交互、创意编程，以及结合现代大语言模型（LLMs，如 OpenAI ChatGPT、Codex 与 GitHub Copilot）构建新型人机协同智能应用的前沿软件平台与范式。

\textbf{【核心知识点与专题大纲】}
\begin{itemize}
    \item \textbf{交互式数据探索分析计算平台}：
    \begin{itemize}
        \item 文档型可执行笔记本环境：Jupyter Notebook、Google Colab、Observable HQ、RStudio。
        \item 云端大数据交互查询平台：云 SQL 交互引擎（Google BigQuery、Snowflake）与分布式内存分析平台。
    \end{itemize}
    \item \textbf{交互式数据可视化生态与低时延渲染反馈}：
    \begin{itemize}
        \item 数据驱动的交互图表展示体系（D3.js、ECharts、Plotly）。
        \item 极低响应延迟的交互反馈闭环：硬件加速图形渲染工具、交互式参数调节画布。
    \end{itemize}
    \item \textbf{创意编程与生成式艺术交互框架}：
    \begin{itemize}
        \item 实时交互式媒体与艺术计算框架（Processing、p5.js、TouchDesigner、Max/MSP）。
        \item 现场即兴现场编码（Live Coding，如 Sonic Pi、TidalCycles 现场音乐生成）。
        \item 人机交互协同生成（人工智能/ML 交互式配对创作体验）。
    \end{itemize}
    \item \textbf{大语言模型（LLM）开发生态与协同应用}：
    \begin{itemize}
        \item 主流大模型编程助手与 API 集成：OpenAI ChatGPT/Codex API、GitHub Copilot、Anthropic Claude、本地开源推理平台（Ollama、vLLM）。
        \item \textbf{辅助数理与工科计算研究}：音频信号时域与傅里叶频域分析交互、数据统计建模与算法原型敏捷验证。
        \item \textbf{赋能人文社科与数字创意探索}：基于数据的交互式新闻叙事、文学脚本生成与音乐和声理论探索。
    \end{itemize}
\end{itemize}

\textbf{【示范性学习成果与能力指标】}
\begin{itemize}
    \item 在 Jupyter Notebook 或 Google Colab 平台中交互式加载并分析大规模复杂数据集，完成数据清洗、多维统计与深度交互图表绘制。
    \item 基于实时现场编程环境（如 Sonic Pi），编写一段能接收用户参数输入并实时调制音频合成器的音乐演奏程序。
    \item 利用主流大模型 API 构建一个具备特定上下文提示工程与函数调用能力的智能助手程序。
    \item 对比分析 AI 辅助平台自动生成的代码与人类专家编写的代码在时间复杂度、边缘边界条件防御以及长期可维护性上的本质差异。
    \item 阐明交互式探索计算与严谨工业级生产数据管道在开发范式、错误处理与吞吐量设计上的核心分歧。
\end{itemize}

\subsection{平台社会责任、伦理合规与专业素养}
\textbf{【学时建议】}：3 领域核心学时（领域核心）。

本单元立足专用平台开发的工程前沿，系统探讨移动端传感数据合规、物联网端侧监控、沉浸式元宇宙伦理、人机协作机器人生理心理安全以及界面多元化包容性（DEI）等重大社会与伦理议题。

\textbf{【核心知识点与专题大纲】}
\begin{itemize}
    \item \textbf{沉浸式扩展现实与增强技术的社会伦理冲击}：VR/AR 头戴显示器沉浸式体验对人类认知模式与现实脱节风险；元宇宙虚拟经济中的知识产权与未成年人心理防护。
    \item \textbf{具身机器人与自动化系统的人机伦理困境}：自动驾驶汽车事故责任归属与“电车难题”抉择；工业协作机器人的人身物理安全防护边界；服务型机器人拟人化带来的隐私边界模糊。
    \item \textbf{多元化、公平性与包容性（DEI）在界面设计中的工程落地}：适老化界面改造、消除对少数群体与残障群体的算法偏见；多语言全球化文化适配与色彩无歧视呈现。
    \item \textbf{端侧设备数据隐私法规与监控防范}：移动定位信息滥用规避、智能家居物联网麦克风与摄像头端侧隐私保护；GDPR 及相关个人信息保护法规在边缘计算设备上的合规实施。
\end{itemize}

\section{专业素养与思维倾向}

专用平台开发面向高度碎片化、快速迭代且异构的软硬件生态，要求计算机专业学生具备以下核心专业素养：
\begin{itemize}
    \item \textbf{自主学习与前沿技术追踪力}：新型硬件架构、专用操作系统与领域专用语言（DSL）层出不穷，工程师必须具备迅速阅读技术规范、脱离既有思维定势快速掌握全新平台开发工具链的元学习能力。
    \item \textbf{严苛约束下的工程创造力}：在电池功耗、内存空间、计算散热及通信带宽受到严厉约束的非常规平台上，能够创造性地设计高效算法与精简数据结构，达成极致性能优化。
    \item \textbf{异构环境敏捷适应性}：能够灵活游刃于完全不同的软硬件生态之间，面对移动端碎机型、嵌入式总线偶发噪声及云端动态漂移环境迅速调整工程实现方案。
    \item \textbf{严谨求实的软硬件联调素养 (Meticulousness in Debugging \& Testing)}：在面对跨时钟域总线、异步网络波动以及多线程并发竞态条件时，保持极端细致的逻辑严密性，善用硬件示波器、底层 Profiler 及系统日志深挖隐蔽 Bug。
\end{itemize}

\section{数学与跨学科基础要求}

专用平台软件开发横跨信号处理、机械运动控制、三维渲染管线与空间几何计算，对以下数学领域提出了明确要求：
\begin{itemize}
    \item \textbf{微积分}：用于机器人动力学中质点加速度与力矩的微分方程推导、物理引擎中位移积分计算以及连续物理信号采集滤波。
    \item \textbf{线性代数}：用于三维空间坐标系旋转平移（齐次坐标矩阵、四元数姿态插值变换）、机器人正逆运动学雅可比矩阵计算，以及现代 GPU 顶点几何渲染流水线变换。
    \item \textbf{概率论与数理统计}：用于物联网传感器噪声过滤（卡尔曼滤波算法）、弱网通信抖动与丢包率概率模型构建、大规模交互用户行为特征分布分析。
    \item \textbf{离散数学与结构}：用于自主导航算法中的图论寻路（$A^*$、Dijkstra）、机器人状态机与游戏行为树模型构建、事件驱动系统优先队列设计。
\end{itemize}

\section{交叉主题与跨学科共享概念}

专用平台开发是多门经典计算机学科的集大成者，与以下知识领域高度交叉：
\begin{itemize}
    \item \textbf{人工智能 (AI)}：端侧轻量化神经网络模型推理（TinyML）、大模型 API 赋能平台交互、机器人环境感知与自主决策规划。
    \item \textbf{计算机图形学 (GIT)}：三维网格几何表达、实时光栅化渲染、自定义片元/顶点着色器编写、虚拟现实沉浸式渲染。
    \item \textbf{人机交互 (HCI)}：触觉与多模态传感交互、以人为本的界面工效学设计、WCAG 无障碍可访问性开发标准。
    \item \textbf{系统建模与仿真}：刚体动力学仿真、高保真虚拟物理环境构建、数字孪生技术。
    \item \textbf{编程语言理论 (FPL)}：特定平台专用领域语言（DSL）、着色器着色语言机制、元编程与跨平台运行时的解释与 JIT 编译。
    \item \textbf{软件工程 (SE)}：模块解耦设计、持续集成与自动化发布流水线（CI/CD）、跨平台代码复用度管理、微服务治理。
    \item \textbf{需求工程}：深入挖掘垂直领域用户特殊场景痛点，在设备物理制约与业务功能体验之间达成最佳工程平衡。
\end{itemize}

\section{课程体系规划与开课建议}

各高校可依据自身师资力量与专业特色方向，将本领域内容规划组合为以下典型专业模块课程：
\begin{enumerate}
    \item \textbf{移动与现代 Web 应用全栈开发 (Mobile and Modern Web App Development, 48 学时)}：
    全面覆盖 Web 前端组件化与云端微服务基础（24h）、移动原生与跨平台架构实战（18h）、数据安全通信与端侧轻量数据库（6h）。侧重从零搭建一套具备云端弹性后端支撑的高性能多端移动商业应用。
    \item \textbf{机器人软件工程与 ROS 系统实战 (Robotics Software Engineering and ROS, 45 学时)}：
    深入讲解 Linux 实时环境与 C++ 开发（9h）、ROS2 节点/话题/服务架构与 Gazebo 仿真（15h）、机器人运动学与 SLAM 自主导航建图（15h）、实体差速底盘或机械臂硬件抓取实验（6h）。
    \item \textbf{嵌入式系统与物联网全栈开发 (Embedded Systems and IoT Development, 45 学时)}：
    涵盖微控制器体系与 C 语言裸机外设编程（12h）、FreeRTOS 任务调度与多任务并发同步（15h）、低功耗无线物联网协议接入（BLE/Wi-Fi/MQTT, 12h）、硬件在环故障诊断实训（6h）。
    \item \textbf{游戏引擎架构与实时交互图形系统 (Game Engine Architecture and Interactive Media, 48$\sim$60 学时)}：
    涵盖现代 GPU 架构与实时着色器编程（18h）、游戏主循环与空间碰撞加速结构（15h）、现代商业引擎核心架构与网络同步对战机制（15h）。
    \item \textbf{交互式数据计算与生成式 AI 编程 (Interactive Data Computing and LLM Programming, 32 学时)}：
    涵盖交互式 Notebook 数据探索与流式可视化（12h）、大模型 API 提示词工程与函数调用架构（12h）、基于 LLM 的智能助手与创意代码生成实战（8h）。
\end{enumerate}

\section{真实能力任务规范}

\begin{taskbox}[任务 1：原生开发与跨平台技术选型决策评估]
\textbf{【任务情境】}：某初创科技企业拟推出一款融合高帧率动画、复杂底层传感器调用及多端快速部署特征的数字化健康服务产品，技术团队在“全原生开发（iOS Swift + Android Kotlin）”与“跨平台框架”之间存在选型分歧。\\
\textbf{【能力陈述】}：综合技术架构与产品设计维度，权衡系统运行性能、包体积、原生硬件接口访问能力、发布周期及人力维护成本，撰写客观详实的架构技术选型论证报告。\\
\textbf{【能力领域】}：应用开发 \quad \textbf{【能力单元】}：系统评估\\
\textbf{【关联知识单元】}：软件工程/工程工具与开发环境、专用平台/专用平台共性特征、专用平台/移动平台\\
\textbf{【达标等级】}：\textbf{解释与评估} \quad \textbf{【学时层级】}：计算机核心
\end{taskbox}

\begin{taskbox}[任务 2：高一致性多端自适应移动应用设计与实现]
\textbf{【任务情境】}：企业核心业务移动客户端需要在各厂商碎片化 Android 设备、多样化分辨率折叠屏以及不同版本 iOS 系统上保持高度一致的用户体验与流畅动画。\\
\textbf{【能力陈述】}：熟练运用现代移动开发语言与 UI 响应式布局体系，优化应用触控响应与内存占用，通过真机云测与自动化 UI 测试确保在跨机型、异构操作系统上稳定达标 60 FPS 刷新率。\\
\textbf{【能力领域】}：软件与应用 \quad \textbf{【能力单元】}：设计、开发与测试\\
\textbf{【关联知识单元】}：人机交互/理解用户、编程语言/面向对象编程、软件开发基础/开发方法与实践、软件工程/软件构造、软件工程/测试与质量保证、专用平台/移动平台\\
\textbf{【达标等级】}：\textbf{开发} \quad \textbf{【学时层级】}：计算机核心
\end{taskbox}

\begin{taskbox}[任务 3：基于引擎将策划蓝图原型重构为高性能跨平台核心系统]
\textbf{【任务情境】}：游戏策划与交互设计师在商业游戏引擎中利用可视化脚本（如 Unreal Blueprints）快速搭建了玩法原型，但在高密度复杂场景下出现严重的 CPU 主线程掉帧，且需移植发布至 PC、游戏主机与移动终端。\\
\textbf{【能力陈述】}：熟练剖析游戏主循环与底层渲染瓶颈，将可视化原型逻辑重构为基于 C++ 的高性能模块，运用空间数据结构与多线程机制优化碰撞计算，确保核心玩法的跨平台无缝兼容与高流畅度。\\
\textbf{【能力领域】}：软件与系统应用 \quad \textbf{【能力单元】}：系统研发、性能评测与调优\\
\textbf{【关联知识单元】}：软件工程/软件架构与设计、软件工程/软件构造、算法/算法策略与范式、系统基础/系统性能、专用平台/专用平台共性特征、专用平台/游戏平台\\
\textbf{【达标等级】}：\textbf{开发} \quad \textbf{【学时层级】}：领域核心
\end{taskbox}

\begin{taskbox}[任务 4：沉浸式头显（HMD）与混合现实新型外设交互可用性与无障碍评估]
\textbf{【任务情境】}：某团队研发了一款配合空间计算头戴式显示器（HMD）使用的手部骨骼力反馈外设，拟面向公共教育与游戏互动场景进行广泛推广。\\
\textbf{【能力陈述】}：制定严谨的用户人机工效测试方案，针对不同手型、行动障碍以及弱视用户群体开展可用性与无障碍合规性评估，针对前庭动晕症（晕动症）与力反馈延迟提出工程改进建议。\\
\textbf{【能力领域】}：系统与人机体验 \quad \textbf{【能力单元】}：需求分析、用户测试与社会适应性\\
\textbf{【关联知识单元】}：人机交互/理解用户、人机交互/无障碍与包容性设计、图形交互/沉浸式技术、图形交互/物理计算、社会伦理/多样性、平等与包容 (DEI)、专用平台/游戏平台\\
\textbf{【达标等级】}：\textbf{应用与评估} \quad \textbf{【学时层级】}：领域核心
\end{taskbox}

\begin{taskbox}[任务 5：基于演进商业需求构建高安全与高弹性 Web 业务系统]
\textbf{【任务情境】}：某企业数字化交易平台随着业务扩张面临跨国数据合规审计与日益频繁的 Web 漏洞攻击风险，亟需重构底层业务架构。\\
\textbf{【能力陈述】}：熟练运用全栈现代编程语言与云服务组件，实施防范 XSS、CSRF 与 SQL 注入的纵深安全防御，构建多可用区弹性微服务底座，确保系统在流量洪峰下兼顾安全与高可用。\\
\textbf{【能力领域】}：应用工程 \quad \textbf{【能力单元】}：系统设计与安全评估\\
\textbf{【关联知识单元】}：体系结构/性能与能效、网络通信/网络空间安全、操作系统/安全防护与保护机制、系统基础/系统安全、软件工程/软件架构与设计、专用平台/专用平台共性特征、专用平台/Web 平台、社会伦理/隐私保护\\
\textbf{【达标等级】}：\textbf{开发} \quad \textbf{【学时层级】}：计算机核心
\end{taskbox}

\begin{taskbox}[任务 6：面向移动计算生态设计与交付高性能移动业务 API]
\textbf{【任务情境】}：公司需要重构移动应用数据层交互通道，以往臃肿的桌面端通用 REST 接口在移动弱网环境下导致严重的客户端请求堆叠与电池电量过耗。\\
\textbf{【能力陈述】}：遵循移动平台最佳工程规范，设计轻量级、低冗余、支持批量聚合与断点续传的移动专属 API 接口规范，编写高可靠移动端数据网络层 SDK 并实现完善的本地离线缓存支持。\\
\textbf{【能力领域】}：应用软件工程 \quad \textbf{【能力单元】}：接口研发与性能评测\\
\textbf{【关联知识单元】}：编程语言/硬件接口、操作系统/文件系统 API 与底层实现、软件工程/软件架构与设计、软件工程/工程工具与开发环境、专用平台/专用平台共性特征、专用平台/移动平台\\
\textbf{【达标等级】}：\textbf{开发} \quad \textbf{【学时层级】}：计算机核心
\end{taskbox}

\begin{taskbox}[任务 7：自主智能机器人软硬件交互系统端到端设计与实现]
\textbf{【任务情境】}：针对自动化工厂物流中周转物料的自动寻引与精确装卸需求，团队需要开发一款具备自主避障定位与末端抓取的机器人系统。\\
\textbf{【能力陈述】}：准确梳理工业场景核心任务需求，设计基于 ROS2 架构的软件节点通信拓扑，打通激光雷达建图、底盘逆运动学解算与末端机械臂动作控制，实现机器人全自动化运行。\\
\textbf{【能力领域】}：系统应用工程 \quad \textbf{【能力单元】}：系统设计、集成与评估\\
\textbf{【关联知识单元】}：算法/算法策略与范式、软件工程/工程工具与开发环境、软件工程/软件架构与设计、专用平台/专用平台共性特征、专用平台/机器人平台\\
\textbf{【达标等级】}：\textbf{解释与应用} \quad \textbf{【学时层级】}：领域核心
\end{taskbox}

\begin{taskbox}[任务 8：面向严苛硬件指标的高可靠嵌入式系统架构设计与协议研发]
\textbf{【任务情境】}：医疗设备研发团队需开发一款便携式多参数生理监护仪，要求系统在毫瓦级低功耗下持续采样心电信号，且软硬件系统绝不允许发生阻塞死锁。\\
\textbf{【能力陈述】}：依据硬件芯片物理规格书，在轻量实时操作系统（RTOS）上设计多任务确定性调度体系，实现 SPI 传感器总线高频通信与看门狗自动故障恢复机制，编写完整的板级支持包（BSP）。\\
\textbf{【能力领域】}：系统应用工程 \quad \textbf{【能力单元】}：系统架构研发、软硬件部署与集成\\
\textbf{【关联知识单元】}：编程语言/硬件接口、软件工程/软件架构与设计、软件工程/工程工具与开发环境、系统基础/系统性能、专用平台/专用平台共性特征、专用平台/嵌入式平台\\
\textbf{【达标等级】}：\textbf{开发} \quad \textbf{【学时层级】}：领域核心
\end{taskbox}

\begin{taskbox}[任务 9：大规模分布式 Web 站点与线上资产持续运维与性能优化]
\textbf{【任务情境】}：某跨国电商企业在“大促”活动前夕，线上多套 Web 站点与核心交易网关面临瞬时百万级并发写入，底层数据库与反向代理服务器偶发超时告警。\\
\textbf{【能力陈述】}：深入分析系统网络吞吐、单跳链路与 NoSQL 分布式存储集群瓶颈，部署 CDN 静态资产边缘缓存、微前端动态路由分发与读写分离集群，建立自动化监控与灾备熔断机制。\\
\textbf{【能力领域】}：系统运维与工程应用 \quad \textbf{【能力单元】}：系统集成、监控与性能调优\\
\textbf{【关联知识单元】}：数据管理/NoSQL 系统与非关系型存储、软件工程/工程工具与开发环境、软件工程/软件架构与设计、网络通信/单跳通信与局域网、操作系统/文件系统 API 与底层实现、专用平台/专用平台共性特征、专用平台/Web 平台\\
\textbf{【达标等级】}：\textbf{应用与调优} \quad \textbf{【学时层级】}：计算机核心
\end{taskbox}

\begin{taskbox}[任务 10：跨职能团队协同交付企业级高可用全栈平台解决方案]
\textbf{【任务情境】}：在现代数字化敏捷研发流程中，涉及前端工程师、后端微服务架构师、移动客户端工程师、UI/UX 设计师以及产品运营经理等多角色的紧密协作。\\
\textbf{【能力陈述】}：展现良好的团队协作与技术沟通能力，综合梳理前后端交互协议、数据实体约束与客户端交互边界，制定标准化 API 契约与持续集成（CI）流水线，高质量按期交付全栈产品。\\
\textbf{【能力领域】}：工程管理与软件系统 \quad \textbf{【能力单元】}：团队协作、部署与交付\\
\textbf{【关联知识单元】}：软件工程/项目管理与度量、软件工程/软件架构与设计、软件工程/团队协作、专用平台/专用平台共性特征、专用平台/Web 平台、专用平台/移动平台\\
\textbf{【达标等级】}：\textbf{解释与协作} \quad \textbf{【学时层级】}：领域通用
\end{taskbox}

\section{指导委员会与主要贡献者}

\begin{itemize}
    \item \textbf{工作组主席}：
    \begin{itemize}
        \item Christian Servin (El Paso Community College, 美国)
    \end{itemize}
    \item \textbf{委员会核心委员}：
    \begin{itemize}
        \item Sherif G. Aly (The American University in Cairo, 埃及)
        \item Yoonsik Cheon (The University of Texas at El Paso, 美国)
        \item Eric Eaton (University of Pennsylvania, 美国)
        \item Claudia L. Guevara (Jochen Schweizer mydays Holding GmbH, 德国)
        \item Larry Heimann (Carnegie Mellon University, 美国)
        \item Amruth N. Kumar (Ramapo College of New Jersey, 美国)
        \item R. Tyler Pirtle (Google, 美国)
        \item Michael Scott (Falmouth University, 英国)
    \end{itemize}
    \item \textbf{主要贡献学者与行业专家}：
    \begin{itemize}
        \item Orlando Gordillo (NASA, 美国)
        \item Sean R. Piotrowski (Rider University, 美国)
        \item Mark O'Neil (Blackboard Inc., 美国)
        \item John DiGennaro (Qwickly, 美国)
        \item Rory K. Summerley (Falmouth University, 英国)
    \end{itemize}
\end{itemize}

\begin{strip}
\section{附录：核心主题与技能要求等级对照表}

\begingroup
\zihao{-5}\renewcommand{\arraystretch}{0.90}\setlength{\tabcolsep}{3pt}
\noindent\begin{tabularx}{\textwidth}{p{0.17\textwidth} Y c c c}
\toprule
\textbf{知识单元} & \textbf{核心知识点与具体主题} & \textbf{认知等级} & \textbf{分类} & \textbf{学时} \\
\midrule
\multirow{4}{*}{\makecell[l]{专用平台/跨平台\\共性技术}} 
& 专用平台生态全景（Web、移动端、机器人、嵌入式、游戏平台） & 解释 & 计算机核心 & 0.5 \\
& 专用平台 API 编程契约 vs. 传统桌面通用软件构建 & 应用 & 计算机核心 & 0.5 \\
& 跨平台开发语言图谱（Kotlin、Swift、C\#、C++、TypeScript 等） & 应用 & 计算机核心 & 0.5 \\
& 严苛资源受限（算力、内存、通信、功耗、可用性）下的工程构建 & 应用 & 计算机核心 & 0.5 \\
\midrule
\multirow{4}{*}{\makecell[l]{专用平台/Web 平台\\开发体系}} 
& 现代 Web 编程语言标准与前端组件化元框架 & 应用 & 计算机核心 & 1 \\
& SaaS 多租户架构与 W3C/WCAG 无障碍合规性规范 & 解释 & 计算机核心 & 0.5 \\
& Web 纵深安全防御机制（XSS、CSRF、CSP 与跨域通信） & 评估 & 计算机核心 & 0.5 \\
& 云原生弹性伸缩、无服务器与关系/NoSQL 存储选型 & 应用 & 计算机核心 & 1 \\
\midrule
\multirow{4}{*}{\makecell[l]{专用平台/移动计算\\平台开发}} 
& 移动端原生开发与跨平台开发框架架构选型对比 & 评估 & 计算机核心 & 0.5 \\
& 移动平台硬性软硬件约束（触控设计、功耗与性能权衡、沙箱安全） & 解释 & 计算机核心 & 0.5 \\
& 移动端特有使能技术（LBS 位置感知、多传感器融合、AR 空间计算） & 应用 & 计算机核心 & 1 \\
& 现代异步非阻塞编程范式（协程、async/await）与弱网容错测试 & 应用 & 计算机核心 & 1 \\
\midrule
\multirow{3}{*}{\makecell[l]{专用平台/机器人\\开发平台}} 
& 机器人平台类型、板载嵌入式计算、传感器与伺服驱动执行器 & 解释 & 领域核心 & 1 \\
& 机器人操作系统 ROS/ROS2 体系架构与软硬件虚实仿真互联 & 应用 & 领域核心 & 1 \\
& 正逆运动学、刚体动力学基础、SLAM 室内建图与自主路径规划 & 应用 & 领域核心 & 1 \\
\midrule
\multirow{3}{*}{\makecell[l]{专用平台/嵌入式与\\物联网平台}} 
& 嵌入式微控制器特性、硬实时/软实时时延确定性与功能安全分析 & 解释 & 领域核心 & 1 \\
& 物理感知外设硬件接口驱动（GPIO、ADC/DAC、SPI、I2C、CAN） & 应用 & 领域核心 & 1 \\
& 实时操作系统 RTOS 任务调度、低功耗休眠管理与硬件在环调试 & 应用 & 领域核心 & 1 \\
\midrule
\multirow{4}{*}{\makecell[l]{专用平台/游戏引擎\\互动平台}} 
& 游戏硬件演进、显示与外设跟踪接口、无障碍与绿色计算规范 & 解释 & 领域核心 & 1 \\
& GPU 现代并行架构、TBDR 与立即渲染流水线对比、着色器编程 & 应用 & 领域核心 & 1 \\
& 实时仿真算法（游戏主循环、八叉树/BVH 空间划分、刚体碰撞检测） & 应用 & 领域核心 & 1 \\
& 商业游戏引擎架构、跨平台移植调优与多端网络同步预测回滚 & 应用 & 领域核心 & 1 \\
\midrule
\multirow{2}{*}{\makecell[l]{专用平台/交互计算\\与大模型平台}} 
& 交互式数据分析探索平台与低时延动态可视化 & 应用 & 领域核心 & 1 \\
& 大语言模型（LLMs）开发生态、API 集成与人机协同编程实践 & 应用 & 领域核心 & 1 \\
\midrule
\makecell[l]{专用平台/社会伦理\\与专业素养} 
& 沉浸式技术现实冲击、机器人伦理安全、适老化 DEI 与端侧隐私合规 & 评估 & 领域核心 & 3 \\
\bottomrule
\end{tabularx}
\endgroup
\end{strip}

\chapter{系统基础}

\section{领域导言与演进背景}

\begin{preamblebox}[知识领域概述与变革]
计算机系统是由底层硬件与上层系统软件共同构建的综合计算基础设施，是所有现代应用软件运行的物理与逻辑载体。在当今数字化社会中，计算机系统已成为支撑人类生产与生活不可或缺的中枢基石。因此，系统学习计算机系统全栈知识、熟练掌握系统设计与工程构建技能，并深刻洞察其底层运行原理与核心设计哲学，是培养计算机专业拔尖人才的核心要务。

在计算机科学专业课程体系中，系统知识的传授通常跨越操作系统、并行与分布式系统、计算机网络、计算机体系结构与组织以及软件工程等多门核心骨干课程。正如其名称所示，“系统基础 (SF)”知识领域旨在提炼这些系统类核心课程共有且通用的基础概念与思维模式，以高度凝练、前后贯通的统一视角建立宏观全景，为学生后续深入学习各专用领域的特定机制与策略奠定坚实的系统思维基石。

\textbf{自 CS2013 以来的核心变革：}
\begin{enumerate}
    \item \textbf{增设两大全新核心单元}：正式增设“系统安全”与“系统设计”两个独立的知识单元，将现代安全防御与宏观架构设计准则确立为必修基石。
    \item \textbf{重组并创立“系统性能”单元}：废弃原 CS2013 中的“临近性”与“虚拟化与隔离”单元，将其核心精髓整合提升为全新的“系统性能”单元。
    \item \textbf{重组并创立“性能评估”单元}：废弃原 CS2013 中的“评估”与“定量评估”单元，全面重构为现代化的“性能评估”单元。
    \item \textbf{计算范式演进为“计算机系统概述”}：将原“计算范式”单元调整更名为“计算机系统概述”，精简陈旧内容，并吸收并入了原“跨层通信”单元的核心主题。
    \item \textbf{更名并深化“系统基本概念”}：将原“状态与状态转换”更名为“系统基本概念”，系统性补充了有限状态机（FSM）与时序系统建模等理论内容。
    \item \textbf{优化调整并行性主题}：精简了原“并行性”单元中的部分内容，重点保留并强化了简易应用级并行处理（如请求级并发、多线程多服务器模型及流水线）。
    \item \textbf{废弃独立跨层通信单元}：将原“跨层通信”相关主题有机并入“计算机系统概述”单元。
    \item \textbf{废弃原双重评估单元}：将原“评估”与“定量评估”合并提炼至“性能评估”单元。
    \item \textbf{整合临近性与虚拟化主题}：将原“临近性”与“虚拟化与隔离”的议题无缝融入“系统性能”单元。
    \item \textbf{废弃独立并行单元并归并基础概念}：取消独立的“并行性”单元，将其基础部分融入“系统基本概念”单元。
    \item \textbf{重命名并升格“系统可靠性”}：将原“基于冗余的可靠性”重命名并拓展为具有全生命周期视角的“系统可靠性”单元。
\end{enumerate}
\end{preamblebox}

\section{核心学时分配建议}

本知识领域在“向日葵模型”下的核心学时推荐如下表所示（计算机核心共 18 学时，领域核心共 9 学时，总计 27 学时）：

\begin{table}[H]
\centering
\setlength{\tabcolsep}{3.5pt}\fontsize{8pt}{9.6pt}\selectfont\begin{tabularx}{\columnwidth}{Y c c}
\toprule
\textbf{知识单元} & \textbf{计算机核心} & \textbf{领域核心} \\
\midrule
计算机系统概述 & 3 & -- \\
系统基本概念 & 4 & -- \\
资源分配与调度 & 1 & 2 \\
系统性能 & 2 & 2 \\
性能评估 & 2 & 2 \\
系统可靠性 & 2 & 1 \\
系统安全 & 2 & 1 \\
系统设计 & 2 & 1 \\
\midrule
\textbf{合计} & \textbf{18 学时} & \textbf{9 学时} \\
\bottomrule
\end{tabularx}
\end{table}

\section{知识单元与详细大纲}

\subsection{计算机系统概述}
\textbf{【学时建议】}：3 计算机核心学时。

本单元建立对现代计算系统分层架构与软硬件协同设计的宏观全局认知，阐释从基础物理逻辑到高层应用程序接口的逐级抽象演进。

\subsubsection{核心主题}
\begin{itemize}
    \item \textbf{计算机基本构建块与组件}：逻辑门电路、触发器、寄存器、总线与内部互连结构；现代处理器的经典划分（数据通路 Datapath + 控制逻辑 Control + 存储体系 Memory）。
    \item \textbf{作为计算范式的底层硬件}：基本数字逻辑单元；布尔代数逻辑表达式、卡诺图化简、积之和式（与或式，Sum-of-Products）。
    \item \textbf{程序设计抽象与接口规范}：软件抽象层级划分、应用程序编程接口（API）、底层系统库与动态链接库调用机制。
    \item \textbf{应用服务与操作系统服务的边界辨析}：用户态与内核态特权级分离、内核系统服务职责边界、远程过程调用（RPC）核心思想。
    \item \textbf{应用层与操作系统交互机制}：系统调用时序、陷阱指令、上下文保存与恢复、跨特权级数据拷贝开销。
    \item \textbf{流水线技术基础原理}：指令与任务级流水线核心概念、重叠执行阶段划分与吞吐提升。
    \item \textbf{系统扩展性基本概念}：强扩展（Strong Scaling，算力扩充加速单任务求解）与弱扩展（Weak Scaling，资源扩充支持更大问题规模）的本质差异。
\end{itemize}

\subsubsection{示范性学习成果与能力目标}
\begin{enumerate}
    \item \textbf{[解释]} 详细阐明计算机的基本硬件构建模块及其在计算机体系结构历史演进过程中的关键作用。
    \item \textbf{[应用/开发]} 运用基本逻辑门与触发器等硬件设计基础构建块，设计并实现可解决基础计算任务的组合与时序逻辑电路（例如全加器、计数器）。
    \item \textbf{[评估]} 熟练运用原理图捕获、逻辑综合与数字仿真工具链，对数字逻辑电路设计的逻辑功能与时序特性进行验证与性能评估。
    \item \textbf{[解释]} 深刻阐释现代计算系统如何基于“关注点分离”设计哲学构建层级体系，通过严谨规范的接口协议向上层屏蔽底层实现的复杂细节。
    \item \textbf{[解释]} 深入阐明底层硬件、操作系统内核、虚拟机运行时与应用层软件在计算本质上均是提供额外解释与加工能力的多层递进处理实体。
    \item \textbf{[解释]} 阐明在复杂分层体系结构中，底层故障检测、跨层异常状态向上汇报以及分级差错恢复与容错降级的运行机制。
    \item \textbf{[开发]} 综合运用分层设计方法、差错检测与自愈恢复技术，编写实现一个能够跨层反馈并优雅处理错误状态的典型软件系统（如可靠的 TCP 客户端/服务器通信系统）。
    \item \textbf{[应用/评估]} 熟练运用现代程序跟踪分析、单步调试及断点分析工具，在具有深层调用栈的分层系统中精准定位并修复复杂软件缺陷。
    \item \textbf{[解释/评估]} 透彻理解强扩展与弱扩展的本质机理，即问题求解性能如何受制于问题规模与计算资源配比，并能结合现实云计算或超级计算工程实例进行深度阐述。
\end{enumerate}

\subsection{系统基本概念}
\textbf{【学时建议】}：4 计算机核心学时。

本单元提炼计算系统中最本质的数学与物理抽象，涵盖连续与离散、状态与状态机、时钟时序以及串行与并行计算范式。

\subsubsection{核心主题}
\begin{itemize}
    \item \textbf{系统范式对比}：数字系统与模拟系统、离散系统与连续系统的本质差异与现实工业对应。
    \item \textbf{布尔逻辑与化简}：基本逻辑门、复合逻辑表达式求值、布尔代数恒等律与真值表化简。
    \item \textbf{时钟、状态与时序控制}：系统时钟信号、触发边沿、时钟周期节拍、状态保持与时序同步机制。
    \item \textbf{状态与状态转移}：离散系统状态空间定义、初始状态、中间态、终止状态以及状态全生命周期演化模型。
    \item \textbf{有限状态机理论}：确定性有限状态自动机（DFA）与非确定性有限状态自动机（NFA）的形式化定义、状态转移表与状态转移图。
    \item \textbf{逻辑电路与存储体系}：无记忆组合逻辑电路与含反馈记忆时序逻辑电路的辨析；锁存器、寄存器、静态与动态存储器单元。
    \item \textbf{作为状态机的计算机与协议}：通用计算机体系结构作为指令解释状态机；网络通信协议（如 TCP 三次握手）作为协同状态机。
    \item \textbf{串行与并行处理范式辨析}：顺序计算流与并发计算流的本质区别。
    \item \textbf{应用级串行处理}：单线程执行模型、程序计数器顺序推进与阻塞式 I/O。
    \item \textbf{简易应用级并行处理}：
    \begin{itemize}
        \item 请求级并行（Web 服务、C/S 架构客户端-服务器端交互、分布式集群服务）。
        \item 单服务器单线程轮询事件循环模型。
        \item 多服务器多线程并发处理架构。
        \item 任务级流水线并行及其重叠加速机制。
    \end{itemize}
\end{itemize}

\subsubsection{示范性学习成果与能力目标}
\begin{enumerate}
    \item \textbf{[解释]} 深刻辨析数字系统与模拟系统、离散系统与连续系统的根本区别，并能列举工程实际与日常生活中的典型实例进行说明。
    \item \textbf{[解释]} 将通用的计算过程形式化建模为具备已知配置集合（状态空间）并在不同配置间依照确定规则发生跃迁（状态转移）的动态演化系统。
    \item \textbf{[解释]} 严密剖析输出仅取决于当前输入的无状态系统与具有历史记忆依赖的有状态系统在架构设计、并发控制与横向扩展上的本质异同。
    \item \textbf{[开发]} 针对具体工程控制问题（如多相位交通信号灯时序控制器、通信数据流序列模式匹配器），建立形式化的有限状态机状态转移模型与规范。
    \item \textbf{[解释]} 将通用计算机体系结构深刻解释为一个逐条取指、译码并依据操作码更新内部寄存器与内存状态的确定性状态机。
    \item \textbf{[解释]} 阐明通用计算机软件程序或网络通信协议如何通过状态机进行精确无歧义建模，并论证同一计算逻辑可存在多种等价的状态表示方式。
    \item \textbf{[分析/评估]} 根据给定的状态机模型规范（如 TCP 连接生命周期管理状态机），严密推导并绘制系统在各种正常与异常输入场景下的时序演变行为曲线。
    \item \textbf{[开发]} 针对同一计算问题，分别编写其单线程串行版本和简单多线程/多进程并行版本的程序实现。
    \item \textbf{[评估]} 在不同输入规模负载下，评测并对比串行程序与并行程序的运行耗时，准确计算并量化其实际达成的并行加速比。
    \item \textbf{[解释/评估]} 在执行时间轴图上可视化展示并发操作与并行事件在物理时间上的同时发生，清晰阐述发掘并发潜力如何有效压缩任务的端到端运行耗时。
\end{enumerate}

\subsection{资源分配与调度}
\textbf{【学时建议】}：1 计算机核心学时，2 领域核心学时。

本单元深入剖析计算机有限物理与虚拟资源的分配机制，研究兼顾吞吐量、低延迟、公平性与防饥饿的经典调度算法及其工程权衡。

\subsubsection{核心主题}
\begin{itemize}
    \item \textbf{多维度计算机资源分类}：CPU 处理器时间配额、物理/虚拟内存空间、磁盘/固态硬盘存储 I/O 带宽、网络传输带宽及专用硬件加速器。
    \item \textbf{经典调度算法全景}：
    \begin{itemize}
        \item 先来先服务调度。
        \item 基于优先级的调度（静态优先级、动态优先级调整）。
        \item 最早截止时间优先调度（EDF，实时系统）。
        \item 公平份额调度与完全公平调度器（CFS）。
        \item 抢占式调度与非抢占式调度机制。
    \end{itemize}
    \item \textbf{常见调度算法的优缺点权衡}：周转时间、系统吞吐率、首响响应时间、调度决策开销、算法复杂度、资源利用率以及线程饥饿和优先级反转的防范。
\end{itemize}

\subsubsection{示范性学习成果与能力目标}
\begin{enumerate}
    \item \textbf{[解释]} 深刻阐释有限的计算机软硬件资源（CPU 份额、物理内存、持久化存储以及网络带宽）如何通过细粒度配额分配与准入控制有效供给多个竞争实体共享。
    \item \textbf{[解释]} 详细阐述主流调度算法（FCFS、轮转 RR、抢占式优先级调度、公平分享调度）的核心数据结构、工作原理与调度时序推进机制。
    \item \textbf{[评估]} 全面分析并评估常见调度算法在不同工作负载特征下的优劣势，深入权衡吞吐量、响应延迟、调度器开销与公平性之间的设计取舍。
    \item \textbf{[开发/评估]} 编程实现典型调度算法原型，搭建合成测试工作负载环境，对其调度延迟、吞吐指标与公平性表现进行严密的量化评测与对比。
\end{enumerate}

\subsection{系统性能}
\textbf{【学时建议】}：2 计算机核心学时，2 领域核心学时。\\
\textbf{【跨领域交叉引用】}：体系结构/Memory Management（内存管理）, 操作系统/Virtual Memory（虚拟内存）。

本单元从软硬件结合视角剖析物理时延鸿沟与局部性原理，揭示层次存储、多级缓存以及虚拟化与隔离机制对系统性能的关键赋能。

\subsubsection{核心主题}
\begin{itemize}
    \item \textbf{计算机系统中的物理与逻辑时延}：
    \begin{itemize}
        \item 物理光速极限与纳秒级计算尺度：真空光速约为每纳秒 30 厘米（1 英尺），相比之下 1~GHz 处理器的主频周期即为 1 纳秒。
        \item 数量级访存鸿沟对比：CPU 内部寄存器（$<1$~ns）、L1/L2/L3 高速缓存（$1\sim 10$~ns）、主存 DRAM 访问（$50\sim 100$~ns）、固态硬盘与机械磁盘 I/O（微秒至毫秒级）、跨局域网与广域网远程网络通信（毫秒至百毫秒级）。
    \end{itemize}
    \item \textbf{多级缓存机制与局部性原理}：时间局部性与空间局部性对处理器执行流水线及全栈系统吞吐的决定性影响。
    \item \textbf{全栈缓存与缓存一致性}：
    \begin{itemize}
        \item 缓存抽象的普遍性：数据库查询缓存、操作系统虚拟文件系统页缓存、分布式键值缓存以及多核体系结构硬件缓存。
        \item 缓存一致性挑战与维护机制（写穿透、写回、总线监听与目录式一致性协议）。
    \end{itemize}
    \item \textbf{分层存储系统与平均访存时间}：处理器分层存储体系结构引论；平均访存时间经典计算公式：$\text{AMAT} = \text{Hit Time} + \text{Miss Rate} \times \text{Miss Penalty}$。
    \item \textbf{虚拟化与隔离的底层设计动因}：多租户运行下的安全边界防护、单应用故障隔离以及可预测的服务质量保障。
    \item \textbf{间接层 设计哲学}：“计算机科学领域的任何问题都可以通过增加一个间接层来解决”；以虚拟内存对物理内存的解耦与重映射管理为例深入分析。
    \item \textbf{虚拟内存与虚拟机的实现机制}：页表、多级页表、转换后援缓冲器（TLB 快表）、缺页异常处理机制；虚拟机监控器、硬件辅助虚拟化技术（如 Intel VT-x）与容器化轻量级隔离技术。
\end{itemize}

\subsubsection{示范性学习成果与能力目标}
\begin{enumerate}
    \item \textbf{[解释]} 深刻阐释时间局部性与空间局部性原理在决定计算机底层硬件微架构、算法运行效率与系统软件性能中的基石作用。
    \item \textbf{[计算/评估]} 熟练运用平均访存时间（AMAT）公式进行严密量化推导，从存储容量、命中率/缺失率、访问时延与硅片开销四个维度深度评估存储层次结构的性能权衡。
    \item \textbf{[解释]} 深入阐明在多个并发程序与互不信任的执行环境共享底层物理计算资源的场景下，实施强隔离与安全保护的核心必要性。
    \item \textbf{[解释]} 详细阐述间接重映射抽象如何为并发运行的多道程序与隔离环境创造出拥有“专属独占物理机器与连续内存资源”的透明逻辑幻象。
    \item \textbf{[评估]} 针对运行在相互隔离的独立虚拟机或容器环境中的两个应用实例，实测并量化其在共用 CPU/内存/磁盘高压负载下的相互干扰度，科学评估系统性能隔离机制的有效性。
\end{enumerate}

\subsection{性能评估与评测}
\textbf{【学时建议】}：2 计算机核心学时，2 领域核心学时。

本单元建立现代计算机系统的科学量化评估方法学，涵盖核心指标体系构建、基准测试集设计、性能分析定律、数量级估算以及性能诊断工具链。

\subsubsection{核心主题}
\begin{itemize}
    \item \textbf{系统性能核心度量指标体系}：系统吞吐量、端到端响应时间、长尾延迟（Tail Latency, P99/P99.9 分位数）、资源饱和度与能效比（吞吐/瓦特）。
    \item \textbf{工作负载建模与代表性基准测试集}：标准合成负载建模、代表性工业基准测试套件（如 SPEC CPU、TPC-C、YCSB 等），以及性能测量数据的规范化采集、统计滤波与分布特征分析。
    \item \textbf{CPI (Cycles per Instruction) 性能方程}：$\text{CPU Time} = I \times \text{CPI} \times T_{\text{clk}}$；作为深刻理解指令集架构演进、处理器流水线深度与层次存储系统组织权衡的定量分析工具。
    \item \textbf{阿姆达尔定律 (Amdahl's Law)}：计算任务中无法被加速或无法并行的串行瓶颈部分，严格限制了增加算力资源所能获得的整体最大加速比上限：$S(s) = \frac{1}{(1-p) + \frac{p}{s}}$。
    \item \textbf{指导定量评估的解析建模工具}：排队论模型直觉（Little 定律、M/M/1 队列模型）、均值分析与系统瓶颈理论推演。
    \item \textbf{数量级分析与大 O 渐进记号}：Order of Magnitude 分析法、大 O 渐进时间/空间复杂度分析及其在工业级大规模系统资源预算规划中的应用。
    \item \textbf{快路径与慢路径分离}：系统架构中的“快路径”与“慢路径”设计哲学（让常见关键路径极度精炼高效，慢路径保障完备性）。
    \item \textbf{底层系统微架构事件对宏观性能的损耗}：流水线冒险停顿、多级缓存缺失、缺页异常、TLB 缺失刷新与分支预测错误惩罚。
    \item \textbf{分层系统与多租户异构平台评测挑战}：复杂分层调用栈、并发多租户工作负载与异构计算平台的非线性相互作用，及其给客观定量评测带来的工程挑战。
    \item \textbf{微基准测试 的常见陷阱}：编译器激进优化剔除关键测试代码、系统冷启动与缓存预热偏差、操作系统动态调频（DVFS）干扰与多核伪共享（False Sharing）。
\end{itemize}

\subsubsection{示范性学习成果与能力目标}
\begin{enumerate}
    \item \textbf{[解释]} 深入阐明计算机系统体系结构各层组件（硬件微架构、操作系统内核调度、语言运行时、网络栈）如何协同促进系统端到端性能的提升。
    \item \textbf{[解释]} 阐明在特定工业场景（如高频量化交易系统、大规模交互式 Web 搜索、大数据离线批处理计算）下选用特定性能指标（如平均时延 vs 尾延迟 P99、批处理总吞吐 vs 单任务执行时间）的判定依据与业务意义。
    \item \textbf{[解释/评估]} 辩证阐释基准测试集在客观评价系统性能方面的核心指导价值及其固有的代表性偏差与过度优化风险。
    \item \textbf{[解释/分析]} 准确阐述阿姆达尔定律的数学公式与物理内涵，深入讨论其在强扩展与大规模异构并行计算环境下的适用边界与工程局限。
    \item \textbf{[应用/评估]} 熟练运用极限分析法或粗略数量级心算法，在复杂系统方案规划阶段快速定量估算出目标系统性能指标的理论上下界。
    \item \textbf{[应用]} 熟练运用现代系统级性能剖析与诊断工具链（如 perf、eBPF、火焰图 FlameGraph、gprof 等），对目标系统的执行热点、调用链路与资源消耗进行精细化测量与画像。
    \item \textbf{[开发/评估]} 针对典型主流工业级计算系统（如 Linux 操作系统内核、分布式大数据分析引擎 Spark），自主设计并端到端执行一套规范严密的性能基准评测实验。
    \item \textbf{[评估]} 在复杂分层计算系统环境中开展受控对照基准实验，精确量化某一系统底层关键参数（如并发线程池规模、缓存淘汰策略、TCP 拥塞算法）调整对整体宏观性能产生的实际影响。
\end{enumerate}

\subsection{系统可靠性}
\textbf{【学时建议】}：2 计算机核心学时，1 领域核心学时。

本单元立足全生命周期的系统高可用与容错工程，辨析缺陷、差错与故障的级联传导，掌握基于校验、冗余编码与状态复制的容错机制。

\subsubsection{核心主题}
\begin{itemize}
    \item \textbf{软件缺陷与物理故障辨析}：程序逻辑缺陷、系统运行时差错与底层硬件物理故障（Faults，如坏内存位翻转、瞬态电磁干扰）的本质机理与级联演变路径。
    \item \textbf{基于冗余保障高可靠性机制之一：状态校验与重试}：校验和、超时与自适应重传、状态回滚与幂等操作设计。
    \item \textbf{基于冗余保障高可靠性机制之二：信息冗余编码}：检错与纠错编码（海明码 ECC、循环冗余校验 CRC、前向纠错码 FEC）。
    \item \textbf{基于冗余保障高可靠性机制之三：空间与硬件冗余}：双机热备、磁盘阵列镜像容错（RAID-1/RAID-5/6）、多副本状态机复制。
    \item \textbf{系统可靠性与高可用其他工程途径}：健康心跳探测、看门狗定时器、服务熔断与限流降级、自动化故障转移以及混沌工程故障主动注入测试。
\end{itemize}

\subsubsection{示范性学习成果与能力目标}
\begin{enumerate}
    \item \textbf{[解释]} 深入辨析软件程序编写缺陷、操作系统运行时差错、底层物理硬件故障以及运行时异常（Exceptions，如除零异常、越界中断）之间的层次传导关系与本质区别。
    \item \textbf{[解释]} 深刻阐释在系统工程中对故障进行“早感知检测”、“局部隔离容忍”以及“无缝状态自愈与恢复”的完整防御纵深与关键实现手段。
    \item \textbf{[解释]} 阐释纠错编码技术（ECC）在易失性动态内存（DRAM）、持久化存储介质以及网络物理信道通信中防范比特翻转衰减与噪声干扰的关键保护作用。
    \item \textbf{[应用]} 运用经典信息冗余算法（如海明编码、奇偶校验），设计并编程实现具备数据位差错检测与自修复能力的数据编解码软件模块。
    \item \textbf{[评估]} 从数据存储空间冗余开销、算法计算复杂度，以及编码、检错与纠错在不同系统层级上的时间延迟代价等维度，多角度系统对比不同纠错机制的综合优劣与工程选型策略。
\end{enumerate}

\subsection{系统安全}
\textbf{【学时建议】}：2 计算机核心学时，1 领域核心学时。

本单元研究计算机系统面临的软硬件安全威胁与攻击向量，阐述基于密码学底座、最小特权架构设计与全方位纵深防御的现代系统防护工程。

\subsubsection{核心主题}
\begin{itemize}
    \item \textbf{常见系统安全威胁与攻击向量}：计算机病毒、木马与蠕虫；分布式拒绝服务攻击（DDoS）；内存缓冲区溢出与特权提升漏洞；恶意窃听与中间人攻击（MITM）。
    \item \textbf{系统级安全防御对策体系}：
    \begin{itemize}
        \item 现代密码学基石：对称与非对称数据加密算法、数字证书与 PKI 信任链、不可逆单向哈希函数与数字签名技术。
        \item 安全体系结构设计：最小特权原则、可信计算基（TCB）、强制与自主访问控制（MAC/DAC）、硬件隔离安全区（如 Intel SGX、ARM TrustZone）。
        \item 动态防御与安全审计：入侵检测与防御系统（IDS/IPS）、下一代防火墙（NGFW）、安全审计与系统日志不可篡改追踪。
    \end{itemize}
\end{itemize}

\subsubsection{示范性学习成果与能力目标}
\begin{enumerate}
    \item \textbf{[解释]} 详细阐明计算系统面临的常见安全威胁及其攻击机理，并能结合业界真实网络空间安全攻击案例进行深入剖析与安全推演。
    \item \textbf{[解释/评估]} 全面阐述应对各类系统安全漏洞与威胁的主流系统级防护对策，深刻理解密码学加密、沙箱隔离、可信执行环境与访问控制列表在全栈安全加固中的协同机制与部署架构。
\end{enumerate}

\subsection{系统设计}
\textbf{【学时建议】}：2 计算机核心学时，1 领域核心学时。

本单元培养学生的宏观系统架构设计与工程权衡思维，确立高可用、高扩展、强健壮与安全的通用设计准则，并深度剖析业界经典工业级系统架构。

\subsubsection{核心主题}
\begin{itemize}
    \item \textbf{系统设计核心通用准则}：
    \begin{itemize}
        \item 活跃性：系统保证终将发生期望进展，不发生死锁、活锁或无限饥饿。
        \item 安全性：系统保证绝不发生未定义的非法破坏状态，数据不变性与完整性始终受控。
        \item 鲁棒性/健壮性：在面对畸形输入、部分节点瘫痪或硬件瞬态失效时保持优雅降级。
        \item 可扩展性：系统算力随资源横向堆叠或纵向升级呈现近线性扩容。
        \item 系统安全性：防御未授权访问、恶意篡改与隐私数据泄露。
    \end{itemize}
    \item \textbf{业界代表性工业级系统架构全景剖析}：
    \begin{itemize}
        \item Apache HTTP Server：多进程 Prefork 与多线程 Worker 高性能并发连接服务模型。
        \item Apache Spark：基于内存 RDD（弹性分布式数据集）的有向无环图（DAG）分布式任务调度与执行引擎。
        \item Linux 操作系统：基于虚拟文件系统（VFS）、统一设备抽象模型与模块化单内核的工业级稳健架构设计。
    \end{itemize}
\end{itemize}

\subsubsection{示范性学习成果与能力目标}
\begin{enumerate}
    \item \textbf{[解释]} 深入阐明现代计算系统设计中的通用评价准则（包括活跃性、安全性、鲁棒性、高可扩展性与系统安全性）的形式化内涵与工程度量。
    \item \textbf{[开发]} 给定特定应用系统的业务功能性需求与关键非功能性服务等级协议（SLO/SLA 约束），设计并输出高水准、组件接口清晰的顶层系统宏观架构设计方案。
    \item \textbf{[解释/评估]} 深入剖析并阐述业界代表性复杂软件系统（如 Linux 内核、Spark 分布式计算框架、高性能 Web 服务器）的底层体系架构、核心关键决策与其背后的工程权衡取舍。
\end{enumerate}

\section{专业素养与思维倾向}

\begin{itemize}
    \item \textbf{工程严谨性}：在学习、剖析、量化评估与调试大型计算系统时，必须始终对微观软硬件细节、跨层接口边界、时序竞态条件以及潜在的故障诱因保持高度的敏锐度与严密性，养成“差之毫厘，失之千里”的严谨工程品格。
    \item \textbf{动态适应性与权衡思维}：在架构设计与技术选型中展现出高度的灵活性与辩证思维。深刻认知世上不存在放之四海而皆准的“银弹”系统设计；面对迥异的业务需求、硬件约束与工作负载特征，能够审时度势，在吞吐量、响应时延、一致性、开发复杂度与成本之间做出最科学合理的工程设计决策。
\end{itemize}

\section{数学基础要求}

\subsection{必修数学基础}
\begin{itemize}
    \item \textbf{离散数学}：
    \begin{itemize}
        \item 集合论与二元关系（等价关系、偏序关系）。
        \item 基础图论（树结构、有向无环图 DAG、拓扑排序、图的连通性分析）。
        \item 数理逻辑基础（布尔代数运算、命题逻辑、真值表推导）。
    \end{itemize}
    \item \textbf{线性代数}：
    \begin{itemize}
        \item 基础矩阵与向量运算（矩阵乘法、状态转移矩阵、基变换与空间映射）。
    \end{itemize}
    \item \textbf{概率与数理统计}：
    \begin{itemize}
        \item 随机变量及其分布特性。
        \item 贝叶斯定理（Bayes' Theorem）与条件概率。
        \item 数学期望与方差/标准差。
        \item 累积分布函数（CDF）与概率密度函数（PDF），用于延迟分位数（P95/P99/P99.9）评测分析。
    \end{itemize}
\end{itemize}

\subsection{进阶建议数学基础}
\begin{itemize}
    \item \textbf{排队论基础}：M/M/1 排队模型、Little 定律、服务到达率与服务速率平衡，用于系统服务延迟与队列积压分析。
    \item \textbf{随机过程初步}：离散时间马尔可夫链（DTMC）与连续时间马尔可夫链（CTMC）状态转移平衡方程。
\end{itemize}

\section{跨领域共享概念与系统思维}

本知识领域作为多门骨干系统课程的公约数底座，与以下专业核心知识领域深度交织：
\begin{itemize}
    \item \textbf{网络与通信 (NC)}：协议分层模型、套接字 API 抽象、TCP/UDP 传输状态机、端到端可靠性与拥塞控制。
    \item \textbf{操作系统 (OS)}：进程与线程并发模型、虚拟内存分页管理、多级调度器设计、系统调用与特权级转换。
    \item \textbf{计算机体系结构与组织 (AR)}：微架构流水线、多级层次缓存、指令级并行、存储器一致性与软硬件接口边界。
    \item \textbf{并行与分布式计算 (PDC)}：多核并发控制、无锁同步、分布式状态机复制、强扩展与弱扩展评测。
\end{itemize}

\textbf{系统思维跨学科贯通核心维度：}
\begin{itemize}
    \item \textbf{软硬件协同设计}：打通算法抽象、编译运行时与底层微架构，破除单层优化瓶颈。
    \item \textbf{分层与抽象屏蔽}：通过严密的 API 与契约屏蔽复杂性，实现关注点分离。
    \item \textbf{状态与状态机模型}：统一利用有限状态机刻画时序控制流、通信交互协议与机器执行轨迹。
    \item \textbf{资源虚拟化与隔离}：以间接层提供透明独占抽象，通过沙箱机制实现多租户强安全与 QoS 隔离。
    \item \textbf{性能瓶颈与局部性}：洞悉光速物理极限与层级时延鸿沟，深度利用时空局部性与缓存体系攻克“访存墙”。
    \item \textbf{容错与冗余设计}：承认硬件物理与软件逻辑故障的必然性，基于校验、重试、信息编码与副本机制构筑高可用防御纵深。
    \item \textbf{全生命周期系统安全性}：遵循最小特权与深度防御原则，融合密码学机制构筑系统可信根基。
\end{itemize}

\section{真实能力任务规范}

\begin{taskbox}[任务 1：阐述系统设计的通用设计准则]
\textbf{【任务情境】}：针对企业拟建设的特定高并发业务系统（如大规模 C/S 客户端-服务器架构应用或分布式数据库集群），召开架构方案技术评审会。\\
\textbf{【能力陈述】}：面对给定的业务应用场景，系统性阐述并论证在设计此类计算系统时必须重点考量的核心通用准则（包括活跃性、安全性、鲁棒性、可扩展性与系统安全性），阐明各设计准则间的冲突与工程权衡取舍，形成严谨的设计规范书。\\
\textbf{【能力领域】}：系统 \quad \textbf{【能力单元】}：需求分析 / 架构设计 / 技术文档规范\\
\textbf{【关联知识单元】}：系统基础/系统设计、系统基础/系统可靠性、系统基础/系统安全\\
\textbf{【核心要求】}：计算机核心 \quad \textbf{【达标等级】}：\textbf{解释}
\end{taskbox}

\begin{taskbox}[任务 2：剖析与对比不同资源调度算法机制]
\textbf{【任务情境】}：操作系统内核或分布式容器调度集群需要为多租户多任务高并发场景合理分配关键计算资源（CPU、内存、网络带宽）。\\
\textbf{【能力陈述】}：给定主流调度算法族（如静态/动态优先级调度、最早截止时间优先 EDF、抢占式轮转 RR 与非抢占式调度等），深入剖析各算法的底层决策逻辑与调度时序，并从吞吐率、响应延迟、开销及抗饥饿能力等维度全面评估其各自的优缺点与适用场景。\\
\textbf{【能力领域】}：系统 / 计算理论 \quad \textbf{【能力单元】}：需求分析 / 技术文档 / 性能评估\\
\textbf{【关联知识单元】}：系统基础/资源分配与调度、系统基础/系统性能\\
\textbf{【核心要求】}：计算机核心 / 领域核心 \quad \textbf{【达标等级】}：\textbf{解释与评估}
\end{taskbox}

\begin{taskbox}[任务 3：面向特定应用场景选择并实现调度算法]
\textbf{【任务情境】}：针对具有特定延迟敏感性约束的特定业务环境（如实时音视频流转码分发系统、高并发 API 聚合网关或批处理数据清洗平台），默认通用调度策略无法满足特定 SLA/SLO 约束。\\
\textbf{【能力陈述】}：深入分析目标应用场景的负载特征与性能目标，精准选型或创新设计针对性的定制化调度算法，编程实现该调度引擎原型，并利用合成工作负载对其吞吐量、平均周转时间及尾延迟进行量化评测与调优。\\
\textbf{【能力领域】}：系统 / 应用 \quad \textbf{【能力单元】}：需求分析 / 架构设计 / 核心开发 / 质量测试 / 技术文档 / 综合评估\\
\textbf{【关联知识单元】}：系统基础/资源分配与调度、系统基础/系统性能、系统基础/系统性能评估、系统基础/系统设计\\
\textbf{【核心要求】}：领域核心 \quad \textbf{【达标等级】}：\textbf{应用、评估与开发}
\end{taskbox}

\begin{taskbox}[任务 4：端到端高可用系统架构设计与工程研发]
\textbf{【任务情境】}：从零研发一套满足生产环境严苛工业标准的高性能、高可用分布式服务系统（如轻量级分布式键值数据库或高并发异步 Web 服务器）。\\
\textbf{【能力陈述】}：面对给定的业务场景需求及多维度非功能性质量约束（高吞吐量、低延迟、99.99\% 高可用性容灾与多层级安全防护），主导自顶向下的高内聚低耦合系统架构设计，端到端编码实现核心功能模块，实施全套单元测试与集成压测，顺利完成系统交付验收。\\
\textbf{【能力领域】}：系统 / 应用 \quad \textbf{【能力单元】}：需求工程 / 架构设计 / 核心开发 / 系统测试 / 生产部署 / 模块集成 / 技术文档 / 性能评测 / 验收交付\\
\textbf{【关联知识单元】}：系统基础/系统基础概念、系统基础/系统设计、系统基础/系统可靠性、系统基础/系统安全、系统基础/系统性能评估\\
\textbf{【核心要求】}：领域核心 \quad \textbf{【达标等级】}：\textbf{应用、评估与开发}
\end{taskbox}

\begin{taskbox}[任务 5：评测指定系统的多维综合性能指标]
\textbf{【任务情境】}：针对已上线的工业级计算系统或开源大数据处理组件，需客观评定其在极限负载下的处理能力与抗压边界。\\
\textbf{【能力陈述】}：针对目标被测系统，综合运用理论性能上限解析推演与实证基准测试实验，对其核心性能指标（计算正确性、并发吞吐量 QPS、平均响应时延以及 P95/P99/P99.9 尾延迟等）开展严密的受控测量，并输出详实的量化评估分析报告。\\
\textbf{【能力领域】}：系统 / 计算理论 \quad \textbf{【能力单元】}：质量测试 / 部署调优 / 实验文档 / 性能评估\\
\textbf{【关联知识单元】}：系统基础/系统性能、系统基础/系统性能评估\\
\textbf{【核心要求】}：计算机核心 / 领域核心 \quad \textbf{【达标等级】}：\textbf{应用与评估}
\end{taskbox}

\begin{taskbox}[任务 6：精准诊断与定位复杂系统的性能瓶颈]
\textbf{【任务情境】}：处于生产环境中的大型分布式计算系统遭遇突发性性能恶化或系统抖动，扩容算力无法换来线性的吞吐增长。\\
\textbf{【能力陈述】}：针对在真实部署环境中运行的复杂系统，穿透全栈调用链路，熟练运用系统级性能剖析工具与理论分析推导，精准排查并断定限制系统吞吐与时延的核心瓶颈（如内存访存墙、CPU 算力耗尽、网络/磁盘 I/O 拥塞或锁争用开销）。\\
\textbf{【能力领域】}：系统 / 计算理论 \quad \textbf{【能力单元】}：架构设计 / 系统测试 / 环境部署 / 技术文档 / 性能评估\\
\textbf{【关联知识单元】}：系统基础/系统性能、系统基础/系统性能评估、系统基础/系统设计、系统基础/计算机系统概览\\
\textbf{【核心要求】}：领域核心 \quad \textbf{【达标等级】}：\textbf{应用、评估与开发}
\end{taskbox}

\begin{taskbox}[任务 7：基于深度工具链排查并修复分层系统软件缺陷]
\textbf{【任务情境】}：在具备复杂分层架构的软件系统中，偶发深层次跨层异常、并发竞态死锁或未知内存越界破坏，导致服务间歇性崩溃。\\
\textbf{【能力陈述】}：基于目标系统的故障异常日志与运行环境，熟练运用高级程序跟踪分析、单步调试、核心转储 Core Dump 与内存检测工具链，在庞大的调用栈中迅速定位深层次代码缺陷根源，并实施高质量的代码重构与回归验证测试。\\
\textbf{【能力领域】}：软件 / 系统 / 应用 \quad \textbf{【能力单元】}：核心开发 / 质量测试 / 文档记录 / 缺陷评估 / 运维保障\\
\textbf{【关联知识单元】}：系统基础/系统设计、系统基础/系统可靠性、系统基础/系统性能评估\\
\textbf{【核心要求】}：计算机核心 / 领域核心 \quad \textbf{【达标等级】}：\textbf{应用、评估与开发}
\end{taskbox}

\begin{taskbox}[任务 8：将本地系统平滑部署至现代云原生环境]
\textbf{【任务情境】}：研发团队需将原本在本地物理单机上开发测试完备的业务系统，向弹性公共云或私有云基础设施进行无缝云原生化迁移与多节点部署。\\
\textbf{【能力陈述】}：给定可在本地机器正常运行的系统程序，进行容器化封装镜像构建，编排配置容器网络拓扑、持久化数据卷与资源限额，顺利将其部署至基于 Kubernetes 或云平台容器编排引擎的集群环境中，达成高可用弹性运行。\\
\textbf{【能力领域】}：软件 / 系统 \quad \textbf{【能力单元】}：生产部署 / 模块集成 / 技术文档规范\\
\textbf{【关联知识单元】}：系统基础/计算机系统概览\\
\textbf{【核心要求】}：计算机核心 \quad \textbf{【达标等级】}：\textbf{应用}
\end{taskbox}

\begin{taskbox}[任务 9：大型系统的日常高可用运维与架构持续演进]
\textbf{【任务情境】}：面对工业级在线系统的生产常态化运行，业务流量急剧攀升，用户新功能需求与故障反馈纷至沓来，系统面临可用性与演进压力。\\
\textbf{【能力陈述】}：针对长期上线运转的大型系统，负责其全天候运维监控、告警响应、故障自愈与容灾灾备；并以用户反馈与生产监控数据为驱动，规划制定系统架构重构与平滑升级方案，推动复杂系统持续健壮演进。\\
\textbf{【能力领域】}：软件 / 系统 \quad \textbf{【能力单元】}：需求演进 / 架构重构 / 软件开发 / 质量测试 / 自动化部署 / 模块集成 / 文档演进 / 运行评估 / 生产运维 / 工程管理 / 用户验收 / 社会议题调适 / 持续改进\\
\textbf{【关联知识单元】}：系统基础/计算机系统概览、系统基础/系统基础概念、系统基础/系统设计、系统基础/系统可靠性、系统基础/系统安全、系统基础/系统性能评估\\
\textbf{【核心要求】}：领域核心 \quad \textbf{【达标等级】}：\textbf{应用、评估与开发}
\end{taskbox}

\section{课程设计与教学实施建议}

\subsection{导论型课程：计算机系统导论 (40 学时)}
本门课程面向大一第二学期或大二第一学期的计算机科学专业本科生，旨在建立系统思维的宏观基石：
\begin{itemize}
    \item \textbf{学时规划分配}：
    \begin{itemize}
        \item 计算机系统概述 --- 2 学时
        \item 系统基本概念 --- 6 学时
        \item 资源分配与调度 --- 4 学时
        \item 系统性能 --- 6 学时
        \item 性能评估 --- 6 学时
        \item 系统可靠性 --- 4 学时
        \item 系统设计 --- 6 学时
        \item 系统安全 --- 6 学时
    \end{itemize}
    \item \textbf{先修课程要求}：
    \begin{itemize}
        \item 离散数学中的集合与关系、基础图论与基本布尔逻辑。
        \item 线性代数中的基础矩阵与向量运算。
        \item 概率与数理统计中的随机变量、贝叶斯定理、数学期望与方差、累积分布函数（CDF）与概率密度函数（PDF）。
    \end{itemize}
    \item \textbf{达成能力陈述}：
    完成本课程学习的学生应能够：(1) 深刻理解计算机系统的全栈核心概念；(2) 掌握在性能、可靠性与安全性维度设计计算系统的关键工程准则；(3) 熟练根据官方技术文档，在真实环境中成功部署并科学评测代表性复杂工业系统（如 MySQL 数据库与 Apache Spark 计算框架）；(4) 具备独立设计并编程实现简易计算系统的工程能力（如交互式命令行程序、高并发轻量级 Web 服务器、简易持久化存储系统）。
\end{itemize}

\subsection{进阶高级课程：高级计算机系统架构与设计 (40--45 学时)}
本门课程面向大三或大四本科生，侧重大型分布式系统设计、微架构协同调优与前沿工程实战：
\begin{itemize}
    \item \textbf{学时规划分配}：
    \begin{itemize}
        \item 系统设计深度专题 --- 10 学时
        \item 操作系统交叉单元 (OS 核心单元深化) --- 2 学时
        \item 网络通信交叉单元 (NC 核心单元深化) --- 2 学时
        \item 并行分布式计算交叉单元 (PDC 核心单元深化) --- 2 学时
        \item 体系结构交叉单元 (AR 核心单元深化) --- 2 学时
        \item 系统可靠性与容错机制 --- 10 学时
        \item 系统性能优化与评测进阶 --- 6 学时
        \item 系统安全与纵深防御架构 --- 6 学时
    \end{itemize}
    \item \textbf{先修课程要求}：
    \begin{itemize}
        \item 排队论基础与随机过程初步。
        \item 系统基础（SF）导论型先修课程（或系统核心课程组合）。
    \end{itemize}
    \item \textbf{达成能力陈述}：
    完成本门高级进阶课程的学生应能够：(1) 对计算机系统设计的核心体系架构原则建立深刻的本质认知；(2) 将这些核心原则深度映射至经典系统（如 Linux 操作系统内核、SQL 数据库引擎与 TCP/IP 协议栈）以及现代前沿系统（如 Hadoop/Spark 分布式计算框架、大规模分布式存储与云原生基础设施）的工业级设计中；(3) 具备独立设计、研发落地与系统调优复杂工业级计算机系统的核心工程实战能力（如从零构建微型文件系统、高并发异步 I/O 网络服务器）。
\end{itemize}

\section{指导委员会与工作组成员}

\begin{itemize}
    \item \textbf{工作组主席}：
    \begin{itemize}
        \item 向乔，厦门大学
    \end{itemize}
    \item \textbf{工作组成员}：
    \begin{itemize}
        \item Doug Lea，纽约州立大学奥斯威戈分校, 美国纽约
        \item Monica D. Anderson，阿拉巴马大学, 美国阿拉巴马
        \item Matthias Hauswirth，卢加诺大学, 瑞士卢加诺
        \item 翟恩南，阿里巴巴集团, 中国杭州
        \item 刘宇彤，上海交通大学, 中国上海
    \end{itemize}
    \item \textbf{重要贡献者}：
    \begin{itemize}
        \item Michael S. Kirkpatrick，詹姆斯麦迪逊大学, 美国弗吉尼亚
        \item 孔令和，上海交通大学, 中国上海
    \end{itemize}
\end{itemize}

\begin{strip}
\section{附录：核心主题与技能要求等级对照表}

\noindent\begin{tabularx}{\textwidth}{p{0.18\textwidth} Y c c c}
\toprule
\textbf{知识单元} & \textbf{核心知识点与具体主题} & \textbf{认知等级} & \textbf{分类} & \textbf{学时} \\
\midrule
\multirow{3}{*}{\makecell[l]{系统基础/计算机系统\\概述}} 
& 计算机基本组件（门、触发器、寄存器、总线；数据通路+控制+存储器）与体系结构历史演进 & 解释 & 计算机核心 & \multirow{3}{*}{3} \\
& 作为计算范式的硬件逻辑、布尔表达式、逻辑化简与与或式（SOP）；编程抽象、接口与库调用 & 解释 & 计算机核心 & \\
& 应用程序与操作系统服务边界、RPC；系统调用与特权级交互；流水线原理与强/弱扩展性 & 解释 & 计算机核心 & \\
\midrule
\multirow{3}{*}{\makecell[l]{系统基础/系统基本\\概念}} 
& 数字 vs 模拟系统、离散 vs 连续系统对比；基本逻辑门、时钟节拍与布尔代数化简 & 应用 & 计算机核心 & \multirow{3}{*}{4} \\
& 状态与状态转移生命周期；有限状态机模型（DFA/NFA）；组合逻辑与时序逻辑 & 应用 & 计算机核心 & \\
& 计算机与网络协议作为状态机抽象；串行单线程 vs 应用级并行处理（请求级、多线程多服务器、流水线） & 应用 & 计算机核心 & \\
\midrule
\multirow{2}{*}{\makecell[l]{系统基础/资源分配\\与调度}} 
& 异构计算机资源分类（CPU 份额、内存空间、磁盘存储 I/O、网络传输带宽） & 解释 & 计算机核心 & 1 \\
& 经典调度算法机制（FCFS、优先级调度、公平共享、抢占式/非抢占式、EDF）及其优缺点综合权衡 & 解释 & 领域核心 & 2 \\
\midrule
\multirow{3}{*}{\makecell[l]{系统基础/系统性能}} 
& 物理时延层次（光速极限与纳秒时钟、内存 vs 磁盘 vs 网络传输时延鸿沟） & 应用 & 计算机核心 & \multirow{3}{*}{\makecell[c]{2 \\ (CS:2 \\ KA:2)}} \\
& 缓存机制与时间/空间局部性；跨层缓存与缓存一致性（数据库、OS、分布式系统、微架构） & 应用 & 计算机核心 & \\
& 处理器层次存储体系与 AMAT 公式；虚拟化与隔离动因；间接层哲学与虚拟内存/虚拟机机制 & 应用 & 领域核心 & \\
\midrule
\multirow{3}{*}{\makecell[l]{系统基础/性能评估}} 
& 性能指标体系（吞吐量、延迟、P99 尾延迟、资源饱和度）；工作负载建模与基准测试集 & 评估 & 计算机核心 & \multirow{3}{*}{\makecell[c]{2 \\ (CS:2 \\ KA:2)}} \\
& CPI 性能方程；阿姆达尔定律数学推导与并行极限；数量级分析（大 O 记号）与快慢路径分离 & 评估 & 计算机核心 & \\
& 底层微架构事件损耗（流水线停顿、缓存缺失、缺页异常）；分层平台评测挑战与微基准测试陷阱 & 评估 & 领域核心 & \\
\midrule
\multirow{2}{*}{\makecell[l]{系统基础/系统可靠性}} 
& 软件缺陷、系统差错与物理故障的本质区别与级联传导；检错与重试机制 & 评估 & 计算机核心 & 2 \\
& 基于冗余编码（ECC、CRC、FEC）与副本镜像复制的高可用机制；心跳、看门狗、熔断与故障转移 & 评估 & 领域核心 & 1 \\
\midrule
\multirow{2}{*}{\makecell[l]{系统基础/系统安全}} 
& 常见系统安全威胁全景（病毒、蠕虫、DDoS 攻击、恶意窃听、特权提升漏洞） & 评估 & 计算机核心 & 2 \\
& 系统级安全防御对策（现代密码学算法、安全架构与最小特权原则、IDS/IPS 与防火墙防护） & 评估 & 领域核心 & 1 \\
\midrule
\multirow{2}{*}{\makecell[l]{系统基础/系统设计}} 
& 系统设计核心通用衡量准则（活跃性、安全性、鲁棒性、高可扩展性与系统安全性） & 开发 & 计算机核心 & 2 \\
& 业界工业级代表性系统架构全景深度剖析（Apache HTTP Server、Apache Spark、Linux 内核） & 开发 & 领域核心 & 1 \\
\bottomrule
\end{tabularx}
\end{strip}

\chapter{课程体系规划与全球教学实践}

\section{三大能力领域的课程打包整合}

\begin{preamblebox}[从孤立课程走向能力集聚]
在 CS2023 指导纲要中，知识领域（KA）是内容切片的逻辑集合，而非孤立的高校课程。为了避免传统教学中“各自为政、知识脱节”的弊病，工作组围绕产业界与研究界的三大核心能力支柱，提出了跨领域的课程打包与贯穿整合范式：
\end{preamblebox}

\subsection{软件核心能力领域}
涵盖软件构建与理论分析的全生命周期，深度整合以下四大知识领域：
\begin{itemize}
    \item \textbf{软件开发基础 (SDF)}：编程思维、算法设计初步、防御性编码与调试。
    \item \textbf{算法与复杂性 (AL)}：高阶抽象数据类型、图算法、NPC 理论与算法伦理。
    \item \textbf{编程语言理论基础 (FPL)}：类型系统、形式语义、函数式/面向对象范式与运行时垃圾回收。
    \item \textbf{软件工程 (SE)}：大型代码仓演进、CI/CD 持续集成、重构、敏捷迭代与架构质量。
\end{itemize}
\textbf{【典型整合课程】}：《高级软件构造与算法设计》、《面向对象系统分析与形式化验证》、《企业级软件架构演化实训》。

\subsection{系统核心能力领域}
聚焦软硬件交界面及大规模分布式基础设施，深度整合以下七大知识领域：
\begin{itemize}
    \item \textbf{系统基础 (SF)}：系统思维、软硬件协同、状态机与资源抽象。
    \item \textbf{体系结构与组织 (AR)}：微架构、流水线、ILP、缓存一致性与 DSA/PIM 加速器。
    \item \textbf{操作系统 (OS)}：进程线程并发、虚拟内存、驱动架构与安全隔离。
    \item \textbf{并行与分布式计算 (PDC)}：OpenMP/MPI 并行、GPU 异构计算、分布式共识。
    \item \textbf{网络与通信 (NC)}：协议栈、TCP/QUIC 拥塞控制、BGP 路由与数据中心 Spine-Leaf 网络。
    \item \textbf{计算机安全 (SEC)}：内存防护、权限提升防御、威胁建模与系统韧性。
    \item \textbf{数据管理 (DM)}：底层存储引擎、WAL 预写日志、事务 ACID 与 LSM-Tree 存储结构。
\end{itemize}
\textbf{【典型整合课程】}：《现代计算机系统软硬件全栈设计》、《高并发分布式云原生系统》、《操作系统与网络协议栈内核实训》。

\subsection{应用核心能力领域}
面向真实应用场景，深度整合以下六大知识领域：
\begin{itemize}
    \item \textbf{人工智能 (AI)}：机器学习基础、深度神经网络、大语言模型（LLM）工程与强化学习。
    \item \textbf{图形学与交互技术 (GIT)}：现代实时渲染管线、物理仿真、XR/VR/AR 空间计算。
    \item \textbf{人机交互 (HCI)}：以人为本的设计（HCD）、可用性工程度量与多模态交互体验。
    \item \textbf{专用平台开发 (SPD)}：移动生态、物联网微控制器与边缘计算。
    \item \textbf{计算机安全 (SEC)} 与 \textbf{数据管理 (DM)}：应用层隐私注入、SQL 审计与多源数据管道。
\end{itemize}
\textbf{【典型整合课程】}：《智能具身感知与人机协同系统》、《空间计算与三维数字孪生引擎》、《数据驱动智能应用工程》。

\section{高校专业培养方案制定建议}

针对不同高校的定位与办学特长，CS2023 提出因校制宜的培养方案规划路径：
\begin{enumerate}
    \item \textbf{软件与工程主导型培养方案}：夯实编程底座（SDF）与算法深度（AL），在大三、大四重点强化软件工程工业级流水线，注重代码评审、自动化测试覆盖率与高内聚低耦合系统重构。
    \item \textbf{系统与底层架构主导型培养方案}：打通“数字逻辑 $\to$ 体系结构 $\to$ 操作系统 $\to$ 网络 $\to$ 分布式系统”全链路，强调动手编写简易 CPU、实现微内核与自研网络协议栈，培养具备底层破局能力的高精尖系统架构人才。
    \item \textbf{智能与交叉应用主导型培养方案}：紧扣现代 AI、大模型与沉浸式人机交互，前置连续数学基础（线性代数与概率优化），让学生在具体垂直行业（如智慧医疗、自动驾驶、金融科技）中开展实战课题。
\end{enumerate}

\section{全域高阶能力任务规范}

\begin{taskbox}[跨领域重大任务 1：高可用分布式智能交易结算平台架构设计]
\textbf{【跨领域融合】}：Systems + Software + Applications（涵盖 PDC, OS, NC, SEC, DM, SE）\\
\textbf{【能力陈述】}：针对千万级 TPS 实时交易结算场景，综合运用分库分表、分布式事务（Saga/两阶段提交）、Raft 共识日志强同步、零信任网络加密及全链路高可用监控，设计系统整体架构，制定故障自动容灾转移演练方案。\\
\textbf{【达标等级】}：\textbf{开发与评估}
\end{taskbox}

\begin{taskbox}[跨领域重大任务 2：面向低功耗边缘端的大模型轻量化推理引擎]
\textbf{【跨领域融合】}：Systems + Applications（涵盖 AR, AI, SPD, SF）\\
\textbf{【能力陈述】}：在算力与电量受限的嵌入式/移动硬件平台上部署大语言模型，通过权重量化（INT4/FP8）、算子融合、剪枝及硬件特定 NPU 汇编级指令调优，实现毫秒级首字延迟与超低功耗运行。\\
\textbf{【达标等级】}：\textbf{开发}
\end{taskbox}

\section{全球计算机教育前沿实践与时代思考}

\subsection{社会价值与伦理融合实践}
\begin{itemize}
    \item \textbf{无障碍辅助技术教学}：将 WCAG 标准、无障碍屏幕阅读器、色障友好交互作为前端与系统设计的核心质量指标。
    \item \textbf{科技向善}：引导学生参与非营利组织开源项目，将算法算力用于环境保护、贫困消减与公共卫生。
    \item \textbf{全球伦理视野与南北半球互通}：不仅探讨欧美发达经济体的技术法律框架，更深入审视发展中国家在数据主权、数字鸿沟与算法剥削上的严峻现实。
\end{itemize}

\subsection{全球化视野下的计算机教育模式}
\begin{itemize}
    \item \textbf{中国计算机教育实践}：探讨中国“101 计划”、新工科建设、拔尖创新人才培养基地以及产教融合协同育人机制，为全球超大规模本科计算机教学质量提升提供中国方案。
    \item \textbf{通识博雅院校 与社区学院模式}：小班研讨、个性化项目驱动与低门槛衔接通道经验互鉴。
    \item \textbf{拉美、阿拉伯世界与非洲的计算机教育崛起}：依托开源课程、低成本硬件与远程数字化赋能跨越式发展。
\end{itemize}

\subsection{专业建设与面向未来的思考}
\begin{itemize}
    \item \textbf{CS + X 跨学科复合专业}：计算生物学、数字法学、计算金融等新兴交叉学位建设范式。
    \item \textbf{智能时代的教学资源与教材变革}：从静态纸质教材向可交互 Jupyter/Colab 在线工坊、生成式 AI 交互助教演进。
\end{itemize}


\end{document}
